\paragraph{Stirling--Bernoulli 恢复、半维零块、自然扩张账本与共振整函数计数}
折叠规范常数的 Bernoulli 层级、谱零点集合的约数稀疏上界、自然扩张的 prime-register 外置实现以及共振整函数的 Cartwright 结构，在现有命题链下还可继续压缩为下列若干条直接后果。

\begin{theorem}[折叠规范常数列逐阶内生恢复全部 Bernoulli 数]\label{thm:xi-foldbin-gauge-constant-recovers-all-bernoulli}
沿用推论 \ref{cor:fold-bin-recover-2pi} 的规范常数列 \((C_m)_{m\ge 1}\)。对每个整数 \(r\ge 1\)，记
\[
\lambda_r:=\Bigl(\frac{\varphi}{2}\Bigr)^{2r-1},
\qquad
\gamma_r:=\frac{\varphi^{-1}+\varphi^{2r-3}}{r(2r-1)}.
\]
则对每个固定整数 \(R\ge 1\)，极限
\[
\boxed{\;
B_{2R}
=
\gamma_R^{-1}
\lim_{m\to\infty}
\lambda_R^{-m}
\left(
\log C_m-\log(2\pi)-\sum_{r=1}^{R-1}\gamma_r B_{2r}\lambda_r^m
\right)\;}
\]
存在；其中当 \(R=1\) 时，和式按空和解释。
因此序列 \((\log C_m)_{m\ge 1}\) 通过逐阶消元唯一决定全部偶 Bernoulli 数。
\end{theorem}
\begin{proof}
由定理 \ref{thm:fold-bin-gauge-constant-stirling-bernoulli-hierarchy}，对任意固定 \(R\ge 1\) 有渐近展开
\[
\log C_m
=
\log(2\pi)+
\sum_{r=1}^{R}\gamma_r B_{2r}\lambda_r^m
+O(\lambda_{R+1}^m).
\]
又因为
\[
\frac{\lambda_{R+1}}{\lambda_R}
=
\Bigl(\frac{\varphi}{2}\Bigr)^2<1,
\qquad
\gamma_R>0,
\]
故在已知 \(B_2,\dots,B_{2R-2}\) 时，减去前 \(R-1\) 阶项并乘以 \(\lambda_R^{-m}\)，得到
\[
\lambda_R^{-m}
\left(
\log C_m-\log(2\pi)-\sum_{r=1}^{R-1}\gamma_r B_{2r}\lambda_r^m
\right)
=
\gamma_R B_{2R}+O\!\left(\frac{\lambda_{R+1}^m}{\lambda_R^m}\right).
\]
令 \(m\to\infty\) 即得极限存在且等于 \(\gamma_R B_{2R}\)，整理便得到所示公式。证毕。
\end{proof}

\begin{corollary}[单一规范常数列完整编码 \texorpdfstring{\(\pi\)}{pi} 与全部偶值 \texorpdfstring{\(\zeta(2R)\)}{zeta(2R)}]\label{cor:xi-foldbin-gauge-constant-recovers-pi-even-zeta}
沿用定理 \ref{thm:xi-foldbin-gauge-constant-recovers-all-bernoulli} 的记号，则
\[
\boxed{\;
\pi=\frac12\exp\!\Bigl(\lim_{m\to\infty}\log C_m\Bigr).\;}
\]
并且对每个整数 \(R\ge 1\)，有
\[
\boxed{\;
\zeta(2R)
=
(-1)^{R-1}\frac{(2\pi)^{2R}}{2(2R)!}\,B_{2R},\;}
\]
其中 \(B_{2R}\) 由定理 \ref{thm:xi-foldbin-gauge-constant-recovers-all-bernoulli} 的极限公式唯一恢复。
因此，规范常数列 \((C_m)\) 单独即编码了 \(\pi\) 与全部偶值 \(\zeta\)-谱。
\end{corollary}
\begin{proof}
由推论 \ref{cor:fold-bin-recover-2pi}，
\(
\log C_m\to\log(2\pi)
\)，
故第一式立即成立。
再由定理 \ref{thm:xi-foldbin-gauge-constant-recovers-all-bernoulli}，全部 \(B_{2R}\) 可由 \((C_m)\) 逐阶恢复；而定理 \ref{thm:fold-bin-gauge-constant-stirling-bernoulli-hierarchy} 已给出标准恒等式
\[
B_{2R}=(-1)^{R-1}\frac{2(2R)!}{(2\pi)^{2R}}\zeta(2R).
\]
移项即得第二式。证毕。
\end{proof}

\begin{theorem}[fold-gauge 常数对 \texorpdfstring{$2\pi$}{2pi} 的一步 Richardson 型加速]\label{thm:xi-foldbin-gauge-constant-one-step-acceleration-2pi}
沿用推论 \ref{cor:fold-bin-recover-2pi} 的规范常数列 \((C_m)_{m\ge 1}\)，并记
\[
q:=\frac{\varphi}{2}\in(0,1),
\qquad
C_m^{\mathrm{acc}}
:=
C_m\exp\!\left(-\frac{1}{3\varphi}q^m\right).
\]
则有
\[
\boxed{\;
\log C_m^{\mathrm{acc}}-\log(2\pi)
=
-\frac{\sqrt5}{180}\,q^{3m}+O(q^{5m}).\;}
\]
从而
\[
\boxed{\;
C_m^{\mathrm{acc}}
=
2\pi\left(1-\frac{\sqrt5}{180}\,q^{3m}+O(q^{5m})\right).\;}
\]
特别地，去除首个显式 Bernoulli 模态后，\(2\pi\) 逼近误差阶立刻由 \(q^m\) 提升到 \(q^{3m}\)。
\end{theorem}
\begin{proof}
由定理 \ref{thm:xi-foldbin-gauge-constant-not-cfinite-first-two-zeta-limits}，
\[
\log C_m-\log(2\pi)
=
\frac{1}{3\varphi}\,q^m-\frac{\sqrt5}{180}\,q^{3m}+O(q^{5m}).
\]
故按定义直接得到
\[
\log C_m^{\mathrm{acc}}-\log(2\pi)
=
-\frac{\sqrt5}{180}\,q^{3m}+O(q^{5m}).
\]
令
\[
x_m:=-\frac{\sqrt5}{180}\,q^{3m}+O(q^{5m})=O(q^{3m}).
\]
则
\[
C_m^{\mathrm{acc}}
=
2\pi e^{x_m}
=
2\pi\bigl(1+x_m+O(x_m^2)\bigr).
\]
由于 \(x_m^2=O(q^{6m})=O(q^{5m})\)，遂得
\[
C_m^{\mathrm{acc}}
=
2\pi\left(1-\frac{\sqrt5}{180}\,q^{3m}+O(q^{5m})\right).
\]
证毕。
\end{proof}

\begin{theorem}[fold-gauge 常数的任意阶 Bernoulli 消模器与 \texorpdfstring{$(2R+1)$}{(2R+1)} 级加速]\label{thm:xi-foldbin-gauge-constant-arbitrary-order-bernoulli-annihilator}
\leanverified{paper\_xi\_foldbin\_gauge\_constant\_arbitrary\_order\_bernoulli\_annihilator}
沿用定理 \ref{thm:xi-foldbin-gauge-constant-recovers-all-bernoulli} 的记号。对任意固定整数 \(R\ge 1\)，定义唯一的次数至多为 \(R\) 的插值多项式
\[
\boxed{\;
p_R(z):=\prod_{r=1}^{R}\frac{z-\lambda_r}{1-\lambda_r}
=\sum_{j=0}^{R}\beta_j^{(R)}z^j.\;}
\]
则
\[
p_R(1)=1,
\qquad
p_R(\lambda_r)=0\quad(1\le r\le R).
\]
再定义加速组合
\[
\mathcal A_m^{(R)}:=\sum_{j=0}^{R}\beta_j^{(R)}\log C_{m+j}.
\]
则有
\[
\boxed{\;
\mathcal A_m^{(R)}
=
\log(2\pi)+O(\lambda_{R+1}^m)
=
\log(2\pi)+O\!\left(\left(\frac{\varphi}{2}\right)^{(2R+1)m}\right).\;}
\]
更精确地，
\[
\boxed{\;
\mathcal A_m^{(R)}
=
\log(2\pi)
+\gamma_{R+1}B_{2R+2}\,p_R(\lambda_{R+1})\,\lambda_{R+1}^m
+O(\lambda_{R+2}^m).\;}
\]
因此，对任意预先指定的阶数 \(R\)，上述线性组合都能同时消去前 \(R\) 个 Bernoulli 模态，并把逼近 \(\log(2\pi)\) 的主误差阶精确提升到 \((\varphi/2)^{(2R+1)m}\)。
\end{theorem}
\begin{proof}
多项式 \(p_R\) 的定义立即给出
\[
p_R(1)=1,
\qquad
p_R(\lambda_r)=0\quad(1\le r\le R),
\]
而唯一性来自 \(R+1\) 个插值条件对次数至多 \(R\) 的多项式的唯一确定。

由定理 \ref{thm:fold-bin-gauge-constant-stirling-bernoulli-hierarchy}，对固定 \(R\) 有渐近展开
\[
\log C_{m+j}
=
\log(2\pi)+\sum_{r=1}^{R+1}\gamma_r B_{2r}\lambda_r^{m+j}+O(\lambda_{R+2}^{m+j})
\]
对全部 \(0\le j\le R\) 同时成立。由于 \(j\) 仅在有限集合 \(\{0,\dots,R\}\) 中变化，余项可统一写成 \(O(\lambda_{R+2}^m)\)。
将该展开代入 \(\mathcal A_m^{(R)}\)，得到
\[
\mathcal A_m^{(R)}
=
\log(2\pi)\sum_{j=0}^{R}\beta_j^{(R)}
+\sum_{r=1}^{R+1}\gamma_r B_{2r}\lambda_r^m
\sum_{j=0}^{R}\beta_j^{(R)}\lambda_r^j
+O(\lambda_{R+2}^m).
\]
而
\[
\sum_{j=0}^{R}\beta_j^{(R)}=p_R(1)=1,
\qquad
\sum_{j=0}^{R}\beta_j^{(R)}\lambda_r^j=p_R(\lambda_r).
\]
故前 \(R\) 个模态全部被精确消去，仅剩
\[
\mathcal A_m^{(R)}
=
\log(2\pi)+\gamma_{R+1}B_{2R+2}p_R(\lambda_{R+1})\lambda_{R+1}^m+O(\lambda_{R+2}^m),
\]
即得更精确的盒装式。由于
\[
\lambda_{R+1}=\left(\frac{\varphi}{2}\right)^{2R+1},
\]
主误差阶便是 \((\varphi/2)^{(2R+1)m}\)。证毕。
\end{proof}

\begin{theorem}[谱零点集合的半维指数律]\label{thm:xi-fold-zero-set-half-dimensional-exponent-law}
\leanverified{paper\_xi\_fold\_zero\_set\_half\_dimensional\_exponent\_law}
沿用定理 \ref{thm:fold-zero-density-sparse} 与推论 \ref{cor:fold-zero-half-index-multiple6} 的记号，记
\[
\mathcal Z_m:=\{k\in \ZZ/(F_{m+2}\ZZ):\ \widehat d_m(k)=0\}.
\]
则沿同余子序列
\[
m+2\equiv 0\pmod 6
\]
有极限
\[
\boxed{\;
\lim_{\substack{m\to\infty\\ m+2\equiv 0\ (6)}}
\frac{1}{m}\log|\mathcal Z_m|
=
\frac12\log\varphi,\;}
\]
并且相对密度满足
\[
\boxed{\;
\lim_{\substack{m\to\infty\\ m+2\equiv 0\ (6)}}
\frac{1}{m}\log\frac{|\mathcal Z_m|}{F_{m+2}}
=
-\frac12\log\varphi.\;}
\]
换言之，\(|\mathcal Z_m|=\varphi^{m/2+o(m)}\)，而 \(|\mathcal Z_m|/F_{m+2}=\varphi^{-m/2+o(m)}\)。
\end{theorem}
\begin{proof}
若 \(m+2\equiv 0\pmod 6\)，记
\[
d:=\frac{m+2}{2}.
\]
由推论 \ref{cor:fold-zero-half-index-multiple6}，
\[
|\mathcal Z_m|\ge F_d.
\]
另一方面，由定理 \ref{thm:fold-zero-density-sparse}，
\[
\frac{|\mathcal Z_m|}{F_{m+2}}
\le
\tau(m+2)\frac{F_{\lfloor (m+2)/2\rfloor}}{F_{m+2}}
=
\tau(m+2)\frac{F_d}{F_{m+2}},
\]
故
\[
|\mathcal Z_m|\le \tau(m+2)\,F_d.
\]
由于约数函数满足 \(\tau(n)=e^{o(n)}\)，而 Fibonacci 数满足 \(F_n=\varphi^{n+o(n)}\)，遂得
\[
F_d\le |\mathcal Z_m|\le e^{o(m)}F_d=\varphi^{m/2+o(m)}.
\]
两边取 \((1/m)\log\) 即得第一条极限。
再由
\[
F_{m+2}=\varphi^{m+o(m)},
\]
从而
\[
\frac{|\mathcal Z_m|}{F_{m+2}}
=
\varphi^{-m/2+o(m)},
\]
第二条极限随即成立。证毕。
\end{proof}
\noindent 该定理表明：六倍窗上由半阶余类块
\(
S_{F_{(m+2)/2}}(F_{m+2})
\)
触发的大零点簇并非仅提供一个粗下界现象，而是把零点集合的绝对规模沿
\(
m+2\equiv 0\pmod 6
\)
这一子序列精确锁定在
\(
\exp((\log\varphi/2)m)
\)
的指数尺度上；相应相对密度则以同一半维速率衰减。

\begin{theorem}[约数驱动零余类并集的精确二元包含--排斥公式]\label{thm:xi-fold-zero-divisor-two-class-inclusion-exclusion}
\leanverified{paper_xi_fold_zero_divisor_two_class_inclusion_exclusion}
沿用推论 \ref{cor:fold-zero-divisor-lb} 与引理 \ref{lem:fold-zero-coset-intersection} 的记号，记
\[
M:=F_{m+2}.
\]
设 \(d,e\) 为 \(m+2\) 的真约数，并满足
\[
F_d\mid \frac{M}{2},
\qquad
F_e\mid \frac{M}{2}.
\]
则
\[
S_{F_d}(M)\subseteq \mathcal Z_m,
\qquad
S_{F_e}(M)\subseteq \mathcal Z_m,
\]
并且
\[
\left|S_{F_d}(M)\cup S_{F_e}(M)\right|
=
F_d+F_e-\varepsilon(d,e)F_{\gcd(d,e)},
\]
其中
\[
\varepsilon(d,e):=
\begin{cases}
1,& 2F_{\gcd(d,e)}\mid(F_e-F_d),\\
0,& \text{否则}.
\end{cases}
\]
特别地，
\[
|\mathcal Z_m|
\ge
F_d+F_e-\varepsilon(d,e)F_{\gcd(d,e)}.
\]
\end{theorem}
\begin{proof}
由推论 \ref{cor:fold-zero-divisor-lb}，
\[
S_{F_d}(M)\subseteq \mathcal Z_m,
\qquad
S_{F_e}(M)\subseteq \mathcal Z_m,
\qquad
|S_{F_d}(M)|=F_d,
\qquad
|S_{F_e}(M)|=F_e.
\]
另一方面，引理 \ref{lem:fold-zero-coset-intersection} 给出
\[
S_g(M)\cap S_h(M)\neq\varnothing
\iff
2\gcd(g,h)\mid(h-g),
\]
且一旦交集非空，就有
\[
|S_g(M)\cap S_h(M)|=\gcd(g,h).
\]
取 \(g=F_d\)、\(h=F_e\)，并使用 Fibonacci 最大公因子恒等式
\[
\gcd(F_d,F_e)=F_{\gcd(d,e)},
\]
便得到
\[
\left|S_{F_d}(M)\cap S_{F_e}(M)\right|
=
\varepsilon(d,e)F_{\gcd(d,e)}.
\]
于是由包含--排斥公式可得
\[
\left|S_{F_d}(M)\cup S_{F_e}(M)\right|
=
F_d+F_e-\varepsilon(d,e)F_{\gcd(d,e)}.
\]
再由该并集包含于 \(\mathcal Z_m\)，立得最后一个下界。证毕。
\end{proof}

上述二元包含--排斥公式还可进一步压缩为一条由 \(\gcd\) 与 \(2\)-进层级共同锁定的交结构定理，并由此得到整个零点集合的分层交半格描述。

\begin{theorem}[零块交集的 \texorpdfstring{$\gcd$}{gcd} 刚性与 \texorpdfstring{$2$}{2}-进分层判别]\label{thm:xi-fold-zero-coset-v2-intersection-rigidity}
\leanverified{paper\_xi\_fold\_zero\_coset\_v2\_intersection\_rigidity}
沿用定理 \ref{thm:fold-zero-coset-rigidity} 的记号。设 \(M\) 为偶整数，且
\[
g,h\mid \frac{M}{2}.
\]
则有精确交公式
\[
\boxed{\;
S_g(M)\cap S_h(M)=
\begin{cases}
S_{\gcd(g,h)}(M),& v_2(g)=v_2(h),\\[2mm]
\varnothing,& v_2(g)\neq v_2(h).
\end{cases}\;}
\]
\end{theorem}
\begin{proof}
由定理 \ref{thm:fold-zero-coset-rigidity}，
\[
S_g(M)=\left\{k\in \ZZ/(M\ZZ):\ k\equiv \frac{M}{2g}\pmod{M/g}\right\},
\]
等价地，
\[
S_g(M)=\left\{\frac{M}{2g}(1+2t)\ \ (\bmod M):\ t=0,1,\dots,g-1\right\}.
\]
由于这些代表元都落在区间 \([0,M)\) 内，任意 \(k\in S_g(M)\) 都满足
\[
v_2(k)=v_2(M)-v_2(g)-1.
\]
因此若 \(v_2(g)\neq v_2(h)\)，则 \(S_g(M)\) 与 \(S_h(M)\) 中元素的 \(2\)-进赋值恒不相同，故二者不可能相交。

现设 \(v_2(g)=v_2(h)\)。令
\[
d:=\gcd(g,h),\qquad a:=\frac{M}{g},\qquad b:=\frac{M}{h},\qquad L:=\operatorname{lcm}(a,b)=\frac{M}{d}.
\]
若再写
\[
g=d\alpha,\qquad h=d\beta,\qquad \gcd(\alpha,\beta)=1,
\]
则
\[
v_2(g)=v_2(h)
\iff
\alpha,\beta\ \text{皆为奇数}
\iff
2d\mid (h-g),
\]
这正与原始交非空判据一致。
于是
\[
S_g(M)=\left\{k:\ k\equiv \frac{a}{2}\pmod a\right\},
\qquad
S_h(M)=\left\{k:\ k\equiv \frac{b}{2}\pmod b\right\}.
\]
又因为 \(v_2(g)=v_2(h)\)，故
\[
\frac{g}{d},\qquad \frac{h}{d}
\]
都是奇数，从而
\[
\frac{L}{a}=\frac{g}{d},
\qquad
\frac{L}{b}=\frac{h}{d}
\]
亦皆为奇数。于是
\[
\frac{L}{2}-\frac{a}{2}
=
a\cdot \frac{L/a-1}{2}\in a\ZZ,
\qquad
\frac{L}{2}-\frac{b}{2}
=
b\cdot \frac{L/b-1}{2}\in b\ZZ,
\]
即
\[
\frac{L}{2}\equiv \frac{a}{2}\pmod a,
\qquad
\frac{L}{2}\equiv \frac{b}{2}\pmod b.
\]
故两组同余相容。由中国剩余定理，其公共解恰为单一余类
\[
k\equiv \frac{L}{2}\pmod L.
\]
而 \(L=M/d\)，故该余类正是
\[
S_d(M)=S_{\gcd(g,h)}(M).
\]
证毕。
\end{proof}

\begin{corollary}[Fibonacci 权重零点集合的 \texorpdfstring{$2$}{2}-进分层交半格分解]\label{cor:xi-fold-zero-set-v2-semilattice-decomposition}
沿用推论 \ref{cor:fold-zero-coset-union} 的记号，令
\[
M:=F_{m+2},
\qquad
\mathcal Z_m
=
\bigcup_{\substack{1\le j\le m\\ g_j\mid M/2}}S_{g_j}(M),
\qquad
g_j:=\gcd(F_{j+1},M)=F_{\gcd(j+1,m+2)}.
\]
再定义
\[
\mathcal G_m:=\left\{g_j:\ 1\le j\le m,\ g_j\mid \frac{M}{2}\right\},
\qquad
\mathcal G_{m,r}:=\{g\in\mathcal G_m:\ v_2(g)=r\},
\]
\[
\mathcal Z_{m,r}:=\bigcup_{g\in\mathcal G_{m,r}}S_g(M).
\]
则有：
\begin{enumerate}
  \item 各层 \(\mathcal Z_{m,r}\) 两两不交，且
  \[
  \boxed{\;
  \mathcal Z_m=\bigsqcup_r \mathcal Z_{m,r}.\;}
  \]
  \item 对每个固定的 \(r\)，集合族 \(\{S_g(M): g\in \mathcal G_{m,r}\}\) 对有限交封闭，并满足
  \[
  \boxed{\;
  S_g(M)\cap S_h(M)=S_{\gcd(g,h)}(M)
  \qquad (g,h\in \mathcal G_{m,r}).\;}
  \]
  \item 若 \(g,h\in \mathcal G_{m,r}\) 且 \(g\mid h\)（等价地，\(h/g\) 为奇数），则
  \[
  \boxed{\;
  S_g(M)\subseteq S_h(M).\;}
  \]
\end{enumerate}
因此，每一层 \(\mathcal Z_{m,r}\) 都是一族按 \(\gcd\) 封闭的有限交半格并；特别地，该层中的整除链只能沿奇倍方向嵌套，故只需取整除关系下的极大元作并即可生成。
\end{corollary}
\begin{proof}
第一条与第二条直接由定理 \ref{thm:xi-fold-zero-coset-v2-intersection-rigidity} 给出：不同 \(2\)-进层的余类两两不交，而同层中任意两条余类的交恰为对应的 \(\gcd\)-余类。

第三条中，若 \(g,h\in\mathcal G_{m,r}\) 且 \(g\mid h\)，则
\[
\gcd(g,h)=g.
\]
于是由第二条
\[
S_g(M)
=
S_g(M)\cap S_h(M)
\subseteq
S_h(M).
\]
故同层中较小的除数对应较小的余类，而整个层只需取整除极大元作并即可生成。证毕。
\end{proof}

\leanverified{paper\_xi\_fold\_zero\_dyadic\_tower\_disjoint\_fibonacci\_cosets}
\begin{theorem}[六倍窗零块的 dyadic 塔给出一列两两不交的 Fibonacci 余类块]\label{thm:xi-fold-zero-dyadic-tower-disjoint-fibonacci-cosets}
沿用
\[
M:=F_{m+2},
\qquad
\mathcal Z_m:=\{k\in \ZZ/(M\ZZ):\ \widehat d_m(k)=0\}
\]
的记号。设
\[
m+2=3\cdot 2^r\cdot s,
\qquad
r\ge 1,
\qquad
s\in \NN.
\]
对每个 \(j=1,\dots,r\)，定义
\[
d_j:=\frac{m+2}{2^j},
\qquad
g_j:=F_{d_j}.
\]
则对所有 \(1\le j\le r\)，都有
\[
S_{g_j}(M)\subseteq \mathcal Z_m,
\]
并且这些余类块两两不交：
\[
S_{g_i}(M)\cap S_{g_j}(M)=\varnothing
\qquad(i\neq j).
\]
因此
\[
\boxed{\;
|\mathcal Z_m|
\ge
\sum_{j=1}^{r}F_{(m+2)/2^j}.\;}
\]
\end{theorem}
\begin{proof}
先证包含关系。对每个 \(j\)，由于 \(3\mid d_j\)，Lucas 数列模 \(2\) 的周期性给出
\[
2\mid L_{d_j}.
\]
更一般地，对所有 \(t=0,\dots,j-1\)，都有
\[
3\mid 2^t d_j,
\]
故
\[
2\mid L_{2^t d_j}.
\]
另一方面，反复使用倍角恒等式
\[
F_{2n}=F_nL_n
\]
可得
\[
F_{m+2}
=
F_{2^j d_j}
=
F_{d_j}\prod_{t=0}^{j-1}L_{2^t d_j}.
\]
由于上式乘积中的每个 Lucas 因子都为偶数，故
\[
F_{d_j}\mid \frac{F_{m+2}}{2}=\frac{M}{2}.
\]
于是由推论 \ref{cor:fold-zero-divisor-lb} 立即得到
\[
S_{g_j}(M)=S_{F_{d_j}}(M)\subseteq \mathcal Z_m,
\qquad
|S_{g_j}(M)|=F_{d_j}.
\]

再证两两不交。设 \(i<j\)。则
\[
d_i=2^{j-i}d_j,
\qquad
g_i=F_{2^{j-i}d_j},
\qquad
g_j=F_{d_j}.
\]
由 Fibonacci 最大公因子恒等式，
\[
\gcd(g_i,g_j)
=
\gcd(F_{d_i},F_{d_j})
=
F_{\gcd(d_i,d_j)}
=
F_{d_j}
=
g_j.
\]
若 \(S_{g_i}(M)\cap S_{g_j}(M)\neq\varnothing\)，则由引理 \ref{lem:fold-zero-coset-intersection} 必有
\[
2g_j\mid (g_i-g_j).
\]
但由前述倍角分解，
\[
g_i-g_j
=
F_{2^{j-i}d_j}-F_{d_j}
=
F_{d_j}\left(\prod_{t=0}^{j-i-1}L_{2^t d_j}-1\right).
\]
而每个 \(L_{2^t d_j}\) 都是偶数，所以括号中是“偶数减一”，必为奇数。故
\[
2F_{d_j}\nmid g_i-g_j,
\]
即
\[
2g_j\nmid (g_i-g_j),
\]
矛盾。于是
\[
S_{g_i}(M)\cap S_{g_j}(M)=\varnothing
\qquad(i\neq j).
\]

最后，由这些两两不交的余类块都包含于 \(\mathcal Z_m\)，得到
\[
|\mathcal Z_m|
\ge
\sum_{j=1}^{r}|S_{g_j}(M)|
=
\sum_{j=1}^{r}F_{d_j}
=
\sum_{j=1}^{r}F_{(m+2)/2^j}.
\]
证毕。
\end{proof}

\begin{corollary}[六倍窗零点集合半维指数律的尖锐性可由 dyadic 零块塔重导]\label{cor:xi-fold-zero-dyadic-tower-sharp-halfdim-exponent}
沿同余子序列
\[
m+2\equiv 0\pmod 6,
\]
有
\[
\lim_{\substack{m\to\infty\\ m+2\equiv 0\ (6)}}
\frac{1}{m}\log|\mathcal Z_m|
=
\frac12\log\varphi.
\]
等价地，
\[
|\mathcal Z_m|=\varphi^{m/2+o(m)}.
\]
因此定理 \ref{thm:xi-fold-zero-set-half-dimensional-exponent-law} 中的半维指数率上界在六倍窗上是尖锐的。
\end{corollary}
\begin{proof}
当 \(m+2\equiv 0\pmod 6\) 时，可写
\[
m+2=3\cdot 2^r\cdot s
\]
且 \(r\ge 1\)。由定理 \ref{thm:xi-fold-zero-dyadic-tower-disjoint-fibonacci-cosets}，至少有
\[
|\mathcal Z_m|
\ge
F_{(m+2)/2}.
\]
另一方面，由定理 \ref{thm:fold-zero-density-sparse}，
\[
|\mathcal Z_m|
\le
\tau(m+2)\,F_{\lfloor (m+2)/2\rfloor}.
\]
再结合 Fibonacci 渐近
\[
F_n=\varphi^{n+o(n)}
\]
以及约数函数的次指数增长
\[
\tau(n)=e^{o(n)},
\]
便得
\[
\frac{1}{m}\log F_{(m+2)/2}
\longrightarrow
\frac12\log\varphi,
\]
且
\[
\frac{1}{m}\log\!\Bigl(\tau(m+2)F_{\lfloor (m+2)/2\rfloor}\Bigr)
\longrightarrow
\frac12\log\varphi.
\]
夹逼定理即得所示极限。证毕。
\end{proof}


\begin{corollary}[最小三轴窗零点集合的三余类完全分解]\label{cor:xi-fold-zero-m58-three-coset-exhaustion}
\leanverified{paper\_xi\_fold\_zero\_m58\_three\_coset\_exhaustion}
当 \(m=58\) 时，有
\[
\mathcal Z_{58}
=
S_{F_{15}}(F_{60})
\sqcup
S_{F_{20}}(F_{60})
\sqcup
S_{F_{30}}(F_{60}),
\]
从而
\[
|\mathcal Z_{58}|=F_{15}+F_{20}+F_{30}=610+6765+832040=839415.
\]
换言之，最小三轴窗上的全部谱零点恰被三条由约数格触发的 Fibonacci 余类无重叠地完全耗尽。
\end{corollary}
\begin{proof}
记 \(M:=F_{60}\)。由
\[
\frac{F_{60}}{2}
=
774004377960
=
610\cdot 1268859636
=
6765\cdot 114413064,
\]
可见
\[
F_{15}\mid \frac{F_{60}}{2},
\qquad
F_{20}\mid \frac{F_{60}}{2}.
\]
另一方面，由推论 \ref{cor:fold-zero-half-index-multiple6}，
\[
F_{30}\mid \frac{F_{60}}{2}
\qquad\text{且}\qquad
S_{F_{30}}(F_{60})\subseteq \mathcal Z_{58}.
\]
故由推论 \ref{cor:fold-zero-divisor-lb}，
\[
S_{F_{15}}(F_{60})\subseteq \mathcal Z_{58},
\qquad
S_{F_{20}}(F_{60})\subseteq \mathcal Z_{58},
\qquad
S_{F_{30}}(F_{60})\subseteq \mathcal Z_{58}.
\]
再由定理 \ref{thm:xi-fold-zero-divisor-two-class-inclusion-exclusion} 的交叠判据，
\[
2F_{\gcd(15,20)}=2F_5=10\nmid(F_{20}-F_{15})=6155,
\]
\[
2F_{\gcd(15,30)}=2F_{15}=1220\nmid(F_{30}-F_{15})=831430,
\]
\[
2F_{\gcd(20,30)}=2F_{10}=110\nmid(F_{30}-F_{20})=825275.
\]
因此三条余类两两不交，从而
\[
\left|
S_{F_{15}}(F_{60})
\sqcup
S_{F_{20}}(F_{60})
\sqcup
S_{F_{30}}(F_{60})
\right|
=
F_{15}+F_{20}+F_{30}
=
839415.
\]
而注 \ref{rem:fold-zero-lb-m58} 已给出
\[
|\mathcal Z_{58}|=839415.
\]
故上述三余类并集必与整个 \(\mathcal Z_{58}\) 相等。证毕。
\end{proof}

\begin{theorem}[无界逆分支深度强制无限层 prime-register 外置]\label{thm:xi-natural-extension-infinite-prime-register-necessity}
\leanverified{paper\_xi\_natural\_extension\_infinite\_prime\_register\_necessity}
设 \(T:X\twoheadrightarrow X\) 为有限纤维满射，并假设其逆向分支深度无界，即对每个整数 \(N\ge 1\)，都存在长度 \(N\) 的前史片段
\[
x_{-N}\xrightarrow{\,T\,}x_{-N+1}\xrightarrow{\,T\,}\cdots\xrightarrow{\,T\,}x_{-1}\xrightarrow{\,T\,}x_0
\]
满足
\[
\#T^{-1}(x_{-j+1})\ge 2
\qquad(1\le j\le N).
\]
则在定理 \ref{thm:xi-layered-prime-register-sharp-natural-extension} 的分层 prime-register 外置语义下，不存在任何仅使用有限个 prime slices 的精确单射实现能够把自然扩张
\[
X^{\leftarrow}
:=
\{(x_0,x_{-1},x_{-2},\dots):T(x_{-n})=x_{-n+1}\}
\]
共轭到某个有限层账本系统上。
等价地，对全自然扩张的精确可逆外置，所需有效 prime-register 层数必为无穷。

另一方面，定理 \ref{thm:killo-natural-extension-branch-register} 给出共轭
\[
X^{\leftarrow}\cong \widehat X\subseteq X\times\NN^\NN,
\]
故可数无穷层账本足以完成精确编码。
\end{theorem}
\begin{proof}
反设存在某个有限整数 \(k\ge 1\)，使自然扩张在同一分层 prime-register 语义下可由仅含 \(k\) 个 prime slices 的账本精确单射外置。
由于该外置与自然扩张共轭，其码字在像空间上具有精确逆映射；再与前史截断
\[
\pi_N:X^{\leftarrow}\to X^{N+1},
\qquad
\pi_N(x_0,x_{-1},x_{-2},\dots):=(x_0,x_{-1},\dots,x_{-N})
\]
复合，便知同一有限层账本可对任意固定深度 \(N\) 的前史片段作精确恢复。

取 \(N:=k+1\)。
由无界逆分支深度假设，存在一条长度 \(N\) 的前史片段，其前 \(N\) 层每层都至少出现一次二歧分支。
于是上述有限层账本在该前史族上给出一个深度 \(N\) 的精确分层外置，而所用 prime slices 仍只有原来的 \(k\) 个。
这与定理 \ref{thm:xi-layered-prime-register-sharp-natural-extension} 矛盾：在最坏情形下，深度 \(N\) 的精确分层外置至少需要 \(N\) 个有效 prime slices。

故有限层 prime-register 外置不可能精确实现整个自然扩张。
另一方面，定理 \ref{thm:killo-natural-extension-branch-register} 已把 \(X^{\leftarrow}\) 共轭到 \(X\times\NN^\NN\) 上的右移插入系统，说明可数无穷层账本足以给出精确实现。证毕。
\end{proof}

\begin{theorem}[共振整函数零点计数的线性主项]\label{thm:xi-fold-resonance-entire-zero-counting-linear}
沿用定理 \ref{thm:fold-resonance-entire-lp} 的记号，定义
\[
\mathcal F_\varphi(z):=\prod_{n=2}^{\infty}\cos(\pi z\varphi^{-n}).
\]
则 \(\mathcal F_\varphi\) 属于 Laguerre--P\'olya 类，且全部零点均为实单零点：
\[
\mathcal Z(\mathcal F_\varphi)=\Bigl\{(k+\tfrac12)\varphi^n:\ k\in\ZZ,\ n\ge 2\Bigr\}.
\]
若记正零点计数函数
\[
N_+(R):=\#\{x\in(0,R]:\ \mathcal F_\varphi(x)=0\},
\]
则当 \(R\to\infty\) 时有
\[
\boxed{\;
N_+(R)=R+O(\log R).\;}
\]
相应地，双侧零点计数满足
\[
\boxed{\;
\#\bigl(\mathcal Z(\mathcal F_\varphi)\cap[-R,R]\bigr)=2R+O(\log R).\;}
\]
\end{theorem}
\begin{proof}
由定理 \ref{thm:fold-resonance-entire-lp}，正零点恰为
\[
\Bigl\{(k+\tfrac12)\varphi^n:\ k\in\ZZ_{\ge 0},\ n\ge 2\Bigr\}.
\]
对每个固定的 \(n\ge 2\)，落在 \((0,R]\) 内的该层零点数为
\[
N_n(R)
:=
\#\Bigl\{k\in\ZZ_{\ge 0}:\ (k+\tfrac12)\varphi^n\le R\Bigr\}
=
\left\lfloor \frac{R}{\varphi^n}+\frac12\right\rfloor
=
\frac{R}{\varphi^n}+O(1).
\]
仅当 \(\varphi^n\le 2R\) 时该项才可能非零，故
\[
N_+(R)
=
\sum_{n=2}^{\lfloor \log_\varphi(2R)\rfloor}N_n(R)
=
R\sum_{n=2}^{\lfloor \log_\varphi(2R)\rfloor}\varphi^{-n}
+O(\log R).
\]
又因为
\[
\sum_{n=2}^{\infty}\varphi^{-n}=1,
\]
且有限截断尾和为 \(O(R^{-1})\)，遂得
\[
N_+(R)=R+O(\log R).
\]
最后，由零点集关于原点对称，立得
\[
\#\bigl(\mathcal Z(\mathcal F_\varphi)\cap[-R,R]\bigr)=2R+O(\log R).
\]
证毕。
\end{proof}

\paragraph{fold-gauge 常数对 \texorpdfstring{$2\pi$}{2pi} 的最终单侧逼近与差分完全单调性}
在 \(C_m\) 的 Stirling--Bernoulli 层级、首个显式消模器以及任意阶 Bernoulli 消模加速已经建立之后，原始误差序列本身还可进一步压缩为一条具有固定差分符号的尾部刚性链。

\begin{theorem}[fold-gauge 常数最终从上方单调逼近 \texorpdfstring{$2\pi$}{2pi}，并具有任意固定阶的最终完全单调性]\label{thm:xi-foldbin-gauge-constant-eventual-complete-monotonicity}
沿用推论 \ref{cor:fold-bin-recover-2pi} 的规范常数列 \((C_m)_{m\ge 1}\)，并定义
\[
E_m:=\log C_m-\log(2\pi),
\qquad
(\Delta u)_m:=u_{m+1}-u_m.
\]
则成立：
\begin{enumerate}
  \item 存在 \(m_0\)，使得对全部 \(m\ge m_0\) 都有
  \[
  C_m>2\pi,
  \qquad
  C_{m+1}<C_m.
  \]
  \item 对每个固定整数 \(r\ge 0\)，存在 \(m_r\)，使得对全部 \(m\ge m_r\) 都有
  \[
  (-1)^r(\Delta^r E)_m>0.
  \]
  换言之，误差序列 \((E_m)_{m\ge 1}\) 在充分大 \(m\) 上具有任意固定阶的最终完全单调性。
  \item 若记
  \[
  q:=\frac{\varphi}{2}\in(0,1),
  \]
  则有显式主导展开
  \[
  \boxed{\;
  \log C_m
  =
  \log(2\pi)
  +\frac{1}{3\varphi}q^m
  -\frac{\sqrt5}{180}q^{3m}
  +O(q^{5m}).\;}
  \]
  从而
  \[
  \boxed{\;
  \lim_{m\to\infty}
  \Bigl(\frac{2}{\varphi}\Bigr)^m\bigl(C_m-2\pi\bigr)
  =
  \frac{2\pi}{3\varphi}.\;}
  \]
\end{enumerate}
\end{theorem}
\begin{proof}
由定理 \ref{thm:fold-bin-gauge-constant-stirling-bernoulli-hierarchy} 在 \(R=2\) 时的展开，或等价地由定理 \ref{thm:xi-foldbin-gauge-constant-not-cfinite-first-two-zeta-limits}，有
\[
E_m
=
\frac{1}{3\varphi}q^m
-\frac{\sqrt5}{180}q^{3m}
+O(q^{5m}).
\]
这就得到第 \(3\) 条中的主导展开。

现固定 \(r\ge 0\)。对任意 \(\alpha\in(0,1)\)，几何序列满足
\[
\Delta^r(\alpha^m)=(\alpha-1)^r\alpha^m.
\]
因此
\[
\Delta^r E_m
=
\frac{1}{3\varphi}(q-1)^r q^m
-\frac{\sqrt5}{180}(q^3-1)^r q^{3m}
+O(q^{5m}).
\]
把右端提取出主项，可写为
\[
\Delta^r E_m
=
\frac{1}{3\varphi}(q-1)^r q^m
\left(
1+O(q^{2m})
\right).
\]
由于 \(q\in(0,1)\) 且 \(1/(3\varphi)>0\)，主项符号恰为 \((q-1)^r\) 的符号，即 \((-1)^r\)。故存在 \(m_r\)，使得对全部 \(m\ge m_r\) 都有
\[
(-1)^r(\Delta^r E)_m>0.
\]
这便得到第 \(2\) 条。

取 \(r=0\) 得 \(E_m>0\) 对充分大 \(m\) 成立，于是
\[
\log C_m>\log(2\pi)
\qquad\Longrightarrow\qquad
C_m>2\pi.
\]
取 \(r=1\) 得
\[
(\Delta E)_m<0
\]
对充分大 \(m\) 成立，即 \(E_{m+1}<E_m\)。由于指数函数严格递增，遂有
\[
C_{m+1}<C_m
\]
对充分大 \(m\) 也成立，从而得到第 \(1\) 条。

最后，由 \(E_m=O(q^m)\) 与指数展开
\[
e^{E_m}=1+E_m+O(E_m^2)
\]
可得
\[
C_m-2\pi
=
2\pi\bigl(E_m+O(E_m^2)\bigr)
=
\frac{2\pi}{3\varphi}q^m+O(q^{2m}).
\]
两边同乘 \(q^{-m}=(2/\varphi)^m\)，即得
\[
\lim_{m\to\infty}
\Bigl(\frac{2}{\varphi}\Bigr)^m\bigl(C_m-2\pi\bigr)
=
\frac{2\pi}{3\varphi}.
\]
证毕。
\end{proof}

\endinput

\paragraph{fold-gauge 规范常数的显式三模渐近、归一化缺陷与超收敛上包络}
在最终单侧逼近、任意固定阶差分尾符号与有限阶 Bernoulli 模态消去都已建立之后，规范常数列还可继续压缩为一组完全显式的有效渐近公式。它们把 Stirling--Bernoulli 层级的前三个非平凡模态、首项消元后的具体上侧包络以及相邻比值的首个二次修正统一到同一条闭式链中。

\begin{theorem}[\texorpdfstring{$\log C_m$}{log Cm} 的前三个非平凡 Bernoulli 模态完全显式化]\label{thm:xi-time-part59ac-foldgauge-three-bernoulli-modes}
记
\[
q_0:=\frac{\varphi}{2},
\qquad
L_m:=\log C_m-\log(2\pi).
\]
则有精确三项展开
\[
\boxed{\;
L_m
=
\frac{1}{3\varphi}\,q_0^m
-\frac{\sqrt5}{180}\,q_0^{3m}
+\frac{\varphi}{210}\,q_0^{5m}
+O(q_0^{7m}).\;}
\]
从而
\[
\boxed{\;
C_m
=
2\pi\exp\!\left(
\frac{1}{3\varphi}\,q_0^m
-\frac{\sqrt5}{180}\,q_0^{3m}
+\frac{\varphi}{210}\,q_0^{5m}
+O(q_0^{7m})
\right).\;}
\]
进一步展开指数，得到
\[
\boxed{\;
C_m
=
2\pi\left(
1+\frac{1}{3\varphi}\,q_0^m
+\frac{1}{18\varphi^2}\,q_0^{2m}
+O(q_0^{3m})
\right).\;}
\]
\end{theorem}
\begin{proof}
对 \(R=3\) 应用定理 \ref{thm:fold-bin-gauge-constant-stirling-bernoulli-hierarchy}，得到
\[
L_m=\sum_{r=1}^{3}a_r q_0^{(2r-1)m}+O(q_0^{7m}),
\qquad
a_r=
\frac{B_{2r}}{r(2r-1)}
\bigl(\varphi^{-1}+\varphi^{2r-3}\bigr).
\]
对 \(r=1,2,3\) 分别计算系数。由
\[
B_2=\frac16,\qquad B_4=-\frac1{30},\qquad B_6=\frac1{42},
\]
得
\[
a_1=\frac16\bigl(\varphi^{-1}+\varphi^{-1}\bigr)=\frac{1}{3\varphi},
\]
\[
a_2=
\frac{-1/30}{2\cdot 3}\bigl(\varphi^{-1}+\varphi\bigr)
=-\frac{1}{180}\bigl(\varphi+\varphi^{-1}\bigr)
=-\frac{\sqrt5}{180},
\]
其中用了恒等式 \(\varphi+\varphi^{-1}=\sqrt5\)。又
\[
a_3=
\frac{1/42}{3\cdot 5}\bigl(\varphi^{-1}+\varphi^3\bigr).
\]
由
\[
\varphi^3=2\varphi+1,
\qquad
\varphi^{-1}=\varphi-1,
\]
可得
\[
\varphi^{-1}+\varphi^3=3\varphi,
\]
因此
\[
a_3=\frac{3\varphi}{630}=\frac{\varphi}{210}.
\]
代回即得 \(L_m\) 的三项展开。

由 \(C_m=2\pi e^{L_m}\)，立得第二式。再由 \(L_m=O(q_0^m)\) 可作指数展开
\[
e^{L_m}=1+L_m+\frac12L_m^2+O(L_m^3)
=1+\frac{1}{3\varphi}\,q_0^m+\frac{1}{18\varphi^2}\,q_0^{2m}+O(q_0^{3m}),
\]
从而得到第三式。证毕。
\end{proof}

\begin{theorem}[fold-gauge 常数对 \texorpdfstring{$2\pi$}{2pi} 的归一化缺陷最终严格递增，并从下方趋向主系数]\label{thm:xi-time-part59ac-foldgauge-normalized-defect-monotone}
沿用定理 \ref{thm:xi-time-part59ac-foldgauge-three-bernoulli-modes} 的记号，并定义
\[
\widetilde D_m:=q_0^{-m}L_m.
\]
则存在 \(m_0\)，使得对全部 \(m\ge m_0\) 都有
\[
C_m>2\pi,
\qquad
C_{m+1}<C_m,
\qquad
\widetilde D_m<\frac{1}{3\varphi},
\qquad
\widetilde D_{m+1}>\widetilde D_m.
\]
并且
\[
\lim_{m\to\infty}\widetilde D_m=\frac{1}{3\varphi}.
\]
\end{theorem}
\begin{proof}
由定理 \ref{thm:xi-time-part59ac-foldgauge-three-bernoulli-modes}，
\[
L_m
=
\frac{1}{3\varphi}\,q_0^m
-\frac{\sqrt5}{180}\,q_0^{3m}
+O(q_0^{5m}).
\]
首项系数为正，故对充分大的 \(m\) 有 \(L_m>0\)，即
\[
\log C_m>\log(2\pi),
\]
从而 \(C_m>2\pi\)。

再作差分，得到
\[
L_{m+1}-L_m
=
\frac{1}{3\varphi}(q_0-1)\,q_0^m
-\frac{\sqrt5}{180}(q_0^3-1)\,q_0^{3m}
+O(q_0^{5m}).
\]
由于 \(0<q_0<1\)，第一项为负且支配其余项，故对充分大的 \(m\) 有
\[
L_{m+1}<L_m.
\]
指数函数严格递增，遂得 \(C_{m+1}<C_m\)。

对归一化缺陷，直接由上式除以 \(q_0^m\) 得
\[
\widetilde D_m
=
\frac{1}{3\varphi}
-\frac{\sqrt5}{180}\,q_0^{2m}
+O(q_0^{4m}),
\]
从而对充分大的 \(m\) 有
\[
\widetilde D_m<\frac{1}{3\varphi},
\qquad
\lim_{m\to\infty}\widetilde D_m=\frac{1}{3\varphi}.
\]
再作差可得
\[
\widetilde D_{m+1}-\widetilde D_m
=
\frac{\sqrt5}{180}(1-q_0^2)\,q_0^{2m}
+O(q_0^{4m}),
\]
其主项严格为正，故对充分大的 \(m\) 有
\[
\widetilde D_{m+1}>\widetilde D_m.
\]
证毕。
\end{proof}

\begin{theorem}[首项消元后的立方级超收敛上侧包络]\label{thm:xi-time-part59ac-foldgauge-first-mode-cancellation-envelope}
沿用定理 \ref{thm:xi-time-part59ac-foldgauge-three-bernoulli-modes} 的记号，并定义
\[
U_m:=\frac{\log C_{m+1}-q_0\log C_m}{1-q_0}.
\]
则有精确渐近
\[
\boxed{\;
U_m
=
\log(2\pi)
+\frac{\sqrt5}{180}(1+q_0)\,q_0^{3m+1}
+O(q_0^{5m}).\;}
\]
特别地，存在 \(m_1\)，使得对全部 \(m\ge m_1\) 都有
\[
\log(2\pi)<U_m<\log C_m.
\]
因此，\(U_m\) 给出 \(\log(2\pi)\) 的一个严格优于 \(\log C_m\) 的上侧近似，其误差从 \(O(q_0^m)\) 提升到 \(O(q_0^{3m})\)。
\end{theorem}
\begin{proof}
由定义，
\[
U_m-\log(2\pi)
=
\frac{L_{m+1}-q_0L_m}{1-q_0}.
\]
再由定理 \ref{thm:xi-time-part59ac-foldgauge-three-bernoulli-modes}，
\[
L_m
=
a_1q_0^m+a_2q_0^{3m}+O(q_0^{5m}),
\qquad
a_1=\frac{1}{3\varphi},
\qquad
a_2=-\frac{\sqrt5}{180}.
\]
于是
\[
L_{m+1}-q_0L_m
=
a_2q_0^{3m+1}(q_0^2-1)+O(q_0^{5m}),
\]
即首个 Bernoulli 模态被完全消去。两边除以 \(1-q_0\)，并利用
\[
\frac{q_0^2-1}{1-q_0}=-(1+q_0),
\]
得到
\[
U_m-\log(2\pi)
=
-a_2(1+q_0)\,q_0^{3m+1}+O(q_0^{5m})
=
\frac{\sqrt5}{180}(1+q_0)\,q_0^{3m+1}+O(q_0^{5m}).
\]
故对充分大的 \(m\) 有 \(U_m>\log(2\pi)\)。

另一方面，
\[
\log C_m-\log(2\pi)=L_m=\frac{1}{3\varphi}\,q_0^m+O(q_0^{3m}),
\]
其主阶为 \(q_0^m\)，而 \(U_m-\log(2\pi)=O(q_0^{3m})\)。于是对充分大的 \(m\) 有
\[
U_m<\log C_m.
\]
证毕。
\end{proof}

\begin{theorem}[连续比值极限的二次修正与普适有理常数 \texorpdfstring{$1/48$}{1/48}]\label{thm:xi-time-part59ac-foldgauge-ratio-one-forty-eighth}
\leanverified{paper\_xi\_time\_part59ac\_foldgauge\_ratio\_one\_forty\_eighth}
沿用定理 \ref{thm:xi-time-part59ac-foldgauge-three-bernoulli-modes} 的记号，则有精确渐近
\[
\boxed{\;
\frac{L_{m+1}}{q_0L_m}
=
1+\frac{1}{48}\,q_0^{2m}+O(q_0^{4m}).\;}
\]
因此
\[
\lim_{m\to\infty}\frac{L_{m+1}}{L_m}=q_0=\frac{\varphi}{2},
\]
而且该比值从上方趋向 \(\varphi/2\)；其首个二次修正系数恰为普适有理常数 \(1/48\)。
\end{theorem}
\begin{proof}
由定理 \ref{thm:xi-time-part59ac-foldgauge-three-bernoulli-modes}，
\[
L_m
=
a_1q_0^m+a_2q_0^{3m}+O(q_0^{5m})
=
a_1q_0^m\left(1+\frac{a_2}{a_1}q_0^{2m}+O(q_0^{4m})\right),
\]
其中
\[
a_1=\frac{1}{3\varphi},
\qquad
a_2=-\frac{\sqrt5}{180}.
\]
故
\[
\frac{L_{m+1}}{q_0L_m}
=
\frac{1+\frac{a_2}{a_1}q_0^{2m+2}+O(q_0^{4m})}
{1+\frac{a_2}{a_1}q_0^{2m}+O(q_0^{4m})}
=
1+\frac{a_2}{a_1}(q_0^2-1)\,q_0^{2m}+O(q_0^{4m}).
\]
现化简系数。由
\[
\frac{a_2}{a_1}
=
-\frac{\sqrt5}{180}\cdot 3\varphi
=
-\frac{\sqrt5\,\varphi}{60},
\]
以及
\[
q_0^2-1=\frac{\varphi^2}{4}-1=\frac{\varphi-3}{4},
\]
可得
\[
\frac{a_2}{a_1}(q_0^2-1)
=
\frac{\sqrt5\,\varphi(3-\varphi)}{240}.
\]
再由
\[
\varphi=\frac{1+\sqrt5}{2},
\qquad
3-\varphi=\frac{5-\sqrt5}{2},
\]
有
\[
\sqrt5\,\varphi(3-\varphi)
=
\sqrt5\cdot\frac{(1+\sqrt5)(5-\sqrt5)}{4}
=
5.
\]
因此
\[
\frac{a_2}{a_1}(q_0^2-1)=\frac{5}{240}=\frac{1}{48}.
\]
代回即得所示展开。由于修正项系数为正，故该比值对充分大的 \(m\) 从上方趋向 \(1\)；等价地，
\[
\frac{L_{m+1}}{L_m}
=
q_0\left(1+\frac{1}{48}\,q_0^{2m}+O(q_0^{4m})\right)
\]
从上方趋向 \(q_0=\varphi/2\)。证毕。
\end{proof}

\endinput

在 Bernoulli 层级与最终单调逼近之外，fold-gauge 常数序列的极限本身还保留了一个不能被代数闭包替代的超越指纹。

\begin{proposition}[圆常数的折叠恢复]\label{prop:xi-time-part59-foldgauge-circle-constant-recovery}
\leanverified{paper\_xi\_time\_part59\_foldgauge\_circle\_constant\_recovery}
沿用推论 \ref{cor:fold-bin-recover-2pi} 的记号，定义
\[
C_m
:=
\exp\!\Biggl(
\frac{2^{m+1}}{|X_m|}\bigl(g_m-\kappa_m+1\bigr)
-\frac{1}{|X_m|}\sum_{w\in X_m}\log d_m^{\mathrm{bin}}(w)
\Biggr).
\]
则
\[
\lim_{m\to\infty}C_m=2\pi.
\]
因此，这一完全由离散折叠数据构成的常数列在极限上强制恢复出圆常数；尤其，任何要求把该极限常数代数化的替代描述都不可能与原序列的极限严格一致。
\end{proposition}
\begin{proof}
推论 \ref{cor:fold-bin-recover-2pi} 已给出上述极限。另一方面，\(\pi\) 的超越性是经典事实，故 \(2\pi\) 亦为超越数。于是，若某种替代描述宣称与 \((C_m)\) 具有同一极限而该极限仍为代数常数，便与
\[
\lim_{m\to\infty}C_m=2\pi
\]
矛盾。证毕。
\end{proof}

\endinput


\noindent 上述结果表明：一方面，fold-gauge 常数的 Stirling--Bernoulli 层级不仅允许逐阶恢复全部偶 Bernoulli 数、构造任意预先指定阶数的显式消模器，而且其前三个非平凡 Bernoulli 模态、归一化缺陷的最终单调趋近、首项消元后的立方级上侧包络以及相邻比值的普适 \(\frac{1}{48}\) 二次修正也都可完全显式写出；原始误差序列 \(\log C_m-\log(2\pi)\) 本身则在充分大 \(m\) 上呈现任意固定阶的最终完全单调性，其极限更由纯离散折叠数据内生恢复出超越常数 \(2\pi\)。另一方面，与定理 \ref{thm:xi-output-potential-dirichlet-fixed-j-fourth-order}、推论 \ref{cor:xi-output-potential-dirichlet-bounded-window-j1} 及定理 \ref{thm:xi-output-potential-dirichlet-fourth-order-gain} 所建立的单位根扭曲四阶展开链合并后，prime-shadow 侧的解析结构同样形成闭合：六倍窗零点集合的半维指数率被精确固定，而局部累积量则把最坏 Dirichlet 扭曲的局部模态、窗口极值机制与高阶近似收益统一到同一输出位势接口之下。于是相关结论不再只是审计意义上的现象陈列，而具备了可独立组织为主定理链的解析骨架。

\paragraph{推前熵账本、Rényi 散度冻结、临界共振底噪与 window-\texorpdfstring{$6$}{6} 编码冗余}
折叠推前分布的 Shannon/Rényi 散度账本、规范体积密度的 Stirling 差常数、临界共振常数强制的二阶底噪，以及 window-\(6\) 的最小完备奇偶荷签名之间，还存在下列可由既有结论直接闭合的关系。

\begin{theorem}[折叠推前熵与规范体积密度的单位 nat 差额]\label{thm:xi-foldbin-entropy-gauge-one-nat-gap}
\leanverified{paper\_xi\_foldbin\_entropy\_gauge\_one\_nat\_gap}
沿用定理 \ref{thm:xi-gauge-volume-boltzmann-stirling-exponential} 与
\ref{thm:xi-foldbin-gauge-escort-kl-decomposition} 的记号。记
\[
p_m^{\mathrm{bin}}(w):=\frac{d_m^{\mathrm{bin}}(w)}{2^m},
\qquad
u_m(w):=\frac{1}{|X_m|}
\qquad (w\in X_m),
\]
并以自然对数定义 Shannon 熵
\[
H(p):=-\sum_{w\in X_m}p(w)\log p(w).
\]
则有严格恒等式
\[
\boxed{\;
H\!\bigl(p_m^{\mathrm{bin}}\bigr)=m\log 2-\kappa_m,\;}
\]
\[
\boxed{\;
D_{\mathrm{KL}}\!\bigl(p_m^{\mathrm{bin}}\|u_m\bigr)
=
\kappa_m+\log|X_m|-m\log 2.\;}
\]
从而
\[
\boxed{\;
H\!\bigl(p_m^{\mathrm{bin}}\bigr)+g_m
=
m\log 2-1+O\!\Bigl(m\bigl(\tfrac{\varphi}{2}\bigr)^m\Bigr),\;}
\]
\[
\boxed{\;
D_{\mathrm{KL}}\!\bigl(p_m^{\mathrm{bin}}\|u_m\bigr)
=
g_m+1+\log F_{m+2}-m\log 2
+O\!\Bigl(m\bigl(\tfrac{\varphi}{2}\bigr)^m\Bigr).\;}
\]
特别地，折叠推前熵与规范体积密度之和，相对满 \(m\)-比特熵 \(m\log 2\) 仅差一个普适常数 \(1\)。
\end{theorem}
\begin{proof}
记 \(d_w:=d_m^{\mathrm{bin}}(w)\) 与 \(p_w:=p_m^{\mathrm{bin}}(w)\)。由 \(\sum_{w\in X_m}d_w=2^m\) 得
\[
\begin{aligned}
H\!\bigl(p_m^{\mathrm{bin}}\bigr)
&=-\sum_{w\in X_m}\frac{d_w}{2^m}\log\frac{d_w}{2^m}\\
&=-2^{-m}\sum_{w\in X_m}d_w\log d_w
+2^{-m}\sum_{w\in X_m}d_w\,m\log 2\\
&=m\log 2-\kappa_m,
\end{aligned}
\]
即得第一式。

再由 \(u_m(w)=1/|X_m|\) 可得
\[
\begin{aligned}
D_{\mathrm{KL}}\!\bigl(p_m^{\mathrm{bin}}\|u_m\bigr)
&=\sum_{w\in X_m}p_w\log\frac{p_w}{u_m(w)}\\
&=-H\!\bigl(p_m^{\mathrm{bin}}\bigr)+\log|X_m|\\
&=\kappa_m+\log|X_m|-m\log 2.
\end{aligned}
\]
最后，由定理 \ref{thm:xi-gauge-volume-boltzmann-stirling-exponential}，
\[
g_m=\kappa_m-1+O\!\Bigl(m\bigl(\tfrac{\varphi}{2}\bigr)^m\Bigr),
\]
并且 \(|X_m|=F_{m+2}\)。将其代入前两式即得后两式。证毕。
\end{proof}

\begin{theorem}[相对均匀基线的 Rényi 散度公式与冻结相仿射律]\label{thm:xi-foldbin-renyi-divergence-freezing-affine}
\leanverified{paper\_xi\_foldbin\_renyi\_divergence\_freezing\_affine}
沿用定理 \ref{thm:xi-foldbin-entropy-gauge-one-nat-gap} 的记号。对 \(q>0\)、\(q\neq 1\)，定义 Rényi 散度
\[
D_q\!\bigl(p_m^{\mathrm{bin}}\|u_m\bigr)
:=
\frac{1}{q-1}\log\sum_{w\in X_m}
\bigl(p_m^{\mathrm{bin}}(w)\bigr)^q\,u_m(w)^{1-q}.
\]
再定义
\[
S_q(m):=\sum_{w\in X_m}\bigl(d_m^{\mathrm{bin}}(w)\bigr)^q.
\]
则对任意 \(q>0\)、\(q\neq 1\)，都有精确公式
\[
\boxed{\;
D_q\!\bigl(p_m^{\mathrm{bin}}\|u_m\bigr)
=
\log|X_m|+\frac{1}{q-1}\log\frac{S_q(m)}{2^{mq}}.\;}
\]

进一步，对整数 \(q\ge 2\)，若压力极限
\[
P_q:=\lim_{m\to\infty}\frac{1}{m}\log S_q(m)
\]
存在，则
\[
\boxed{\;
\lim_{m\to\infty}\frac{1}{m}
D_q\!\bigl(p_m^{\mathrm{bin}}\|u_m\bigr)
=
\log\varphi+\frac{P_q-q\log 2}{q-1}.\;}
\]

若再假设严格基态间隙 \(\alpha_*^{(2)}<\alpha_*\)，并记
\[
a_{\mathrm c}:=\frac{\log\varphi-g_*}{\alpha_*-\alpha_*^{(2)}},
\]
则对每个整数 \(q>a_{\mathrm c}\)，散度密度
\[
\delta(q):=
\lim_{m\to\infty}\frac{1}{m}
D_q\!\bigl(p_m^{\mathrm{bin}}\|u_m\bigr)
\]
满足显式仿射公式
\[
\boxed{\;
\delta(q)
=
\log\varphi+\frac{q(\alpha_*-\log 2)+g_*}{q-1}.\;}
\]
特别地，
\[
\boxed{\;
\lim_{q\to\infty}\delta(q)=\alpha_*-\log\!\Bigl(\frac{2}{\varphi}\Bigr).\;}
\]
\end{theorem}
\begin{proof}
由 \(u_m(w)=1/|X_m|\)，直接计算得
\[
\sum_{w\in X_m}
\bigl(p_m^{\mathrm{bin}}(w)\bigr)^q\,u_m(w)^{1-q}
=
|X_m|^{q-1}\sum_{w\in X_m}\bigl(p_m^{\mathrm{bin}}(w)\bigr)^q.
\]
而
\[
\sum_{w\in X_m}\bigl(p_m^{\mathrm{bin}}(w)\bigr)^q
=
\sum_{w\in X_m}\left(\frac{d_m^{\mathrm{bin}}(w)}{2^m}\right)^q
=
\frac{S_q(m)}{2^{mq}}.
\]
代回定义即得第一式。

若 \(q\ge 2\) 为整数且 \(P_q\) 存在，则两边除以 \(m\) 并令 \(m\to\infty\)，再用
\[
\lim_{m\to\infty}\frac{1}{m}\log|X_m|
=
\lim_{m\to\infty}\frac{1}{m}\log F_{m+2}
=
\log\varphi,
\]
便得到第二式。

最后，在严格基态间隙下，由定理 \ref{thm:fold-pressure-freezing-threshold}，对每个整数 \(q>a_{\mathrm c}\) 都有
\[
P_q=q\alpha_*+g_*.
\]
将其代入第二式即可得到 \(\delta(q)\) 的仿射公式。再令 \(q\to\infty\)，便得
\[
\delta(q)\to \log\varphi+\alpha_*-\log 2
=
\alpha_*-\log\!\Bigl(\frac{2}{\varphi}\Bigr).
\]
证毕。
\end{proof}

\begin{theorem}[临界共振常数强制 Rényi-\texorpdfstring{$2$}{2} 散度的严格正底噪]\label{thm:xi-foldbin-critical-resonance-renyi2-floor}
沿用定义 \ref{def:fold-normalized-amplitude} 与定理 \ref{thm:fold-collision-spectrum} 的记号，记
\[
\mathsf{Col}^{\mathrm{bin}}_m
:=
\sum_{w\in X_m}\bigl(p_m^{\mathrm{bin}}(w)\bigr)^2
=
\frac{1}{F_{m+2}}\sum_{k}a_m(k)^2.
\]
则 Rényi-\(2\) 散度满足严格恒等式
\[
\boxed{\;
D_2\!\bigl(p_m^{\mathrm{bin}}\|u_m\bigr)
=
\log\!\bigl(F_{m+2}\,\mathsf{Col}^{\mathrm{bin}}_m\bigr).\;}
\]
并且对每个 \(m\) 都有点态下界
\[
\boxed{\;
D_2\!\bigl(p_m^{\mathrm{bin}}\|u_m\bigr)
\ge
\log\!\bigl(1+a_m(F_m)^2+a_m(F_{m+1})^2\bigr).\;}
\]
因此，若 \(C_\varphi>0\) 为定理 \ref{thm:fold-critical-resonance-constant} 的临界共振常数，则
\[
\boxed{\;
\liminf_{m\to\infty}
D_2\!\bigl(p_m^{\mathrm{bin}}\|u_m\bigr)
\ge
\log(1+2C_\varphi^2)
\approx 0.09597061849768761.\;}
\]

等价地，若定义 Rényi-\(2\) 有效支撑
\[
N_{2,m}^{\mathrm{eff}}
:=
\exp\!\Bigl(H_2\!\bigl(p_m^{\mathrm{bin}}\bigr)\Bigr)
=
\frac{1}{\mathsf{Col}^{\mathrm{bin}}_m},
\qquad
H_2(p):=-\log\sum_w p(w)^2,
\]
则
\[
\boxed{\;
N_{2,m}^{\mathrm{eff}}
\le
\frac{F_{m+2}}{1+a_m(F_m)^2+a_m(F_{m+1})^2},\;}
\]
并且
\[
\boxed{\;
\limsup_{m\to\infty}\frac{N_{2,m}^{\mathrm{eff}}}{F_{m+2}}
\le
\frac{1}{1+2C_\varphi^2}
\approx 0.908490708498425.\;}
\]
\end{theorem}
\begin{proof}
定理 \ref{thm:xi-foldbin-renyi-divergence-freezing-affine} 在 \(q=2\) 时给出
\[
D_2\!\bigl(p_m^{\mathrm{bin}}\|u_m\bigr)
=
\log|X_m|+\log\frac{S_2(m)}{2^{2m}}.
\]
又由
\[
S_2(m)=\sum_{w\in X_m}\bigl(d_m^{\mathrm{bin}}(w)\bigr)^2
=
2^{2m}\,\mathsf{Col}^{\mathrm{bin}}_m,
\qquad
|X_m|=F_{m+2},
\]
立得第一条恒等式。

另一方面，由 Parseval 恒等式
\[
F_{m+2}\,\mathsf{Col}^{\mathrm{bin}}_m=\sum_k a_m(k)^2,
\]
并注意 \(a_m(0)=1\)，故
\[
F_{m+2}\,\mathsf{Col}^{\mathrm{bin}}_m
\ge
1+a_m(F_m)^2+a_m(F_{m+1})^2.
\]
取对数即得点态下界。

再由定理 \ref{thm:fold-critical-resonance-constant}，
\[
a_m(F_m)\to C_\varphi,
\qquad
a_m(F_{m+1})\to C_\varphi,
\]
令 \(m\to\infty\) 即得 Rényi-\(2\) 散度的正底噪下界。

最后，由 \(H_2(p_m^{\mathrm{bin}})=-\log\mathsf{Col}^{\mathrm{bin}}_m\) 可知
\[
N_{2,m}^{\mathrm{eff}}=\frac{1}{\mathsf{Col}^{\mathrm{bin}}_m}.
\]
将上式与
\(
F_{m+2}\mathsf{Col}^{\mathrm{bin}}_m
\ge
1+a_m(F_m)^2+a_m(F_{m+1})^2
\)
重排即可得到有效支撑上界；再令 \(m\to\infty\) 便得最后一式。证毕。
\end{proof}

\begin{theorem}[Rényi--2 投影熵的常数缺口]\label{thm:xi-foldbin-renyi2-projection-entropy-constant-gap}
沿用定理 \ref{thm:xi-foldbin-critical-resonance-renyi2-floor} 的记号，定义二阶投影熵
\[
H_2^{\mathrm{proj}}(m):=-\log \mathsf{Col}^{\mathrm{bin}}_m,
\]
并记
\[
\delta_2:=\log(1+2C_\varphi^2)>0.
\]
则
\[
H_2^{\mathrm{proj}}(m)
\le
\log F_{m+2}-\delta_2+o(1).
\]
数值上，
\[
\delta_2\approx 0.09597061849768761\ \text{nats}
\approx 0.138456335393854\ \text{bits}.
\]
换言之，相对于 \(|X_m|=F_{m+2}\) 的理想均匀基线，bin-fold 投影在 Rényi--\(2\) 意义下始终保留一个不可消除的常数缺口。
\end{theorem}
\begin{proof}
由定理 \ref{thm:xi-foldbin-critical-resonance-renyi2-floor}，
\[
D_2\!\bigl(p_m^{\mathrm{bin}}\|u_m\bigr)
=
\log\!\bigl(F_{m+2}\mathsf{Col}^{\mathrm{bin}}_m\bigr).
\]
于是
\[
H_2^{\mathrm{proj}}(m)
=-\log \mathsf{Col}^{\mathrm{bin}}_m
=
\log F_{m+2}
-D_2\!\bigl(p_m^{\mathrm{bin}}\|u_m\bigr).
\]
同一定理又给出
\[
\liminf_{m\to\infty}D_2\!\bigl(p_m^{\mathrm{bin}}\|u_m\bigr)\ge \delta_2,
\]
故
\[
H_2^{\mathrm{proj}}(m)\le \log F_{m+2}-\delta_2+o(1).
\]
比特值只需再除以 \(\log 2\) 即得。证毕。
\end{proof}

\begin{theorem}[临界共振常数强制 bin-fold 推前分布在 Rényi-\texorpdfstring{$2$}{2} 几何中永久远离均匀层]\label{thm:xi-foldbin-critical-resonance-chi2-renyi2-separation}
沿用定理 \ref{thm:xi-foldbin-critical-resonance-renyi2-floor} 的记号。则有精确恒等式
\[
\boxed{\;
\chi^2\!\bigl(p_m^{\mathrm{bin}}\|u_m\bigr)
=
F_{m+2}\left(\mathsf{Col}^{\mathrm{bin}}_m-\frac{1}{F_{m+2}}\right).\;}
\]
进一步，
\[
\boxed{\;
\liminf_{m\to\infty}\chi^2\!\bigl(p_m^{\mathrm{bin}}\|u_m\bigr)
\ge
2C_\varphi^2
\approx 0.1007267225141179.\;}
\]
等价地，
\[
\boxed{\;
\liminf_{m\to\infty}D_2\!\bigl(p_m^{\mathrm{bin}}\|u_m\bigr)
\ge
\log(1+2C_\varphi^2)
\approx 0.09597061849768761.\;}
\]
因此，bin-fold 推前分布在 \(\chi^2\) 与 Rényi-\(2\) 意义下都不可能渐近均匀化。
\end{theorem}
\begin{proof}
第一条恒等式正是命题 \ref{prop:fold-bin-chi2-col} 的内容。另一方面，由定理 \ref{thm:xi-foldbin-critical-resonance-renyi2-floor} 的点态下界，
\[
F_{m+2}\mathsf{Col}^{\mathrm{bin}}_m
\ge
1+a_m(F_m)^2+a_m(F_{m+1})^2,
\]
故
\[
\chi^2\!\bigl(p_m^{\mathrm{bin}}\|u_m\bigr)
=
F_{m+2}\mathsf{Col}^{\mathrm{bin}}_m-1
\ge
a_m(F_m)^2+a_m(F_{m+1})^2.
\]
再由定理 \ref{thm:fold-critical-resonance-constant}，
\[
a_m(F_m)\to C_\varphi,
\qquad
a_m(F_{m+1})\to C_\varphi,
\]
便得到
\[
\liminf_{m\to\infty}\chi^2\!\bigl(p_m^{\mathrm{bin}}\|u_m\bigr)\ge 2C_\varphi^2.
\]
最后，由
\[
D_2\!\bigl(p_m^{\mathrm{bin}}\|u_m\bigr)
=
\log\!\bigl(F_{m+2}\mathsf{Col}^{\mathrm{bin}}_m\bigr)
=
\log\!\Bigl(1+\chi^2\!\bigl(p_m^{\mathrm{bin}}\|u_m\bigr)\Bigr),
\]
取下极限即得 Rényi-\(2\) 版本。证毕。
\end{proof}

\leanverified{paper\_xi\_window6\_renyi\_divergence\_parity\_charge\_redundancy}
\begin{proposition}[window-\texorpdfstring{$6$}{6} 的 Rényi 散度闭式与最小完备奇偶荷签名的熵冗余]\label{prop:xi-window6-renyi-divergence-parity-charge-redundancy}
取 \(m=6\)。由推论 \ref{cor:fold-bin-m6-fiber-hist}，
\[
\#\{w\in X_6:\ d_6^{\mathrm{bin}}(w)=2\}=8,\qquad
\#\{w\in X_6:\ d_6^{\mathrm{bin}}(w)=3\}=4,\qquad
\#\{w\in X_6:\ d_6^{\mathrm{bin}}(w)=4\}=9,
\]
且 \(|X_6|=21\)。因此
\[
p_6^{\mathrm{bin}}
=
\underbrace{\tfrac{1}{32},\dots,\tfrac{1}{32}}_{8\text{ 次}},
\underbrace{\tfrac{3}{64},\dots,\tfrac{3}{64}}_{4\text{ 次}},
\underbrace{\tfrac{1}{16},\dots,\tfrac{1}{16}}_{9\text{ 次}}.
\]
并记 \(u_6(w)\equiv 1/21\) 为 \(X_6\) 上的均匀分布。
对任意 \(q>0\)、\(q\neq 1\)，Rényi 散度曲线具有闭式
\[
\boxed{\;
D_q\!\bigl(p_6^{\mathrm{bin}}\|u_6\bigr)
=
\log 21+\frac{1}{q-1}
\log\frac{8\cdot 2^q+4\cdot 3^q+9\cdot 4^q}{64^q}.\;}
\]

特别地，Shannon 熵与 KL 散度为
\[
\boxed{\;
H\!\bigl(p_6^{\mathrm{bin}}\bigr)
=
\frac{37}{8}\log 2-\frac{3}{16}\log 3
\approx 2.9998159059644762,\;}
\]
\[
\boxed{\;
D_{\mathrm{KL}}\!\bigl(p_6^{\mathrm{bin}}\|u_6\bigr)
=
\log 21-\frac{37}{8}\log 2+\frac{3}{16}\log 3
\approx 0.04470653175894679.\;}
\]

若把 Shannon 熵改以 bit 为单位写成
\[
H_{\mathrm{bit}}\!\bigl(p_6^{\mathrm{bin}}\bigr)
:=
\frac{H(p_6^{\mathrm{bin}})}{\log 2},
\]
则
\[
\boxed{\;
H_{\mathrm{bit}}\!\bigl(p_6^{\mathrm{bin}}\bigr)
=
\frac{37}{8}-\frac{3}{16}\log_2 3
\approx 4.327819531114783.\;}
\]

另一方面，由定理 \ref{thm:xi-foldbin-parity-charge-split-exact-minrank}，window-\(6\) 的最小完备奇偶荷签名秩恰为 \(21\)。因此其相对于输出 Shannon 熵的不可压缩冗余为
\[
\boxed{\;
\operatorname{Red}_6
:=
21-H_{\mathrm{bit}}\!\bigl(p_6^{\mathrm{bin}}\bigr)
=
\frac{131}{8}+\frac{3}{16}\log_2 3
\approx 16.672180468885216\ \mathrm{bits}.\;}
\]
\end{proposition}
\begin{proof}
由 window-\(6\) 直方图，直接有
\[
S_q(6)=8\cdot 2^q+4\cdot 3^q+9\cdot 4^q
\qquad(q>0).
\]
再由定理 \ref{thm:xi-foldbin-renyi-divergence-freezing-affine} 的精确公式，
\[
D_q\!\bigl(p_6^{\mathrm{bin}}\|u_6\bigr)
=
\log|X_6|+\frac{1}{q-1}\log\frac{S_q(6)}{2^{6q}},
\]
而 \(|X_6|=21\)，故得第一条闭式。

Shannon 熵直接按三类纤维质量求和：
\[
\begin{aligned}
H\!\bigl(p_6^{\mathrm{bin}}\bigr)
&=
8\cdot \frac{1}{32}\log 32
+4\cdot \frac{3}{64}\log \frac{64}{3}
+9\cdot \frac{1}{16}\log 16\\
&=
\frac{37}{8}\log 2-\frac{3}{16}\log 3.
\end{aligned}
\]
于是由定理 \ref{thm:xi-foldbin-entropy-gauge-one-nat-gap} 的
\[
D_{\mathrm{KL}}\!\bigl(p_6^{\mathrm{bin}}\|u_6\bigr)
=
\log|X_6|-H\!\bigl(p_6^{\mathrm{bin}}\bigr)
\]
即可得到 KL 散度的闭式。再除以 \(\log 2\) 即得 \(H_{\mathrm{bit}}(p_6^{\mathrm{bin}})\)。

最后，由定理 \ref{thm:xi-foldbin-parity-charge-split-exact-minrank}，最小完备奇偶荷签名秩为 \(21\)，故
\[
\operatorname{Red}_6
=
21-\frac{H(p_6^{\mathrm{bin}})}{\log 2}
=
21-\left(\frac{37}{8}-\frac{3}{16}\log_2 3\right),
\]
整理即得所示公式。证毕。
\end{proof}

\paragraph{审计偶数窗上的中心坍缩、饱和判据与熵--体积--碰撞恒等链}
本段把附录中已建立的 bin-fold 退化谱、规范群与 Stirling 展开结果，回收为六条可直接调用的终端结论：它们分别刻画审计偶数窗上的中心坍缩与 Abel 满秩分离、最大纤维组合上界的精确饱和判据、任意纤维多重度的普适 Fibonacci 上界及其末位强化、\(\kappa_m\)--KL--碰撞概率之间的恒等链、规范群体积的三因子乘法分解，以及最小退化扇区单独强制的显式碰撞下界。

\begin{theorem}[审计偶数窗上中心坍缩与 Abel 满秩的严格分离]\label{thm:xi-foldbin-audited-even-center-collapse-abel-full-rank}
令
\[
G_m^{\mathrm{bin}}
:=
\Aut(\Fold_m^{\mathrm{bin}})
\cong
\prod_{w\in X_m}S_{d_m^{\mathrm{bin}}(w)}.
\]
若
\(
m\in\{6,8,10,12\}
\)，则
\[
\boxed{\;
\bigl(G_m^{\mathrm{bin}}\bigr)^{\mathrm{ab}}
\cong
(\ZZ/2\ZZ)^{|X_m|}.\;}
\]
更强地，
\[
\boxed{\;
Z\!\bigl(G_6^{\mathrm{bin}}\bigr)\cong (\ZZ/2\ZZ)^8,\qquad
Z\!\bigl(G_m^{\mathrm{bin}}\bigr)=1\ \ (m=8,10,12).\;}
\]
并且在 \(m=6\) 时，boundary sheet-parity 规范子群
\[
Z_{\partial}\cong (\ZZ/2\ZZ)^3\hookrightarrow Z\!\bigl(G_6^{\mathrm{bin}}\bigr)
\]
只是整个中心中的一个规范三维坐标子格；事实上
\[
\boxed{\;
Z\!\bigl(G_6^{\mathrm{bin}}\bigr)/Z_{\partial}\cong (\ZZ/2\ZZ)^5.\;}
\]
\leanpartial{paper\_xi\_foldbin\_audited\_even\_center\_collapse\_abel\_full\_rank}{requested Lean signature quantifies over arbitrary Prop inputs with no hypotheses, so the conjunction is not derivable as stated}
\end{theorem}
\begin{proof}
由定理 \ref{thm:xi-audited-even-abel-center-fibonacci-splitting}，
\[
\bigl(G_m^{\mathrm{bin}}\bigr)^{\mathrm{ab}}
\cong
(\ZZ/2\ZZ)^{F_m}\oplus(\ZZ/2\ZZ)^{D_{2m-2}}.
\]
而
\(
F_m+D_{2m-2}=F_m+F_{m+1}=F_{m+2}=|X_m|
\)，故得到
\[
\bigl(G_m^{\mathrm{bin}}\bigr)^{\mathrm{ab}}\cong (\ZZ/2\ZZ)^{|X_m|}.
\]

另一方面，定理 \ref{thm:xi-foldbin-center-solvable-fiber-criterion} 给出
\[
Z\!\bigl(G_m^{\mathrm{bin}}\bigr)\cong (\ZZ/2\ZZ)^{B_m},
\qquad
B_m:=\#\{w\in X_m:\ d_m^{\mathrm{bin}}(w)=2\}.
\]
由表 \ref{tab:fold_bin_degeneracy_spectrum_m6_m12} 的直方图，
\[
m=6:\ 2:8,\,3:4,\,4:9,
\]
故 \(B_6=8\)；而
\[
m=8:\ 3:21,\,5:11,\,6:23,\qquad
m=10:\ 5:55,\,8:52,\,9:37,\qquad
m=12:\ 8:144,\,12:85,\,13:148,
\]
均无 \(d=2\) 纤维，因此 \(B_8=B_{10}=B_{12}=0\)。从而
\[
Z\!\bigl(G_6^{\mathrm{bin}}\bigr)\cong (\ZZ/2\ZZ)^8,
\qquad
Z\!\bigl(G_m^{\mathrm{bin}}\bigr)=1\ \ (m=8,10,12).
\]

最后，由定理 \ref{thm:xi-window6-central-noncentral-anomaly-splitting}，
\(
Z_{\partial}\cong (\ZZ/2\ZZ)^3
\)
是
\(
Z(G_6^{\mathrm{bin}})
\)
的规范直和因子，并且
\[
Z\!\bigl(G_6^{\mathrm{bin}}\bigr)/Z_{\partial}\cong (\ZZ/2\ZZ)^5.
\]
证毕。
\end{proof}

\begin{theorem}[bin-fold 最大纤维的精确饱和判据与最终严格次 Fibonacci 律]\label{thm:xi-foldbin-maxfiber-saturation-final-subfibonacci}
\leanverified{paper\_xi\_foldbin\_maxfiber\_saturation\_final\_subfibonacci}
令 \(K=K(m)\) 为满足
\[
F_{K+1}\le 2^m-1<F_{K+2}
\]
的唯一整数，并置 \(L:=K-m\)。
定义合法尾部集合
\[
\mathcal T_m
:=
\Bigl\{
(c_{m+1},\dots,c_K)\in\{0,1\}^{L}:\ c_kc_{k+1}=0
\Bigr\},
\]
以及最大尾和值
\[
S_{\max}(m)
:=
\max\Bigl\{
\sum_{k=m+1}^{K}c_kF_{k+1}:\ (c_{m+1},\dots,c_K)\in\mathcal T_m
\Bigr\}.
\]
则最大纤维满足
\[
\boxed{\;
d_{\max}^{\mathrm{bin}}(m)
:=
\max_{w\in X_m}d_m^{\mathrm{bin}}(w)
=
d_m^{\mathrm{bin}}(0^m)
=
\#\Bigl\{
c\in\mathcal T_m:\ \sum_{k=m+1}^{K}c_kF_{k+1}\le 2^m-1
\Bigr\}.\;}
\]
并且，记组合上界 \(F_{L+2}=|\mathcal T_m|\)，则
\[
\boxed{\;
d_{\max}^{\mathrm{bin}}(m)=F_{L+2}
\iff
S_{\max}(m)\le 2^m-1.\;}
\]
进一步，存在 \(m_0\) 使得对所有 \(m\ge m_0\) 都有
\[
\boxed{\;
d_{\max}^{\mathrm{bin}}(m)\le F_{L+2}-1.\;}
\]
\end{theorem}
\begin{proof}
由命题 \ref{prop:fold-bin-digit-dp}，对任意 \(w\in X_m\)，
\(
d_m^{\mathrm{bin}}(w)
\)
恰等于满足相邻约束、边界约束与阈值约束的尾部位串数目。
取 \(w=0^m\) 时，\(V_m(w)=0\) 且 \(w_m=0\)，因而边界约束消失，得到
\[
d_m^{\mathrm{bin}}(0^m)
=
\#\Bigl\{
c\in\mathcal T_m:\ \sum_{k=m+1}^{K}c_kF_{k+1}\le 2^m-1
\Bigr\}.
\]
再由推论 \ref{cor:fold-bin-saturation}，最大纤维确由 \(0^m\) 取得，故第一条公式成立。

同一推论又给出
\[
d_{\max}^{\mathrm{bin}}(m)\le F_{L+2}=|\mathcal T_m|,
\]
且该上界饱和当且仅当阈值约束对全部合法尾部都冗余，也即
\[
2^m-1\ge S_{\max}(m).
\]
这就证明了第二条等价判据。

最后，由注 \ref{rem:fold-bin-finite}，上述组合上界的饱和只能发生有限次。
因此存在 \(m_0\) 使得对一切 \(m\ge m_0\)，恒有
\(
d_{\max}^{\mathrm{bin}}(m)<F_{L+2}
\)；
由于两边都是整数，便得到
\[
d_{\max}^{\mathrm{bin}}(m)\le F_{L+2}-1.
\]
证毕。
\end{proof}

\begin{conjecture}[bin-fold 最大纤维组合上界的精确饱和集是有限孤立集]\label{conj:xi-foldbin-maxfiber-saturation-sporadic-set}
沿用定理 \ref{thm:xi-foldbin-maxfiber-saturation-final-subfibonacci} 的记号。生成表
\ref{tab:fold_binary_saturation_points} 对扫描窗口 \(m\le 400\) 给出的审计结果表明，组合上界
\[
d_{\max}^{\mathrm{bin}}(m)\le F_{L+2}
\]
的饱和点恰为
\[
\{1,2,3,4,5,7,12\}.
\]
据此提出猜想
\[
\boxed{\;
d_{\max}^{\mathrm{bin}}(m)=F_{L+2}
\iff
m\in\{1,2,3,4,5,7,12\}.\;}
\]
等价地，
\[
\boxed{\;
2^m-1\ge S_{\max}(m)
\iff
m\in\{1,2,3,4,5,7,12\}.\;}
\]
若此猜想成立，则定理 \ref{thm:xi-foldbin-maxfiber-saturation-final-subfibonacci} 中“饱和只能发生有限次”的结论可进一步精化为一个完全显式的有限异常集判定。
\end{conjecture}

\begin{theorem}[bin-fold 纤维多重度的普适 Fibonacci 上界与末位强化]\label{thm:xi-foldbin-universal-fibonacci-upper-bound-lastbit-improvement}
\leanverified{paper\_xi\_foldbin\_universal\_fibonacci\_upper\_bound\_lastbit\_improvement}
沿用定理 \ref{thm:xi-foldbin-maxfiber-saturation-final-subfibonacci} 的记号，记
\[
L:=K(m)-m.
\]
则对任意 \(w\in X_m\)，都有
\[
\boxed{\;
d_m^{\mathrm{bin}}(w)\le F_{L+2}=F_{K(m)-m+2}.\;}
\]
若进一步满足 \(w_m=1\)，则更强地有
\[
\boxed{\;
d_m^{\mathrm{bin}}(w)\le F_{L+1}=F_{K(m)-m+1}.\;}
\]
\end{theorem}
\begin{proof}
由命题 \ref{prop:fold-bin-digit-dp}，对任意 \(w\in X_m\)，
\(
d_m^{\mathrm{bin}}(w)
\)
恰等于满足
\[
c_kc_{k+1}=0,\qquad
w_mc_{m+1}=0,\qquad
V_m(w)+\sum_{k=m+1}^{K}c_kF_{k+1}\le 2^m-1
\]
的尾部位串 \((c_{m+1},\dots,c_K)\) 的数目。

先忽略阈值约束。长度 \(L\) 的二进制串中不含相邻 \(1\) 的总数是标准 Fibonacci 计数 \(F_{L+2}\)，故
\[
d_m^{\mathrm{bin}}(w)\le F_{L+2}.
\]

若 \(w_m=1\)，则边界约束强制 \(c_{m+1}=0\)。当 \(L=0\) 时尾部为空串，结论平凡成立；当 \(L\ge 1\) 时，剩余自由尾部是长度 \(L-1\) 的无相邻 \(1\) 串，其数目为
\[
F_{(L-1)+2}=F_{L+1}.
\]
再加入阈值约束只会进一步减少可行尾部，故
\[
d_m^{\mathrm{bin}}(w)\le F_{L+1}.
\]
证毕。
\end{proof}

\begin{theorem}[\texorpdfstring{$\kappa_m$}{kappa_m}、KL 散度与碰撞概率的精确恒等链]\label{thm:xi-foldbin-kappa-kl-collision-identity-chain}
\leanverified{paper\_xi\_foldbin\_kappa\_kl\_collision\_identity\_chain}
令
\[
p_m^{\mathrm{bin}}(w):=\frac{d_m^{\mathrm{bin}}(w)}{2^m},
\qquad
u_m(w):=\frac{1}{|X_m|}
\qquad (w\in X_m).
\]
则有严格恒等式
\[
\boxed{\;
D_{\mathrm{KL}}\!\bigl(p_m^{\mathrm{bin}}\|u_m\bigr)
=
\kappa_m-\log\!\Bigl(\frac{2^m}{|X_m|}\Bigr).\;}
\]
因而
\[
\boxed{\;
\log\!\Bigl(\frac{2^m}{|X_m|}\Bigr)
\le
\kappa_m
\le
\log\!\bigl(2^m\mathsf{Col}^{\mathrm{bin}}_m\bigr).\;}
\]
又由
\[
\chi^2\!\bigl(p_m^{\mathrm{bin}}\|u_m\bigr)
=
|X_m|\,\mathsf{Col}^{\mathrm{bin}}_m-1,
\]
可把上界改写为
\[
\kappa_m
\le
\log\!\Bigl(\frac{2^m}{|X_m|}\Bigr)
\;+\;
\log\!\Bigl(1+\chi^2\!\bigl(p_m^{\mathrm{bin}}\|u_m\bigr)\Bigr).
\]
并且两端取等都当且仅当 \(p_m^{\mathrm{bin}}=u_m\)。
\end{theorem}
\begin{proof}
第一条恒等式正是定理 \ref{thm:xi-foldbin-entropy-gauge-one-nat-gap} 的
\[
D_{\mathrm{KL}}\!\bigl(p_m^{\mathrm{bin}}\|u_m\bigr)
=
\kappa_m+\log|X_m|-m\log 2
\]
的改写。
由于 \(D_{\mathrm{KL}}\ge 0\)，立得
\[
\kappa_m\ge \log\!\Bigl(\frac{2^m}{|X_m|}\Bigr),
\]
且等号成立当且仅当 \(p_m^{\mathrm{bin}}=u_m\)。

另一方面，由定理 \ref{thm:xi-foldbin-gauge-escort-kl-decomposition}，
\[
\kappa_m
=
\log\!\bigl(2^m\mathsf{Col}^{\mathrm{bin}}_m\bigr)
-D_{\mathrm{KL}}\!\bigl(p_m^{\mathrm{bin}}\|\widetilde p_m\bigr),
\]
其中 \(\widetilde p_m(w)=(p_m^{\mathrm{bin}}(w))^2/\mathsf{Col}^{\mathrm{bin}}_m\)。
于是
\[
\kappa_m\le \log\!\bigl(2^m\mathsf{Col}^{\mathrm{bin}}_m\bigr),
\]
且等号成立当且仅当 \(p_m^{\mathrm{bin}}=\widetilde p_m\)，即所有 \(p_m^{\mathrm{bin}}(w)\) 在 \(X_m\) 上相等，从而 \(p_m^{\mathrm{bin}}=u_m\)。

最后，由命题 \ref{prop:fold-bin-chi2-col}，
\[
1+\chi^2\!\bigl(p_m^{\mathrm{bin}}\|u_m\bigr)
=
|X_m|\,\mathsf{Col}^{\mathrm{bin}}_m,
\]
故
\[
\log\!\bigl(2^m\mathsf{Col}^{\mathrm{bin}}_m\bigr)
=
\log\!\Bigl(\frac{2^m}{|X_m|}\Bigr)
\;+\;
\log\!\Bigl(1+\chi^2\!\bigl(p_m^{\mathrm{bin}}\|u_m\bigr)\Bigr),
\]
从而得到所述恒等链。证毕。
\end{proof}

\begin{theorem}[规范群体积的三因子精确乘法分解]\label{thm:xi-foldbin-gauge-volume-three-factor-product}
\leanverified{paper\_xi\_foldbin\_gauge\_volume\_three\_factor\_product}
沿用
\(
G_m^{\mathrm{bin}}:=\Aut(\Fold_m^{\mathrm{bin}})
\)
的记号。
则存在实数 \(r_m\) 满足
\[
0\le r_m\le \frac{|X_m|}{12},
\]
使得
\[
\boxed{\;
|G_m^{\mathrm{bin}}|
=
\exp\!\bigl(2^m(\kappa_m-1)\bigr)\,
(2\pi)^{|X_m|/2}\,
\Bigl(\prod_{w\in X_m}d_m^{\mathrm{bin}}(w)\Bigr)^{1/2}
e^{r_m}.\;}
\]
因此，规范群体积在主熵项之外，精确分裂为 Gaussian 半密度因子与纤维谱几何均值因子，而剩余乘法误差严格受
\(
\exp(|X_m|/12)
\)
控制。
\end{theorem}
\begin{proof}
由定理 \ref{thm:fold-bin-gauge-volume-stirling-second-order}，
\[
g_m
=
\kappa_m-1+\frac{1}{2^{m+1}}
\Bigl(
|X_m|\log(2\pi)+\sum_{w\in X_m}\log d_m^{\mathrm{bin}}(w)
\Bigr)+R_m,
\]
其中
\[
0\le R_m\le \frac{|X_m|}{12\cdot 2^m}.
\]
两边同乘 \(2^m\)，并记
\(
r_m:=2^mR_m
\)，则
\[
\log|G_m^{\mathrm{bin}}|
=
2^m(\kappa_m-1)
+\frac12\Bigl(
|X_m|\log(2\pi)+\sum_{w\in X_m}\log d_m^{\mathrm{bin}}(w)
\Bigr)
+r_m,
\]
且 \(0\le r_m\le |X_m|/12\)。
对该恒等式指数化即得所示乘法分解。证毕。
\end{proof}

\begin{theorem}[审计偶数窗上最小退化扇区强制的显式碰撞下界]\label{thm:xi-foldbin-audited-even-minsector-collision-explicit-floor}
若
\(
m\in\{6,8,10,12\}
\)，则
\[
\boxed{\;
\mathsf{Col}^{\mathrm{bin}}_m
\ge
\frac{1}{2^{2m}}
\left(
F_mF_{m/2}^{\,2}
+
\frac{\bigl(2^m-F_mF_{m/2}\bigr)^2}{F_{m+1}}
\right).\;}
\]
等价地，
\[
\boxed{\;
\chi^2\!\bigl(p_m^{\mathrm{bin}}\|u_m\bigr)
\ge
\frac{F_{m+2}}{2^{2m}}
\left(
F_mF_{m/2}^{\,2}
+
\frac{\bigl(2^m-F_mF_{m/2}\bigr)^2}{F_{m+1}}
\right)-1.\;}
\]
并且等号成立当且仅当所有非最小退化纤维具有同一尺寸。
\end{theorem}
\begin{proof}
记
\[
n:=|X_m|=F_{m+2},
\qquad
r:=\bigl|\mathcal S_{m,d_{\min}(m)}\bigr|=F_m,
\qquad
a:=d_{\min}(m)=F_{m/2},
\]
其中第一、第二、第三个等式分别来自稳定类型计数律、定理 \ref{thm:fold-bin-min-degeneracy-fib-audited} 与同一定理。
再由推论 \ref{cor:fold-bin-minsector-complement-maxfiber-audited}，
\[
n-r=F_{m+1}.
\]

把全部纤维大小写为
\[
\underbrace{a,\dots,a}_{r\text{ 次}},\ e_1,\dots,e_{n-r},
\]
则
\[
\sum_{i=1}^{n-r}e_i=2^m-ra=2^m-F_mF_{m/2}.
\]
于是
\[
\sum_{w\in X_m}\bigl(d_m^{\mathrm{bin}}(w)\bigr)^2
=
ra^2+\sum_{i=1}^{n-r}e_i^2.
\]
对后者应用 Cauchy--Schwarz 不等式，
\[
\sum_{i=1}^{n-r}e_i^2
\ge
\frac{\bigl(\sum_{i=1}^{n-r}e_i\bigr)^2}{n-r}
=
\frac{\bigl(2^m-F_mF_{m/2}\bigr)^2}{F_{m+1}}.
\]
因此
\[
\sum_{w\in X_m}\bigl(d_m^{\mathrm{bin}}(w)\bigr)^2
\ge
F_mF_{m/2}^{\,2}
+
\frac{\bigl(2^m-F_mF_{m/2}\bigr)^2}{F_{m+1}}.
\]
再由
\(
\mathsf{Col}^{\mathrm{bin}}_m
=
2^{-2m}\sum_w(d_m^{\mathrm{bin}}(w))^2
\)
得到碰撞下界。
最后由命题 \ref{prop:fold-bin-chi2-col} 的恒等式
\(
\chi^2(p_m^{\mathrm{bin}}\|u_m)=|X_m|\mathsf{Col}^{\mathrm{bin}}_m-1
\)
即得第二条公式。

Cauchy--Schwarz 取等当且仅当
\(
e_1=\cdots=e_{n-r}
\)，也就是全部非最小退化纤维等大。证毕。
\end{proof}

\begin{theorem}[已审计偶数窗存在完全线性的 Fibonacci 低预算相]\label{thm:xi-foldbin-audited-even-fibonacci-linear-low-budget-phase}
若
\[
m\in\{6,8,10,12\}
\]
且整数预算 \(B\ge 0\) 满足
\[
2^B\le F_{m/2},
\]
则每一条 bin-fold 纤维都超过预算阈值，因而
\[
\boxed{\;
C_m^{\mathrm{bin}}(B)=2^B\,F_{m+2},\qquad
\mathrm{Succ}_m^{\mathrm{bin}}(B)=2^{B-m}F_{m+2}.\;}
\]
换言之，在已审计偶数窗上存在一段完全线性的低预算相；其斜率与截距仅由
\(
F_{m/2}
\)
与
\(
F_{m+2}
\)
两条 Fibonacci 尺度共同控制，而与更高退化层的细节无关。
\end{theorem}
\begin{proof}
由定理 \ref{thm:fold-bin-min-degeneracy-fib-audited}，当 \(m\in\{6,8,10,12\}\) 时，对每个 \(w\in X_m\) 都有
\[
d_m^{\mathrm{bin}}(w)\ge d_{\min}(m)=F_{m/2}.
\]
若 \(2^B\le F_{m/2}\)，则对所有 \(w\in X_m\) 皆有
\[
\min\bigl(d_m^{\mathrm{bin}}(w),2^B\bigr)=2^B.
\]
因此
\[
C_m^{\mathrm{bin}}(B)
=
\sum_{w\in X_m}\min\bigl(d_m^{\mathrm{bin}}(w),2^B\bigr)
=
\sum_{w\in X_m}2^B
=
2^B|X_m|.
\]
再由 \(|X_m|=F_{m+2}\)，便得到
\[
C_m^{\mathrm{bin}}(B)=2^B F_{m+2}.
\]
最后两边除以 \(2^m\)，即得
\[
\mathrm{Succ}_m^{\mathrm{bin}}(B)=2^{B-m}F_{m+2}.
\]
证毕。
\end{proof}

\paragraph{审计偶数窗的奇偶荷截面、可解阴影与纯奇偶荷阶段}
在审计偶数窗上的中心坍缩、Abel 满秩、规范体积与低预算相已经确定之后，还可把
\(
G_m^{\mathrm{bin}}
\)
的可解阴影层进一步压缩为下列几条直接结论。它们分别给出阿贝尔化的一个规范奇偶荷截面、最大可解商的显式相图、中心坍缩与可解非交换坍缩的不同步、纯奇偶荷阶段的半单可解残核、任意可解阴影相对完整规范复杂度的指数级稀薄性，以及
\(
\kappa_m
\)
对规范体积的双侧压缩。

\begin{theorem}[奇偶荷截面的中心交公式]\label{thm:xi-foldbin-even-window-parity-section-center-intersection}
对每个审计偶数窗
\[
m\in\{6,8,10,12\},
\]
沿用定理 \ref{thm:xi-foldbin-audited-even-center-collapse-abel-full-rank} 的记号，并在每个因子
\(
\mathfrak S_{d_m^{\mathrm{bin}}(w)}
\)
中任选一条换位
\(
\tau_w
\)。
定义
\[
E_m:=\bigl\langle \tau_w:\ w\in X_m\bigr\rangle\subset G_m^{\mathrm{bin}}.
\]
则有
\[
E_m\cong (\ZZ/2\ZZ)^{|X_m|},
\]
并且复合映射
\[
E_m\hookrightarrow G_m^{\mathrm{bin}}\twoheadrightarrow
\bigl(G_m^{\mathrm{bin}}\bigr)^{\mathrm{ab}}
\]
是同构。此外，
\[
E_m\cap Z\!\bigl(G_m^{\mathrm{bin}}\bigr)\cong (\ZZ/2\ZZ)^{B_m},
\qquad
B_m:=\#\{w\in X_m:\ d_m^{\mathrm{bin}}(w)=2\}.
\]
特别地，
\[
B_6=8,
\qquad
B_8=B_{10}=B_{12}=0.
\]
\end{theorem}
\begin{proof}
不同的 \(w\) 对应于直积分解
\[
G_m^{\mathrm{bin}}
\cong
\prod_{w\in X_m}\mathfrak S_{d_m^{\mathrm{bin}}(w)}
\]
中的不同坐标因子，因此所选换位两两可交换，且每个 \(\tau_w\) 都满足 \(\tau_w^2=1\)。于是
\[
E_m\cong (\ZZ/2\ZZ)^{|X_m|}.
\]

又由定理 \ref{thm:fold-bin-min-degeneracy-fib-audited}，审计偶数窗上每条 bin-fold 纤维都满足
\(
d_m^{\mathrm{bin}}(w)\ge d_{\min}(m)=F_{m/2}\ge 2
\)。
因此每个因子的符号映射都给出满同态
\[
\operatorname{sgn}_w:\mathfrak S_{d_m^{\mathrm{bin}}(w)}\twoheadrightarrow \ZZ/2\ZZ.
\]
把这些符号映射相乘，得到
\[
\prod_{w\in X_m}\operatorname{sgn}_w:
G_m^{\mathrm{bin}}
\twoheadrightarrow
(\ZZ/2\ZZ)^{|X_m|}.
\]
而每个 \(\tau_w\) 在第 \(w\) 个坐标上给出唯一的奇元，在其余坐标上为单位元，所以该映射限制到 \(E_m\) 后恰是标准坐标同构。再结合定理 \ref{thm:xi-foldbin-audited-even-center-collapse-abel-full-rank} 的
\[
\bigl(G_m^{\mathrm{bin}}\bigr)^{\mathrm{ab}}
\cong
(\ZZ/2\ZZ)^{|X_m|},
\]
便知复合映射
\(
E_m\hookrightarrow G_m^{\mathrm{bin}}\twoheadrightarrow (G_m^{\mathrm{bin}})^{\mathrm{ab}}
\)
为同构。

最后，由定理 \ref{thm:xi-foldbin-audited-even-center-collapse-abel-full-rank}，
\[
Z\!\bigl(G_m^{\mathrm{bin}}\bigr)\cong (\ZZ/2\ZZ)^{B_m}.
\]
直积群的中心等于各因子中心的直积，而
\[
Z(\mathfrak S_d)=
\begin{cases}
\mathfrak S_2,& d=2,\\
1,& d\ge 3.
\end{cases}
\]
故 \(\tau_w\) 落在中心中当且仅当 \(d_m^{\mathrm{bin}}(w)=2\)。于是
\[
E_m\cap Z\!\bigl(G_m^{\mathrm{bin}}\bigr)
=
\bigl\langle \tau_w:\ d_m^{\mathrm{bin}}(w)=2\bigr\rangle
\cong
(\ZZ/2\ZZ)^{B_m}.
\]
最后再由表 \ref{tab:fold_bin_degeneracy_spectrum_m6_m12} 的审计直方图读出
\(
B_6=8
\)
与
\(
B_8=B_{10}=B_{12}=0
\)。
证毕。
\end{proof}
\leanverified{paper\_xi\_foldbin\_even\_window\_parity\_section\_center\_intersection}

\begin{theorem}[审计偶数窗最大可解商的显式相图]\label{thm:xi-foldbin-audited-even-maximal-solvable-quotient}
对
\[
m\in\{6,8,10,12\},
\]
记
\(
Q_m^{\mathrm{sol}}
\)
为
\(
G_m^{\mathrm{bin}}
\)
的最大可解商。定义
\[
q(d):=
\begin{cases}
\mathfrak S_d,& d\le 4,\\
\ZZ/2\ZZ,& d\ge 5.
\end{cases}
\]
则有自然同构
\[
Q_m^{\mathrm{sol}}
\cong
\prod_{w\in X_m}q\!\bigl(d_m^{\mathrm{bin}}(w)\bigr).
\]
特别地，
\[
Q_6^{\mathrm{sol}}
\cong
\mathfrak S_2^{\,8}\times \mathfrak S_3^{\,4}\times \mathfrak S_4^{\,9},
\]
\[
Q_8^{\mathrm{sol}}
\cong
\mathfrak S_3^{\,21}\times (\ZZ/2\ZZ)^{34},
\]
\[
Q_{10}^{\mathrm{sol}}
\cong
(\ZZ/2\ZZ)^{144},
\qquad
Q_{12}^{\mathrm{sol}}
\cong
(\ZZ/2\ZZ)^{377}.
\]
\end{theorem}
\begin{proof}
对称群的标准结构表明：
\[
\mathfrak S_d\ \text{可解}
\iff
d\le 4,
\]
而当 \(d\ge 5\) 时，
\(
A_d
\)
是非阿贝尔单群，因此
\[
\mathfrak S_d/A_d\cong \ZZ/2\ZZ
\]
是 \(\mathfrak S_d\) 的最大可解商。故上式中定义的 \(q(d)\) 正是因子 \(\mathfrak S_d\) 的最大可解商。

现在令
\[
\phi:
G_m^{\mathrm{bin}}
\cong
\prod_{w\in X_m}\mathfrak S_{d_m^{\mathrm{bin}}(w)}
\longrightarrow H
\]
为任意到可解群 \(H\) 的同态。则其在每个坐标因子上的限制
\[
\phi_w:\mathfrak S_{d_m^{\mathrm{bin}}(w)}\to H
\]
具有可解像，因此都必经由 \(q(d_m^{\mathrm{bin}}(w))\)。由直积的泛性质，\(\phi\) 必经由
\[
\prod_{w\in X_m}q\!\bigl(d_m^{\mathrm{bin}}(w)\bigr).
\]
这就证明了该直积正是
\(
G_m^{\mathrm{bin}}
\)
的最大可解商。

最后只需把表 \ref{tab:fold_bin_degeneracy_spectrum_m6_m12} 的审计直方图代入。对
\(
m=6
\)
有
\[
2:8,\qquad 3:4,\qquad 4:9,
\]
故
\[
Q_6^{\mathrm{sol}}
\cong
\mathfrak S_2^{\,8}\times \mathfrak S_3^{\,4}\times \mathfrak S_4^{\,9}.
\]
对
\(
m=8
\)
有
\[
3:21,\qquad 5:11,\qquad 6:23,
\]
故
\[
Q_8^{\mathrm{sol}}
\cong
\mathfrak S_3^{\,21}\times (\ZZ/2\ZZ)^{11+23}
=
\mathfrak S_3^{\,21}\times (\ZZ/2\ZZ)^{34}.
\]
对
\(
m=10
\)
有
\[
5:55,\qquad 8:52,\qquad 9:37,
\]
全部因子都贡献符号商，所以
\[
Q_{10}^{\mathrm{sol}}\cong (\ZZ/2\ZZ)^{55+52+37}=(\ZZ/2\ZZ)^{144}.
\]
同理，对
\(
m=12
\)
的直方图
\[
8:144,\qquad 12:85,\qquad 13:148
\]
可得
\[
Q_{12}^{\mathrm{sol}}\cong (\ZZ/2\ZZ)^{144+85+148}=(\ZZ/2\ZZ)^{377}.
\]
证毕。
\end{proof}
\leanverified{paper\_xi\_foldbin\_audited\_even\_maximal\_solvable\_quotient}

\begin{corollary}[中心坍缩与可解非交换坍缩不同步]\label{cor:xi-foldbin-center-solvable-collapse-two-stage}
在审计偶数窗上，中心坍缩与可解非交换坍缩是两件不同层级的事件。
更具体地，
\[
Z\!\bigl(G_m^{\mathrm{bin}}\bigr)=1
\qquad (m=8,10,12),
\]
但
\[
Q_8^{\mathrm{sol}}
\cong
\mathfrak S_3^{\,21}\times (\ZZ/2\ZZ)^{34}
\]
仍然保留非阿贝尔可解因子，而
\[
Q_{10}^{\mathrm{sol}}
\cong
(\ZZ/2\ZZ)^{144},
\qquad
Q_{12}^{\mathrm{sol}}
\cong
(\ZZ/2\ZZ)^{377}
\]
才首次退化为纯奇偶荷格。
因此，在已审计偶数窗中，“中心已死而可解非交换尚存”的现象恰发生在 \(m=8\)。
\end{corollary}
\begin{proof}
由定理 \ref{thm:xi-foldbin-audited-even-center-collapse-abel-full-rank}，
\[
Z\!\bigl(G_m^{\mathrm{bin}}\bigr)=1
\qquad (m=8,10,12).
\]
另一方面，定理 \ref{thm:xi-foldbin-audited-even-maximal-solvable-quotient} 给出
\[
Q_8^{\mathrm{sol}}
\cong
\mathfrak S_3^{\,21}\times (\ZZ/2\ZZ)^{34},
\]
故在 \(m=8\) 时最大可解商仍含非阿贝尔可解因子；而在
\(
m=10,12
\)
时，最大可解商都已坍缩为纯二元奇偶荷格。证毕。
\end{proof}
\leanverified{paper\_xi\_foldbin\_center\_solvable\_collapse\_two\_stage}

\leanverified{paper\_xi\_foldbin\_pure\_parity\_solvable\_residual\_kernel}
\begin{theorem}[纯奇偶荷阶段的半单可解残核]\label{thm:xi-foldbin-pure-parity-solvable-residual-kernel}
设分辨率 \(m\) 满足
\[
d_m^{\mathrm{bin}}(w)\ge 5
\qquad (\forall\, w\in X_m).
\]
定义正规子群
\[
N_m:=\prod_{w\in X_m}A_{d_m^{\mathrm{bin}}(w)}
\triangleleft
G_m^{\mathrm{bin}}.
\]
则有短正合列
\[
1\longrightarrow N_m
\longrightarrow G_m^{\mathrm{bin}}
\longrightarrow (\ZZ/2\ZZ)^{|X_m|}
\longrightarrow 1.
\]
并且：
\begin{enumerate}
\item \(N_m\) 是 \(G_m^{\mathrm{bin}}\) 的可解残核，即所有到可解群同态之核的交；
\item \(N_m\) 是 \(|X_m|\) 个非阿贝尔单群的直积；
\item 任意到可解群的同态都必经由奇偶荷格
\(
(\ZZ/2\ZZ)^{|X_m|}
\)。
\end{enumerate}
特别地，表 \ref{tab:fold_bin_degeneracy_spectrum_m6_m12} 表明该结论适用于
\[
m=10,\qquad m=12.
\]
\end{theorem}
\begin{proof}
在假设
\(
d_m^{\mathrm{bin}}(w)\ge 5
\)
下，每个交错群
\(
A_{d_m^{\mathrm{bin}}(w)}
\)
都是非阿贝尔单群，并满足
\[
\mathfrak S_{d_m^{\mathrm{bin}}(w)}/A_{d_m^{\mathrm{bin}}(w)}
\cong
\ZZ/2\ZZ.
\]
因此
\[
G_m^{\mathrm{bin}}/N_m
\cong
\prod_{w\in X_m}
\mathfrak S_{d_m^{\mathrm{bin}}(w)}/A_{d_m^{\mathrm{bin}}(w)}
\cong
(\ZZ/2\ZZ)^{|X_m|},
\]
从而得到所示短正合列。

再令
\(
\phi:G_m^{\mathrm{bin}}\to H
\)
为任意到可解群 \(H\) 的同态。其在每个坐标因子上的限制
\[
\phi_w:\mathfrak S_{d_m^{\mathrm{bin}}(w)}\to H
\]
具有可解像，而
\(
A_{d_m^{\mathrm{bin}}(w)}
\)
是 \(\mathfrak S_{d_m^{\mathrm{bin}}(w)}\) 的唯一非平凡非阿贝尔单正规子群，因此必落在 \(\ker\phi_w\) 中。对全部 \(w\) 同时成立，于是
\[
N_m\subseteq \ker\phi.
\]
这说明任意到可解群的同态都经由
\(
G_m^{\mathrm{bin}}/N_m\cong (\ZZ/2\ZZ)^{|X_m|}
\)，故 \(N_m\) 正是可解残核。

第二条由 \(N_m\) 的定义与每个 \(A_{d_m^{\mathrm{bin}}(w)}\) 的非阿贝尔单性立即得出；第三条只是第一条在商层的重述。最后，由表 \ref{tab:fold_bin_degeneracy_spectrum_m6_m12} 可见 \(m=10,12\) 时所有纤维大小都至少为 \(5\)，故这两个审计偶数窗均处于纯奇偶荷阶段。证毕。
\end{proof}

\begin{theorem}[任意可解阴影相对于规范复杂度的指数稀薄律]\label{thm:xi-foldbin-solvable-shadow-exponential-thinning}
设分辨率 \(m\) 满足
\[
d_m^{\mathrm{bin}}(w)\ge 5
\qquad (\forall\, w\in X_m),
\]
并令 \(Q\) 为 \(G_m^{\mathrm{bin}}\) 的任意非平凡可解商。则有
\[
\log |Q|
\le
|X_m|\log 2
=
F_{m+2}\log 2.
\]
另一方面，
\[
\log\bigl|G_m^{\mathrm{bin}}\bigr|
=
2^m g_m
=
2^m
\left(
m\log\!\frac{2}{\varphi}
-1-\frac{\log\varphi}{1+\varphi^2}
+O\!\left(m\left(\frac{\varphi}{2}\right)^m\right)
\right).
\]
因此
\[
\frac{\log|G_m^{\mathrm{bin}}|}{\log|Q|}
\ge
\frac{2^m}{F_{m+2}\log 2}
\left(
m\log\!\frac{2}{\varphi}
-1-\frac{\log\varphi}{1+\varphi^2}
+O\!\left(m\left(\frac{\varphi}{2}\right)^m\right)
\right).
\]
沿任意一条满足上述假设的共尾子列，进一步有
\[
\frac{\log|G_m^{\mathrm{bin}}|}{\log|Q|}
\ge
\frac{\sqrt5}{\varphi^2\log 2}
\left(
m\log\!\frac{2}{\varphi}
-1-\frac{\log\varphi}{1+\varphi^2}
+o(m)
\right)
\left(\frac{2}{\varphi}\right)^m.
\]
\end{theorem}
\begin{proof}
由定理 \ref{thm:xi-foldbin-pure-parity-solvable-residual-kernel}，任意可解商都必经由奇偶荷格
\[
(\ZZ/2\ZZ)^{|X_m|}.
\]
因此
\[
|Q|\le 2^{|X_m|},
\]
也就是
\[
\log |Q|\le |X_m|\log 2.
\]
再由稳定类型计数律 \(|X_m|=F_{m+2}\)，得到第一条公式。

另一方面，定理 \ref{thm:xi-foldbin-gauge-density-exact-thermodynamic-constant} 给出
\[
g_m
=
m\log\!\Bigl(\frac{2}{\varphi}\Bigr)
-1-\frac{\log\varphi}{1+\varphi^2}
+O\!\Bigl(m\Bigl(\frac{\varphi}{2}\Bigr)^m\Bigr).
\]
乘以 \(2^m\) 便得到
\[
\log|G_m^{\mathrm{bin}}|=2^m g_m
\]
的所示展开。把这两个估计合并，便得到第二条不等式。

最后，沿任意满足假设的共尾子列，Binet 公式给出
\[
F_{m+2}
=
\frac{\varphi^{m+2}}{\sqrt5}+o(\varphi^m),
\]
因此
\[
\frac{2^m}{F_{m+2}\log 2}
=
\frac{\sqrt5}{\varphi^2\log 2}
\left(\frac{2}{\varphi}\right)^m
\bigl(1+o(1)\bigr).
\]
与前式相乘后，主项来自
\(
m\log(2/\varphi)
\)，而全部误差可统一吸收到括号中的 \(o(m)\) 中，故得到最后一式。证毕。
\end{proof}

因此，一旦进入纯奇偶荷阶段，任何可解证书所能承载的全部信息量都至多线性于
\(
F_{m+2}
\)，而完整规范复杂度的主尺度却是
\(
2^m\log(2/\varphi)
\)；
可解阴影相对于完整 gauge 体积因此只是指数级稀薄的投影。

\input{subsubsec__xi-time-protocol-conclusions-part60ag-binfold-kappa-compression-pureparity}

\endinput


\endinput

\paragraph{规范群熵账本、有限群代数障碍与偶 \texorpdfstring{$\zeta$}{zeta} 反演}
在规范群体积的二阶 Stirling 展开、审计偶数窗的最小退化律、约数触发零余类并的碰撞削减机制以及 fold-gauge 常数的 Stirling--Bernoulli 层级已经建立之后，还可进一步压缩出下列五条直接结论。

\begin{theorem}[规范群熵的精确账本恒等式]\label{thm:xi-time-part60aa-gauge-entropy-ledger-exact-identity}
沿用定理 \ref{thm:fold-bin-gauge-volume-stirling-second-order} 的记号，定义
\[
p_m^{\mathrm{bin}}(w):=\frac{d_m^{\mathrm{bin}}(w)}{2^m}
\qquad (w\in X_m),
\]
并记其 Shannon 熵为
\[
H\!\bigl(p_m^{\mathrm{bin}}\bigr):=
-\sum_{w\in X_m}p_m^{\mathrm{bin}}(w)\log p_m^{\mathrm{bin}}(w).
\]
则规范群体积密度 \(g_m\) 满足严格恒等式
\[
\boxed{\;
g_m
=
m\log 2
-H\!\bigl(p_m^{\mathrm{bin}}\bigr)
-1
+\frac{|X_m|\log(2\pi)+\sum_{w\in X_m}\log d_m^{\mathrm{bin}}(w)}{2^{m+1}}
+R_m,\;}
\]
其中余项满足
\[
\boxed{\;
0\le R_m\le \frac{|X_m|}{12\cdot 2^m}.\;}
\]
特别地，
\[
\boxed{\;
g_m
=
m\log 2
-H\!\bigl(p_m^{\mathrm{bin}}\bigr)
-1
+O\!\Bigl(\frac{|X_m|\,m}{2^m}\Bigr),\;}
\]
并且
\[
\boxed{\;
\kappa_m=m\log 2-H\!\bigl(p_m^{\mathrm{bin}}\bigr).\;}
\]
因此，\(\kappa_m\) 恰是 \(g_m\) 的零阶主项。
\end{theorem}
\begin{proof}
由定理 \ref{thm:fold-bin-gauge-volume-stirling-second-order}，
\[
g_m
=
\kappa_m-1+\frac{|X_m|\log(2\pi)+\sum_{w\in X_m}\log d_m^{\mathrm{bin}}(w)}{2^{m+1}}
+R_m,
\]
且
\[
0\le R_m\le \frac{|X_m|}{12\cdot 2^m}.
\]

另一方面，由
\[
p_m^{\mathrm{bin}}(w)=\frac{d_m^{\mathrm{bin}}(w)}{2^m},
\qquad
\sum_{w\in X_m}d_m^{\mathrm{bin}}(w)=2^m,
\]
可直接计算
\[
\begin{aligned}
H\!\bigl(p_m^{\mathrm{bin}}\bigr)
&=
-\sum_{w\in X_m}\frac{d_m^{\mathrm{bin}}(w)}{2^m}
\log\frac{d_m^{\mathrm{bin}}(w)}{2^m}\\
&=
-\frac{1}{2^m}\sum_{w\in X_m}d_m^{\mathrm{bin}}(w)\log d_m^{\mathrm{bin}}(w)
+\frac{m\log 2}{2^m}\sum_{w\in X_m}d_m^{\mathrm{bin}}(w)\\
&=
-\kappa_m+m\log 2.
\end{aligned}
\]
即
\[
\kappa_m=m\log 2-H\!\bigl(p_m^{\mathrm{bin}}\bigr).
\]
把该式代回 \(g_m\) 的 Stirling 展开，即得第一条盒装恒等式。

最后，由 \(1\le d_m^{\mathrm{bin}}(w)\le 2^m\) 可知
\[
0\le \sum_{w\in X_m}\log d_m^{\mathrm{bin}}(w)\le |X_m|\,m\log 2,
\]
再结合余项界，便得到
\[
g_m
=
m\log 2-H\!\bigl(p_m^{\mathrm{bin}}\bigr)-1
+O\!\Bigl(\frac{|X_m|\,m}{2^m}\Bigr).
\]
证毕。
\end{proof}

\begin{theorem}[连续容量曲线的对数 Stieltjes 表示与规范体积密度的普适常数差]\label{thm:xi-time-part60aa-log-stieltjes-capacity-gauge-density}
沿用定理 \ref{thm:xi-time-part60aa-gauge-entropy-ledger-exact-identity} 的记号，并定义连续容量曲线
\[
\mathcal C_m^{\mathrm{cont}}(T):=\sum_{w\in X_m}\min\bigl(d_m^{\mathrm{bin}}(w),T\bigr)
\qquad(T\ge 0).
\]
则有精确 Stieltjes 表示
\[
\boxed{\;
\kappa_m
=
1+\frac{1}{2^m}\int_0^\infty \log T\,\dd \mathcal C_m^{\mathrm{cont}}(T).\;}
\]
进一步，规范体积密度满足精确恒等式
\[
\boxed{\;
g_m
=
\frac{1}{2^m}\int_0^\infty \log T\,\dd \mathcal C_m^{\mathrm{cont}}(T)
+\frac{|X_m|\log(2\pi)+\sum_{w\in X_m}\log d_m^{\mathrm{bin}}(w)}{2^{m+1}}
+R_m,\;}
\]
其中
\[
0\le R_m\le \frac{|X_m|}{12\cdot 2^m}.
\]
因此
\[
\boxed{\;
g_m
=
\frac{1}{2^m}\int_0^\infty \log T\,\dd \mathcal C_m^{\mathrm{cont}}(T)
+O\!\Bigl(m\Bigl(\frac{\varphi}{2}\Bigr)^m\Bigr),\;}
\]
并且
\[
\boxed{\;
\kappa_m-g_m
=
1+O\!\Bigl(m\Bigl(\frac{\varphi}{2}\Bigr)^m\Bigr).\;}
\]
换言之，\(\kappa_m\) 恰是连续容量曲线的对数 Stieltjes 积分再加上一枚普适常数 \(1\)，而 \(g_m\) 与该积分之间只差指数小校正。
\end{theorem}
\begin{proof}
对每个固定 \(w\in X_m\)，函数
\[
T\longmapsto \min\bigl(d_m^{\mathrm{bin}}(w),T\bigr)
\]
是有界变差函数，其 Stieltjes 测度正是区间 \([0,d_m^{\mathrm{bin}}(w)]\) 上的 Lebesgue 测度。故
\[
\int_0^\infty \log T\,\dd\min\bigl(d_m^{\mathrm{bin}}(w),T\bigr)
=
\int_0^{d_m^{\mathrm{bin}}(w)}\log T\,\dd T
=
d_m^{\mathrm{bin}}(w)\log d_m^{\mathrm{bin}}(w)-d_m^{\mathrm{bin}}(w).
\]
对全部 \(w\in X_m\) 求和，得到
\[
\int_0^\infty \log T\,\dd\mathcal C_m^{\mathrm{cont}}(T)
=
\sum_{w\in X_m}
\Bigl(
d_m^{\mathrm{bin}}(w)\log d_m^{\mathrm{bin}}(w)-d_m^{\mathrm{bin}}(w)
\Bigr).
\]
由于
\[
\sum_{w\in X_m}d_m^{\mathrm{bin}}(w)=2^m,
\]
而
\[
\kappa_m=\frac{1}{2^m}\sum_{w\in X_m}d_m^{\mathrm{bin}}(w)\log d_m^{\mathrm{bin}}(w),
\]
遂有
\[
\frac{1}{2^m}\int_0^\infty \log T\,\dd\mathcal C_m^{\mathrm{cont}}(T)=\kappa_m-1,
\]
即第一条盒装公式。

再将
\[
\kappa_m-1=\frac{1}{2^m}\int_0^\infty \log T\,\dd\mathcal C_m^{\mathrm{cont}}(T)
\]
代回定理 \ref{thm:xi-time-part60aa-gauge-entropy-ledger-exact-identity} 的精确账本恒等式，便得到 \(g_m\) 的精确表达。

最后，由 \(|X_m|=F_{m+2}\) 且
\[
0\le \sum_{w\in X_m}\log d_m^{\mathrm{bin}}(w)\le |X_m|\,m\log 2,
\]
可把显示校正项与 \(R_m\) 一并吸收到
\[
O\!\Bigl(\frac{|X_m|\,m}{2^m}\Bigr)
=
O\!\Bigl(m\Bigl(\frac{\varphi}{2}\Bigr)^m\Bigr)
\]
中，从而得到两条渐近式。证毕。
\end{proof}

\begin{theorem}[审计偶数窗上的纤维群胚代数不可能是有限群代数]\label{thm:xi-time-part60aa-audited-even-groupoid-no-finite-group-algebra}
对审计偶数窗
\[
m\in\{6,8,10,12\},
\]
纤维群胚代数
\[
\CC[\mathcal R_m^{\mathrm{bin}}]
\cong
\bigoplus_{w\in X_m}M_{d_m^{\mathrm{bin}}(w)}(\CC)
\]
不可能与任何有限群 \(H\) 的复群代数 \(\CC[H]\) 同构。更强地，不存在任何非零代数同态
\[
\chi:\CC[\mathcal R_m^{\mathrm{bin}}]\to\CC.
\]
\end{theorem}
\begin{proof}
由命题 \ref{prop:fold-bin-groupoid-algebra-wedderburn}，
\[
\CC[\mathcal R_m^{\mathrm{bin}}]
\cong
\bigoplus_{w\in X_m}M_{d_m^{\mathrm{bin}}(w)}(\CC).
\]
而定理 \ref{thm:fold-bin-min-degeneracy-fib-audited} 在审计偶数窗上给出
\[
d_m^{\mathrm{bin}}(w)\ge d_{\min}(m)=F_{m/2}\ge 2
\qquad(\forall\,w\in X_m).
\]
因此上述 Wedderburn 分解中不存在任何 \(M_1(\CC)\) 块。

现证不存在非零代数同态 \(\chi:\CC[\mathcal R_m^{\mathrm{bin}}]\to\CC\)。若此类 \(\chi\) 存在，则其在某个直和因子
\[
M_{d_m^{\mathrm{bin}}(w)}(\CC)
\]
上的限制必非零。由于矩阵代数 \(M_n(\CC)\) 为简单代数，非零代数同态的核只能为零，故该限制必为单射。但当 \(n\ge 2\) 时，\(M_n(\CC)\) 非交换，而其像落在交换代数 \(\CC\) 中，矛盾。故不存在非零代数同态。

若反设存在有限群 \(H\) 使
\[
\CC[\mathcal R_m^{\mathrm{bin}}]\cong \CC[H],
\]
则群代数上的增广映射
\[
\varepsilon:\CC[H]\to\CC,
\qquad
\varepsilon\!\left(\sum_{h\in H}a_h h\right):=\sum_{h\in H}a_h
\]
给出一个非零代数同态。经由上述同构传回 \(\CC[\mathcal R_m^{\mathrm{bin}}]\)，便得到矛盾。故
\[
\CC[\mathcal R_m^{\mathrm{bin}}]\not\cong \CC[H]
\]
对一切有限群 \(H\) 都成立。证毕。
\end{proof}

\begin{theorem}[约数触发零余类并给出的规范群熵显式上界]\label{thm:xi-time-part60aa-divisor-zero-class-gauge-entropy-upper-bound}
记
\[
M:=F_{m+2}.
\]
若存在真约数 \(d\mid (m+2)\)、\(d<m+2\)，满足
\[
F_d\mid \frac{M}{2},
\]
则规范群体积密度 \(g_m\) 满足严格上界
\[
\boxed{\;
g_m
\le
m\log 2
-1
+\log\!\Bigl(1-\frac{F_d}{M}\Bigr)
+\frac{|X_m|\log(2\pi)+\sum_{w\in X_m}\log d_m^{\mathrm{bin}}(w)}{2^{m+1}}
+\frac{|X_m|}{12\cdot 2^m}.\;}
\]
特别地，
\[
\boxed{\;
g_m
\le
m\log 2
-1
+\log\!\Bigl(1-\frac{F_d}{M}\Bigr)
+O\!\Bigl(\frac{|X_m|\,m}{2^m}\Bigr).\;}
\]
\end{theorem}
\begin{proof}
由定理 \ref{thm:xi-time-part60aa-gauge-entropy-ledger-exact-identity}，
\[
g_m
=
m\log 2
-H\!\bigl(p_m^{\mathrm{bin}}\bigr)
-1
+\frac{|X_m|\log(2\pi)+\sum_{w\in X_m}\log d_m^{\mathrm{bin}}(w)}{2^{m+1}}
+R_m,
\]
其中
\[
0\le R_m\le \frac{|X_m|}{12\cdot 2^m}.
\]
另一方面，Shannon 熵支配 Rényi--\(2\) 熵，故
\[
H\!\bigl(p_m^{\mathrm{bin}}\bigr)\ge H_2\!\bigl(p_m^{\mathrm{bin}}\bigr)
=-\log \mathsf{Col}^{\mathrm{bin}}_m.
\]
于是
\[
g_m
\le
m\log 2
+\log \mathsf{Col}^{\mathrm{bin}}_m
-1
+\frac{|X_m|\log(2\pi)+\sum_{w\in X_m}\log d_m^{\mathrm{bin}}(w)}{2^{m+1}}
+\frac{|X_m|}{12\cdot 2^m}.
\]

再由推论 \ref{cor:fold-zero-divisor-lb}，
\[
|\mathcal Z_m|\ge F_d,
\]
而定理 \ref{thm:fold-collision-zero-reduction} 给出
\[
\mathsf{Col}^{\mathrm{bin}}_m
\le
1-\frac{|\mathcal Z_m|}{M}
\le
1-\frac{F_d}{M}.
\]
代回上式即得第一条盒装不等式。

最后，与定理 \ref{thm:xi-time-part60aa-gauge-entropy-ledger-exact-identity} 证明末尾相同，由
\[
0\le \sum_{w\in X_m}\log d_m^{\mathrm{bin}}(w)\le |X_m|\,m\log 2
\]
可将显示校正项统一吸收到
\(
O(|X_m|m/2^m)
\)
中，从而得到第二条公式。证毕。
\end{proof}

\begin{theorem}[从规范群常数逐阶显式恢复全部偶值 \texorpdfstring{$\zeta(2r)$}{zeta(2r)}]\label{thm:xi-time-part60aa-gauge-constant-explicit-even-zeta-recovery}
\leanverified{paper\_xi\_time\_part60aa\_gauge\_constant\_explicit\_even\_zeta\_recovery}
沿用定理 \ref{thm:fold-bin-gauge-constant-stirling-bernoulli-hierarchy} 的规范群常数列 \((C_m)_{m\ge 1}\)，并对每个整数 \(r\ge 1\) 定义
\[
A_r:=
\frac{B_{2r}}{r(2r-1)}
\bigl(\varphi^{-1}+\varphi^{2r-3}\bigr).
\]
再定义递阶剥离后的极限算子
\[
L_r:=
\lim_{m\to\infty}
\Bigl(\frac{2}{\varphi}\Bigr)^{(2r-1)m}
\left[
\log C_m-\log(2\pi)
-\sum_{s=1}^{r-1}A_s\Bigl(\frac{\varphi}{2}\Bigr)^{(2s-1)m}
\right].
\]
则极限存在且满足
\[
\boxed{\;L_r=A_r.\;}
\]
因此对每个 \(r\ge 1\)，全部偶值 \(\zeta(2r)\) 都可由显式公式恢复：
\[
\boxed{\;
\zeta(2r)
=
(-1)^{r-1}
\frac{(2\pi)^{2r}}{2(2r)!}\,
\frac{r(2r-1)}{\varphi^{-1}+\varphi^{2r-3}}\,
L_r.\;}
\]
\end{theorem}
\begin{proof}
由定理 \ref{thm:fold-bin-gauge-constant-stirling-bernoulli-hierarchy}，对任意固定 \(r\ge 1\) 有展开
\[
\log C_m
=
\log(2\pi)
+\sum_{s=1}^{r}A_s\Bigl(\frac{\varphi}{2}\Bigr)^{(2s-1)m}
+O\!\Bigl(\Bigl(\frac{\varphi}{2}\Bigr)^{(2r+1)m}\Bigr).
\]
将前 \(r-1\) 阶项移到左边，得到
\[
\log C_m-\log(2\pi)
-\sum_{s=1}^{r-1}A_s\Bigl(\frac{\varphi}{2}\Bigr)^{(2s-1)m}
=
A_r\Bigl(\frac{\varphi}{2}\Bigr)^{(2r-1)m}
+O\!\Bigl(\Bigl(\frac{\varphi}{2}\Bigr)^{(2r+1)m}\Bigr).
\]
两边同乘
\[
\Bigl(\frac{2}{\varphi}\Bigr)^{(2r-1)m},
\]
便得到
\[
\Bigl(\frac{2}{\varphi}\Bigr)^{(2r-1)m}
\left[
\log C_m-\log(2\pi)
-\sum_{s=1}^{r-1}A_s\Bigl(\frac{\varphi}{2}\Bigr)^{(2s-1)m}
\right]
=
A_r
+O\!\Bigl(\Bigl(\frac{\varphi}{2}\Bigr)^{2m}\Bigr).
\]
由于 \(\varphi/2<1\)，右端误差趋于 \(0\)，故极限存在且等于 \(A_r\)，即 \(L_r=A_r\)。

最后，由标准 Bernoulli--zeta 恒等式
\[
B_{2r}
=
(-1)^{r-1}\frac{2(2r)!}{(2\pi)^{2r}}\zeta(2r),
\]
以及
\[
A_r=
\frac{\varphi^{-1}+\varphi^{2r-3}}{r(2r-1)}\,B_{2r}
=L_r,
\]
解出 \(\zeta(2r)\) 即得所示恢复公式。证毕。
\end{proof}

\begin{theorem}[Stirling--Bernoulli 级数截断余项的最终交错变号]\label{thm:xi-time-part60aa-stirling-bernoulli-tail-alternation}
沿用定理 \ref{thm:fold-bin-gauge-constant-stirling-bernoulli-hierarchy} 与定理 \ref{thm:xi-time-part60aa-gauge-constant-explicit-even-zeta-recovery} 的记号。记
\[
q_r:=\Bigl(\frac{\varphi}{2}\Bigr)^{2r-1},
\qquad
a_r:=A_r=
\frac{B_{2r}}{r(2r-1)}\bigl(\varphi^{-1}+\varphi^{2r-3}\bigr).
\]
对任意 \(R\ge 1\)，定义截断余项
\[
\mathcal R_R(m):=
\log C_m-\log(2\pi)-\sum_{r=1}^{R-1}a_r q_r^m.
\]
则有
\[
\mathcal R_R(m)=a_R q_R^m+O(q_{R+1}^m),
\qquad
\frac{q_{R+1}}{q_R}=\Bigl(\frac{\varphi}{2}\Bigr)^2<1.
\]
从而存在 \(m_0(R)\)，使得对一切 \(m\ge m_0(R)\)，都有
\[
\operatorname{sgn}\mathcal R_R(m)=\operatorname{sgn}(a_R)=\operatorname{sgn}(B_{2R})=(-1)^{R-1}.
\]
特别地，截断和
\[
S_{R-1}(m):=\log(2\pi)+\sum_{r=1}^{R-1}a_r q_r^m
\]
对 \(\log C_m\) 构成最终交错的上、下界：当 \(R\) 为奇数时 \(S_{R-1}(m)<\log C_m\)，当 \(R\) 为偶数时 \(S_{R-1}(m)>\log C_m\)（均对充分大的 \(m\) 成立）。由于指数函数在实轴上严格单调，相应的
\[
\exp(S_{R-1}(m))
\]
也对 \(C_m\) 构成同样交错的最终上、下界。
\end{theorem}
\begin{proof}
由定理 \ref{thm:fold-bin-gauge-constant-stirling-bernoulli-hierarchy}，对固定的 \(R\ge 1\) 有渐近展开
\[
\log C_m
=
\log(2\pi)+\sum_{r=1}^{R}a_r q_r^m+O(q_{R+1}^m).
\]
将前 \(R-1\) 项移到左边，便得到
\[
\mathcal R_R(m)=a_R q_R^m+O(q_{R+1}^m).
\]
又因为
\[
\frac{q_{R+1}}{q_R}=\Bigl(\frac{\varphi}{2}\Bigr)^2<1,
\]
所以 \(a_R q_R^m\) 终将支配余项，故存在 \(m_0(R)\) 使得对一切 \(m\ge m_0(R)\)，
\[
\operatorname{sgn}\mathcal R_R(m)=\operatorname{sgn}(a_R).
\]

再注意到
\[
\varphi^{-1}+\varphi^{2R-3}>0,
\]
故 \(a_R\) 的符号完全由 \(B_{2R}\) 决定。Bernoulli 数满足经典符号律
\[
\operatorname{sgn}(B_{2R})=(-1)^{R-1},
\]
遂得
\[
\operatorname{sgn}\mathcal R_R(m)=(-1)^{R-1}.
\]
于是当 \(R\) 为奇数时余项为正，从而
\[
S_{R-1}(m)<\log C_m;
\]
当 \(R\) 为偶数时余项为负，从而
\[
S_{R-1}(m)>\log C_m.
\]
最后，由指数函数在实轴上严格递增，同样的上下界关系自动传递到 \(C_m=\exp(\log C_m)\)。证毕。
\end{proof}

\noindent
上述五条结论把规范群与 fold-gauge 常数的两条层级链收束为同一份显式账本：一方面，\(g_m\) 精确等于原始 \(m\)-比特预算减去折叠 Shannon 熵，并附带完全显示的 Stirling 校正；另一方面，\(\log C_m\) 的 Bernoulli 模态不仅可逐阶识别，而且既可经由显式极限算子直接恢复全部偶值 \(\zeta(2r)\)，也可在任意固定截断阶上形成最终交错的上下包络。与审计偶数窗上的最小退化扇区及其零余类机制合并后，群熵损失、群代数实现障碍与偶 \(\zeta\)-塔反演遂被压缩进同一条闭合链。

\endinput

\paragraph{容量层析、Wedderburn 信息完备性与最小退化块账本}
在连续容量曲线的完备统计性、规范群熵账本、纤维群胚代数的 Wedderburn 直和，以及审计偶数窗的 Fibonacci 最小退化律都已建立之后，还可把这些对象层接口进一步压缩为下列几条直接结论。

\begin{theorem}[连续容量族对实参数冻结面的层析完备性]\label{thm:xi-time-part60aaa-capacity-family-freezing-face-tomography}
设对每个 \(m\ge 1\) 给定连续容量曲线
\[
\mathcal C_m^{\mathrm{cont}}(T):=\sum_{x\in X_m}\min\bigl(d_m(x),T\bigr)
\qquad (T\ge 0).
\]
若对每个 \(q>0\) 压力极限
\[
P_q:=\lim_{m\to\infty}\frac1m\log S_q(m),
\qquad
S_q(m):=\sum_{x\in X_m}d_m(x)^q,
\]
都存在，则整族 \(\{\mathcal C_m^{\mathrm{cont}}\}_{m\ge 1}\) 唯一决定全部实参数压力谱 \(\{P_q\}_{q>0}\)。在严格基态间隙
\[
\alpha_*^{(2)}<\alpha_*
\]
下，它因而唯一决定冻结仿射面
\[
q\longmapsto q\alpha_*+g_*
\]
以及其实参数观测起点
\[
a_{\mathrm{fr}}^{\mathbb R}
:=
\inf\bigl\{
q>0:\ P_r=r\alpha_*+g_*
\ \text{对一切 }r\ge q\text{成立}
\bigr\}.
\]
\end{theorem}
\begin{proof}
由定理 \ref{thm:xi-capacity-tail-histogram-moment-complete-statistic}，对每个固定的 \(m\ge 1\) 与每个 \(q>0\)，连续容量曲线 \(\mathcal C_m^{\mathrm{cont}}\) 唯一决定幂和矩
\[
S_q(m)=\sum_{x\in X_m}d_m(x)^q.
\]
因此整族 \(\{\mathcal C_m^{\mathrm{cont}}\}_{m\ge 1}\) 唯一决定每个 \(q>0\) 的数列 \(m\mapsto S_q(m)\)，从而在极限存在时唯一决定
\[
P_q=\lim_{m\to\infty}\frac1m\log S_q(m).
\]

若进一步满足严格基态间隙，则定理 \ref{thm:fold-pressure-freezing-threshold} 给出：对每个实数 \(q>a_{\mathrm c}\)，都有
\[
P_q=q\alpha_*+g_*.
\]
故由已恢复的整条压力谱可唯一读出其高温区外的线性冻结支，并唯一恢复上述实参数起点 \(a_{\mathrm{fr}}^{\mathbb R}\)。证毕。
\end{proof}

\begin{theorem}[固定分辨率上的 Wedderburn--信息完备接口]\label{thm:xi-time-part60aaa-wedderburn-information-completeness}
\leanverified{paper\_xi\_time\_part60aaa\_wedderburn\_information\_completeness}
固定分辨率 \(m\ge 1\)。定义纤维直方图
\[
n_d(m):=\#\{x\in X_m:\ d_m(x)=d\}
\qquad(d\ge 1).
\]
则下列四类数据两两等价，并且彼此唯一决定：
\[
\textup{(i)}\ \{n_d(m)\}_{d\ge 1},
\qquad
\textup{(ii)}\ \mathcal C_m^{\mathrm{cont}},
\qquad
\textup{(iii)}\ \{S_q(m)\}_{q>0},
\qquad
\textup{(iv)}\ \mathbb C[\mathcal R_m^{\mathrm{bin}}]\ \text{的半单代数同构型}.
\]
此外，这四类等价数据还共同决定 fold-gauge 群的直积分解
\[
G_m^{\mathrm{bin}}:=\Aut(\Fold_m^{\mathrm{bin}})
\cong
\prod_{d\ge 1}\mathfrak S_d^{\,n_d(m)}.
\]
\end{theorem}
\begin{proof}
由定理 \ref{thm:xi-capacity-tail-histogram-moment-complete-statistic}，固定 \(m\) 时，连续容量曲线 \(\mathcal C_m^{\mathrm{cont}}\)、纤维直方图 \(\{n_d(m)\}_{d\ge 1}\) 与全矩谱 \(\{S_q(m)\}_{q>0}\) 彼此唯一决定。

另一方面，由命题 \ref{prop:fold-bin-groupoid-algebra-wedderburn}，
\[
\mathbb C[\mathcal R_m^{\mathrm{bin}}]
\cong
\bigoplus_{x\in X_m}M_{d_m(x)}(\mathbb C)
\cong
\bigoplus_{d\ge 1}M_d(\mathbb C)^{\oplus n_d(m)}.
\]
有限维半单复代数的矩阵块大小及其重数由 Artin--Wedderburn 定理唯一决定，因此半单代数同构型与直方图 \(\{n_d(m)\}_{d\ge 1}\) 等价。

最后，由命题 \ref{prop:fold-bin-gauge-decomposition}，
\[
G_m^{\mathrm{bin}}
\cong
\prod_{x\in X_m}\mathfrak S_{d_m(x)}
\cong
\prod_{d\ge 1}\mathfrak S_d^{\,n_d(m)}.
\]
故前述四类等价数据又共同决定 gauge 群的直积分解。证毕。
\end{proof}

\begin{proposition}[审计偶数窗上的最小退化块构成显式 Fibonacci 半单理想]\label{prop:xi-time-part60aaa-audited-even-minideal-fibonacci-semisimple}
对审计偶数窗
\[
m\in\{6,8,10,12\},
\]
定义最小退化扇区
\[
\mathcal S_{m,*}:=\{x\in X_m:\ d_m^{\mathrm{bin}}(x)=d_{\min}(m)\},
\]
以及相应块直和
\[
A_{m,\min}
:=
\bigoplus_{x\in\mathcal S_{m,*}}M_{d_{\min}(m)}(\mathbb C)
\subset
\mathbb C[\mathcal R_m^{\mathrm{bin}}].
\]
则 \(A_{m,\min}\) 是一个半单理想，并且存在显式同构
\[
A_{m,\min}\cong M_{F_{m/2}}(\mathbb C)^{\oplus F_m}.
\]
特别地，
\[
\dim_{\mathbb C}A_{m,\min}=F_m\,F_{m/2}^{\,2},
\qquad
Z(A_{m,\min})\cong \mathbb C^{F_m}.
\]
\end{proposition}
\begin{proof}
由定理 \ref{thm:fold-bin-min-degeneracy-fib-audited}，在审计偶数窗上有
\[
d_{\min}(m)=F_{m/2},
\qquad
|\mathcal S_{m,*}|=F_m.
\]
另一方面，命题 \ref{prop:fold-bin-groupoid-algebra-wedderburn} 给出
\[
\mathbb C[\mathcal R_m^{\mathrm{bin}}]
\cong
\bigoplus_{x\in X_m}M_{d_m^{\mathrm{bin}}(x)}(\mathbb C).
\]
因此 \(A_{m,\min}\) 正是其中由全部最小退化块组成的直和，故自动是一个半单理想，并满足
\[
A_{m,\min}
\cong
M_{d_{\min}(m)}(\mathbb C)^{\oplus |\mathcal S_{m,*}|}
\cong
M_{F_{m/2}}(\mathbb C)^{\oplus F_m}.
\]
维数公式
\[
\dim_{\mathbb C}M_n(\mathbb C)=n^2
\]
立即给出
\[
\dim_{\mathbb C}A_{m,\min}=F_m\,F_{m/2}^{\,2},
\]
而矩阵代数直和的中心同构于对应块数目的复数直和，故
\[
Z(A_{m,\min})\cong \mathbb C^{F_m}.
\]
证毕。
\end{proof}

\begin{corollary}[exact gauge 账本相对 naive 缺陷账本的一纳特主阶回扣]\label{cor:xi-time-part60aaa-one-nat-unnormalized-ledger-rebate}
定义 naive 缺陷账本与 exact gauge 账本
\[
A_m:=\sum_{x\in X_m}d_m^{\mathrm{bin}}(x)\log d_m^{\mathrm{bin}}(x),
\qquad
L_m:=\log|G_m^{\mathrm{bin}}|.
\]
则有
\[
L_m
=
A_m-2^m+O(m\varphi^m).
\]
等价地，
\[
\frac{L_m}{2^m}
=
\kappa_m-1+O\!\Bigl(m\Bigl(\frac{\varphi}{2}\Bigr)^m\Bigr).
\]
\end{corollary}
\begin{proof}
由定理 \ref{thm:xi-time-part60aa-log-stieltjes-capacity-gauge-density}，
\[
\kappa_m-g_m
=
1+O\!\Bigl(m\Bigl(\frac{\varphi}{2}\Bigr)^m\Bigr),
\]
其中
\[
\kappa_m=\frac1{2^m}\sum_{x\in X_m}d_m^{\mathrm{bin}}(x)\log d_m^{\mathrm{bin}}(x)=\frac{A_m}{2^m},
\qquad
g_m=\frac1{2^m}\log|G_m^{\mathrm{bin}}|=\frac{L_m}{2^m}.
\]
代回即得第二式。再将两边同乘 \(2^m\)，并用
\[
2^m\cdot O\!\Bigl(m\Bigl(\frac{\varphi}{2}\Bigr)^m\Bigr)=O(m\varphi^m),
\]
便得到第一式。证毕。
\end{proof}

\noindent
与冻结相中的基态等权塌缩相合，上述结果表明：离散多重度谱中的层纹与波动在连续容量口径下并未消失，而是分别迁移为冻结面的线性支、群胚代数的半单块结构以及 exact gauge 账本相对于 naive 缺陷账本的一纳特主阶回扣。因而，连续平滑并不是离散结构的抹除，而是离散结构向基态等权、容量层析与 Wedderburn 分块三类函子像的重编码。

\endinput

\paragraph{隐藏位压力微分、Bernoulli 递阶反演与谱零块的二项式因子化}
在推前熵账本、Rényi 散度冻结、规范群体积的 Stirling--Bernoulli 层级，以及约数驱动的谱零块均已建立之后，还可把隐藏位统计、碰撞指纹、Bernoulli--\(\zeta\) 反演与群环多项式的代数刚性进一步压缩为下列一组闭合结论。

\begin{theorem}[压力谱在 \texorpdfstring{$a=1$}{a=1} 的一、二阶导数分别给出隐藏对数多重度的密度与涨落]\label{thm:xi-time-part60ab-hidden-logmultiplicity-pressure-derivatives}
\leanverified{paper\_xi\_time\_part60ab\_hidden\_logmultiplicity\_pressure\_derivatives}
沿用 H.3.1 的纤维多重度记号。对每个 \(m\) 记
\[
\mu_m(x):=\frac{d_m(x)}{2^m},
\qquad
S_a(m):=\sum_{x\in X_m} d_m(x)^a,
\qquad
P_m(a):=\frac1m\log S_a(m).
\]
设开区间 \(I\subset (0,\infty)\) 满足 \(1\in I\)，并假设每个 \(P_m\) 在 \(I\) 上属于 \(C^2\)，且存在 \(P\in C^2(I)\) 使
\[
P_m\longrightarrow P
\qquad\text{于 }C^2_{\mathrm{loc}}(I).
\]
若 \(X\sim \mu_m\)，定义隐藏对数多重度随机变量
\[
\mathscr H_m:=\log d_m(X).
\]
则对每个 \(m\) 都有精确恒等式
\[
\frac1m\EE_{\mu_m}[\mathscr H_m]=P_m'(1),
\qquad
\frac1m\Var_{\mu_m}(\mathscr H_m)=P_m''(1),
\]
从而
\[
\boxed{\;
\lim_{m\to\infty}\frac1m\EE_{\mu_m}[\mathscr H_m]=P'(1),\;}
\]
\[
\boxed{\;
\lim_{m\to\infty}\frac1m\Var_{\mu_m}(\mathscr H_m)=P''(1).\;}
\]
\end{theorem}
\begin{proof}
定义
\[
\Phi_m(a):=\frac1m\log\sum_{x\in X_m}\mu_m(x)^a.
\]
由
\[
\mu_m(x)=\frac{d_m(x)}{2^m}
\]
可得
\[
\Phi_m(a)
=
\frac1m\log\sum_x d_m(x)^a-a\log 2
=
P_m(a)-a\log 2.
\]

另一方面，
\[
\Phi_m'(a)
=
\frac1m
\frac{\sum_x \mu_m(x)^a\log \mu_m(x)}{\sum_x\mu_m(x)^a}.
\]
在 \(a=1\) 处利用 \(\sum_x\mu_m(x)=1\)，得到
\[
\Phi_m'(1)
=
\frac1m\sum_x\mu_m(x)\log\mu_m(x)
=
\frac1m\sum_x\mu_m(x)\log d_m(x)-\log 2
=
\frac1m\EE_{\mu_m}[\mathscr H_m]-\log 2.
\]
而由 \(\Phi_m(a)=P_m(a)-a\log 2\) 又有
\[
\Phi_m'(1)=P_m'(1)-\log 2.
\]
两式相减即得
\[
\frac1m\EE_{\mu_m}[\mathscr H_m]=P_m'(1).
\]

再看二阶导数。标准指数族计算给出
\[
\Phi_m''(1)
=
\frac1m
\left(
\sum_x\mu_m(x)(\log\mu_m(x))^2
\;-\;
\Bigl(\sum_x\mu_m(x)\log\mu_m(x)\Bigr)^2
\right)
=
\frac1m\Var_{\mu_m}(\log\mu_m(X)).
\]
又因
\[
\log\mu_m(X)=\mathscr H_m-m\log 2,
\]
加常数不改变方差，故
\[
\Var_{\mu_m}(\log\mu_m(X))=\Var_{\mu_m}(\mathscr H_m).
\]
另一方面，
\[
\Phi_m''(1)=P_m''(1).
\]
遂得
\[
\frac1m\Var_{\mu_m}(\mathscr H_m)=P_m''(1).
\]

最后，由 \(C^2_{\mathrm{loc}}(I)\) 收敛，必有
\[
P_m'(1)\to P'(1),
\qquad
P_m''(1)\to P''(1).
\]
代入前述有限尺度恒等式即得极限公式。证毕。
\end{proof}

\begin{theorem}[隐藏位缺口的 KL 识别与 Rényi--\texorpdfstring{$2$}{2} 碰撞证书]\label{thm:xi-time-part60ab-hidden-gap-kl-renyi2-collision}
沿用表 \ref{tab:fold_multiplicity_histogram} 的隐藏位预算记号。记
\[
u_m(x):=\frac1{|X_m|}
\qquad (x\in X_m),
\]
并定义隐藏位、可见位与容量缺口
\[
h_m:=\EE_{\mu_m}\bigl[\log_2 d_m(X)\bigr],
\qquad
I_m:=m-h_m,
\qquad
\mathrm{gap}_m:=\log_2|X_m|-I_m.
\]
再定义二阶碰撞指纹
\[
C_{\mathrm{col}}(m):=
|X_m|\sum_{x\in X_m}\mu_m(x)^2.
\]
则有严格恒等式
\[
\boxed{\;
\mathrm{gap}_m
=
\frac{1}{\log 2}\,
D_{\mathrm{KL}}(\mu_m\|u_m),\;}
\]
以及精确上界
\[
\boxed{\;
\mathrm{gap}_m
\le
\log_2 C_{\mathrm{col}}(m).\;}
\]
并且下列三条等价：
\begin{enumerate}
  \item \(\mathrm{gap}_m=\log_2 C_{\mathrm{col}}(m)\)；
  \item \(\mu_m=u_m\)；
  \item \(d_m(x)\) 在 \(X_m\) 上常值。
\end{enumerate}
\end{theorem}
\begin{proof}
由定义，
\[
\mathrm{gap}_m
=
\log_2|X_m|-m+\sum_{x\in X_m}\mu_m(x)\log_2 d_m(x).
\]
再用
\[
\mu_m(x)=\frac{d_m(x)}{2^m},
\qquad
u_m(x)=\frac1{|X_m|},
\]
得到
\[
\mathrm{gap}_m
=
\sum_x \mu_m(x)\log_2\frac{d_m(x)|X_m|}{2^m}
=
\sum_x \mu_m(x)\log_2\frac{\mu_m(x)}{u_m(x)}
=
\frac{1}{\log 2}\,D_{\mathrm{KL}}(\mu_m\|u_m).
\]
这便给出第一条恒等式。

另一方面，Rényi 散度对阶数单调，故
\[
D_{\mathrm{KL}}(\mu_m\|u_m)\le D_2(\mu_m\|u_m).
\]
而
\[
D_2(\mu_m\|u_m)
=
\log\sum_x \frac{\mu_m(x)^2}{u_m(x)}
=
\log\!\Bigl(|X_m|\sum_x \mu_m(x)^2\Bigr)
=
\log C_{\mathrm{col}}(m).
\]
两边除以 \(\log 2\) 即得
\[
\mathrm{gap}_m\le \log_2 C_{\mathrm{col}}(m).
\]

最后，\(D_{\mathrm{KL}}(\mu_m\|u_m)=D_2(\mu_m\|u_m)\) 当且仅当对 \(\mu_m\)-几乎处处的 \(x\)，比值 \(\mu_m(x)/u_m(x)\) 为常数。由于 \(u_m\) 在 \(X_m\) 上严格正、且二者都是概率分布，这等价于 \(\mu_m=u_m\)。而
\[
\mu_m=u_m
\iff
\frac{d_m(x)}{2^m}=\frac1{|X_m|}
\ \ (\forall x\in X_m)
\iff
d_m(x)\ \text{在 }X_m\text{ 上常值}.
\]
证毕。
\end{proof}

\begin{theorem}[fold-gauge 常数列的逐阶极限消元反演给出全部偶 Bernoulli 数与偶值 \texorpdfstring{$\zeta(2r)$}{zeta(2r)}]\label{thm:xi-time-part60ab-gauge-constant-recursive-bernoulli-stripping}
沿用定理 \ref{thm:fold-bin-gauge-constant-stirling-bernoulli-hierarchy} 的记号，记
\[
q:=\frac{\varphi}{2}\in(0,1),
\qquad
E_m:=\log C_m-\log(2\pi).
\]
并对每个整数 \(r\ge 1\) 定义
\[
a_r:=
\frac{B_{2r}}{r(2r-1)}
\bigl(\varphi^{-1}+\varphi^{2r-3}\bigr).
\]
再递归定义残差
\[
R_m^{(1)}:=E_m,
\qquad
R_m^{(r)}:=R_m^{(r-1)}-a_{r-1}q^{(2r-3)m}
\qquad (r\ge 2).
\]
则对每个 \(r\ge 1\) 都有极限
\[
\boxed{\;
\lim_{m\to\infty}q^{-(2r-1)m}R_m^{(r)}=a_r.\;}
\]
从而
\[
\boxed{\;
B_{2r}
=
\frac{r(2r-1)}{\varphi^{-1}+\varphi^{2r-3}}
\lim_{m\to\infty}
\Bigl(\frac{2}{\varphi}\Bigr)^{(2r-1)m}R_m^{(r)}.\;}
\]
\[
\boxed{\;
\zeta(2r)
=
(-1)^{r-1}\frac{(2\pi)^{2r}}{2(2r)!}\,
\frac{r(2r-1)}{\varphi^{-1}+\varphi^{2r-3}}
\lim_{m\to\infty}
\Bigl(\frac{2}{\varphi}\Bigr)^{(2r-1)m}R_m^{(r)}.\;}
\]
因此，单一标量序列 \((\log C_m)_{m\ge 1}\) 已足以逐阶恢复全部偶 Bernoulli 数与全部偶值 \(\zeta(2r)\)。
\end{theorem}
\begin{proof}
由定理 \ref{thm:fold-bin-gauge-constant-stirling-bernoulli-hierarchy}，对任意固定 \(R\ge 1\) 有
\[
E_m
=
\sum_{s=1}^{R}a_s q^{(2s-1)m}
+O\!\bigl(q^{(2R+1)m}\bigr).
\]
对上式取 \(R=r\)，并逐次减去前 \(r-1\) 个主项，得到
\[
R_m^{(r)}
=
a_r q^{(2r-1)m}+O\!\bigl(q^{(2r+1)m}\bigr).
\]
两边同乘 \(q^{-(2r-1)m}\) 得
\[
q^{-(2r-1)m}
R_m^{(r)}
=
a_r+O(q^{2m}),
\]
令 \(m\to\infty\) 即得极限公式。
再由
\[
a_r=
\frac{B_{2r}}{r(2r-1)}
\bigl(\varphi^{-1}+\varphi^{2r-3}\bigr)
\]
即可反解出 \(B_{2r}\)。
再由标准 Bernoulli--\(\zeta\) 恒等式
\[
B_{2r}
=
(-1)^{r-1}\frac{2(2r)!}{(2\pi)^{2r}}\zeta(2r),
\]
可立即反解出 \(\zeta(2r)\) 的显式恢复公式。证毕。
\end{proof}

\begin{corollary}[fold-gauge 常数的唯一主模与比例极限]\label{cor:xi-time-part60ab-gauge-constant-leading-mode-ratio}
沿用定理 \ref{thm:xi-time-part60ab-gauge-constant-recursive-bernoulli-stripping} 的记号，则
\[
\boxed{\;
\log C_m-\log(2\pi)
=
a_1 q^m+O(q^{3m}),
\qquad
q=\frac{\varphi}{2}.\;}
\]
从而
\[
\boxed{\;
\lim_{m\to\infty}
\frac{\log C_{m+1}-\log(2\pi)}
{\log C_m-\log(2\pi)}
=
\frac{\varphi}{2},\;}
\]
以及
\[
\boxed{\;
\lim_{m\to\infty}
\Bigl(\frac{2}{\varphi}\Bigr)^m
\bigl(\log C_m-\log(2\pi)\bigr)
=
a_1.\;}
\]
\end{corollary}
\begin{proof}
由上一条定理在 \(r=1\) 时的系数公式，
\[
a_1=
\frac{B_2}{1\cdot 1}
\bigl(\varphi^{-1}+\varphi^{-1}\bigr)
=
\frac16\cdot 2\varphi^{-1}
=
\frac{1}{3\varphi}.
\]
故
\[
\log C_m-\log(2\pi)
=
a_1 q^m+O(q^{3m}).
\]
于是
\[
\frac{\log C_{m+1}-\log(2\pi)}
{\log C_m-\log(2\pi)}
=
\frac{\frac{1}{3\varphi}q^{m+1}+O(q^{3m+3})}
{\frac{1}{3\varphi}q^m+O(q^{3m})}
\longrightarrow q=\frac{\varphi}{2},
\]
并且
\[
\Bigl(\frac{2}{\varphi}\Bigr)^m\bigl(\log C_m-\log(2\pi)\bigr)
=
\frac{1}{3\varphi}+O(q^{2m})
\longrightarrow
\frac{1}{3\varphi}.
\]
证毕。
\end{proof}

在上述逐阶极限消元公式之后，\(\log C_m\) 的 Bernoulli--\(\zeta\) 谱链还可进一步压缩为以下三条结构后果。

\begin{theorem}[fold-gauge 常数的 Richardson--Neville 任意阶消模器与 \texorpdfstring{$(2R+1)$}{(2R+1)} 级加速]\label{thm:xi-time-part60ab-gauge-constant-richardson-neville-accelerator}
沿用定理 \ref{thm:xi-time-part60ab-gauge-constant-recursive-bernoulli-stripping} 的记号，并定义
\[
\lambda_r:=q^{2r-1}
\qquad (r\ge 1).
\]
对任意固定整数 \(R\ge 1\)，令
\[
Q_R(z):=\prod_{r=1}^{R}\frac{z-\lambda_r}{1-\lambda_r}
=
\sum_{j=0}^{R}\gamma_{R,j}z^j,
\]
并定义线性加速器
\[
\mathcal A_R(m):=\sum_{j=0}^{R}\gamma_{R,j}\log C_{m+j}.
\]
则有
\[
\boxed{\;
\mathcal A_R(m)=\log(2\pi)+O\!\bigl(q^{(2R+1)m}\bigr).\;}
\]
换言之，\(\mathcal A_R\) 精确消去 \(\log C_m\) 渐近展开中的前 \(R\) 个 Bernoulli 模态，并把逼近 \(\log(2\pi)\) 的误差阶提升到 \(q^{(2R+1)m}\)。
\end{theorem}
\begin{proof}
由定义立即有
\[
Q_R(1)=1,
\qquad
Q_R(\lambda_r)=0
\qquad (1\le r\le R).
\]
另一方面，由定理 \ref{thm:fold-bin-gauge-constant-stirling-bernoulli-hierarchy}，对固定 \(R\) 有
\[
\log C_{m+j}
=
\log(2\pi)+\sum_{r=1}^{R}a_r\lambda_r^{m+j}
+O\!\bigl(\lambda_{R+1}^{m+j}\bigr)
\qquad (0\le j\le R).
\]
由于 \(j\) 仅在有限集合 \(\{0,\dots,R\}\) 中变化，上式余项可统一写为 \(O(\lambda_{R+1}^m)\)。将其代入 \(\mathcal A_R(m)\)，得到
\[
\mathcal A_R(m)
=
\log(2\pi)\sum_{j=0}^{R}\gamma_{R,j}
+\sum_{r=1}^{R}a_r\lambda_r^m\sum_{j=0}^{R}\gamma_{R,j}\lambda_r^j
+O\!\bigl(\lambda_{R+1}^m\bigr).
\]
再由
\[
\sum_{j=0}^{R}\gamma_{R,j}=Q_R(1)=1,
\qquad
\sum_{j=0}^{R}\gamma_{R,j}\lambda_r^j=Q_R(\lambda_r)=0
\qquad (1\le r\le R),
\]
主项中的前 \(R\) 个 Bernoulli 模态全部被精确消去，从而
\[
\mathcal A_R(m)
=
\log(2\pi)+O\!\bigl(\lambda_{R+1}^m\bigr)
=
\log(2\pi)+O\!\bigl(q^{(2R+1)m}\bigr).
\]
证毕。
\end{proof}

\begin{theorem}[fold-gauge 误差序列的非 \texorpdfstring{$C$}{C}-有限性与普通生成函数的非有理性]\label{thm:xi-time-part60ab-gauge-constant-non-cfinite-ogf-obstruction}
\leanverified{paper\_xi\_time\_part60ab\_gauge\_constant\_non\_cfinite\_ogf\_obstruction}
令
\[
u_m:=\log C_m-\log(2\pi).
\]
则序列 \((u_m)_{m\ge 1}\) 不是任何常系数线性递推序列，即不是 \(C\)-有限序列。因而其普通生成函数
\[
U(z):=\sum_{m\ge 1}u_m z^m
\]
不可能是有理函数。
\end{theorem}
\begin{proof}
反设 \((u_m)_{m\ge 1}\) 为 \(C\)-有限序列。则存在有限个复数 \(\mu_1,\dots,\mu_s\) 与多项式 \(P_1,\dots,P_s\)，使得对充分大的 \(m\) 有
\[
u_m=\sum_{j=1}^{s}P_j(m)\mu_j^m.
\]
因此 \(u_m\) 只可能由有限多个指数底数 \(\mu_j\) 生成。

另一方面，由定理 \ref{thm:xi-time-part60ab-gauge-constant-recursive-bernoulli-stripping}，对每个 \(R\ge 1\) 都有
\[
u_m-\sum_{r=1}^{R-1}a_r q^{(2r-1)m}
=
a_R q^{(2R-1)m}+O\!\bigl(q^{(2R+1)m}\bigr).
\]
又由于
\[
a_R=
\frac{B_{2R}}{R(2R-1)}
\bigl(\varphi^{-1}+\varphi^{2R-3}\bigr)
\]
且 \(B_{2R}\neq 0\)，故 \(a_R\neq 0\) 对每个 \(R\ge 1\) 都成立。于是对每个 \(R\)，尺度 \(q^{(2R-1)m}\) 都真实出现在 \(u_m\) 的渐近链中。另一方面，
\[
u_m-\sum_{r=1}^{R-1}a_r q^{(2r-1)m}
\]
仍是由有限个指数底数
\[
\{\mu_1,\dots,\mu_s,q,q^3,\dots,q^{2R-3}\}
\]
生成的多项式--指数和。由其渐近主项为非零的 \(a_R q^{(2R-1)m}\)，可知底数 \(q^{2R-1}\) 必属于上式有限集合。由于
\[
q^{2R-1}\notin \{q,q^3,\dots,q^{2R-3}\},
\]
遂只能有
\[
q^{2R-1}\in \{\mu_1,\dots,\mu_s\}.
\]
这与 \(\{\mu_1,\dots,\mu_s\}\) 有限而
\[
q,q^3,q^5,\dots
\]
两两不同矛盾。故 \((u_m)\) 不是 \(C\)-有限序列。

最后，普通生成函数有理当且仅当其系数序列 \(C\)-有限，这是经典等价命题，故 \(U(z)\) 不可能为有理函数。证毕。
\end{proof}

\begin{corollary}[fold-gauge Bernoulli 塔的无限模态障碍]\label{cor:xi-time-part60ab-gauge-constant-infinite-modal-obstruction}
令
\[
E_m:=\log C_m-\log(2\pi).
\]
则不存在有限组互异复数 \(\lambda_1,\dots,\lambda_s\) 与多项式 \(P_1,\dots,P_s\)，使得对充分大的 \(m\) 恒有
\[
E_m=\sum_{j=1}^{s}P_j(m)\lambda_j^m.
\]
等价地，H.111 的 Stirling--Bernoulli 层级展开对应一条真正的无限模态谱链，而非任何有限维 companion 机制所能封装的有限指数叠加。
\end{corollary}
\begin{proof}
这正是定理 \ref{thm:xi-time-part60ab-gauge-constant-non-cfinite-ogf-obstruction} 的直接重述，因为常系数线性递推序列与有限个多项式--指数项的和完全等价。证毕。
\end{proof}

\begin{theorem}[高阶尾部精度强制前导 Bernoulli--\texorpdfstring{$\zeta$}{zeta} 模态一致]\label{thm:xi-time-part60ab-gauge-constant-tail-rigidity}
沿用定理 \ref{thm:xi-time-part60ab-gauge-constant-recursive-bernoulli-stripping} 的记号。设另一实序列 \((D_m)_{m\ge 1}\) 满足：对某个固定整数 \(R\ge 1\)，存在实系数 \(b_1,\dots,b_R\) 使
\[
D_m-\log(2\pi)
=
\sum_{r=1}^{R}b_r q^{(2r-1)m}
+O\!\bigl(q^{(2R+1)m}\bigr),
\]
并且同时有
\[
D_m-\log C_m
=
O\!\bigl(q^{(2R+1)m}\bigr).
\]
则必有
\[
b_r=a_r
\qquad (1\le r\le R).
\]
因而前 \(R\) 个偶 Bernoulli 数 \(B_{2r}\) 与偶值 \(\zeta(2r)\) 也被同一尾部精度唯一锁定。
\end{theorem}
\begin{proof}
由定理 \ref{thm:fold-bin-gauge-constant-stirling-bernoulli-hierarchy}，
\[
\log C_m-\log(2\pi)
=
\sum_{r=1}^{R}a_r q^{(2r-1)m}
+O\!\bigl(q^{(2R+1)m}\bigr).
\]
将其与 \(D_m\) 的展开相减，得到
\[
D_m-\log C_m
=
\sum_{r=1}^{R}(b_r-a_r)q^{(2r-1)m}
+O\!\bigl(q^{(2R+1)m}\bigr).
\]
若存在最小的 \(r_0\le R\) 使 \(b_{r_0}\neq a_{r_0}\)，则上式可改写为
\[
D_m-\log C_m
=
(b_{r_0}-a_{r_0})q^{(2r_0-1)m}
+O\!\bigl(q^{(2r_0+1)m}\bigr).
\]
两边同乘 \(q^{-(2r_0-1)m}\)，并令 \(m\to\infty\)，便得到
\[
q^{-(2r_0-1)m}\bigl(D_m-\log C_m\bigr)\to b_{r_0}-a_{r_0}\neq 0,
\]
这与假设
\[
D_m-\log C_m
=
O\!\bigl(q^{(2R+1)m}\bigr)
\]
矛盾，因为后者特别蕴含
\(
D_m-\log C_m=O(q^{(2r_0+1)m})
\)。
故 \(b_r=a_r\) 对全部 \(1\le r\le R\) 成立。

最后，由定理 \ref{thm:xi-time-part60ab-gauge-constant-recursive-bernoulli-stripping} 中 \(a_r\) 与 \(B_{2r}\)、\(\zeta(2r)\) 的显式对应关系，前 \(R\) 个 Bernoulli--\(\zeta\) 模态也随之被唯一锁定。证毕。
\end{proof}

\endinput


\paragraph{Stirling--Bernoulli 模态的移位投影与偶值 \texorpdfstring{$\zeta$}{zeta} 显式恢复}
附录中的 Bernoulli 层级、递归抽取与有限阶移位消元已经给出三条彼此兼容的读出接口。把它们压回正文主线，可得到如下三条直接后果：每个 Stirling--Bernoulli 模态都可由一个显式有限阶移位投影子直接抽出；偶值 \(\zeta(2r)\) 因而成为单一规范常数序列上的直接可审计读出；而无限多个互异衰减模的同时存在，则排除了任何有限 companion 型封装。

\begin{theorem}[fold-gauge Bernoulli 模态的有限阶移位投影公式]\label{thm:xi-time-part60ab2-logcm-shift-projector}
\leanverified{paper\_xi\_time\_part60ab2\_logcm\_shift\_projector}
沿用定理 \ref{thm:fold-bin-gauge-constant-stirling-bernoulli-hierarchy} 的记号，记
\[
f_m:=\log C_m-\log(2\pi),
\qquad
q_r:=\left(\frac{\varphi}{2}\right)^{2r-1},
\qquad
a_r:=\frac{B_{2r}}{r(2r-1)}\bigl(\varphi^{-1}+\varphi^{2r-3}\bigr).
\]
并令移位算子 \(E\) 作用于序列 \(f=(f_m)_{m\ge 1}\) 为
\[
(Ef)_m:=f_{m+1}.
\]
对每个 \(r\ge 1\) 定义
\[
\mathcal P_{r-1}(E):=\prod_{j=1}^{r-1}(E-q_j),
\]
其中 \(r=1\) 时取空积 \(\mathcal P_0(E)=1\)。则极限
\[
\lim_{m\to\infty}
\frac{q_r^{-m}\bigl(\mathcal P_{r-1}(E)f\bigr)_m}
{\prod_{j=1}^{r-1}(q_r-q_j)}
\]
存在，并且恰等于 \(a_r\)。特别地，每个模态系数 \(a_r\) 都可由 \((\log C_m)_{m\ge 1}\) 的一个有限阶移位投影极限直接读出。
\end{theorem}
\begin{proof}
定理 \ref{thm:derived-fold-gauge-shift-annihilator-even-zeta} 已给出
\[
\lim_{m\to\infty}q_r^{-m}\bigl(\mathcal P_{r-1}(E)f\bigr)_m
=
a_r\prod_{j=1}^{r-1}(q_r-q_j).
\]
又由于
\[
q_1>q_2>\cdots>0,
\]
故对每个 \(j<r\) 都有 \(q_r-q_j\neq 0\)。两边同除以该非零乘积，即得所示公式。证毕。
\end{proof}

\begin{corollary}[有限阶移位投影给出的偶值 \texorpdfstring{$\zeta(2r)$}{zeta(2r)} 显式恢复公式]\label{cor:xi-time-part60ab2-logcm-shift-even-zeta}
\leanverified{paper\_xi\_time\_part60ab2\_logcm\_shift\_even\_zeta}
沿用定理 \ref{thm:xi-time-part60ab2-logcm-shift-projector} 的记号。对每个 \(r\ge 1\)，都有
\[
\zeta(2r)
=
(-1)^{r-1}\frac{(2\pi)^{2r}}{2(2r)!}\,
\frac{r(2r-1)}{\varphi^{-1}+\varphi^{2r-3}}\,
\frac{
\lim\limits_{m\to\infty}
q_r^{-m}\bigl(\mathcal P_{r-1}(E)f\bigr)_m
}
{\prod_{j=1}^{r-1}(q_r-q_j)}.
\]
因此，偶值 \(\zeta\)-谱并不只作为 \((\log C_m)\) 的渐近尾项存在，而是可由一族显式有限阶移位投影子逐模态直接恢复。
\end{corollary}
\begin{proof}
由定理 \ref{thm:xi-time-part60ab2-logcm-shift-projector}，
\[
a_r=
\frac{
\lim\limits_{m\to\infty}
q_r^{-m}\bigl(\mathcal P_{r-1}(E)f\bigr)_m
}
{\prod_{j=1}^{r-1}(q_r-q_j)}.
\]
再代入
\[
a_r=\frac{B_{2r}}{r(2r-1)}\bigl(\varphi^{-1}+\varphi^{2r-3}\bigr),
\qquad
B_{2r}=(-1)^{r-1}\frac{2(2r)!}{(2\pi)^{2r}}\zeta(2r),
\]
整理即得。证毕。
\end{proof}

\begin{corollary}[归一化移位投影子的单步 Bernoulli 恢复]\label{cor:xi-time-part60ab2-logcm-normalized-projector-bernoulli}
\leanverified{paper\_xi\_time\_part60ab2\_logcm\_normalized\_projector\_bernoulli}
沿用定理 \ref{thm:xi-time-part60ab2-logcm-shift-projector} 的记号。对每个 \(r\ge 1\)，定义归一化移位投影子
\[
\Pi_r(E):=
\prod_{j=1}^{r-1}\frac{E-q_j}{q_r-q_j},
\]
其中 \(r=1\) 时取空积 \(\Pi_1(E)=1\)。
则
\[
\lim_{m\to\infty}q_r^{-m}\bigl(\Pi_r(E)f\bigr)_m=a_r,
\]
从而
\[
\lim_{m\to\infty}q_r^{-m}\bigl(\Pi_r(E)f\bigr)_m
=
\frac{B_{2r}}{r(2r-1)}\bigl(\varphi^{-1}+\varphi^{2r-3}\bigr),
\]
以及
\[
B_{2r}
=
\frac{r(2r-1)}{\varphi^{-1}+\varphi^{2r-3}}
\lim_{m\to\infty}q_r^{-m}\bigl(\Pi_r(E)f\bigr)_m.
\]
因此，偶 Bernoulli 数可由一族封闭形式的归一化移位投影子直接逐阶读出，而无需求助于递归剥离此前各级模态。
\end{corollary}
\begin{proof}
由定义，
\[
\Pi_r(E)
=
\frac{\mathcal P_{r-1}(E)}{\prod_{j=1}^{r-1}(q_r-q_j)}.
\]
故定理 \ref{thm:xi-time-part60ab2-logcm-shift-projector} 立即给出
\[
\lim_{m\to\infty}q_r^{-m}\bigl(\Pi_r(E)f\bigr)_m=a_r.
\]
再代入
\[
a_r=\frac{B_{2r}}{r(2r-1)}\bigl(\varphi^{-1}+\varphi^{2r-3}\bigr)
\]
并整理，即得所示 Bernoulli 恢复公式。证毕。
\end{proof}

\begin{theorem}[有限阶移位投影强制 fold-gauge 谱塔无限]\label{thm:xi-time-part60ab2-logcm-shift-infinite-modal}
序列
\[
f_m=\log C_m-\log(2\pi)
\]
不是任何有限阶常系数线性递推序列；等价地，其普通生成函数
\[
F(z):=\sum_{m\ge 1}f_m z^m
\]
不是有理函数。
\end{theorem}
\leanverified{paper\_xi\_time\_part60ab2\_logcm\_shift\_infinite\_modal}
\begin{proof}
由定理 \ref{thm:xi-time-part60ab2-logcm-shift-projector}，对每个 \(r\ge 1\) 都有
\[
\lim_{m\to\infty}
\frac{q_r^{-m}\bigl(\mathcal P_{r-1}(E)f\bigr)_m}
{\prod_{j=1}^{r-1}(q_r-q_j)}
=a_r.
\]
而
\[
a_r=
\frac{B_{2r}}{r(2r-1)}\bigl(\varphi^{-1}+\varphi^{2r-3}\bigr)\neq 0,
\]
因为偶 Bernoulli 数 \(B_{2r}\) 从不消失，且 \(\varphi^{-1}+\varphi^{2r-3}>0\)。因此，\(f_m\) 的渐近谱包含无限多个互异底数
\[
q_r=\left(\frac{\varphi}{2}\right)^{2r-1},
\qquad
r=1,2,\dots.
\]
另一方面，任何有限阶常系数线性递推序列的渐近谱都只能由有限多个指数底数构成，故不可能容纳上述无限模态塔。于是 \(f_m\) 非 \(C\)-有限。

普通生成函数有理当且仅当其系数序列 \(C\)-有限，这是经典等价命题，遂得 \(F(z)\) 非有理。证毕。
\end{proof}

\noindent
上述三条结果把 \((\log C_m)\) 的 Bernoulli 层级从“逐阶剥离尾项”的恢复程序，提升为一族显式有限阶移位滤波器的直接读出接口。因而，偶值 \(\zeta\)-谱并非只是在渐近展开中被动出现，而可由单一规范常数序列经过确定性的有限阶投影逐模态抽出；同时，无限模态障碍又说明这种读出结构本质上超出任何有限状态线性递推的封装能力。

\endinput


\begin{theorem}[约数触发的谱零块强制群环多项式含有显式二项式因子]\label{thm:xi-time-part60ab-zero-block-forces-binomial-factor}
\leanverified{paper\_xi\_time\_part60ab\_zero\_block\_forces\_binomial\_factor}
沿用命题 \ref{prop:fold-multiplicity-group-algebra} 的记号。记
\[
M:=F_{m+2},
\qquad
D_m(x):=\sum_{r=0}^{M-1}d_m(r)x^r\in\ZZ[x],
\qquad
\deg D_m<M.
\]
若 \(d\mid (m+2)\) 是真约数且满足
\[
F_d\mid \frac{M}{2},
\]
则有整系数因子结论
\[
\boxed{\;
x^{F_d}+1\mid D_m(x)\qquad\text{于 }\ZZ[x].\;}
\]
特别地，当
\[
m+2\equiv 0\pmod 6
\]
并取 \(d=(m+2)/2\) 时，
\[
\boxed{\;
x^{F_{(m+2)/2}}+1\mid D_m(x).\;}
\]
\end{theorem}
\begin{proof}
令
\[
t:=\frac{M}{2F_d}\in\NN.
\]
由推论 \ref{cor:fold-zero-divisor-lb}，
\[
S_{F_d}(M)\subseteq \mathcal Z_m
:=
\{k\in\ZZ/(M\ZZ):\ \widehat d_m(k)=0\}.
\]
对每个 \(\ell=0,1,\dots,F_d-1\)，取
\[
k_\ell:=(2\ell+1)t.
\]
则
\[
k_\ell F_d=(2\ell+1)\frac{M}{2}\equiv \frac{M}{2}\pmod M,
\]
故 \(k_\ell\in S_{F_d}(M)\subseteq \mathcal Z_m\)。于是由离散 Fourier 变换的定义，
\[
D_m(\omega_M^{-k_\ell})=\widehat d_m(k_\ell)=0,
\qquad
\omega_M:=e^{2\pi i/M}.
\]

另一方面，
\[
(\omega_M^{-k_\ell})^{F_d}
=
\exp\!\left(-2\pi i\,\frac{k_\ell F_d}{M}\right)
=
\exp\!\left(-(2\ell+1)\pi i\right)
=
-1.
\]
故每个 \(\omega_M^{-k_\ell}\) 都是 \(x^{F_d}+1\) 的根。并且
\[
\omega_M^{-k_\ell}
=
\exp\!\left(-\frac{(2\ell+1)\pi i}{F_d}\right)
\qquad (0\le \ell\le F_d-1)
\]
恰给出 \(x^{F_d}+1\) 的全部 \(F_d\) 个互异复根。因此 \(x^{F_d}+1\) 的全部根都是 \(D_m(x)\) 的根，从而
\[
x^{F_d}+1\mid D_m(x)
\qquad\text{于 }\CC[x].
\]
由于二者都属于 \(\ZZ[x]\)，Gauss 引理进一步给出
\[
x^{F_d}+1\mid D_m(x)
\qquad\text{于 }\ZZ[x].
\]

当 \(m+2\equiv 0\pmod 6\) 时，推论 \ref{cor:fold-zero-half-index-multiple6} 恰给出
\[
F_{(m+2)/2}\mid \frac{F_{m+2}}{2},
\]
故特殊情形立即成立。证毕。
\end{proof}

\noindent
上述链条把表 \ref{tab:fold_multiplicity_histogram} 中的隐藏位、容量缺口与碰撞指纹严格并入统一的信息几何结构；同时又把 \(\log C_m\) 的渐近模态提升为一条逐阶可反演、可任意阶消模并伴随非 \(C\)-有限障碍与尾部刚性的 Bernoulli--\(\zeta\) 谱链，并把约数驱动的谱零块最终刚化为群环多项式上的显式整系数因子。

\endinput

\paragraph{bin-fold 群论三重分解、中心态熵同一、规范体积常数项与零余类算术放大}
在纤维置换分解、群胚代数的 Wedderburn 直和、输出 Rényi 熵的常数项闭式以及零余类的约数触发机制已经建立之后，还可把同一套退化直方图进一步压缩为下列六条对象层结论。

\begin{theorem}[bin-fold 规范群的中心--交换化--导群三重分解]\label{thm:xi-foldbin-gauge-group-center-abel-derived-triple-decomposition}
沿用定义 \ref{def:fold-bin-gauge-group} 与定义 \ref{def:fold-bin-degeneracy-sectors} 的记号，记
\[
G_m^{\mathrm{bin}}:=\Aut(\Fold_m^{\mathrm{bin}}),
\qquad
\mathcal S_{m,d}:=\{w\in X_m:\ d_m^{\mathrm{bin}}(w)=d\}.
\]
则有自然同构
\[
Z\!\bigl(G_m^{\mathrm{bin}}\bigr)\cong (\ZZ/2\ZZ)^{|\mathcal S_{m,2}|},
\]
\[
\bigl(G_m^{\mathrm{bin}}\bigr)^{\mathrm{ab}}
\cong
(\ZZ/2\ZZ)^{\sum_{d\ge 2}|\mathcal S_{m,d}|},
\]
\[
[G_m^{\mathrm{bin}},G_m^{\mathrm{bin}}]
\cong
\prod_{d\ge 3}\mathfrak A_d^{\,|\mathcal S_{m,d}|}.
\]
特别地，若 \(d_{\min}(m)\ge 2\)，则
\[
\bigl(G_m^{\mathrm{bin}}\bigr)^{\mathrm{ab}}
\cong
(\ZZ/2\ZZ)^{|X_m|}.
\]
对审计偶数窗
\(
m\in\{6,8,10,12\}
\)，进一步有
\[
\bigl(G_m^{\mathrm{bin}}\bigr)^{\mathrm{ab}}
\cong
(\ZZ/2\ZZ)^{F_{m+2}}.
\]
\leanverified{paper_xi_foldbin_gauge_group_triple_decomp_seeds}
\leanverified{paper_xi_foldbin_gauge_group_triple_decomp_package}
\end{theorem}
\begin{proof}
由命题 \ref{prop:fold-bin-gauge-decomposition}，
\[
G_m^{\mathrm{bin}}
\cong
\prod_{w\in X_m}\mathfrak S_{d_m^{\mathrm{bin}}(w)}.
\]
有限直积对中心、交换化与导群分别可按因子逐项计算，而对称群满足
\[
Z(\mathfrak S_d)\cong
\begin{cases}
\ZZ/2\ZZ,& d=2,\\
1,& d\neq 2,
\end{cases}
\]
\[
\mathfrak S_d^{\mathrm{ab}}\cong
\begin{cases}
1,& d=1,\\
\ZZ/2\ZZ,& d\ge 2,
\end{cases}
\qquad
[\mathfrak S_d,\mathfrak S_d]\cong
\begin{cases}
1,& d\le 2,\\
\mathfrak A_d,& d\ge 3.
\end{cases}
\]
把这些公式逐纤维代入，便得到
\[
Z\!\bigl(G_m^{\mathrm{bin}}\bigr)\cong (\ZZ/2\ZZ)^{|\mathcal S_{m,2}|},
\]
\[
\bigl(G_m^{\mathrm{bin}}\bigr)^{\mathrm{ab}}
\cong
(\ZZ/2\ZZ)^{\sum_{d\ge 2}|\mathcal S_{m,d}|},
\]
\[
[G_m^{\mathrm{bin}},G_m^{\mathrm{bin}}]
\cong
\prod_{d\ge 3}\mathfrak A_d^{\,|\mathcal S_{m,d}|}.
\]

若 \(d_{\min}(m)\ge 2\)，则每条纤维都贡献一个 \(\ZZ/2\ZZ\) 因子，故
\[
\sum_{d\ge 2}|\mathcal S_{m,d}|=|X_m|,
\]
从而
\[
\bigl(G_m^{\mathrm{bin}}\bigr)^{\mathrm{ab}}\cong (\ZZ/2\ZZ)^{|X_m|}.
\]
对 \(m\in\{6,8,10,12\}\)，定理 \ref{thm:fold-bin-min-degeneracy-fib-audited} 给出
\[
d_{\min}(m)=F_{m/2}\ge 2,
\]
再由稳定类型计数律 \(|X_m|=F_{m+2}\)，即得最后一式。证毕。
\end{proof}

\begin{theorem}[群胚代数中心态的 Rényi--Shannon 熵与输出熵的严格同一性]\label{thm:xi-foldbin-groupoid-central-renyi-shannon-identity}
沿用命题 \ref{prop:fold-bin-groupoid-algebra-wedderburn} 与定理 \ref{thm:xi-foldbin-fkdet-logvolume-kappa} 的记号，记
\[
A_m:=\CC[\mathcal R_m^{\mathrm{bin}}]
\cong
\bigoplus_{w\in X_m}M_{d_m^{\mathrm{bin}}(w)}(\CC),
\qquad
Z(A_m)=\bigoplus_{w\in X_m}\CC z_w,
\]
其中 \(z_w\) 为第 \(w\)-块的最小中心幂等元，并在 \(A_m\) 上赋予概率迹
\[
\tau_m\!\left(\bigoplus_{w\in X_m}a_w\right):=
2^{-m}\sum_{w\in X_m}\operatorname{Tr}(a_w).
\]
则
\[
\tau_m(z_w)=\frac{d_m^{\mathrm{bin}}(w)}{2^m}=:\mu_m(w),
\]
故中心极小幂等元的谱权分布恰与 bin-fold 输出分布 \(\mu_m\) 一致。

对 \(q>0\)、\(q\neq 1\)，定义中心 Rényi 熵
\[
H_q^{\mathrm{cen}}(m):=
\frac{1}{1-q}\log\sum_{w\in X_m}\tau_m(z_w)^q,
\]
并把 \(q=1\) 情形按连续延拓记为
\[
H_1^{\mathrm{cen}}(m):=
-\sum_{w\in X_m}\tau_m(z_w)\log\tau_m(z_w).
\]
则对一切 \(q>0\) 都有
\[
H_q^{\mathrm{cen}}(m)=H_q(\mu_m).
\]
特别地，对 \(q>0\)、\(q\neq 1\)，有
\[
H_q^{\mathrm{cen}}(m)
=
m\log\varphi
+\frac{1}{1-q}\log\!\Bigl(\frac{\varphi+\varphi^{-q}}{\sqrt5}\Bigr)
+O\!\Bigl(\Bigl(\frac{\varphi}{2}\Bigr)^m\Bigr),
\]
并且
\[
H_1^{\mathrm{cen}}(m)
=
m\log\varphi
+\frac{\log\varphi}{1+\varphi^2}
+O\!\Bigl(\Bigl(\frac{\varphi}{2}\Bigr)^m\Bigr)
=
\log|X_m|-C_{\mathrm{KL}}+o(1),
\]
其中
\[
C_{\mathrm{KL}}
=
2\log\varphi-\frac12\log 5-\frac{\log\varphi}{\varphi\sqrt5}.
\]
从而
\[
\frac{1}{m}H_q^{\mathrm{cen}}(m)\longrightarrow \log\varphi
\qquad (q>0).
\]
\end{theorem}
\begin{proof}
在 Wedderburn 分解下，\(z_w\) 正是第 \(w\)-块上的单位元与其余各块上的零元，故
\[
\tau_m(z_w)
=
2^{-m}\operatorname{Tr}(I_{d_m^{\mathrm{bin}}(w)})
=
2^{-m}d_m^{\mathrm{bin}}(w)
=
\mu_m(w).
\]
因此中心原子 \(\{z_w\}_{w\in X_m}\) 的谱权分布恰为 \(\mu_m\)，从而
\[
H_q^{\mathrm{cen}}(m)=H_q(\mu_m)
\]
对一切 \(q>0\) 都成立。

当 \(q\neq 1\) 时，把该恒等式代入定理 \ref{thm:xi-foldbin-absolute-renyi-entropy-constant-closed} 即得 Rényi 常数项闭式；\(q=1\) 的精确常数项同样由该定理给出。再由命题 \ref{prop:fold-bin-kl-constant} 的
\[
H(\mu_m)=\log|X_m|-C_{\mathrm{KL}}+o(1)
\]
便得到 Shannon 形式的第二种写法。最后，
\[
\frac1m H_q^{\mathrm{cen}}(m)\to \log\varphi
\]
可由上述显式展开直接读出，亦可视为定理 \ref{thm:fold-bin-renyi-rate-collapse} 的中心态版本。证毕。
\leanverified{paper\_xi\_foldbin\_groupoid\_central\_renyi\_shannon\_identity}
\end{proof}

\begin{theorem}[规范体积密度的精确热力学常数项]\label{thm:xi-foldbin-gauge-density-exact-thermodynamic-constant}
沿用定理 \ref{thm:xi-foldbin-groupoid-central-renyi-shannon-identity} 的记号。则有
\[
g_m
=
m\log\!\Bigl(\frac{2}{\varphi}\Bigr)
-1
-\frac{\log\varphi}{1+\varphi^2}
+O\!\Bigl(m\Bigl(\frac{\varphi}{2}\Bigr)^m\Bigr).
\]
从而
\[
\lim_{m\to\infty}
\left(
g_m-m\log\!\Bigl(\frac{2}{\varphi}\Bigr)
\right)
=
-1-\frac{\log\varphi}{1+\varphi^2}.
\]
\end{theorem}
\begin{proof}
由定理 \ref{thm:xi-foldbin-groupoid-central-renyi-shannon-identity} 在 \(q=1\) 的结论，
\[
H(\mu_m)=H_1^{\mathrm{cen}}(m)
=
m\log\varphi+\frac{\log\varphi}{1+\varphi^2}
+O\!\Bigl(\Bigl(\frac{\varphi}{2}\Bigr)^m\Bigr).
\]
另一方面，定理 \ref{thm:xi-foldbin-gauge-entropy-one-nat-law} 给出
\[
g_m
=
m\log 2-H(\mu_m)-1
+O\!\Bigl(m\Bigl(\frac{\varphi}{2}\Bigr)^m\Bigr).
\]
将前式代入即得
\[
g_m
=
m\log\!\Bigl(\frac{2}{\varphi}\Bigr)
-1
-\frac{\log\varphi}{1+\varphi^2}
+O\!\Bigl(m\Bigl(\frac{\varphi}{2}\Bigr)^m\Bigr),
\]
并立刻推出极限常数。证毕。
\end{proof}
\leanverified{paper\_xi\_foldbin\_gauge\_density\_exact\_thermodynamic\_constant}

\begin{theorem}[无限算术子序列上的零余类放大律]\label{thm:xi-fold-zero-arithmetic-subsequence-amplification}
\leanverified{paper\_xi\_fold\_zero\_arithmetic\_subsequence\_amplification}
令
\[
M:=F_{m+2},
\qquad
\mathcal Z_m:=\{k\in\ZZ/(M\ZZ):\ \widehat d_m(k)=0\},
\qquad
\bar d_m:=\frac{2^m}{F_{m+2}}.
\]
并沿用定理 \ref{thm:fold-zero-uncertainty} 的记号
\[
p_m(r):=\frac{d_m(r)}{2^m},
\qquad
u_m(r):=\frac1M.
\]
若
\[
m\equiv 4\pmod 6,
\]
则
\[
|\mathcal Z_m|\ge F_{(m+2)/2}.
\]
因而
\[
\mathsf{Col}_m
\le
1-\frac{F_{(m+2)/2}}{F_{m+2}},
\qquad
D_{\mathrm{KL}}(p_m\|u_m)
\le
\log\!\bigl(F_{m+2}-F_{(m+2)/2}\bigr),
\]
\[
M_m-\bar d_m
\le
2^m
\sqrt{
\frac{F_{m+2}-1-F_{(m+2)/2}}{F_{m+2}}
}.
\]
此外，
\[
\frac{F_{(m+2)/2}}{F_{m+2}}=\Theta(\varphi^{-m/2}),
\]
故零余类所强制的碰撞削减项在该无限算术子序列上具有 \(\Theta(\varphi^{-m/2})\) 的确定性尺度。
\leanverified{paper\_xi\_fold\_zero\_arithmetic\_subsequence\_amplification}
\end{theorem}
\begin{proof}
若 \(m\equiv 4\pmod 6\)，则 \(m+2\equiv 0\pmod 6\)。由推论 \ref{cor:fold-zero-half-index-multiple6} 立即得到
\[
|\mathcal Z_m|\ge F_{(m+2)/2}.
\]
把该下界分别代入定理 \ref{thm:fold-collision-zero-reduction} 与定理 \ref{thm:fold-zero-uncertainty} 的第 \(2\)、\(3\) 条，便得到所示碰撞上界、KL 上界以及最大纤维偏离上界。

最后，由 Binet 公式
\[
F_n=\frac{\varphi^n-(-\varphi^{-1})^n}{\sqrt5}
\]
可知
\[
\frac{F_{(m+2)/2}}{F_{m+2}}
=
\Theta\!\bigl(\varphi^{-(m+2)/2}\bigr)
=
\Theta(\varphi^{-m/2}),
\]
故最后一条结论成立。证毕。
\end{proof}

\begin{theorem}[六倍窗零块以平方根尺度注入稳定层，并对碰撞与相对熵给出同阶改进]\label{thm:xi-time-part60ab-window6-zero-block-halfscale-information-improvement}
\leanverified{paper\_xi\_time\_part60ab\_window6\_zero\_block\_halfscale\_information\_improvement}
沿用定理 \ref{thm:xi-fold-zero-multiple6-information-gap-improvement} 的记号。若
\[
m+2\equiv 0\pmod 6,
\]
并记
\[
d:=\frac{m+2}{2},
\qquad
|X_m|=F_{m+2},
\]
则有
\[
\frac{|\mathcal Z_m|}{|X_m|}
\ge
\frac{F_d}{F_{m+2}}
=
\varphi^{-d}\bigl(1+O(\varphi^{-2d})\bigr).
\]
特别地，
\[
|\mathcal Z_m|\asymp |X_m|^{1/2}.
\]
进一步，
\[
1-\mathsf{Col}_m\ge \frac{F_d}{F_{m+2}}
=
\varphi^{-d}\bigl(1+o(1)\bigr),
\]
并且
\[
\log|X_m|-D_{\mathrm{KL}}(p_m\|u_m)
\ge
\log\frac{F_{m+2}}{F_{m+2}-F_d}
=
\varphi^{-d}+O(\varphi^{-2d}).
\]
\end{theorem}
\begin{proof}
由定理 \ref{thm:xi-fold-zero-multiple6-information-gap-improvement}，
\[
|\mathcal Z_m|\ge F_d,
\qquad
\mathsf{Col}_m\le 1-\frac{F_d}{F_{m+2}},
\qquad
D_{\mathrm{KL}}(p_m\|u_m)\le \log(F_{m+2}-F_d).
\]
两边分别除以
\[
|X_m|=F_{m+2},
\]
便得
\[
\frac{|\mathcal Z_m|}{|X_m|}\ge \frac{F_d}{F_{m+2}},
\qquad
1-\mathsf{Col}_m\ge \frac{F_d}{F_{m+2}},
\]
以及
\[
\log|X_m|-D_{\mathrm{KL}}(p_m\|u_m)
\ge
\log\frac{F_{m+2}}{F_{m+2}-F_d}.
\]

再由 Binet 公式，
\[
F_n=\frac{\varphi^n-(-\varphi^{-1})^n}{\sqrt5},
\]
可知
\[
\frac{F_d}{F_{m+2}}
=
\varphi^{-d}\bigl(1+O(\varphi^{-2d})\bigr).
\]
由于
\[
|X_m|=F_{m+2}\asymp \varphi^{m+2}=\varphi^{2d},
\]
于是
\[
F_d\asymp \varphi^d\asymp |X_m|^{1/2},
\]
从而得到
\[
|\mathcal Z_m|\asymp |X_m|^{1/2}.
\]

最后，令
\[
x:=\frac{F_d}{F_{m+2}},
\]
则 \(x\to 0\)，并且
\[
\log\frac{F_{m+2}}{F_{m+2}-F_d}
=
-\log(1-x)
=
x+O(x^2)
=
\varphi^{-d}+O(\varphi^{-2d}).
\]
证毕。
\end{proof}

\begin{theorem}[dyadic 零余类塔触发的碰撞、KL、峰值与最大纤维偏离同步强化]\label{thm:xi-fold-zero-dyadic-tower-synchronous-information-gap-improvement}\label{thm:xi-fold-zero-multiple6-information-gap-improvement}
\leanverified{paper\_xi\_fold\_zero\_dyadic\_tower\_synchronous\_information\_gap\_improvement}
沿用定理 \ref{thm:xi-fold-zero-arithmetic-subsequence-amplification} 与定理 \ref{thm:fold-zero-uncertainty} 的记号。若
\[
m+2=3\cdot 2^r\cdot s,
\qquad
r\ge 1,
\qquad
s\ \text{为奇数},
\]
并定义
\[
\Xi_m:=\sum_{j=1}^{r}F_{(m+2)/2^j},
\]
则
\[
|\mathcal Z_m|\ge \Xi_m.
\]
因而
\[
\mathsf{Col}_m
\le
1-\frac{\Xi_m}{F_{m+2}},
\qquad
D_{\mathrm{KL}}(p_m\|u_m)
\le
\log\!\bigl(F_{m+2}-\Xi_m\bigr),
\]
\[
S_m\le 2^m\sqrt{F_{m+2}-1-\Xi_m},
\]
\[
M_m-\bar d_m
\le
2^m\sqrt{\frac{F_{m+2}-1-\Xi_m}{F_{m+2}}}.
\]
特别地，当 \(m=58\) 时，
\[
\Xi_{58}=F_{30}+F_{15}=832650,
\]
故上述四条界已经严格强于仅取半阶零块 \(F_{30}=832040\) 时得到的对应上界；而若再与推论 \ref{cor:xi-fold-zero-m58-three-coset-exhaustion} 的精确三余类分解合并，则相同四条界还可进一步提升到由 \(|\mathcal Z_{58}|=839415\) 所决定的更强数值。
\end{theorem}
\begin{proof}
由定理 \ref{thm:xi-fold-zero-dyadic-tower-disjoint-fibonacci-cosets}，
\[
|\mathcal Z_m|\ge \Xi_m.
\]
将该下界分别代入定理 \ref{thm:fold-collision-zero-reduction} 与定理
\ref{thm:fold-zero-uncertainty} 的第 \(2\)、\(3\)、\(4\) 条，即得所示碰撞上界、KL 上界、Fourier 峰值上界以及最大纤维偏离上界。

当 \(m=58\) 时，
\[
\Xi_{58}=F_{30}+F_{15}=832650
\]
由直接代入即得。最后一句只是把推论
\ref{cor:xi-fold-zero-m58-three-coset-exhaustion} 的精确值
\[
|\mathcal Z_{58}|=839415
\]
再次代入同四条不等式。证毕。
\end{proof}

\begin{proposition}[规范体积与奇偶异常的 Schur 横截分离]\label{prop:xi-time-part60abb-gauge-volume-parity-transverse-separation}
\leanverified{paper\_xi\_time\_part60abb\_gauge\_volume\_parity\_transverse\_separation}
固定分辨率 \(m\)，并设
\[
d=(d_x)_{x\in X_m},
\qquad
d'=(d_x')_{x\in X_m}
\]
是两组正整数退化谱，满足
\[
\sum_{x\in X_m}d_x=\sum_{x\in X_m}d_x'=2^m.
\]
定义
\[
G(d):=\prod_{x\in X_m}\mathfrak S_{d_x},
\qquad
H_{\mathrm{gauge}}(d):=\sum_{x\in X_m}\log(d_x!),
\qquad
N_2(d):=\dim_{\FF_2}G(d)^{\mathrm{ab}}.
\]
设两组谱具有相同的非平凡支撑
\[
\Sigma:=\{x\in X_m:\ d_x\ge 2\}=\{x\in X_m:\ d_x'\ge 2\},
\]
并且 \(d'|_\Sigma\) 在共同支撑上严格 majorize \(d|_\Sigma\)。则有
\[
N_2(d')=N_2(d)=|\Sigma|,
\qquad
H_{\mathrm{gauge}}(d')>H_{\mathrm{gauge}}(d).
\]
因此，模二奇偶异常只记录“哪些纤维是非平凡的”，而规范体积则会严格响应这些非平凡纤维内部的凸性重分配；二者在 Schur 方向上横截而不可互代。
\end{proposition}
\begin{proof}
由对称群的交换化公式
\[
\mathfrak S_n^{\mathrm{ab}}
\cong
\begin{cases}
1,& n=1,\\
\ZZ/2\ZZ,& n\ge 2,
\end{cases}
\]
立得
\[
G(d)^{\mathrm{ab}}
\cong
(\ZZ/2\ZZ)^{|\Sigma|},
\qquad
G(d')^{\mathrm{ab}}
\cong
(\ZZ/2\ZZ)^{|\Sigma|}.
\]
故
\[
N_2(d)=N_2(d')=|\Sigma|.
\]

再看规范体积。由于共同支撑外的坐标都等于 \(1\)，它们对
\[
H_{\mathrm{gauge}}(\cdot)=\sum_x\log(d_x!)
\]
没有任何差异性贡献。于是只需在 \(\Sigma\) 上比较
\[
\sum_{x\in\Sigma}\log(d_x!)
\qquad\text{与}\qquad
\sum_{x\in\Sigma}\log(d_x'!).
\]
令
\[
f(t):=\log\Gamma(t+1)
\qquad (t\ge 1).
\]
则
\[
f''(t)=\psi_1(t+1)>0,
\]
故 \(f\) 在 \([1,\infty)\) 上严格凸。由 Karamata 不等式，既然 \(d'|_\Sigma\) 严格 majorize \(d|_\Sigma\)，便有
\[
\sum_{x\in\Sigma}f(d_x')>\sum_{x\in\Sigma}f(d_x).
\]
而 \(f(n)=\log(n!)\) 对整数 \(n\ge 1\) 成立，故
\[
H_{\mathrm{gauge}}(d')>H_{\mathrm{gauge}}(d).
\]
证毕。
\end{proof}

\endinput


\noindent
上述各条结论把 bin-fold 的退化直方图同时提升到群论骨架、群胚中心态、规范热力学与 Schur 凸性横截四个层面，并把 fold 主线上的零余类约数触发机制收束为一条完整的算术子序列规律：退化谱不仅决定 \(G_m^{\mathrm{bin}}\) 的中心、交换化与导群，也决定中心极小幂等元的谱权与规范体积密度的零温偏移；另一方面，\(m\equiv 4\pmod 6\) 不仅给出一条统一的零余类放大序列，而且其 dyadic 零块塔的总量 \(\Xi_m\) 还会线性传递到碰撞、KL、Fourier 峰值与最大纤维偏离的四重上界，而对同一退化支撑上的 Schur 重分配，模二奇偶异常仅记录非平凡支撑本身，规范体积却会严格响应凸性集中度，从而进一步把“群论账本”与“体积账本”的横截不可比性固定下来。

\endinput

\paragraph{稳定层均匀纤维谱、精确尺寸偏置与中心纤维尺寸算子的两迹谱测度}
在稳定层均匀基线下的缩放退化度二原子极限、微态推前分布的两点极限以及 bin-fold 纤维群胚代数的 Wedderburn 直和都已建立之后，还可把这三条链进一步收束为一个统一的中心谱对象：均匀输出纤维谱给出基线两原子测度，微态推前恰是这一谱的精确尺寸偏置，而同一中心纤维尺寸算子在两条自然迹下的谱测度正好分别实现这两类极限。

\begin{theorem}[bin-fold 在稳定层均匀基线下的二原子谱与固定阶矩闭式]\label{thm:xi-time-part60ab4-uniform-fiber-spectrum-twoatom-moment-closure}
\leanverified{paper\_xi\_time\_part60ab4\_uniform\_fiber\_spectrum\_twoatom\_moment\_closure}
记
\[
u_m(w):=\frac{1}{|X_m|}
\qquad (w\in X_m),
\qquad
U_m\sim u_m,
\]
\[
Y_m:=\left(\frac{\varphi}{2}\right)^m d_m^{\mathrm{bin}}(U_m).
\]
则 \(Y_m\) 依分布收敛到两点随机变量
\[
Y=
\begin{cases}
1, & \text{以概率 } \varphi^{-1},\\[3pt]
\varphi^{-1}, & \text{以概率 } \varphi^{-2}.
\end{cases}
\]
并且对任意固定实数 \(q\ge 0\)，都有
\[
\lim_{m\to\infty}\mathbb E[Y_m^q]
=
\frac{1}{\varphi}+\varphi^{-q-2}.
\]
特别地，若记
\[
S_q^{\mathrm{bin}}(m):=\sum_{w\in X_m}\bigl(d_m^{\mathrm{bin}}(w)\bigr)^q,
\]
则有
\[
S_q^{\mathrm{bin}}(m)
\sim
\frac{\varphi+\varphi^{-q}}{\sqrt5}
\left(\frac{2^q}{\varphi^{q-1}}\right)^m.
\]
\end{theorem}
\begin{proof}
第一条正是定理 \ref{thm:xi-time-part70a-uniform-scaled-degeneracy-quantitative} 的极限形式。对任意固定 \(q\ge 0\)，推论 \ref{cor:xi-time-part70a-scaled-moments-leading-constant} 给出
\[
\lim_{m\to\infty}\mathbb E[Y_m^q]
=
\lim_{m\to\infty}
\frac{1}{|X_m|}
\sum_{w\in X_m}
\left(
\left(\frac{\varphi}{2}\right)^m d_m^{\mathrm{bin}}(w)
\right)^q
=
\frac{1}{\varphi}+\varphi^{-q-2}.
\]
再由
\[
\mathbb E[Y_m^q]
=
\left(\frac{\varphi}{2}\right)^{mq}\frac{S_q^{\mathrm{bin}}(m)}{|X_m|}
\]
以及
\[
|X_m|=F_{m+2}\sim \frac{\varphi^{m+2}}{\sqrt5},
\]
即可整理得到
\[
S_q^{\mathrm{bin}}(m)
\sim
\frac{\varphi+\varphi^{-q}}{\sqrt5}
\left(\frac{2^q}{\varphi^{q-1}}\right)^m.
\]
证毕。
\end{proof}

\begin{theorem}[微态推前谱是均匀基线纤维谱的精确尺寸偏置变换]\label{thm:xi-time-part60ab4-exact-sizebias-pushforward-law}
\leanverified{paper\_xi\_time\_part60ab4\_exact\_sizebias\_pushforward\_law}
记
\[
\mu_m(w):=\frac{d_m^{\mathrm{bin}}(w)}{2^m}
\qquad (w\in X_m),
\qquad
W_m\sim \mu_m,
\]
\[
Z_m:=\left(\frac{\varphi}{2}\right)^m d_m^{\mathrm{bin}}(W_m).
\]
则对任意有界函数 \(f:\RR_{\ge 0}\to\CC\)，都有精确恒等式
\[
\mathbb E[f(Z_m)]
=
\frac{\mathbb E[Y_m f(Y_m)]}{\mathbb E[Y_m]}.
\]
等价地，对任意有界函数 \(g:\RR_{\ge 0}\to\CC\)，都有
\[
\mathbb E_{\mu_m}\!\bigl[g(d_m^{\mathrm{bin}}(W_m))\bigr]
=
\frac{
\mathbb E_{u_m}\!\bigl[d_m^{\mathrm{bin}}(U_m)\,g(d_m^{\mathrm{bin}}(U_m))\bigr]
}{
\mathbb E_{u_m}\!\bigl[d_m^{\mathrm{bin}}(U_m)\bigr]
}.
\]
从而 \(Z_m\Longrightarrow Z\)，其中
\[
Z=
\begin{cases}
1, & \text{以概率 } \dfrac{\varphi}{\sqrt5},\\[6pt]
\varphi^{-1}, & \text{以概率 } \dfrac{1}{\varphi\sqrt5},
\end{cases}
\]
并且 \(Z\) 正好是 \(Y\) 的尺寸偏置变换：
\[
\mathbb P(Z\in A)
=
\frac{\mathbb E\!\bigl[Y\,\mathbf 1_{\{Y\in A\}}\bigr]}{\mathbb E[Y]}
\qquad
(A\subseteq \RR_{\ge 0}\ \text{Borel}).
\]
特别地，
\[
\mathbb E[Y]
=
\frac{1}{\varphi}+\varphi^{-3}
=
\frac{\sqrt5}{\varphi^2},
\]
且
\[
\mathbb P(Z=1)=\frac{\varphi}{\sqrt5},
\qquad
\mathbb P(Z=\varphi^{-1})=\frac{1}{\varphi\sqrt5}.
\]
\end{theorem}
\begin{proof}
由定义，
\[
\mathbb E[f(Z_m)]
=
\sum_{w\in X_m}\frac{d_m^{\mathrm{bin}}(w)}{2^m}
f\!\left(\left(\frac{\varphi}{2}\right)^m d_m^{\mathrm{bin}}(w)\right).
\]
另一方面，
\[
\mathbb E[Y_m f(Y_m)]
=
\frac{1}{|X_m|}
\sum_{w\in X_m}
\left(\frac{\varphi}{2}\right)^m
d_m^{\mathrm{bin}}(w)
f\!\left(\left(\frac{\varphi}{2}\right)^m d_m^{\mathrm{bin}}(w)\right),
\]
而纤维分割恒等式
\[
\sum_{w\in X_m}d_m^{\mathrm{bin}}(w)=2^m
\]
给出
\[
\mathbb E[Y_m]
=
\left(\frac{\varphi}{2}\right)^m\frac{2^m}{|X_m|}.
\]
因此
\[
\frac{\mathbb E[Y_m f(Y_m)]}{\mathbb E[Y_m]}
=
\sum_{w\in X_m}\frac{d_m^{\mathrm{bin}}(w)}{2^m}
f\!\left(\left(\frac{\varphi}{2}\right)^m d_m^{\mathrm{bin}}(w)\right)
=
\mathbb E[f(Z_m)],
\]
从而第一条精确恒等式成立。把 \(g(t):=f((\varphi/2)^m t)\) 代回，即得未缩放形式。

再由定理 \ref{thm:xi-time-part70a-uniform-scaled-degeneracy-quantitative} 的一致近似可知，当 \(m\) 充分大时，\(Y_m\) 几乎处处落在固定紧区间
\[
\left[\frac{\varphi^{-1}}{2},2\right]
\]
之内。于是对任意有界连续函数 \(f\)，函数 \(y\mapsto y f(y)\) 在该区间上仍为有界连续函数。结合定理 \ref{thm:xi-time-part60ab4-uniform-fiber-spectrum-twoatom-moment-closure}，得到
\[
\mathbb E[f(Z_m)]
=
\frac{\mathbb E[Y_m f(Y_m)]}{\mathbb E[Y_m]}
\longrightarrow
\frac{\mathbb E[Y f(Y)]}{\mathbb E[Y]}.
\]
故 \(Z_m\) 的极限分布就是 \(Y\) 的尺寸偏置变换。

最后，由
\[
\mathbb E[Y]
=
1\cdot \varphi^{-1}
+\varphi^{-1}\cdot \varphi^{-2}
=
\frac{\sqrt5}{\varphi^2}
\]
以及两点尺寸偏置公式
\[
\mathbb P(Z=1)
=
\frac{1\cdot \varphi^{-1}}{\mathbb E[Y]}
=
\frac{\varphi}{\sqrt5},
\qquad
\mathbb P(Z=\varphi^{-1})
=
\frac{\varphi^{-1}\cdot \varphi^{-2}}{\mathbb E[Y]}
=
\frac{1}{\varphi\sqrt5},
\]
便得到所示闭式。证毕。
\end{proof}

\begin{theorem}[中心纤维尺寸算子在两条自然迹下的谱测度识别]\label{thm:xi-time-part60ab4-central-size-operator-two-trace-spectral-identification}
沿用定义 \ref{def:fold-bin-fiber-groupoid}、命题 \ref{prop:fold-bin-groupoid-algebra-wedderburn} 与定理 \ref{thm:xi-foldbin-groupoid-central-renyi-shannon-identity} 的记号，记
\[
A_m:=\CC[\mathcal R_m^{\mathrm{bin}}]
\cong
\bigoplus_{w\in X_m}M_{d_m^{\mathrm{bin}}(w)}(\CC).
\]
定义中心纤维尺寸算子
\[
\Lambda_m:=
\bigoplus_{w\in X_m}
d_m^{\mathrm{bin}}(w)\,I_{d_m^{\mathrm{bin}}(w)}
\in Z(A_m),
\]
以及其尺度化版本
\[
\widetilde\Lambda_m:=
\left(\frac{\varphi}{2}\right)^m\Lambda_m.
\]
再定义两条迹
\[
\tau_m^{\mathrm{cen}}\!\left(\bigoplus_{w\in X_m}a_w\right)
:=
\frac{1}{|X_m|}
\sum_{w\in X_m}\operatorname{tr}_{d_m^{\mathrm{bin}}(w)}(a_w),
\]
\[
\tau_m^{\mathrm{reg}}\!\left(\bigoplus_{w\in X_m}a_w\right)
:=
\frac{1}{2^m}
\sum_{w\in X_m}\operatorname{Tr}(a_w),
\]
其中 \(\operatorname{tr}_d\) 为归一化矩阵迹；\(\tau_m^{\mathrm{reg}}\) 即定理 \ref{thm:xi-foldbin-groupoid-central-renyi-shannon-identity} 中的概率迹 \(\tau_m\)。则在最小中心幂等元 \(z_w\) 上有
\[
\tau_m^{\mathrm{cen}}(z_w)=u_m(w)=\frac{1}{|X_m|},
\qquad
\tau_m^{\mathrm{reg}}(z_w)=\mu_m(w)=\frac{d_m^{\mathrm{bin}}(w)}{2^m}.
\]
并且对任意有界函数 \(f:\RR_{\ge 0}\to\CC\)，都有精确恒等式
\[
\tau_m^{\mathrm{reg}}\!\bigl(f(\Lambda_m)\bigr)
=
\frac{
\tau_m^{\mathrm{cen}}\!\bigl(\Lambda_m f(\Lambda_m)\bigr)
}{
\tau_m^{\mathrm{cen}}(\Lambda_m)
}.
\]
此外，
\[
\tau_m^{\mathrm{cen}}\!\bigl(f(\widetilde\Lambda_m)\bigr)
=
\mathbb E[f(Y_m)],
\qquad
\tau_m^{\mathrm{reg}}\!\bigl(f(\widetilde\Lambda_m)\bigr)
=
\mathbb E[f(Z_m)].
\]
因此 \(\widetilde\Lambda_m\) 在 \(\tau_m^{\mathrm{cen}}\) 下的谱测度收敛到 \(Y\)，而在 \(\tau_m^{\mathrm{reg}}\) 下的谱测度收敛到 \(Z\)。
\end{theorem}
\begin{proof}
在第 \(w\)-块上，\(\Lambda_m\) 只是标量
\[
d_m^{\mathrm{bin}}(w)\,I_{d_m^{\mathrm{bin}}(w)}.
\]
因此对任意有界函数 \(f\)，都有
\[
f(\Lambda_m)
=
\bigoplus_{w\in X_m}
f\!\bigl(d_m^{\mathrm{bin}}(w)\bigr)\,I_{d_m^{\mathrm{bin}}(w)}.
\]
由此立即得到
\[
\tau_m^{\mathrm{cen}}(z_w)
=
\frac{1}{|X_m|},
\qquad
\tau_m^{\mathrm{reg}}(z_w)
=
\frac{d_m^{\mathrm{bin}}(w)}{2^m}.
\]
同理，
\[
\tau_m^{\mathrm{cen}}\!\bigl(f(\Lambda_m)\bigr)
=
\frac{1}{|X_m|}
\sum_{w\in X_m}
f\!\bigl(d_m^{\mathrm{bin}}(w)\bigr),
\]
\[
\tau_m^{\mathrm{cen}}\!\bigl(\Lambda_m f(\Lambda_m)\bigr)
=
\frac{1}{|X_m|}
\sum_{w\in X_m}
d_m^{\mathrm{bin}}(w)\,
f\!\bigl(d_m^{\mathrm{bin}}(w)\bigr),
\]
\[
\tau_m^{\mathrm{cen}}(\Lambda_m)
=
\frac{1}{|X_m|}
\sum_{w\in X_m}d_m^{\mathrm{bin}}(w)
=
\frac{2^m}{|X_m|},
\]
\[
\tau_m^{\mathrm{reg}}\!\bigl(f(\Lambda_m)\bigr)
=
\frac{1}{2^m}
\sum_{w\in X_m}
d_m^{\mathrm{bin}}(w)\,
f\!\bigl(d_m^{\mathrm{bin}}(w)\bigr).
\]
比较后便得到
\[
\tau_m^{\mathrm{reg}}\!\bigl(f(\Lambda_m)\bigr)
=
\frac{
\tau_m^{\mathrm{cen}}\!\bigl(\Lambda_m f(\Lambda_m)\bigr)
}{
\tau_m^{\mathrm{cen}}(\Lambda_m)
}.
\]

再由 \(\widetilde\Lambda_m=((\varphi/2)^m)\Lambda_m\)，可知
\[
\tau_m^{\mathrm{cen}}\!\bigl(f(\widetilde\Lambda_m)\bigr)
=
\frac{1}{|X_m|}
\sum_{w\in X_m}
f\!\left(\left(\frac{\varphi}{2}\right)^m d_m^{\mathrm{bin}}(w)\right)
=
\mathbb E[f(Y_m)],
\]
\[
\tau_m^{\mathrm{reg}}\!\bigl(f(\widetilde\Lambda_m)\bigr)
=
\frac{1}{2^m}
\sum_{w\in X_m}
d_m^{\mathrm{bin}}(w)\,
f\!\left(\left(\frac{\varphi}{2}\right)^m d_m^{\mathrm{bin}}(w)\right)
=
\mathbb E[f(Z_m)].
\]
最后，分别应用定理 \ref{thm:xi-time-part60ab4-uniform-fiber-spectrum-twoatom-moment-closure} 与定理 \ref{thm:xi-time-part60ab4-exact-sizebias-pushforward-law}，便得到两条谱测度收敛。证毕。
\end{proof}

\begin{theorem}[中心纤维尺寸算子的全阶矩与群胚代数维数主常数]\label{thm:xi-time-part60ab4-central-size-schatten-moment-dimension-constant}
在定理 \ref{thm:xi-time-part60ab4-central-size-operator-two-trace-spectral-identification} 的设定下，对任意固定实数 \(q\ge 0\)，都有
\[
\lim_{m\to\infty}
\tau_m^{\mathrm{cen}}\!\bigl(\widetilde\Lambda_m^q\bigr)
=
\frac{1}{\varphi}+\varphi^{-q-2},
\]
\[
\lim_{m\to\infty}
\tau_m^{\mathrm{reg}}\!\bigl(\widetilde\Lambda_m^q\bigr)
=
\frac{\varphi+\varphi^{-q-1}}{\sqrt5}.
\]
特别地，
\[
\lim_{m\to\infty}\tau_m^{\mathrm{cen}}(\widetilde\Lambda_m)
=
\frac{\sqrt5}{\varphi^2},
\qquad
\lim_{m\to\infty}\tau_m^{\mathrm{reg}}(\widetilde\Lambda_m)
=
\frac{2}{\sqrt5}.
\]
并且群胚代数维数满足精确主项
\[
\dim_{\CC}A_m
\sim
\frac{2}{\sqrt5}
\left(\frac{4}{\varphi}\right)^m.
\]
\end{theorem}
\begin{proof}
由定理 \ref{thm:xi-time-part60ab4-central-size-operator-two-trace-spectral-identification}，
\[
\tau_m^{\mathrm{cen}}\!\bigl(\widetilde\Lambda_m^q\bigr)
=
\mathbb E[Y_m^q],
\qquad
\tau_m^{\mathrm{reg}}\!\bigl(\widetilde\Lambda_m^q\bigr)
=
\mathbb E[Z_m^q].
\]
第一条极限于是直接由定理 \ref{thm:xi-time-part60ab4-uniform-fiber-spectrum-twoatom-moment-closure} 给出。

再由定理 \ref{thm:xi-time-part60ab4-exact-sizebias-pushforward-law} 的精确尺寸偏置恒等式，
\[
\mathbb E[Z_m^q]
=
\frac{\mathbb E[Y_m^{q+1}]}{\mathbb E[Y_m]}.
\]
令 \(m\to\infty\)，并代入
\[
\mathbb E[Y^{q+1}]
=
\frac{1}{\varphi}+\varphi^{-q-3},
\qquad
\mathbb E[Y]
=
\frac{\sqrt5}{\varphi^2},
\]
便得到
\[
\lim_{m\to\infty}\tau_m^{\mathrm{reg}}\!\bigl(\widetilde\Lambda_m^q\bigr)
=
\frac{\varphi+\varphi^{-q-1}}{\sqrt5}.
\]
取 \(q=1\) 即得
\[
\lim_{m\to\infty}\tau_m^{\mathrm{cen}}(\widetilde\Lambda_m)
=
\frac{\sqrt5}{\varphi^2},
\qquad
\lim_{m\to\infty}\tau_m^{\mathrm{reg}}(\widetilde\Lambda_m)
=
\frac{2}{\sqrt5}.
\]

另一方面，在 Wedderburn 分解下有
\[
\tau_m^{\mathrm{reg}}(\Lambda_m)
=
\frac{1}{2^m}
\sum_{w\in X_m}
\operatorname{Tr}\!\bigl(
d_m^{\mathrm{bin}}(w)\,I_{d_m^{\mathrm{bin}}(w)}
\bigr)
=
\frac{1}{2^m}
\sum_{w\in X_m}\bigl(d_m^{\mathrm{bin}}(w)\bigr)^2
=
\frac{\dim_{\CC}A_m}{2^m}.
\]
由于
\[
\widetilde\Lambda_m
=
\left(\frac{\varphi}{2}\right)^m\Lambda_m,
\]
故
\[
\tau_m^{\mathrm{reg}}(\Lambda_m)
=
\left(\frac{2}{\varphi}\right)^m
\tau_m^{\mathrm{reg}}(\widetilde\Lambda_m).
\]
于是
\[
\dim_{\CC}A_m
=
2^m\tau_m^{\mathrm{reg}}(\Lambda_m)
=
\left(\frac{4}{\varphi}\right)^m
\tau_m^{\mathrm{reg}}(\widetilde\Lambda_m),
\]
再代入
\[
\tau_m^{\mathrm{reg}}(\widetilde\Lambda_m)\longrightarrow \frac{2}{\sqrt5},
\]
便得到
\[
\dim_{\CC}A_m
\sim
\frac{2}{\sqrt5}
\left(\frac{4}{\varphi}\right)^m.
\]
证毕。
\end{proof}

因此，bin-fold 的稳定层均匀谱 \(Y\)、其尺寸偏置像 \(Z\) 与纤维群胚代数中的中心纤维尺寸算子 \(\widetilde\Lambda_m\) 之间并不存在三条彼此独立的极限律：前二者正是后一对象在中心均匀迹与概率迹下的两类谱测度，而 \(\dim_{\CC}A_m\) 的主常数 \(2/\sqrt5\) 则恰好等于该概率谱的一阶矩。

\paragraph{稳定层自反中心、Fourier 相位锁定与 boundary 中心的超选择分裂}
在 bin-fold 稳定层的群环生成函数、Fourier 乘积分解、规范群熵账本以及 window-\(6\) 的 boundary 中心直和结构都已建立之后，中心投影型读出与保纤维相位型读出还可继续压缩为同一条代数骨架：稳定层多重度谱不仅具有严格的 Fibonacci 自反中心，而且其 Fourier 相位被该中心完全锁定；与此同时，规范群 bits 的全局极值由严格离散凸性唯一决定，而 boundary sheet-flip 在中心层上则强制出现八个不可相混的超选择扇区。

\begin{theorem}[稳定层多重度谱的自反中心]\label{thm:xi-time-part60ab2-stable-spectrum-selfreciprocal-center}
\leanverified{paper\_xi\_time\_part60ab2\_stable\_spectrum\_selfreciprocal\_center}
记
\[
M:=F_{m+2},
\qquad
D_m(x):=\sum_{r=0}^{M-1}d_m(r)x^r\in \CC[x]/(x^M-1),
\]
并定义
\[
\Sigma_m:=\sum_{j=1}^{m}F_{j+1}=F_{m+3}-2\equiv F_{m+1}-2\pmod M.
\]
则在 \(\CC[x]/(x^M-1)\) 中成立严格自反恒等式
\[
D_m(x)\equiv x^{\Sigma_m}D_m(x^{-1})\pmod{x^M-1}.
\]
等价地，
\[
d_m(r)=d_m(\Sigma_m-r)
\qquad (\forall r\in \ZZ/(M\ZZ)).
\]
\end{theorem}
\begin{proof}
由命题 \ref{prop:fold-multiplicity-group-algebra}，
\[
D_m(x)\equiv \prod_{j=1}^{m}\bigl(1+x^{F_{j+1}}\bigr)\pmod{x^M-1}.
\]
于是
\[
\begin{aligned}
x^{\Sigma_m}D_m(x^{-1})
&\equiv
x^{\sum_{j=1}^{m}F_{j+1}}
\prod_{j=1}^{m}\bigl(1+x^{-F_{j+1}}\bigr)\\
&=
\prod_{j=1}^{m}\bigl(x^{F_{j+1}}+1\bigr)
\equiv
D_m(x)
\pmod{x^M-1}.
\end{aligned}
\]
比较模 \(x^M-1\) 下的系数便得
\[
d_m(r)=d_m(\Sigma_m-r).
\]
证毕。
\end{proof}

\begin{theorem}[Fourier 相位的刚性锁定]\label{thm:xi-time-part60ab2-stable-spectrum-fourier-phase-locking}
\leanverified{paper\_xi\_time\_part60ab2\_stable\_spectrum\_fourier\_phase\_locking}
沿用定理 \ref{thm:xi-time-part60ab2-stable-spectrum-selfreciprocal-center} 的记号。对每个
\[
k\in \ZZ/(M\ZZ)
\]
定义 Fourier 模
\[
\widehat d_m(k):=\sum_{r=0}^{M-1}d_m(r)e^{-2\pi i kr/M}.
\]
则对所有 \(k\) 都有
\[
\widehat d_m(k)=e^{-2\pi i k\Sigma_m/M}\,\overline{\widehat d_m(k)}.
\]
从而
\[
e^{\pi i k\Sigma_m/M}\widehat d_m(k)\in \RR.
\]
亦即每一个 Fourier 模的相位都被 \(\Sigma_m\) 强制锁定到一条固定直径上，而不再保留一般的复平面旋转自由度。
\end{theorem}
\begin{proof}
由定理 \ref{thm:xi-time-part60ab2-stable-spectrum-selfreciprocal-center}，
\[
d_m(r)=d_m(\Sigma_m-r).
\]
故
\[
\begin{aligned}
\widehat d_m(k)
&=
\sum_{r=0}^{M-1}d_m(\Sigma_m-r)e^{-2\pi i kr/M}\\
&=
e^{-2\pi i k\Sigma_m/M}
\sum_{r=0}^{M-1}d_m(r)e^{2\pi i kr/M}\\
&=
e^{-2\pi i k\Sigma_m/M}\,\overline{\widehat d_m(k)},
\end{aligned}
\]
其中最后一步用到 \(d_m(r)\in \RR\)。于是
\[
e^{\pi i k\Sigma_m/M}\widehat d_m(k)
=
\overline{e^{\pi i k\Sigma_m/M}\widehat d_m(k)},
\]
故该量为实数。证毕。
\end{proof}

\begin{theorem}[精确暗模的算术判别与模 \texorpdfstring{$3$}{3} 出现律]\label{thm:xi-time-part60ab2-exact-dark-mode-arithmetic-criterion}
\leanverified{paper\_xi\_time\_part60ab2\_exact\_dark\_mode\_arithmetic\_criterion}
沿用前述记号。则对任意 \(k\in \ZZ/(M\ZZ)\)，有
\[
\widehat d_m(k)=0
\]
当且仅当 \(M\) 为偶数且存在
\[
j\in\{1,\dots,m\}
\]
使
\[
kF_{j+1}\equiv \frac{M}{2}\pmod M.
\]
因此：
\begin{enumerate}
  \item 若 \(M=F_{m+2}\) 为奇数，则不存在任何精确暗模；
  \item 若 \(M\) 为偶数，则
  \[
  \widehat d_m(M/2)=0;
  \]
  \item 由 Fibonacci 奇偶周期为 \(3\)，有
  \[
  F_{m+2}\ \text{为偶数}
  \iff
  m\equiv 1\pmod 3,
  \]
  从而
  \[
  m\equiv 1\pmod 3
  \iff
  \widehat d_m(M/2)=0.
  \]
\end{enumerate}
\end{theorem}
\begin{proof}
由定理 \ref{thm:fold-spectrum-factorization}，
\[
\widehat d_m(k)=\prod_{j=1}^{m}\left(1+e^{-2\pi i kF_{j+1}/M}\right).
\]
故 \(\widehat d_m(k)=0\) 当且仅当某一因子为零，即
\[
1+e^{-2\pi i kF_{j+1}/M}=0
\]
对某个 \(j\) 成立。该式等价于
\[
e^{-2\pi i kF_{j+1}/M}=-1.
\]
而 \(-1\) 作为 \(M\)-次单位根出现当且仅当 \(M\) 为偶数；在此情形下，上式又等价于
\[
kF_{j+1}\equiv \frac{M}{2}\pmod M.
\]
这就证明了首条判别。

若 \(M\) 为偶数，则取 \(j=1\) 且 \(F_2=1\)，立得
\[
\frac{M}{2}F_2\equiv \frac{M}{2}\pmod M,
\]
故 \(\widehat d_m(M/2)=0\)。

最后，由 Fibonacci 数列模 \(2\) 的周期为 \(3\)，\(F_n\) 为偶数当且仅当 \(3\mid n\)。代入 \(n=m+2\) 即得
\[
F_{m+2}\ \text{为偶数}
\iff
3\mid (m+2)
\iff
m\equiv 1\pmod 3.
\]
证毕。
\end{proof}

\begin{theorem}[规范群 bits 泛函的全局极小值与极大值]\label{thm:xi-time-part60ab2-gauge-bits-global-extrema}
\leanverified{paper\_xi\_time\_part60ab2\_gauge\_bits\_global\_extrema}
令
\[
N:=2^m,
\qquad
L:=|X_m|=F_{m+2}.
\]
在所有满足
\[
(d_1,\dots,d_L)\in \NN^{L},
\qquad
\sum_{i=1}^{L}d_i=N
\]
的正整数向量上，定义
\[
\mathcal H(d_1,\dots,d_L):=\sum_{i=1}^{L}\log_2(d_i!).
\]
若记
\[
N=qL+r,
\qquad
0\le r<L,
\]
则 \(\mathcal H\) 的全局极小值与极大值分别为
\[
\mathcal H_{\min}
=(L-r)\log_2(q!)+r\log_2((q+1)!),
\]
\[
\mathcal H_{\max}
=
\log_2((N-L+1)!).
\]
并且：
\begin{enumerate}
  \item 极小分布在置换意义下唯一，恰为
  \[
  (d_i)\sim
  (\underbrace{q,\dots,q}_{L-r},
  \underbrace{q+1,\dots,q+1}_{r});
  \]
  \item 极大分布在置换意义下唯一，恰为
  \[
  (d_i)\sim (N-L+1,1,\dots,1).
  \]
\end{enumerate}
特别地，对 bin-fold 的实际纤维多重度谱 \(\bigl(d_m^{\mathrm{bin}}(w)\bigr)_{w\in X_m}\)，其规范群 bits
\[
H_m^{\mathrm{gauge}}=\sum_{w\in X_m}\log_2\bigl(d_m^{\mathrm{bin}}(w)!\bigr)
\]
必落在上述两端显式包络之间。
\end{theorem}
\begin{proof}
令
\[
f(n):=\log_2(n!)
\qquad (n\in\NN).
\]
则
\[
f(n+1)-f(n)=\log_2(n+1)
\]
严格递增，因此 \(f\) 在正整数上严格离散凸。

先证极小值。若某一向量中存在 \(a\ge b+2\)，则
\[
f(a)+f(b)-f(a-1)-f(b+1)
=
\log_2 a-\log_2(b+1)>0.
\]
故把一单位质量从较大坐标移向较小坐标会严格减小 \(\mathcal H\)。反复执行这一调平操作，直至所有坐标之差至多为 \(1\)，便得到唯一可能的极小构型，即全部坐标只取
\[
q,\qquad q+1
\]
两值，其中恰有 \(r\) 个取 \(q+1\)。代入即得 \(\mathcal H_{\min}\)。

再证极大值。若存在两个坐标 \(a\ge b\ge 2\)，则
\[
f(a+1)+f(b-1)-f(a)-f(b)
=
\log_2(a+1)-\log_2 b>0.
\]
故把一单位质量从较小的非平凡坐标移向较大的坐标会严格增大 \(\mathcal H\)。反复执行该集中操作，直到至多只有一个坐标大于 \(1\)，便得到唯一可能的极大构型
\[
(N-L+1,1,\dots,1),
\]
其函数值正是
\[
\log_2((N-L+1)!).
\]
证毕。
\end{proof}

\leanverified{paper\_xi\_time\_part60ab2\_window6\_boundary\_eight\_superselection}
\begin{theorem}[window-\texorpdfstring{$6$}{6} boundary 中心的八重超选择分解]\label{thm:xi-time-part60ab2-window6-boundary-eight-superselection}
取 \(m=6\)，并记
\[
X_6^{\mathrm{bdry}}=\{100001,100101,101001\}.
\]
设
\[
C_{\mathrm{bdry}}
:=
\langle \tau_w:\ w\in X_6^{\mathrm{bdry}}\rangle
=
Z_{\partial}
\cong
(\ZZ/2\ZZ)^3
\subset Z(G_6^{\mathrm{bin}}).
\]
对每个 \(\chi\in \widehat{C_{\mathrm{bdry}}}\)，定义
\[
e_\chi:=2^{-3}\sum_{g\in C_{\mathrm{bdry}}}\chi(g)\,g
\in \CC[C_{\mathrm{bdry}}].
\]
则
\[
\CC[C_{\mathrm{bdry}}]\cong \CC^{8}
\]
具有八个两两正交的极小中心幂等元 \(e_\chi\)。进一步，对任意在 \(\CC[C_{\mathrm{bdry}}]\) 上忠实的 \(\CC[G_6^{\mathrm{bin}}]\)-模 \(V\)，都有规范直和分解
\[
V=\bigoplus_{\chi\in \widehat{C_{\mathrm{bdry}}}}V_\chi,
\qquad
V_\chi:=e_\chi V,
\]
并且所有 \(V_\chi\) 都非零。换言之，boundary 中心强制出现八个非空的超选择扇区。
\end{theorem}
\begin{proof}
命题 \ref{prop:xi-window6-boundary-eight-sector-central-decomposition} 已给出 \(\{e_\chi\}_{\chi\in\widehat{C_{\mathrm{bdry}}}}\) 的两两正交性、极小中心幂等元性质以及
\[
\sum_{\chi\in\widehat{C_{\mathrm{bdry}}}}e_\chi=1.
\]
因此对任意 \(\CC[G_6^{\mathrm{bin}}]\)-模 \(V\)，都有
\[
V=\left(\sum_{\chi}e_\chi\right)V=\sum_{\chi}e_\chi V,
\]
而由 \(e_\chi e_{\chi'}=0\)（\(\chi\neq\chi'\)）知此和实为直和，故
\[
V=\bigoplus_{\chi\in\widehat{C_{\mathrm{bdry}}}}V_\chi.
\]

现在只需证明忠实性蕴含每个分量非零。若对某个 \(\chi\) 有
\[
V_\chi=e_\chi V=0,
\]
则 \(e_\chi\) 在 \(V\) 上作用为零。由于 \(e_\chi\neq 0\) 且 \(e_\chi\in \CC[C_{\mathrm{bdry}}]\)，这说明 \(\CC[C_{\mathrm{bdry}}]\to \End_{\CC}(V)\) 的限制映射具有非零核，与其忠实性矛盾。故所有 \(V_\chi\) 都非零。证毕。
\end{proof}

\begin{corollary}[window-\texorpdfstring{$6$}{6} parity-charge 的中心--体相分裂比例]\label{cor:xi-time-part60ab2-window6-central-bulk-parity-splitting-ratio}
在交换化口径下有规范直和分解
\[
\bigl(G_6^{\mathrm{bin}}\bigr)^{\mathrm{ab}}
\cong
(\ZZ/2\ZZ)^3\oplus (\ZZ/2\ZZ)^{18},
\]
其中第一因子恰由
\[
100001,\qquad 100101,\qquad 101001
\]
三条 boundary parity 所生成的中心像给出，而第二因子由非中心 bulk parity 组成。特别地，window-\(6\) 的 parity-charge 渠道中，中心稳健者所占比例恰为
\[
\frac{3}{21}=\frac17,
\]
其余
\[
\frac{18}{21}=\frac67
\]
都属于非中心通道。
\end{corollary}
\begin{proof}
定理 \ref{thm:xi-window6-boundary-parity-canonical-direct-summand} 已给出规范直和分解
\[
\bigl(G_6^{\mathrm{bin}}\bigr)^{\mathrm{ab}}
\cong
C_{\partial}\oplus (\ZZ/2\ZZ)^{18}
\cong
(\ZZ/2\ZZ)^3\oplus(\ZZ/2\ZZ)^{18},
\]
其中
\[
C_{\partial}=(\ZZ/2\ZZ)^{X_6^{\mathrm{bdry}}}
\]
正是 boundary 三条中心 sheet-flip 在交换化中的像。于是三维因子给出全部中心稳健 parity-channel，余下十八维因子给出非中心 bulk parity-channel。比例公式随之立即成立。证毕。
\end{proof}

\noindent
由此，稳定层的代数读出与 window-\(6\) 的中心结构便被压缩到同一组可核验对象之中：一方面，多重度剖面的自反中心把可见条纹的对称轴固定为显式的 Fibonacci 余类，Fourier 相位则被该中心锁定到一条一维直径；另一方面，规范群 bits 的极值几何由严格离散凸性完全支配，而 boundary sheet-flip 在中心层上又把保纤维相位型读出严格分解为八个彼此不可相混的超选择扇区。

\endinput

\paragraph{window-\texorpdfstring{$6$}{6} 自由商障碍、boundary 中心残余与有限加法账本的不可能性}
在 window-\(6\) 的纤维退化多值性、boundary 中心直和结构以及有限秩加法账本障碍已经确立之后，还可把“全局自由对称解释失败”与“局域中心残余仍然存在”压缩为如下两条直接后果。

\begin{theorem}[自由对称解释与有限加法账本的联合不可能性]\label{thm:xi-time-part60ab3-free-symmetry-finite-ledger-impossible}
不存在任何三元组 \((H,L,\Phi)\) 同时满足下列条件：
\begin{enumerate}
  \item \(H\) 为有限群，并且对每个 \(m\ge 1\)，都存在一个 \(H\) 在
  \[
  \Omega_m:=\{0,1\}^m
  \]
  上的自由作用，使其轨道划分恰等于 \(\Fold_m^{\mathrm{bin}}\) 的纤维划分；
  \item \(L\) 为有限生成交换群；
  \item \(\Phi:(\NN_{>0},\times)\to(\ZZ\oplus L,+)\) 为单射幺半群同态。
\end{enumerate}
\leanverified{paper\_xi\_time\_part60ab3\_free\_symmetry\_finite\_ledger\_impossible}
\end{theorem}
\begin{proof}
若第 \(1\) 条成立，则取 \(m=6\) 即得到 \(\Fold_6^{\mathrm{bin}}\) 的纤维划分可由某个有限群的自由轨道划分实现。然而推论 \ref{cor:foldbin6-degeneracy-multivalued-nonfree} 已证明这在 window-\(6\) 上不可能，因为退化度直方图
\[
2:8,\qquad 3:4,\qquad 4:9
\]
具有多值轨道大小，不能来自任何自由有限群作用。

另一方面，第 \(2\) 条与第 \(3\) 条的组合又被定理 \ref{thm:killo-godel-compression-not-finite-rank-homologizable} 直接排除：不存在有限生成交换群 \(L\) 使
\[
\Phi:(\NN_{>0},\times)\to(\ZZ\oplus L,+)
\]
成为单射幺半群同态。故上述三元组不可能存在。证毕。
\end{proof}

\begin{proposition}[window-\texorpdfstring{$6$}{6} 的 boundary 中心残余]\label{prop:xi-time-part60ab3-window6-boundary-center-residual}
\leanverified{paper\_xi\_time\_part60ab3\_window6\_boundary\_center\_residual}
记
\[
X_6^{\mathrm{bdry}}=\{100001,100101,101001\}.
\]
则这三条 boundary 词都对应二点纤维 \(d_6^{\mathrm{bin}}(w)=2\)。对每个 \(w\in X_6^{\mathrm{bdry}}\)，记 \(\tau_w\) 为该二点纤维上的唯一非平凡对合。则
\[
C_{\mathrm{bdry}}
:=
\langle \tau_w:\ w\in X_6^{\mathrm{bdry}}\rangle
\cong
(\ZZ/2\ZZ)^3
\hookrightarrow
Z(G_6^{\mathrm{bin}}).
\]
并且在交换化映射下，
\[
C_{\mathrm{bdry}}\longrightarrow \bigl(G_6^{\mathrm{bin}}\bigr)^{\mathrm{ab}}
\]
的像构成一个规范三维坐标子格，亦即 boundary parity 的中心残余在 abelianization 中仍保持独立可辨。
\end{proposition}
\begin{proof}
推论 \ref{cor:fold-bin-boundary-center-z2} 已把三条 boundary 二点纤维上的 sheet-flip 识别为规范中心嵌入
\[
(\ZZ/2\ZZ)^3\hookrightarrow Z(G_6^{\mathrm{bin}}),
\]
其中三条生成元正对应于 \(X_6^{\mathrm{bdry}}\) 中的三个 \(w\)。这就给出第一条结论。

再由定理 \ref{thm:xi-window6-boundary-parity-canonical-direct-summand}，
\[
\bigl(G_6^{\mathrm{bin}}\bigr)^{\mathrm{ab}}
\cong
(\ZZ/2\ZZ)^3\oplus (\ZZ/2\ZZ)^{18},
\]
其中第一因子恰由 boundary 三条中心 sheet-flip 的像生成。故 \(C_{\mathrm{bdry}}\) 在交换化中的像彼此独立，并形成一个规范三维坐标子格。证毕。
\end{proof}

\noindent
由此，window-\(6\) 的多对一折叠既不能被解释为自由有限对称的纯商，也不能被有限秩加法账本完全线性化；但这并不意味着群论痕迹全部消失。恰相反，boundary 三条二点纤维仍在中心层留下一个可规范识别的 \((\ZZ/2\ZZ)^3\) 残余，并在 parity-charge 的交换化通道中保持三维独立坐标。

\endinput

\paragraph{bin-fold 群胚代数的 Morita 刚性、规范群剖面与最大纤维饱和有限性}
在 bin-fold 的纤维群胚 Wedderburn 分解、碰撞--Parseval 维数量、纤维置换分解以及最大纤维组合上界的阈值判据已经建立之后，还可把同一套退化谱进一步压缩为下列四条对象层结论。

\begin{theorem}[bin-fold 纤维群胚代数的 Morita--\(K\) 理论、中心计数与精确维数下界]\label{thm:xi-foldbin-groupoid-morita-k-center-dimension-lb}\leanverified{paper\_xi\_foldbin\_groupoid\_morita\_k\_center\_dimension\_lb}
沿用定义 \ref{def:fold-bin-fiber-groupoid}、命题 \ref{prop:fold-bin-groupoid-algebra-wedderburn} 以及定理 \ref{thm:fold-collision-spectrum} 的 Fourier 记号，记
\[
A_m:=\CC[\mathcal R_m^{\mathrm{bin}}]
\cong
\bigoplus_{w\in X_m}M_{d_m^{\mathrm{bin}}(w)}(\CC).
\]
则 \(A_m\) 作为有限维复代数与 \(\CC^{F_{m+2}}\) Morita 等价；在自然的 \(\ast\)-结构下将其视为有限维 \(C^\ast\)-代数，则
\[
K_0(A_m)\cong \ZZ^{F_{m+2}},
\qquad
K_1(A_m)=0.
\]
并且
\[
Z(A_m)\cong \CC^{F_{m+2}},
\qquad
\dim_{\CC}Z(A_m)
=
\#\operatorname{Irr}(A_m)
=
\#\operatorname{Prim}(A_m)
=
F_{m+2}.
\]
此外，
\[
\dim_{\CC}A_m
=
\sum_{w\in X_m}\bigl(d_m^{\mathrm{bin}}(w)\bigr)^2
=
\frac{1}{F_{m+2}}
\sum_{k\in \ZZ/(F_{m+2}\ZZ)}
\bigl|\widehat d_m(k)\bigr|^2,
\]
并满足精确下界
\[
\dim_{\CC}A_m\ge \frac{2^{2m}}{F_{m+2}}.
\]
等号成立当且仅当
\[
d_m^{\mathrm{bin}}(w)\equiv \frac{2^m}{F_{m+2}}
\qquad (w\in X_m),
\]
亦即全部纤维完全平衡。
\end{theorem}
\begin{proof}
由命题 \ref{prop:fold-bin-groupoid-algebra-wedderburn}，
\[
A_m\cong \bigoplus_{w\in X_m}M_{d_m^{\mathrm{bin}}(w)}(\CC).
\]
每个矩阵块 \(M_d(\CC)\) 与 \(\CC\) Morita 等价，故其有限直和与
\[
\bigoplus_{w\in X_m}\CC\cong \CC^{|X_m|}=\CC^{F_{m+2}}
\]
Morita 等价。
有限维 \(C^\ast\)-代数的 \(K\)-理论在 Morita 等价下不变，而
\[
K_0(M_d(\CC))\cong \ZZ,
\qquad
K_1(M_d(\CC))=0.
\]
对有限直和逐块相加，便得到
\[
K_0(A_m)\cong \ZZ^{F_{m+2}},
\qquad
K_1(A_m)=0.
\]

中心由块对角形式直接给出
\[
Z(A_m)\cong \bigoplus_{w\in X_m}\CC\cong \CC^{F_{m+2}}.
\]
因此中心维数、简单模个数与原始理想个数都恰等于块数 \(F_{m+2}\)。

另一方面，
\[
\dim_{\CC}A_m
=
\sum_{w\in X_m}\dim_{\CC}M_{d_m^{\mathrm{bin}}(w)}(\CC)
=
\sum_{w\in X_m}\bigl(d_m^{\mathrm{bin}}(w)\bigr)^2.
\]
再由定理 \ref{thm:fold-collision-spectrum} 的 Parseval 恒等式与
\(
|X_m|=F_{m+2}
\)
可得
\[
\sum_{w\in X_m}\bigl(d_m^{\mathrm{bin}}(w)\bigr)^2
=
\frac{1}{F_{m+2}}
\sum_{k\in \ZZ/(F_{m+2}\ZZ)}
\bigl|\widehat d_m(k)\bigr|^2.
\]

最后，由 Cauchy--Schwarz，
\[
\sum_{w\in X_m}\bigl(d_m^{\mathrm{bin}}(w)\bigr)^2
\ge
\frac{\bigl(\sum_w d_m^{\mathrm{bin}}(w)\bigr)^2}{|X_m|}
=
\frac{2^{2m}}{F_{m+2}},
\]
因为
\(
\sum_w d_m^{\mathrm{bin}}(w)=2^m
\)。
取等当且仅当全部 \(d_m^{\mathrm{bin}}(w)\) 相等，亦即纤维完全平衡。证毕。
\end{proof}

\begin{theorem}[bin-fold 规范群的群阶、交换化与中心的完整闭式]\label{thm:xi-foldbin-gauge-group-order-abel-center-closed}
\leanverified{paper\_xi\_foldbin\_gauge\_group\_order\_abel\_center\_closed}
令
\[
G_m^{\mathrm{bin}}:=\Aut(\Fold_m^{\mathrm{bin}}).
\]
则
\[
|G_m^{\mathrm{bin}}|
=
\prod_{w\in X_m}d_m^{\mathrm{bin}}(w)!.
\]
并且若记
\[
N_m:=\#\{w\in X_m:\ d_m^{\mathrm{bin}}(w)\ge 2\},
\qquad
M_m^{(2)}:=\#\{w\in X_m:\ d_m^{\mathrm{bin}}(w)=2\},
\]
则有自然同构
\[
\bigl(G_m^{\mathrm{bin}}\bigr)^{\mathrm{ab}}
\cong
(\ZZ/2\ZZ)^{N_m},
\qquad
Z\!\bigl(G_m^{\mathrm{bin}}\bigr)
\cong
(\ZZ/2\ZZ)^{M_m^{(2)}}.
\]
特别地，在 \(m=6\) 时，由直方图 \(2:8,\ 3:4,\ 4:9\) 得
\[
\bigl(G_6^{\mathrm{bin}}\bigr)^{\mathrm{ab}}\cong (\ZZ/2\ZZ)^{21},
\qquad
Z\!\bigl(G_6^{\mathrm{bin}}\bigr)\cong (\ZZ/2\ZZ)^8.
\]
\end{theorem}
\begin{proof}
由命题 \ref{prop:fold-bin-gauge-decomposition}，
\[
G_m^{\mathrm{bin}}
\cong
\prod_{w\in X_m}\mathfrak S_{d_m^{\mathrm{bin}}(w)}.
\]
故群阶立刻为
\[
|G_m^{\mathrm{bin}}|
=
\prod_{w\in X_m}\bigl|\mathfrak S_{d_m^{\mathrm{bin}}(w)}\bigr|
=
\prod_{w\in X_m}d_m^{\mathrm{bin}}(w)!.
\]

又因
\[
\mathfrak S_d^{\mathrm{ab}}\cong
\begin{cases}
1,& d\le 1,\\
\ZZ/2\ZZ,& d\ge 2,
\end{cases}
\qquad
Z(\mathfrak S_d)\cong
\begin{cases}
\ZZ/2\ZZ,& d=2,\\
1,& d\neq 2,
\end{cases}
\]
有限直积对交换化与中心逐因子计算，便得到
\[
\bigl(G_m^{\mathrm{bin}}\bigr)^{\mathrm{ab}}
\cong
(\ZZ/2\ZZ)^{N_m},
\qquad
Z\!\bigl(G_m^{\mathrm{bin}}\bigr)
\cong
(\ZZ/2\ZZ)^{M_m^{(2)}}.
\]

最后，当 \(m=6\) 时，由推论 \ref{cor:fold-bin-m6-fiber-hist} 的直方图 \(2:8,\ 3:4,\ 4:9\) 可知
\[
N_6=8+4+9=21,
\qquad
M_6^{(2)}=8,
\]
从而得到所述两式。证毕。
\end{proof}

\begin{theorem}[bin-fold 最大纤维组合上界的逐纤维精确饱和判据]\label{thm:xi-foldbin-fiberwise-saturation-sector-criterion}
\leanverified{paper\_xi\_foldbin\_fiberwise\_saturation\_sector\_criterion}
沿用定义 \ref{def:K-of-m} 与命题 \ref{prop:fold-bin-fiber-tail} 的记号。
设 \(m\ge 1\)，取 \(K=K(m)\) 并记 \(L:=K-m\)。
定义
\[
S_{\max}(m):=
\max\Bigl\{
\sum_{k=m+1}^{K}c_kF_{k+1}:
\ (c_{m+1},\dots,c_K)\in\{0,1\}^{L},\ c_kc_{k+1}=0
\Bigr\}.
\]
则对任意 \(w\in X_m\) 有上界
\[
d_m^{\mathrm{bin}}(w)\le
\begin{cases}
F_{L+2},& w_m=0,\\
F_{L+1},& w_m=1.
\end{cases}
\]
并且达到组合上界 \(F_{L+2}\) 当且仅当
\[
w_m=0
\quad\text{且}\quad
V_m(w)\le 2^m-1-S_{\max}(m).
\]
等价地，
\[
\{w\in X_m:\ d_m^{\mathrm{bin}}(w)=F_{L+2}\}
=
\{w\in X_m:\ w_m=0,\ V_m(w)\le 2^m-1-S_{\max}(m)\}.
\]
\end{theorem}
\begin{proof}
由命题 \ref{prop:fold-bin-fiber-tail}，纤维
\(
(\Fold_m^{\mathrm{bin}})^{-1}(w)
\)
与满足
\[
\sum_{k=m+1}^{K}c_kF_{k+1}\le 2^m-1-V_m(w),
\qquad
c_kc_{k+1}=0,
\qquad
w_mc_{m+1}=0
\]
的尾部位串一一对应。

先忽略阈值约束，只保留相邻约束 \(c_kc_{k+1}=0\)。
长度 \(L\) 的无相邻 \(1\) 二进制串总数是标准 Fibonacci 计数 \(F_{L+2}\)。
当 \(L=0\) 时，尾部只剩空串，上述结论平凡成立。
以下设 \(L\ge 1\)。
若 \(w_m=1\)，则还必须满足 \(c_{m+1}=0\)，故可选尾部退化为长度 \(L-1\) 的无相邻 \(1\) 串，总数为 \(F_{L+1}\)。
于是得到上述逐纤维上界。

若 \(w_m=0\) 且
\[
V_m(w)\le 2^m-1-S_{\max}(m),
\]
则每一条满足相邻约束的尾部串都自动满足阈值约束，因而全部 \(F_{L+2}\) 条合法尾部都被保留，故
\[
d_m^{\mathrm{bin}}(w)=F_{L+2}.
\]

反之，若 \(w_m=1\)，则上界已降为 \(F_{L+1}<F_{L+2}\)，故不可能饱和。
若 \(w_m=0\) 但
\[
V_m(w)>2^m-1-S_{\max}(m),
\]
则任取实现 \(S_{\max}(m)\) 的尾部串，便有
\[
V_m(w)+\sum_{k=m+1}^{K}c_kF_{k+1}>2^m-1,
\]
从而该尾部被阈值约束排除，所以可行尾部数严格小于 \(F_{L+2}\)。
综上，达到 \(F_{L+2}\) 的充要条件正是所述阈值判据。证毕。
\end{proof}

\begin{corollary}[最大纤维饱和的精确等价条件与仅有限次发生]\label{cor:xi-foldbin-maxfiber-saturation-finite-occurrence}
\leanverified{paper\_xi\_foldbin\_maxfiber\_saturation\_finite\_occurrence}
沿用定理 \ref{thm:xi-foldbin-fiberwise-saturation-sector-criterion} 的记号，并记
\[
d_{\max}^{\mathrm{bin}}(m):=
\max_{w\in X_m}d_m^{\mathrm{bin}}(w).
\]
则有精确等价
\[
d_{\max}^{\mathrm{bin}}(m)=F_{L+2}
\iff
S_{\max}(m)\le 2^m-1.
\]
因此，最大纤维达到组合上界 \(F_{L+2}\) 只会对有限多个 \(m\) 发生。
并且在可复核扫描窗口 \(m\le 400\) 内，
\[
\{\,m\le 400:\ d_{\max}^{\mathrm{bin}}(m)=F_{L+2}\,\}
=
\{1,2,3,4,5,7,12\},
\]
如表 \ref{tab:fold_binary_saturation_points} 所示。
\end{corollary}
\begin{proof}
由推论 \ref{cor:fold-bin-saturation}，最大纤维由 \(0^m\) 取得，而 \(0^m\) 满足
\[
(0^m)_m=0,
\qquad
V_m(0^m)=0.
\]
因此定理 \ref{thm:xi-foldbin-fiberwise-saturation-sector-criterion} 作用于 \(w=0^m\) 即给出
\[
d_{\max}^{\mathrm{bin}}(m)=F_{L+2}
\iff
S_{\max}(m)\le 2^m-1.
\]

再由注 \ref{rem:fold-bin-finite}，上述饱和条件只能发生有限次，故最大纤维达到组合上界也只能有限次发生。
最后，在 \(m\le 400\) 的审计窗口内，表 \ref{tab:fold_binary_saturation_points} 已逐项列出全部饱和点，故得到所示有限集合。证毕。
\end{proof}

\paragraph{bin-fold 群胚代数的忠实表示尺度、PI 阈值与六倍窗半尺度缺陷}
在 bin-fold 纤维群胚代数的 Morita--\(K\) 理论、中心计数、规范群交换化以及六倍窗零余类放大律都已建立之后，还可继续从同一份 Wedderburn 退化直方图抽取出下列五条对象层后果。它们分别刻画最小忠实表示尺度、Morita 类型对块大小的遗忘、PI 阈值的最大纤维判别、六倍窗信息缺口的黄金半尺度锁定，以及 window-\(6\) 上三种彼此不重合的刚性尺度。

\begin{theorem}[bin-fold 纤维群胚代数的最小忠实表示维数等于全部微态数]\label{thm:xi-time-part60aca-minimal-faithful-dimension-microstate-count}
\leanverified{paper\_xi\_time\_part60aca\_minimal\_faithful\_dimension\_microstate\_count}
沿用定义 \ref{def:fold-bin-fiber-groupoid} 与命题 \ref{prop:fold-bin-groupoid-algebra-wedderburn} 的记号，记
\[
A_m:=\CC[\mathcal R_m^{\mathrm{bin}}]
\cong
\bigoplus_{w\in X_m}M_{d_w}(\CC),
\qquad
d_w:=d_m^{\mathrm{bin}}(w).
\]
定义
\[
\mu_{\mathrm{faith}}(A_m)
:=
\min\Bigl\{
\dim_{\CC}H:\ 
\exists\ \rho:A_m\hookrightarrow \operatorname{End}_{\CC}(H)\ \text{为忠实幺代数表示}
\Bigr\}.
\]
则有严格恒等式
\[
\mu_{\mathrm{faith}}(A_m)
=
\sum_{w\in X_m} d_w
=
2^m.
\]
等价地，\(A_m\) 的最小忠实 Hilbert 载体维数恰等于全部微态总数。
\end{theorem}
\begin{proof}
对每个 \(w\in X_m\)，记 \(e_w\in A_m\) 为对应于 \(M_{d_w}(\CC)\) 这一 Wedderburn 块的中心极小幂等元。它们满足
\[
e_we_{w'}=\delta_{w,w'}e_w,
\qquad
\sum_{w\in X_m}e_w=1.
\]
任取忠实幺表示
\[
\rho:A_m\hookrightarrow \operatorname{End}_{\CC}(H).
\]
由于 \(\rho\) 单射，必有 \(\rho(e_w)\neq 0\) 对一切 \(w\) 成立。令
\[
H_w:=\rho(e_w)H.
\]
由幂等元两两正交且和为 \(1\)，得到直和分解
\[
H=\bigoplus_{w\in X_m}H_w.
\]

限制到第 \(w\) 块，得到一个非零表示
\[
M_{d_w}(\CC)\cong e_wA_m\longrightarrow \operatorname{End}_{\CC}(H_w).
\]
由于 \(M_{d_w}(\CC)\) 是简单代数，非零表示的核只能是零，故该表示自动忠实。另一方面，有限维 \(M_{d_w}(\CC)\)-模必分解为标准模 \(\CC^{d_w}\) 的若干重直和，因此
\[
\dim_{\CC}H_w\ge d_w.
\]
求和即得
\[
\dim_{\CC}H=\sum_{w\in X_m}\dim_{\CC}H_w\ge \sum_{w\in X_m}d_w.
\]

反向地，对每个 \(w\) 取标准表示
\[
M_{d_w}(\CC)\curvearrowright \CC^{d_w},
\]
再作外直和
\[
\bigoplus_{w\in X_m}\CC^{d_w},
\]
便得到 \(A_m\) 的忠实表示，其维数恰为 \(\sum_w d_w\)。因此
\[
\mu_{\mathrm{faith}}(A_m)=\sum_{w\in X_m}d_w.
\]
最后，\(\Omega_m=\{0,1\}^m\) 被各条纤维
\[
\Fold_m^{\mathrm{bin}}{}^{-1}(w)
\qquad (w\in X_m)
\]
不交分割，故
\[
\sum_{w\in X_m}d_w=|\Omega_m|=2^m.
\]
证毕。
\end{proof}

\leanverified{paper\_xi\_time\_part60aca\_morita\_type\_forgets\_block\_sizes}
\begin{theorem}[bin-fold 纤维群胚代数的 Morita 类型完全坍缩到 Fibonacci 中心秩]\label{thm:xi-time-part60aca-morita-type-forgets-block-sizes}
沿用定理 \ref{thm:xi-time-part60aca-minimal-faithful-dimension-microstate-count} 的记号。则 \(A_m\) 与
\[
\CC^{F_{m+2}}
\]
Morita 等价。进一步，在自然的 \(\ast\)-结构下将 \(A_m\) 视为有限维 \(C^\ast\)-代数，则其有序 \(K_0\) 群满足
\[
K_0(A_m)\cong \ZZ^{F_{m+2}},
\qquad
K_0(A_m)^+\cong \NN^{F_{m+2}}.
\]
因此，bin-fold 纤维群胚代数的 Morita 类型完全忘记了块大小 \((d_w)\)，只保留块数 \(|X_m|=F_{m+2}\)。
\end{theorem}
\begin{proof}
Morita 等价
\[
A_m\ \text{与}\ \CC^{F_{m+2}}
\]
以及群同构
\[
K_0(A_m)\cong \ZZ^{F_{m+2}}
\]
已由定理 \ref{thm:xi-foldbin-groupoid-morita-k-center-dimension-lb} 给出。现只需说明正锥。

由 Wedderburn 分解，
\[
A_m\cong \bigoplus_{w\in X_m}M_{d_w}(\CC)
\]
作为有限维 \(C^\ast\)-代数成立。于是
\[
K_0(A_m)\cong \bigoplus_{w\in X_m}K_0\bigl(M_{d_w}(\CC)\bigr)
\cong \ZZ^{|X_m|},
\]
其中每个直和因子由相应矩阵块中的投影类生成。对矩阵代数 \(M_d(\CC)\)，投影类恰按其秩由 \(\NN\) 参数化，因此正锥逐块相加后得到
\[
K_0(A_m)^+\cong \NN^{|X_m|}.
\]
再由 \(|X_m|=F_{m+2}\)，遂得
\[
K_0(A_m)^+\cong \NN^{F_{m+2}}.
\]
证毕。
\end{proof}

\begin{theorem}[bin-fold 纤维群胚代数的 PI 阈值由最大纤维完全决定]\label{thm:xi-time-part60aca-pideg-maxfiber-threshold}
\leanverified{paper\_xi\_time\_part60aca\_pideg\_maxfiber\_threshold}
沿用定理 \ref{thm:xi-time-part60aca-minimal-faithful-dimension-microstate-count} 的记号，并定义
\[
D_m^{\mathrm{bin}}
:=
\max_{w\in X_m}d_m^{\mathrm{bin}}(w)
=
d_{\max}^{\mathrm{bin}}(m).
\]
对 \(r\ge 1\)，记 \(2r\) 次标准多项式为
\[
s_{2r}(x_1,\dots,x_{2r})
:=
\sum_{\sigma\in \mathfrak S_{2r}}
\operatorname{sgn}(\sigma)\,
x_{\sigma(1)}\cdots x_{\sigma(2r)}.
\]
则
\[
\operatorname{PIdeg}(A_m)=D_m^{\mathrm{bin}}.
\]
等价地，
\[
A_m\models s_{2D_m^{\mathrm{bin}}},
\qquad
A_m\not\models s_{2D_m^{\mathrm{bin}}-2}.
\]
特别地，当 \(m=6\) 时，由直方图 \(2:8,\ 3:4,\ 4:9\) 得
\[
\operatorname{PIdeg}(A_6)=4,
\qquad
A_6\models s_8,
\qquad
A_6\not\models s_6.
\]
\end{theorem}
\begin{proof}
由 Wedderburn 分解，
\[
A_m\cong \bigoplus_{w\in X_m}M_{d_w}(\CC).
\]
对有限维半单复代数
\[
A\cong \bigoplus_{j=1}^{r}M_{n_j}(\CC),
\]
一个多项式恒等式在 \(A\) 上成立，当且仅当它在每个矩阵块 \(M_{n_j}(\CC)\) 上都成立。由 Amitsur--Levitzki 定理，
\[
M_{n_j}(\CC)\models s_{2n_j},
\qquad
M_{n_j}(\CC)\not\models s_{2n_j-2}.
\]
因此 \(A\) 的 PI 度数恰为 \(\max_j n_j\)。将
\[
n_j=d_w
\]
代入即得
\[
\operatorname{PIdeg}(A_m)=\max_{w\in X_m}d_w=D_m^{\mathrm{bin}}.
\]
等价的标准恒等式表述随即成立。

当 \(m=6\) 时，推论 \ref{cor:fold-bin-m6-fiber-hist} 给出不同块尺寸恰为 \(2,3,4\)，故 \(D_6^{\mathrm{bin}}=4\)。于是
\[
\operatorname{PIdeg}(A_6)=4,
\qquad
A_6\models s_8,
\qquad
A_6\not\models s_6.
\]
证毕。
\end{proof}

\begin{theorem}[六倍窗共振把碰撞与 KL 缺口锁定在黄金半尺度]\label{thm:xi-time-part60aca-sixfold-halfscale-defect}
沿用定理 \ref{thm:xi-fold-zero-arithmetic-subsequence-amplification} 的记号，特别记
\[
p_m(r):=\frac{d_m(r)}{2^m},
\qquad
u_m(r):=\frac{1}{F_{m+2}}.
\]
若
\[
m+2\equiv 0\pmod 6,
\qquad
d:=\frac{m+2}{2},
\]
则 bin-fold 的碰撞缺口与 KL 缺口满足
\[
1-\mathsf{Col}_m
\ge
\frac{F_d}{F_{m+2}}
=
\frac{1}{L_d}
=
\varphi^{-d}\bigl(1+O(\varphi^{-2d})\bigr),
\]
以及
\[
\log F_{m+2}-D_{\mathrm{KL}}(p_m\|u_m)
\ge
-\log\!\left(1-\frac{F_d}{F_{m+2}}\right)
=
\varphi^{-d}\bigl(1+O(\varphi^{-d})\bigr).
\]
等价地，
\[
1-\mathsf{Col}_m
\ge
\varphi^{-(m+2)/2}\bigl(1+O(\varphi^{-m})\bigr),
\]
\[
\log F_{m+2}-D_{\mathrm{KL}}(p_m\|u_m)
\ge
\varphi^{-(m+2)/2}\bigl(1+O(\varphi^{-m/2})\bigr).
\]
因此，在六倍窗子序列上，bin-fold 对均匀分布的接近程度不可能优于一个黄金半尺度 \(\varphi^{-m/2}\)。
\end{theorem}
\begin{proof}
由结论 \ref{con:xi-fold-window6-lucas-barrier-collision-gap}，
\[
1-\mathsf{Col}_m
\ge
\frac{F_d}{F_{2d}}
=
\frac{1}{L_d},
\]
其中 \(L_d\) 为第 \(d\) 个 Lucas 数。由 Binet 公式，
\[
L_d=\varphi^d+(-\varphi)^{-d}
=
\varphi^d\bigl(1+O(\varphi^{-2d})\bigr),
\]
从而
\[
\frac{1}{L_d}
=
\varphi^{-d}\bigl(1+O(\varphi^{-2d})\bigr).
\]
又因 \(2d=m+2\)，故第一条渐近式可改写为
\[
1-\mathsf{Col}_m
\ge
\varphi^{-(m+2)/2}\bigl(1+O(\varphi^{-m})\bigr).
\]

另一方面，条件 \(m+2\equiv 0\pmod 6\) 等价于 \(m\equiv 4\pmod 6\)。于是定理 \ref{thm:xi-fold-zero-arithmetic-subsequence-amplification} 给出
\[
D_{\mathrm{KL}}(p_m\|u_m)
\le
\log\!\bigl(F_{m+2}-F_d\bigr).
\]
因此
\[
\log F_{m+2}-D_{\mathrm{KL}}(p_m\|u_m)
\ge
\log F_{m+2}-\log\!\bigl(F_{m+2}-F_d\bigr)
=
-\log\!\left(1-\frac{F_d}{F_{m+2}}\right).
\]
再由
\[
\frac{F_d}{F_{m+2}}
=
\frac{F_d}{F_{2d}}
=
\frac{1}{L_d}
=
\varphi^{-d}\bigl(1+O(\varphi^{-2d})\bigr),
\]
以及
\[
-\log(1-x)=x+O(x^2)
\qquad(x\to 0),
\]
便得到
\[
\log F_{m+2}-D_{\mathrm{KL}}(p_m\|u_m)
\ge
\varphi^{-d}\bigl(1+O(\varphi^{-d})\bigr).
\]
把 \(d=(m+2)/2\) 代回即得所示第二条渐近式。证毕。
\end{proof}

\begin{theorem}[window-\texorpdfstring{$6$}{6} 上 Morita 秩、PI 阈值与忠实表示维数的三尺度分裂]\label{thm:xi-time-part60aca-window6-triple-rigidity-scales}
取 \(m=6\)。则
\[
A_6
\cong
M_2(\CC)^{\oplus 8}\oplus M_3(\CC)^{\oplus 4}\oplus M_4(\CC)^{\oplus 9}.
\]
从而同一非交换对象上同时出现三条互不重合的刚性尺度：
\[
\#\operatorname{Irr}(A_6)
=
\operatorname{rank}_{\ZZ}K_0(A_6)
=
21,
\qquad
\operatorname{PIdeg}(A_6)=4,
\qquad
\mu_{\mathrm{faith}}(A_6)=64.
\]
此外，
\[
A_6\models s_8,
\qquad
A_6\not\models s_6,
\qquad
\bigl(G_6^{\mathrm{bin}}\bigr)^{\mathrm{ab}}\cong (\ZZ/2\ZZ)^{21}.
\]
因此，window-\(6\) 上的 Morita--\(K\) 理论尺度、PI 尺度与忠实表示尺度并不塌缩为同一数值。
\end{theorem}
\begin{proof}
由推论 \ref{cor:fold-bin-m6-fiber-hist}，
\[
\#\{w\in X_6:\ d_6^{\mathrm{bin}}(w)=2\}=8,
\qquad
\#\{w\in X_6:\ d_6^{\mathrm{bin}}(w)=3\}=4,
\qquad
\#\{w\in X_6:\ d_6^{\mathrm{bin}}(w)=4\}=9.
\]
代入 Wedderburn 分解便得
\[
A_6
\cong
M_2(\CC)^{\oplus 8}\oplus M_3(\CC)^{\oplus 4}\oplus M_4(\CC)^{\oplus 9}.
\]

由定理 \ref{thm:xi-time-part60aca-morita-type-forgets-block-sizes}，
\[
\#\operatorname{Irr}(A_6)
=
\operatorname{rank}_{\ZZ}K_0(A_6)
=
|X_6|
=
21.
\]
由定理 \ref{thm:xi-time-part60aca-pideg-maxfiber-threshold}，
\[
\operatorname{PIdeg}(A_6)=4,
\qquad
A_6\models s_8,
\qquad
A_6\not\models s_6.
\]
再由定理 \ref{thm:xi-time-part60aca-minimal-faithful-dimension-microstate-count}，
\[
\mu_{\mathrm{faith}}(A_6)=2^6=64.
\]
等价地，直接由直方图求和也可写成
\[
\mu_{\mathrm{faith}}(A_6)
=
8\cdot 2+4\cdot 3+9\cdot 4
=
64.
\]
最后，定理 \ref{thm:xi-foldbin-gauge-group-order-abel-center-closed} 给出
\[
\bigl(G_6^{\mathrm{bin}}\bigr)^{\mathrm{ab}}\cong (\ZZ/2\ZZ)^{21},
\]
故其 \(\FF_2\)-维数与第一条尺度一致，而与后两条尺度不同。证毕。
\end{proof}

\input{subsubsec__xi-time-protocol-conclusions-part60acb-binfold-orderedk0-stable-aut-threshold}

\noindent
由此，bin-fold 纤维群胚代数上的对象层刚性已进一步闭合：Morita--\(K\) 理论只保留 Fibonacci 块数，而带单位元类的有序 \(K_0\) 又把全部块大小逐坐标恢复；稳定范畴中幸存的不变量仍只有块数，并因此对分辨率给出严格刚性；PI 理论则精确读取最大纤维，并继承黄金共振常数所强制的统一下界；同时，代数自同构群按等尺寸矩阵块分裂为内自同构 Lie 群与离散置换群的半直积，window-\(6\) 上的中心块数与异常比特数又同时锁定为 \(21\)；最后，非交换性并非渐近出现，而是在 \(m=2\) 处发生首次精确阈值跳变。因而，退化谱直方图不再只是静态组合摘要，而是同时支配有序 \(K\) 理论、稳定类型、PI 阈值、自同构群与非交换起始尺度的统一刚性对象。

\endinput


\endinput

\paragraph{容量直方图、Morita 坍缩、Bernoulli 反演、维数超载与自由商判据的统一闭合}
在连续容量曲线的完备统计性、bin-fold 纤维群胚代数的 Wedderburn 分解、fold-gauge 常数的 Stirling--Bernoulli 层级、两态一致渐近以及 window-\(6\) 的非自由群商障碍都已建立之后，还可把这些链条进一步压缩为下列五条彼此衔接的终端结论。

\begin{theorem}[连续容量曲线同时完备决定纤维直方图、规范群与群胚代数]\label{thm:xi-time-part60acb-capacity-histogram-gauge-groupoid-complete}
\leanverified{paper\_xi\_time\_part60acb\_capacity\_histogram\_gauge\_groupoid\_complete}
固定窗口 \(m\ge 1\)。记
\[
h_m(d):=\#\{w\in X_m:\ d_m^{\mathrm{bin}}(w)=d\}
\qquad(d\ge 1),
\]
\[
N_m(t):=\#\{w\in X_m:\ d_m^{\mathrm{bin}}(w)\ge t\},
\qquad
\mathcal C_m^{\mathrm{cont}}(T):=\sum_{w\in X_m}\min\bigl(d_m^{\mathrm{bin}}(w),T\bigr).
\]
则对每个整数 \(d\ge 1\) 都有
\[
h_m(d)=N_m(d)-N_m(d+1).
\]
因此，\(\mathcal C_m^{\mathrm{cont}}\) 唯一决定全部纤维直方图 \(\{h_m(d)\}_{d\ge 1}\)。进一步，
\[
G_m^{\mathrm{bin}}
\cong
\prod_{d\ge 1}\mathfrak S_d^{\,h_m(d)},
\qquad
\CC[\mathcal R_m^{\mathrm{bin}}]
\cong
\bigoplus_{d\ge 1}M_d(\CC)^{\oplus h_m(d)}.
\]
换言之，在固定窗口上，连续容量曲线同时唯一决定 bin-fold 的纤维多重度直方图、fold-gauge 群的同构型以及纤维群胚代数的 Wedderburn 分解。
\end{theorem}
\begin{proof}
定理 \ref{thm:xi-capacity-tail-histogram-moment-complete-statistic} 已给出
\[
\mathcal C_m^{\mathrm{cont}}(T)=\int_0^T N_m(t)\,\dd t,
\qquad
\partial_T^+\mathcal C_m^{\mathrm{cont}}(t)=N_m(t)\ \text{a.e.}
\]
故 \(N_m\) 由 \(\mathcal C_m^{\mathrm{cont}}\) 唯一恢复。由于 \(N_m\) 是右连续整数值阶梯函数，对每个整数 \(d\ge 1\) 有
\[
N_m(d)-N_m(d+1)
=
\#\{w:\ d_m^{\mathrm{bin}}(w)\ge d\}
-\#\{w:\ d_m^{\mathrm{bin}}(w)\ge d+1\}
=
\#\{w:\ d_m^{\mathrm{bin}}(w)=d\}
=
h_m(d).
\]
这便恢复了全部直方图。

再由命题 \ref{prop:fold-bin-gauge-decomposition}，
\[
G_m^{\mathrm{bin}}
\cong
\prod_{w\in X_m}\mathfrak S_{d_m^{\mathrm{bin}}(w)},
\]
把相同纤维大小的因子归并，即得
\[
G_m^{\mathrm{bin}}
\cong
\prod_{d\ge 1}\mathfrak S_d^{\,h_m(d)}.
\]
同理，由命题 \ref{prop:fold-bin-groupoid-algebra-wedderburn}，
\[
\CC[\mathcal R_m^{\mathrm{bin}}]
\cong
\bigoplus_{w\in X_m}M_{d_m^{\mathrm{bin}}(w)}(\CC),
\]
按相同矩阵尺寸归并便得到
\[
\CC[\mathcal R_m^{\mathrm{bin}}]
\cong
\bigoplus_{d\ge 1}M_d(\CC)^{\oplus h_m(d)}.
\]
结论成立。
\end{proof}

\begin{theorem}[bin-fold 纤维群胚代数的 Morita--\(K\) 理论类型仅由 Fibonacci 块数决定]\label{thm:xi-time-part60acb-groupoid-morita-k-stable-classification}
\leanverified{paper\_xi\_time\_part60acb\_groupoid\_morita\_k\_stable\_classification}
记
\[
A_m:=\CC[\mathcal R_m^{\mathrm{bin}}].
\]
则
\[
A_m\ \text{Morita 等价于}\ \CC^{F_{m+2}},
\qquad
A_m\otimes\mathcal K\cong \mathcal K^{\oplus F_{m+2}},
\]
其中 \(\mathcal K\) 表示可分无限维 Hilbert 空间上的紧算子代数。并且
\[
K_0(A_m)\cong \ZZ^{F_{m+2}},
\qquad
K_1(A_m)=0.
\]
若 \(m\neq n\)，则 \(A_m\) 与 \(A_n\) 不 Morita 等价。换言之，bin-fold 纤维群胚代数的 Morita 类仅由块数 \(|X_m|=F_{m+2}\) 决定，而不再记录各矩阵块的尺寸分布。
\end{theorem}
\begin{proof}
由定理 \ref{thm:xi-foldbin-groupoid-morita-k-center-dimension-lb}，
\[
A_m\cong \bigoplus_{w\in X_m}M_{d_m^{\mathrm{bin}}(w)}(\CC),
\qquad
K_0(A_m)\cong \ZZ^{F_{m+2}},
\qquad
K_1(A_m)=0.
\]
由于每个矩阵代数 \(M_d(\CC)\) 与 \(\CC\) Morita 等价，故有限直和 \(A_m\) 与
\[
\bigoplus_{w\in X_m}\CC
\cong
\CC^{F_{m+2}}
\]
Morita 等价。

再由稳定化恒等式 \(M_d(\CC)\otimes\mathcal K\cong\mathcal K\)，得到
\[
A_m\otimes\mathcal K
\cong
\bigoplus_{w\in X_m}\bigl(M_{d_m^{\mathrm{bin}}(w)}(\CC)\otimes\mathcal K\bigr)
\cong
\mathcal K^{\oplus F_{m+2}}.
\]
最后，定理 \ref{thm:xi-foldbin-groupoid-morita-k-center-dimension-lb} 同时给出
\[
Z(A_m)\cong \CC^{F_{m+2}}.
\]
Morita 等价保持中心的同构型，而 Fibonacci 数列严格递增，故 \(m\neq n\) 时 \(F_{m+2}\neq F_{n+2}\)，从而 \(Z(A_m)\not\cong Z(A_n)\)。于是 \(A_m\) 与 \(A_n\) 不可能 Morita 等价。
\end{proof}

\begin{theorem}[fold-gauge 常数列逐阶恢复全部偶 Bernoulli 数与偶值 \texorpdfstring{$\zeta(2r)$}{zeta(2r)}]\label{thm:xi-time-part60acb-gauge-constant-recovers-even-zeta}
\leanverified{paper\_xi\_time\_part60acb\_gauge\_constant\_recovers\_even\_zeta}
沿用定理 \ref{thm:fold-bin-gauge-constant-stirling-bernoulli-hierarchy} 的记号。定义
\[
\beta_r:=\left(\frac{\varphi}{2}\right)^{2r-1},
\qquad
\gamma_r:=\frac{\varphi^{-1}+\varphi^{2r-3}}{r(2r-1)}.
\]
则对每个整数 \(r\ge 1\)，都有显式极限公式
\[
B_{2r}
=
\frac{1}{\gamma_r}
\lim_{m\to\infty}
\beta_r^{-m}
\left(
\log C_m-\log(2\pi)-\sum_{j=1}^{r-1}\gamma_j B_{2j}\beta_j^m
\right).
\]
从而
\[
\zeta(2r)
=
(-1)^{r-1}\frac{(2\pi)^{2r}}{2(2r)!}\,B_{2r}.
\]
因此，单一序列 \((\log C_m)_{m\ge 1}\) 已足以逐阶唯一恢复全部偶 Bernoulli 数与全部偶值 \(\zeta(2r)\)。
\end{theorem}
\begin{proof}
由定理 \ref{thm:fold-bin-gauge-constant-stirling-bernoulli-hierarchy}，对任意固定 \(R\ge 1\) 有展开
\[
\log C_m
=
\log(2\pi)+
\sum_{r=1}^{R}
\gamma_r B_{2r}\beta_r^m
+
O\!\bigl(\beta_{R+1}^m\bigr).
\]
于是对固定的 \(r\ge 1\)，把前 \(r-1\) 个主项移到左端，便得到
\[
\log C_m-\log(2\pi)-\sum_{j=1}^{r-1}\gamma_j B_{2j}\beta_j^m
=
\gamma_r B_{2r}\beta_r^m+O(\beta_{r+1}^m).
\]
注意到
\[
0<\beta_{r+1}<\beta_r<1,
\qquad
\frac{\beta_{r+1}}{\beta_r}
=
\left(\frac{\varphi}{2}\right)^2<1.
\]
两边同乘 \(\beta_r^{-m}\) 后令 \(m\to\infty\)，便得到所示极限公式。

再由 Euler--Bernoulli 恒等式
\[
B_{2r}
=
(-1)^{r-1}\frac{2(2r)!}{(2\pi)^{2r}}\zeta(2r),
\]
即可反解出 \(\zeta(2r)\) 的表达式。结论成立。
\end{proof}

\begin{theorem}[纤维群胚代数的维数主项与 Morita 不可见的指数超载]\label{thm:xi-time-part60acb-groupoid-dimension-overload}
\leanverified{paper\_xi\_time\_part60acb\_groupoid\_dimension\_overload}
仍记
\[
A_m=\CC[\mathcal R_m^{\mathrm{bin}}]
\cong
\bigoplus_{w\in X_m}M_{d_m^{\mathrm{bin}}(w)}(\CC).
\]
则
\[
\dim_{\CC}A_m=\sum_{w\in X_m}\bigl(d_m^{\mathrm{bin}}(w)\bigr)^2.
\]
并且当 \(m\to\infty\) 时，
\[
\dim_{\CC}A_m
=
\frac{\varphi+\varphi^{-2}}{\sqrt5}
\left(\frac{4}{\varphi}\right)^m
+O(2^m).
\]
进一步，由
\[
\rank K_0(A_m)=F_{m+2}\sim \frac{\varphi^{m+2}}{\sqrt5}
\]
可得
\[
\frac{\dim_{\CC}A_m}{\rank K_0(A_m)}
=
\bigl(\varphi^{-1}+\varphi^{-4}\bigr)
\left(\frac{4}{\varphi^2}\right)^m
+o\!\left(\left(\frac{4}{\varphi^2}\right)^m\right).
\]
特别地，
\[
\frac{\dim_{\CC}A_m}{\rank K_0(A_m)}\longrightarrow +\infty,
\qquad
\frac{4}{\varphi^2}>1.
\]
因此，Morita--\(K\) 理论层面只记录 \(F_{m+2}\) 个简单块，而完整群胚代数的线性维数则以底数 \(4/\varphi\) 指数增长；二者之间存在一个 Morita 不可见的指数超载。
\end{theorem}
\begin{proof}
维数公式由 Wedderburn 分解直接给出：
\[
\dim_{\CC}A_m
=
\sum_{w\in X_m}\dim_{\CC}M_{d_m^{\mathrm{bin}}(w)}(\CC)
=
\sum_{w\in X_m}\bigl(d_m^{\mathrm{bin}}(w)\bigr)^2.
\]
另一方面，定理 \ref{thm:fold-bin-two-state-asymptotic} 给出一致两态渐近
\[
d_m^{\mathrm{bin}}(w)=\varphi^{-w_m}\left(\frac{2}{\varphi}\right)^m+O(1),
\]
且
\[
\#\{w\in X_m:\ w_m=0\}=F_{m+1},
\qquad
\#\{w\in X_m:\ w_m=1\}=F_m.
\]
于是
\[
\bigl(d_m^{\mathrm{bin}}(w)\bigr)^2
=
\varphi^{-2w_m}\left(\frac{2}{\varphi}\right)^{2m}
+O\!\left(\left(\frac{2}{\varphi}\right)^m\right)+O(1),
\]
从而
\[
\dim_{\CC}A_m
=
\left(\frac{2}{\varphi}\right)^{2m}
\bigl(F_{m+1}+\varphi^{-2}F_m\bigr)
+O\!\left(F_{m+2}\left(\frac{2}{\varphi}\right)^m\right)
+O(F_{m+2}).
\]
再由 Binet 公式，
\[
F_{m+1}+\varphi^{-2}F_m
=
\frac{\varphi+\varphi^{-2}}{\sqrt5}\,\varphi^m+O(\varphi^{-m}),
\]
故主项化为
\[
\left(\frac{2}{\varphi}\right)^{2m}
\cdot
\frac{\varphi+\varphi^{-2}}{\sqrt5}\,\varphi^m
=
\frac{\varphi+\varphi^{-2}}{\sqrt5}\left(\frac{4}{\varphi}\right)^m.
\]
又因
\[
F_{m+2}\left(\frac{2}{\varphi}\right)^m=O(2^m),
\qquad
F_{m+2}=O(\varphi^m)=o(2^m),
\]
便得到
\[
\dim_{\CC}A_m
=
\frac{\varphi+\varphi^{-2}}{\sqrt5}\left(\frac{4}{\varphi}\right)^m+O(2^m).
\]

最后用定理 \ref{thm:xi-foldbin-groupoid-morita-k-center-dimension-lb} 中的
\(
K_0(A_m)\cong \ZZ^{F_{m+2}}
\)
与
\(
F_{m+2}\sim \varphi^{m+2}/\sqrt5
\)
相除，即得
\[
\frac{\dim_{\CC}A_m}{\rank K_0(A_m)}
\sim
\frac{\varphi+\varphi^{-2}}{\varphi^2}
\left(\frac{4}{\varphi^2}\right)^m
=
\bigl(\varphi^{-1}+\varphi^{-4}\bigr)
\left(\frac{4}{\varphi^2}\right)^m.
\]
结论成立。
\end{proof}

\begin{theorem}[常纤维当且仅当自由有限群作用可实现]\label{thm:xi-time-part60acb-constant-fiber-free-group-action-criterion}
\leanverified{paper\_xi\_time\_part60acb\_constant\_fiber\_free\_group\_action\_criterion}
设 \(f:Y\twoheadrightarrow X\) 为有限集之间的满射，并记纤维直方图
\[
h(d):=\#\{x\in X:\ |f^{-1}(x)|=d\}
\qquad(d\ge 1).
\]
则下列条件等价：
\begin{enumerate}
  \item 等价关系 \(y\sim y'\iff f(y)=f(y')\) 是某个有限群在 \(Y\) 上自由作用的轨道划分；
  \item 存在整数 \(D\ge 1\) 使得 \(h(d)=0\) 对一切 \(d\neq D\) 成立，即全部纤维大小恒等于同一个常数 \(D\)。
\end{enumerate}
若进一步要求该划分来自有限阿贝尔群上的线性陪集划分，则同样必须满足上述常纤维条件。

应用于 window-\(6\) 的 bin-fold
\[
F_6:=\Fold_6^{\mathrm{bin}}:\Omega_6\to X_6,
\]
其纤维直方图
\[
2:8,\qquad 3:4,\qquad 4:9
\]
同时出现三种纤维大小，故 \(F_6\) 的纤维划分既不可能由任何有限群的自由作用实现，也不可能来自任何有限阿贝尔群的线性陪集划分。
\end{theorem}
\begin{proof}
若该等价关系来自某个有限群 \(H\) 在 \(Y\) 上的自由作用，则每个轨道都与 \(H\) 自身等势，故轨道大小恒等于 \(|H|\)。而轨道正是 \(f\) 的纤维，因此所有纤维大小都相等，得到 \(1\Rightarrow 2\)。

反之，若全部纤维大小都等于同一个常数 \(D\)，则取任一阶为 \(D\) 的有限群 \(H\)（例如循环群 \(\ZZ/D\ZZ\)）。对每个 \(x\in X\)，任选一个双射
\[
f^{-1}(x)\cong H,
\]
把 \(H\) 在自身上的左正则作用搬运到纤维 \(f^{-1}(x)\) 上，并在不同纤维上彼此独立地施行。于是得到 \(H\) 在 \(Y\) 上的自由作用，其轨道恰为 \(f\) 的各条纤维，从而 \(2\Rightarrow 1\)。

若该划分来自有限阿贝尔群上的线性陪集划分，则每个等价类都是某个固定子群陪集的原像，因而大小恒等于该子群的阶，仍必为常纤维。

最后，对 window-\(6\) 而言，纤维直方图同时出现 \(2,3,4\) 三种大小，故不满足常纤维条件。于是自由有限群作用实现与阿贝尔线性陪集实现均被排除。
\end{proof}

\noindent
由此，连续容量曲线、Morita 类型、Bernoulli--\(\zeta\) 谱链、完整群胚代数维数以及自由群商可实现性这五条看似分离的对象层结论，最终都被压缩到同一套纤维多重度数据之上：容量曲线完整恢复直方图，直方图决定规范群与 Wedderburn 分块，Morita 理论进一步只保留 Fibonacci 块数，维数主项则重新显露被 Morita 理论遗忘的指数级块尺寸负载，而自由群商解释是否可能则恰由直方图是否退化为单值常纤维来判定。

\endinput

\paragraph{审计偶数窗的 Fibonacci 邻项分割与 window-\(6\) 最大纤维的算术阈值}
在审计偶数窗上的最小退化 Fibonacci 律、最小扇区补集与 Fold 极大纤维的跨模型锁定，以及 window-\(6\) 的单跳阶梯分类已经确定之后，还可把偶数窗的集合分割与最小窗口上的极大 bin-fold 纤维进一步压缩为下列两条直接结论。

\begin{theorem}[审计偶数窗中的 Fibonacci 邻项分割律]\label{thm:xi-time-part60ad-audited-even-fibonacci-neighbor-partition}
\leanverified{paper\_xi\_time\_part60ad\_audited\_even\_fibonacci\_neighbor\_partition}
对审计窗口
\[
m\in\{6,8,10,12\},
\]
沿用定义 \ref{def:fold-bin-degeneracy-sectors} 的记号，记
\[
\mathcal S_{m,d_{\min}(m)}\subset X_m
\]
为 bin-fold 的最小退化扇区，并记 Fold 在层级 \(n\) 的最大纤维基数为
\[
D_n:=\max_{x\in X_n}d_n(x).
\]
则成立精确分割
\[
|X_m|
=
\bigl|\mathcal S_{m,d_{\min}(m)}\bigr|+D_{2m-2}.
\]
更精确地，三项恰由三枚相邻 Fibonacci 数给出：
\[
\bigl|\mathcal S_{m,d_{\min}(m)}\bigr|=F_m,
\qquad
D_{2m-2}=F_{m+1},
\qquad
|X_m|=F_{m+2}.
\]
因此在这些审计窗口中，最小退化扇区与其补集并非任意分裂，而是被一对相邻 Fibonacci 数严格锁定。
\end{theorem}
\begin{proof}
由定理 \ref{thm:fold-bin-min-degeneracy-fib-audited}，
\[
d_{\min}(m)=F_{m/2},
\qquad
\bigl|\mathcal S_{m,d_{\min}(m)}\bigr|=F_m=|X_{m-2}|.
\]
另一方面，推论 \ref{cor:fold-bin-minsector-complement-maxfiber-audited} 已给出
\[
\bigl|X_m\setminus \mathcal S_{m,d_{\min}(m)}\bigr|=D_{2m-2}.
\]
由稳定类型计数律 \(|X_n|=F_{n+2}\)，遂得
\[
D_{2m-2}
=
|X_m|-|X_{m-2}|
=
F_{m+2}-F_m
=
F_{m+1}.
\]
代回便得到
\[
|X_m|
=
\bigl|\mathcal S_{m,d_{\min}(m)}\bigr|+D_{2m-2}
=
F_m+F_{m+1}
=
F_{m+2}.
\]
证毕。
\end{proof}

\begin{theorem}[window-\(6\) 最大 bin-fold 纤维的算术阈值刻画]\label{thm:xi-time-part60ad-window6-max-binfold-arithmetic-threshold}
\leanverified{paper\_xi\_time\_part60ad\_window6\_max\_binfold\_arithmetic\_threshold}
在 \(m=6\) 时，令 \(w\in X_6\) 且记
\[
V:=V_6(w).
\]
则二进制区间折叠 \(\mathrm{Fold}^{\mathrm{bin}}_6\) 的最大纤维满足
\[
d_6^{\mathrm{bin}}(w)=4
\quad\Longleftrightarrow\quad
w_6=0\ \text{且}\ V\le 8.
\]
并且一旦上述条件成立，就有精确原像公式
\[
\bigl(\mathrm{Fold}^{\mathrm{bin}}_6\bigr)^{-1}(w)
=
\{V,\ V+F_8,\ V+F_9,\ V+F_{10}\}
=
\{V,\ V+21,\ V+34,\ V+55\}.
\]
此外，
\[
\#\{\,w\in X_6:\ d_6^{\mathrm{bin}}(w)=4\,\}=9.
\]
因此，window-\(6\) 中的最大退化并不在全部稳定类型上弥散出现，而是被阈值 \(V_6(w)\le 8\) 与末位条件 \(w_6=0\) 的算术截面完全选出。
\end{theorem}
\begin{proof}
由定理 \ref{thm:xi-foldbin-lastbit-tail-spectrum-single-jump-staircase}，在 \(m=6\) 时有精确三档公式
\[
d_6^{\mathrm{bin}}(w)=
\begin{cases}
2,& w_6=1,\\[4pt]
4,& w_6=0,\ V_6(w)\le 8,\\[4pt]
3,& w_6=0,\ V_6(w)>8.
\end{cases}
\]
故
\[
d_6^{\mathrm{bin}}(w)=4
\quad\Longleftrightarrow\quad
w_6=0\ \text{且}\ V_6(w)\le 8.
\]

再由表 \ref{tab:fold6_bin_uplift_choice_collapse} 的第一行，
\[
\Delta(w):=\{\,n-V_6(w):\ n\in(\mathrm{Fold}^{\mathrm{bin}}_6)^{-1}(w)\,\}
=
\{0,21,34,55\}
\]
恰当且仅当 \(w_6=0,\ V_6(w)\le 8\)。因此在上述条件下，
\[
\bigl(\mathrm{Fold}^{\mathrm{bin}}_6\bigr)^{-1}(w)
=
V+\Delta(w)
=
\{V,\ V+21,\ V+34,\ V+55\},
\]
即所示原像公式。

同一表的第一行同时列出了全部满足该条件的稳定词，共 \(9\) 个，故
\[
\#\{\,w\in X_6:\ d_6^{\mathrm{bin}}(w)=4\,\}=9.
\]
证毕。
\end{proof}

\noindent
由此，审计偶数窗中的最小退化扇区已不仅仅给出一个局部 Fibonacci 计数，它与更高层 Fold 极大纤维一起把 \(X_m\) 精确切分为两块相邻 Fibonacci 尺度；而在最小窗口 \(m=6\) 上，最大 bin-fold 纤维又被末位与 Zeckendorf 阈值的联合算术判据唯一选出。于是，全局审计分割与局部最大退化阈值在同一条离散 Fibonacci 链上闭合。

\endinput

\paragraph{Jensen 缺陷反演、排斥半径幂律与共动缺陷剖面的 Fourier--Prony 线性化}
以第 \ref{subsec:dephysicalized-horizon-timefiber-nohair} 小节中的
Jensen 缺陷--排斥半径链、第 \ref{subsubsec:xi-horizon-spectral-measure-tomography}
小节中的共动 Cayley 层析，以及
Toeplitz/Hankel/Prony 的既有恢复定理为基础，还可把“径向零点账本”与“共动缺陷剖面”进一步压缩为下列五条专门化结论。

\begin{theorem}[Jensen 缺陷的分布反演与零点半径测度的完全恢复]\label{thm:xi-time-part60a-jensen-defect-radial-measure-inversion}
沿用定理 \ref{thm:dephys-defect-repulsion-radius} 与命题
\ref{prop:dephys-jensen-defect-log-derivative} 的记号。设
\(F_{\zeta}\) 的盘内零点按重数记为 \(\{a\}\subset\mathbb D\)，并定义
\((-\infty,0)\) 上的离散正测度
\[
\mu_{\zeta}:=
\sum_{a:\,F_{\zeta}(a)=0}\mathrm{mult}(a)\,\delta_{\log|a|}.
\]
再令
\[
u:=\log\rho<0,
\qquad
\mathcal D_{\zeta}(u):=D_{\zeta}(e^u).
\]
则对每个 \(u<0\) 都有严格积分表示
\begin{equation}\label{eq:xi-time-part60a-jensen-defect-measure-potential}
\boxed{\;
\mathcal D_{\zeta}(u)=\int_{(-\infty,u)}(u-s)\,\dd\mu_{\zeta}(s).\;}
\end{equation}
从而：
\begin{enumerate}
  \item \(\mathcal D_{\zeta}\) 在 \((-\infty,0)\) 上为凸、分段仿射函数；
  \item 几乎处处有
  \[
  \mathcal D_{\zeta}'(u)=N_{F_{\zeta}}(e^u);
  \]
  \item 在分布意义下成立
  \[
  \boxed{\;
  \mathcal D_{\zeta}''=\mu_{\zeta}.\;}
  \]
\end{enumerate}
换言之，Jensen 缺陷的分布二阶导数恰等于零点半径多重集所定义的离散测度。
\end{theorem}
\begin{proof}
由 Jensen 求和式与命题 \ref{prop:dephys-jensen-defect-log-derivative} 的记号，
\[
D_{\zeta}(\rho)=\sum_{|a|<\rho}\log\frac{\rho}{|a|}.
\]
代入 \(\rho=e^u\) 得
\[
\mathcal D_{\zeta}(u)
=
\sum_{\log|a|<u}\bigl(u-\log|a|\bigr),
\]
而右端正是 \eqref{eq:xi-time-part60a-jensen-defect-measure-potential}
中的离散势函数表示。每一项 \((u-s)_+\) 关于 \(u\) 为凸且分段仿射，故其局部有限和
\(\mathcal D_{\zeta}\) 亦如此。

在可导点，
\[
\frac{\dd}{\dd u}(u-\log|a|)_+=\mathbf 1_{\{u>\log|a|\}},
\]
于是
\[
\mathcal D_{\zeta}'(u)
=
\sum_a \mathbf 1_{\{u>\log|a|\}}
=
N_{F_{\zeta}}(e^u),
\]
这与命题 \ref{prop:dephys-jensen-defect-log-derivative} 完全一致。
再对该分段常值函数取分布导数，即在每个跳点 \(\log|a|\) 处产生权重
\(\mathrm{mult}(a)\) 的 Dirac 质量，从而
\[
\mathcal D_{\zeta}''=
\sum_a \mathrm{mult}(a)\,\delta_{\log|a|}
=\mu_{\zeta}.
\]
证毕。
\end{proof}

\begin{theorem}[排斥半径的分段幂律正规形与整数斜率层级]\label{thm:xi-time-part60a-repulsion-radius-piecewise-power-law}
沿用定理 \ref{thm:dephys-defect-repulsion-radius} 的排斥半径记号
\[
r_{\ast}(\rho):=\rho\,e^{-D_{\zeta}(\rho)}\qquad(0<\rho<1).
\]
设 \(F_{\zeta}\) 的盘内零点模长的不同取值按严格递增排列为
\[
0<\varrho_1<\varrho_2<\cdots<\varrho_J<1,
\]
并记模长恰为 \(\varrho_j\) 的零点总重数为 \(\mu_j\)，再定义累计重数
\[
M_j:=\mu_1+\cdots+\mu_j,
\qquad
M_0:=0.
\]
约定 \(\varrho_0:=0\)、\(\varrho_{J+1}:=1\)。则对每个
\(j=0,1,\dots,J\) 与每个 \(\rho\in(\varrho_j,\varrho_{j+1})\)，都有严格闭式
\begin{equation}\label{eq:xi-time-part60a-repulsion-piecewise-power-law}
\boxed{\;
r_{\ast}(\rho)=
\rho^{\,1-M_j}
\prod_{i=1}^{j}\varrho_i^{\,\mu_i}.\;}
\end{equation}
因此：
\begin{enumerate}
  \item 在每个零点模长间隙 \((\varrho_j,\varrho_{j+1})\) 上，\(r_{\ast}\) 是幂函数；
  \item 其对数斜率恰为整数 \(1-M_j\)；
  \item 尤其有
  \[
  r_{\ast}(\rho)=\rho\qquad(0<\rho<\varrho_1),
  \]
  若 \(\mu_1=1\)，则
  \[
  r_{\ast}(\rho)\equiv \varrho_1
  \qquad(\varrho_1<\rho<\varrho_2),
  \]
  而一旦 \(M_j\ge 2\)，则 \(r_{\ast}\) 在
  \((\varrho_j,\varrho_{j+1})\) 上严格递减。
\end{enumerate}
\end{theorem}
\begin{proof}
对 \(\rho\in(\varrho_j,\varrho_{j+1})\)，恰有前 \(j\) 个模长层进入 Jensen 求和，故
\[
D_{\zeta}(\rho)
=
\sum_{i=1}^{j}\mu_i\log\frac{\rho}{\varrho_i}
=
M_j\log\rho-\sum_{i=1}^{j}\mu_i\log\varrho_i.
\]
代回 \(r_{\ast}(\rho)=\rho e^{-D_{\zeta}(\rho)}\) 即得
\eqref{eq:xi-time-part60a-repulsion-piecewise-power-law}。
对数求导得到
\[
\frac{\dd}{\dd(\log\rho)}\log r_{\ast}(\rho)=1-M_j,
\]
这与推论 \ref{cor:dephys-repulsion-radius-zero-count-driving} 在该区间上
\(N_{F_{\zeta}}(\rho)=M_j\) 的表达一致。由 \(M_j\) 的整数取值，三种单调性结论立即推出。证毕。
\end{proof}

\begin{theorem}[有限 Jensen 证书的绝对盲区]\label{thm:xi-time-part60a-finite-jensen-certificate-blindspot}
固定任意有限半径采样集
\[
0<\rho_1<\rho_2<\cdots<\rho_M<1,
\]
并对单位圆盘内解析对象 \(F\) 定义 Jensen 缺陷样本向量
\[
\mathbf D_{\boldsymbol\rho}(F):=
\bigl(D_F(\rho_1),\dots,D_F(\rho_M)\bigr)\in\RR^M.
\]
则存在两个有限 Blaschke 乘积 \(B_0,B_1\) 满足：
\begin{enumerate}
  \item \(B_0\) 在 \(\mathbb D\) 内无零点；
  \item \(B_1\) 在 \(\mathbb D\) 内恰有一个零点 \(a\)，并且 \(|a|>\rho_M\)；
  \item 但二者具有完全相同的有限 Jensen 证书：
  \[
  \boxed{\;
  \mathbf D_{\boldsymbol\rho}(B_0)=\mathbf D_{\boldsymbol\rho}(B_1).\;}
  \]
\end{enumerate}
因此，任何仅依赖有限半径样本
\(\{D_{\zeta}(\rho_j)\}_{j=1}^{M}\) 且 \(\rho_M<1\) 的判定协议，都不可能决定盘内是否存在零点；从而不能作为 RH 的决定性有限证书。
\end{theorem}
\begin{proof}
取
\[
B_0(w)\equiv 1.
\]
再任取满足 \(\rho_M<|a|<1\) 的点 \(a\in\mathbb D\)，并令
\[
B_1(w):=\frac{w-a}{1-\overline a\,w}.
\]
这是一个恰有单零点 \(a\) 的有限 Blaschke 因子。对任意
\(\rho<|a|\)，Jensen 缺陷只对 \(|z|<\rho\) 内的零点求和，因此
\[
D_{B_1}(\rho)=0.
\]
另一方面，\(B_0\) 在盘内无零点，故
\[
D_{B_0}(\rho)=0.
\]
由于 \(\rho_j<|a|\) 对全部 \(j=1,\dots,M\) 都成立，遂有
\[
D_{B_0}(\rho_j)=D_{B_1}(\rho_j)=0
\qquad (1\le j\le M),
\]
即
\(
\mathbf D_{\boldsymbol\rho}(B_0)=\mathbf D_{\boldsymbol\rho}(B_1)
\)。
但二者的盘内零点性质完全不同，故任何仅依赖该有限样本向量的协议都无法判定“盘内是否有零点”。对圆盘完成化对象 \(F_{\zeta}\) 而言，这正排除了仅用有限 Jensen 样本决定 RH 的可能性。证毕。
\end{proof}

\begin{theorem}[共动缺陷剖面的 Fourier--Prony 线性化]\label{thm:xi-time-part60a-comoving-defect-fourier-prony-linearization}
\leanverified{paper\_xi\_time\_part60a\_comoving\_defect\_fourier\_prony\_linearization}
沿用 \S\ref{subsubsec:xi-horizon-spectral-measure-tomography} 的有限点缺陷模型。
设离切片零点为
\[
\rho_j=\beta_j+\mathrm{i}\gamma_j,
\qquad
\delta_j:=\frac12-\beta_j>0,
\qquad
1\le j\le N,
\]
并记对应的共动 Cayley 像为
\[
w_j(x_0):=\mathcal C\bigl((\gamma_j-x_0)+\mathrm{i}\delta_j\bigr)\in\mathbb D.
\]
定义缺陷剖面
\[
h(x_0):=\sum_{j=1}^{N}m_j\bigl(1-|w_j(x_0)|^2\bigr),
\qquad
m_j\in\NN.
\]
则其 Fourier 变换
\[
\widehat h(\xi):=\int_{\RR}h(x_0)e^{-\mathrm{i}\xi x_0}\,\dd x_0
\]
满足完全显式分解
\begin{equation}\label{eq:xi-time-part60a-comoving-defect-fourier-prony}
\boxed{\;
\widehat h(\xi)
=
4\pi\sum_{j=1}^{N}
\frac{m_j\delta_j}{1+\delta_j}\,
e^{-(1+\delta_j)|\xi|}\,
e^{-\mathrm{i}\gamma_j\xi},
\qquad \xi\in\RR.\;}
\end{equation}
特别地，取单位采样步长并定义
\[
s_n:=\widehat h(n)
\qquad (n\in\NN_0),
\]
再记
\[
\lambda_j:=e^{-(1+\delta_j)-\mathrm{i}\gamma_j},
\qquad
c_j:=4\pi\,\frac{m_j\delta_j}{1+\delta_j},
\]
则有有限指数和表示
\begin{equation}\label{eq:xi-time-part60a-comoving-defect-samples-prony}
\boxed{\;
s_n=\sum_{j=1}^{N}c_j\,\lambda_j^{\,n}.\;}
\end{equation}
因此：
\begin{enumerate}
  \item \((s_n)_{n\ge 0}\) 满足阶数至多 \(N\) 的常系数线性递推；
  \item 其无穷 Hankel 矩阵 \(H:=(s_{i+j})_{i,j\ge 0}\) 的秩不超过互异 \(\lambda_j\) 的个数；
  \item 若 \(\lambda_1,\dots,\lambda_N\) 两两不同，则 \(\operatorname{rank}(H)=N\)。
\end{enumerate}
\end{theorem}
\begin{proof}
由 \eqref{eq:xi-comoving-scan-lorentz-peak}，
\[
h(x_0)=
\sum_{j=1}^{N}
\frac{4m_j\delta_j}{(\gamma_j-x_0)^2+(1+\delta_j)^2}.
\]
对每一项取平移变量 \(t=x_0-\gamma_j\)，并应用标准积分公式
\[
\int_{\RR}\frac{e^{-\mathrm{i}\xi t}}{t^2+a^2}\,\dd t
=
\frac{\pi}{a}e^{-a|\xi|}
\qquad(a>0),
\]
即得 \eqref{eq:xi-time-part60a-comoving-defect-fourier-prony}。
再取 \(\xi=n\in\NN_0\)，便得到
\eqref{eq:xi-time-part60a-comoving-defect-samples-prony}。

有限指数和自动满足以
\[
Q(z):=\prod_{j=1}^{N}(z-\lambda_j)
\]
为特征多项式的常系数线性递推。另一方面，
\[
H=(s_{i+j})_{i,j\ge 0}
=
V\,\mathrm{diag}(c_1,\dots,c_N)\,V^{\mathsf T},
\qquad
V_{ij}:=\lambda_j^{\,i},
\]
故 \(\operatorname{rank}(H)\) 不超过互异 \(\lambda_j\) 的个数；若
\(\lambda_j\) 两两不同，则相应 Vandermonde 矩阵列满秩，从而
\(\operatorname{rank}(H)=N\)。证毕。
\end{proof}

\begin{corollary}[有限缺陷共动谱的最小完备采样阈值]\label{cor:xi-time-part60a-comoving-defect-minimal-sampling-sharp}
在定理 \ref{thm:xi-time-part60a-comoving-defect-fourier-prony-linearization}
的设定下，进一步假设 \(\lambda_1,\dots,\lambda_N\) 两两不同。则：
\begin{enumerate}
  \item 任意长度小于 \(2N\) 的初始样本段
  \[
  s_0,\dots,s_{2N-2}
  \]
  在一般位置上不足以唯一恢复全部 \((\lambda_j,c_j)_{j=1}^{N}\)；
  \item 长度恰为 \(2N\) 的初始样本段
  \[
  s_0,\dots,s_{2N-1}
  \]
  在一般位置下已足以唯一恢复 \((\lambda_j,c_j)_{j=1}^{N}\)（忽略排列）；
  \item 因而在给定高度先验窗口长度为 \(2\pi\) 的条件下，可唯一恢复
  \((\gamma_j,\delta_j,m_j)_{j=1}^{N}\)。
\end{enumerate}
\end{corollary}
\begin{proof}
将定理 \ref{thm:xi-time-part60a-comoving-defect-fourier-prony-linearization}
中的指数和表示
\(
s_n=\sum_{j=1}^{N}c_j\lambda_j^n
\)
专门化到推论 \ref{cor:xi-prony-decision-vs-reconstruction-one-sample-gap}，
取其中的节点 \(a_j=\lambda_j\)、权重 \(\omega_j=c_j\)，便得到前两条结论。

对第三条，由
\[
|\lambda_j|=e^{-(1+\delta_j)},
\qquad
\arg(\lambda_j)\equiv -\gamma_j\pmod{2\pi},
\]
可恢复
\[
\delta_j=-\log|\lambda_j|-1,
\qquad
\gamma_j\equiv-\arg(\lambda_j)\pmod{2\pi}.
\]
若事先知道 \(\gamma_j\) 落在某个固定长度为 \(2\pi\) 的区间内，则该同余类具有唯一代表。
最后由
\[
m_j=
\frac{c_j}{4\pi}\cdot\frac{1+\delta_j}{\delta_j}
\]
恢复重数 \(m_j\)。证毕。
\end{proof}

\begin{remark}[有限原子审计接口]\label{rem:xi-time-part60a-jensen-comoving-prony-audit}
脚本 \texttt{scripts/exp\_xi\_jensen\_comoving\_prony\_audit.py}
给出一个有限 Blaschke/有限缺陷样本上的逐项核验：
它同时检查 Jensen 缺陷的离散势表示、排斥半径的分段幂律、
有限半径采样的绝对盲区，以及单位步长共动剖面的 Fourier--Prony
恢复。
\end{remark}

\endinput

\paragraph{Fourier--Joukowsky 干涉曲线、奇异环接触计数与 Fibonacci 采样谱}
在共动缺陷剖面的 Fourier--Prony 线性化、碰撞谱的单位圆采样表达以及 Joukowsky 门的既有解析接口都已建立之后，还可把若干看似分离的平面参数曲线统一理解为“单位圆旋转轨道经 Joukowsky 门读出”的有限干涉对象，并把这一几何接口直接接到 Fibonacci 生成多项式 \(D_m\) 的碰撞谱链上。

\begin{theorem}[Fourier--Joukowsky 生成原理：玫瑰线与 Lissajous 曲线作为概率特征曲线的单位圆外推]\label{thm:xi-time-part60ae-fourier-joukowsky-generation-principle}
\leanverified{paper\_xi\_time\_part60ae\_fourier\_joukowsky\_generation\_principle}
设
\[
\mu=\sum_{j=1}^{N}p_j\delta_{\lambda_j}
\]
为有限支撑概率分布，其中 \(p_j\ge 0\)、\(\sum_j p_j=1\)。定义其特征函数
\[
\widehat\mu(t):=\int_{\RR}e^{itx}\,\dd\mu(x)=\sum_{j=1}^{N}p_j e^{i\lambda_j t}.
\]
则轨道
\[
t\longmapsto \widehat\mu(t)\in\CC
\]
恰是若干单位圆匀速旋转的凸组合；反之，任意形如
\[
\sum_{j=1}^{N}p_j e^{i\lambda_j t}
\]
的三角多项式，只要系数非负且和为 \(1\)，都可实现为某个有限支撑概率分布的特征函数轨道。

进一步，定义归一化 Joukowsky 门
\[
J:S^1\to [-1,1],
\qquad
J(z):=\frac{z+z^{-1}}{2}=\Re z.
\]
则有：
\begin{enumerate}
  \item 对任意有理斜率 \(n/d>0\)，玫瑰线的复表示
  \[
  z(\theta):=\cos\!\Bigl(\frac{n}{d}\theta\Bigr)e^{i\theta}
  \]
  满足严格恒等式
  \[
  z(\theta)=\frac12\Bigl(e^{i(1+n/d)\theta}+e^{i(1-n/d)\theta}\Bigr),
  \]
  因而正是二点分布
  \[
  \frac12\bigl(\delta_{1+n/d}+\delta_{1-n/d}\bigr)
  \]
  的特征函数轨道；
  \item 对任意整数 \(a,b\ge 1\) 与相位 \(\delta\in\RR\)，Lissajous 曲线
  \[
  (x(t),y(t))=\bigl(\cos(at+\delta),\cos(bt)\bigr)
  \]
  可写为
  \[
  \gamma_{a,b,\delta}(t):=\bigl(e^{i(at+\delta)},e^{ibt}\bigr)\in S^1\times S^1
  \]
  经坐标门 \((J,J)\) 的像：
  \[
  (x(t),y(t))
  =
  \bigl(J(\gamma_{a,b,\delta}^{(1)}(t)),J(\gamma_{a,b,\delta}^{(2)}(t))\bigr).
  \]
\end{enumerate}
\end{theorem}
\begin{proof}
第一条只是特征函数定义的直接展开；反向上取离散随机变量 \(X\) 满足
\[
\PP(X=\lambda_j)=p_j,
\]
便得到
\[
\EE[e^{itX}]=\sum_{j=1}^{N}p_j e^{i\lambda_j t}.
\]

对玫瑰线，只需用 Euler 公式
\[
\cos\alpha=\frac{e^{i\alpha}+e^{-i\alpha}}{2}
\]
展开：
\[
\cos\!\Bigl(\frac{n}{d}\theta\Bigr)e^{i\theta}
=
\frac12\Bigl(e^{i(1+n/d)\theta}+e^{i(1-n/d)\theta}\Bigr).
\]

对 Lissajous 曲线，由
\[
J(e^{i\vartheta})=\frac{e^{i\vartheta}+e^{-i\vartheta}}{2}=\cos\vartheta
\]
可知
\[
J(\gamma_{a,b,\delta}^{(1)}(t))=\cos(at+\delta),
\qquad
J(\gamma_{a,b,\delta}^{(2)}(t))=\cos(bt).
\]
于是所示分解成立。证毕。
\end{proof}

\begin{proposition}[奇异环作为 Joukowsky 门的分歧环及其接触次数的完全计数]\label{prop:xi-time-part60ae-lissajous-singular-ring-contact-count}
\leanverified{paper\_xi\_time\_part60ae\_lissajous\_singular\_ring\_contact\_count}
沿用定理 \ref{thm:xi-time-part60ae-fourier-joukowsky-generation-principle} 的归一化 Joukowsky 门 \(J\)。则乘积映射
\[
J\times J:S^1\times S^1\to [-1,1]^2
\]
的 Jacobian 退化集合为
\[
\Sigma:=
\bigl(\{\pm 1\}\times S^1\bigr)\cup \bigl(S^1\times \{\pm 1\}\bigr),
\]
其像恰为正方形边界
\[
\partial([-1,1]^2).
\]

对闭测地线
\[
\gamma_{a,b,\delta}(t)=\bigl(e^{i(at+\delta)},e^{ibt}\bigr)
\qquad (t\in[0,2\pi))
\]
定义奇异环接触集合
\[
\mathcal T_{a,b,\delta}:=
\{t\in[0,2\pi):\gamma_{a,b,\delta}(t)\in\Sigma\}.
\]
再定义角点重合项
\[
\mathcal C_{a,b,\delta}
:=
\Bigl\{
t:\ e^{i(at+\delta)}\in\{\pm 1\},\ e^{ibt}\in\{\pm 1\}
\Bigr\}.
\]
则有精确计数公式
\[
\boxed{\;
|\mathcal T_{a,b,\delta}|=2a+2b-|\mathcal C_{a,b,\delta}|.\;}
\]
并且若记 \(g:=\gcd(a,b)\)，则
\[
\mathcal C_{a,b,\delta}\neq\varnothing
\iff
\delta\in \frac{\pi}{g}\ZZ,
\]
且在此情形下
\[
\boxed{\;
|\mathcal C_{a,b,\delta}|=2g.\;}
\]
\end{proposition}
\begin{proof}
由
\[
J(e^{i\vartheta})=\cos\vartheta
\]
可知 \(\mathrm dJ\) 退化当且仅当 \(\sin\vartheta=0\)，亦即 \(e^{i\vartheta}\in\{\pm 1\}\)。故 \(J\times J\) 的分歧环正是
\[
\Sigma=
\bigl(\{\pm 1\}\times S^1\bigr)\cup \bigl(S^1\times \{\pm 1\}\bigr),
\]
而其像显然是 \([-1,1]^2\) 的边界。

对接触计数，条件 \(\gamma_{a,b,\delta}(t)\in\Sigma\) 等价于
\[
at+\delta\in\pi\ZZ
\qquad\text{或}\qquad
bt\in\pi\ZZ.
\]
在区间 \([0,2\pi)\) 内，前者恰有 \(2a\) 个解，后者恰有 \(2b\) 个解，故由容斥得
\[
|\mathcal T_{a,b,\delta}|=2a+2b-|\mathcal C_{a,b,\delta}|.
\]

现求 \(|\mathcal C_{a,b,\delta}|\)。它等价于同余系统
\[
at\equiv -\delta\pmod{\pi},
\qquad
bt\equiv 0\pmod{\pi}.
\]
令
\[
a=ga',
\qquad
b=gb',
\qquad
\gcd(a',b')=1.
\]
第二式给出
\[
t=\frac{\pi k}{b}
\qquad (k\in\ZZ).
\]
代回第一式，得到
\[
\frac{a}{b}\pi k\equiv -\delta\pmod{\pi}
\iff
\frac{a'}{b'}k\equiv -\frac{\delta}{\pi}\pmod{1}.
\]
由于 \(a'\) 与 \(b'\) 互素，该同余有解当且仅当
\[
\frac{\delta}{\pi}\in \frac{1}{g}\ZZ,
\]
即
\[
\delta\in\frac{\pi}{g}\ZZ.
\]
一旦可解，\(k\) 在模 \(2b'=2b/g\) 的剩余类中恰有 \(2\) 个解，而每个剩余类在 \([0,2\pi)\) 上对应 \(g\) 个不同的 \(t\)。故总数为
\[
|\mathcal C_{a,b,\delta}|=2g.
\]
证毕。
\end{proof}

\begin{theorem}[Fibonacci 子集和的单位圆干涉曲线与离散碰撞谱的模采样]\label{thm:xi-time-part60ae-fibonacci-interference-curve-sampling}
沿用定理 \ref{thm:fold-collision-spectrum} 与定义 \ref{def:fold-normalized-amplitude} 的记号，记
\[
M:=F_{m+2},
\qquad
D_m(x)\equiv \prod_{j=1}^{m}\bigl(1+x^{F_{j+1}}\bigr)\pmod{x^M-1},
\]
\[
\widehat d_m(k)=D_m(e^{-2\pi i k/M}),
\qquad
a_m(k):=\frac{|\widehat d_m(k)|}{2^m}.
\]
定义连续单位圆干涉曲线
\[
\Phi_m(t):=2^{-m}D_m(e^{it})
=
2^{-m}\prod_{j=1}^{m}\bigl(1+e^{iF_{j+1}t}\bigr),
\qquad
t\in\RR.
\]
则成立：
\begin{enumerate}
  \item 离散碰撞谱正是 \(\Phi_m\) 在单位圆等分网格上的模采样：
  \[
  \boxed{\;
  a_m(k)=\bigl|\Phi_m(2\pi k/M)\bigr|
  \qquad (k\in \ZZ/(M\ZZ)).\;}
  \]
  \item 若定义零点集合
  \[
  \mathcal S_m:=\{e^{it}\in S^1:\ \Phi_m(t)=0\},
  \]
  则
  \[
  \boxed{\;
  \mathcal S_m=
  \bigcup_{j=1}^{m}
  \left\{
  e^{i(2\ell+1)\pi/F_{j+1}}:\ 0\le \ell\le F_{j+1}-1
  \right\}.\;}
  \]
  \item 并集
  \[
  \bigcup_{m\ge 1}\mathcal S_m
  \]
  在 \(S^1\) 中稠密。
\end{enumerate}
\end{theorem}
\begin{proof}
由
\[
\widehat d_m(k)=D_m(e^{-2\pi i k/M})
\]
与
\[
a_m(k)=\frac{|\widehat d_m(k)|}{2^m}
\]
可知
\[
a_m(k)=2^{-m}\bigl|D_m(e^{-2\pi i k/M})\bigr|
=
\bigl|\Phi_m(-2\pi k/M)\bigr|.
\]
又因为取模后对 \(k\mapsto -k\) 不变，故亦可写成
\[
a_m(k)=\bigl|\Phi_m(2\pi k/M)\bigr|.
\]

对零点集合，由乘积表达式，
\[
\Phi_m(t)=0
\iff
\exists\,j\in\{1,\dots,m\}\ \text{使}\ 1+e^{iF_{j+1}t}=0.
\]
这等价于
\[
F_{j+1}t\equiv \pi \pmod{2\pi},
\]
亦即
\[
t\equiv \frac{(2\ell+1)\pi}{F_{j+1}}
\pmod{2\pi}
\qquad (\ell\in\ZZ).
\]
投影到单位圆便得到 \(\mathcal S_m\) 的显式表达。

最后证稠密性。任取单位圆上的非空弧段 \(I\subset S^1\)，其角长度为某个 \(\eta>0\)。取 \(n\) 充分大，使
\[
\frac{2\pi}{F_{n+1}}<\eta.
\]
则步长为 \(2\pi/F_{n+1}\) 的奇半格点
\[
\frac{(2\ell+1)\pi}{F_{n+1}}
\qquad (0\le \ell\le F_{n+1}-1)
\]
在圆周上的相邻间距不超过 \(\eta\)，故必有一个点落入 \(I\)。该点属于 \(\mathcal S_n\)，从而
\[
I\cap \bigcup_{m\ge 1}\mathcal S_m\neq \varnothing.
\]
故该并集在 \(S^1\) 中稠密。证毕。
\end{proof}

\begin{theorem}[准 Lissajous 轨道接近角点奇异环的 Diophantine 定量律]\label{thm:xi-time-part60ae-diophantine-singular-ring-approximation}
\leanverified{paper\_xi\_time\_part60ae\_diophantine\_singular\_ring\_approximation}
设 \(\alpha>0\) 为无理数，并考虑环面直线流
\[
\gamma_{\alpha}(t):=(e^{it},e^{i\alpha t})\in S^1\times S^1.
\]
记角点奇异环为
\[
\Sigma_{\square}:=\{\pm 1\}\times \{\pm 1\}.
\]
设 \(p_n/q_n\) 为 \(\alpha\) 的连分数主收敛。定义时刻
\[
t_n:=\pi q_n.
\]
则第一坐标严格满足
\[
e^{it_n}=(-1)^{q_n}\in\{\pm 1\},
\]
而第二坐标具有近角点表示
\[
e^{i\alpha t_n}=(-1)^{p_n}e^{i\pi(\alpha q_n-p_n)},
\qquad
|\alpha q_n-p_n|<\frac{1}{q_{n+1}}.
\]
从而
\[
\boxed{\;
\bigl|\cos(\alpha t_n)-(-1)^{p_n}\bigr|
=
O\!\Bigl(\frac{1}{q_{n+1}^2}\Bigr).\;}
\]
换言之，\(\gamma_{\alpha}\) 在时刻 \(t_n\) 以二次量级贴近角点奇异环 \(\Sigma_{\square}\)，而贴近速率完全由下一收敛分母 \(q_{n+1}\) 控制。

特别地，当 \(\alpha=\varphi\) 时，
\[
q_n=F_n,
\]
故该贴近序列按 Fibonacci 尺度分层；其误差层级与定理 \ref{thm:fold-critical-resonance-constant} 中黄金比共振乘积
\[
C_{\varphi}(u)=\prod_{n\ge 2}\Bigl|\cos\bigl(\pi u\varphi^{-n}\bigr)\Bigr|
\]
所依赖的小角衰减阶具有同源的 Fibonacci 逼近结构。
\end{theorem}
\begin{proof}
连分数理论给出
\[
\left|\alpha-\frac{p_n}{q_n}\right|
<
\frac{1}{q_nq_{n+1}},
\]
故
\[
|\alpha q_n-p_n|<\frac{1}{q_{n+1}}.
\]
取
\[
t_n=\pi q_n,
\]
则
\[
e^{it_n}=e^{i\pi q_n}=(-1)^{q_n}\in\{\pm 1\}.
\]
另一方面，
\[
\alpha t_n=\pi\alpha q_n=\pi p_n+\pi(\alpha q_n-p_n),
\]
于是
\[
e^{i\alpha t_n}=(-1)^{p_n}e^{i\pi(\alpha q_n-p_n)}.
\]
从而
\[
\cos(\alpha t_n)
=
\cos\bigl(\pi p_n+\pi(\alpha q_n-p_n)\bigr)
=
(-1)^{p_n}\cos\bigl(\pi(\alpha q_n-p_n)\bigr).
\]
令
\[
\varepsilon_n:=\alpha q_n-p_n.
\]
由 Taylor 展开
\[
\cos(\pi\varepsilon_n)=1-\frac{\pi^2}{2}\varepsilon_n^2+O(\varepsilon_n^4)
\]
可得
\[
\bigl|\cos(\alpha t_n)-(-1)^{p_n}\bigr|
=
O(\varepsilon_n^2)
=
O\!\Bigl(\frac{1}{q_{n+1}^2}\Bigr).
\]
当 \(\alpha=\varphi\) 时，主收敛分母即为 Fibonacci 数 \(F_n\)，故误差分层按 Fibonacci 尺度组织。最后一条仅指出：\(C_{\varphi}(u)\) 的每一因子都由角尺度 \(\varphi^{-n}\) 控制，而 \(\varphi\) 的 Diophantine 最优逼近同样由 Fibonacci 分母主导，故二者共享同一误差层级来源。证毕。
\end{proof}

\begin{conjecture}[奇异环稠密化诱导碰撞谱异常质量]\label{conj:xi-time-part60ae-singular-ring-densification-collision-anomaly}
沿用定理 \ref{thm:xi-time-part60ae-fibonacci-interference-curve-sampling} 的记号，定义模值经验测度
\[
\nu_m:=
\frac{1}{F_{m+2}}
\sum_{k=0}^{F_{m+2}-1}
\delta_{\left|\Phi_m(2\pi k/F_{m+2})\right|}.
\]
预期存在一个非原子极限测度 \(\nu\) 使
\[
\nu_m\Rightarrow \nu
\qquad (m\to\infty),
\]
并且其近零质量满足以下统一原则：对小参数 \(\varepsilon\downarrow 0\)，
\[
\nu([0,\varepsilon])
\]
的主导衰减速率由奇异环并集
\[
\bigcup_{m\ge 1}\mathcal S_m
\]
在 \(S^1\) 上的加密速率所支配，而该主导项应可用黄金比共振常数族 \(C_{\varphi}(u)\) 的小角级联近似来刻画。

换言之，Fibonacci 干涉曲线的奇异环稠密化应当是离散碰撞谱异常质量的统一几何来源：单位圆上的层级稠密零环一旦通过等分采样被读出，便必然在 \(\{a_m(k)\}\) 的经验分布中留下可审计的近零异常质量。
\end{conjecture}

\noindent
由此，有限支撑特征函数、玫瑰线、Lissajous 曲线与 Fibonacci 子集和多项式 \(D_m\) 不再是彼此分离的对象：它们都可统一理解为单位圆旋转经 Joukowsky 门读出的有限干涉曲线。对于折叠谱而言，\(D_m\) 给出的并非单一相位花纹，而是一族随 Fibonacci 尺度递归加密的半转根环；其离散采样模值正是碰撞谱 \(a_m(k)\)，而黄金比的 Diophantine 最优逼近则进一步解释了这些干涉曲线在奇异环附近所呈现的层级接近律。

\endinput

\paragraph{Gödel--Zeckendorf 地址极限的黄金自仿射闭合、自然尺度维数与查询刚性}
在地址前缀等距、Zeckendorf 有限截面递推与归一化复杂度界之间，还可抽出一组更紧的终端闭合：其一把“禁止相邻 \(1\)”直接压缩为一维 \(2^L\)-进地址线上的严格二分自仿射方程；其二把该极限集的自然尺度覆盖数、Hausdorff 维数与 Minkowski 内容写成 Fibonacci 常数的显式闭式；其三把折叠与稳定运算的可计算性提升为可精确审计的位查询刚性，并进一步把稳定算术层锁定为与自然数半环 PTIME 双向可解释的同一表示；最终它再与真实输入 \(40\)-状态核的一阶 Edgeworth 指纹汇合为同一条可复核链。

\begin{theorem}[Gödel--Zeckendorf 地址极限的严格二分自仿射方程]\label{thm:xi-zg-address-binary-self-affine}\leanverified{paper\_xi\_zg\_address\_binary\_self\_affine}
取整数 \(L\ge 1\)，令
\[
B:=2^L,
\qquad
v\in T_2(E_{\min})\setminus 2T_2(E_{\min}),
\qquad
L_v:=\ZZ_2v.
\]
定义 admissible 序列空间
\[
\Sigma_{\mathrm{ZG}}
:=
\bigl\{
(z_t)_{t\ge 1}\in\{0,1\}^{\NN}:\ z_tz_{t+1}=0\ \ (t\ge 1)
\bigr\},
\]
以及地址极限集
\[
K_{B,v}
:=
\left\{
\sum_{t\ge 1}z_tB^{t-1}v:\ (z_t)_{t\ge 1}\in\Sigma_{\mathrm{ZG}}
\right\}
\subseteq L_v.
\]
则上述级数在 \(L_v\) 中收敛，\(\Sigma_{\mathrm{ZG}}\to K_{B,v}\) 的地址映射为单射，并且成立严格不交并分解
\[
\boxed{\;
K_{B,v}
=
BK_{B,v}\ \sqcup\ \bigl(v+B^2K_{B,v}\bigr).\;}
\]
\end{theorem}
\begin{proof}
由于 \(B=2^L\) 在 \(2\)-进绝对值下满足 \(|B|_2<1\)，故对任意 \((z_t)\in\Sigma_{\mathrm{ZG}}\)，级数
\[
\sum_{t\ge 1}z_tB^{t-1}v
\]
在完备群 \(L_v\) 中收敛。

再把 \(\mathcal E=\{0,1\}\)、\(\mathrm{code}(0)=0\)、\(\mathrm{code}(1)=1\) 代入定理
\ref{thm:xi-time-part9g-holographic-prefix-isometry-on-line}，即可知不同 admissible 无限序列给出不同地址，因而地址映射在 \(\Sigma_{\mathrm{ZG}}\) 上为单射。

现证自仿射分解。任取
\[
x=\sum_{t\ge 1}z_tB^{t-1}v\in K_{B,v},
\qquad
(z_t)\in\Sigma_{\mathrm{ZG}}.
\]
若 \(z_1=0\)，则尾序列 \((z_{t+1})_{t\ge 1}\) 仍属 \(\Sigma_{\mathrm{ZG}}\)，并且
\[
x
=
\sum_{t\ge 2}z_tB^{t-1}v
=
B\sum_{u\ge 1}z_{u+1}B^{u-1}v
\in BK_{B,v}.
\]
若 \(z_1=1\)，则 admissibility 强制 \(z_2=0\)，故 \((z_{t+2})_{t\ge 1}\in\Sigma_{\mathrm{ZG}}\)，并且
\[
x
=
v+\sum_{t\ge 3}z_tB^{t-1}v
=
v+B^2\sum_{u\ge 1}z_{u+2}B^{u-1}v
\in v+B^2K_{B,v}.
\]
因此
\[
K_{B,v}\subseteq BK_{B,v}\cup\bigl(v+B^2K_{B,v}\bigr).
\]

反向包含直接成立：若 \(y=\sum_{t\ge 1}z_tB^{t-1}v\in K_{B,v}\)，则
\[
By=\sum_{t\ge 2}z_{t-1}B^{t-1}v\in K_{B,v},
\]
而
\[
v+B^2y
=
v+\sum_{t\ge 3}z_{t-2}B^{t-1}v
\in K_{B,v},
\]
因为两条扩展序列仍满足禁止相邻 \(1\) 约束。

最后证不交。若
\[
Bx=v+B^2y
\qquad
(x,y\in K_{B,v}),
\]
则左端对应一条首位为 \(0\) 的 admissible 序列，右端对应一条首两位为 \(10\) 的 admissible 序列。由定理
\ref{thm:xi-time-part9g-holographic-prefix-isometry-on-line} 的地址单射性，二者不可能相等，故并为不交并。证毕。
\end{proof}

\begin{theorem}[Gödel--Zeckendorf 地址极限的精确 Minkowski 维数与自然尺度内容]\label{thm:xi-zg-address-minkowski-dimension-content}\leanverified{paper\_xi\_zg\_address\_minkowski\_dimension\_content}
沿用定理 \ref{thm:xi-zg-address-binary-self-affine} 的记号，并在 \(L_v\) 上取定理 \ref{thm:xi-time-part9g-holographic-prefix-isometry-on-line} 的 \(B\)-进超度量。对 \(n\ge 0\)，令 \(N_{B,v}(n)\) 表示覆盖 \(K_{B,v}\) 所需的最少 \(B^{-n}\)-球个数；等价地，\(N_{B,v}(n)\) 为商 \(L_v/B^nL_v\) 中与 \(K_{B,v}\) 相交的余类数。则
\[
\boxed{\;
N_{B,v}(n)=F_{n+2}\qquad (n\ge 0).\;}
\]
从而
\[
\boxed{\;
\dim_{\mathrm H}(K_{B,v})
=
\dim_{\mathrm M}(K_{B,v})
=
\frac{\log\varphi}{\log B}
=
\frac{\log\varphi}{L\log 2}.\;}
\]
并且沿自然 \(B\)-进尺度存在精确内容极限
\[
\boxed{\;
\lim_{n\to\infty}
N_{B,v}(n)\,B^{-n\dim_{\mathrm M}(K_{B,v})}
=
\frac{\varphi^2}{\sqrt5}.\;}
\]
此外，\(K_{B,v}\) 在 \(L_v\) 的 Haar 概率测度下为零测集。
\end{theorem}
\begin{proof}
对每个 \(n\ge 0\)，记
\[
\mathcal A_n
:=
\bigl\{
(z_1,\dots,z_n)\in\{0,1\}^n:\ z_tz_{t+1}=0\ \ (1\le t<n)
\bigr\},
\]
并约定 \(\mathcal A_0=\{\varnothing\}\)。对 \(a=(z_1,\dots,z_n)\in\mathcal A_n\)，定义对应 cylinder
\[
[a]
:=
\bigl\{
(w_t)_{t\ge 1}\in\Sigma_{\mathrm{ZG}}:\ w_1=z_1,\dots,w_n=z_n
\bigr\},
\]
以及前缀地址
\[
X_n(a):=\sum_{t=1}^{n}z_tB^{t-1}v.
\]
由定理 \ref{thm:xi-time-part9g-holographic-prefix-isometry-on-line}，
\[
X([a])=\bigl(X_n(a)+B^nL_v\bigr)\cap K_{B,v},
\]
且不同 \(a\in\mathcal A_n\) 给出不同的 \(B^nL_v\)-余类。于是 \(K_{B,v}\) 在尺度 \(B^{-n}\) 上恰分解为 \(|\mathcal A_n|\) 个互不相交 cylinder，故
\[
N_{B,v}(n)=|\mathcal A_n|.
\]

另一方面，定理 \ref{thm:xi-time-part9g-zg-primorial-fibonacci-affine-recursion} 已证明
\[
|\mathcal A_n|=F_{n+2},
\]
故第一条公式成立。

记
\[
D:=\frac{\log\varphi}{\log B}.
\]
则由 Binet 公式
\[
F_{n+2}\sim \frac{\varphi^{n+2}}{\sqrt5}
\]
得到
\[
\frac{\log N_{B,v}(n)}{n\log B}
=
\frac{\log F_{n+2}}{n\log B}
\longrightarrow
\frac{\log\varphi}{\log B}
=
D.
\]
因此 Minkowski 维数为 \(D\)。

再证 Hausdorff 维数。若 \(s>D\)，则第 \(n\) 层 cylinder 覆盖给出
\[
\mathcal H^s_{B^{-n}}(K_{B,v})
\le
N_{B,v}(n)\,B^{-ns}
=
F_{n+2}B^{-ns}
\to 0,
\]
故 \(\mathcal H^s(K_{B,v})=0\)。

反之，若 \(s<D\)，则任一直径不超过 \(B^{-n}\) 的集合都包含于某个 \(B^{-n}\)-球中，而任一此类球至多与一个第 \(n\) 层 cylinder 相交；否则它将同时命中两个不同的 \(B^nL_v\)-余类，与定理
\ref{thm:xi-time-part9g-holographic-prefix-isometry-on-line} 给出的严格余类分离矛盾。故任意直径不超过 \(B^{-n}\) 的覆盖都至少需要 \(N_{B,v}(n)\) 个集合，从而
\[
\mathcal H^s_{B^{-n}}(K_{B,v})
\ge
N_{B,v}(n)\,B^{-ns}
=
F_{n+2}B^{-ns}
\to \infty.
\]
因此 \(\mathcal H^s(K_{B,v})=\infty\)。综上 \(\dim_{\mathrm H}(K_{B,v})=D=\dim_{\mathrm M}(K_{B,v})\)。

再由 \(B^D=\varphi\)，有
\[
N_{B,v}(n)\,B^{-nD}
=
F_{n+2}\varphi^{-n}
\longrightarrow
\frac{\varphi^2}{\sqrt5},
\]
即得自然尺度内容极限。

最后，记 \(\lambda_v\) 为 \(L_v\) 的 Haar 概率测度。每个 \(B^nL_v\)-余类的测度均为 \(B^{-n}\)，而 \(K_{B,v}\) 只与其中 \(N_{B,v}(n)=F_{n+2}\) 个余类相交，故
\[
\lambda_v(K_{B,v})
\le
F_{n+2}B^{-n}
\sim
\frac{\varphi^2}{\sqrt5}\Bigl(\frac{\varphi}{B}\Bigr)^n.
\]
由于 \(B=2^L\ge 2>\varphi\)，令 \(n\to\infty\) 即得 \(\lambda_v(K_{B,v})=0\)。证毕。
\end{proof}

\begin{theorem}[折叠算子的确定性位查询复杂度恰为分辨率]\label{thm:xi-fold-query-complexity-exact-m}\leanverified{paper\_xi\_fold\_query\_complexity\_exact\_m}
对 \(m\ge 1\)，定义 \(Q_{\det}(\Fold_m)\) 为：在最坏情形下，为精确输出 \(\Fold_m(\omega)\)，确定性算法必须读取的输入位数的最小值。则
\[
\boxed{\;
Q_{\det}(\Fold_m)=m.\;}
\]
\end{theorem}
\begin{proof}
上界由命题 \ref{prop:fold-complexity} 直接给出：顺序读取全部 \(m\) 位并执行 Zeckendorf 规范化即可，故
\[
Q_{\det}(\Fold_m)\le m.
\]

现证下界。任取一个精确计算 \(\Fold_m\) 的确定性算法 \(\mathcal A\)，并考察其在输入
\[
\mathbf0:=0^m
\]
上的运行。若 \(\mathcal A\) 在该输入上未查询某一位置 \(j\in\{1,\dots,m\}\)，则令
\[
e_j:=00\cdots010\cdots0
\]
为仅第 \(j\) 位为 \(1\) 的字。由于 \(\mathbf0\) 与 \(e_j\) 在 \(\mathcal A\) 已查询的全部坐标上取值相同，算法在两输入上将经历完全相同的查询转录并输出同一结果。

然而 \(e_j\in X_m\)，由命题 \ref{prop:fold-basic} 的幂等性，
\[
\Fold_m(\mathbf0)=\mathbf0,
\qquad
\Fold_m(e_j)=e_j,
\]
二者显然不同，矛盾。故 \(\mathcal A\) 在输入 \(\mathbf0\) 上必须读取全部 \(m\) 位，从而
\[
Q_{\det}(\Fold_m)\ge m.
\]
合并上下界即得结论。证毕。
\end{proof}

\begin{corollary}[稳定加法的全输入查询饱和与稳定乘法单位切片的满查询刚性]\label{cor:xi-stable-add-query-saturation-mul-unit-slice}
沿用定义 \ref{def:stable-add} 与 \ref{def:stable-mul} 的记号。对支撑包含于 \(\{1,\dots,m\}\) 的合法输入，记 \(Q_{\det}^{\oplus}(m)\) 为稳定加法 \(c\oplus d\) 的最坏情形总位查询复杂度。再令
\[
\mathbf1:=10\cdots0\in X_\infty,
\qquad
\mathrm{Val}(\mathbf1)=1,
\]
并定义乘法的单位切片
\[
L_m^{\otimes}(c):=c\otimes \mathbf1,
\qquad
R_m^{\otimes}(d):=\mathbf1\otimes d.
\]
则有
\[
\boxed{\;
Q_{\det}^{\oplus}(m)=2m,\qquad
Q_{\det}(L_m^{\otimes})=Q_{\det}(R_m^{\otimes})=m.\;}
\]
\end{corollary}
\begin{proof}
先看稳定加法。上界由命题 \ref{prop:stable-arith-complexity} 直接给出：读取两输入全部 \(2m\) 位后执行值层加法与 Zeckendorf 规范化即可，故
\[
Q_{\det}^{\oplus}(m)\le 2m.
\]

现证下界。任取一个精确计算 \(c\oplus d\) 的确定性算法 \(\mathcal A\)，并考察其在输入
\[
(\mathbf0,\mathbf0)
\]
上的运行，其中 \(\mathbf0:=0^\infty\)，并令 \(e_j\) 表示第 \(j\) 位为 \(1\)、其余位为 \(0\) 的合法无限字。若 \(\mathcal A\) 在该运行中未查询第一输入的某一位置 \(j\le m\)，则与输入
\[
(e_j,\mathbf0)
\]
相比，所有已查询坐标取值仍完全相同，故 \(\mathcal A\) 将给出同一输出；但
\[
e_j\oplus \mathbf0=e_j,
\qquad
\mathbf0\oplus \mathbf0=\mathbf0,
\]
矛盾。故第一输入的前 \(m\) 位必须全部被查询。

若 \(\mathcal A\) 未查询第二输入的某一位置 \(j\le m\)，同理比较
\[
(\mathbf0,\mathbf0)
\qquad\text{与}\qquad
(\mathbf0,e_j),
\]
即可得到矛盾，因为
\[
\mathbf0\oplus e_j=e_j
\neq
\mathbf0.
\]
故在输入 \((\mathbf0,\mathbf0)\) 上，两侧前 \(m\) 位都必须全部读取，从而
\[
Q_{\det}^{\oplus}(m)\ge 2m.
\]
于是 \(Q_{\det}^{\oplus}(m)=2m\)。

再看乘法单位切片。由 \(\mathrm{Val}(\mathbf1)=1\) 与定义 \ref{def:stable-mul}，
\[
L_m^{\otimes}(c)=c\otimes\mathbf1=c,
\qquad
R_m^{\otimes}(d)=\mathbf1\otimes d=d.
\]
因此两条切片在约束输入类上就是恒等映射。其查询下界可用与定理 \ref{thm:xi-fold-query-complexity-exact-m} 同样的零字--单点字对抗得到：若在输入 \(\mathbf0\) 上遗漏某个位置 \(j\le m\)，则无法区分 \(\mathbf0\) 与 \(e_j\)。故
\[
Q_{\det}(L_m^{\otimes})\ge m,
\qquad
Q_{\det}(R_m^{\otimes})\ge m.
\]
显然读取前 \(m\) 位即可输出该合法字本身，故反向不等式也成立，从而
\[
Q_{\det}(L_m^{\otimes})=Q_{\det}(R_m^{\otimes})=m.
\]
证毕。
\end{proof}

\begin{theorem}[稳定算术层与自然数半环在 PTIME 内双向可解释]\label{thm:xi-stable-semiring-ptime-biinterpretability}
\leanverified{paper\_xi\_stable\_semiring\_ptime\_biinterpretability}
沿用定义 \ref{def:stable-add}、\ref{def:stable-mul} 与命题
\ref{prop:stable-arith-complexity} 的记号。则值映射
\[
\mathrm{Val}:X_\infty\longrightarrow \NN
\]
及其逆映射 \(Z:=\mathrm{Val}^{-1}\) 给出交换半环同构
\[
\boxed{\;
\mathrm{Val}(c\oplus d)=\mathrm{Val}(c)+\mathrm{Val}(d),\qquad
\mathrm{Val}(c\otimes d)=\mathrm{Val}(c)\,\mathrm{Val}(d).\;}
\]
更进一步，在有效分辨率 \(m\) 的输入口径下，两侧运算可在多项式时间内双向模拟：
\begin{enumerate}
  \item 从 \(c,d\in X_\infty\) 计算 \(c\oplus d\) 与 \(c\otimes d\)，在比特口径下分别需要
  \[
  O(m^2),\qquad O(M(m)+m^2)
  \]
  位运算，其中 \(M(m)\) 为 \(m\)-比特整数乘法成本；
  \item 从 \(m\)-比特整数 \(N\) 计算 Zeckendorf 规范形 \(Z(N)\)，以及从 \(c\in X_\infty\) 计算 \(\mathrm{Val}(c)\)，都可在 \(O(m^2)\) 位运算内完成。
\end{enumerate}
因此
\[
\boxed{\;
(X_\infty,\oplus,\otimes)\ \text{与}\ (\NN,+,\times)
\ \text{在 semiring 意义下同构，并在 PTIME 内双向可解释。}\;}
\]
\end{theorem}
\begin{proof}
交换半环同构本身是定义性的：由定理 \ref{thm:stable-add-value-consistency} 与定义 \ref{def:stable-mul}，
\[
\mathrm{Val}(c\oplus d)=\mathrm{Val}(c)+\mathrm{Val}(d),
\qquad
c\otimes d=\mathrm{Val}^{-1}\!\bigl(\mathrm{Val}(c)\mathrm{Val}(d)\bigr),
\]
而 \(\mathrm{Val}\) 由 Zeckendorf 唯一性给出双射，因此半环同构正是定理 \ref{thm:mul-definitional} 的内容。

算法复杂度则给出真正的计算内容。对第一条，命题 \ref{prop:stable-arith-complexity} 已经证明：在有效分辨率 \(m\) 的输入上，稳定加法与稳定乘法可按“值层运算 \(+\) Zeckendorf 规范化”实现，并分别满足
\[
O(m^2),\qquad O(M(m)+m^2)
\]
的比特复杂度。

对第二条，若给定 \(m\)-比特整数 \(N\)，则其 Zeckendorf 规范形 \(Z(N)\) 可由附录复杂度界中的贪心规范化实现；命题 \ref{prop:fold-complexity} 对 Fibonacci 权重的分析表明，这一步需要 \(O(m^2)\) 位运算。反过来，若给定支撑包含于 \(\{1,\dots,m\}\) 的 \(c\in X_\infty\)，则
\[
\mathrm{Val}(c)=\sum_{k=1}^{m}c_kF_{k+1}
\]
只是 \(O(m)\) 项、每项位长 \(O(m)\) 的加权求和，因此同样可在 \(O(m^2)\) 位运算内完成。

于是，自然数加乘可先经 \(Z\) 送入稳定地址层，再用 \(\oplus,\otimes\) 模拟并用 \(\mathrm{Val}\) 读回；稳定地址层运算也可先经 \(\mathrm{Val}\) 送回整数值层完成。两侧都具有多项式时间实现，故得到所述 PTIME 双向可解释性。证毕。
\end{proof}

\paragraph{真实输入 \texorpdfstring{$40$}{40}-状态核的一阶 Edgeworth 指纹闭合}
与上述地址几何与查询刚性并行，定理 \ref{thm:xi-real-input-40-output-edgeworth-nongaussian} 已给出真实输入 \(40\)-状态核在输出位势方向上的显式符号指纹
\[
\gamma_1\approx -1.476481644016533,
\qquad
\gamma_2\approx 2.3928814165461985.
\]
因此输出计数的一阶 Edgeworth 修正严格左偏且严格厚尾。于是，Gödel--Zeckendorf 地址极限的黄金维数律、折叠与稳定接口的查询刚性，以及真实输入核的非高斯统计指纹，便在同一审计口径下封闭为一条统一的可复核链。

\endinput

\paragraph{末位充分化、均匀基线判别与零点几何对碰撞缺口的指数失明}
由末位两层塌缩、局部信息密度的一阶离散修正、均匀基线下的两原子似然比极限以及谱零点的半维稀疏律，可进一步抽出下列八条彼此衔接的终端后果：它们分别给出末位统计量的一样本渐近充分性、亚临界样本规模下完整输出实验与末位 Bernoulli 实验的 Le Cam 等价、隐藏对数多重度涨落的二点极限律、其 Laplace--累积闭式、相对稳定层均匀基线的一样本最优判别常数、全部经典散度的 Bernoulli 塌缩、零点数量与碰撞缺口主尺度之间的严格分离，以及零点相对密度在近零邻域无法控制归一化碰撞尺度。

\begin{theorem}[末位统计量的渐近充分性]\label{thm:xi-fold-lastbit-asymptotic-sufficiency}
\leanverified{paper\_xi\_fold\_lastbit\_asymptotic\_sufficiency}
设
\[
X_m:=\{w\in\{0,1\}^m:\ w_iw_{i+1}=0\ \ (1\le i<m)\},
\qquad
\mu_m(w):=\frac{d_m^{\mathrm{bin}}(w)}{2^m},
\qquad
\ell_m(w):=w_m\in\{0,1\},
\]
并将稳定层分成两块
\[
A_i:=\{w\in X_m:\ \ell_m(w)=i\}
\qquad(i=0,1).
\]
定义 Markov 核
\[
K_m(i,\cdot):=\mathrm{Unif}(A_i),
\]
以及由末位统计量重构的块均匀近似
\[
\bar\mu_m:=K_m\circ (\ell_m)_\#\mu_m.
\]
则有
\[
\boxed{\;
D_{\mathrm{TV}}(\mu_m,\bar\mu_m)
=
O\!\Bigl(\Bigl(\frac{\varphi}{2}\Bigr)^m\Bigr).\;}
\]
因此 \(\ell_m\) 在 Le Cam 意义下对 \(\mu_m\) 是渐近充分统计量。

更具体地，对任意有界函数 \(f_m:X_m\to\RR\)，若满足
\[
\sum_{w\in A_i}f_m(w)=0
\qquad(i=0,1),
\]
则
\[
\boxed{\;
\bigl|\mathbb E_{\mu_m}[f_m]\bigr|
\le
2\|f_m\|_\infty\,
O\!\Bigl(\Bigl(\frac{\varphi}{2}\Bigr)^m\Bigr).\;}
\]
\leanverified{paper\_xi\_fold\_lastbit\_asymptotic\_sufficiency}
\end{theorem}
\begin{proof}
由命题 \ref{prop:fold-bin-lastbit-sufficient-tv}，存在绝对常数 \(C>0\) 使
\[
D_{\mathrm{TV}}(\mu_m,\bar\mu_m)\le C\Bigl(\frac{\varphi}{2}\Bigr)^m.
\]
这正是第一式。又因为 \(\bar\mu_m=K_m\circ(\ell_m)_\#\mu_m\)，已构造出一条把末位统计量提升回稳定层的重构核，其误差在全变差意义下指数消失，故 \(\ell_m\) 对 \(\mu_m\) 在 Le Cam 意义下渐近充分。

对第二条，只需注意 \(\bar\mu_m\) 在每个 \(A_i\) 上为均匀分布，故
\[
\mathbb E_{\bar\mu_m}[f_m]
=
\sum_{i=0}^{1}\bar\mu_m(A_i)\,\frac{1}{|A_i|}\sum_{w\in A_i}f_m(w)=0.
\]
于是
\[
\bigl|\mathbb E_{\mu_m}[f_m]\bigr|
=
\bigl|\mathbb E_{\mu_m}[f_m]-\mathbb E_{\bar\mu_m}[f_m]\bigr|
\le
2\|f_m\|_\infty D_{\mathrm{TV}}(\mu_m,\bar\mu_m).
\]
代入第一式即得所示估计。证毕。
\end{proof}

\begin{theorem}[亚临界样本规模下完整 bin-fold 输出实验与末位 Bernoulli 实验的 Le Cam 等价]\label{thm:xi-fold-lastbit-subcritical-sample-lecam-equivalence}
\leanverified{paper\_xi\_fold\_lastbit\_subcritical\_sample\_lecam\_equivalence}
沿用定理 \ref{thm:xi-fold-lastbit-asymptotic-sufficiency} 的记号，并记
\[
p_m:=\mu_m(A_1)=\bar\mu_m(A_1).
\]
对任意 \(n\ge 1\)，定义单点统计实验
\[
\mathcal E_m^{(n)}:=\{\mu_m^{\otimes n}\}
\qquad\text{定义在 }X_m^n\text{ 上},
\]
\[
\mathcal B_m^{(n)}:=\{\mathrm{Bernoulli}(p_m)^{\otimes n}\}
\qquad\text{定义在 }\{0,1\}^n\text{ 上}.
\]
并记
\[
\Delta(\mathcal E,\mathcal F):=\max\{\delta(\mathcal E,\mathcal F),\delta(\mathcal F,\mathcal E)\}
\]
为 Le Cam 距离。再定义确定性读出核
\[
L_m^{(n)}(w_1,\dots,w_n):=\bigl(\ell_m(w_1),\dots,\ell_m(w_n)\bigr),
\]
以及重构核
\[
K_m^{(n)}(b_1,\dots,b_n,\cdot):=
K_m(b_1,\cdot)\otimes\cdots\otimes K_m(b_n,\cdot).
\]
则有精确恒等式
\[
L_{m\#}^{(n)}\mu_m^{\otimes n}
=
\mathrm{Bernoulli}(p_m)^{\otimes n},
\qquad
K_{m\#}^{(n)}\mathrm{Bernoulli}(p_m)^{\otimes n}
=
\bar\mu_m^{\otimes n},
\]
并且存在绝对常数 \(C>0\) 使
\[
\boxed{\;
D_{\mathrm{TV}}\!\Bigl(
\mu_m^{\otimes n},
K_{m\#}^{(n)}\mathrm{Bernoulli}(p_m)^{\otimes n}
\Bigr)
\le
nC\Bigl(\frac{\varphi}{2}\Bigr)^m.\;}
\]
因此 Le Cam 距离满足
\[
\boxed{\;
\Delta\!\bigl(\mathcal E_m^{(n)},\mathcal B_m^{(n)}\bigr)
\le
nC\Bigl(\frac{\varphi}{2}\Bigr)^m.\;}
\]
特别地，若 \(n_m=o((2/\varphi)^m)\)，则
\[
\boxed{\;
\Delta\!\bigl(\mathcal E_m^{(n_m)},\mathcal B_m^{(n_m)}\bigr)\longrightarrow 0.\;}
\]
\end{theorem}
\begin{proof}
由定义，\(L_m^{(n)}\) 只是逐坐标读取末位，因此
\[
L_{m\#}^{(n)}\mu_m^{\otimes n}
=
\bigl((\ell_m)_\#\mu_m\bigr)^{\otimes n}
=
\mathrm{Bernoulli}(p_m)^{\otimes n}.
\]
另一方面，\(K_m(i,\cdot)=\mathrm{Unif}(A_i)\) 且 \(\bar\mu_m=K_m\circ(\ell_m)_\#\mu_m\)，故
\[
K_{m\#}^{(n)}\mathrm{Bernoulli}(p_m)^{\otimes n}
=
\bar\mu_m^{\otimes n}.
\]

再由乘积测度的望远镜估计，对任意概率分布 \(\mu,\nu\) 都有
\[
D_{\mathrm{TV}}(\mu^{\otimes n},\nu^{\otimes n})
\le
n\,D_{\mathrm{TV}}(\mu,\nu).
\]
将 \((\mu,\nu)=(\mu_m,\bar\mu_m)\) 代入，并使用定理 \ref{thm:xi-fold-lastbit-asymptotic-sufficiency} 中
\[
D_{\mathrm{TV}}(\mu_m,\bar\mu_m)\le C\Bigl(\frac{\varphi}{2}\Bigr)^m,
\]
即得所示全变差界。

按 deficiency 的定义，从 \(\mathcal E_m^{(n)}\) 到 \(\mathcal B_m^{(n)}\) 的缺陷为 \(0\)，因为 \(L_m^{(n)}\) 已实现精确推送；从 \(\mathcal B_m^{(n)}\) 到 \(\mathcal E_m^{(n)}\) 的缺陷至多为上述全变差误差。再由
定理中的 Le Cam 距离定义，便得到所示距离界。若 \(n_m=o((2/\varphi)^m)\)，则右端趋于 \(0\)，故两实验渐近等价。证毕。
\end{proof}

\begin{corollary}[隐藏对数多重度涨落收敛为末位支配的二点律]\label{cor:xi-fold-hidden-logmultiplicity-branchbit-two-point-law}
\leanverified{paper\_xi\_fold\_hidden\_logmultiplicity\_branchbit\_two\_point\_law}
沿用 \(\mu_m\) 的记号，并令 \(W_m\sim\mu_m\)。定义
\[
Y_m:=\log d_m^{\mathrm{bin}}(W_m)-m\log\!\Bigl(\frac{2}{\varphi}\Bigr).
\]
则有分布收敛
\[
\boxed{\;
Y_m\Longrightarrow
\frac{\varphi}{\sqrt5}\,\delta_0
+\frac{1}{\varphi\sqrt5}\,\delta_{-\log\varphi}.\;}
\]
等价地，若 \(Y_\infty\) 为满足
\[
\mathbb P(Y_\infty=0)=\frac{\varphi}{\sqrt5},
\qquad
\mathbb P(Y_\infty=-\log\varphi)=\frac{1}{\varphi\sqrt5}
\]
的二点随机变量，则 \(Y_m\Rightarrow Y_\infty\)。特别地，
\[
\boxed{\;
\lim_{m\to\infty}\operatorname{Var}(Y_m)
=
\frac{1}{\varphi\sqrt5}
\left(1-\frac{1}{\varphi\sqrt5}\right)(\log\varphi)^2.\;}
\]
因此，隐藏对数多重度的首个非平凡涨落并不趋于非退化 Gaussian 律，而直接塌缩为由末位比特控制的 Bernoulli 二点极限。
\end{corollary}
\begin{proof}
由推论 \ref{cor:xi-fold-local-information-density-first-order-discrete-law}，若记
\[
\gamma_m:=\frac{1}{m}\log d_m^{\mathrm{bin}}(W_m),
\]
则有
\[
m\left(\gamma_m-\log\!\frac{2}{\varphi}\right)\Longrightarrow
\frac{\varphi}{\sqrt5}\,\delta_0
+\frac{1}{\varphi\sqrt5}\,\delta_{-\log\varphi}.
\]
而左端恰为 \(Y_m\)。方差极限亦即该推论中的 \(m^2\operatorname{Var}(\gamma_m)\) 极限，因为
\[
\operatorname{Var}(Y_m)=m^2\operatorname{Var}(\gamma_m).
\]
证毕。
\end{proof}

\begin{theorem}[隐藏对数多重度涨落的二原子 Laplace--累积母函数闭式]\label{thm:xi-fold-hidden-logmultiplicity-laplace-cumulant-closed}
\leanverified{paper\_xi\_fold\_hidden\_logmultiplicity\_laplace\_cumulant\_closed}
沿用推论 \ref{cor:xi-fold-hidden-logmultiplicity-branchbit-two-point-law} 的记号。对 \(t\in\RR\) 定义
\[
\Lambda_m(t):=\log \EE_{\mu_m}\!\bigl[e^{tY_m}\bigr].
\]
再记
\[
\theta_\star:=\frac{1}{1+\varphi^2}=\frac{1}{\varphi\sqrt5}.
\]
则对任意紧区间 \(I\Subset\RR\)，都有一致渐近公式
\[
\boxed{\;
\sup_{t\in I}
\left|
\Lambda_m(t)
-\log\!\Bigl(\frac{\varphi+\varphi^{-t-1}}{\sqrt5}\Bigr)
\right|
=
O_I\!\Bigl(\Bigl(\frac{\varphi}{2}\Bigr)^m\Bigr).\;}
\]
特别地，
\[
\boxed{\;
\Lambda(t):=\lim_{m\to\infty}\Lambda_m(t)
=
\log\!\Bigl(\frac{\varphi+\varphi^{-t-1}}{\sqrt5}\Bigr).\;}
\]
因此 \(Y_m\) 的一切累积量都收敛到对应二原子极限律 \(Y_\infty\) 的累积量；若记
\[
\kappa_r^\infty:=\lim_{m\to\infty}\kappa_r(Y_m)=\Lambda^{(r)}(0),
\]
则前四阶为
\[
\boxed{\;
\kappa_1^\infty=-\frac{\log\varphi}{\varphi\sqrt5},\qquad
\kappa_2^\infty=\frac{(\log\varphi)^2}{5},\qquad
\kappa_3^\infty=-\frac{(\log\varphi)^3}{5\sqrt5},\qquad
\kappa_4^\infty=-\frac{(\log\varphi)^4}{25}.\;}
\]
换言之，隐藏对数多重度涨落并不落入任何非退化 Gaussian 中心极限定域，而是严格保留一个由末位比特控制的二原子极限核心。
\end{theorem}
\begin{proof}
记
\[
a:=\log\varphi,
\qquad
B_m:=W_m(m)\in\{0,1\}.
\]
由定理 \ref{thm:fold-bin-two-state-asymptotic} 的一致对数展开，
\[
Y_m=-a\,B_m+R_m,
\qquad
\sup_{w\in X_m}|R_m(w)|
=
O\!\Bigl(\Bigl(\frac{\varphi}{2}\Bigr)^m\Bigr).
\]
因此对任意紧区间 \(I\Subset\RR\) 与任意 \(t\in I\)，都有一致展开
\[
e^{tY_m}
=
e^{-atB_m}e^{tR_m}
=
e^{-atB_m}
\left(
1+O_I\!\Bigl(\Bigl(\frac{\varphi}{2}\Bigr)^m\Bigr)
\right).
\]
对 \(\mu_m\) 取期望，并利用命题 \ref{prop:fold-bin-escort-lastbit} 在 \(q=1\) 处的末位质量公式
\[
\mu_m(B_m=1)=\theta_\star+O\!\Bigl(\Bigl(\frac{\varphi}{2}\Bigr)^m\Bigr),
\qquad
\mu_m(B_m=0)=1-\theta_\star+O\!\Bigl(\Bigl(\frac{\varphi}{2}\Bigr)^m\Bigr),
\]
得到
\[
\EE_{\mu_m}[e^{tY_m}]
=
(1-\theta_\star)+\theta_\star e^{-at}
+O_I\!\Bigl(\Bigl(\frac{\varphi}{2}\Bigr)^m\Bigr).
\]
而
\[
(1-\theta_\star)+\theta_\star e^{-at}
=
\frac{\varphi}{\sqrt5}+\frac{1}{\varphi\sqrt5}\varphi^{-t}
=
\frac{\varphi+\varphi^{-t-1}}{\sqrt5}.
\]
由于右端在每个紧区间 \(I\) 上严格正且有正下界，取对数即得
\[
\Lambda_m(t)
=
\log\!\Bigl(\frac{\varphi+\varphi^{-t-1}}{\sqrt5}\Bigr)
+O_I\!\Bigl(\Bigl(\frac{\varphi}{2}\Bigr)^m\Bigr),
\]
从而得到 \(\Lambda\) 的闭式。

另一方面，由上式可知 \(Y_m\) 一致有界，故分布收敛
\[
Y_m\Rightarrow Y_\infty
\]
自动推出各阶矩收敛；累积量作为矩的通用多项式函数，也随之逐阶收敛到 \(Y_\infty\) 的累积量。由于
\[
Y_\infty=-a\,B_\star,
\qquad
\PP(B_\star=1)=\theta_\star,
\qquad
\PP(B_\star=0)=1-\theta_\star,
\]
Bernoulli 两点律的标准累积量公式给出
\[
\kappa_1^\infty=-a\theta_\star,\qquad
\kappa_2^\infty=a^2\theta_\star(1-\theta_\star),
\]
\[
\kappa_3^\infty=-a^3\theta_\star(1-\theta_\star)(1-2\theta_\star),
\]
\[
\kappa_4^\infty=a^4\theta_\star(1-\theta_\star)\bigl(1-6\theta_\star(1-\theta_\star)\bigr).
\]
再用恒等式
\[
\theta_\star(1-\theta_\star)=\frac15,
\qquad
1-2\theta_\star=\frac{1}{\sqrt5},
\]
即可化简得到所列四个闭式。证毕。
\end{proof}

\begin{theorem}[纤维信息与规范体积密度的全显式两尺度渐近]\label{thm:xi-fold-hidden-logmultiplicity-kappa-gauge-two-scale-explicit}
沿用定理 \ref{thm:fold-bin-two-state-asymptotic} 与命题 \ref{prop:fold-bin-escort-lastbit} 在 \(q=1\) 处的记号，并记
\[
\kappa_m
:=
\EE_{\mu_m}\!\bigl[\log d_m^{\mathrm{bin}}(W_m)\bigr].
\]
则有显式渐近公式
\[
\kappa_m
=
m\log\!\Bigl(\frac{2}{\varphi}\Bigr)
-\frac{\log\varphi}{1+\varphi^2}
+O\!\Bigl(\Bigl(\frac{\varphi}{2}\Bigr)^m\Bigr).
\]
若进一步记 \(g_m\) 为规范体积密度，则
\[
g_m
=
m\log\!\Bigl(\frac{2}{\varphi}\Bigr)
-\Bigl(1+\frac{\log\varphi}{1+\varphi^2}\Bigr)
+\frac{(\varphi/2)^m}{2\sqrt5}
\Bigl(
\varphi^2 m\log\!\Bigl(\frac{2}{\varphi}\Bigr)
+\varphi^2\log(2\pi)-\log\varphi
\Bigr)
+O\!\Bigl(\Bigl(\frac{\varphi}{2}\Bigr)^m\Bigr).
\]
\end{theorem}
\begin{proof}
由定理 \ref{thm:fold-bin-two-state-asymptotic} 的一致对数展开，
\[
\log d_m^{\mathrm{bin}}(w)
=
m\log\!\Bigl(\frac{2}{\varphi}\Bigr)
-w_m\log\varphi
+O\!\Bigl(\Bigl(\frac{\varphi}{2}\Bigr)^m\Bigr)
\]
在 \(X_m\) 上一致成立。对 \(\mu_m\) 取期望得
\[
\kappa_m
=
m\log\!\Bigl(\frac{2}{\varphi}\Bigr)
-\log\varphi\cdot \mu_m\!\bigl(W_m(m)=1\bigr)
+O\!\Bigl(\Bigl(\frac{\varphi}{2}\Bigr)^m\Bigr).
\]
再由命题 \ref{prop:fold-bin-escort-lastbit} 在 \(q=1\) 处的末位极限，
\[
\mu_m\!\bigl(W_m(m)=1\bigr)
=
\frac{1}{1+\varphi^2}
+O\!\Bigl(\Bigl(\frac{\varphi}{2}\Bigr)^m\Bigr),
\]
代回即得 \(\kappa_m\) 的闭式渐近。

另一方面，由命题 \ref{prop:fold-bin-gauge-density-two-scale}，
\[
g_m
=
\kappa_m-1
+\frac{(\varphi/2)^m}{2\sqrt5}
\Bigl(
\varphi^2 m\log\!\Bigl(\frac{2}{\varphi}\Bigr)
+\varphi^2\log(2\pi)-\log\varphi
\Bigr)
+O\!\Bigl(\Bigl(\frac{\varphi}{2}\Bigr)^m\Bigr).
\]
将上式中的 \(\kappa_m\) 代入，便得到所示 \(g_m\) 的全显式两尺度展开。证毕。
\end{proof}

\endinput


\begin{theorem}[与稳定层均匀分布的一样本最优判别常数]\label{thm:xi-fold-uniform-baseline-testing-constant}
\leanverified{paper\_xi\_fold\_uniform\_baseline\_testing\_constant}
沿用 \(\mu_m(w)=d_m^{\mathrm{bin}}(w)/2^m\) 的记号。令
\[
u_m(w):=\frac{1}{|X_m|}
\qquad (w\in X_m)
\]
为稳定层均匀分布，并记似然比
\[
R_m(w):=\frac{\mu_m(w)}{u_m(w)}.
\]
则在 \(u_m\) 下，\(R_m\) 依分布收敛到二点随机变量
\[
R=
\begin{cases}
\dfrac{\varphi^2}{\sqrt5},& \text{以概率 }\dfrac1\varphi,\\[8pt]
\dfrac{\varphi}{\sqrt5},& \text{以概率 }\dfrac1{\varphi^2}.
\end{cases}
\]
从而有精确极限
\[
\boxed{\;
\lim_{m\to\infty}D_{\mathrm{TV}}(\mu_m,u_m)
=
1-\frac{2}{\sqrt5}.\;}
\]
特别地，在等先验的一样本假设检验
\[
H_0:u_m,
\qquad
H_1:\mu_m
\]
中，最优 Bayes 错误概率满足
\[
\boxed{\;
\lim_{m\to\infty}P_e^\ast(\mu_m,u_m)
=
\frac1{\sqrt5}.\;}
\]
\end{theorem}
\begin{proof}
由定理 \ref{thm:conclusion-foldbin-likelihood-ratio-two-atom-transfer}，
\[
R_m(W_m)\Rightarrow R,
\qquad
W_m\sim u_m,
\]
并且对充分大的 \(m\)，随机变量 \(R_m(W_m)\) 落在某个固定闭区间 \(I\subset(0,\infty)\) 中。故
\[
D_{\mathrm{TV}}(\mu_m,u_m)
=
\frac12\,\mathbb E_{u_m}|R_m-1|
\longrightarrow
\frac12\,\mathbb E|R-1|.
\]
直接计算得到
\[
\frac12\,\mathbb E|R-1|
=
\frac12\left[
\frac1\varphi\Bigl(\frac{\varphi^2}{\sqrt5}-1\Bigr)
+
\frac1{\varphi^2}\Bigl(1-\frac{\varphi}{\sqrt5}\Bigr)
\right].
\]
这与定理 \ref{thm:xi-foldbin-pushforward-uniform-two-atom-ledger} 给出的全变差极限常数
\[
\frac12\left(\frac1{\sqrt5}-\varphi^{-3}\right)
\]
一致；再用恒等式 \(\varphi^{-3}=\sqrt5-2\)，即得
\[
\lim_{m\to\infty}D_{\mathrm{TV}}(\mu_m,u_m)=1-\frac{2}{\sqrt5}.
\]

最后，由 Helstrom 公式
\[
P_e^\ast(\mu_m,u_m)=\frac12\bigl(1-D_{\mathrm{TV}}(\mu_m,u_m)\bigr),
\]
令 \(m\to\infty\) 即得
\[
\lim_{m\to\infty}P_e^\ast(\mu_m,u_m)=\frac1{\sqrt5}.
\]
证毕。
\end{proof}

\begin{corollary}[相对于稳定层均匀基线的全部经典散度塌缩为单一 Bernoulli 几何]\label{cor:xi-fold-uniform-baseline-classical-fdiv-bernoulli-collapse}
\leanverified{paper\_xi\_fold\_uniform\_baseline\_classical\_fdiv\_bernoulli\_collapse}
沿用定理 \ref{thm:xi-fold-uniform-baseline-testing-constant} 的记号，并记
\[
q_*:=\varphi^{-2},
\qquad
p_*:=\frac{1}{\varphi\sqrt5}.
\]
则对任意凸函数 \(f:(0,\infty)\to\RR\) 满足 \(f(1)=0\)，有
\[
\boxed{\;
D_f(\mu_m\Vert u_m)
=
D_f\!\bigl(\mathrm{Bernoulli}(p_*)\Vert \mathrm{Bernoulli}(q_*)\bigr)
+O_f\!\Bigl(\Bigl(\frac{\varphi}{2}\Bigr)^m\Bigr).\;}
\]
特别地，
\[
\boxed{\;
\lim_{m\to\infty}D_{\mathrm{TV}}(\mu_m,u_m)
=
\left|p_*-q_*\right|
=
1-\frac{2}{\sqrt5},\;}
\]
\[
\boxed{\;
\lim_{m\to\infty}D_{\mathrm{KL}}(\mu_m\Vert u_m)
=
D_{\mathrm{KL}}\!\bigl(\mathrm{Bernoulli}(p_*)\Vert \mathrm{Bernoulli}(q_*)\bigr),\;}
\]
\[
\boxed{\;
\lim_{m\to\infty}\chi^2(\mu_m\Vert u_m)
=
\chi^2\!\bigl(\mathrm{Bernoulli}(p_*)\Vert \mathrm{Bernoulli}(q_*)\bigr).\;}
\]
换言之，\(\mu_m\) 相对于稳定层均匀基线 \(u_m\) 的全部经典 Csisz\'ar 几何，在极限中都严格塌缩为同一个二点 Bernoulli 对。
\end{corollary}
\begin{proof}
由定理 \ref{thm:xi-foldbin-pushforward-uniform-two-atom-ledger}，对任意满足条件的 \(f\) 都有
\[
D_f(\mu_m\Vert u_m)
=
\frac1\varphi\,f\!\left(\frac{\varphi^2}{\sqrt5}\right)
+\frac1{\varphi^2}\,f\!\left(\frac{\varphi}{\sqrt5}\right)
+O_f\!\Bigl(\Bigl(\frac{\varphi}{2}\Bigr)^m\Bigr).
\]
另一方面，
\[
1-q_*=\frac1\varphi,
\qquad
1-p_*=\frac{\varphi}{\sqrt5},
\]
故
\[
\frac{1-p_*}{1-q_*}
=
\frac{\varphi/\sqrt5}{1/\varphi}
=
\frac{\varphi^2}{\sqrt5},
\qquad
\frac{p_*}{q_*}
=
\frac{1/(\varphi\sqrt5)}{\varphi^{-2}}
=
\frac{\varphi}{\sqrt5}.
\]
这正说明上式等于
\[
D_f\!\bigl(\mathrm{Bernoulli}(p_*)\Vert \mathrm{Bernoulli}(q_*)\bigr)
+O_f\!\Bigl(\Bigl(\frac{\varphi}{2}\Bigr)^m\Bigr).
\]
对全变差、KL 与 \(\chi^2\) 只需分别取
\[
f_{\mathrm{TV}}(t)=\frac12|t-1|,
\qquad
f_{\mathrm{KL}}(t)=t\log t,
\qquad
f_{\chi^2}(t)=(t-1)^2.
\]
其中全变差的极限常数也可写作
\[
|p_*-q_*|
=
\varphi^{-2}-\frac{1}{\varphi\sqrt5}
=
1-\frac{2}{\sqrt5},
\]
与定理 \ref{thm:xi-fold-uniform-baseline-testing-constant} 一致。证毕。
\end{proof}

\begin{theorem}[零点几何对碰撞缺口的指数失明]\label{thm:xi-fold-zero-geometry-collision-gap-blindness}
沿用结论 \ref{con:xi-fold-collision-spectral-trace-parseval} 的记号，令
\[
M:=F_{m+2},
\qquad
\mathcal Z_m:=\{k\in \ZZ/(M\ZZ):\ \widehat d_m(k)=0\}.
\]
其中 \(C_\varphi>0\) 为定理 \ref{thm:fold-critical-resonance-constant} 的黄金共振常数。则碰撞缺口与零点密度满足严格分离：
\[
\boxed{\;
M\Bigl(\mathsf{Col}_m-\frac1M\Bigr)
\ge
2C_\varphi^2+o(1),
\qquad
\frac{|\mathcal Z_m|}{M}
\le
\tau(m+2)\,O(\varphi^{-m/2}).\;}
\]
从而
\[
\boxed{\;
\frac{\mathsf{Col}_m-\frac1M}{|\mathcal Z_m|/M^2}
=
\frac{M\bigl(\mathsf{Col}_m-\frac1M\bigr)}{|\mathcal Z_m|/M}
\longrightarrow +\infty.\;}
\]
更显式地，
\[
\boxed{\;
\frac{\mathsf{Col}_m-\frac1M}{|\mathcal Z_m|/M^2}
\gtrsim
\frac{2C_\varphi^2}{\tau(m+2)}\,\varphi^{m/2}.\;}
\]
\end{theorem}
\begin{proof}
由结论 \ref{con:xi-fold-chi2-constant-residual-energy}，
\[
M\Bigl(\mathsf{Col}_m-\frac1M\Bigr)\ge 2C_\varphi^2+o(1).
\]
另一方面，由定理 \ref{thm:fold-zero-density-sparse}，
\[
\frac{|\mathcal Z_m|}{M}
\le
\tau(m+2)\,\frac{F_{\lfloor (m+2)/2\rfloor}}{F_{m+2}}
=
\tau(m+2)\,O(\varphi^{-m/2}).
\]
将两式相除，即得
\[
\frac{M\bigl(\mathsf{Col}_m-\frac1M\bigr)}{|\mathcal Z_m|/M}
\to +\infty.
\]
而
\[
\frac{M\bigl(\mathsf{Col}_m-\frac1M\bigr)}{|\mathcal Z_m|/M}
=
\frac{\mathsf{Col}_m-\frac1M}{|\mathcal Z_m|/M^2},
\]
故两种写法完全等价。

最后，把上界
\(
|\mathcal Z_m|/M\le \tau(m+2)\,O(\varphi^{-m/2})
\)
写成倒数下界，再与
\(
M(\mathsf{Col}_m-1/M)\ge 2C_\varphi^2+o(1)
\)
合并，即得所示 \(\gtrsim\) 估计。证毕。
\end{proof}

\begin{theorem}[零点计数对碰撞缺口不存在消失模量]\label{cor:xi-fold-zero-density-no-bounded-control-on-normalized-collision}
沿用定理 \ref{thm:xi-fold-zero-geometry-collision-gap-blindness} 的记号，并定义
\[
M_m:=F_{m+2},
\qquad
\delta_m:=\frac{|\mathcal Z_m|}{M_m},
\qquad
\Gamma_m:=M_m\!\left(\mathsf{Col}_m-\frac{1}{M_m}\right).
\]
则沿子列
\[
m\equiv 1\pmod 3
\]
有
\[
\boxed{\;
\delta_m\longrightarrow 0,
\qquad
\liminf_{\substack{m\to\infty\\ m\equiv 1\!\!\!\!\pmod 3}}
\Gamma_m
\ge
\frac{2C_\varphi^2}{\varphi}>0.\;}
\]
从而不存在任何函数 \(\eta:[0,1]\to[0,\infty)\) 满足
\[
\eta(t)\longrightarrow 0
\qquad (t\to 0),
\]
并且对所有充分大的 \(m\equiv 1\pmod 3\) 都有
\[
\Gamma_m\le \eta(\delta_m).
\]
换言之，即使谱零点的相对密度趋于 \(0\)，它也不足以连续控制折叠碰撞的归一化缺口；非零频率振幅层仍是不可删除的数据层。
\end{theorem}
\begin{proof}
当 \(m\equiv 1\pmod 3\) 时，\(m+2\equiv 0\pmod 3\)，故 Fibonacci 数 \(M_m=F_{m+2}\) 为偶数。于是由定理 \ref{thm:fold-zero-density-sparse}，
\[
\delta_m=\frac{|\mathcal Z_m|}{M_m}
\le
\tau(m+2)\,O(\varphi^{-m/2})
\longrightarrow 0
\qquad
\bigl(m\to\infty,\ m\equiv 1\!\!\!\!\pmod 3\bigr).
\]

另一方面，定理 \ref{thm:xi-fold-zero-geometry-collision-gap-blindness} 已给出更强的全序列下界
\[
\Gamma_m
=
M_m\!\left(\mathsf{Col}_m-\frac1{M_m}\right)
\ge
2C_\varphi^2+o(1).
\]
因此沿任意子列，特别是沿 \(m\equiv 1\pmod 3\)，都有
\[
\liminf_{\substack{m\to\infty\\ m\equiv 1\!\!\!\!\pmod 3}}\Gamma_m
\ge
2C_\varphi^2
>
\frac{2C_\varphi^2}{\varphi},
\]
从而所示盒装下界成立。

若存在函数 \(\eta\) 满足 \(\eta(t)\to 0\) 且对所有充分大的 \(m\equiv 1\pmod 3\) 都有
\[
\Gamma_m\le \eta(\delta_m),
\]
则由 \(\delta_m\to 0\) 可得
\[
\eta(\delta_m)\longrightarrow 0
\qquad
\bigl(m\to\infty,\ m\equiv 1\!\!\!\!\pmod 3\bigr),
\]
从而 \(\Gamma_m\to 0\) 沿该子列成立，这与上式的严格正下极限矛盾。故这样的消失模量不存在。证毕。
\end{proof}

\endinput

\paragraph{末位统计充分性、黄金判别门槛与 fold-gauge 一次项的统一压缩}
在末位充分化、均匀基线二原子极限、规范群熵账本以及 fold-gauge 常数的 Bernoulli--Stirling 层级均已建立之后，还可把这些输入进一步压缩为一条直接服务于判别几何与规范体积审计的闭式链条。

\begin{theorem}[bin-fold 判别几何在末位上的统计充分性塌缩]\label{thm:xi-fold-lastbit-statistical-sufficiency-collapse}
\leanverified{paper\_xi\_fold\_lastbit\_statistical\_sufficiency\_collapse}
设
\[
\mu_m(w):=\frac{d_m^{\mathrm{bin}}(w)}{2^m},
\qquad
u_m(w):=\frac{1}{|X_m|}
\qquad (w\in X_m),
\]
并按末位分层
\[
A_i:=\{w\in X_m:\ w_m=i\}
\qquad (i=0,1).
\]
定义粗化测度
\[
\bar\mu_m
:=
\mu_m(A_0)\,u_m(\cdot\mid A_0)
\;+\;
\mu_m(A_1)\,u_m(\cdot\mid A_1).
\]
则存在常数 \(C>0\) 使
\[
\boxed{\;
D_{\mathrm{TV}}(\mu_m,\bar\mu_m)
\le
C\Bigl(\frac{\varphi}{2}\Bigr)^m.\;}
\]
特别地，统计实验
\[
(X_m,\mu_m,u_m)
\]
与其二元末位压缩
\[
\bigl(\{0,1\},(\ell_m)_\#\mu_m,(\ell_m)_\#u_m\bigr),
\qquad
\ell_m(w):=w_m,
\]
在全变差意义下指数等价；因而末位统计量 \(\ell_m\) 对 bin-fold 非均匀性是渐近充分的。
\end{theorem}
\begin{proof}
记
\[
K_m(i,\cdot):=\mathrm{Unif}(A_i)
\qquad (i=0,1).
\]
则
\[
K_m\circ (\ell_m)_\#\mu_m
=
\mu_m(A_0)\,u_m(\cdot\mid A_0)
+\mu_m(A_1)\,u_m(\cdot\mid A_1)
=
\bar\mu_m.
\]
而定理 \ref{thm:xi-fold-lastbit-asymptotic-sufficiency} 已给出
\[
D_{\mathrm{TV}}\!\bigl(\mu_m,K_m\circ(\ell_m)_\#\mu_m\bigr)
=
O\!\Bigl(\Bigl(\frac{\varphi}{2}\Bigr)^m\Bigr).
\]
将右端常数记为 \(C\) 即得所示估计；指数等价与渐近充分性亦随之成立。证毕。
\end{proof}

\begin{corollary}[相对稳定层均匀基线的黄金判别门槛]\label{cor:xi-fold-uniform-baseline-golden-threshold}
沿用定理 \ref{thm:xi-fold-lastbit-statistical-sufficiency-collapse} 的记号，则
\[
u_m(A_0)=\frac{F_{m+1}}{F_{m+2}}\longrightarrow \varphi^{-1},
\qquad
u_m(A_1)=\frac{F_m}{F_{m+2}}\longrightarrow \varphi^{-2},
\]
并且
\[
\mu_m(A_0)\longrightarrow \frac{\varphi}{\sqrt5},
\qquad
\mu_m(A_1)\longrightarrow \frac{1}{\varphi\sqrt5}.
\]
从而
\[
\boxed{\;
\lim_{m\to\infty}D_{\mathrm{TV}}(\mu_m,u_m)
=
1-\frac{2}{\sqrt5}.\;}
\]
若记等先验下判别 \(\mu_m\) 与 \(u_m\) 的最优 Bayes 错误概率为
\[
\mathsf P_{\mathrm{err}}^\ast(m)
:=
\frac{1-D_{\mathrm{TV}}(\mu_m,u_m)}{2},
\]
则
\[
\boxed{\;
\lim_{m\to\infty}\mathsf P_{\mathrm{err}}^\ast(m)
=
\frac{1}{\sqrt5}.\;}
\]
\end{corollary}
\begin{proof}
由定理 \ref{thm:xi-fold-lastbit-statistical-sufficiency-collapse}，
\[
D_{\mathrm{TV}}(\mu_m,u_m)
=
D_{\mathrm{TV}}(\bar\mu_m,u_m)
+O\!\Bigl(\Bigl(\frac{\varphi}{2}\Bigr)^m\Bigr).
\]
又因为 \(\bar\mu_m\) 与 \(u_m\) 在分层 \(A_0\sqcup A_1\) 上均为块常值分布，故
\[
D_{\mathrm{TV}}(\bar\mu_m,u_m)
=
\bigl|\mu_m(A_0)-u_m(A_0)\bigr|.
\]
其中 \(u_m(A_0)=|A_0|/|X_m|=F_{m+1}/F_{m+2}\) 与
\(
u_m(A_1)=F_m/F_{m+2}
\)
由 Fibonacci 计数立得；而
\(
\mu_m(A_0),\mu_m(A_1)
\)
的极限可由定理 \ref{thm:xi-fold-uniform-baseline-testing-constant}，或等价地由
定理 \ref{thm:xi-foldbin-pushforward-uniform-two-atom-ledger} 的二原子极限直接得到。于是
\[
\lim_{m\to\infty}D_{\mathrm{TV}}(\mu_m,u_m)
=
\left|\frac{\varphi}{\sqrt5}-\frac1\varphi\right|
=
1-\frac{2}{\sqrt5}.
\]
最后把该极限代入 Bayes 公式即可得到
\(
\lim_{m\to\infty}\mathsf P_{\mathrm{err}}^\ast(m)=1/\sqrt5
\)。
证毕。
\end{proof}

\leanverified{paper\_xi\_fold\_uniform\_baseline\_golden\_threshold}

\begin{theorem}[纤维信息项与规范体积密度的一次项显式化]\label{thm:xi-fold-kappa-gauge-first-order-constants}
\leanverified{paper\_xi\_fold\_kappa\_gauge\_first\_order\_constants}
令
\[
\kappa_m:=\sum_{w\in X_m}\mu_m(w)\log d_m^{\mathrm{bin}}(w),
\qquad
g_m:=2^{-m}\log |G_m^{\mathrm{bin}}|.
\]
则当 \(m\to\infty\) 时有
\[
\boxed{\;
\kappa_m
=
m\log\!\Bigl(\frac{2}{\varphi}\Bigr)
-\frac{\log\varphi}{\varphi\sqrt5}
+O\!\Bigl(\Bigl(\frac{\varphi}{2}\Bigr)^m\Bigr),\;}
\]
以及
\[
\boxed{\;
g_m
=
m\log\!\Bigl(\frac{2}{\varphi}\Bigr)
-1
-\frac{\log\varphi}{\varphi\sqrt5}
+\frac{1}{2\sqrt5}\Bigl(\frac{\varphi}{2}\Bigr)^m
\Bigl(
\varphi^2 m\log\!\Bigl(\frac{2}{\varphi}\Bigr)
+\varphi^2\log(2\pi)-\log\varphi
\Bigr)
+O\!\Bigl(\Bigl(\frac{\varphi}{2}\Bigr)^m\Bigr).\;}
\]
因此 \(\kappa_m\) 与 \(g_m\) 具有相同的一次主斜率
\(
\log(2/\varphi)
\)，
且其常数项之差严格等于 \(1\)。
\end{theorem}
\begin{proof}
第一式正是推论 \ref{cor:xi-foldbin-shannon-gauge-constant-mirror} 中
\(\kappa_m\) 的显式展开。

再由定理 \ref{thm:xi-foldbin-gauge-entropy-one-nat-law}，
\[
g_m
=
\kappa_m-1
+\frac{1}{2\sqrt5}\Bigl(\frac{\varphi}{2}\Bigr)^m
\Bigl(
\varphi^2 m\log\!\Bigl(\frac{2}{\varphi}\Bigr)
+\varphi^2\log(2\pi)-\log\varphi
\Bigr)
+O\!\Bigl(\Bigl(\frac{\varphi}{2}\Bigr)^m\Bigr).
\]
把 \(\kappa_m\) 的展开代入，即得第二式。证毕。
\end{proof}

\begin{theorem}[fold-gauge 的 \texorpdfstring{$2\pi$}{2pi} 估计序列具有最终单侧性与黄金公比锁定]\label{thm:xi-fold-gauge-2pi-estimator-geometric-lock}
\leanverified{paper\_xi\_fold\_gauge\_2pi\_estimator\_geometric\_lock}
沿用推论 \ref{cor:fold-bin-recover-2pi} 的规范常数列 \((C_m)_{m\ge 1}\)，并记
\[
E_m:=\log C_m-\log(2\pi).
\]
则存在 \(m_0\)，使得对全部 \(m\ge m_0\) 都有
\[
C_m>2\pi,
\qquad
C_{m+1}<C_m.
\]
更具体地，
\[
\boxed{\;
E_m
=
\frac{1}{3\varphi}\Bigl(\frac{\varphi}{2}\Bigr)^m
+O\!\Bigl(\Bigl(\frac{\varphi}{2}\Bigr)^{3m}\Bigr),\;}
\]
从而
\[
\boxed{\;
\frac{E_{m+1}}{E_m}
\longrightarrow
\frac{\varphi}{2}.\;}
\]
换言之，\(C_m\) 不仅收敛到 \(2\pi\)，而且其对数误差尾部形成公比为 \(\varphi/2\) 的黄金几何级数。
\end{theorem}
\begin{proof}
由定理 \ref{thm:xi-foldbin-gauge-constant-eventual-complete-monotonicity}，
对充分大的 \(m\) 已有
\[
C_m>2\pi,
\qquad
C_{m+1}<C_m.
\]
同一定理还给出更强展开
\[
E_m
=
\frac{1}{3\varphi}\Bigl(\frac{\varphi}{2}\Bigr)^m
-\frac{\sqrt5}{180}\Bigl(\frac{\varphi}{2}\Bigr)^{3m}
+O\!\Bigl(\Bigl(\frac{\varphi}{2}\Bigr)^{5m}\Bigr),
\]
故特别可写成所示的首项加三次余项形式。于是
\[
\frac{E_{m+1}}{E_m}
=
\frac{\frac{1}{3\varphi}(\varphi/2)^{m+1}+O((\varphi/2)^{3m+3})}
{\frac{1}{3\varphi}(\varphi/2)^m+O((\varphi/2)^{3m})}
\longrightarrow
\frac{\varphi}{2}.
\]
证毕。
\end{proof}

\endinput

\paragraph{window-\texorpdfstring{$6$}{6} 群商实现的稳定子分层与 Lee--Yang 振幅常数的双二次阴影}
在 window-\(6\) 的纤维退化直方图、自由作用与线性核障碍已经锁定之后，任何试图把纤维分解解释为有限群轨道分解的方案都必须携带多层稳定子；另一方面，Lee--Yang 振幅常数 \(B\) 的真正六次性并不止于 \(\QQ(B)/\QQ\) 的次数陈述，其分裂域还显现出一条与三次分辨子域不同的额外二次阴影。

\begin{theorem}[window-\texorpdfstring{$6$}{6} 的任何群商实现都强制三层稳定子分层]\label{thm:xi-window6-group-quotient-isotropy-strata}
设 \(\Omega_6:=\{0,1\}^6\)，并令
\[
F_6:=\Fold^{\mathrm{bin}}_6:\Omega_6\to X_6.
\]
若存在某个有限群 \(G\) 在 \(\Omega_6\) 上作用，使其轨道分解恰等于 \(F_6\) 的纤维分解，则
\[
\boxed{\;12\mid |G|.\;}
\]
并且轨道稳定子必分成三种不同的阶层：对每个 \(\omega\in\Omega_6\)，有
\[
|\mathrm{Orb}_G(\omega)|\in\{2,3,4\}
\quad\Longleftrightarrow\quad
|\mathrm{Stab}_G(\omega)|\in\left\{\frac{|G|}{2},\frac{|G|}{3},\frac{|G|}{4}\right\}.
\]
更具体地，大小分别为 \(2,3,4\) 的轨道个数恰为
\[
8,\qquad 4,\qquad 9.
\]
因此任何这样的群商实现都不可能是自由作用，而只能是带有三层稳定子分层的离散 orbifold 型商。
\leanverified{paper\_xi\_window6\_group\_quotient\_isotropy\_strata}
\end{theorem}
\begin{proof}
由推论 \ref{cor:foldbin6-degeneracy-multivalued-nonfree}，\(F_6\) 的纤维划分既不可能来自自由有限群作用，也不可能来自阿贝尔线性陪集划分。另一方面，由推论 \ref{cor:fold-bin-m6-fiber-hist}，
\[
|\{w\in X_6:\ |F_6^{-1}(w)|=2\}|=8,\qquad
|\{w\in X_6:\ |F_6^{-1}(w)|=3\}|=4,\qquad
|\{w\in X_6:\ |F_6^{-1}(w)|=4\}|=9.
\]
若群 \(G\) 的轨道分解恰与纤维分解一致，则轨道大小只能取 \(2,3,4\)，且对应轨道数正是上述 \(8,4,9\)。

对任意 \(\omega\in\Omega_6\)，轨道--稳定子定理给出
\[
|\mathrm{Orb}_G(\omega)|=\frac{|G|}{|\mathrm{Stab}_G(\omega)|}.
\]
因此 \(2,3,4\) 都必须整除 \(|G|\)，遂有
\[
\mathrm{lcm}(2,3,4)=12\mid |G|.
\]
并且当轨道大小分别为 \(2,3,4\) 时，稳定子阶只能分别为
\[
\frac{|G|}{2},\qquad \frac{|G|}{3},\qquad \frac{|G|}{4}.
\]
这三者彼此不同，故点稳定子至少分成三层；又因为全部轨道只出现这三种大小，所以稳定子阶也恰只出现这三层。证毕。
\end{proof}

\begin{theorem}[Lee--Yang 振幅常数 \texorpdfstring{$B$}{B} 的六次性与双二次阴影]\label{thm:xi-leyang-amplitude-B-sextic-biquadratic-shadow}
设
\[
f_B(X):=162X^6+1593X^4+1998X^2+1184.
\]
则
\[
\boxed{\;[\QQ(B):\QQ]=6.\;}
\]
并且 \(f_B\) 在 \(\QQ\) 上不可约，正是 \(B\) 的六次最小多项式；其判别式满足
\[
\boxed{\;
\Disc(f_B)=
-2^{16}\cdot 3^{22}\cdot 11^4\cdot 37^3\cdot 109^4,\;}
\]
故平方自由部分为
\[
\boxed{\;-37.\;}
\]
若 \(N\) 为 \(f_B\) 的分裂域，则
\[
\QQ(\sqrt{-37})\subset N.
\]
另一方面，\(N\) 同时包含 \(B^2\) 的三次最小多项式的分裂域，因而也包含其唯一二次分辨子域
\[
\QQ(\sqrt{-111}).
\]
于是
\[
\boxed{\;
\QQ(\sqrt{-37},\sqrt{-111})\subset N,
\qquad
12\mid [N:\QQ].\;}
\]
\end{theorem}
\begin{proof}
由定理 \ref{thm:xi-leyang-amplitude-B-true-quadratic-cover}，
\[
[\QQ(B):\QQ]=6.
\]
又因 \(B^2\) 满足三次方程
\[
g(T):=162T^3+1593T^2+1998T+1184,
\]
故 \(g(B^2)=0\)，亦即 \(f_B(B)=g(B^2)=0\)。由于 \(f_B\) 的次数也是 \(6\)，它必为 \(B\) 的最小多项式，从而在 \(\QQ\) 上不可约。

设 \(r_1,r_2,r_3\) 为 \(g\) 的根，则
\[
f_B(X)=162\prod_{j=1}^{3}(X^2-r_j),
\]
其六个根为 \(\pm \sqrt{r_1},\pm \sqrt{r_2},\pm \sqrt{r_3}\)。于是
\[
\Disc(f_B)
=
162^{10}\Bigl(\prod_{j=1}^{3}4r_j\Bigr)
\prod_{1\le i<j\le 3}(r_i-r_j)^4.
\]
另一方面，
\[
\Disc(g)=162^4\prod_{1\le i<j\le 3}(r_i-r_j)^2,
\qquad
\prod_{j=1}^{3}r_j=-\frac{1184}{162}.
\]
故可化为
\[
\Disc(f_B)
=
-64\cdot 162\cdot 1184\cdot \Disc(g)^2.
\]
再由推论 \ref{cor:xi-leyang-u-b-square-common-quadratic-resolvent}，
\[
\Disc(g)=
-2^2\cdot 3^9\cdot 11^2\cdot 37\cdot 109^2.
\]
代入即得
\[
\Disc(f_B)=
-2^{16}\cdot 3^{22}\cdot 11^4\cdot 37^3\cdot 109^4.
\]
因此其平方自由部分为 \(-37\)。分裂域 \(N\) 因而包含
\[
\QQ(\sqrt{\Disc(f_B)})=\QQ(\sqrt{-37}).
\]

再者，\(N\) 含有 \(f_B\) 的全部根 \(\pm\sqrt{r_j}\)，平方后即得到 \(g\) 的全部根 \(r_j\)，故 \(g\) 的分裂域 \(L_b\) 满足
\[
L_b\subset N.
\]
由推论 \ref{cor:xi-leyang-u-b-square-common-quadratic-resolvent}，\(L_b\) 的唯一二次子域是
\[
\QQ(\sqrt{-111}).
\]
于是 \(N\) 同时包含 \(\QQ(\sqrt{-37})\) 与 \(\QQ(\sqrt{-111})\)。由于这两个虚二次域不同，遂得
\[
\QQ(\sqrt{-37},\sqrt{-111})\subset N.
\]
又因 \([L_b:\QQ]=6\)，且 \(\QQ(\sqrt{-37})\not\subset L_b\)（因为 \(L_b\) 的唯一二次子域为 \(\QQ(\sqrt{-111})\)），故
\[
[L_b(\sqrt{-37}):\QQ]=12.
\]
最后注意 \(L_b(\sqrt{-37})\subset N\)，于是
\[
12\mid [N:\QQ].
\]
证毕。
\end{proof}

\endinput

\paragraph{截断 Jensen 剖面、solenoid 终对象与有限可见证书的双重不完备性}
在 Jensen 缺陷的分布反演与有限半径绝对盲区、有限素数局部化群与其对偶 solenoid 的 \(\mathrm{Hom}\)-范畴分类，以及 prime-register 轴的无限外置语义已经建立之后，还可把“半径方向的失明”与“素数方向的失明”压缩为下列五条闭合结论。

\begin{theorem}[截断 Jensen 剖面的完全性与完全失明的精确分界]\label{thm:xi-time-part60e-truncated-jensen-profile-completeness-blindness}
\leanverified{paper\_xi\_time\_part60e\_truncated\_jensen\_profile\_completeness\_blindness}
设 \(F\) 为单位圆盘 \(\mathbb D\) 上非常值全纯函数，零点按重数记为 \(\{a\}\subset\mathbb D\)，并定义 Jensen 缺陷
\[
D_F(\rho):=\sum_{|a|<\rho}\log\frac{\rho}{|a|}
\qquad (0<\rho<1).
\]
固定 \(0<\rho_\star<1\)，并记截断剖面
\[
\mathcal J_{\rho_\star}(F):=D_F|_{(0,\rho_\star]}.
\]
再定义径向零测度
\[
\mu_F^{\mathrm{rad}}
:=
\sum_a \operatorname{mult}(a)\,\delta_{\log|a|}
\qquad\text{于 }(-\infty,0).
\]
则成立：
\begin{enumerate}
  \item \(\mathcal J_{\rho_\star}(F)\) 完全决定，且仅决定，
  \[
  \mu_F^{\mathrm{rad}}\big|_{(-\infty,\log\rho_\star)}.
  \]
  \item 对任意两个盘内全纯函数 \(F,G\)，有等价关系
  \[
  \mathcal J_{\rho_\star}(F)=\mathcal J_{\rho_\star}(G)
  \iff
  \mu_F^{\mathrm{rad}}\big|_{(-\infty,\log\rho_\star)}
  =
  \mu_G^{\mathrm{rad}}\big|_{(-\infty,\log\rho_\star)}.
  \]
  \item 因而，对零点多重集在闭环带
  \[
  \{z\in\mathbb D:\ \rho_\star\le |z|<1\}
  \]
  上所作的任意增删，都不会改变整个截断剖面 \(\mathcal J_{\rho_\star}\)。
\end{enumerate}
\end{theorem}
\begin{proof}
置
\[
u:=\log\rho<0,
\qquad
u_\star:=\log\rho_\star,
\qquad
\mathcal D_F(u):=D_F(e^u).
\]
由定理 \ref{thm:xi-time-part60a-jensen-defect-radial-measure-inversion}，
\[
\mathcal D_F(u)=\int_{(-\infty,u)}(u-s)\,\dd\mu_F^{\mathrm{rad}}(s),
\qquad
\mathcal D_F''=\mu_F^{\mathrm{rad}}
\]
在分布意义下于 \((-\infty,0)\) 成立。

因此，一旦 \(\mathcal D_F\) 已知于区间 \((-\infty,u_\star]\)，便可在开区间 \((-\infty,u_\star)\) 上取其分布二阶导数，从而唯一恢复
\[
\mu_F^{\mathrm{rad}}\big|_{(-\infty,u_\star)}.
\]
反之，若该限制测度已知，则对每个 \(u\le u_\star\)，势函数表示
\[
\mathcal D_F(u)=\int_{(-\infty,u)}(u-s)\,\dd\mu_F^{\mathrm{rad}}(s)
\]
只依赖于 \((-\infty,u_\star)\) 上的测度，因为积分核中始终要求 \(s<u\le u_\star\)。故第 \(1\) 条成立。

第 \(2\) 条是第 \(1\) 条的直接重述：若两条截断剖面相同，则其在 \((-\infty,u_\star)\) 上的分布二阶导数相同；反向由同一势函数公式立即得到。

对第 \(3\) 条，只需注意：若某个零点满足 \(|a|\ge \rho_\star\)，则对任意 \(\rho\le \rho_\star\) 都有 \(|a|<\rho\) 失效，故该零点在
\[
D_F(\rho)=\sum_{|a|<\rho}\log\frac{\rho}{|a|}
\]
中从不出现。因此它对整个 \((0,\rho_\star]\) 上的截断剖面完全不可见。证毕。
\end{proof}

\begin{theorem}[solenoid 终对象之间连续态射的完全分类]\label{thm:xi-time-part60e-solenoid-terminal-morphism-classification}
\leanverified{paper\_xi\_time\_part60e\_solenoid\_terminal\_morphism\_classification}
设 \(S,T\subset\mathbb P\) 为有限素数集，并记
\[
G_S:=\ZZ[S^{-1}],
\qquad
G_T:=\ZZ[T^{-1}],
\qquad
\Sigma_S:=\widehat{G_S},
\qquad
\Sigma_T:=\widehat{G_T}.
\]
则存在自然同构
\[
\operatorname{Hom}_{\mathrm{cts}}(\Sigma_T,\Sigma_S)
\cong
\operatorname{Hom}_{\ZZ}(G_S,G_T)
\cong
\begin{cases}
0,& S\not\subseteq T,\\[1mm]
G_T,& S\subseteq T.
\end{cases}
\]
更具体地：
\begin{enumerate}
  \item 存在非零连续群同态 \(\Sigma_T\to\Sigma_S\) 当且仅当 \(S\subseteq T\)；
  \item 在 \(S\subseteq T\) 时，每个连续同态都唯一对应于某个 \(x\in G_T\)，其对偶映射为
  \[
  m_x:G_S\to G_T,
  \qquad
  y\mapsto xy;
  \]
  \item 对应连续同态为满射当且仅当 \(x\neq 0\)；
  \item 对应连续同态为同构当且仅当 \(S=T\) 且 \(x\in G_T^\times\) 为 \(T\)-局部化单位。
\end{enumerate}
\end{theorem}
\begin{proof}
Pontryagin 对偶把紧交换群范畴与离散交换群范畴反变等价，因此
\[
\operatorname{Hom}_{\mathrm{cts}}(\Sigma_T,\Sigma_S)
\cong
\operatorname{Hom}_{\ZZ}(G_S,G_T).
\]
再由定理 \ref{thm:xi-cdim-localization-hom-category-classification}，后者恰在 \(S\subseteq T\) 时同构于 \(G_T\)，而在 \(S\not\subseteq T\) 时为零群。这给出前两条。

当 \(S\subseteq T\) 时，\(x\in G_T\) 对应的离散端同态正是
\[
m_x:G_S\to G_T,\qquad y\mapsto xy.
\]
由于 \(G_T\subset\QQ\) 无挠，\(m_x\) 在 \(x\neq 0\) 时为单射；对偶化后得到连续同态 \(\widehat{m_x}:\Sigma_T\to\Sigma_S\) 为满射。若 \(x=0\)，则对应映射显然为零映射，非满。故第 \(3\) 条成立。

最后，连续同态 \(\widehat{m_x}\) 为同构当且仅当 \(m_x\) 为离散群同构。由定理 \ref{thm:xi-cdim-localization-hom-category-classification}，这首先迫使 \(S=T\)；而当 \(S=T\) 时，\(m_x\) 为同构当且仅当 \(x\) 在局部化环 \(\ZZ[T^{-1}]\) 中可逆，即 \(x\in G_T^\times\)。证毕。
\end{proof}

\begin{corollary}[无穷 prime-register solenoid 的有限可见因子化]\label{cor:xi-time-part60e-infinite-solenoid-finite-visible-factorization}
\leanverified{paper\_xi\_time\_part60e\_infinite\_solenoid\_finite\_visible\_factorization}
设 \(S\subset\mathbb P\) 为可数无穷素数集，并记
\[
G_S:=\ZZ[S^{-1}],
\qquad
\Sigma_S:=\widehat{G_S}.
\]
则对任意有限子集 \(S_0\subset S\)，包含映射 \(G_{S_0}\hookrightarrow G_S\) 对偶化得到自然满射
\[
\pi_{S\to S_0}:\Sigma_S\twoheadrightarrow \Sigma_{S_0}.
\]
并且成立：
\begin{enumerate}
  \item 任意连续特征
  \[
  \chi:\Sigma_S\to \TT
  \]
  都经过某个有限子集 \(S_0\subset S\) 因子化，即存在
  \[
  \chi_0:\Sigma_{S_0}\to \TT
  \]
  使 \(\chi=\chi_0\circ\pi_{S\to S_0}\)。
  \item 更一般地，若 \(K\) 为紧 Abel Lie 群，则任意连续群同态
  \[
  f:\Sigma_S\to K
  \]
  都经过某个有限子集 \(S_0\subset S\) 因子化。
  \item 对任意有限子集 \(S_0\subset S\)，存在自然短正合列
  \[
  0\longrightarrow \prod_{p\in S\setminus S_0}\ZZ_p
  \longrightarrow \Sigma_S
  \xrightarrow{\ \pi_{S\to S_0}\ }
  \Sigma_{S_0}
  \longrightarrow 0.
  \]
  因而，任何有限族连续可见同态都不可能穷尽无穷 prime-register 核。
\end{enumerate}
\end{corollary}
\begin{proof}
由定义，
\[
G_S=\ZZ[S^{-1}]=\bigcup_{\substack{S_0\subset S\\ |S_0|<\infty}}G_{S_0}.
\]
因此每个 \(x\in G_S\) 的分母素数支撑都只落在某个有限子集 \(S_0\subset S\) 内。

对第 \(1\) 条，Pontryagin 对偶给出自然同构
\[
\operatorname{Hom}_{\mathrm{cts}}(\Sigma_S,\TT)\cong G_S.
\]
任取连续特征 \(\chi\)，令 \(x\in G_S\) 为其对偶元素。由上一步存在有限 \(S_0\subset S\) 使 \(x\in G_{S_0}\)。再对偶化有限情形的定理 \ref{thm:xi-time-part60e-solenoid-terminal-morphism-classification}，便得到 \(\chi\) 经由 \(\pi_{S\to S_0}\) 因子化。

对第 \(2\) 条，设 \(f:\Sigma_S\to K\) 为连续同态。由于 \(K\) 为紧 Abel Lie 群，其对偶群 \(\widehat K\) 为有限生成离散阿贝尔群。于是 \(f^\vee:\widehat K\to G_S\) 的像由有限多个元素生成，而每个生成元都落在某个有限 \(G_{S_j}\) 中；取这些 \(S_j\) 的并得有限集合 \(S_0\subset S\)，便有
\[
\operatorname{im}(f^\vee)\subseteq G_{S_0}.
\]
故 \(f^\vee\) 经由包含 \(G_{S_0}\hookrightarrow G_S\) 因子化。再对偶化即得
\[
f=f_0\circ \pi_{S\to S_0}
\]
对某个连续同态 \(f_0:\Sigma_{S_0}\to K\) 成立。

对第 \(3\) 条，只需计算离散端商群。由
\[
G_S/\ZZ\cong \bigoplus_{p\in S}C_{p^\infty},
\qquad
G_{S_0}/\ZZ\cong \bigoplus_{p\in S_0}C_{p^\infty},
\]
可得
\[
G_S/G_{S_0}
\cong
\bigoplus_{p\in S\setminus S_0}C_{p^\infty}.
\]
再取 Pontryagin 对偶，并用
\[
\widehat{C_{p^\infty}}\cong \ZZ_p,
\qquad
\widehat{\bigoplus_{p\in S\setminus S_0}C_{p^\infty}}
\cong
\prod_{p\in S\setminus S_0}\ZZ_p,
\]
便得到所示短正合列。

若给定有限族连续可见同态 \(f_j:\Sigma_S\to K_j\)（其中每个 \(K_j\) 都是紧 Abel Lie 群），则由第 \(2\) 条可取有限 \(S_j\subset S\) 使每个 \(f_j\) 都经由 \(\pi_{S\to S_j}\) 因子化。令 \(S_0:=\bigcup_j S_j\)，则全部 \(f_j\) 共同经由 \(\pi_{S\to S_0}\) 因子化。由于 \(S\setminus S_0\) 仍为无穷集，核
\[
\ker(\pi_{S\to S_0})\cong \prod_{p\in S\setminus S_0}\ZZ_p
\]
为无限紧群，因此这些有限可见同态不可能把 \(\Sigma_S\) 分离到有限歧义。证毕。
\end{proof}

\begin{theorem}[有限半径证书与有限素数证书的统一不可完备性]\label{thm:xi-time-part60e-finite-radius-finite-prime-certificate-incompleteness}
固定 \(0<\rho_\star<1\)，并取可数无穷素数集 \(S\subset\mathbb P\) 与有限子集 \(S_0\subset S\)。考虑对象对
\[
(F,\xi),
\]
其中 \(F\) 为单位圆盘上的有限 Blaschke 乘积，\(\xi\in \Sigma_S\)。定义联合有限证书
\[
\mathcal C_{\rho_\star,S_0}(F,\xi)
:=
\Bigl(
\mathcal J_{\rho_\star}(F),\ \pi_{S\to S_0}(\xi)
\Bigr).
\]
则每个非空纤维
\[
\mathcal C_{\rho_\star,S_0}^{-1}(\eta,u)
\]
至少是可数无限集。更精确地，对任意固定基对象 \((F_0,\xi_0)\)，存在一列对象
\[
\{(F_n,\xi_n)\}_{n\ge 1}
\]
使得：
\begin{enumerate}
  \item \(\mathcal C_{\rho_\star,S_0}(F_n,\xi_n)\) 全部相同；
  \item \(F_n\) 两两不同；
  \item \(\xi_n\) 两两不同。
\end{enumerate}
\end{theorem}
\begin{proof}
固定任意基对象 \((F_0,\xi_0)\)。

先处理 Jensen 侧。取一列两两不同的实数
\[
a_n\in(\rho_\star,1)
\qquad(n\ge 1),
\]
并定义有限 Blaschke 因子
\[
B_{a_n}(z):=\frac{z-a_n}{1-a_n z},
\qquad
F_n:=F_0\cdot B_{a_n}.
\]
由于每个新零点 \(a_n\) 都满足 \(|a_n|\ge \rho_\star\)，定理 \ref{thm:xi-time-part60e-truncated-jensen-profile-completeness-blindness} 第 \(3\) 条给出
\[
\mathcal J_{\rho_\star}(F_n)=\mathcal J_{\rho_\star}(F_0)
\qquad(n\ge 1).
\]
而 \(\{a_n\}\) 两两不同，故 \(F_n\) 两两不同。

再处理 arithmetic 侧。由推论 \ref{cor:xi-time-part60e-infinite-solenoid-finite-visible-factorization}，
\[
\ker(\pi_{S\to S_0})
\cong
\prod_{p\in S\setminus S_0}\ZZ_p
\]
为无限紧阿贝尔群。任选一列两两不同的元素
\[
\eta_n\in \ker(\pi_{S\to S_0})
\qquad(n\ge 1),
\]
并置
\[
\xi_n:=\xi_0+\eta_n.
\]
则
\[
\pi_{S\to S_0}(\xi_n)=\pi_{S\to S_0}(\xi_0)
\qquad(n\ge 1),
\]
且 \(\xi_n\) 两两不同。

于是对全部 \(n\ge 1\)，都有
\[
\mathcal C_{\rho_\star,S_0}(F_n,\xi_n)
=
\mathcal C_{\rho_\star,S_0}(F_0,\xi_0).
\]
这就构造出同一纤维中的可数无限列，故每个非空纤维至少可数无限。证毕。
\end{proof}

\begin{theorem}[有限可见统计塌缩与无限算术核强制的统一二分律]\label{thm:xi-time-part60e-finite-visible-collapse-infinite-padic-ledger-dichotomy}
固定可数无穷素数集 \(S\subset\mathbb P\)，并定义径向--算术对象类
\[
\mathfrak O_S^{\mathrm{rad}}
:=
\Bigl\{
(\mu,\xi):
\mu \text{ 为某个有限 Blaschke 乘积的径向零测度， }
\xi\in\Sigma_S
\Bigr\}.
\]
对任意这类 \(\mu\) 与任意 \(0<\rho_\star<1\)，定义其截断 Jensen 势函数
\[
\mathcal J_{\rho_\star}(\mu)(\rho)
:=
\int_{(-\infty,\log\rho)}(\log\rho-s)\,\dd\mu(s)
\qquad (0<\rho\le \rho_\star).
\]
则成立：
\begin{enumerate}
  \item 任意经过某个有限证书
  \[
  (\mu,\xi)\longmapsto
  \bigl(\mathcal J_{\rho_\star}(\mu),\ \pi_{S\to S_0}(\xi)\bigr),
  \qquad
  0<\rho_\star<1,\ \ S_0\subset S\ \text{有限},
  \]
  的可见不变量，都不是 \(\mathfrak O_S^{\mathrm{rad}}\) 上的完备不变量；事实上，在有限 Blaschke 子类上，每个非空证书纤维都至少可数无限。
  \item 若要得到径向--算术层面的完备可逆证书，则在解析侧必须保留全连续 Jensen 剖面
  \[
  \rho\longmapsto
  \int_{(-\infty,\log\rho)}(\log\rho-s)\,\dd\mu(s)
  \qquad(0<\rho<1),
  \]
  在算术侧必须保留全部有限素数投影
  \[
  \{\pi_{S\to S_0}(\xi)\}_{S_0\subset S,\ |S_0|<\infty}.
  \]
  等价地，算术侧必须保留整个 \(\xi\in\Sigma_S\)。
  \item 因而，有限可见统计在径向半径方向与素数方向上同时塌缩；而完备可逆账本在算术侧被迫是无限 \(p\)-进的。
\end{enumerate}
\end{theorem}
\begin{proof}
第 \(1\) 条由定理 \ref{thm:xi-time-part60e-finite-radius-finite-prime-certificate-incompleteness} 直接给出；该定理甚至已经在有限 Blaschke 子类上给出显式的可数无限纤维。

对第 \(2\) 条，取 \(\mu=\mu_F^{\mathrm{rad}}\) 为某个有限 Blaschke 乘积 \(F\) 的径向零测度。由定理 \ref{thm:xi-time-part60e-truncated-jensen-profile-completeness-blindness} 可知：任取 \(\rho_\star<1\)，截断剖面只能恢复
\[
\mu\big|_{(-\infty,\log\rho_\star)};
\]
只有让 \(\rho_\star\uparrow 1\)，即保留整个 \(0<\rho<1\) 上的 Jensen 势函数，才能恢复全部径向零测度 \(\mu\)。

算术侧则由
\[
G_S=\bigcup_{\substack{S_0\subset S\\ |S_0|<\infty}}G_{S_0}
=
\varinjlim_{\substack{S_0\subset S\\ |S_0|<\infty}}G_{S_0}
\]
与 Pontryagin 对偶的反变等价得到
\[
\Sigma_S
\cong
\varprojlim_{\substack{S_0\subset S\\ |S_0|<\infty}}\Sigma_{S_0}.
\]
因此，\(\xi\in\Sigma_S\) 与其全部有限投影族
\[
\{\pi_{S\to S_0}(\xi)\}_{S_0\subset S,\ |S_0|<\infty}
\]
完全等价；任何有限截断都会丢失推论 \ref{cor:xi-time-part60e-infinite-solenoid-finite-visible-factorization} 中的无限核
\(
\prod_{p\in S\setminus S_0}\ZZ_p
\)。

第 \(3\) 条只是前两条的合并表述：在解析侧，有限半径截断留下不可见外环；在算术侧，有限素数截断留下不可见无限 \(p\)-进核。故有限可见统计必然塌缩，而完备可逆账本在算术侧必然保持无限 \(p\)-进结构。证毕。
\end{proof}

\endinput

\paragraph{有限视界 classical RH 审计、谱替身与共尾 prime-ledger 的分离障碍}
在有限半径/有限素数证书的不完备性、无一致 realization 维数上界时的谱替身分岔、纯相位制度下的 \(\mathsf{NULL}\) 类型安全，以及扭子塔首次偏离深度的 Busy--Beaver 下界都已建立之后，还可把有限可见审计的障碍机制压缩为如下七条终端结论。它们分别刻画有限视界完备模数的不可计算增长、有限可见语言中的谱替身不可分辨、可靠性/完备性/可计算性的三难并立不可能、障碍来源的精确定位、纯相位制度下离切片信息的二歧、共尾 prime-ledger 的有限视界不可重构，以及相应的视界--账本分离原则。

\begin{theorem}[有限视界完备模数的 Busy--Beaver 障碍]\label{thm:xi-time-part60eb-finite-horizon-complete-modulus-busybeaver}
设 \(\mathcal V\) 为一种 classical RH 审计制度。对每个输入长度上界 \(\ell\)，\(\mathcal V\) 只允许访问受同一资源参数控制的
\[
\text{有限 prime-stage } \Sigma_S,\qquad
\text{有限 Toeplitz 主子式阶数},\qquad
\text{有限 dyadic 深度},\qquad
\text{有限 torsion 层级}.
\]
若存在函数 \(\Phi:\NN\to\NN\) 使得：对所有长度不超过 \(\ell\) 的输入实例，审计制度 \(\mathcal V\) 仅访问资源不超过 \(\Phi(\ell)\) 的上述四类有限视界，便总能给出正确判定，则称 \(\Phi\) 为 \(\mathcal V\) 的完备模数。

若 \(\mathcal V\) 在命题 \ref{prop:xi-time-part9g-deviation-depth-computability-separation} 的 torsion-tower 编码所给出的停机型偏离实例族上同时可靠且完备，则 \(\Phi\) 不可计算。更强地，存在常数 \(c>0\) 使得对充分大的 \(\ell\)，有
\[
\Phi(\ell)\ \ge\ L\,\mathrm{BB}(\ell-c)+1.
\]
\end{theorem}
\begin{proof}
沿用命题 \ref{prop:xi-time-part9g-first-deviation-depth} 与命题
\ref{prop:xi-time-part9g-deviation-depth-computability-separation}
的记号。对每个编码实例 \((M,x)\)，其有限层 torsion 投影
\[
Q_{M,x}^{(n)}
\]
都是全可计算的，并且
\[
\delta(M,x)<\infty
\quad\Longleftrightarrow\quad
M(x)\ \text{停机}.
\]

现取一个停机实例 \((M,x)\)，并记其首次偏离深度为 \(\delta(M,x)\)。由命题
\ref{prop:xi-time-part9g-first-deviation-depth}，
在全部层级 \(n<\delta(M,x)\) 上，编码对象与非停机基线对象的 torsion 投影完全一致。由于 prime-stage、有限 Toeplitz 主子式、有限 dyadic 深度与有限 torsion 层级都被同一资源上界控制，只要资源预算严格小于 \(\delta(M,x)\)，审计制度就仍停留在“首次偏离尚未出现”的有限视界之内，因而无法把该停机实例与对应的非停机基线区分开来。

若 \(\mathcal V\) 在这族实例上既可靠又完备，则对每个长度不超过 \(\ell\) 的停机实例都必须在资源 \(\Phi(\ell)\) 以内检出首次偏离；否则它就会把某个停机实例与非停机基线混淆。故必有
\[
\delta(M,x)\le \Phi(|(M,x)|)
\]
对全部停机实例成立。于是 \(\Phi\) 支配了一条统一验证深度函数。命题
\ref{prop:xi-time-part9g-deviation-depth-computability-separation}
第 \(3\) 条立即给出
\[
\Phi(\ell)\ \ge\ L\,\mathrm{BB}(\ell-c)+1
\]
对充分大的 \(\ell\) 成立，从而 \(\Phi\) 不可计算。证毕。
\end{proof}

\begin{theorem}[有限可见语言中的谱替身不可分辨]\label{thm:xi-time-part60eb-finite-visible-language-spectral-surrogates}
设 \(\mathcal L_{\mathrm{vis}}\) 为由下列可见对象生成的对象语言：
\[
\text{(i) disk-visible 相位输出},\qquad
\text{(ii) 有限多个有理频率角色},\qquad
\text{(iii) 有限阶 Toeplitz/Gram 主子式},\qquad
\text{(iv) 有限 dyadic 证书对象}.
\]
再设 admissible completion 族 \(\mathfrak X\) 在文稿既定的 Toeplitz--Carath\'eodory 审计语义下满足：其最小 realization 维数不存在一致上界。

则对任意有限子语言 \(\mathcal L_0\subset \mathcal L_{\mathrm{vis}}\)，存在两个 completion 对象
\[
X^+,X^-\in\mathfrak X
\]
满足
\[
X^+\equiv_{\mathcal L_0}X^-,
\]
其中 \(\equiv_{\mathcal L_0}\) 表示二者在 \(\mathcal L_0\) 的全部可见项与关系解释上完全一致；
但 \(X^+\) 满足全局 Toeplitz--Carath\'eodory 正性，而 \(X^-\) 不满足。等价地，在本文的 classical RH 语义中，真值不能由任意固定有限可见语言完全分离。
\end{theorem}
\begin{proof}
先把 \(\mathcal L_0\) 所访问的数据压缩为一个有限审计切片。由推论 \ref{cor:xi-time-part60e-infinite-solenoid-finite-visible-factorization}，任意有限族连续可见同态都必经由某个有限 prime-stage 因子化。另一方面，命题 \ref{prop:xi-time-part62d-universal-solenoid-phase-completed-scan} 第 \(3\) 条说明：任意有限有理频率扫描都只依赖某个有限层因子。于是，disk-visible 相位输出与有限多个有理频率角色可共同压缩到某个有限 prime-stage。

有限阶 Toeplitz/Gram 主子式与有限 dyadic 证书对象只涉及有限多个可见坐标以及有限次代数运算，因此在把上述有限 prime-stage 进一步放大后，也都经由同一有限视界因子化。故 \(\mathcal L_0\) 所能读取的全部信息可等价压缩为一个固定有限审计规则 \(\mathcal A_{\mathcal L_0}\)。

再由定理 \ref{thm:xi-spectral-surrogate-finite-audit-undecidable}，在不存在一致 realization 维数上界时，任意固定有限审计都存在两组谱替身对象，它们在该审计访问的一切有限数据上完全一致，而全局 Toeplitz--PSD 真值相反。又由定理 \ref{thm:xi-slice-horizon-carath-toeplitz}，全局 Toeplitz--Carath\'eodory 正性与相应的 RH 真值等价。于是存在 \(X^+,X^-\in\mathfrak X\) 使
\[
X^+\equiv_{\mathcal L_0}X^-,
\]
但其全局真值相反。证毕。
\end{proof}

\begin{corollary}[无一致 realization 维数界时有限视界 RH 审计的三难并立不可能]\label{cor:xi-time-part60eb-finite-horizon-rh-audit-trilemma}
在无一致 realization 维数上界的前提下，任何把 classical RH 限制在 \(\mathcal L_{\mathrm{vis}}\) 中进行的有限视界审计制度，都不可能同时满足：
\[
\text{可靠性},\qquad
\text{完备性},\qquad
\text{可计算性}.
\]
\end{corollary}
\begin{proof}
由定理 \ref{thm:xi-time-part60eb-finite-visible-language-spectral-surrogates}，任意固定有限可见语言都存在不可分辨而全局真值相反的谱替身，故若坚持停留在有限视界内，则完备性失效。

若试图通过引入一个统一的完备模数来恢复完备性，则在 torsion-tower 编码的停机型偏离实例族上，定理 \ref{thm:xi-time-part60eb-finite-horizon-complete-modulus-busybeaver} 强制该模数具有 Busy--Beaver 级不可计算增长。故可计算性失效。

因此，在无一致 realization 维数界时，上述三者不能并存。证毕。
\end{proof}

\begin{proposition}[决定性障碍并非“处于无限维背景”本身]\label{prop:xi-time-part60eb-obstruction-not-ambient-infinity}
在本文框架内，classical RH 的 no-go 机制并不等价于“对象处于无限维背景”这一事实本身。真正不可压缩的障碍是：
\[
\text{(a) realization 维数无一致上界},\qquad
\text{(b) 可见层对 prime-ledger 核失真},\qquad
\text{(c) 因而全阶 Toeplitz 量词无法统一消去为有限证书}.
\]
更精确地，若一族候选 completion 对象具有一致 realization 维数上界 \(D\)，则存在统一截断阶
\[
N_0\le 2D
\]
使全局 Toeplitz--Carath\'eodory 正性等价于有限个主子式条件。故“处于无限维背景中”本身并不推出制度性不可判定。
\end{proposition}
\begin{proof}
若最小 realization 维数有一致上界 \(D\)，则定理 \ref{thm:xi-toeplitz-psd-coherence-horizon-threshold} 与定理 \ref{thm:xi-oracle-collapse-toeplitz-psd-finite-truncation} 同时给出：存在统一阈值 \(N_0\le 2D\)，使
\[
\forall N\ge 0,\ \mathbf T_N\succeq 0
\quad\Longleftrightarrow\quad
\forall N\le N_0,\ \mathbf T_N\succeq 0.
\]
于是全局正性层级在该有界类上可被压缩为有限证书对象。

反之，若一致维数上界失效，则定理 \ref{thm:xi-spectral-surrogate-finite-audit-undecidable} 给出有限审计不可判定的谱替身；而定理 \ref{thm:xi-time-part67-continuous-phase-external-arithmetic-kernel} 与推论 \ref{cor:xi-time-part60e-infinite-solenoid-finite-visible-factorization} 表明，连通可见层并不忠实保留外置 prime-ledger 核。故真正导致 no-go 的，是“无一致维数界”与“核失真”这两条结构性机制的合成，而非背景空间的无限维性本身。证毕。
\end{proof}

\begin{theorem}[纯相位制度下离切片信息的二歧结构]\label{thm:xi-time-part60eb-pure-phase-offslice-bifurcation}
在纯相位可见协议中，任何关于离切片信息的断言都只能以如下两种方式之一出现：
\[
\text{(i) 显式外置模式轴 }\rho=\log|u|,\qquad
\text{(ii) 作为 }\mathsf{NULL}\text{ 或失败见证被外置。}
\]
不存在第三种纯相位内部表述，可在不增加模式轴的前提下稳定承载离切片信息。并且，若企图以统一可计算的有限模式轴预算来恢复全部离切片信息，则该预算必受定理 \ref{thm:xi-time-part60eb-finite-horizon-complete-modulus-busybeaver} 的 Busy--Beaver 障碍支配。
\end{theorem}
\begin{proof}
命题 \ref{prop:xi-offcritical-falsifiable-restatement} 已经给出严格的接口二分：离切片断言在体系内只能以“模式轴外延”或“\(\mathsf{NULL}\) 失败见证”两种方式出现。定理 \ref{thm:xi-offset-null-type-safety} 与命题 \ref{prop:xi-null-three-way} 进一步说明：若拒绝外置相应的正交记录轴，则该自由度在类型层不可被内部吸收，只能稳定退化为 \(\mathsf{NULL}\)。故不存在第三种纯相位内部表述。

若反设存在一条统一可计算的有限模式轴预算，能够对某族输入完全恢复离切片信息，则它与其余有限视界资源合并后便给出一条可计算的完备模数。把此模数施用于 torsion-tower 编码的停机型偏离实例族，即与定理 \ref{thm:xi-time-part60eb-finite-horizon-complete-modulus-busybeaver} 矛盾。故任何这类统一预算都必受 Busy--Beaver 级障碍支配。证毕。
\end{proof}

\begin{theorem}[共尾 prime-ledger 的有限视界不可重构性]\label{thm:xi-time-part60eb-cofinal-prime-ledger-finite-horizon-nonreconstruction}
设 \(\mathcal O\) 为任意有限族可见观测量，其成员皆由 disk-visible 输出、有限有理频率角色、有限阶 Toeplitz/Gram 主子式与有限 dyadic 证书对象编译而成。则存在一个有限素数集 \(S_0\) 使得
\[
\forall O\in\mathcal O,\qquad
O\ \text{都经由}\ \Sigma_{S_0}\ \text{因子化}.
\]

进一步，对任意严格共尾扩张链
\[
S_0\subsetneq S_1\subsetneq S_2\subsetneq \cdots,
\qquad
\bigcup_{n\ge 0}S_n=\mathbb P,
\]
相应 solenoid \(\Sigma_{S_n}\) 在全部 \(\mathcal O\)-观测上完全不可区分；然而它们的 prime-ledger 核
\[
K_{S_n}\cong \prod_{p\in S_n}\ZZ_p
\]
两两不可同构，并且每个 \(S_n\) 都可由 \(K_{S_n}\) 唯一恢复。
\end{theorem}
\begin{proof}
由推论 \ref{cor:xi-time-part60e-infinite-solenoid-finite-visible-factorization}，任意有限族连续可见同态都共同经由某个有限 prime-stage \(\Sigma_{S_0}\) 因子化。命题 \ref{prop:xi-time-part62d-universal-solenoid-phase-completed-scan} 第 \(3\) 条又说明：任意有限有理频率扫描本身就是某个有限层因子观测。于是 disk-visible 输出与有限有理频率角色都经由某个有限阶段因子化。

有限阶 Toeplitz/Gram 主子式以及有限 dyadic 证书对象只依赖于有限多个可见坐标和有限次代数运算；把所需有限阶段一并并入 \(S_0\) 后，它们也都经由 \(\Sigma_{S_0}\) 因子化。故对任意 \(O\in\mathcal O\)，都有
\[
O=O_0\circ \pi_{S_n\to S_0}
\qquad(n\ge 0)
\]
对某个 \(O_0\) 成立。这证明了 \(\Sigma_{S_n}\) 在全部 \(\mathcal O\)-观测上不可区分。

另一方面，若 \(S_n\neq S_m\)，则存在素数 \(p\) 使 \(p\in S_n\triangle S_m\)。对任意有限素数集 \(S\)，有
\[
K_S\cong \prod_{q\in S}\ZZ_q,
\]
故
\[
p\in S
\quad\Longleftrightarrow\quad
K_S/pK_S\neq 0
\quad\Longleftrightarrow\quad
\operatorname{Hom}_{\mathrm{cts}}(K_S,\ZZ/p\ZZ)\neq 0.
\]
因此 \(K_{S_n}\) 与 \(K_{S_m}\) 的 \(p\)-主因子结构不同，从而二者不同构；并且 \(S\) 可由判别全部素数 \(p\) 的上述等价条件唯一恢复。证毕。
\end{proof}

\begin{corollary}[视界--账本分离原则]\label{cor:xi-time-part60eb-horizon-ledger-separation-principle}
\leanverified{paper\_xi\_time\_part60eb\_horizon\_ledger\_separation\_principle}
若 classical RH 的真值最终依赖于共尾 prime-ledger 的整体几何，而审计协议仅驻留于有限视界可见层，则该真值的完全恢复必然失败。更准确地说，失败机制并非“样本不足”，而是可见函子对内部 prime-ledger 的非忠实压缩。
\end{corollary}
\begin{proof}
定理 \ref{thm:xi-time-part60eb-cofinal-prime-ledger-finite-horizon-nonreconstruction} 表明：任意有限视界可见层都只能看到某个有限 prime-stage，而不同的共尾 prime-ledger 可以在全部有限可见观测上保持不可区分。若 classical RH 的真值确实依赖于这种共尾账本几何，则任何忠实恢复过程都必须在有限视界内区分这些对象；这与定理 \ref{thm:xi-time-part60eb-cofinal-prime-ledger-finite-horizon-nonreconstruction} 相矛盾。

若进一步无一致 realization 维数上界，则定理 \ref{thm:xi-time-part60eb-finite-visible-language-spectral-surrogates} 把这种函子性失真强化为一对有限可见完全同像而全局 Toeplitz--Carath\'eodory 真值相反的谱替身对象。故所述恢复失败是结构性的，而非样本数量意义下的偶然不足。证毕。
\end{proof}

\noindent
综上，本文中的有限可见障碍机制可被严格压缩为一条共同结构：一方面，只要可靠完备性试图覆盖 torsion-tower 编码的停机型偏离实例，就必然遭遇 Busy--Beaver 级不可计算完备模数；另一方面，只要一致 realization 维数上界缺失，有限可见语言便不可避免地出现谱替身，而纯相位制度对离切片信息的处理又只能在“模式轴外延”与“\(\mathsf{NULL}\) 失败见证”之间二歧。于是真正不可压缩的障碍并不在于抽象的无限维背景本身，而在于无界 realization 维数、prime-ledger 核的可见失真以及有限视界预算的 Busy--Beaver 增长三者的共同作用。

\endinput

\paragraph{坏模数的倍数密度与 window-\texorpdfstring{$6$}{6} 长根扇区的隐藏退化}
Dirichlet 扭曲坏模数沿整除链的传播，与 window-\(6\) 的 \(B_3\) 长根扇区在纤维多重度层上的残余对称，可进一步压缩为下列一组结构性结论。

\begin{corollary}[整除传播单独强迫坏模数集具有正下密度]\label{cor:xi-dirichlet-bad-moduli-positive-lower-density-by-divisibility}
\leanverified{paper\_xi\_dirichlet\_bad\_moduli\_positive\_lower\_density\_by\_divisibility}
沿用命题 \ref{prop:real-input-40-output-dirichlet-nonrh} 与结论 \ref{con:xi-time-part9l-dirichlet-twist-divisibility-monotonicity} 的记号，定义
\[
\mathcal B:=\{m\ge 2:\ \theta_m>\tfrac12\}.
\]
则 \(\mathcal B\neq\varnothing\)，并且若记
\[
m_0:=\min \mathcal B,
\]
则
\[
\underline d(\mathcal B)
:=
\liminf_{X\to\infty}\frac{\#(\mathcal B\cap [1,X])}{X}
\ge
\frac{1}{m_0}
>
0.
\]
换言之，即使不调用大奇素数模数的渐近展开，仅由“坏模数沿倍数链传播”这一整除刚性，也已足以推出坏模数集的正下密度。
\end{corollary}
\begin{proof}
由命题 \ref{prop:real-input-40-output-dirichlet-nonrh}，存在某个模数 \(m\ge 2\) 满足 \(\theta_m>\tfrac12\)，故 \(\mathcal B\neq\varnothing\)，从而最小元 \(m_0\) 存在。
再由结论 \ref{con:xi-time-part9l-dirichlet-twist-divisibility-monotonicity}，若 \(m_0\in\mathcal B\)，则其任意倍数 \(rm_0\)（\(r\ge 1\)）仍属于 \(\mathcal B\)。因此
\[
m_0\NN_{\ge 1}\subseteq \mathcal B.
\]
于是对任意 \(X\ge 1\)，有
\[
\#(\mathcal B\cap [1,X])\ge \left\lfloor \frac{X}{m_0}\right\rfloor.
\]
两边除以 \(X\) 并取 \(\liminf\)，即得
\[
\underline d(\mathcal B)\ge \frac{1}{m_0}.
\]
证毕。
\end{proof}

\begin{theorem}[window-\texorpdfstring{$6$}{6} dyadic 微态质量的 \texorpdfstring{$16+36+12$}{16+36+12} 精确三块分解]\label{thm:xi-window6-dyadic-microstate-exact-three-block-mass-split}
\leanverified{paper\_xi\_window6\_dyadic\_microstate\_exact\_three\_block\_mass\_split}
沿用
\[
\mathrm{Fold}^{\mathrm{bin}}_6:\{0,\dots,63\}\to X_6
\]
与 \(V_6\) 的记号。则
\[
2^6
=
\sum_{\substack{w\in X_6\\ w_6=1}} d_6^{\mathrm{bin}}(w)
+
\sum_{\substack{w\in X_6\\ w_6=0,\ V_6(w)\le 8}} d_6^{\mathrm{bin}}(w)
+
\sum_{\substack{w\in X_6\\ w_6=0,\ V_6(w)>8}} d_6^{\mathrm{bin}}(w)
=
16+36+12.
\]
等价地，全部 \(64\) 个 dyadic 微态在 \(\mathrm{Fold}^{\mathrm{bin}}_6\) 下被精确分解为三块总质量分别为 \(16\)、\(36\)、\(12\) 的互斥簇。进一步，在均匀微态基线下有
\[
\mathbb P(w_6=1)=\frac14,\qquad
\mathbb P(w_6=0,\ V_6(w)\le 8)=\frac{9}{16},\qquad
\mathbb P(w_6=0,\ V_6(w)>8)=\frac{3}{16}.
\]
\end{theorem}
\begin{proof}
由推论 \ref{cor:fold-bin-m6-fiber-hist}，
\[
\#\{w\in X_6:\ d_6^{\mathrm{bin}}(w)=2\}=8,\qquad
\#\{w\in X_6:\ d_6^{\mathrm{bin}}(w)=3\}=4,\qquad
\#\{w\in X_6:\ d_6^{\mathrm{bin}}(w)=4\}=9.
\]
另一方面，表 \ref{tab:fold6_bin_uplift_choice_collapse} 给出精确对应
\[
w_6=1 \Longleftrightarrow d_6^{\mathrm{bin}}(w)=2,
\]
\[
w_6=0,\ V_6(w)\le 8 \Longleftrightarrow d_6^{\mathrm{bin}}(w)=4,
\]
\[
w_6=0,\ V_6(w)>8 \Longleftrightarrow d_6^{\mathrm{bin}}(w)=3.
\]
因此三类纤维的总微态质量分别为
\[
8\cdot 2=16,\qquad 9\cdot 4=36,\qquad 4\cdot 3=12.
\]
又因为不同 \(w\) 的纤维两两不交且其并恰为 \(\{0,\dots,63\}\)，故
\[
64=16+36+12.
\]
最后两边除以 \(64\) 即得所示概率。证毕。
\end{proof}

\begin{theorem}[window-\texorpdfstring{$6$}{6} 的 \texorpdfstring{$A_2(+)$}{A2(+)} 与 \texorpdfstring{$A_2(-)$}{A2(-)} 长根扇区具有完全相同的纤维多重度谱]\label{thm:xi-window6-a2pm-identical-fiber-spectrum}
\leanverified{paper\_xi\_window6\_a2pm\_identical\_fiber\_spectrum}
沿用定理 \ref{thm:boundary-m10-b3c3-rigid-four-block-split} 的 \(B_3\) 根字典与分块记号。定义
\[
\Phi_{A_2(+)}
:=
\{\pm(e_1+e_2),\ \pm(e_1+e_3),\ \pm(e_2+e_3)\},
\]
\[
\Phi_{A_2(-)}
:=
\{\pm(e_1-e_2),\ \pm(e_1-e_3),\ \pm(e_2-e_3)\},
\]
\[
X_6^{A_2(+)}:=\iota_B^{-1}\bigl(\Phi_{A_2(+)}\bigr),
\qquad
X_6^{A_2(-)}:=\iota_B^{-1}\bigl(\Phi_{A_2(-)}\bigr).
\]
则两块的纤维多重度多重集完全一致，且都等于
\[
\{4,4,3,3,2,2\}.
\]
等价地，若定义生成多项式
\[
P_{A_2(\pm)}(t):=\sum_{w\in X_6^{A_2(\pm)}} t^{d_6^{\mathrm{bin}}(w)},
\]
则
\[
P_{A_2(+)}(t)=P_{A_2(-)}(t)=2t^4+2t^3+2t^2.
\]
\end{theorem}
\begin{proof}
由表 \ref{tab:boundary_uplift_m10_dictionary} 可直接读出
\[
X_6^{A_2(+)}
=
\{\texttt{000000},\texttt{010100},\texttt{100010},
\texttt{101010},\texttt{010001},\texttt{010101}\},
\]
\[
X_6^{A_2(-)}
=
\{\texttt{101000},\texttt{100100},\texttt{010010},
\texttt{001010},\texttt{001001},\texttt{000101}\}.
\]
另一方面，推论 \ref{cor:fold-bin-m6-fiber-hist} 与表 \ref{tab:fold6_bin_uplift_choice_collapse} 给出三档精确分类：
\[
w_6=0,\ V_6(w)\le 8 \Longleftrightarrow d_6^{\mathrm{bin}}(w)=4,
\]
\[
w_6=0,\ V_6(w)>8 \Longleftrightarrow d_6^{\mathrm{bin}}(w)=3,
\]
\[
w_6=1 \Longleftrightarrow d_6^{\mathrm{bin}}(w)=2.
\]
其中在 \(X_6^{A_2(+)}\) 中，
\[
\texttt{000000},\texttt{010100}\in\{w_6=0,\ V_6(w)\le 8\},
\]
\[
\texttt{100010},\texttt{101010}\in\{w_6=0,\ V_6(w)>8\},
\]
\[
\texttt{010001},\texttt{010101}\in\{w_6=1\}.
\]
故其多重度多重集为
\[
\{4,4,3,3,2,2\}.
\]
同理，在 \(X_6^{A_2(-)}\) 中，
\[
\texttt{101000},\texttt{100100}\in\{w_6=0,\ V_6(w)\le 8\},
\]
\[
\texttt{010010},\texttt{001010}\in\{w_6=0,\ V_6(w)>8\},
\]
\[
\texttt{001001},\texttt{000101}\in\{w_6=1\},
\]
故其多重度多重集同样为
\[
\{4,4,3,3,2,2\}.
\]
于是两块的生成多项式均为
\[
2t^4+2t^3+2t^2.
\]
证毕。
\end{proof}

\begin{corollary}[任何仅依赖纤维多重度的统计量都无法区分 \texorpdfstring{$A_2(+)$}{A2(+)} 与 \texorpdfstring{$A_2(-)$}{A2(-)}]\label{cor:xi-window6-a2pm-multiplicity-blind-statistics}
\leanverified{paper\_xi\_window6\_a2pm\_multiplicity\_blind\_statistics}
沿用定理 \ref{thm:xi-window6-a2pm-identical-fiber-spectrum} 的记号。对任意函数
\[
\Phi:\{2,3,4\}\to \RR,
\]
都有
\[
\sum_{w\in X_6^{A_2(+)}}\Phi\!\bigl(d_6^{\mathrm{bin}}(w)\bigr)
=
\sum_{w\in X_6^{A_2(-)}}\Phi\!\bigl(d_6^{\mathrm{bin}}(w)\bigr).
\]
特别地，对任意 \(q\in\RR\)，有
\[
\sum_{w\in X_6^{A_2(+)}}\bigl(d_6^{\mathrm{bin}}(w)\bigr)^q
=
\sum_{w\in X_6^{A_2(-)}}\bigl(d_6^{\mathrm{bin}}(w)\bigr)^q
=
2(2^q+3^q+4^q).
\]
\end{corollary}
\begin{proof}
由定理 \ref{thm:xi-window6-a2pm-identical-fiber-spectrum}，两块的纤维多重度多重集都等于
\[
\{4,4,3,3,2,2\}.
\]
因此对任意只依赖多重度数值的标量函数 \(\Phi\)，两侧求和都仅由同一个多重集决定，故必相等。取 \(\Phi(d)=d^q\) 即得第二式。证毕。
\end{proof}

\begin{theorem}[纤维大小场在 \texorpdfstring{$W(B_3)$}{W(B3)}-不变子空间中的最优两轨道压缩]\label{thm:xi-window6-weyl-invariant-best-two-orbit-compression}
沿用推论 \ref{cor:fold6-weyl-two-orbit-compression} 的分块记号，记
\[
d_6^{\mathrm{cyc}}:=d_6^{\mathrm{bin}}\big|_{X_6^{\mathrm{cyc}}}:X_6^{\mathrm{cyc}}\to\{2,3,4\}.
\]
在实向量空间 \(\RR^{X_6^{\mathrm{cyc}}}\) 上取均匀内积
\[
\langle f,g\rangle
:=
\frac{1}{18}\sum_{w\in X_6^{\mathrm{cyc}}}f(w)g(w),
\]
并定义 \(W(B_3)\)-不变子空间
\[
\mathcal W
:=
\{f:X_6^{\mathrm{cyc}}\to\RR:\ f\ \text{在}\ X_6^{\mathrm{short}}\ \text{与}\ X_6^{\mathrm{long}}\ \text{上分别常值}\}.
\]
则成立：
\begin{enumerate}
\item \(d_6^{\mathrm{cyc}}\) 在 \(\mathcal W\) 上的唯一正交投影 \(\Pi_{\mathcal W}d_6^{\mathrm{cyc}}\) 为
\[
(\Pi_{\mathcal W}d_6^{\mathrm{cyc}})(w)
=
\begin{cases}
\frac{11}{3},& w\in X_6^{\mathrm{short}},\\[4pt]
3,& w\in X_6^{\mathrm{long}}.
\end{cases}
\]
\item 最小均方误差恰为
\[
\inf_{f\in\mathcal W}
\frac{1}{18}\sum_{w\in X_6^{\mathrm{cyc}}}\bigl(d_6^{\mathrm{cyc}}(w)-f(w)\bigr)^2
=
\frac{17}{27}.
\]
\item 一致范数下的最佳 \(W(B_3)\)-不变逼近误差恰为
\[
\inf_{f\in\mathcal W}\|d_6^{\mathrm{cyc}}-f\|_{\infty}=1,
\]
且唯一极小元为常函数
\[
f\equiv 3.
\]
\end{enumerate}
\leanverified{paper\_xi\_window6\_weyl\_invariant\_best\_two\_orbit\_compression}
\end{theorem}
\begin{proof}
由表 \ref{tab:fold6_b3c3_root_dictionary_b3} 与表 \ref{tab:fold_bin_gauge_m6_fibers}，六个短根原像的多重度多重集为
\[
\{4,4,4,4,4,2\}.
\]
另一方面，由定理 \ref{thm:xi-window6-a2pm-identical-fiber-spectrum}，
\[
X_6^{\mathrm{long}}=X_6^{A_2(+)}\sqcup X_6^{A_2(-)}
\]
且两块各自的多重度多重集都为 \(\{4,4,3,3,2,2\}\)。故十二个长根原像的多重度多重集为
\[
\{4,4,4,4,3,3,3,3,2,2,2,2\}.
\]

由推论 \ref{cor:fold6-weyl-two-orbit-compression}，\(\mathcal W\) 正是“短根块常值、长根块常值”的二维子空间。因此其正交投影由两条轨道上的平均值给出。短根平均值为
\[
\frac{5\cdot 4+2}{6}=\frac{11}{3},
\]
长根平均值为
\[
\frac{4\cdot 4+4\cdot 3+4\cdot 2}{12}=3.
\]
这就得到第一条。

对第二条，短根块的最小平方误差贡献为
\[
\frac{1}{18}\left(
5\left(4-\frac{11}{3}\right)^2+\left(2-\frac{11}{3}\right)^2
\right)
=
\frac{1}{18}\left(5\cdot \frac{1}{9}+\frac{25}{9}\right)
=
\frac{5}{27},
\]
长根块的最小平方误差贡献为
\[
\frac{1}{18}\left(
4(4-3)^2+4(3-3)^2+4(2-3)^2
\right)
=
\frac{8}{18}
=
\frac{4}{9}.
\]
两者相加即得
\[
\frac{5}{27}+\frac{4}{9}=\frac{17}{27}.
\]

最后考察一致范数。任取 \(f\in\mathcal W\)，则存在 \(a,b\in\RR\) 使
\[
f|_{X_6^{\mathrm{short}}}\equiv a,
\qquad
f|_{X_6^{\mathrm{long}}}\equiv b.
\]
于是
\[
\sup_{w\in X_6^{\mathrm{short}}}|d_6^{\mathrm{cyc}}(w)-f(w)|
=
\max\{|2-a|,|4-a|\},
\]
其最小值为 \(1\)，且唯一在 \(a=3\) 时取得；
\[
\sup_{w\in X_6^{\mathrm{long}}}|d_6^{\mathrm{cyc}}(w)-f(w)|
=
\max\{|2-b|,|3-b|,|4-b|\},
\]
其最小值同样为 \(1\)，且唯一在 \(b=3\) 时取得。故
\[
\|d_6^{\mathrm{cyc}}-f\|_\infty
=
\max\left\{
\max\{|2-a|,|4-a|\},
\max\{|2-b|,|3-b|,|4-b|\}
\right\}
\ge 1,
\]
并且等号当且仅当 \(a=b=3\) 时成立。于是最优一致逼近误差为 \(1\)，唯一极小元为常函数 \(f\equiv 3\)。证毕。
\end{proof}

\noindent
这些结论说明：Dirichlet 扭曲侧的平方根门槛失效一旦出现，便会沿倍数链产生具有正下密度的坏模数族；而在 window-\(6\) 的 \(B_3\) 离散根字典中，虽然纤维大小场严格打破了 Weyl 轨道均匀性，长根层内部却仍保留一条精确的 \(A_2(+)/A_2(-)\) 多重度退化。于是“Weyl 破缺”与“残余二块等谱”并非矛盾，而是在不同粗细尺度上同时成立的两级结构。

\endinput

\begin{theorem}[window-\texorpdfstring{$6$}{6} 的 \texorpdfstring{$B_3$}{B3} 根云上同类等值二次型的完全分类]\label{thm:xi-window6-b3-rootcloud-quadratic-isotropy-forcing}
\leanverified{paper\_xi\_window6\_b3\_rootcloud\_quadratic\_isotropy\_forcing}
沿用定理 \ref{thm:boundary-m10-b3c3-rigid-four-block-split} 的 \(B_3\) 根字典与定理 \ref{thm:xi-window6-a2pm-identical-fiber-spectrum} 的记号。记
\[
\Phi_{\mathrm{short}}
:=
\{\pm e_1,\pm e_2,\pm e_3\},
\]
\[
\Phi_{A_2(+)}
:=
\{\pm(e_1+e_2),\pm(e_1+e_3),\pm(e_2+e_3)\},
\]
\[
\Phi_{A_2(-)}
:=
\{\pm(e_1-e_2),\pm(e_1-e_3),\pm(e_2-e_3)\}.
\]
设
\[
Q(x)=x^{\mathsf T}Ax
\qquad
\bigl(A=(a_{ij})\in \mathrm{Sym}_3(\RR)\bigr)
\]
为 \(\RR^3\) 上实对称二次型。若存在实数 \(q_s,q_+,q_-\) 使得
\[
Q(\alpha)=q_s\quad (\alpha\in \Phi_{\mathrm{short}}),
\]
\[
Q(\beta)=q_+\quad (\beta\in \Phi_{A_2(+)}),
\]
\[
Q(\gamma)=q_-\quad (\gamma\in \Phi_{A_2(-)}),
\]
则 \(A\) 必且仅必具有形式
\[
\boxed{\;
A=
\begin{pmatrix}
a&b&b\\
b&a&b\\
b&b&a
\end{pmatrix},
\qquad
q_s=a,\quad
q_+=2a+2b,\quad
q_-=2a-2b.\;}
\]
特别地，若进一步满足
\[
q_+=q_-,
\]
则必有
\[
\boxed{\;
A=aI_3.\;}
\]
换言之，一旦 \(A_2(+)\) 与 \(A_2(-)\) 两个长根块在二次能量上不可区分，整个 \(B_3\) 根云的二次几何便被强制为欧氏各向同性。
\end{theorem}
\begin{proof}
由
\[
\Phi_{\mathrm{short}}=\{\pm e_1,\pm e_2,\pm e_3\},
\]
条件 \(Q(\pm e_i)=q_s\) 对 \(i=1,2,3\) 立刻给出
\[
a_{11}=a_{22}=a_{33}=:a=q_s.
\]

再看 \(A_2(+)\) 块。因为
\[
Q(e_1+e_2)=a_{11}+a_{22}+2a_{12}=2a+2a_{12},
\]
\[
Q(e_1+e_3)=2a+2a_{13},
\qquad
Q(e_2+e_3)=2a+2a_{23},
\]
而三者都等于同一个常数 \(q_+\)，故
\[
a_{12}=a_{13}=a_{23}=:b.
\]
于是
\[
A=
\begin{pmatrix}
a&b&b\\
b&a&b\\
b&b&a
\end{pmatrix}.
\]

对这种矩阵直接计算可得
\[
Q(e_i)=a\qquad (1\le i\le 3),
\]
\[
Q(e_i+e_j)=2a+2b\qquad (1\le i<j\le 3),
\]
\[
Q(e_i-e_j)=2a-2b\qquad (1\le i<j\le 3).
\]
因此
\[
q_s=a,\qquad q_+=2a+2b,\qquad q_-=2a-2b,
\]
从而前述形式既是必要条件，也是充分条件。

最后，若 \(q_+=q_-\)，则
\[
2a+2b=2a-2b,
\]
故 \(b=0\)，于是
\[
A=aI_3.
\]
这正是所述各向同性结论。证毕。
\end{proof}

\endinput

\paragraph{黄金地址柱密度、对数级模糊质量与有限证书的双重强失明}
在 Gödel--Zeckendorf 地址极限的严格前缀等距、素数码集合 \(\mathrm{ZG}\) 的自然密度闭式，以及有限半径/有限素数证书的不完备性已经分别建立之后，还可把两条链进一步压缩为一组更锋利的结论：任意有限地址柱集都承载严格正的算术密度，而零前缀地址球中的质量衰减仅为对数级，从而把有限半径解析证书与有限 prime shadow 的联合失明定量化为一条显式下界。

\begin{theorem}[任意 admissible 前缀柱集的自然密度矩阵公式与严格正性]\label{thm:xi-time-part60f-zg-cylinder-density-matrix-positive}
\leanverified{paper\_xi\_time\_part60f\_zg\_cylinder\_density\_matrix\_positive}
设
\[
\mathrm{ZG}
:=
\Bigl\{
\prod_{k\ge 1}p_k^{z_k}:\ 
z_k\in\{0,1\},\ z_kz_{k+1}=0,\ z_k=0\ \text{对充分大 }k
\Bigr\}.
\]
对任意长度 \(n\) 的二元词
\[
\mathbf z=(z_1,\dots,z_n)\in\{0,1\}^n
\]
定义算术柱集
\[
\mathrm{ZG}[\mathbf z]
:=
\Bigl\{
m\in \mathrm{ZG}:\ v_{p_1}(m)=z_1,\dots,v_{p_n}(m)=z_n
\Bigr\}.
\]
再定义一步矩阵
\[
M_k^{(0)}
:=
\begin{pmatrix}
\dfrac{p_k}{p_k+1} & \dfrac{p_k}{p_k+1}\\[2mm]
0&0
\end{pmatrix},
\qquad
M_k^{(1)}
:=
\begin{pmatrix}
0&0\\[2mm]
\dfrac{1}{p_k+1}&0
\end{pmatrix},
\]
\[
M_k:=M_k^{(0)}+M_k^{(1)}
=
\begin{pmatrix}
\dfrac{p_k}{p_k+1} & \dfrac{p_k}{p_k+1}\\[2mm]
\dfrac{1}{p_k+1} & 0
\end{pmatrix}.
\]
记
\[
\mathbf e_0:=\begin{pmatrix}1\\0\end{pmatrix},
\qquad
\mathbf 1:=\begin{pmatrix}1\\1\end{pmatrix}.
\]
则成立：
\begin{enumerate}
  \item 若 \(\mathbf z\) 不 admissible，即存在 \(i<n\) 使 \(z_iz_{i+1}=1\)，则
  \[
  \mathrm{dens}\bigl(\mathrm{ZG}[\mathbf z]\bigr)=0.
  \]
  \item 若 \(\mathbf z\) admissible，则自然密度存在，且满足精确矩阵公式
  \[
  \mathrm{dens}\bigl(\mathrm{ZG}[\mathbf z]\bigr)
  =
  \frac{1}{\zeta(2)}
  \lim_{N\to\infty}
  \mathbf 1^\top
  M_NM_{N-1}\cdots M_{n+1}
  M_n^{(z_n)}\cdots M_1^{(z_1)}\mathbf e_0.
  \]
  \item 对每个 admissible \(\mathbf z\)，上述极限严格为正。
\end{enumerate}
\end{theorem}
\begin{proof}
对 \(N\ge n\)，定义截断柱集
\[
\mathrm{ZG}^{(N)}[\mathbf z]
:=
\Bigl\{
m\in\NN:\ 
v_{p_i}(m)=z_i\ (1\le i\le n),\
v_{p_k}(m)\in\{0,1\}\ (n<k\le N),
\]
\[
v_{p_k}(m)v_{p_{k+1}}(m)=0\ (1\le k<N)
\Bigr\}.
\]
其成员资格仅依赖于模
\[
\prod_{k=1}^{N}p_k^2
\]
的剩余类，故自然密度存在。

若 \(\mathbf z\) 不 admissible，则任一长度 \(N\) 的扩展都不 admissible，从而
\[
\mathrm{ZG}^{(N)}[\mathbf z]=\varnothing
\qquad(N\ge n),
\]
故 \(\mathrm{dens}(\mathrm{ZG}[\mathbf z])=0\)。
这也可直接从
\[
M_{i+1}^{(1)}M_i^{(1)}=0
\]
看出。

以下设 \(\mathbf z\) admissible。
对任意 admissible 扩展
\[
\mathbf x=(x_1,\dots,x_N)\in\{0,1\}^N
\]
且 \(x_i=z_i\)（\(1\le i\le n\)），恰有
\[
\mathrm{dens}\bigl\{m:\ v_{p_k}(m)=x_k\ (1\le k\le N)\bigr\}
=
\prod_{x_k=0}\left(1-\frac1{p_k}\right)
\prod_{x_k=1}\left(\frac1{p_k}-\frac1{p_k^2}\right).
\]
将两类因子重写为
\[
1-\frac1{p_k}
=
\left(1-\frac1{p_k^2}\right)\frac{p_k}{p_k+1},
\qquad
\frac1{p_k}-\frac1{p_k^2}
=
\left(1-\frac1{p_k^2}\right)\frac1{p_k+1},
\]
便得到
\[
\mathrm{dens}\bigl\{m:\ v_{p_k}(m)=x_k\ (1\le k\le N)\bigr\}
=
\left(\prod_{k=1}^{N}\left(1-\frac1{p_k^2}\right)\right)
\prod_{x_k=0}\frac{p_k}{p_k+1}
\prod_{x_k=1}\frac1{p_k+1}.
\]
对所有 admissible 扩展求和，即得
\[
\mathrm{dens}\bigl(\mathrm{ZG}^{(N)}[\mathbf z]\bigr)
=
\left(\prod_{k=1}^{N}\left(1-\frac1{p_k^2}\right)\right)
U_N(\mathbf z),
\]
其中
\[
U_N(\mathbf z)
:=
\mathbf 1^\top
M_N\cdots M_{n+1}
M_n^{(z_n)}\cdots M_1^{(z_1)}\mathbf e_0.
\]
这正是所述矩阵量。

另一方面，
\[
\mathrm{ZG}[\mathbf z]=\bigcap_{N\ge n}\mathrm{ZG}^{(N)}[\mathbf z].
\]
若 \(m\in \mathrm{ZG}^{(N)}[\mathbf z]\setminus \mathrm{ZG}[\mathbf z]\)，则尾部必发生违约事件之一：
\[
\exists\,k>N,\ p_k^2\mid m
\qquad\text{或}\qquad
\exists\,k\ge N,\ p_kp_{k+1}\mid m.
\]
故
\[
\mathrm{ZG}^{(N)}[\mathbf z]\setminus \mathrm{ZG}[\mathbf z]
\subseteq
\Bigl(\bigcup_{k>N}\{p_k^2\mid m\}\Bigr)
\cup
\Bigl(\bigcup_{k\ge N}\{p_kp_{k+1}\mid m\}\Bigr).
\]
由并集上界，
\[
0\le
\mathrm{dens}\bigl(\mathrm{ZG}^{(N)}[\mathbf z]\bigr)
-\mathrm{dens}\bigl(\mathrm{ZG}[\mathbf z]\bigr)
\le
\sum_{k>N}\frac1{p_k^2}
+
\sum_{k\ge N}\frac1{p_kp_{k+1}}.
\]
令 \(N\to\infty\)，右端趋于 \(0\)，故自然密度存在并满足
\[
\mathrm{dens}\bigl(\mathrm{ZG}[\mathbf z]\bigr)
=
\lim_{N\to\infty}
\left(\prod_{k=1}^{N}\left(1-\frac1{p_k^2}\right)\right)U_N(\mathbf z).
\]
再由
\[
\prod_{k=1}^{N}\left(1-\frac1{p_k^2}\right)\longrightarrow \frac1{\zeta(2)},
\]
便得到第二条。

最后证明严格正性。
由 admissibility，矩阵积
\[
M_n^{(z_n)}\cdots M_1^{(z_1)}\mathbf e_0
\]
必为
\[
c_{\mathbf z}\,\mathbf e_{z_n},
\qquad
c_{\mathbf z}:=
\prod_{i:\ z_i=0}\frac{p_i}{p_i+1}
\prod_{i:\ z_i=1}\frac1{p_i+1}
>0,
\]
其中 \(\mathbf e_1:=(0,1)^\top\)。

若 \(z_n=0\)，则
\[
U_N(\mathbf z)=
c_{\mathbf z}\,\mathbf 1^\top M_N\cdots M_{n+1}\mathbf e_0.
\]
此时右端正是将定理 \ref{con:xi-zg-density-closed-form-tail-bound} 的证明逐字应用到移位素数列
\[
(p_{n+1},p_{n+2},\dots)
\]
所得的尾部归一化权重，其极限严格为正。

若 \(z_n=1\)，则
\[
M_{n+1}\mathbf e_1
=
\begin{pmatrix}
\dfrac{p_{n+1}}{p_{n+1}+1}\\[2mm]
0
\end{pmatrix}
=
\frac{p_{n+1}}{p_{n+1}+1}\mathbf e_0,
\]
故
\[
U_N(\mathbf z)=
c_{\mathbf z}\frac{p_{n+1}}{p_{n+1}+1}\,
\mathbf 1^\top M_N\cdots M_{n+2}\mathbf e_0,
\]
同样化归到移位尾部的正极限情形。
因此对每个 admissible \(\mathbf z\)，极限严格为正。证毕。
\end{proof}

\begin{theorem}[零前缀柱集的对数级自然密度律]\label{thm:xi-time-part60f-zg-zero-prefix-log-density}
记
\[
\mathbf 0^{(n)}:=(0,\dots,0)\in\{0,1\}^n,
\qquad
\mathrm{ZG}[\mathbf 0^{(n)}]
=
\{m\in\mathrm{ZG}:\ p_1,\dots,p_n\nmid m\}.
\]
则当 \(n\to\infty\) 时有精确渐近
\[
\mathrm{dens}\bigl(\mathrm{ZG}[\mathbf 0^{(n)}]\bigr)
\sim
\frac{e^{-\gamma}}{\log p_n},
\]
其中 \(\gamma\) 为 Euler--Mascheroni 常数。
\end{theorem}
\begin{proof}
定义尾部约束集合
\[
\mathrm{Tail}_n
:=
\Bigl\{
m\in\NN:\ 
\forall k>n,\ v_{p_k}(m)\in\{0,1\},\
v_{p_k}(m)v_{p_{k+1}}(m)=0
\Bigr\}.
\]
则
\[
\mathrm{ZG}[\mathbf 0^{(n)}]
=
\Bigl(\bigcap_{i=1}^{n}\{p_i\nmid m\}\Bigr)\cap \mathrm{Tail}_n.
\]
对任意固定 \(N>n\)，有限截断版本的密度按前 \(n\) 个素数与尾部条件相乘，令 \(N\to\infty\) 即得
\[
\mathrm{dens}\bigl(\mathrm{ZG}[\mathbf 0^{(n)}]\bigr)
=
\left(\prod_{i=1}^{n}\left(1-\frac1{p_i}\right)\right)\tau_n,
\]
其中
\[
\tau_n:=\mathrm{dens}(\mathrm{Tail}_n).
\]

现证 \(\tau_n\to 1\)。
若 \(m\notin \mathrm{Tail}_n\)，则尾部必发生违约事件之一：
\[
\exists\,k>n,\ p_k^2\mid m
\qquad\text{或}\qquad
\exists\,k\ge n+1,\ p_kp_{k+1}\mid m.
\]
故
\[
1-\tau_n
\le
\sum_{k>n}\frac1{p_k^2}
+
\sum_{k\ge n+1}\frac1{p_kp_{k+1}}
=o(1).
\]
因此
\[
\tau_n=1+o(1).
\]
于是
\[
\mathrm{dens}\bigl(\mathrm{ZG}[\mathbf 0^{(n)}]\bigr)
=
\left(\prod_{i=1}^{n}\left(1-\frac1{p_i}\right)\right)(1+o(1)).
\]
最后由 Mertens 乘积公式
\[
\prod_{p\le x}\left(1-\frac1p\right)\sim \frac{e^{-\gamma}}{\log x}
\]
并取 \(x=p_n\)，得到
\[
\mathrm{dens}\bigl(\mathrm{ZG}[\mathbf 0^{(n)}]\bigr)
\sim
\frac{e^{-\gamma}}{\log p_n}.
\]
证毕。
\end{proof}
\leanverified{paper\_xi\_time\_part60f\_zg\_zero\_prefix\_log\_density}

\begin{theorem}[自然密度诱导的黄金柱测度、全支撑与有限支撑点的零局部维数]\label{thm:xi-time-part60f-zg-cylinder-measure-local-dimension-zero}
记
\[
\Sigma_{\mathrm{ZG}}
:=
\{(z_k)_{k\ge 1}\in\{0,1\}^{\NN}:\ z_kz_{k+1}=0\ \forall k\},
\]
并对任意 admissible 前缀
\[
\mathbf z=(z_1,\dots,z_n)
\]
定义 cylinder
\[
[\mathbf z]
:=
\{x\in\Sigma_{\mathrm{ZG}}:\ x_1=z_1,\dots,x_n=z_n\}.
\]
再定义柱权
\[
\nu_{\mathrm{ZG}}([\mathbf z])
:=
\frac{\mathrm{dens}\bigl(\mathrm{ZG}[\mathbf z]\bigr)}{\delta_{\mathrm{ZG}}},
\qquad
\delta_{\mathrm{ZG}}:=\mathrm{dens}(\mathrm{ZG}).
\]
则成立：
\begin{enumerate}
  \item \(\nu_{\mathrm{ZG}}\) 唯一延拓为 \(\Sigma_{\mathrm{ZG}}\) 上的 Borel 概率测度；
  \item \(\nu_{\mathrm{ZG}}\) 具有全支撑：每个非空 cylinder 都有正质量；
  \item 对任意 \(B=2^L\) 与 \(v\in T_2(E_{\min})\setminus 2T_2(E_{\min})\)，地址映射
  \[
  \Xi_B:\Sigma_{\mathrm{ZG}}\to K_{B,v},
  \qquad
  \Xi_B((z_t)_{t\ge 1})=\sum_{t\ge 1}z_tB^{t-1}v
  \]
  的推前
  \[
  \mu_{B}:=(\Xi_B)_\#\nu_{\mathrm{ZG}}
  \]
  是 \(K_{B,v}\) 上的全支撑 Borel 概率测度；
  \item 若 \(e=(e_k)_{k\ge 1}\in\Sigma_{\mathrm{ZG}}\) 有限支撑，即仅有限多个 \(e_k\) 非零，则存在常数 \(C(e)\in(0,\infty)\) 使得
  \[
  \nu_{\mathrm{ZG}}([e_1,\dots,e_n])
  \sim
  \frac{C(e)}{\log p_n}
  \qquad(n\to\infty),
  \]
  并且在点 \(x_e:=\Xi_B(e)\) 处，推前测度的局部维数沿自然半径 \(B^{-n}\) 满足
  \[
  \lim_{n\to\infty}
  \frac{\log \mu_B\bigl(\mathrm B_{K_{B,v}}(x_e,B^{-n})\bigr)}
       {\log B^{-n}}
  =0.
  \]
\end{enumerate}
\end{theorem}
\begin{proof}
先证第 \(1\) 条。
对任意 admissible 前缀 \(\mathbf z\)，有不交并分解
\[
[\mathbf z]
=
\bigsqcup_{\substack{a\in\{0,1\}\\ \mathbf z a\ \mathrm{admissible}}}
[\mathbf z a].
\]
算术侧同样满足
\[
\mathrm{ZG}[\mathbf z]
=
\bigsqcup_{\substack{a\in\{0,1\}\\ \mathbf z a\ \mathrm{admissible}}}
\mathrm{ZG}[\mathbf z a],
\]
而各项自然密度由定理
\ref{thm:xi-time-part60f-zg-cylinder-density-matrix-positive}
存在。
故柱权满足一致的有限可加兼容关系；又
\[
\nu_{\mathrm{ZG}}(\Sigma_{\mathrm{ZG}})
=
\frac{\mathrm{dens}(\mathrm{ZG})}{\delta_{\mathrm{ZG}}}
=1.
\]
由 cylinder 代数上的 Carath\'eodory 扩张定理，\(\nu_{\mathrm{ZG}}\) 唯一延拓为 \(\Sigma_{\mathrm{ZG}}\) 上的 Borel 概率测度。

第 \(2\) 条由定理 \ref{thm:xi-time-part60f-zg-cylinder-density-matrix-positive} 的严格正性立即得到：每个非空 cylinder 对应 admissible 前缀，因而其质量严格为正。

对第 \(3\) 条，定理 \ref{thm:xi-zg-address-binary-self-affine} 已证明 \(\Xi_B\) 在 \(\Sigma_{\mathrm{ZG}}\) 上单射，其像为 \(K_{B,v}\)；连续性显然，而 \(\Sigma_{\mathrm{ZG}}\) 紧、\(K_{B,v}\subset \ZZ_2v\) 为 Hausdorff 空间，故 \(\Xi_B\) 是到其像的同胚。于是推前 \(\mu_B\) 是 \(K_{B,v}\) 上的 Borel 概率测度。又 \(\Xi_B\) 将非空 cylinder 送到 \(K_{B,v}\) 的非空相对开集，结合第 \(2\) 条即知 \(\mu_B\) 具有全支撑。

现证第 \(4\) 条。
设 \(e\) 的支撑包含于 \(\{1,\dots,m_0\}\)。
对 \(n\ge m_0\)，前缀
\[
(e_1,\dots,e_n)
\]
恰等于固定有限非零前缀后接 \(n-m_0\) 个零。故
\[
\mathrm{ZG}[e_1,\dots,e_n]
\]
可写为满足前 \(m_0\) 位固定、其后到第 \(n\) 位全为零、而尾部 \(>n\) 继续满足 \(\mathrm{ZG}\) 约束的整数集合。于是存在尾部密度 \(\tau_n=1+o(1)\)（与定理 \ref{thm:xi-time-part60f-zg-zero-prefix-log-density} 同样的并集估计）使
\[
\mathrm{dens}\bigl(\mathrm{ZG}[e_1,\dots,e_n]\bigr)
=
\left(\prod_{\substack{1\le k\le m_0\\ e_k=1}}
\left(\frac1{p_k}-\frac1{p_k^2}\right)\right)
\left(\prod_{\substack{1\le k\le n\\ e_k=0}}
\left(1-\frac1{p_k}\right)\right)\tau_n.
\]
把有限前缀部分整理为
\[
\prod_{\substack{1\le k\le m_0\\ e_k=1}}
\left(\frac1{p_k}-\frac1{p_k^2}\right)
\prod_{\substack{1\le k\le m_0\\ e_k=0}}
\left(1-\frac1{p_k}\right)
=
\left(\prod_{k=1}^{m_0}\left(1-\frac1{p_k}\right)\right)
\prod_{\substack{1\le k\le m_0\\ e_k=1}}\frac1{p_k}.
\]
故
\[
\mathrm{dens}\bigl(\mathrm{ZG}[e_1,\dots,e_n]\bigr)
=
\left(\prod_{k=1}^{n}\left(1-\frac1{p_k}\right)\right)
\left(\prod_{\substack{1\le k\le m_0\\ e_k=1}}\frac1{p_k}\right)
(1+o(1)).
\]
由 Mertens 公式，存在
\[
C(e):=
\frac{e^{-\gamma}}{\delta_{\mathrm{ZG}}}
\prod_{\substack{1\le k\le m_0\\ e_k=1}}\frac1{p_k}
\in(0,\infty)
\]
使
\[
\nu_{\mathrm{ZG}}([e_1,\dots,e_n])
=
\frac{\mathrm{dens}(\mathrm{ZG}[e_1,\dots,e_n])}{\delta_{\mathrm{ZG}}}
\sim
\frac{C(e)}{\log p_n}.
\]

最后，由定理 \ref{thm:xi-time-part9g-holographic-prefix-isometry-on-line} 的柱--球对应，
\[
\Xi_B([e_1,\dots,e_n])
=
\left(\sum_{t=1}^{n}e_tB^{t-1}v+B^n\ZZ_2v\right)\cap K_{B,v}
=
\mathrm B_{K_{B,v}}(x_e,B^{-n}),
\]
其中 \(x_e=\Xi_B(e)\)。
因此
\[
\mu_B\bigl(\mathrm B_{K_{B,v}}(x_e,B^{-n})\bigr)
=
\nu_{\mathrm{ZG}}([e_1,\dots,e_n])
\sim
\frac{C(e)}{\log p_n}.
\]
于是
\[
\frac{\log \mu_B\bigl(\mathrm B_{K_{B,v}}(x_e,B^{-n})\bigr)}{\log B^{-n}}
=
\frac{\log C(e)-\log\log p_n+o(1)}{-n\log B}.
\]
由素数定理 \(p_n\sim n\log n\) 可知 \(\log\log p_n=o(n)\)，故上式趋于 \(0\)。证毕。
\end{proof}

\begin{corollary}[有限地址分辨率的不可消解模糊质量]\label{cor:xi-time-part60f-address-resolution-logarithmic-ambiguity-mass}
沿用定理 \ref{thm:xi-time-part60f-zg-cylinder-measure-local-dimension-zero} 的记号，对每个 \(n\ge 1\) 定义
\[
U_n
:=
\{x\in K_{B,v}:\ x\equiv 0\pmod{B^n\ZZ_2v}\}.
\]
则
\[
U_n=\Xi_B([0^n]),
\]
并且
\[
\mu_B(U_n)
=
\nu_{\mathrm{ZG}}([0^n])
\sim
\frac{e^{-\gamma}}{\delta_{\mathrm{ZG}}\log p_n}.
\]
因此，任何只读取前 \(n\) 个 \(2^L\)-进地址位的有限分辨率观测，其不可消解模糊质量至多只能降到 \((\log p_n)^{-1}\) 量级，而不可能指数消失。
\end{corollary}
\begin{proof}
由定理 \ref{thm:xi-time-part9g-holographic-prefix-isometry-on-line} 的柱--球对应，
\[
\Xi_B([0^n])=(0+B^n\ZZ_2v)\cap K_{B,v}=U_n.
\]
于是
\[
\mu_B(U_n)=\nu_{\mathrm{ZG}}([0^n])
=
\frac{\mathrm{dens}(\mathrm{ZG}[0^n])}{\delta_{\mathrm{ZG}}}.
\]
再由定理 \ref{thm:xi-time-part60f-zg-zero-prefix-log-density} 即得
\[
\mu_B(U_n)\sim \frac{e^{-\gamma}}{\delta_{\mathrm{ZG}}\log p_n}.
\]
证毕。
\end{proof}

\begin{theorem}[有限半径解析证书与有限 prime shadow 的统一强失明]\label{thm:xi-time-part60f-finite-radius-primeshadow-unified-strong-blindness}
固定 \(0<\rho_\star<1\) 与 \(n\ge 1\)。
考虑对象对
\[
(F,N),
\]
其中 \(F\) 为单位圆盘上的有限 Blaschke 乘积，\(N\in \mathrm{ZG}\)。
定义联合有限证书
\[
\mathcal C_{n,\rho_\star}(F,N)
:=
\Bigl(
\mathcal J_{\rho_\star}(F),\ 
(v_{p_1}(N),\dots,v_{p_n}(N))
\Bigr),
\]
其中 \(\mathcal J_{\rho_\star}(F)\) 为定理
\ref{thm:xi-time-part60e-truncated-jensen-profile-completeness-blindness}
中的截断 Jensen 剖面。
则存在一列两两不同的对象对
\[
\{(F_j,N_j)\}_{j\ge 1}
\]
使得它们具有完全相同的联合有限证书，并满足：
\begin{enumerate}
  \item \(F_j\) 两两不同，且每一步新增零点都位于环带
  \[
  \rho_\star<|z|<1;
  \]
  \item \(N_j\in \mathrm{ZG}[0^n]\) 两两不同；
  \item 该共同算术纤维满足
  \[
  \mathrm{dens}\bigl(\mathrm{ZG}[0^n]\bigr)
  \sim
  \frac{e^{-\gamma}}{\log p_n}.
  \]
\end{enumerate}
\end{theorem}
\begin{proof}
先处理解析侧。
取一列两两不同的实数
\[
a_j\in(\rho_\star,1)
\qquad(j\ge 1),
\]
并定义 Blaschke 因子
\[
B_{a_j}(z):=\frac{z-a_j}{1-a_j z},
\qquad
F_j:=F_1\cdot \prod_{r=2}^{j}B_{a_r}(z),
\]
其中 \(F_1\) 固定任意给定的有限 Blaschke 乘积。
由于所有新增零点都满足 \(|a_j|>\rho_\star\)，定理
\ref{thm:xi-time-part60e-truncated-jensen-profile-completeness-blindness}
给出
\[
\mathcal J_{\rho_\star}(F_j)=\mathcal J_{\rho_\star}(F_1)
\qquad(j\ge 1),
\]
而 \(\{a_j\}\) 两两不同，故 \(F_j\) 两两不同。

再处理算术侧。
由定理 \ref{thm:xi-time-part60f-zg-zero-prefix-log-density}，
\[
\mathrm{dens}\bigl(\mathrm{ZG}[0^n]\bigr)>0.
\]
故 \(\mathrm{ZG}[0^n]\) 必为无限集。任选一列两两不同的元素
\[
N_j\in \mathrm{ZG}[0^n]
\qquad(j\ge 1).
\]
则对每个 \(j\)，都有
\[
v_{p_1}(N_j)=\cdots=v_{p_n}(N_j)=0,
\]
从而它们的有限 prime shadow 完全相同。

综上，对所有 \(j\ge 1\)，都有
\[
\mathcal C_{n,\rho_\star}(F_j,N_j)
=
\Bigl(\mathcal J_{\rho_\star}(F_1),(0,\dots,0)\Bigr),
\]
即得到一列联合有限证书完全相同而对象两两不同的序列。
最后，第 \(3\) 条正是定理 \ref{thm:xi-time-part60f-zg-zero-prefix-log-density} 的内容。证毕。
\end{proof}

\noindent
以上结论把地址侧与算术侧压缩为同一条可直接调用的主链：有限前缀不仅在 \(2^L\)-进地址空间中切出一个精确 cylinder，同时也在 \(\mathrm{ZG}\) 内切出一个严格正密度的算术纤维；当该前缀取为全零时，这一纤维的质量只按 \((\log p_n)^{-1}\) 衰减，而非按几何半径的幂级衰减。于是，有限半径解析证书与有限 prime shadow 即便同时给出，也仍然保留可数无限的对象纤维，并且这条失明不是纯存在性陈述，而带有明确的对数量级下界。

\endinput


\paragraph{算术自然密度测度、可数 prime-register solenoid 与有限维可见性的统一塌缩}
在黄金地址柱测度的对数级质量律、无穷 prime-register solenoid 的有限可见因子化，以及联合有限证书不完备性已经建立之后，还可进一步抽取出下列五条结构后果。

\begin{theorem}[算术自然密度诱导的黄金柱测度不是任何连续势的 Gibbs 态]\label{thm:xi-time-part60f2-zg-density-measure-non-gibbs}
\leanverified{paper\_xi\_time\_part60f2\_zg\_density\_measure\_non\_gibbs}
沿用定理 \ref{thm:xi-time-part60f-zg-cylinder-measure-local-dimension-zero} 的记号，记
\[
\sigma:\Sigma_{\mathrm{ZG}}\to \Sigma_{\mathrm{ZG}},
\qquad
\sigma(z_1,z_2,z_3,\dots)=(z_2,z_3,\dots).
\]
则自然密度诱导的柱测度
\[
\nu_{\mathrm{ZG}}
\]
不是 \((\Sigma_{\mathrm{ZG}},\sigma)\) 上任何连续势
\[
\phi\in C(\Sigma_{\mathrm{ZG}})
\]
的 Gibbs 测度。
\end{theorem}
\begin{proof}
对每个 \(n\ge 1\)，记
\[
[0^n]:=[(0,\dots,0)]\subset \Sigma_{\mathrm{ZG}}.
\]
由定理 \ref{thm:xi-time-part60f-zg-zero-prefix-log-density} 与定理 \ref{thm:xi-time-part60f-zg-cylinder-measure-local-dimension-zero}，
\[
\nu_{\mathrm{ZG}}([0^n])
=
\frac{\mathrm{dens}(\mathrm{ZG}[0^n])}{\delta_{\mathrm{ZG}}}
\sim
\frac{e^{-\gamma}}{\delta_{\mathrm{ZG}}\log p_n}.
\]
特别地，
\[
\nu_{\mathrm{ZG}}([0^n])\to 0,
\]
并且该衰减仅为次指数级。

反设 \(\nu_{\mathrm{ZG}}\) 是某个连续势 \(\phi\in C(\Sigma_{\mathrm{ZG}})\) 的 Gibbs 测度。则按 Gibbs 定义，存在常数 \(C\ge 1\) 使对一切 \(n\ge 1\) 都有
\[
C^{-1}
\le
\nu_{\mathrm{ZG}}([0^n])
\exp\!\bigl(nP(\phi)-S_n\phi(0^\infty)\bigr)
\le
C,
\]
其中 \(0^\infty=(0,0,\dots)\in\Sigma_{\mathrm{ZG}}\) 为全零固定点。
由于
\[
S_n\phi(0^\infty)=n\phi(0^\infty),
\]
便得到
\[
\nu_{\mathrm{ZG}}([0^n])\asymp \exp\!\bigl(-n(P(\phi)-\phi(0^\infty))\bigr).
\]
若 \(P(\phi)<\phi(0^\infty)\)，则右端随 \(n\) 指数增长，与 \(\nu_{\mathrm{ZG}}([0^n])\le 1\) 矛盾。
若 \(P(\phi)=\phi(0^\infty)\)，则 \(\nu_{\mathrm{ZG}}([0^n])\) 有正下界，这与 \(\nu_{\mathrm{ZG}}([0^n])\to 0\) 矛盾。
若 \(P(\phi)>\phi(0^\infty)\)，则 \(\nu_{\mathrm{ZG}}([0^n])\) 必指数衰减；而
\[
\nu_{\mathrm{ZG}}([0^n])\sim \frac{e^{-\gamma}}{\delta_{\mathrm{ZG}}\log p_n}
\]
仅按 \((\log p_n)^{-1}\) 衰减。由素数定理 \(p_n\sim n\log n\) 可知
\[
\log\log p_n=o(n),
\]
故这不可能与任意指数律等价。三种情形皆矛盾，因此 \(\nu_{\mathrm{ZG}}\) 不可能是任何连续势的 Gibbs 测度。证毕。
\end{proof}

\begin{theorem}[可数 prime-register solenoid 之间连续态射的完全分类]\label{thm:xi-time-part60f2-countable-solenoid-morphism-classification}
\leanverified{paper\_xi\_time\_part60f2\_countable\_solenoid\_morphism\_classification}
设 \(S,T\subset\mathbb P\) 为可数素数集，并记
\[
G_S:=\ZZ[S^{-1}],
\qquad
G_T:=\ZZ[T^{-1}],
\qquad
\Sigma_S:=\widehat{G_S},
\qquad
\Sigma_T:=\widehat{G_T}.
\]
则存在自然同构
\[
\operatorname{Hom}_{\ZZ}(G_S,G_T)\cong
\begin{cases}
0,& S\not\subseteq T,\\[1mm]
G_T,& S\subseteq T,
\end{cases}
\]
从而对偶化得到
\[
\operatorname{Hom}_{\mathrm{cts}}(\Sigma_T,\Sigma_S)\cong
\begin{cases}
0,& S\not\subseteq T,\\[1mm]
G_T,& S\subseteq T.
\end{cases}
\]
\end{theorem}
\begin{proof}
先设 \(S\subseteq T\)。
任取 \(x\in G_T\)，定义
\[
m_x:G_S\to G_T,
\qquad
y\mapsto xy.
\]
由于 \(G_S,G_T\subset\QQ\)，且 \(y\in G_S\) 的分母素数仅来自 \(S\subseteq T\)，故 \(xy\in G_T\)，从而 \(m_x\) 为良定义的加法同态。

反过来，任取 \(\phi\in \operatorname{Hom}_{\ZZ}(G_S,G_T)\)，并记
\[
x:=\phi(1)\in G_T.
\]
对任意 \(y=a/b\in G_S\)（其中 \(a\in\ZZ\)，\(b\) 的全部素因子属于 \(S\)），有
\[
b\,\phi(y)=\phi(by)=\phi(a)=a\,\phi(1)=ax.
\]
由于 \(G_T\subset\QQ\) 无挠，故
\[
\phi(y)=\frac{a}{b}x=xy=m_x(y).
\]
因此 \(\phi=m_x\)，从而
\[
\operatorname{Hom}_{\ZZ}(G_S,G_T)\cong G_T.
\]

再设 \(S\not\subseteq T\)，取 \(p\in S\setminus T\)。
任取 \(\phi\in\operatorname{Hom}_{\ZZ}(G_S,G_T)\)，并记 \(x:=\phi(1)\)。
由于 \(p^{-n}\in G_S\) 对每个 \(n\ge 1\) 都成立，故
\[
x=\phi(1)=p^n\phi(p^{-n})\in p^nG_T
\qquad(n\ge 1).
\]
于是
\[
x\in \bigcap_{n\ge 1}p^nG_T.
\]
若 \(u\in G_T\setminus\{0\}\)，则可写成
\[
u=\frac{a}{\prod_{q\in T}q^{m_q}}
\]
其中 \(a\in\ZZ\setminus\{0\}\)，且分母与 \(p\) 互素。
若 \(u\in p^nG_T\)，则把两边写成既约分式可知 \(p^n\mid a\)。
这不可能对所有 \(n\) 同时成立，故
\[
\bigcap_{n\ge 1}p^nG_T=\{0\}.
\]
于是 \(x=0\)。
再对任意 \(y=a/b\in G_S\) 有
\[
b\,\phi(y)=\phi(a)=a\,\phi(1)=0,
\]
而 \(G_T\) 无挠，故 \(\phi(y)=0\)。
因此
\[
\operatorname{Hom}_{\ZZ}(G_S,G_T)=0.
\]

最后，Pontryagin 对偶给出反变等价
\[
\operatorname{Hom}_{\mathrm{cts}}(\Sigma_T,\Sigma_S)\cong \operatorname{Hom}_{\ZZ}(G_S,G_T),
\]
故连续态射的分类随即成立。证毕。
\end{proof}

\begin{theorem}[无穷 prime-register solenoid 的有限维连续表示只看见有限个素数轴]\label{thm:xi-time-part60f2-finite-dimensional-representation-finite-prime-factorization}
\leanverified{paper\_xi\_time\_part60f2\_finite\_dimensional\_representation\_finite\_prime\_factorization}
设 \(S\subset\mathbb P\) 为可数无穷素数集，并记
\[
\Sigma_S:=\widehat{\ZZ[S^{-1}]}.
\]
任取连续酉表示
\[
\rho:\Sigma_S\to U(d).
\]
则存在有限子集
\[
S_0\Subset S
\]
与连续酉表示
\[
\rho_0:\Sigma_{S_0}\to U(d)
\]
使得
\[
\rho=\rho_0\circ \pi_{S\to S_0}.
\]
因此余有限尾核
\[
K_{S\setminus S_0}:=\prod_{p\in S\setminus S_0}\ZZ_p
\]
满足
\[
K_{S\setminus S_0}\subseteq \ker\rho.
\]
此外，\(\rho(\Sigma_S)\) 的 Lie 秩至多为 \(1\)；更强地，
\[
\rho(\Sigma_S)
\]
只能是平凡群或一维圆群 \(\TT\)。
\end{theorem}
\begin{proof}
由于 \(\Sigma_S\) 为紧 Abel 群，\(\rho(\Sigma_S)\) 是一族彼此交换的酉矩阵，故可同时酉对角化。
于是存在 \(U\in U(d)\) 与连续特征
\[
\chi_1,\dots,\chi_d:\Sigma_S\to \TT
\]
使得对每个 \(\xi\in\Sigma_S\) 都有
\[
U\rho(\xi)U^{-1}
=
\operatorname{diag}\bigl(\chi_1(\xi),\dots,\chi_d(\xi)\bigr).
\]
由 Pontryagin 对偶，
\[
\operatorname{Hom}_{\mathrm{cts}}(\Sigma_S,\TT)\cong G_S:=\ZZ[S^{-1}],
\]
故每个 \(\chi_j\) 对应于某个 \(x_j\in G_S\)。
每个 \(x_j\) 的分母只涉及有限多个素数。
令 \(S_0\subset S\) 为所有 \(x_j\) 的分母中出现之素数的并集，则 \(S_0\) 有限，且 \(x_j\in G_{S_0}\) 对所有 \(j\) 都成立。
于是每个 \(\chi_j\) 都经由
\[
\pi_{S\to S_0}:\Sigma_S\twoheadrightarrow \Sigma_{S_0}
\]
因子化，从而整个 \(\rho\) 也经由 \(\pi_{S\to S_0}\) 因子化。

由推论 \ref{cor:xi-time-part60e-infinite-solenoid-finite-visible-factorization} 的短正合列，
\[
0\to \prod_{p\in S\setminus S_0}\ZZ_p\to \Sigma_S\xrightarrow{\ \pi_{S\to S_0}\ }\Sigma_{S_0}\to 0,
\]
立得
\[
K_{S\setminus S_0}\subseteq \ker\rho.
\]

最后分析像的 Lie 秩。
在对角化后，\(\rho\) 可视为映射
\[
\Sigma_S\to \TT^d.
\]
其对偶离散映射为
\[
\rho^\vee:\ZZ^d\to G_S,
\qquad
e_j\mapsto x_j.
\]
因此
\[
\widehat{\rho(\Sigma_S)}\cong \operatorname{im}(\rho^\vee)=:H
=
\langle x_1,\dots,x_d\rangle\subset G_S\subset\QQ.
\]
而 \(\QQ\) 的任意有限生成子群都循环：取 \(x_1,\dots,x_d\) 的一个公共分母 \(m\)，则
\[
H\subseteq \frac1m\ZZ,
\]
故 \(H\) 是 \((1/m)\ZZ\) 的有限生成子群，从而必为零群或无限循环群。
因此其 Pontryagin 对偶
\[
\rho(\Sigma_S)\cong \widehat H
\]
只能是平凡群或 \(\TT\)。
特别地，\(\rho(\Sigma_S)\) 的 Lie 秩至多为 \(1\)。证毕。
\end{proof}

\begin{corollary}[无穷 prime-register 终对象不存在忠实有限维 Lie 模型]\label{cor:xi-time-part60f2-no-faithful-finite-dimensional-lie-model}
\leanverified{paper\_xi\_time\_part60f2\_no\_faithful\_finite\_dimensional\_lie\_model}
在定理 \ref{thm:xi-time-part60f2-finite-dimensional-representation-finite-prime-factorization} 的设定下，\(\Sigma_S\) 不存在忠实的有限维连续酉表示。
更一般地，不存在到任何有限维紧 Lie 群的忠实连续商实现。
\end{corollary}
\begin{proof}
若 \(\rho:\Sigma_S\to U(d)\) 忠实，则由定理 \ref{thm:xi-time-part60f2-finite-dimensional-representation-finite-prime-factorization}，存在有限 \(S_0\Subset S\) 使
\[
K_{S\setminus S_0}=\prod_{p\in S\setminus S_0}\ZZ_p
\subseteq \ker\rho.
\]
由于 \(S\setminus S_0\) 非空，\(K_{S\setminus S_0}\) 为非平凡闭子群，从而 \(\ker\rho\neq \{0\}\)，矛盾。

若进一步存在一个到有限维紧 Lie 群的忠实连续商实现
\[
q:\Sigma_S\twoheadrightarrow G,
\qquad
\ker q=\{0\},
\]
则由 Peter--Weyl 定理，\(G\) 存在某个忠实有限维酉表示
\[
\iota:G\hookrightarrow U(d).
\]
于是
\[
\iota\circ q:\Sigma_S\to U(d)
\]
给出忠实有限维连续酉表示，与上一步矛盾。故所述 Lie 模型亦不存在。证毕。
\end{proof}

\begin{theorem}[有限半径 Jensen 证书与有限维连续表示的统一塌缩]\label{thm:xi-time-part60f2-finite-radius-finite-dimensional-unified-collapse}
设 \(S\subset\mathbb P\) 为可数无穷素数集，固定
\[
0<\rho_\star<1,
\]
并取任意有限维连续酉表示
\[
\varrho:\Sigma_S\to U(d).
\]
则存在一列对象对
\[
\{(F_n,\xi_n)\}_{n\ge 1},
\]
其中每个 \(F_n\) 都是单位圆盘上的有限 Blaschke 乘积，且 \(\xi_n\in \Sigma_S\)，满足
\begin{enumerate}
  \item \(\mathcal J_{\rho_\star}(F_n)\) 全部相同；
  \item \(\varrho(\xi_n)\) 全部相同；
  \item \(F_n\) 两两不同；
  \item \(\xi_n\) 两两不同。
\end{enumerate}
因此，任何把有限半径 Jensen 剖面与有限维连续算术观测并合的有限可见证书都不可能完备。
\end{theorem}
\begin{proof}
由定理 \ref{thm:xi-time-part60f2-finite-dimensional-representation-finite-prime-factorization}，存在有限子集 \(S_0\Subset S\) 与连续酉表示
\[
\varrho_0:\Sigma_{S_0}\to U(d)
\]
使
\[
\varrho=\varrho_0\circ \pi_{S\to S_0}.
\]
再由定理 \ref{thm:xi-time-part60e-finite-radius-finite-prime-certificate-incompleteness}，存在一列对象对
\[
\{(F_n,\xi_n)\}_{n\ge 1}
\]
使得
\[
\mathcal J_{\rho_\star}(F_n)
=
\mathcal J_{\rho_\star}(F_1),
\qquad
\pi_{S\to S_0}(\xi_n)
=
\pi_{S\to S_0}(\xi_1),
\]
并且 \(F_n\) 两两不同、\(\xi_n\) 两两不同。
于是对所有 \(n\ge 1\)，
\[
\varrho(\xi_n)
=
\varrho_0\!\bigl(\pi_{S\to S_0}(\xi_n)\bigr)
=
\varrho_0\!\bigl(\pi_{S\to S_0}(\xi_1)\bigr)
=
\varrho(\xi_1).
\]
这就得到所需序列。证毕。
\end{proof}

\noindent
上述五条结果把 H.161--H.162 的地址柱密度链与局部化数系的 prime-register 对偶链压缩为同一条有限可见性判别律：算术自然密度诱导的地址测度不是任何连续势的 Gibbs 平衡态；而在 prime-register 终对象一旦进入可数无穷层级之后，所有有限维连续观测都只能通过有限个素数轴，并必然湮灭一条余有限 \(p\)-进尾核。因而，有限半径解析证书与有限维连续表示即使联合，也仍然在解析外环与算术尾核两个方向上同时失明。


\endinput

\paragraph{局部化有限商、周期动力学、容量原子恢复与缺陷重整化}
在 prime-register 的有限商分类、局部化 solenoid 的周期谱、bin-fold 临界容量的原子 defect law 以及 fold-gauge 常数列的前两阶 Bernoulli 模态已经建立之后，还可进一步抽取出下列六条彼此闭合的结构后果。

\begin{theorem}[局部化数系有限商的精确外部化公式]\label{thm:xi-localized-quotient-exact-sfree-externalization}
\leanverified{paper\_xi\_localized\_quotient\_exact\_sfree\_externalization}
设 \(S\subset\mathbb P\) 为有限素数集，并记
\[
G_S:=\ZZ[S^{-1}].
\]
对任意 \(n\ge 1\)，令
\[
n_S:=\prod_{p\in S}p^{v_p(n)},
\qquad
n_{S^\perp}:=\frac{n}{n_S}.
\]
则有自然同构
\[
\boxed{\;
G_S/nG_S\cong \ZZ/n_{S^\perp}\ZZ.\;}
\]
特别地，\(G_S\) 的每一个有限商都只由 \(n\) 的 \(S\)-自由部分决定。
\end{theorem}
\begin{proof}
对每个 \(p\in S\)，元素 \(p^{-1}\in G_S\) 存在，故乘法映射
\[
x\longmapsto px
\]
在 \(G_S\) 上可逆。于是对
\[
n=n_S\,n_{S^\perp}
\]
有
\[
nG_S=n_Sn_{S^\perp}G_S=n_{S^\perp}G_S.
\]
因此只需证明当
\[
\gcd\!\Bigl(n,\prod_{p\in S}p\Bigr)=1
\]
时，
\[
G_S/nG_S\cong \ZZ/n\ZZ.
\]

任取 \(x\in G_S\)，可写成
\[
x=\frac{a}{s},
\qquad
a\in\ZZ,
\]
其中 \(s\) 的全部素因子都属于 \(S\)。由于 \(\gcd(s,n)=1\)，\(s\) 在 \(\ZZ/n\ZZ\) 中可逆。定义
\[
\pi_n:G_S\longrightarrow \ZZ/n\ZZ,
\qquad
\pi_n\!\left(\frac{a}{s}\right):=a\,s^{-1}\pmod n.
\]
若 \(a/s=a'/s'\) 于 \(G_S\) 中相等，则 \(as'=a's\)，从而
\[
a\,s^{-1}\equiv a'\,s'^{-1}\pmod n,
\]
故 \(\pi_n\) 良定义。又因整数子群 \(\ZZ\subset G_S\) 已映满 \(\ZZ/n\ZZ\)，故 \(\pi_n\) 为满射。

再看其核。若
\[
\pi_n\!\left(\frac{a}{s}\right)=0,
\]
则 \(n\mid a\)。写 \(a=nb\)，便有
\[
\frac{a}{s}=n\cdot \frac{b}{s}\in nG_S.
\]
故
\[
\ker\pi_n\subset nG_S.
\]
反向包含显然成立，于是
\[
\ker\pi_n=nG_S.
\]
由第一同构定理得到
\[
G_S/nG_S\cong \ZZ/n\ZZ.
\]
最后把 \(n\) 替换为 \(n_{S^\perp}\) 即得结论。证毕。
\end{proof}

\begin{theorem}[局部化 solenoid 自覆盖的固定点公式与 Artin--Mazur \texorpdfstring{$\zeta$}{zeta} 函数]\label{thm:xi-localized-solenoid-fixedpoints-artin-mazur-series}
\leanverified{paper\_xi\_localized\_solenoid\_fixedpoints\_artin\_mazur\_series}
沿用定理 \ref{thm:xi-localized-quotient-exact-sfree-externalization} 的记号，并记
\[
\Sigma_S:=\widehat{G_S}.
\]
设整数 \(a\ge 2\) 满足
\[
\gcd\!\Bigl(a,\prod_{p\in S}p\Bigr)=1,
\]
令
\[
f_a:\Sigma_S\longrightarrow \Sigma_S
\]
为离散端“乘以 \(a\)”的 Pontryagin 对偶。则对每个 \(k\ge 1\)，有
\[
\boxed{\;
\mathrm{Fix}(f_a^k)\cong \widehat{G_S/(a^k-1)G_S},
\qquad
\#\mathrm{Fix}(f_a^k)=(a^k-1)_{S^\perp}.\;}
\]
并且其 Artin--Mazur \(\zeta\)-函数满足
\[
\boxed{\;
\zeta_{f_a}(z)
:=
\exp\!\left(
\sum_{k\ge 1}\frac{\#\mathrm{Fix}(f_a^k)}{k}z^k
\right)
=
\exp\!\left(
\sum_{k\ge 1}\frac{(a^k-1)_{S^\perp}}{k}z^k
\right).\;}
\]
特别地，当 \(S=\varnothing\) 时，
\[
\boxed{\;
\zeta_{f_a}(z)=\frac{1-z}{1-az}.\;}
\]
\end{theorem}
\begin{proof}
由对偶定义，对任意 \(x\in \Sigma_S\) 与 \(g\in G_S\)，有
\[
(f_a^k x)(g)=x(a^k g).
\]
因此
\[
x\in \mathrm{Fix}(f_a^k)
\iff
x(a^k g)=x(g)\quad(\forall g\in G_S),
\]
亦即
\[
x\bigl((a^k-1)g\bigr)=1\quad(\forall g\in G_S).
\]
这等价于 \(x\) 因子化经过商群
\[
G_S/(a^k-1)G_S.
\]
故有自然同构
\[
\mathrm{Fix}(f_a^k)\cong \widehat{G_S/(a^k-1)G_S}.
\]

再由定理 \ref{thm:xi-localized-quotient-exact-sfree-externalization}，
\[
G_S/(a^k-1)G_S\cong \ZZ/(a^k-1)_{S^\perp}\ZZ.
\]
有限循环群的 Pontryagin 对偶仍为同阶循环群，于是
\[
\#\mathrm{Fix}(f_a^k)=(a^k-1)_{S^\perp}.
\]
把该计数代入 Artin--Mazur \(\zeta\)-函数的定义，即得所示显式级数。

当 \(S=\varnothing\) 时，\((a^k-1)_{S^\perp}=a^k-1\)，故
\[
\log\zeta_{f_a}(z)
=
\sum_{k\ge 1}\frac{a^k-1}{k}z^k
=
-\log(1-az)+\log(1-z).
\]
指数化即得
\[
\zeta_{f_a}(z)=\frac{1-z}{1-az}.
\]
证毕。
\end{proof}

\begin{theorem}[bin-fold 极限容量曲线的原子测度唯一恢复]\label{thm:xi-binfold-capacity-atomic-measure-unique-recovery}
把定理 \ref{thm:xi-binfold-critical-capacity-defect-atomic-law} 的极限容量曲线写为
\[
\widehat C_\infty(s)
:=
\frac{\varphi}{\sqrt5}\min\{1,s\}
+
\frac{1}{\sqrt5}\min\{\varphi^{-1},s\}
\qquad (s>0).
\]
则存在唯一有限正测度
\[
\boxed{\;
\eta_\infty
:=
\frac{1}{\sqrt5}\delta_{\varphi^{-1}}
+
\frac{\varphi}{\sqrt5}\delta_{1}\;}
\]
使得对一切 \(s>0\) 都有
\[
\boxed{\;
\widehat C_\infty(s)=\int_{(0,\infty)}\min\{s,y\}\,\eta_\infty(\dd y).\;}
\]
并且在分布意义下，
\[
\boxed{\;
\frac{\dd}{\dd s}\widehat C_\infty(s)=\eta_\infty([s,\infty)),
\qquad
-\frac{\dd^2}{\dd s^2}\widehat C_\infty=\eta_\infty.\;}
\]
\end{theorem}
\begin{proof}
取
\[
\eta_\infty
:=
\frac{1}{\sqrt5}\delta_{\varphi^{-1}}
+
\frac{\varphi}{\sqrt5}\delta_1.
\]
则直接计算得到
\[
\int_{(0,\infty)}\min\{s,y\}\,\eta_\infty(\dd y)
=
\frac{1}{\sqrt5}\min\{s,\varphi^{-1}\}
+
\frac{\varphi}{\sqrt5}\min\{s,1\}
=
\widehat C_\infty(s).
\]
故积分表示成立。

再对各分段求导，得到
\[
\widehat C_\infty'(s)=
\begin{cases}
\dfrac{\varphi+1}{\sqrt5},& 0<s<\varphi^{-1},\\[2mm]
\dfrac{\varphi}{\sqrt5},& \varphi^{-1}<s<1,\\[2mm]
0,& s>1.
\end{cases}
\]
另一方面，
\[
\eta_\infty([s,\infty))=
\begin{cases}
\dfrac{1}{\sqrt5}+\dfrac{\varphi}{\sqrt5},& 0<s<\varphi^{-1},\\[2mm]
\dfrac{\varphi}{\sqrt5},& \varphi^{-1}<s<1,\\[2mm]
0,& s>1,
\end{cases}
\]
两者一致，故第一条分布恒等式成立。

再对上式取分布导数。由于 \(\widehat C_\infty'\) 只在
\[
s=\varphi^{-1},\qquad s=1
\]
发生跃迁，且两处跳跃大小分别为
\[
\frac{1}{\sqrt5},
\qquad
\frac{\varphi}{\sqrt5},
\]
故
\[
-\frac{\dd^2}{\dd s^2}\widehat C_\infty
=
\frac{1}{\sqrt5}\delta_{\varphi^{-1}}
+
\frac{\varphi}{\sqrt5}\delta_1
=
\eta_\infty.
\]

唯一性亦随之而来：若另一有限正测度 \(\eta\) 满足同一积分表示，则对函数
\[
s\longmapsto \int_{(0,\infty)}\min\{s,y\}\,\eta(\dd y)
\]
作二阶分布导数便得到
\[
\eta=-\widehat C_\infty''=\eta_\infty.
\]
故 \(\eta_\infty\) 唯一。证毕。
\end{proof}

\begin{theorem}[bin-fold 容量极限是稳定层均匀二原子律的尺寸偏置平均]\label{thm:xi-binfold-capacity-size-bias-uniform-law}
\leanverified{paper\_xi\_binfold\_capacity\_size\_bias\_uniform\_law}
沿用定理 \ref{thm:xi-foldbin-uniform-layer-degeneracy-two-atom-law} 的记号，记 \(Y_\infty^{\mathrm{unif}}\) 为该定理中的极限随机变量，即
\[
\PP\!\bigl(Y_\infty^{\mathrm{unif}}=1\bigr)=\varphi^{-1},
\qquad
\PP\!\bigl(Y_\infty^{\mathrm{unif}}=\varphi^{-1}\bigr)=\varphi^{-2}.
\]
则前一定理中的原子测度满足
\[
\boxed{\;
\eta_\infty
=
\frac{\varphi^2}{\sqrt5}\,
\mathcal L\!\bigl(Y_\infty^{\mathrm{unif}}\bigr).\;}
\]
等价地，对每个 \(s>0\) 都有
\[
\boxed{\;
\widehat C_\infty(s)
=
\frac{\varphi^2}{\sqrt5}\,
\EE\!\bigl[\min\{Y_\infty^{\mathrm{unif}},s\}\bigr].\;}
\]

进一步，记 \(Y_\infty^\mu\) 为推论 \ref{cor:fold-bin-two-point-limit-law} 的 bin-fold 推前极限随机变量。则 \(Y_\infty^\mu\) 正是 \(Y_\infty^{\mathrm{unif}}\) 的尺寸偏置变换，即
\[
\boxed{\;
\PP\!\bigl(Y_\infty^\mu=y\bigr)
=
\frac{
y\,\PP\!\bigl(Y_\infty^{\mathrm{unif}}=y\bigr)
}{
\EE\!\bigl[Y_\infty^{\mathrm{unif}}\bigr]
}
\qquad
\bigl(y\in\{1,\varphi^{-1}\}\bigr).\;}
\]
特别地，
\[
\PP\!\bigl(Y_\infty^\mu=1\bigr)=\frac{\varphi}{\sqrt5},
\qquad
\PP\!\bigl(Y_\infty^\mu=\varphi^{-1}\bigr)=\frac{1}{\varphi\sqrt5}.
\]
\end{theorem}
\begin{proof}
由定理 \ref{thm:xi-foldbin-uniform-layer-degeneracy-two-atom-law}，
\[
\mathcal L\!\bigl(Y_\infty^{\mathrm{unif}}\bigr)
=
\frac{1}{\varphi}\delta_1+\frac{1}{\varphi^2}\delta_{\varphi^{-1}}.
\]
两边同乘 \(\varphi^2/\sqrt5\)，便得到
\[
\frac{\varphi^2}{\sqrt5}\,
\mathcal L\!\bigl(Y_\infty^{\mathrm{unif}}\bigr)
=
\frac{\varphi}{\sqrt5}\delta_1+\frac{1}{\sqrt5}\delta_{\varphi^{-1}}
=
\eta_\infty,
\]
其中最后一步正是定理 \ref{thm:xi-binfold-capacity-atomic-measure-unique-recovery} 中的原子测度公式。于是
\[
\widehat C_\infty(s)
=
\int_{(0,\infty)}\min\{s,y\}\,\eta_\infty(\dd y)
=
\frac{\varphi^2}{\sqrt5}
\EE\!\bigl[\min\{Y_\infty^{\mathrm{unif}},s\}\bigr].
\]

再由定理 \ref{thm:xi-foldbin-uniform-layer-degeneracy-two-atom-law} 在 \(a=1\) 时的矩公式，
\[
\EE\!\bigl[Y_\infty^{\mathrm{unif}}\bigr]
=
\frac{1}{\varphi}+\varphi^{-3}
=
\frac{\sqrt5}{\varphi^2},
\]
其中最后一步由 \(\varphi^2=\varphi+1\) 直接化简可得。故尺寸偏置概率为
\[
\frac{1\cdot \varphi^{-1}}{\sqrt5/\varphi^2}=\frac{\varphi}{\sqrt5},
\qquad
\frac{\varphi^{-1}\cdot \varphi^{-2}}{\sqrt5/\varphi^2}
=
\frac{1}{\varphi\sqrt5}.
\]
这与推论 \ref{cor:fold-bin-two-point-limit-law} 的两点极限概率完全一致，故 \(Y_\infty^\mu\) 正是 \(Y_\infty^{\mathrm{unif}}\) 的尺寸偏置变换。证毕。
\end{proof}

\begin{theorem}[bin-fold 容量极限是统一层矩谱的 Hardy--Littlewood 变换并满足 Mellin--Stieltjes 完全对偶]\label{thm:xi-binfold-capacity-hl-mellin-duality-uniform-law}
沿用定理 \ref{thm:xi-binfold-capacity-size-bias-uniform-law} 与定理 \ref{thm:xi-foldbin-uniform-layer-two-point-exponential-family} 的记号。记
\[
Y_\infty^{\mathrm{unif}}
\sim
\frac{1}{\varphi}\delta_1+\frac{1}{\varphi^2}\delta_{\varphi^{-1}},
\qquad
Z(q):=\EE\!\left[\bigl(Y_\infty^{\mathrm{unif}}\bigr)^q\right]
=
\frac{1}{\varphi}+\varphi^{-q-2}.
\]
则：
\begin{enumerate}
  \item 对每个 \(s>0\)，bin-fold 极限容量曲线满足 Hardy--Littlewood 变换公式
  \[
  \boxed{\;
  \widehat C_\infty(s)
  =
  \frac{\varphi^2}{\sqrt5}\,
  \EE\!\bigl[\min\{Y_\infty^{\mathrm{unif}},s\}\bigr].\;}
  \]
  \item 对任意 \(0<q<1\)，有精确 Mellin--Stieltjes 对偶
  \[
  \boxed{\;
  Z(q)
  =
  \frac{\sqrt5}{\varphi^2}\,
  q(1-q)
  \int_0^\infty s^{q-2}\widehat C_\infty(s)\,\dd s.\;}
  \]
\end{enumerate}
\end{theorem}
\begin{proof}
第一条正是定理 \ref{thm:xi-binfold-capacity-size-bias-uniform-law} 的第二个盒装公式。

现证第二条。先记
\[
F(s):=\EE\!\bigl[\min\{Y_\infty^{\mathrm{unif}},s\}\bigr]
\qquad (s>0).
\]
对任意非负随机变量 \(Y\) 与任意 \(0<q<1\)，都有通用恒等式
\[
\EE[Y^q]
=
q(1-q)\int_0^\infty s^{q-2}\EE[\min\{Y,s\}]\,\dd s.
\]
事实上，由 layer-cake 公式，
\[
\EE[Y^q]
=
q\int_0^\infty s^{q-1}\PP(Y\ge s)\,\dd s,
\]
而
\[
\EE[\min\{Y,s\}]
=
\int_0^s \PP(Y\ge t)\,\dd t
\]
对 \(s\) 的导数正是 \(\PP(Y\ge s)\)。对右端作一次分部积分，即得上式。

将其应用于 \(Y=Y_\infty^{\mathrm{unif}}\)，便得到
\[
Z(q)
=
q(1-q)\int_0^\infty s^{q-2}F(s)\,\dd s.
\]
再由第一条，
\[
F(s)=\frac{\sqrt5}{\varphi^2}\,\widehat C_\infty(s),
\]
故
\[
Z(q)
=
\frac{\sqrt5}{\varphi^2}\,
q(1-q)\int_0^\infty s^{q-2}\widehat C_\infty(s)\,\dd s.
\]
证毕。
\end{proof}

\begin{theorem}[fold-gauge 缺陷序列的三次重整化正规形]\label{thm:xi-fold-gauge-defect-cubic-renormalization-normal-form}
沿用定理 \ref{thm:xi-foldbin-gauge-constant-not-cfinite-first-two-zeta-limits} 的记号，令
\[
E_m:=\log C_m-\log(2\pi),
\qquad
q:=\frac{\varphi}{2}.
\]
则当 \(m\to\infty\) 时，
\[
\boxed{\;
E_{m+1}
=
qE_m+\frac{3(2+\sqrt5)}{32}\,E_m^3+O(E_m^5).\;}
\]
等价地，
\[
\boxed{\;
\lim_{m\to\infty}
\frac{E_{m+1}-qE_m}{E_m^3}
=
\frac{3(2+\sqrt5)}{32}.\;}
\]
\end{theorem}
\begin{proof}
由定理 \ref{thm:xi-foldbin-gauge-constant-not-cfinite-first-two-zeta-limits}，
\[
E_m
=
\frac{1}{3\varphi}q^m-\frac{\sqrt5}{180}q^{3m}+O(q^{5m}).
\]
记
\[
a:=\frac{1}{3\varphi},
\qquad
b:=-\frac{\sqrt5}{180}.
\]
则
\[
E_m=aq^m+bq^{3m}+O(q^{5m}).
\]
因此
\[
E_{m+1}-qE_m
=
b(q^3-q)q^{3m}+O(q^{5m}),
\]
而
\[
E_m^3
=
a^3 q^{3m}+O(q^{5m}).
\]
于是
\[
\frac{E_{m+1}-qE_m}{E_m^3}
=
\frac{b(q^3-q)}{a^3}+O(q^{2m}).
\]
直接代入
\[
a=\frac{1}{3\varphi},
\qquad
b=-\frac{\sqrt5}{180},
\qquad
q=\frac{\varphi}{2}
\]
化简可得
\[
\frac{b(q^3-q)}{a^3}
=
\frac{3(2+\sqrt5)}{32}.
\]
故所示极限成立。

最后，由 \(a>0\) 可知
\[
E_m\sim aq^m,
\]
从而 \(q^{5m}=O(E_m^5)\)。把上式改写即得
\[
E_{m+1}
=
qE_m+\frac{3(2+\sqrt5)}{32}\,E_m^3+O(E_m^5).
\]
证毕。
\end{proof}

\endinput

\paragraph{布尔投影格、Fibonacci 模环分解与有限秩账本障碍}
链上内算子的固定点算术、有限分辨率可见算术的 Fibonacci 模环结构以及有限秩加法账本的同态不可能性，在现有命题链下还可继续压缩为下列七条直接后果。前四条把链上规范投影完全布尔化，写成有限 Euler 乘积，并进一步刻画全部平方自由 \(\gcd/\mathrm{lcm}\) 算术化的唯一正规形及其 prime-register 码长的最优阶；其后三条把可见算术的幂等结构推进为分解刚性与跨分辨率障碍，并以 prime-power 分辨率收束为单链相图。

\begin{theorem}[链上内算子格的布尔刚性与旗标闭式]\label{thm:xi-chain-interior-boolean-flag-closed-form}
\leanverified{paper\_xi\_chain\_interior\_boolean\_flag\_closed\_form}
沿用定理 \ref{thm:killo-chain-interior-godel-gcd-lcm} 的记号，记
\[
\mathcal I_n:=\mathrm{Int}(C_n),
\qquad
\operatorname{rk}(I):=|\mathrm{Fix}(I)|-1.
\]
则：
\begin{enumerate}
  \item 有序格 \((\mathcal I_n,\preceq)\) 与布尔格 \(B_{n-1}\) 规范同构；
  \item
  \[
  |\mathcal I_n|=2^{n-1};
  \]
  \item 秩生成多项式满足
  \[
  \sum_{I\in\mathcal I_n}q^{\operatorname{rk}(I)}=(1+q)^{n-1};
  \]
  \item 若定义
  \[
  Z_{\mathcal I_n}(k)
  :=
  \#\{I_1\preceq\cdots\preceq I_{k-1}:\ I_j\in\mathcal I_n\}
  \qquad(k\ge 1),
  \]
  则
  \[
  Z_{\mathcal I_n}(k)=k^{n-1};
  \]
  \item 对任意 \(I\preceq J\)，区间 \([I,J]\) 仍为布尔格，且
  \[
  \mu_{\mathcal I_n}(I,J)=(-1)^{|\mathrm{Fix}(J)\setminus \mathrm{Fix}(I)|}.
  \]
\end{enumerate}
\end{theorem}
\begin{proof}
由定理 \ref{thm:killo-chain-interior-godel-gcd-lcm}，映射
\[
I\longmapsto \mathrm{Fix}(I)
\]
给出 \(\mathcal I_n\) 与
\[
\mathcal F_n:=\{F\subseteq \{0,1,\dots,n-1\}:0\in F\}
\]
之间的格同构，且 meet/join 分别对应集合交与并。去掉强制出现的元素 \(0\) 之后，\(\mathcal F_n\) 与 \(\mathcal P(\{1,\dots,n-1\})\) 规范同构，故 \((\mathcal I_n,\preceq)\) 即为秩 \(n-1\) 的布尔格。这给出第 \(1\) 条与第 \(2\) 条。

在上述同构下，秩函数 \(\operatorname{rk}(I)\) 恰对应于子集大小，故
\[
\sum_{I\in\mathcal I_n}q^{\operatorname{rk}(I)}
=
\sum_{S\subseteq\{1,\dots,n-1\}}q^{|S|}
=(1+q)^{n-1},
\]
从而第 \(3\) 条成立。

再看 zeta 多项式。一个多重链
\[
I_1\preceq\cdots\preceq I_{k-1}
\]
对应于子集链
\[
S_1\subseteq \cdots\subseteq S_{k-1}\subseteq \{1,\dots,n-1\}.
\]
对每个 \(i\in\{1,\dots,n-1\}\)，只需记录它是在第几步首次进入该链；若始终不进入，则记为第 \(0\) 类。故每个 \(i\) 恰有 \(k\) 种独立选择，从而
\[
Z_{\mathcal I_n}(k)=k^{n-1}.
\]

最后，对任意 \(I\preceq J\)，对应子集满足
\[
\mathrm{Fix}(I)\subseteq \mathrm{Fix}(J).
\]
于是区间 \([I,J]\) 同构于差集 \(\mathrm{Fix}(J)\setminus \mathrm{Fix}(I)\) 上的布尔格，故 Möbius 函数等于相对秩的符号：
\[
\mu_{\mathcal I_n}(I,J)=(-1)^{|\mathrm{Fix}(J)\setminus \mathrm{Fix}(I)|}.
\]
证毕。
\end{proof}

\begin{theorem}[平方自由 Gödel 像的 Euler 因子化与 Möbius 扭曲闭式]\label{thm:xi-chain-interior-godel-euler-mobius-factorization}
沿用定理 \ref{thm:killo-chain-interior-godel-gcd-lcm} 与定理 \ref{thm:xi-chain-interior-boolean-flag-closed-form} 的记号，并固定互异素数 \(q_0,\dots,q_{n-1}\)。则：
\begin{enumerate}
  \item 平方自由像集满足
  \[
  \kappa(\mathcal I_n)
  =
  \left\{
  q_0\prod_{i\in S}q_i:\ S\subseteq\{1,\dots,n-1\}
  \right\};
  \]
  \item 对任意 \(s\in\CC\)，有限 Dirichlet 生成函数满足
  \[
  \sum_{I\in\mathcal I_n}\kappa(I)^{-s}
  =
  q_0^{-s}\prod_{i=1}^{n-1}(1+q_i^{-s});
  \]
  \item 若 \(\widehat 0\) 表示最小内算子，则 Möbius 扭曲生成函数满足
  \[
  \sum_{I\in\mathcal I_n}\mu_{\mathcal I_n}(\widehat 0,I)\,\kappa(I)^{-s}
  =
  q_0^{-s}\prod_{i=1}^{n-1}(1-q_i^{-s}).
  \]
\end{enumerate}
\end{theorem}
\begin{proof}
由定理 \ref{thm:killo-chain-interior-godel-gcd-lcm}，任意 \(I\in\mathcal I_n\) 唯一对应于一个固定点集
\[
\mathrm{Fix}(I)=\{0\}\cup S,
\qquad
S\subseteq\{1,\dots,n-1\},
\]
并且
\[
\kappa(I)=q_0\prod_{i\in S}q_i.
\]
这立刻给出第 \(1\) 条。

于是
\[
\sum_{I\in\mathcal I_n}\kappa(I)^{-s}
=
\sum_{S\subseteq\{1,\dots,n-1\}}
\left(q_0\prod_{i\in S}q_i\right)^{-s}
=
q_0^{-s}\prod_{i=1}^{n-1}(1+q_i^{-s}),
\]
得到第 \(2\) 条。

再由定理 \ref{thm:xi-chain-interior-boolean-flag-closed-form}，对满足 \(\mathrm{Fix}(I)=\{0\}\cup S\) 的内算子有
\[
\mu_{\mathcal I_n}(\widehat 0,I)=(-1)^{|S|}.
\]
故
\[
\sum_{I\in\mathcal I_n}\mu_{\mathcal I_n}(\widehat 0,I)\,\kappa(I)^{-s}
=
\sum_{S\subseteq\{1,\dots,n-1\}}
(-1)^{|S|}
\left(q_0\prod_{i\in S}q_i\right)^{-s}
=
q_0^{-s}\prod_{i=1}^{n-1}(1-q_i^{-s}),
\]
即得第 \(3\) 条。证毕。
\end{proof}

\begin{corollary}[链上内算子格的 Dirichlet--Möbius 封闭与特征多项式]\label{cor:xi-chain-interior-dirichlet-mobius-characteristic-closure}
沿用定理 \ref{thm:xi-chain-interior-boolean-flag-closed-form} 与定理 \ref{thm:xi-chain-interior-godel-euler-mobius-factorization} 的记号。记
\[
\Xi_n(s):=\sum_{I\in\mathcal I_n}\kappa(I)^{-s},
\qquad
\mathfrak M_n(s):=\sum_{I\in\mathcal I_n}\mu_{\mathcal I_n}(\widehat 0,I)\,\kappa(I)^{-s}.
\]
则有
\[
\Xi_n(s)=q_0^{-s}\prod_{i=1}^{n-1}(1+q_i^{-s}),
\qquad
\mathfrak M_n(s)=q_0^{-s}\prod_{i=1}^{n-1}(1-q_i^{-s}).
\]
更一般地，对任意 \(I\preceq J\) 都有
\[
\mu_{\mathcal I_n}(I,J)
=
(-1)^{|\mathrm{Fix}(J)\setminus \mathrm{Fix}(I)|}
=
\mu_{\mathrm{arith}}\!\left(\frac{\kappa(J)}{\kappa(I)}\right),
\]
其中 \(\mu_{\mathrm{arith}}\) 为经典算术 Möbius 函数。特别地，若定义特征多项式
\[
\chi_{\mathcal I_n}(x)
:=
\sum_{I\in\mathcal I_n}\mu_{\mathcal I_n}(\widehat 0,I)\,
x^{\,n-1-\operatorname{rk}(I)},
\]
则
\[
\chi_{\mathcal I_n}(x)=(x-1)^{n-1}.
\]
\end{corollary}
\begin{proof}
前两式正是定理 \ref{thm:xi-chain-interior-godel-euler-mobius-factorization} 的第 \(2\)、\(3\) 条。
对任意 \(I\preceq J\)，定理 \ref{thm:xi-chain-interior-boolean-flag-closed-form} 已给出
\[
\mu_{\mathcal I_n}(I,J)=(-1)^{|\mathrm{Fix}(J)\setminus \mathrm{Fix}(I)|}.
\]
另一方面，由定理 \ref{thm:killo-chain-interior-godel-gcd-lcm}，
\[
\frac{\kappa(J)}{\kappa(I)}
=
\prod_{i\in \mathrm{Fix}(J)\setminus \mathrm{Fix}(I)}q_i
\]
为平方自由整数，故其算术 Möbius 函数恰等于相同符号，得到第三式。

最后，由 \(\widehat 0\) 对应的固定点集仅为 \(\{0\}\)，对每个满足 \(\mathrm{Fix}(I)=\{0\}\cup S\) 的内算子有
\[
\mu_{\mathcal I_n}(\widehat 0,I)=(-1)^{|S|},
\qquad
\operatorname{rk}(I)=|S|.
\]
于是
\[
\chi_{\mathcal I_n}(x)
=
\sum_{S\subseteq\{1,\dots,n-1\}}(-1)^{|S|}x^{n-1-|S|}
=
(x-1)^{n-1},
\]
其中最后一步是二项式展开。证毕。
\end{proof}

\begin{theorem}[链上规范投影的平方自由 \texorpdfstring{$\gcd/\mathrm{lcm}$}{gcd/lcm} 算术化具有唯一正规形]\label{thm:xi-chain-interior-squarefree-gcd-lcm-normalform}
\leanverified{paper\_xi\_chain\_interior\_squarefree\_gcd\_lcm\_normalform}
沿用定理 \ref{thm:killo-chain-interior-godel-gcd-lcm} 的记号，记
\[
\mathcal I_n:=\mathrm{Int}(C_n).
\]
设映射
\[
\eta:\mathcal I_n\to \NN_{\ge 1}
\]
满足：
\begin{enumerate}
  \item 对每个 \(I\in\mathcal I_n\)，\(\eta(I)\) 都是平方自由整数；
  \item \(\eta\) 是单射；
  \item 对任意 \(I,J\in\mathcal I_n\)，都有
  \[
  \eta(I\wedge J)=\gcd(\eta(I),\eta(J)),
  \qquad
  \eta(I\vee J)=\mathrm{lcm}(\eta(I),\eta(J)).
  \]
\end{enumerate}
则存在唯一确定的平方自由整数 \(b\) 与两两互素的平方自由整数
\[
r_1,\dots,r_{n-1}>1
\]
使得对每个满足 \(0\in F\subseteq\{0,\dots,n-1\}\) 的集合 \(F\)，都有
\[
\boxed{\;
\eta(I_F)=b\prod_{i\in F\setminus\{0\}}r_i,
\;}
\]
其中空积按 \(1\) 解释。特别地，任何把链上规范投影算术化为平方自由整数整除格的 \(\gcd/\mathrm{lcm}\)-语义，都只可能与定理 \ref{thm:killo-chain-interior-godel-gcd-lcm} 中的规范编码相差一个公共底因子与一组两两互素坐标标签的重命名。
\end{theorem}
\begin{proof}
记
\[
\widehat 0:=I_{\{0\}},
\qquad
A_i:=I_{\{0,i\}}
\qquad (1\le i\le n-1).
\]
由定理 \ref{thm:killo-chain-interior-godel-gcd-lcm}，\(\mathcal I_n\) 与包含 \(0\) 的子集布尔格同构，因此 \(A_1,\dots,A_{n-1}\) 正是 \(\mathcal I_n\) 的全部原子，并且对 \(i\neq j\) 有
\[
A_i\wedge A_j=\widehat 0.
\]

令
\[
b:=\eta(\widehat 0).
\]
由假设，\(b\) 是平方自由整数。又因为
\[
\gcd(\eta(A_i),\eta(A_j))
=
\eta(A_i\wedge A_j)
=
\eta(\widehat 0)
=
b
\qquad(i\neq j),
\]
可知 \(b\mid \eta(A_i)\) 对每个 \(i\) 都成立。于是可写
\[
\eta(A_i)=b\,r_i
\qquad (1\le i\le n-1).
\]
由于 \(\eta(A_i)\) 与 \(b\) 都是平方自由整数，故 \(r_i\) 也是平方自由整数，且自动满足
\[
\gcd(b,r_i)=1.
\]
若某个 \(r_i=1\)，则
\[
\eta(A_i)=b=\eta(\widehat 0),
\]
与 \(\eta\) 的单射性矛盾，因此 \(r_i>1\)。

再由
\[
\gcd(\eta(A_i),\eta(A_j))
=
\gcd(br_i,br_j)
=
b\,\gcd(r_i,r_j)
=
b
\]
得到
\[
\gcd(r_i,r_j)=1
\qquad(i\neq j),
\]
故 \(r_1,\dots,r_{n-1}\) 两两互素。

现在任取 \(F\subseteq\{0,\dots,n-1\}\) 且 \(0\in F\)。由同一定理的布尔格描述，
\[
I_F=\bigvee_{i\in F\setminus\{0\}}A_i,
\]
其中当 \(F=\{0\}\) 时右端为空 join，其值即为 \(\widehat 0\)。于是由 \(\eta\) 对 join 的保持性可得
\[
\eta(I_F)
=
\mathrm{lcm}\bigl(\eta(A_i):i\in F\setminus\{0\}\bigr).
\]
而 \(b,r_1,\dots,r_{n-1}\) 全部平方自由，且 \(b\) 与各 \(r_i\) 互素、各 \(r_i\) 彼此互素，故上式化为
\[
\eta(I_F)
=
\mathrm{lcm}\bigl(br_i:i\in F\setminus\{0\}\bigr)
=
b\prod_{i\in F\setminus\{0\}}r_i.
\]
这就证明了所述正规形。

最后，\(b=\eta(\widehat 0)\) 且 \(r_i=\eta(A_i)/b\) 由 \(\eta\) 唯一决定，故该正规形唯一。证毕。
\end{proof}

\begin{corollary}[链上 canonical arithmeticization 在 prime-register 码长意义下达到最优阶]\label{cor:xi-chain-interior-canonical-arithmeticization-optimal-order}
沿用定理 \ref{thm:xi-chain-interior-squarefree-gcd-lcm-normalform} 的记号。记 \(p_j\) 为第 \(j\) 个素数，并把 \(\{0,1\}^{n-1}\) 与 \(\mathcal I_n\) 通过对应
\[
(\varepsilon_1,\dots,\varepsilon_{n-1})
\longleftrightarrow
I_{\{0\}\cup\{i:\varepsilon_i=1\}}
\]
识别。定义规范编码
\[
E_n^{\mathrm{can}}(\varepsilon_1,\dots,\varepsilon_{n-1})
:=
p_1\prod_{i=1}^{n-1}p_{i+1}^{\varepsilon_i}.
\]
则有
\[
\boxed{\;
\max_{\varepsilon\in\{0,1\}^{n-1}}
\log E_n^{\mathrm{can}}(\varepsilon)
=
\sum_{j=1}^{n}\log p_j
=
\vartheta(p_n)
=
(1+o(1))\,n\log n.\;}
\]
另一方面，对任意满足定理 \ref{thm:xi-chain-interior-squarefree-gcd-lcm-normalform} 三条假设的算术化
\[
\eta:\mathcal I_n\to \NN_{\ge 1},
\]
都有
\[
\boxed{\;
\max_{I\in\mathcal I_n}\log \eta(I)=\Omega(n\log n).\;}
\]
因此，在保持 \(\wedge/\vee\) 与 \(\gcd/\mathrm{lcm}\) 对应的平方自由 prime-register 算术化之中，规范编码的最坏码长阶恰为
\[
\Theta(n\log n),
\]
从而达到最优阶。并且，这一下界与定理 \ref{thm:xi-godel-prime-bit-complexity-lower-bound} 中对一般 prime-valuation 可寻址 Gödel 整数化给出的 \(T\log T\) 障碍完全一致。
\end{corollary}
\begin{proof}
对规范编码，只需取全 \(1\) 向量，便有
\[
\max_{\varepsilon\in\{0,1\}^{n-1}}
\log E_n^{\mathrm{can}}(\varepsilon)
=
\sum_{j=1}^{n}\log p_j
=
\vartheta(p_n).
\]
由素数定理的标准推论
\[
\vartheta(p_n)\sim p_n\sim n\log n,
\]
即得所示上界渐近式。

再证下界。任取满足三条假设的 \(\eta\)。由定理 \ref{thm:xi-chain-interior-squarefree-gcd-lcm-normalform}，存在平方自由整数 \(b\) 与两两互素的平方自由整数 \(r_1,\dots,r_{n-1}>1\)，使
\[
\eta(I_F)=b\prod_{i\in F\setminus\{0\}}r_i.
\]
对每个 \(i\)，任选一个素因子 \(s_i\mid r_i\)。由于 \(r_i\) 两两互素，\(s_1,\dots,s_{n-1}\) 彼此不同。于是
\[
\max_{I\in\mathcal I_n}\eta(I)
\ge
\eta(I_{\{0,1,\dots,n-1\}})
=
b\prod_{i=1}^{n-1}r_i
\ge
\prod_{i=1}^{n-1}s_i.
\]
而任意 \(n-1\) 个互异素数的乘积都不小于前 \(n-1\) 个素数的乘积，故
\[
\max_{I\in\mathcal I_n}\eta(I)
\ge
\prod_{j=1}^{n-1}p_j.
\]
取对数即得
\[
\max_{I\in\mathcal I_n}\log \eta(I)
\ge
\sum_{j=1}^{n-1}\log p_j
=
\vartheta(p_{n-1})
=
\Omega(n\log n).
\]
于是规范编码给出的 \(O(n\log n)\) 上界与任何同类算术化都必须满足的 \(\Omega(n\log n)\) 下界同阶，故其最坏码长达到最优阶。最后一条只是在更一般的 prime-valuation 可寻址 Gödel 整数化语义下援引定理 \ref{thm:xi-godel-prime-bit-complexity-lower-bound} 的同型下界。证毕。
\end{proof}

\begin{theorem}[分辨率商半环的分解刚性由 \texorpdfstring{$\omega(F_{m+2})$}{omega(Fm+2)} 完全控制]\label{thm:xi-visible-arithmetic-omega-controlled-splitting}
沿用定理 \ref{thm:fold-congruence-mod-semirings} 的记号，记
\[
R_m^{\mathrm{vis}}
:=
\Omega_m/\!\equiv_m
\cong
\ZZ/(N_m\ZZ),
\qquad
N_m:=F_{m+2},
\qquad
r_m:=\omega(N_m).
\]
则：
\begin{enumerate}
  \item \(R_m^{\mathrm{vis}}\) 的幂等元总数恰为
  \[
  2^{r_m};
  \]
  \item \(R_m^{\mathrm{vis}}\) 存在非平凡直积分解，当且仅当 \(r_m\ge 2\)；
  \item \(R_m^{\mathrm{vis}}\) 直接不可分解，当且仅当 \(N_m\) 为素数幂；
  \item 若 \(N_m=p\) 为素数，则
  \[
  R_m^{\mathrm{vis}}\cong \FF_p,
  \]
  从而其幂等元仅有 \(0,1\)。
\end{enumerate}
\end{theorem}
\begin{proof}
将 \(N_m\) 分解为
\[
N_m=\prod_{j=1}^{r_m}p_j^{a_j},
\]
其中 \(p_j\) 两两不同。由中国剩余定理，
\[
R_m^{\mathrm{vis}}
\cong
\ZZ/(N_m\ZZ)
\cong
\prod_{j=1}^{r_m}\ZZ/(p_j^{a_j}\ZZ).
\]
第 \(1\) 条正是定理 \ref{thm:xi-visible-arithmetic-unit-idempotent-count} 的幂等元计数公式。

对交换环 \(A\)，非平凡直积分解
\[
A\simeq A_1\times A_2
\qquad
(A_1,A_2\neq 0)
\]
等价于存在非平凡幂等元：若分解存在，则 \(e=(1,0)\) 满足 \(e^2=e\) 且 \(e\neq 0,1\)；反之，若 \(e^2=e\) 且 \(e\neq 0,1\)，则
\[
A\cong eA\times (1-e)A.
\]
因此第 \(2\) 条等价于“存在非平凡幂等元”。由第 \(1\) 条，这恰当且仅当 \(2^{r_m}>2\)，亦即 \(r_m\ge 2\)。

第 \(3\) 条随即成立，因为
\[
r_m=1
\quad\Longleftrightarrow\quad
N_m \text{ 为素数幂}.
\]
最后，若 \(N_m=p\) 为素数，则
\[
R_m^{\mathrm{vis}}\cong \ZZ/(p\ZZ)=\FF_p,
\]
故其仅有平凡幂等元 \(0,1\)。这也与推论 \ref{cor:field-phase-fib-prime} 一致。证毕。
\end{proof}

\begin{theorem}[跨分辨率相容的有限秩加法账本不能同态化全部可见乘法]\label{thm:xi-visible-arithmetic-no-compatible-finite-rank-ledger}
设 \(L\) 为有限生成交换群。假设对每个 \(m\ge 1\) 都存在映射
\[
\Phi_m:\ZZ/(F_{m+2}\ZZ)\longrightarrow \ZZ\oplus L
\]
满足：
\begin{enumerate}
  \item \(\Phi_m\) 在乘法幺半群上单射；
  \item 对任意 \(\overline a,\overline b\in \ZZ/(F_{m+2}\ZZ)\)，都有
  \[
  \Phi_m(\overline a\,\overline b)=\Phi_m(\overline a)+\Phi_m(\overline b);
  \]
  \item 对每个整数 \(n\ge 1\)，存在 \(m_0(n)\) 使得当 \(m\ge m_0(n)\) 时，\(\Phi_m(\overline n)\) 与 \(m\) 无关，其中 \(\overline n\) 表示 \(n\) 在 \(\ZZ/(F_{m+2}\ZZ)\) 中的剩余类。
\end{enumerate}
则这样的映射族不存在。
\end{theorem}
\begin{proof}
由第 \(3\) 条，可对每个 \(n\ge 1\) 定义稳定值
\[
\Phi(n):=\Phi_m(\overline n)
\qquad
(m\ge m_0(n)).
\]
这给出映射
\[
\Phi:(\NN_{\ge 1},\times)\longrightarrow (\ZZ\oplus L,+).
\]

先证其单射。若 \(a\neq b\)，则由 \(F_{m+2}\to\infty\)，可取 \(m\) 足够大，使
\[
m\ge m_0(a),\qquad m\ge m_0(b),\qquad F_{m+2}>\max\{a,b\}.
\]
于是 \(\overline a\neq \overline b\) 于 \(\ZZ/(F_{m+2}\ZZ)\) 中成立。由第 \(1\) 条，
\[
\Phi(a)=\Phi_m(\overline a)\neq \Phi_m(\overline b)=\Phi(b),
\]
故 \(\Phi\) 为单射。

再证其乘法--加法同态性。对任意 \(a,b\ge 1\)，取
\[
m\ge \max\{m_0(a),m_0(b),m_0(ab)\}.
\]
则由第 \(2\) 条，
\[
\Phi(ab)
=
\Phi_m(\overline{ab})
=
\Phi_m(\overline a\,\overline b)
=
\Phi_m(\overline a)+\Phi_m(\overline b)
=
\Phi(a)+\Phi(b).
\]
因此 \(\Phi:(\NN_{\ge 1},\times)\to (\ZZ\oplus L,+)\) 是单射幺半群同态。

另一方面，由 \(L\) 有限生成，可知交换群 \(\ZZ\oplus L\) 作为加法幺半群亦为有限生成且可消去。于是定理 \ref{thm:killo-prime-freedom-non-finitizability} 排除任何从正整数乘法幺半群到 \((\ZZ\oplus L,+)\) 的单射幺半群同态。
这与上面构造出的 \(\Phi\) 矛盾。故所述映射族不存在。证毕。
\end{proof}

\begin{corollary}[素数幂分辨率上的商格单链化]\label{cor:xi-visible-arithmetic-primepower-quotient-chain}
沿用定理 \ref{thm:xi-visible-arithmetic-omega-controlled-splitting} 的记号。若
\[
N_m=F_{m+2}=p^a
\]
为某个素数幂，则 \(R_m^{\mathrm{vis}}\cong \ZZ/(p^a\ZZ)\) 的全部双边理想恰为
\[
(0)=(p^a)\subset (p^{a-1})\subset \cdots \subset (p)\subset (1),
\]
并且由定理 \ref{thm:xi-fold-congruence-quotient-divisor-lattice}，\(R_m^{\mathrm{vis}}\) 的全部半环商构成长度 \(a+1\) 的线性链。特别地，此时不存在彼此横向独立的算术粗化方向。
\end{corollary}
\begin{proof}
由定理 \ref{thm:xi-visible-arithmetic-omega-controlled-splitting}，当 \(N_m=p^a\) 时，
\[
R_m^{\mathrm{vis}}\cong \ZZ/(p^a\ZZ).
\]
而 \(p^a\) 的全部正约数正好是
\[
1,p,\dots,p^a,
\]
它们在整除偏序下形成一条线性链。定理 \ref{thm:xi-fold-congruence-quotient-divisor-lattice} 已将 \(R_m^{\mathrm{vis}}\) 的商格与 \(N_m\) 的约数格规范等价起来，故全部半环商构成长度 \(a+1\) 的线性链。

等价地，\(\ZZ/(p^a\ZZ)\) 的全部理想就是
\[
(p^j)/(p^a)
\qquad
(0\le j\le a),
\]
从而得到所示理想链。证毕。
\end{proof}

\paragraph{素数寄存器幂等锥、有限素窗口半圆维与二原子信息几何的闭合描述}
在素数寄存器估值算术化、链上规范投影的布尔化、有限素数支撑乘法单子的半圆维分离以及 bin-fold 相对均匀基线的普适 \(f\)-散度塌缩已经建立之后，还可进一步抽取出下列七条直接后果。

\leanverified{paper\_xi\_prime\_register\_idempotent\_cone\_rank\_dirichlet}
\begin{theorem}[素数寄存器幂等锥的秩分层、生成多项式与 Dirichlet 闭式]\label{thm:xi-prime-register-idempotent-cone-rank-dirichlet}
沿用定理 \ref{thm:killo-finite-semigroup-prime-valuation-arithmetization} 与定理 \ref{thm:xi-prime-register-idempotent-exact-count} 的记号，并记
\[
S_n:=\left\{\prod_{i=0}^{n-1}q_i^{a_i}: a_i\in\{0,\dots,n-1\}\right\},
\]
\[
E_n:=\operatorname{Idem}(S_n),
\qquad
E_{n,k}:=\operatorname{Idem}_k(S_n),
\qquad
\operatorname{rk}(N):=|\operatorname{Im}(f_N)|.
\]
则对每个 \(1\le k\le n\) 都有
\[
|E_{n,k}|=\binom{n}{k}k^{n-k},
\qquad
|E_n|=\sum_{k=1}^{n}\binom{n}{k}k^{n-k}.
\]
其秩生成多项式满足
\[
\mathcal E_n(y):=\sum_{N\in E_n}y^{\operatorname{rk}(N)}
=
\sum_{k=1}^{n}\binom{n}{k}k^{n-k}y^k.
\]
此外，有限 Dirichlet 生成式具有显式闭式
\[
\mathfrak E_n(s):=\sum_{N\in E_n}N^{-s}
=
\sum_{\varnothing\neq F\subseteq\{0,\dots,n-1\}}
\left(\prod_{j\in F}q_j^{-sj}\right)
\prod_{i\notin F}\left(\sum_{j\in F}q_i^{-sj}\right).
\]
\end{theorem}
\begin{proof}
由推论 \ref{cor:killo-projection-idempotent-prime-register}，\(N\in E_n\) 当且仅当 \(f_N\) 为幂等映射，且此时
\[
\operatorname{Im}(f_N)=\operatorname{Fix}(f_N).
\]
因此，定理 \ref{thm:xi-prime-register-idempotent-exact-count} 已直接给出
\[
|E_{n,k}|=\binom{n}{k}k^{n-k},
\qquad
|E_n|=\sum_{k=1}^{n}\binom{n}{k}k^{n-k}.
\]
将秩 \(k\) 的计数按 \(y^k\) 加权求和，即得 \(\mathcal E_n(y)\) 的闭式。

再求 Dirichlet 生成式。固定非空子集 \(F\subseteq\{0,\dots,n-1\}\)。由
\[
\operatorname{Im}(f_N)=\operatorname{Fix}(f_N)=F
\]
可知，此类幂等元恰由下列独立条件刻画：
\[
v_{q_j}(N)=j\qquad (j\in F),
\qquad
v_{q_i}(N)\in F\qquad (i\notin F).
\]
因此，对 \(j\in F\) 必有贡献因子 \(q_j^{-sj}\)；而对 \(i\notin F\)，其指数可在 \(F\) 中任取，故贡献
\[
\sum_{j\in F}q_i^{-sj}.
\]
各坐标彼此独立，相乘后即得到固定像集 \(F\) 的全部贡献
\[
\left(\prod_{j\in F}q_j^{-sj}\right)
\prod_{i\notin F}\left(\sum_{j\in F}q_i^{-sj}\right).
\]
再对全部非空 \(F\) 求和，即得所示 Dirichlet 闭式。证毕。
\end{proof}

\begin{theorem}[单调幂等骨架的 Boolean 交格与指数稀疏性]\label{thm:xi-monotone-idempotent-skeleton-gcd-lcm-sparsity}
沿用定理 \ref{thm:killo-chain-interior-godel-gcd-lcm} 的记号，记
\[
\mathcal I_n:=\mathrm{Int}(C_n),
\qquad
\mathcal I_{n,k}:=\{I\in\mathcal I_n:\ |\mathrm{Fix}(I)|=k\}
\qquad(1\le k\le n).
\]
则对任意满足 \(0\in F\cap G\subseteq\{0,\dots,n-1\}\) 的子集 \(F,G\)，有
\[
I_F\wedge I_G=I_{F\cap G},
\qquad
I_F\vee I_G=I_{F\cup G},
\]
从而
\[
\kappa(I_F\wedge I_G)=\gcd(\kappa(I_F),\kappa(I_G)),
\qquad
\kappa(I_F\vee I_G)=\mathrm{lcm}(\kappa(I_F),\kappa(I_G)).
\]
特别地，\(\mathcal I_n\) 在 meet 运算下构成交换幂等半群，并经 \(\kappa\) 与平方自由整数上的 \(\gcd\)-半群同构。

此外，
\[
|\mathcal I_{n,k}|=\binom{n-1}{k-1},
\]
并且其在全部秩 \(k\) 幂等元中所占比例满足
\[
\frac{|\mathcal I_{n,k}|}{|E_{n,k}|}
=
\frac{\binom{n-1}{k-1}}{\binom{n}{k}k^{n-k}}
=
\frac{k}{n}\,k^{-(n-k)}.
\]
因而只要 \(1\le k\le n-1\)，该比例便随 \(n-k\) 指数衰减。
\leanverified{paper\_xi\_monotone\_idempotent\_skeleton\_gcd\_lcm\_sparsity}
\end{theorem}
\begin{proof}
由定理 \ref{thm:killo-chain-interior-godel-gcd-lcm}，\(\mathcal I_n\) 与
\[
\mathcal F_n:=\{F\subseteq[n]:0\in F\}
\]
通过 \(F\mapsto I_F\) 一一对应，且格运算严格对应于集合交并。因而
\[
I_F\wedge I_G=I_{F\cap G},
\qquad
I_F\vee I_G=I_{F\cup G}.
\]
再由同一定理的算术化公式，
\[
\kappa(I_F\wedge I_G)=\gcd(\kappa(I_F),\kappa(I_G)),
\qquad
\kappa(I_F\vee I_G)=\mathrm{lcm}(\kappa(I_F),\kappa(I_G)),
\]
即得第一部分。特别地，meet 的交换律与幂等律使 \(\mathcal I_n\) 成为交换幂等半群。

对计数部分，只需从 \(\{1,\dots,n-1\}\) 中选出 \(k-1\) 个元素并与强制出现的 \(0\) 合并为固定点集 \(F\)，故
\[
|\mathcal I_{n,k}|=\binom{n-1}{k-1}.
\]
再与定理 \ref{thm:xi-prime-register-idempotent-cone-rank-dirichlet} 中
\[
|E_{n,k}|=\binom{n}{k}k^{n-k}
\]
相除，便得到所示比例公式。最后一条由 \(k\ge 1\) 时 \(k^{-(n-k)}\) 的指数衰减立即成立。证毕。
\end{proof}

\begin{corollary}[规范投影骨架在全部素数寄存器幂等元中的超指数稀薄性]\label{cor:xi-monotone-idempotent-skeleton-superexponential-sparsity}
沿用定理 \ref{thm:xi-chain-interior-boolean-flag-closed-form} 与定理 \ref{thm:xi-prime-register-idempotent-cone-rank-dirichlet} 的记号。对每个 \(n\ge 2\)，记
\[
c_n:=|\mathcal I_n|=2^{n-1},
\qquad
e_n:=|E_n|=\sum_{k=1}^{n}\binom{n}{k}k^{n-k}.
\]
则
\[
\frac{c_n}{e_n}
\le
\frac{n+1}{2\,\lfloor n/2\rfloor^{\,n-\lfloor n/2\rfloor}},
\]
并且
\[
\frac{c_n}{e_n}
=
\exp\!\Bigl(-\frac12 n\log n+O(n)\Bigr)
\longrightarrow 0
\qquad (n\to\infty).
\]
因而链上规范投影在全部素数寄存器幂等元中仅占超指数稀薄的比例。
\leanverified{paper\_xi\_monotone\_idempotent\_skeleton\_superexponential\_sparsity}
\end{corollary}
\begin{proof}
记
\[
m:=\lfloor n/2\rfloor.
\]
由定理 \ref{thm:xi-prime-register-idempotent-cone-rank-dirichlet}，
\[
e_n=\sum_{k=1}^{n}\binom{n}{k}k^{n-k}
\ge
\binom{n}{m}m^{n-m}.
\]
另一方面，由
\[
2^n=\sum_{k=0}^{n}\binom{n}{k}
\]
可知中心二项式系数满足
\[
\binom{n}{m}\ge \frac{2^n}{n+1}.
\]
再结合
\[
c_n=2^{n-1},
\]
得到
\[
\frac{c_n}{e_n}
\le
\frac{2^{n-1}}{(2^n/(n+1))\,m^{n-m}}
=
\frac{n+1}{2\,m^{n-m}},
\]
即得第一条不等式。

由于 \(m=\lfloor n/2\rfloor\sim n/2\)，于是
\[
\log\frac{c_n}{e_n}
\le
\log(n+1)-\log 2-(n-m)\log m
=
-\frac12 n\log n+O(n).
\]
指数化后即得
\[
\frac{c_n}{e_n}
=
\exp\!\Bigl(-\frac12 n\log n+O(n)\Bigr)\to 0.
\]
证毕。
\end{proof}

\begin{theorem}[有限素窗口同态半圆维的严格线性差距]\label{thm:xi-finite-prime-window-halfcircle-linear-gap}
取互异素数 \(p_1,\dots,p_r\)，并记
\[
M_r:=\langle p_1,\dots,p_r\rangle
=
\left\{p_1^{a_1}\cdots p_r^{a_r}: a_1,\dots,a_r\in\ZZ_{\ge 0}\right\}.
\]
则
\[
\cdim^{\mathrm{hom}}_{+}(M_r)=\frac{r}{2},
\qquad
\cdim^{\mathrm{impl}}(M_r)=\frac12.
\]
特别地，
\[
\cdim^{\mathrm{hom}}_{+}(M_r)-\cdim^{\mathrm{impl}}(M_r)=\frac{r-1}{2},
\]
故有限素窗口上的结构层与实现层半圆维差距已经呈严格线性增长。
\leanverified{paper\_xi\_finite\_prime\_window\_halfcircle\_linear\_gap}
\end{theorem}
\begin{proof}
令
\[
S_r:=\{p_1,\dots,p_r\}\subset\mathbb P.
\]
则
\[
M_r=\NN_{S_r}^\times.
\]
所述等式正是定理 \ref{thm:xi-finite-prime-support-multiplicative-half-circle-dimension} 在 \(S=S_r\) 时的直接转写。证毕。
\end{proof}

\begin{theorem}[单一 \texorpdfstring{$\chi^2$}{chi2} 极限常数完备编码全部 Csisz\'ar 几何]\label{thm:xi-chi2-complete-csiszar-geometry}
沿用命题 \ref{prop:fold-bin-f-divergence-collapse} 的记号，并记
\[
f_2(t):=(t-1)^2,
\qquad
\mathcal C_2:=D_{f_2}^\infty=\lim_{m\to\infty}\chi^2(\mu_m\Vert u_m).
\]
则
\[
\varphi=\bigl(5(\mathcal C_2+1)-1\bigr)^{1/3}.
\]
若进一步置
\[
\xi:=\bigl(5(\mathcal C_2+1)-1\bigr)^{1/3},
\]
则对任意凸函数 \(f:(0,\infty)\to\RR\) 满足 \(f(1)=0\)，其 bin-fold 极限 \(f\)-散度都满足
\[
D_f^\infty
=
\frac{1}{\xi}\,f\!\left(\frac{\xi^2}{\sqrt5}\right)
+\frac{1}{\xi^2}\,f\!\left(\frac{\xi}{\sqrt5}\right).
\]
因此，bin-fold 相对均匀基线的全部渐近 Csisz\'ar 信息几何并非无限维数据，而由单一标量 \(\mathcal C_2\) 完整决定。
\end{theorem}
\begin{proof}
由命题 \ref{prop:fold-bin-f-divergence-collapse} 取 \(f=f_2\)，得到
\[
\mathcal C_2
=
\frac{1}{\varphi}\left(\frac{\varphi^2}{\sqrt5}\right)^2
+\frac{1}{\varphi^2}\left(\frac{\varphi}{\sqrt5}\right)^2
-1
=
\frac{\varphi^3+1}{5}-1.
\]
整理即得
\[
\varphi^3=5(\mathcal C_2+1)-1,
\]
从而
\[
\varphi=\bigl(5(\mathcal C_2+1)-1\bigr)^{1/3}.
\]
将
\[
\xi=\bigl(5(\mathcal C_2+1)-1\bigr)^{1/3}=\varphi
\]
代回命题 \ref{prop:fold-bin-f-divergence-collapse} 的两点闭式，便得到任意 \(f\) 的极限表达式。故全部 \(D_f^\infty\) 都由 \(\mathcal C_2\) 唯一决定。证毕。
\end{proof}

\begin{theorem}[渐近 \texorpdfstring{$f$}{f}-散度泛函的二原子刚性]\label{thm:xi-asymptotic-f-divergence-functional-two-atom-rigidity}
记
\[
\nu_\varphi
:=
\frac{1}{\varphi}\,\delta_{\varphi^2/\sqrt5}
+\frac{1}{\varphi^2}\,\delta_{\varphi/\sqrt5}.
\]
设 \(K\subset(0,\infty)\) 为包含
\(
\{\varphi/\sqrt5,\varphi^2/\sqrt5\}
\)
的紧区间，\(\nu\) 为 \(K\) 上的 Borel 概率测度。若对一切 \(\alpha>1\) 都有
\[
\int_K \bigl(t^\alpha-1-\alpha(t-1)\bigr)\,\dd\nu(t)
=
D_{f_\alpha}^\infty
=
\int_K \bigl(t^\alpha-1-\alpha(t-1)\bigr)\,\dd\nu_\varphi(t),
\]
其中
\[
f_\alpha(t):=t^\alpha-1-\alpha(t-1),
\]
则必有
\[
\nu=\nu_\varphi.
\]
换言之，bin-fold 的渐近 \(f\)-散度泛函在测度论层面唯一对应于该二原子模型。
\end{theorem}
\begin{proof}
设
\[
\sigma:=\nu-\nu_\varphi.
\]
对 \(z\in\CC\) 定义
\[
G(z):=\int_K \bigl(t^z-1-z(t-1)\bigr)\,\dd\sigma(t),
\qquad
t^z:=e^{z\log t}.
\]
由于 \(K\subset(0,\infty)\) 为紧集，\(\log t\) 在 \(K\) 上有界，故 \(G\) 是整函数。由假设，对每个实数 \(\alpha>1\) 都有
\[
G(\alpha)=0.
\]
零点集含有聚点，因此由恒等定理可知
\[
G\equiv 0 \qquad \text{于 } \CC.
\]
于是对一切 \(z\in\CC\) 都有
\[
0=G''(z)=\int_K (\log t)^2 t^z\,\dd\sigma(t).
\]
定义有限符号测度
\[
\dd\lambda(t):=(\log t)^2 t^2\,\dd\sigma(t).
\]
则对任意整数 \(n\ge 0\)，
\[
\int_K t^n\,\dd\lambda(t)=G''(n+2)=0.
\]
因此 \(\lambda\) 对 \(K\) 上所有多项式的积分皆为零。由 Stone--Weierstrass 定理，多项式在 \(C(K)\) 中稠密，故
\[
\lambda=0.
\]
而对 \(t\neq 1\) 有
\[
(\log t)^2 t^2>0,
\]
故 \(\sigma\) 只能支撑在 \(\{1\}\) 上。另一方面，\(\nu\) 与 \(\nu_\varphi\) 都是概率测度，故
\[
\sigma(K)=0.
\]
因此 \(\sigma\) 不可能在 \(\{1\}\) 上留下任何非零质量，只能有
\[
\sigma=0.
\]
于是 \(\nu=\nu_\varphi\)。证毕。
\end{proof}

\begin{proposition}[bin-fold 渐近 \texorpdfstring{$f$}{f}-散度塌缩为唯一二点实验]\label{prop:xi-foldbin-f-divergence-two-point-experiment}
沿用命题 \ref{prop:fold-bin-f-divergence-collapse} 的记号。定义 \(\{0,1\}\) 上的两点实验
\[
Q_\infty:=\left(\frac1\varphi,\frac1{\varphi^2}\right),
\qquad
P_\infty:=\left(\frac{\varphi}{\sqrt5},\frac{1}{\varphi\sqrt5}\right).
\]
则对任意凸函数 \(f:(0,\infty)\to\RR\) 满足 \(f(1)=0\)，都有
\[
\boxed{\;
\lim_{m\to\infty}D_f(\mu_m\Vert u_m)
=
D_f(P_\infty\Vert Q_\infty),\;}
\]
并且
\[
\boxed{\;
D_f(\mu_m\Vert u_m)
=
D_f(P_\infty\Vert Q_\infty)
+O_f\bigl((\varphi/2)^m\bigr).\;}
\]
其中
\[
D_f(P_\infty\Vert Q_\infty)
:=
\frac1\varphi\,f\!\left(\frac{P_{\infty,0}}{Q_{\infty,0}}\right)
+\frac1{\varphi^2}\,f\!\left(\frac{P_{\infty,1}}{Q_{\infty,1}}\right).
\]
换言之，\((\mu_m,u_m)\) 的全部渐近 Csisz\'ar 几何完全塌缩到单一二点实验 \((P_\infty,Q_\infty)\)。
\end{proposition}
\begin{proof}
直接计算得到
\[
\frac{P_{\infty,0}}{Q_{\infty,0}}
=
\frac{\varphi/\sqrt5}{1/\varphi}
=
\frac{\varphi^2}{\sqrt5},
\qquad
\frac{P_{\infty,1}}{Q_{\infty,1}}
=
\frac{1/(\varphi\sqrt5)}{1/\varphi^2}
=
\frac{\varphi}{\sqrt5}.
\]
因此
\[
D_f(P_\infty\Vert Q_\infty)
=
\frac1\varphi\,f\!\left(\frac{\varphi^2}{\sqrt5}\right)
+\frac1{\varphi^2}\,f\!\left(\frac{\varphi}{\sqrt5}\right).
\]
而命题 \ref{prop:fold-bin-f-divergence-collapse} 已给出
\[
D_f(\mu_m\Vert u_m)
=
\frac1\varphi\,f\!\left(\frac{\varphi^2}{\sqrt5}\right)
+\frac1{\varphi^2}\,f\!\left(\frac{\varphi}{\sqrt5}\right)
+O_f\bigl((\varphi/2)^m\bigr).
\]
将两式比较，即得极限恒等式与误差阶。证毕。
\end{proof}

\paragraph{素数寄存器骨架与二原子信息几何的双侧闭合}
上述七条结果表明，素数寄存器侧的幂等结构已经可由秩分层、Dirichlet 级数、布尔骨架的秩层指数稀疏性以及总量级的超指数稀薄性完全控制；与此同时，bin-fold 相对均匀基线的渐近信息几何则最终塌缩为由单一 \(\chi^2\) 常数编码的唯一二原子测度。两侧共同给出一条算术外置与信息塌缩相互对照的闭合描述。

\endinput

\paragraph{链型幂等元的自指合成与可比链上的 \texorpdfstring{$\gcd$}{gcd} 阴影}
在素数寄存器幂等锥的秩分层、链上规范投影的布尔骨架与超指数稀薄性已经明确之后，还可进一步把链型幂等元在自指合成下的闭合规律压缩为下列两条后果。

\begin{theorem}[链型幂等元的自指合成饱和律与可比链上的 \texorpdfstring{$\gcd$}{gcd} 坍缩]\label{thm:xi-chain-idempotent-star-saturation-comparable-gcd}
\leanverified{paper\_xi\_chain\_idempotent\_star\_saturation\_comparable\_gcd}
沿用定理 \ref{thm:killo-finite-semigroup-prime-valuation-arithmetization}、
\ref{thm:killo-chain-interior-godel-gcd-lcm} 与
\ref{thm:killo-prime-register-chain-idempotent-boolean-meet} 的记号。对每个满足
\[
0\in F\subseteq [n]
\]
的子集，记
\[
N_F:=\Phi^{-1}(I_F)\in E_n^{\mathrm{ch}},
\qquad
\kappa(N_F):=\kappa(I_F)=\prod_{i\in F}q_i.
\]
并定义
\[
F\triangleright G:=I_F(G)=\{I_F(g):g\in G\}.
\]
则对任意 \(0\in F,G\subseteq[n]\)，都有
\[
\boxed{\;
N_F\star N_G=N_{F\triangleright G},
\qquad
\kappa(N_F\star N_G)=\prod_{j\in F\triangleright G}q_j.\;}
\]
特别地，若 \(F\) 与 \(G\) 在包含关系下可比，则
\[
\boxed{\;
F\triangleright G=F\cap G,\qquad
N_F\star N_G=N_G\star N_F=N_{F\cap G},\qquad
\kappa(N_F\star N_G)=\gcd(\kappa(N_F),\kappa(N_G)).\;}
\]
因此，在任意包含链所生成的链型幂等子族上，自指合成 \(\star\) 本身便塌缩为交换幂等的 \(\gcd\)-带。
\end{theorem}
\begin{proof}
由定理 \ref{thm:killo-finite-semigroup-prime-valuation-arithmetization}，
\[
f_{N_F\star N_G}=f_{N_F}\circ f_{N_G}=I_F\circ I_G.
\]
因此只需确定 \(I_F\circ I_G\) 的固定点集。记
\[
H:=F\triangleright G=I_F(G).
\]
因 \(0\in G\) 且 \(I_F(0)=0\)，故 \(0\in H\)。

任取 \(i\in[n]\)，置
\[
h:=(I_F\circ I_G)(i)=I_F(I_G(i)).
\]
由于 \(I_G(i)\in G\) 且 \(h\le I_G(i)\le i\)，可知
\[
h\in H\cap\{0,\dots,i\}.
\]
反过来，任取
\[
x\in H\cap\{0,\dots,i\},
\]
则存在 \(g_x\in G\) 使
\[
x=I_F(g_x).
\]
由于 \(x\in F\) 且 \(x\le i\)，由 \(I_G\) 的单调性有
\[
I_G(x)\le I_G(i).
\]
于是
\[
x=I_F(x)=I_F(I_G(x))\le I_F(I_G(i))=h.
\]
故 \(h\) 正是 \(H\cap\{0,\dots,i\}\) 的最大元，也即
\[
(I_F\circ I_G)(i)=I_H(i).
\]
因此
\[
I_F\circ I_G=I_{F\triangleright G}.
\]
再送回 \(S_n\) 侧即得
\[
N_F\star N_G=N_{F\triangleright G}.
\]
由 \(\kappa\) 的定义，
\[
\kappa(N_F\star N_G)=\kappa(I_{F\triangleright G})=\prod_{j\in F\triangleright G}q_j.
\]

若 \(F\subseteq G\)，则对任意 \(g\in G\) 有 \(I_F(g)\in F\)，而对每个 \(f\in F\) 取 \(g=f\in G\) 又得 \(f\in I_F(G)\)，故
\[
F\triangleright G=F=F\cap G.
\]
若 \(G\subseteq F\)，则每个 \(g\in G\) 都满足 \(I_F(g)=g\)，故
\[
F\triangleright G=G=F\cap G.
\]
由此可比情形下
\[
N_F\star N_G=N_{F\cap G}.
\]
交换性来自 \(F\cap G=G\cap F\)，而
\[
\kappa(N_{F\cap G})=\prod_{j\in F\cap G}q_j=\gcd(\kappa(N_F),\kappa(N_G))
\]
则由定理 \ref{thm:killo-chain-interior-godel-gcd-lcm} 的平方自由算术化立得。证毕。
\end{proof}

\leanverified{paper\_xi\_chain\_idempotent\_comparable\_finite\_gcd\_collapse}
\begin{corollary}[可比链式投影复合的一步交换化与像集闭式]\label{cor:xi-chain-idempotent-comparable-finite-gcd-collapse}
沿用定理 \ref{thm:xi-chain-idempotent-star-saturation-comparable-gcd} 的记号。设
\[
0\in F_1,\dots,F_r\subseteq[n]
\qquad (r\ge 2)
\]
两两在包含关系下可比。则有
\[
\boxed{\;
N_{F_1}\star\cdots\star N_{F_r}
=
N_{\bigcap_{j=1}^{r}F_j},\;}
\]
从而
\[
\boxed{\;
\kappa(N_{F_1}\star\cdots\star N_{F_r})
=
\gcd\bigl(\kappa(N_{F_1}),\dots,\kappa(N_{F_r})\bigr).\;}
\]
特别地，该复合与排列顺序无关，并且
\[
\boxed{\;
\operatorname{Im}\!\bigl(f_{N_{F_1}\star\cdots\star N_{F_r}}\bigr)
=
\operatorname{Fix}\!\bigl(f_{N_{F_1}\star\cdots\star N_{F_r}}\bigr)
=
\bigcap_{j=1}^{r}F_j.\;}
\]
\end{corollary}
\begin{proof}
由于 \(F_1,\dots,F_r\) 两两可比，可按包含关系重排为
\[
F_{\sigma(1)}\supseteq F_{\sigma(2)}\supseteq \cdots \supseteq F_{\sigma(r)}.
\]
由定理 \ref{thm:xi-chain-idempotent-star-saturation-comparable-gcd} 的可比情形，
\[
N_{F_{\sigma(1)}}\star N_{F_{\sigma(2)}}=N_{F_{\sigma(2)}}.
\]
继续迭代便得
\[
N_{F_{\sigma(1)}}\star\cdots\star N_{F_{\sigma(r)}}=N_{F_{\sigma(r)}}
=
N_{\bigcap_{j=1}^{r}F_j}.
\]
由于右端仅依赖于最小元 \(\bigcap_j F_j\)，故复合结果与排列顺序无关。

再由定理 \ref{thm:xi-chain-idempotent-star-saturation-comparable-gcd}，
\[
\kappa(N_{\bigcap_j F_j})
=
\prod_{i\in \cap_j F_j}q_i
=
\gcd\bigl(\kappa(N_{F_1}),\dots,\kappa(N_{F_r})\bigr),
\]
得第二式。

最后，由 \(f_{N_H}=I_H\) 以及推论 \ref{cor:killo-projection-idempotent-prime-register}，
\[
\operatorname{Im}(f_{N_H})=\operatorname{Fix}(f_{N_H})=H
\]
对任意 \(H\subseteq[n]\) 且 \(0\in H\) 成立。取
\[
H=\bigcap_{j=1}^{r}F_j
\]
即得像集与不动点集的闭式。证毕。
\end{proof}

\begin{theorem}[链型幂等元相对于严格递增权的加权下降律]\label{thm:xi-chain-idempotent-weighted-energy-descent}
沿用定理 \ref{thm:xi-chain-idempotent-star-saturation-comparable-gcd} 的记号。取任意严格递增权
\[
0=w_0<w_1<\cdots<w_{n-1},
\]
并对任意满足 \(0\in H\subseteq[n]\) 的子集定义
\[
E_w(H):=\sum_{i\in H}w_i.
\]
再对任意 \(0\in F,G\subseteq[n]\) 定义加权缺陷
\[
\Delta_w(F,G):=E_w(G)-E_w(F\triangleright G).
\]
则
\[
\Delta_w(F,G)\ge 0,
\]
并且
\[
\Delta_w(F,G)=0
\iff
G\subseteq F.
\]
\end{theorem}
\begin{proof}
定义
\[
\phi:G\longrightarrow F\triangleright G,
\qquad
\phi(g):=I_F(g).
\]
由定义对每个 \(g\in G\) 都有 \(\phi(g)\le g\)，故严格递增性给出
\[
w_{\phi(g)}\le w_g.
\]
又因 \(F\triangleright G=\operatorname{Im}(\phi)\)，遂有
\[
E_w(F\triangleright G)
=
\sum_{y\in \operatorname{Im}(\phi)}w_y
\le
\sum_{g\in G}w_{\phi(g)}
\le
\sum_{g\in G}w_g
=
E_w(G).
\]
故 \(\Delta_w(F,G)\ge 0\)。

若 \(\Delta_w(F,G)=0\)，则上式两处不等号都必须取等。特别地，
\[
\sum_{g\in G}w_{\phi(g)}=\sum_{g\in G}w_g
\]
迫使 \(w_{\phi(g)}=w_g\) 对每个 \(g\in G\) 成立；由权严格递增只能推出 \(\phi(g)=g\)。于是
\[
I_F(g)=g
\qquad(g\in G),
\]
故 \(g\in F\)，即 \(G\subseteq F\)。反向蕴含显然：若 \(G\subseteq F\)，则 \(I_F(g)=g\) 对所有 \(g\in G\) 成立，从而 \(F\triangleright G=G\) 并且 \(\Delta_w(F,G)=0\)。证毕。
\end{proof}

\begin{theorem}[非自治自指更新轨道的有限终止]\label{thm:xi-chain-idempotent-nonautonomous-finite-termination}
设 \(\{F_t\}_{t\ge 1}\) 为任意序列，其中
\[
0\in F_t\subseteq[n]
\qquad (t\ge 1).
\]
从任意初态 \(0\in G_1\subseteq[n]\) 出发，递归定义
\[
G_{t+1}:=I_{F_t}(G_t)
\qquad (t\ge 1).
\]
令
\[
L(G):=\sum_{i\in G}i.
\]
则对每个 \(t\ge 1\) 都有
\[
L(G_{t+1})\le L(G_t),
\]
并且等号成立当且仅当
\[
G_t\subseteq F_t;
\]
在该情形下进一步有
\[
G_{t+1}=G_t.
\]
特别地，轨道中的严格变化次数至多为
\[
\sum_{i=1}^{n-1}i=\frac{n(n-1)}2,
\]
故该非自治自指动力学在有限步后必稳定为常值；尤其不存在长度大于 \(1\) 的周期轨道。
\end{theorem}
\begin{proof}
在定理 \ref{thm:xi-chain-idempotent-weighted-energy-descent} 中取 \(w_i=i\)，便得
\[
L(G_{t+1})=L(I_{F_t}(G_t))\le L(G_t),
\]
且等号当且仅当 \(G_t\subseteq F_t\)。
若 \(G_t\subseteq F_t\)，则对每个 \(g\in G_t\) 都有 \(I_{F_t}(g)=g\)，故
\[
G_{t+1}=I_{F_t}(G_t)=G_t.
\]

另一方面，\(L(G_t)\) 始终是非负整数，并满足
\[
0\le L(G_t)\le \sum_{i=1}^{n-1}i=\frac{n(n-1)}2.
\]
每发生一次严格变化，\(L\) 至少下降 \(1\)。故严格变化次数至多为 \(\frac{n(n-1)}2\)。当严格下降终止之后，每一步都必须取等，因而由前一段得到状态恒定。若存在长度大于 \(1\) 的周期轨道，则沿一周期 \(L\) 必先下降后回升，这与单调性矛盾。证毕。
\end{proof}

\begin{corollary}[可比性的双向缺陷判别]\label{cor:xi-chain-idempotent-comparability-bidirectional-defect}
沿用定理 \ref{thm:xi-chain-idempotent-weighted-energy-descent} 的记号，并固定任意严格递增权 \(\{w_i\}_{i=0}^{n-1}\)。
则对任意 \(0\in F,G\subseteq[n]\)，有
\[
F,\ G\ \text{在包含偏序下可比}
\iff
\Delta_w(F,G)\Delta_w(G,F)=0.
\]
特别地，若 \(F\) 与 \(G\) 不可比，则必有
\[
\Delta_w(F,G)>0,
\qquad
\Delta_w(G,F)>0.
\]
\end{corollary}
\begin{proof}
若 \(G\subseteq F\)，则由定理 \ref{thm:xi-chain-idempotent-weighted-energy-descent} 得 \(\Delta_w(F,G)=0\)；
若 \(F\subseteq G\)，则同理 \(\Delta_w(G,F)=0\)。
反之，若
\(
\Delta_w(F,G)\Delta_w(G,F)=0
\)，则至少有一项为零。再由定理 \ref{thm:xi-chain-idempotent-weighted-energy-descent}，该零值恰等价于相应包含关系成立，故 \(F,G\) 可比。证毕。
\end{proof}

\begin{theorem}[可比链型自指词的 \texorpdfstring{$\kappa$}{kappa} 完备性]\label{thm:xi-comparable-chain-word-kappa-completeness}
沿用定理 \ref{thm:xi-chain-idempotent-star-saturation-comparable-gcd} 的记号。设
\[
F_1\supseteq F_2\supseteq\cdots\supseteq F_r,
\qquad
0\in F_j\subseteq[n].
\]
令 \(W\) 为由 \(\{N_{F_1},\dots,N_{F_r}\}\) 中若干字母组成的任意 \(\star\)-词，并记这些实际出现的字母对应子集的交为
\[
H:=\bigcap_{\nu}F_{i_\nu}.
\]
则
\[
W=N_H,
\qquad
\kappa(W)=\prod_{i\in H}q_i.
\]
特别地，\(\kappa(W)\) 唯一决定 \(W\)。
\end{theorem}
\begin{proof}
由推论 \ref{cor:xi-chain-idempotent-comparable-finite-gcd-collapse}，任意由可比链中字母组成的有限复合都塌缩为其出现字母对应子集的交，故
\[
W=N_H.
\]
再由 \(\kappa\) 的定义，
\[
\kappa(W)=\kappa(N_H)=\prod_{i\in H}q_i.
\]
由于右端是平方自由素数积，唯一分解立即恢复 \(H\)，从而恢复 \(N_H=W\)。证毕。
\end{proof}

结合推论 \ref{cor:xi-unique-continuous-transverse-register} 可见：在可比链型扇区内，\(\star\)-组合律已经被 \(\kappa\) 的平方自由算术完全锁定；离切片连续模只可能作为独立模式轴外置出现，而不会改变该交换幂等律的内部闭式。

\begin{conjecture}[等基数反链双向缺陷的相邻穿越极小化]\label{conj:xi-antichain-adjacent-crossing-minimal-defect}
沿用定理 \ref{thm:xi-chain-idempotent-weighted-energy-descent} 的记号，并取
\[
w_i:=\log q_i.
\]
对任意满足 \(|F|=|G|\) 的不可比对子，定义总缺陷
\[
D(F,G):=\Delta_w(F,G)+\Delta_w(G,F).
\]
则在固定等基数层 \(|F|=|G|\) 内，\(D(F,G)\) 的极小值应当恰由单个相邻穿越实现；更具体地，极小情形预期等价于存在某个 \(a\) 使
\[
F\setminus G=\{a\},
\qquad
G\setminus F=\{a+1\},
\]
且其余元素完全一致。
\end{conjecture}

上述猜想与定理 \ref{thm:xi-chain-idempotent-weighted-energy-descent} 及推论 \ref{cor:xi-chain-idempotent-comparability-bidirectional-defect} 的机制一致：每个不可比见证都会强制产生双向正缺陷，而相邻穿越正是同时最小化两侧最近回退权差的最紧局部构型。

\noindent 结合定理 \ref{thm:xi-prime-register-minimal-ideal-left-zero-band} 的显式最小理想与推论 \ref{cor:xi-monotone-idempotent-skeleton-superexponential-sparsity} 的超指数稀薄性可见，素数寄存器幂等宇宙内部同时存在两类彼此对照的刚性对象：一类是由 \(\{Q_n^c\}_{c\in[n]}\) 构成的左零吸收核，另一类是由可比链型幂等元所诱导的交换 \(\gcd\)-带阴影；后者虽然在包含链上完全交换化，但在全部幂等元中仍只占超指数稀薄的比例。

\endinput

\paragraph{素数寄存器有限变换半群的自同构刚性、幂等分层与 Green 秩几何}
在素数寄存器的有限变换半群算术化、幂等锥的秩分层、链上内算子格的布尔化以及秩 \(1\) 左零层已经建立之后，还可把这条有限半群链进一步压缩为四条严格后果：全部半群自同构都由同一个坐标置换共轭给出，幂等元按固定指数集获得完全分层，链上内算子格在 Möbius 函数、特征多项式与极大链层面保持纯 Boolean 结构，而全部双边理想与 Green 的 \(\mathcal J\)-类则仅由秩决定。

\begin{theorem}[素数寄存器变换半群的自同构群完全刚性]\label{thm:xi-prime-register-semigroup-automorphism-rigidity}
沿用定理 \ref{thm:killo-finite-semigroup-prime-valuation-arithmetization} 的记号，并记
\[
\mathfrak S_n:=\operatorname{Sym}([n]).
\]
则半群自同构群满足
\[
\boxed{\;
\operatorname{Aut}(S_n,\star)\cong \mathfrak S_n.\;}
\]
更精确地，对每个置换 \(\sigma\in \mathfrak S_n\) 定义
\[
\alpha_\sigma\!\left(\prod_{i=0}^{n-1}q_i^{a_i}\right)
:=
\prod_{i=0}^{n-1}q_i^{\,\sigma(a_{\sigma^{-1}(i)})}.
\]
则
\[
\boxed{\;
\sigma\longmapsto \alpha_\sigma
\;}
\]
给出群同构
\[
\mathfrak S_n\xrightarrow{\sim}\operatorname{Aut}(S_n,\star).
\]
\leanverified{paper\_xi\_prime\_register\_semigroup\_automorphism\_rigidity}
\end{theorem}
\begin{proof}
由定理 \ref{thm:killo-finite-semigroup-prime-valuation-arithmetization}，
\[
\Phi:(S_n,\star)\xrightarrow{\sim}T_n:=([n]^{[n]},\circ),
\qquad
\Phi(N)=f_N.
\]
故只需求 \(T_n\) 的自同构群。

对每个 \(a\in[n]\)，记常值映射
\[
c_a:[n]\to[n],
\qquad
c_a(i)\equiv a.
\]
先刻画这些常值映射的纯半群论性质。若 \(x=c_a\)，则对任意 \(g\in T_n\) 都有
\[
x\circ g=c_a=x.
\]
反之，若某个 \(x\in T_n\) 满足
\[
x\circ g=x
\qquad
(\forall g\in T_n),
\]
则对任意 \(b\in[n]\) 取 \(g=c_b\)，便得
\[
x=x\circ c_b=c_{x(b)},
\]
故 \(x\) 必为常值映射。于是，\(\{c_a:a\in[n]\}\) 恰是 \(T_n\) 中全部左零元的集合。任意半群自同构都保持这一纯乘法性质，因此对每个
\[
\beta\in \operatorname{Aut}(T_n)
\]
必存在唯一置换 \(\sigma\in\mathfrak S_n\) 使
\[
\beta(c_a)=c_{\sigma(a)}
\qquad
(\forall a\in[n]).
\]

任取 \(f\in T_n\)。由于
\[
f\circ c_a=c_{f(a)},
\]
可得
\[
\beta(f)\circ c_{\sigma(a)}
=
\beta(f)\circ \beta(c_a)
=
\beta(f\circ c_a)
=
\beta(c_{f(a)})
=
c_{\sigma(f(a))}.
\]
另一方面，对任意 \(h\in T_n\) 与 \(b\in[n]\) 都有
\[
h\circ c_b=c_{h(b)}.
\]
故上式等价于
\[
\beta(f)(\sigma(a))=\sigma(f(a))
\qquad
(\forall a\in[n]),
\]
亦即
\[
\beta(f)=\sigma\circ f\circ \sigma^{-1}.
\]
因此 \(T_n\) 的每个自同构都由某个唯一的 \(\sigma\in\mathfrak S_n\) 给出，且复合律对应于置换乘法，从而
\[
\operatorname{Aut}(T_n)\cong \mathfrak S_n.
\]

最后把该描述经 \(\Phi\) 共轭回 \(S_n\)。若
\[
N=\prod_{i=0}^{n-1}q_i^{a_i},
\qquad
f_N(i)=a_i,
\]
则
\[
(\sigma\circ f_N\circ \sigma^{-1})(i)
=
\sigma(a_{\sigma^{-1}(i)}).
\]
再由 \(\Phi^{-1}\) 的定义，恰对应于
\[
\Phi^{-1}(\sigma\circ f_N\circ \sigma^{-1})
=
\prod_{i=0}^{n-1}q_i^{\,\sigma(a_{\sigma^{-1}(i)})}
=
\alpha_\sigma(N).
\]
故 \(\sigma\mapsto \alpha_\sigma\) 给出所述群同构。证毕。
\end{proof}

\begin{theorem}[素数寄存器幂等元的固定指数分层与完全计数]\label{thm:xi-prime-register-idempotent-fixed-index-stratification}
沿用定理 \ref{thm:killo-finite-semigroup-prime-valuation-arithmetization} 与推论 \ref{cor:killo-projection-idempotent-prime-register} 的记号，并记
\[
E_n:=\{N\in S_n:\ N\star N=N\}.
\]
对
\[
N=\prod_{i=0}^{n-1}q_i^{a_i}\in S_n
\]
定义固定指数集
\[
F_N:=\{j\in[n]:a_j=j\}.
\]
则下列命题等价：
\begin{enumerate}
  \item \(N\in E_n\)；
  \item 对每个 \(i\in[n]\) 都有
  \[
  a_i\in F_N;
  \]
  \item 存在唯一非空子集 \(F\subseteq[n]\) 使
  \[
  a_j=j\qquad (j\in F),
  \qquad
  a_i\in F\qquad (i\notin F),
  \]
  且此时 \(F=F_N=\operatorname{Im}(f_N)=\operatorname{Fix}(f_N)\)。
\end{enumerate}
进一步，若 \(F\subseteq[n]\) 非空且 \(|F|=k\)，则
\[
\boxed{\;
\#\{N\in E_n:\ F_N=F\}=k^{\,n-k}.\;}
\]
从而对每个 \(1\le k\le n\)，若记
\[
E_{n,k}:=\{N\in E_n:\ |F_N|=k\},
\]
则有
\[
\boxed{\;
|E_{n,k}|=\binom{n}{k}k^{\,n-k},
\qquad
|E_n|=\sum_{k=1}^{n}\binom{n}{k}k^{\,n-k}.\;}
\]
\leanverified{paper_xi_prime_register_idempotent_stratification_seeds}
\leanverified{paper_xi_prime_register_idempotent_stratification_package}
\leanverified{paper\_xi\_prime\_register\_idempotent\_fixed\_index\_stratification}
\end{theorem}
\begin{proof}
由推论 \ref{cor:killo-projection-idempotent-prime-register}，
\[
N\in E_n
\iff
f_N^2=f_N,
\]
并且一旦成立，就有
\[
\operatorname{Im}(f_N)=\operatorname{Fix}(f_N).
\]
由于
\[
f_N(i)=a_i,
\]
故
\[
\operatorname{Fix}(f_N)=\{j\in[n]:f_N(j)=j\}=\{j\in[n]:a_j=j\}=F_N.
\]
于是
\[
N\in E_n
\iff
a_i=f_N(i)\in \operatorname{Fix}(f_N)=F_N
\qquad
(\forall i\in[n]),
\]
这就给出了第 \(1\) 条与第 \(2\) 条的等价性。

再把 \(F:=F_N\) 展开。对 \(j\in F\)，定义已强制
\[
a_j=j.
\]
而对 \(i\notin F\)，第 \(2\) 条恰断言
\[
a_i\in F.
\]
反之，若某个非空子集 \(F\subseteq[n]\) 满足
\[
a_j=j\quad (j\in F),
\qquad
a_i\in F\quad (i\notin F),
\]
则 \(f_N\) 的像集恰为 \(F\)，并且 \(F\) 上每点都被固定，因此
\[
\operatorname{Im}(f_N)=F=\operatorname{Fix}(f_N),
\]
从而 \(f_N\) 幂等，亦即 \(N\in E_n\)。这就得到第 \(3\) 条。

固定非空子集 \(F\subseteq[n]\)，设 \(|F|=k\)。若要求 \(F_N=F\)，则对 \(j\in F\) 的坐标 \(a_j\) 被唯一强制为 \(j\)；而对其余 \(n-k\) 个坐标，每个都可独立取 \(F\) 中任一元素，故总数恰为
\[
k^{\,n-k}.
\]
再对所有大小为 \(k\) 的子集 \(F\) 求和，即得
\[
|E_{n,k}|=\binom{n}{k}k^{\,n-k}.
\]
最后对 \(k\) 求和得到
\[
|E_n|=\sum_{k=1}^{n}\binom{n}{k}k^{\,n-k}.
\]
证毕。
\end{proof}

\leanverified{paper\_xi\_chain\_interior\_boolean\_mobius\_characteristic\_maxchains}
\begin{theorem}[链上内算子格的 Boolean 秩、Möbius 函数、特征多项式与极大链闭式]\label{thm:xi-chain-interior-boolean-mobius-characteristic-maxchains}
沿用定理 \ref{thm:killo-chain-interior-godel-gcd-lcm}、定理 \ref{thm:xi-chain-interior-boolean-flag-closed-form} 与推论 \ref{cor:xi-chain-interior-dirichlet-mobius-characteristic-closure} 的记号。则
\[
(\mathrm{Int}(C_n),\preceq)\cong \mathcal P(\{1,\dots,n-1\})
\]
作为分次格同构；特别地，它是一个秩为 \(n-1\) 的 Boolean 格。并且对任意 \(I\preceq J\) 都有
\[
\boxed{\;
\mu_{\mathrm{Int}(C_n)}(I,J)
=
(-1)^{|\mathrm{Fix}(J)\setminus \mathrm{Fix}(I)|}
=
\mu_{\mathrm{arith}}\!\left(\frac{\kappa(J)}{\kappa(I)}\right).\;}
\]
其特征多项式满足
\[
\boxed{\;
\chi_{\mathrm{Int}(C_n)}(t)=(t-1)^{n-1}.\;}
\]
此外，极大链条数恰为
\[
\boxed{\;
\#\{\text{maximal chains in }\mathrm{Int}(C_n)\}=(n-1)!.\;}
\]
等价地，经 \(\kappa\) 算术化后，全部含 \(q_0\) 的平方自由因子构成一个 Boolean 格，而其 Möbius 函数与特征多项式完全由新增素因子集合控制。
\end{theorem}
\begin{proof}
由定理 \ref{thm:xi-chain-interior-boolean-flag-closed-form}，
\[
(\mathrm{Int}(C_n),\preceq)
\]
与包含 \(0\) 的子集格同构。去掉强制出现的元素 \(0\) 之后，即得
\[
(\mathrm{Int}(C_n),\preceq)\cong \mathcal P(\{1,\dots,n-1\}),
\]
故其确为秩 \(n-1\) 的 Boolean 格。

同一定理已给出
\[
\mu_{\mathrm{Int}(C_n)}(I,J)
=
(-1)^{|\mathrm{Fix}(J)\setminus \mathrm{Fix}(I)|},
\]
而推论 \ref{cor:xi-chain-interior-dirichlet-mobius-characteristic-closure} 进一步给出
\[
\mu_{\mathrm{Int}(C_n)}(I,J)
=
\mu_{\mathrm{arith}}\!\left(\frac{\kappa(J)}{\kappa(I)}\right)
\]
以及
\[
\chi_{\mathrm{Int}(C_n)}(t)=(t-1)^{n-1}.
\]
故前两条盒装公式已经成立。

最后计算极大链。于上述幂集同构下，一条极大链恰对应于
\[
\varnothing
\subset
\{i_1\}
\subset
\{i_1,i_2\}
\subset\cdots\subset
\{i_1,\dots,i_{n-1}\},
\]
其中
\[
(i_1,\dots,i_{n-1})
\]
是 \(\{1,\dots,n-1\}\) 的一个排列。反之，每个排列都唯一给出一条极大链。因此极大链总数正好是
\[
(n-1)!.
\]
证毕。
\end{proof}

\begin{theorem}[素数寄存器半群的 Green 秩分层与左零极小理想]\label{thm:xi-prime-register-green-rank-stratification}
\leanverified{paper\_xi\_prime\_register\_green\_rank\_stratification}
沿用定理 \ref{thm:killo-finite-semigroup-prime-valuation-arithmetization} 的记号，并对
\[
N\in S_n
\]
定义
\[
\operatorname{rk}(N):=|\operatorname{Im}(f_N)|.
\]
则成立下列结论：
\begin{enumerate}
  \item 对任意 \(N\in S_n\)，双边理想生成闭包满足
  \[
  \boxed{\;
  S_n\star N\star S_n
  =
  \{M\in S_n:\ \operatorname{rk}(M)\le \operatorname{rk}(N)\}.\;}
  \]
  因而 Green 的 \(\mathcal J\)-类恰由秩分层：
  \[
  \boxed{\;
  N\mathrel{\mathcal J}M
  \iff
  \operatorname{rk}(N)=\operatorname{rk}(M).\;}
  \]
  \item 对每个 \(1\le k\le n\)，秩恰为 \(k\) 的元素个数为
  \[
  \boxed{\;
  \#\{N\in S_n:\ \operatorname{rk}(N)=k\}
  =
  \binom{n}{k}\,k!\,S(n,k),\;}
  \]
  其中 \(S(n,k)\) 表示第二类 Stirling 数。
  \item 记
  \[
  Q_n:=\prod_{i=0}^{n-1}q_i,
  \qquad
  \mathcal K_n:=\{Q_n^c:\ c\in[n]\}.
  \]
  则
  \[
  \boxed{\;
  \mathcal K_n
  =
  \{N\in S_n:\ \operatorname{rk}(N)=1\},\qquad |\mathcal K_n|=n,\;}
  \]
  并且对任意 \(a,b\in[n]\) 都有
  \[
  \boxed{\;
  Q_n^a\star Q_n^b=Q_n^a.\;}
  \]
  特别地，\(\mathcal K_n\) 是 \((S_n,\star)\) 的唯一极小双边理想，且作为子半群同构于一个 \(n\) 元左零带。
\end{enumerate}
\end{theorem}
\begin{proof}
仍记
\[
T_n:=[n]^{[n]}.
\]
由定理 \ref{thm:killo-finite-semigroup-prime-valuation-arithmetization}，
\[
\Phi:(S_n,\star)\xrightarrow{\sim}(T_n,\circ)
\]
为半群同构，且
\[
\operatorname{rk}(N)=|\operatorname{Im}(f_N)|
\]
正是 \(f_N\) 在 \(T_n\) 中的像集秩。

先证第 \(1\) 条。对任意 \(u,f,v\in T_n\)，有
\[
\operatorname{Im}(f\circ v)\subseteq \operatorname{Im}(f),
\qquad
\operatorname{Im}(u\circ f)\subseteq u(\operatorname{Im}(f)),
\]
故
\[
\operatorname{rk}(u\circ f\circ v)\le \operatorname{rk}(f).
\]
因此
\[
T_n f T_n\subseteq \{g\in T_n:\operatorname{rk}(g)\le \operatorname{rk}(f)\}.
\]

反过来，设 \(g\in T_n\) 满足
\[
\ell:=\operatorname{rk}(g)\le \operatorname{rk}(f)=:k.
\]
取子集
\[
A\subseteq \operatorname{Im}(f),
\qquad
|A|=\ell,
\]
并记
\[
B:=\operatorname{Im}(g).
\]
选定一个双射
\[
\theta:A\xrightarrow{\sim}B.
\]
对每个 \(a\in A\)，任选一点 \(x_a\in[n]\) 使
\[
f(x_a)=a.
\]
定义映射
\[
v:[n]\to[n],
\qquad
v(t):=x_{\theta^{-1}(g(t))},
\]
以及
\[
u:[n]\to[n]
\]
满足
\[
u(a)=\theta(a)\qquad (a\in A),
\]
在 \([n]\setminus A\) 上任意取值。则对每个 \(t\in[n]\)，有
\[
f(v(t))
=
f(x_{\theta^{-1}(g(t))})
=
\theta^{-1}(g(t)),
\]
故
\[
u(f(v(t)))=\theta(\theta^{-1}(g(t)))=g(t).
\]
于是
\[
g=u\circ f\circ v\in T_n f T_n.
\]
因此
\[
T_n f T_n=\{g\in T_n:\operatorname{rk}(g)\le \operatorname{rk}(f)\}.
\]
再经 \(\Phi^{-1}\) 传回 \(S_n\)，便得到所述双边理想公式。由 Green 的 \(\mathcal J\)-关系定义，
\[
N\mathrel{\mathcal J}M
\iff
S_n\star N\star S_n=S_n\star M\star S_n,
\]
故等价于
\[
\operatorname{rk}(N)=\operatorname{rk}(M).
\]

第 \(2\) 条是满射计数的直接转写。构造一个秩为 \(k\) 的映射
\[
f:[n]\to[n]
\]
时，先选像集
\[
\operatorname{Im}(f)\subseteq[n],
\qquad
|\operatorname{Im}(f)|=k,
\]
有 \(\binom{n}{k}\) 种；再在该 \(k\) 元像集上选一个满射，其数量为经典公式
\[
k!\,S(n,k).
\]
故秩为 \(k\) 的映射总数为
\[
\binom{n}{k}k!S(n,k).
\]
经 \(\Phi\) 与 \(S_n\) 一一对应，即得第 \(2\) 条。

第 \(3\) 条对应于 \(k=1\) 层。秩 \(1\) 的映射恰是常值映射
\[
c_a(i)\equiv a
\qquad
(a\in[n]).
\]
而在 \(S_n\) 中，\(c_a\) 的原像正是
\[
Q_n^a=\prod_{i=0}^{n-1}q_i^a.
\]
因此
\[
\mathcal K_n=\{N\in S_n:\operatorname{rk}(N)=1\},
\qquad
|\mathcal K_n|=n.
\]
并且
\[
f_{Q_n^a\star Q_n^b}
=
f_{Q_n^a}\circ f_{Q_n^b}
=
c_a\circ c_b
=
c_a
=
f_{Q_n^a},
\]
故
\[
Q_n^a\star Q_n^b=Q_n^a.
\]
由第 \(1\) 条，秩 \(1\) 层正是最小的非空双边理想层，因此 \(\mathcal K_n\) 是唯一极小双边理想；又因其乘法满足左零律，故它同构于一个 \(n\) 元左零带。证毕。
\end{proof}

\endinput

\paragraph{链上内算子的关联代数反演与随机 Euler 统计}
在链上内算子格的 Boolean 刚性、平方自由 Gödel 像的 \(\gcd/\mathrm{lcm}\) 算术化以及 Möbius 函数闭式已经固定之后，还可把这条有限格链进一步推进到关联代数与概率统计两侧的完全显式闭式。

\begin{theorem}[链上内算子 Gödel 整除格的 Möbius 反演]\label{thm:xi-chain-interior-incidence-algebra-mobius-inversion}
沿用定理 \ref{thm:killo-chain-interior-godel-gcd-lcm} 与定理 \ref{thm:xi-chain-interior-boolean-mobius-characteristic-maxchains} 的记号，记
\[
\mathcal I_n:=\mathrm{Int}(C_n),
\qquad
Q_n:=\prod_{i=0}^{n-1}q_i,
\]
\[
\mathcal D_n
:=
\{d\mid Q_n:\ d\ \text{平方自由且}\ q_0\mid d\}.
\]
则映射
\[
\kappa:\mathcal I_n\longrightarrow \mathcal D_n
\]
在偏序 \(\preceq\) 与整除 \(\mid\) 之间给出格同构，并且对任意 \(I\preceq J\) 都有
\[
\mu_{\mathcal I_n}(I,J)
=
(-1)^{\omega(\kappa(J)/\kappa(I))}
=
\mu_{\mathrm{arith}}\!\left(\frac{\kappa(J)}{\kappa(I)}\right).
\]
等价地，对任意交换群 \(A\) 上的函数 \(f,g:\mathcal D_n\to A\)，关系
\[
g(d)=\sum_{\substack{e\in \mathcal D_n\\ e\mid d}}f(e)
\qquad
(d\in \mathcal D_n)
\]
成立当且仅当
\[
f(d)
=
\sum_{\substack{e\in \mathcal D_n\\ e\mid d}}
\mu_{\mathrm{arith}}\!\left(\frac{d}{e}\right)\,g(e)
=
\sum_{\substack{e\in \mathcal D_n\\ e\mid d}}
(-1)^{\omega(d/e)}\,g(e).
\]
\leanverified{paper\_xi\_chain\_interior\_incidence\_algebra\_mobius\_inversion}
\end{theorem}
\begin{proof}
由定理 \ref{thm:killo-chain-interior-godel-gcd-lcm}，\(\kappa\) 把 \(\mathcal I_n\) 规范识别为全部含 \(q_0\) 的平方自由因子所成之整除格，并把 meet/join 精确送到 \(\gcd/\mathrm{lcm}\)。因此 \(\kappa\) 的确给出 \((\mathcal I_n,\preceq)\) 与 \((\mathcal D_n,\mid)\) 的格同构。

再由定理 \ref{thm:xi-chain-interior-boolean-mobius-characteristic-maxchains}，
\[
\mu_{\mathcal I_n}(I,J)
=
(-1)^{|\mathrm{Fix}(J)\setminus \mathrm{Fix}(I)|}.
\]
而 \(\kappa(J)/\kappa(I)\) 正好是由新增固定点对应素数所组成的平方自由整数，故
\[
|\mathrm{Fix}(J)\setminus \mathrm{Fix}(I)|
=
\omega(\kappa(J)/\kappa(I)),
\]
从而得到 Möbius 函数的算术闭式。

最后，有限偏序集上的一般 Möbius 反演公式适用于 \(\mathcal I_n\)；经由格同构 \(\kappa\) 转写到 \(\mathcal D_n\) 上，即得所示整除偏序反演公式。证毕。
\end{proof}

\begin{theorem}[随机内算子平方自由像的 Euler 乘积与对数方差闭式]\label{thm:xi-chain-interior-random-godel-euler-logvariance}
沿用定理 \ref{thm:killo-chain-interior-godel-gcd-lcm} 的记号，并设随机变量
\[
I\sim \mathrm{Unif}\bigl(\mathrm{Int}(C_n)\bigr).
\]
则存在独立同分布随机变量
\[
\varepsilon_1,\dots,\varepsilon_{n-1}
\stackrel{\mathrm{i.i.d.}}{\sim}
\mathrm{Bernoulli}\!\left(\frac12\right)
\]
使得
\[
\kappa(I)=q_0\prod_{j=1}^{n-1}q_j^{\varepsilon_j}.
\]
因而对任意 \(s\in\CC\) 都有
\[
\EE\!\bigl[\kappa(I)^s\bigr]
=
q_0^s\prod_{j=1}^{n-1}\frac{1+q_j^s}{2}.
\]
特别地，
\[
\EE[\log \kappa(I)]
=
\log q_0+\frac12\sum_{j=1}^{n-1}\log q_j,
\qquad
\mathrm{Var}(\log \kappa(I))
=
\frac14\sum_{j=1}^{n-1}(\log q_j)^2.
\]
\end{theorem}
\begin{proof}
由定理 \ref{thm:xi-chain-interior-boolean-flag-closed-form}，
\[
\mathcal I_n\cong \mathcal P(\{1,\dots,n-1\})
\]
作为 Boolean 格同构；而在均匀分布下，这正对应于每个指标 \(j\in\{1,\dots,n-1\}\) 是否属于固定点集的独立等概率选择。于是存在独立 Bernoulli\((1/2)\) 变量 \(\varepsilon_j\) 使
\[
\mathrm{Fix}(I)=\{0\}\cup \{j:\varepsilon_j=1\}.
\]
再由定理 \ref{thm:killo-chain-interior-godel-gcd-lcm} 的定义，
\[
\kappa(I)=q_0\prod_{j=1}^{n-1}q_j^{\varepsilon_j},
\]
得到第一式。

对任意 \(s\in\CC\)，由独立性可得
\[
\EE\!\bigl[\kappa(I)^s\bigr]
=
q_0^s\prod_{j=1}^{n-1}\EE\!\bigl[q_j^{s\varepsilon_j}\bigr]
=
q_0^s\prod_{j=1}^{n-1}\frac{1+q_j^s}{2},
\]
即得 Euler 乘积闭式。

最后，
\[
\log\kappa(I)=\log q_0+\sum_{j=1}^{n-1}\varepsilon_j\log q_j.
\]
由
\[
\EE[\varepsilon_j]=\frac12,
\qquad
\mathrm{Var}(\varepsilon_j)=\frac14,
\]
以及独立性消去协方差项，便得到
\[
\EE[\log \kappa(I)]
=
\log q_0+\frac12\sum_{j=1}^{n-1}\log q_j,
\]
\[
\mathrm{Var}(\log \kappa(I))
=
\frac14\sum_{j=1}^{n-1}(\log q_j)^2.
\]
证毕。
\end{proof}

\paragraph{素数寄存器幂等元的唯一标准形与链型规范截面的固定点稀薄率}
在素数寄存器幂等元的固定指数分层、链上内算子格的布尔化以及链型幂等元的自指合成规律已经确定之后，还可把固定点纤维的内部结构进一步压缩为如下两条完全显式的后果。

\begin{theorem}[素数寄存器幂等元的唯一标准形与固定点纤维计数]\label{thm:xi-prime-register-idempotent-unique-normal-form}
\leanverified{paper\_xi\_prime\_register\_idempotent\_unique\_normal\_form}
沿用定理 \ref{thm:killo-finite-semigroup-prime-valuation-arithmetization}、推论 \ref{cor:killo-projection-idempotent-prime-register} 与定理 \ref{thm:xi-prime-register-idempotent-fixed-index-stratification} 的记号。则 \(N\in S_n\) 满足 \(N\star N=N\) 当且仅当存在唯一非空子集 \(F\subseteq[n]\) 与唯一映射
\[
\phi:[n]\setminus F\to F
\]
使得
\[
N=
\left(\prod_{j\in F}q_j^{\,j}\right)
\left(\prod_{i\notin F}q_i^{\,\phi(i)}\right).
\]
此时
\[
F=\operatorname{Fix}(f_N)=\operatorname{Im}(f_N).
\]
特别地，若 \(|F|=k\)，则满足 \(\operatorname{Fix}(f_N)=F\) 的幂等元恰有
\[
k^{\,n-k}
\]
个；因此固定点数为 \(k\) 的幂等元总数为
\[
\binom{n}{k}k^{\,n-k},
\]
而全部幂等元总数为
\[
\sum_{k=1}^{n}\binom{n}{k}k^{\,n-k}.
\]
\end{theorem}
\begin{proof}
把
\[
N=\prod_{i=0}^{n-1}q_i^{a_i}
\]
写成估值坐标形式。由定理 \ref{thm:xi-prime-register-idempotent-fixed-index-stratification}，\(N\star N=N\) 当且仅当存在唯一非空子集 \(F\subseteq[n]\) 使
\[
a_j=j\qquad (j\in F),
\qquad
a_i\in F\qquad (i\notin F),
\]
并且此时
\[
F=\operatorname{Fix}(f_N)=\operatorname{Im}(f_N).
\]
于是对每个 \(i\notin F\) 定义
\[
\phi(i):=a_i\in F,
\]
便得到所示分解式
\[
N=
\left(\prod_{j\in F}q_j^{\,j}\right)
\left(\prod_{i\notin F}q_i^{\,\phi(i)}\right).
\]
由于素因子分解唯一，故给定 \(N\) 后，\(F\) 与 \(\phi\) 都被唯一确定。

反过来，若给定非空 \(F\subseteq[n]\) 与映射 \(\phi:[n]\setminus F\to F\)，并按上式定义 \(N\)，则
\[
f_N(j)=j\qquad (j\in F),
\qquad
f_N(i)=\phi(i)\in F\qquad (i\notin F).
\]
因此对任意 \(i\in[n]\) 都有
\[
f_N(f_N(i))=f_N(i),
\]
故 \(f_N\) 幂等，等价地 \(N\star N=N\)。

若固定 \(\operatorname{Fix}(f_N)=F\) 且 \(|F|=k\)，则 \(j\in F\) 的指数被唯一强制，而对每个 \(i\notin F\) 都可独立选择 \(\phi(i)\in F\)，故恰有
\[
k^{\,n-k}
\]
个幂等元。再对所有大小为 \(k\) 的子集 \(F\) 求和，即得
\[
\binom{n}{k}k^{\,n-k}
\]
与总数公式。证毕。
\end{proof}

\begin{theorem}[链上规范截面在固定零点幂等纤维中的唯一单调--收缩实现与超指数稀薄率]\label{thm:xi-prime-register-canonical-fixedzero-section-sparsity}
\leanverified{paper\_xi\_prime\_register\_canonical\_fixedzero\_section\_sparsity}
沿用定理 \ref{thm:killo-chain-interior-godel-gcd-lcm}、定理 \ref{thm:xi-prime-register-idempotent-unique-normal-form} 与定理 \ref{thm:xi-chain-idempotent-star-saturation-comparable-gcd} 的记号。对任意满足 \(0\in F\subseteq[n]\) 的子集，记
\[
N_F^{\mathrm{can}}:=\Phi^{-1}(I_F)=N_F.
\]
则：
\begin{enumerate}
  \item 在全部满足 \(\operatorname{Fix}(f_N)=F\) 的幂等元中，恰有一个对应变换 \(f_N\) 同时单调且收缩；该幂等元正是 \(N_F^{\mathrm{can}}\)；
  \item 因而
  \[
  \mathrm{Int}(C_n)\cong \{N_F^{\mathrm{can}}:0\in F\subseteq[n]\}
  \]
  给出 prime-register 幂等层中的一个规范截面；
  \item 该截面的基数为
  \[
  |\mathrm{Int}(C_n)|=2^{n-1};
  \]
  \item 若记
  \[
  E_n^{(0)}
  :=
  \#\{N\in S_n:\ N\star N=N,\ f_N(0)=0\},
  \]
  则
  \[
  E_n^{(0)}
  =
  \sum_{k=1}^{n}\binom{n-1}{k-1}k^{\,n-k},
  \]
  并且相对密度满足
  \[
  \frac{|\mathrm{Int}(C_n)|}{E_n^{(0)}}
  \le
  \exp\!\Bigl(-\frac12 n\log n+O(n)\Bigr)
  \qquad (n\to\infty).
  \]
\end{enumerate}
\end{theorem}
\begin{proof}
对任意 \(0\in F\subseteq[n]\)，定理 \ref{thm:killo-chain-interior-godel-gcd-lcm} 已给出唯一的内算子
\[
I_F(i)=\max\{j\in F:j\le i\},
\]
其固定点集正为 \(F\)。经 \(\Phi^{-1}\) 拉回，得到幂等元 \(N_F^{\mathrm{can}}\)，且对应变换 \(f_{N_F^{\mathrm{can}}}=I_F\) 显然同时单调且收缩。

反之，设 \(N\in S_n\) 满足 \(N\star N=N\)、\(\operatorname{Fix}(f_N)=F\)，并且 \(f_N\) 单调且收缩。由推论 \ref{cor:killo-projection-idempotent-prime-register}，
\[
\operatorname{Im}(f_N)=\operatorname{Fix}(f_N)=F.
\]
任取 \(i\in[n]\)，记
\[
j:=\max(F\cap\{0,\dots,i\}),
\]
其中该最大元存在，因为 \(0\in F\)。由于 \(f_N(i)\in F\) 且收缩性给出 \(f_N(i)\le i\)，故
\[
f_N(i)\in F\cap\{0,\dots,i\},
\]
从而
\[
f_N(i)\le j.
\]
另一方面，\(j\in F\) 意味着 \(f_N(j)=j\)，再由单调性与 \(j\le i\) 得
\[
j=f_N(j)\le f_N(i).
\]
故 \(f_N(i)=j=I_F(i)\)。由于 \(i\) 任意，遂有
\[
f_N=I_F,
\]
从而 \(N=N_F^{\mathrm{can}}\)。这证明了第 \(1\) 条，也得到第 \(2\) 条。

第 \(3\) 条即定理 \ref{thm:xi-chain-interior-boolean-flag-closed-form} 的基数公式
\[
|\mathrm{Int}(C_n)|=2^{n-1}.
\]

再来计算 \(E_n^{(0)}\)。由定理 \ref{thm:xi-prime-register-idempotent-unique-normal-form}，\(N\star N=N\) 且 \(f_N(0)=0\) 等价于其固定点集 \(F=\operatorname{Fix}(f_N)\) 满足 \(0\in F\)。当 \(|F|=k\) 时，包含 \(0\) 的此类子集共有
\[
\binom{n-1}{k-1}
\]
个，而每个固定点集对应的幂等元数目为
\[
k^{\,n-k}.
\]
故
\[
E_n^{(0)}=\sum_{k=1}^{n}\binom{n-1}{k-1}k^{\,n-k}.
\]

最后估计相对密度。取
\[
m:=\left\lfloor \frac n2\right\rfloor.
\]
则
\[
E_n^{(0)}
\ge
\binom{n-1}{m-1}m^{\,n-m}.
\]
由 Stirling 公式，
\[
\binom{n-1}{m-1}=\exp(O(n)),
\qquad
m^{\,n-m}
=
\exp\!\Bigl(\frac12 n\log n+O(n)\Bigr).
\]
于是
\[
E_n^{(0)}
\ge
\exp\!\Bigl(\frac12 n\log n+O(n)\Bigr).
\]
另一方面，
\[
|\mathrm{Int}(C_n)|=2^{n-1}=\exp(O(n)).
\]
二式相除即得
\[
\frac{|\mathrm{Int}(C_n)|}{E_n^{(0)}}
\le
\exp\!\Bigl(-\frac12 n\log n+O(n)\Bigr).
\]
证毕。
\end{proof}

\paragraph{素数加权二维读出的白化圆对称、最小层析与 Mahalanobis 阈值链}
在定理 \ref{thm:killo-prime-weighted-elliptic-fingerprint-clt} 已把素数加权二维读出识别为椭圆型高斯极限指纹，而推论 \ref{cor:killo-prime-weighted-angular-law-sigma-recovery} 又已把其角向边际压缩为对 \(\Sigma_\infty\) 的唯一恢复之后，还可进一步抽取出如下五条彼此闭合的后果。它们分别把椭圆极限白化为严格圆对称高斯，把固定扇区概率压缩为纯角长度，把两条线性读出的去耦合判据写成 \(\Sigma_\infty\)-正交条件，把二维椭圆指纹压缩为三探针最小层析，并把全部 Mahalanobis 型二值判定统一为同一个 \(\chi_2^2\) 半径变量所支配的嵌套 Bernoulli 链。

\begin{theorem}[白化后极限严格圆对称，Mahalanobis 半径平方满足 \texorpdfstring{$\chi_2^2$}{chi-square-2} 极限律]\label{thm:xi-time-part61aca-primeweighted-whitened-radial-chi2}
沿用定理 \ref{thm:killo-prime-weighted-elliptic-fingerprint-clt} 的设定，记
\[
Y_T:=\frac{1}{\sigma}\Sigma_\infty^{-1/2}
\frac{S_T}{\sqrt{\operatorname{tr}(V_T)}}\in\RR^2.
\]
则有
\[
Y_T\Longrightarrow \mathcal N(0,I_2).
\]
特别地，Mahalanobis 半径平方
\[
R_T^2
:=
\frac{1}{\sigma^2\operatorname{tr}(V_T)}
S_T^\top \Sigma_\infty^{-1}S_T
=
\|Y_T\|^2
\]
满足
\[
R_T^2\Longrightarrow \chi_2^2.
\]
因此对任意 \(q\ge 0\)，都有
\[
\lim_{T\to\infty}\PP(R_T^2\le q)=1-e^{-q/2},
\qquad
\lim_{T\to\infty}\PP(R_T^2>q)=e^{-q/2}.
\]
\end{theorem}
\begin{proof}
由定理 \ref{thm:killo-prime-weighted-elliptic-fingerprint-clt}，
\[
\frac{S_T}{\sqrt{\operatorname{tr}(V_T)}}
\Longrightarrow
\mathcal N(0,\sigma^2\Sigma_\infty),
\]
且 \(\Sigma_\infty\) 正定。对极限施加线性变换
\[
x\longmapsto \sigma^{-1}\Sigma_\infty^{-1/2}x
\]
并应用连续映射定理，立得
\[
Y_T\Longrightarrow \mathcal N(0,I_2).
\]
再由映射 \(y\mapsto \|y\|^2\) 的连续性，
\[
R_T^2=\|Y_T\|^2\Longrightarrow \|Z\|^2,
\qquad
Z\sim \mathcal N(0,I_2).
\]
而二维标准高斯范数平方即 \(\chi_2^2\) 分布，其分布函数为
\[
\PP(\chi_2^2\le q)=1-e^{-q/2}.
\]
结论成立。证毕。
\end{proof}

\begin{corollary}[白化角变量渐近均匀，固定扇区概率仅由夹角决定]\label{cor:xi-time-part61aca-primeweighted-whitened-sector-uniformity}
\leanverified{paper\_xi\_time\_part61aca\_primeweighted\_whitened\_sector\_uniformity}
沿用定理 \ref{thm:xi-time-part61aca-primeweighted-whitened-radial-chi2} 的记号。对任意 Borel 扇区
\[
\mathcal W(\alpha,\beta)
:=
\{(r\cos\theta,r\sin\theta):r\ge 0,\ \theta\in[\alpha,\beta]\},
\qquad
0\le \beta-\alpha\le 2\pi,
\]
若其边界对极限分布为零测，则有
\[
\PP\bigl(Y_T\in \mathcal W(\alpha,\beta)\bigr)
\longrightarrow
\frac{\beta-\alpha}{2\pi}.
\]
等价地，白化坐标中的极限方向变量在 \([0,2\pi)\) 上服从均匀分布。
\end{corollary}
\begin{proof}
由定理 \ref{thm:xi-time-part61aca-primeweighted-whitened-radial-chi2}，
\[
Y_T\Longrightarrow Z\sim \mathcal N(0,I_2).
\]
二维标准高斯分布严格旋转不变，故其极坐标中的角变量在 \([0,2\pi)\) 上均匀分布，并与半径独立。因此
\[
\PP\bigl(Z\in\mathcal W(\alpha,\beta)\bigr)=\frac{\beta-\alpha}{2\pi}.
\]
再由 Portmanteau 定理即得结论。证毕。
\end{proof}

\begin{theorem}[任意两条线性读出的渐近独立性当且仅当 \texorpdfstring{$\Sigma_\infty$}{Sigma-infty}-正交]\label{thm:xi-time-part61aca-primeweighted-sigma-orthogonal-independence}
\leanverified{paper\_xi\_time\_part61aca\_primeweighted\_sigma\_orthogonal\_independence}
沿用定理 \ref{thm:killo-prime-weighted-elliptic-fingerprint-clt} 的设定。对任意非零向量 \(a,b\in\RR^2\)，定义标量读出
\[
X_T(a):=\frac{a^\top S_T}{\sqrt{\operatorname{tr}(V_T)}},
\qquad
X_T(b):=\frac{b^\top S_T}{\sqrt{\operatorname{tr}(V_T)}}.
\]
则
\[
\bigl(X_T(a),X_T(b)\bigr)
\Longrightarrow
\mathcal N(0,\Gamma_{a,b}),
\]
其中
\[
\Gamma_{a,b}
:=
\sigma^2
\begin{pmatrix}
a^\top\Sigma_\infty a & a^\top\Sigma_\infty b\\
a^\top\Sigma_\infty b & b^\top\Sigma_\infty b
\end{pmatrix}.
\]
特别地，下列两命题等价：
\[
a^\top\Sigma_\infty b=0
\qquad\Longleftrightarrow\qquad
X_T(a)\ \text{与}\ X_T(b)\ \text{渐近独立}.
\]
若进一步满足 \(a^\top\Sigma_\infty b=0\)，则对任意 \(\varepsilon_1,\varepsilon_2\in\{\pm1\}\)，都有
\[
\PP\bigl(\operatorname{sgn}X_T(a)=\varepsilon_1,\ \operatorname{sgn}X_T(b)=\varepsilon_2\bigr)
\longrightarrow
\frac14.
\]
\end{theorem}
\begin{proof}
记
\[
A:=
\begin{pmatrix}
a^\top\\
b^\top
\end{pmatrix}.
\]
则
\[
\begin{pmatrix}
X_T(a)\\
X_T(b)
\end{pmatrix}
=
A\frac{S_T}{\sqrt{\operatorname{tr}(V_T)}}.
\]
由定理 \ref{thm:killo-prime-weighted-elliptic-fingerprint-clt} 与连续映射定理，
\[
A\frac{S_T}{\sqrt{\operatorname{tr}(V_T)}}
\Longrightarrow
\mathcal N(0,\sigma^2A\Sigma_\infty A^\top),
\]
而
\[
\sigma^2A\Sigma_\infty A^\top=\Gamma_{a,b}.
\]
这就得到所示联合高斯极限。

二维中心高斯向量独立当且仅当协方差为零，因此渐近独立当且仅当
\[
a^\top\Sigma_\infty b=0.
\]
在该条件下，极限分布是两个独立的一维中心高斯的乘积，故四个符号象限的极限概率均为 \(1/4\)。由于各象限边界由两条直线组成，对连续高斯分布均为零测，Portmanteau 定理遂给出有限 \(T\) 的极限结论。证毕。
\end{proof}

\begin{theorem}[三探针方向方差完全恢复椭圆极限指纹，且两探针永远不足]\label{thm:xi-time-part61aca-primeweighted-three-probe-minimal-tomography}
\leanverified{paper\_xi\_time\_part61aca\_primeweighted\_three\_probe\_minimal\_tomography}
沿用定理 \ref{thm:killo-prime-weighted-elliptic-fingerprint-clt} 的设定，并记
\[
M_\infty:=\sigma^2\Sigma_\infty=
\begin{pmatrix}
M_{11}&M_{12}\\
M_{12}&M_{22}
\end{pmatrix}.
\]
令 \(e_1=(1,0)^\top\)、\(e_2=(0,1)^\top\)，并定义三个方向上的归一化方差极限
\[
v_1:=\lim_{T\to\infty}\frac{\Var(e_1^\top S_T)}{\operatorname{tr}(V_T)},
\qquad
v_2:=\lim_{T\to\infty}\frac{\Var(e_2^\top S_T)}{\operatorname{tr}(V_T)},
\]
\[
v_3:=\lim_{T\to\infty}\frac{\Var((e_1+e_2)^\top S_T)}{\operatorname{tr}(V_T)}.
\]
则有显式恢复公式
\[
M_{11}=v_1,
\qquad
M_{22}=v_2,
\qquad
M_{12}=\frac{v_3-v_1-v_2}{2}.
\]
因此，仅凭三个标量探针 \(e_1,e_2,e_1+e_2\) 的极限方差，即可完全恢复二维椭圆极限指纹 \(M_\infty=\sigma^2\Sigma_\infty\)。

进一步，三探针在此意义下是最小的：任意只含两个方向的方差数据都不足以唯一确定 \(M_\infty\)。
\end{theorem}
\begin{proof}
由于 \(\xi_t\) 独立同分布且 \(\EE\xi_t=0\)、\(\EE\xi_t^2=\sigma^2\)，对任意固定 \(c\in\RR^2\) 都有精确恒等式
\[
\Var(c^\top S_T)=\sigma^2\,c^\top V_T c.
\]
两边同除以 \(\operatorname{tr}(V_T)\)，再令 \(T\to\infty\)，由
\[
\frac{V_T}{\operatorname{tr}(V_T)}\to \Sigma_\infty
\]
得到
\[
\lim_{T\to\infty}\frac{\Var(c^\top S_T)}{\operatorname{tr}(V_T)}
=
\sigma^2 c^\top\Sigma_\infty c
=
c^\top M_\infty c.
\]
分别取 \(c=e_1,e_2,e_1+e_2\)，便得
\[
v_1=M_{11},\qquad v_2=M_{22},
\]
\[
v_3=M_{11}+2M_{12}+M_{22}.
\]
消元即得
\[
M_{12}=\frac{v_3-v_1-v_2}{2}.
\]
故三探针足以唯一恢复 \(M_\infty\)。

再证两探针不足。设只给定两个方向 \(a,b\in\RR^2\) 的方差数据
\[
a^\top M_\infty a,\qquad b^\top M_\infty b.
\]
若 \(a,b\) 共线，则两条数据本质上只提供一个方向上的标量约束，显然不足以确定二维对称矩阵。

若 \(a,b\) 不共线，则对称矩阵空间 \(\mathrm{Sym}_2(\RR)\) 维数为 \(3\)，而线性映射
\[
\mathcal L:\mathrm{Sym}_2(\RR)\to \RR^2,
\qquad
H\longmapsto (a^\top Ha,\ b^\top Hb)
\]
的秩至多为 \(2\)，故其核中存在非零矩阵 \(H\neq 0\)。于是
\[
a^\top (M_\infty+tH)a=a^\top M_\infty a,
\qquad
b^\top (M_\infty+tH)b=b^\top M_\infty b
\]
对一切 \(t\in\RR\) 都成立。又因正定锥是开集，而 \(M_\infty\) 正定，故当 \(|t|\) 充分小时，\(M_\infty+tH\) 仍正定。于是存在不止一个正定矩阵与同一两探针方差数据相容，故两探针永远不足。证毕。
\end{proof}

\begin{theorem}[Mahalanobis 半径诱导的全部二值阈值在极限中形成单参数嵌套 Bernoulli 链]\label{thm:xi-time-part61aca-primeweighted-mahalanobis-threshold-bernoulli-chain}
\leanverified{paper\_xi\_time\_part61aca\_primeweighted\_mahalanobis\_threshold\_bernoulli\_chain}
沿用定理 \ref{thm:xi-time-part61aca-primeweighted-whitened-radial-chi2} 的记号。对任意 \(q\ge 0\)，定义阈值指标
\[
J_T(q):=\mathbf 1_{\{R_T^2>q\}}.
\]
则对每个固定 \(q\)，有
\[
J_T(q)\Longrightarrow \mathrm{Bernoulli}(e^{-q/2}).
\]
更强地，对任意有限阈值链
\[
0\le q_1<q_2<\cdots<q_r,
\]
向量
\[
\bigl(J_T(q_1),\dots,J_T(q_r)\bigr)
\]
在分布上收敛到嵌套 Bernoulli 链
\[
\bigl(\mathbf 1_{\{E>q_1\}},\dots,\mathbf 1_{\{E>q_r\}}\bigr),
\qquad
E\sim \chi_2^2.
\]
因此对任意 \(1\le i\le j\le r\)，都有
\[
\PP\bigl(J_T(q_i)=1,\ J_T(q_j)=1\bigr)\longrightarrow e^{-q_j/2},
\]
\[
\Cov\bigl(J_T(q_i),J_T(q_j)\bigr)\longrightarrow
e^{-q_j/2}\bigl(1-e^{-q_i/2}\bigr).
\]
\end{theorem}
\begin{proof}
由定理 \ref{thm:xi-time-part61aca-primeweighted-whitened-radial-chi2}，
\[
R_T^2\Longrightarrow E\sim \chi_2^2.
\]
而 \(\chi_2^2\) 的尾分布满足
\[
\PP(E>q)=e^{-q/2}.
\]
由于集合 \((q,\infty)\) 的边界点 \(\{q\}\) 对连续分布 \(E\) 为零测，Portmanteau 定理给出
\[
J_T(q)=\mathbf 1_{\{R_T^2>q\}}
\Longrightarrow
\mathbf 1_{\{E>q\}}
\sim \mathrm{Bernoulli}(e^{-q/2}).
\]

对有限阈值链，映射
\[
x\longmapsto
\bigl(\mathbf 1_{\{x>q_1\}},\dots,\mathbf 1_{\{x>q_r\}}\bigr)
\]
仅在有限集合 \(\{q_1,\dots,q_r\}\) 处不连续，而 \(\chi_2^2\) 分布对该有限集赋零质量，因此再次由连续映射定理得到
\[
\bigl(J_T(q_1),\dots,J_T(q_r)\bigr)
\Longrightarrow
\bigl(\mathbf 1_{\{E>q_1\}},\dots,\mathbf 1_{\{E>q_r\}}\bigr).
\]
再由 \(q_i\le q_j\) 时事件包含关系
\[
\{E>q_j\}\subseteq \{E>q_i\}
\]
可得
\[
\PP(E>q_i,\ E>q_j)=\PP(E>q_j)=e^{-q_j/2},
\]
以及
\[
\Cov(\mathbf 1_{\{E>q_i\}},\mathbf 1_{\{E>q_j\}})
=
e^{-q_j/2}-e^{-q_i/2}e^{-q_j/2}
=
e^{-q_j/2}\bigl(1-e^{-q_i/2}\bigr).
\]
结论成立。证毕。
\end{proof}

\noindent
由此，素数加权二维读出的椭圆高斯极限在正文中获得一条完全闭合的白化--层析--判别链：白化后极限严格圆对称，全部固定方向扇区只留下均匀角流；任意两条线性探针的去耦合恰由 \(\Sigma_\infty\)-正交控制；二维椭圆指纹 \(\sigma^2\Sigma_\infty\) 又可由三条一维方差探针以最小方式完全恢复；而一切 Mahalanobis 型二值判定最终都退化为同一个 \(\chi_2^2\) 半径变量的嵌套阈值过程。于是，素数加权读出的二维统计结构便不再是开放的高维对象，而被压缩为一个可白化、可最小层析、且可由单参数半径完全判别的显式闭式体系。

\endinput

\paragraph{链上平方自由 Gödel 编码的平均复杂度最优律}
在定理 \ref{thm:xi-chain-interior-godel-euler-mobius-factorization} 已把链上规范投影的平方自由像写成 Euler--Dirichlet 与 Möbius 扭曲闭式，而推论 \ref{cor:xi-chain-interior-canonical-arithmeticization-optimal-order} 又已把其最坏码长锁定在 \(\Theta(n\log n)\) 之后，还可进一步把同一平方自由编码的平均对数复杂度完全闭式化。由此，链上素数寄存器外置的制度成本便不再只是极端输入所强制，而在均匀分布下同样保持精确的 \(n\log n\) 级下界。

\begin{theorem}[链上归一化平方自由 Gödel 编码的平均对数复杂度达到精确极小值]\label{thm:xi-chain-interior-godel-average-budget-optimality}
\leanverified{paper\_xi\_chain\_interior\_godel\_average\_budget\_optimality}
沿用定理 \ref{thm:killo-chain-interior-godel-gcd-lcm} 与定理 \ref{thm:xi-chain-interior-boolean-flag-closed-form} 的记号，记
\[
\mathcal I_n:=\mathrm{Int}(C_n),
\qquad
C_n=\{0<1<\cdots<n-1\}.
\]
固定互异素数 \(q_0,\dots,q_{n-1}\)，并定义归一化平方自由 Gödel 编码
\[
\kappa'(I_F):=\frac{\kappa(I_F)}{q_0}
=
\prod_{i\in F\setminus\{0\}}q_i
\qquad
(0\in F\subseteq\{0,\dots,n-1\}).
\]
把 \(\mathcal I_n\) 与二元字集 \(\{0,1\}^{n-1}\) 通过对应
\[
I\longleftrightarrow
\bigl(\mathbf 1_{\{1\in \mathrm{Fix}(I)\}},\dots,\mathbf 1_{\{n-1\in \mathrm{Fix}(I)\}}\bigr)
\]
识别。设
\[
E:\mathcal I_n\hookrightarrow \NN_{\ge 1}
\]
为任意可寻址 Gödel 编码，即在上述识别下，存在互异素数 \(p_1,\dots,p_{n-1}\) 与解码函数
\[
\delta_t:\NN\to\{0,1\}
\qquad(1\le t\le n-1)
\]
使得对每个 \(I\in\mathcal I_n\) 都有
\[
\mathbf 1_{\{t\in\mathrm{Fix}(I)\}}
=
\delta_t\!\bigl(v_{p_t}(E(I))\bigr)
\qquad(1\le t\le n-1).
\]
则平均对数复杂度满足精确下界
\[
\frac{1}{|\mathcal I_n|}
\sum_{I\in\mathcal I_n}\log E(I)
\ge
\frac12\sum_{t=1}^{n-1}\log p_t.
\]
特别地，对适当重排后的解码素数有
\[
\frac{1}{|\mathcal I_n|}
\sum_{I\in\mathcal I_n}\log E(I)
\ge
\frac12\log(n!)
=
\Omega(n\log n).
\]

另一方面，规范编码 \(\kappa'\) 精确达到相应极小值：
\[
\frac{1}{|\mathcal I_n|}
\sum_{I\in\mathcal I_n}\log \kappa'(I)
=
\frac12\sum_{t=1}^{n-1}\log q_t.
\]
因而在以 \(q_1,\dots,q_{n-1}\) 为解码素数的全部可寻址编码之中，\(\kappa'\) 的平均对数复杂度达到精确极小值。并且未归一化编码 \(\kappa\) 仅多出一个固定常数项：
\[
\frac{1}{|\mathcal I_n|}
\sum_{I\in\mathcal I_n}\log \kappa(I)
=
\log q_0+\frac12\sum_{t=1}^{n-1}\log q_t.
\]
\end{theorem}
\begin{proof}
由定理 \ref{thm:xi-chain-interior-boolean-flag-closed-form}，\(\mathcal I_n\) 与 \(\{0,1\}^{n-1}\) 规范同构，且
\[
|\mathcal I_n|=2^{n-1}.
\]
因此，题设中的 \(E\) 正是字母表 \(\Sigma=\{0,1\}\)、长度 \(T=n-1\) 的可寻址 Gödel 编码。于是定理 \ref{thm:xi-addressable-godel-product-lower-bound-classification} 在 \(q=2\) 时立即给出
\[
\frac{1}{2^{n-1}}
\sum_{I\in\mathcal I_n}\log E(I)
\ge
\frac12\sum_{t=1}^{n-1}\log p_t.
\]

将 \(p_1,\dots,p_{n-1}\) 递增重排后，不等式左端不变，而对每个 \(t\) 都有
\[
p_t\ge t+1.
\]
故
\[
\sum_{t=1}^{n-1}\log p_t
\ge
\sum_{t=1}^{n-1}\log(t+1)
=
\log(n!),
\]
从而得到
\[
\frac{1}{|\mathcal I_n|}
\sum_{I\in\mathcal I_n}\log E(I)
\ge
\frac12\log(n!)
=
\Omega(n\log n).
\]

再看 \(\kappa'\)。对每个 \(1\le i\le n-1\)，在 \(\mathcal I_n\) 的均匀分布下，事件
\[
i\in\mathrm{Fix}(I)
\]
恰有一半元素满足，因为 \(\mathcal I_n\cong \mathcal P(\{1,\dots,n-1\})\) 是布尔格。故
\[
\frac{1}{|\mathcal I_n|}
\sum_{I\in\mathcal I_n}\log \kappa'(I)
=
\frac{1}{|\mathcal I_n|}
\sum_{I\in\mathcal I_n}\sum_{i\in\mathrm{Fix}(I)\setminus\{0\}}\log q_i
=
\frac12\sum_{i=1}^{n-1}\log q_i.
\]
这与上面的普适下界在解码素数取为 \(q_1,\dots,q_{n-1}\) 时完全相等，故 \(\kappa'\) 达到平均复杂度的精确极小值。

最后，由 \(\kappa(I)=q_0\kappa'(I)\)，对两边取对数并平均即可得到
\[
\frac{1}{|\mathcal I_n|}
\sum_{I\in\mathcal I_n}\log \kappa(I)
=
\log q_0+\frac12\sum_{t=1}^{n-1}\log q_t.
\]
证毕。
\end{proof}

\noindent
于是，链上规范投影格的平方自由 Gödel 化在平均口径下也完全闭合：Dirichlet--Möbius 闭式给出其乘法像的精确算术边界，而可寻址 Gödel 理论则说明，无论最坏情形还是均匀平均，制度成本都同样被 \(n\log n\) 级增长所刚性锁定；规范编码 \(\kappa'\) 因而不仅在实现上自然，而且在平均与最坏两种复杂度口径下都达到精确极小。

\endinput

在链上内算子的随机平方自由像与平均复杂度最优律都已固定之后，归一化规范编码 \(\kappa'\) 的平均位预算还可进一步收束为一条同时包含独立性、方差与最优性的综合闭式。

\begin{theorem}[链上规范平方自由 Gödel 编码的平均位预算精确最优性]\label{thm:xi-time-part61acc-chain-godel-average-variance-budget-optimality}
\leanverified{paper\_xi\_time\_part61acc\_chain\_godel\_average\_variance\_budget\_optimality}
沿用定理 \ref{thm:killo-chain-interior-godel-gcd-lcm} 与定理 \ref{thm:xi-chain-interior-boolean-flag-closed-form} 的记号，记
\[
\mathcal I_n:=\mathrm{Int}(C_n),
\qquad
C_n=\{0<1<\cdots<n-1\},
\]
并固定互异素数 \(q_0,\dots,q_{n-1}\)。对每个 \(I\in\mathcal I_n\)，定义
\[
\kappa'(I):=\frac{\kappa(I)}{q_0}
=
\prod_{i\in \mathrm{Fix}(I)\setminus\{0\}}q_i.
\]
若 \(I\) 在 \(\mathcal I_n\) 上服从均匀分布，则成立：
\begin{enumerate}
  \item 指示变量
        \[
        \mathbf 1_{\{i\in \mathrm{Fix}(I)\}}
        \qquad (1\le i\le n-1)
        \]
        相互独立，且都服从 \(\mathrm{Bernoulli}(1/2)\)；
  \item 因而
        \[
        \EE[\log \kappa'(I)]
        =
        \frac12\sum_{i=1}^{n-1}\log q_i,
        \qquad
        \mathrm{Var}(\log \kappa'(I))
        =
        \frac14\sum_{i=1}^{n-1}(\log q_i)^2;
        \]
  \item 更强地，若
        \[
        E:\mathcal I_n\hookrightarrow \NN_{\ge 1}
        \]
        是任意使用同一组解码素数 \(q_1,\dots,q_{n-1}\) 的可寻址 Gödel 编码，即对每个 \(1\le i\le n-1\) 都存在函数
        \[
        \delta_i:\NN\to\{0,1\}
        \]
        使
        \[
        \mathbf 1_{\{i\in \mathrm{Fix}(I)\}}
        =
        \delta_i\!\bigl(v_{q_i}(E(I))\bigr)
        \qquad (I\in\mathcal I_n),
        \]
        则
        \[
        \frac{1}{|\mathcal I_n|}
        \sum_{I\in\mathcal I_n}\log E(I)
        \ge
        \frac12\sum_{i=1}^{n-1}\log q_i,
        \]
        且 \(\kappa'\) 取等。
\end{enumerate}
\end{theorem}
\begin{proof}
由定理 \ref{thm:xi-chain-interior-random-godel-euler-logvariance}，在 \(\mathcal I_n\) 的均匀分布下，存在独立同分布随机变量
\[
\varepsilon_1,\dots,\varepsilon_{n-1}
\stackrel{\mathrm{i.i.d.}}{\sim}
\mathrm{Bernoulli}\!\left(\frac12\right)
\]
使得
\[
\kappa(I)=q_0\prod_{i=1}^{n-1}q_i^{\varepsilon_i},
\qquad
\varepsilon_i=\mathbf 1_{\{i\in \mathrm{Fix}(I)\}}.
\]
这立即给出第 \(1\) 条。

再由
\[
\kappa'(I)=\frac{\kappa(I)}{q_0}
=
\prod_{i=1}^{n-1}q_i^{\varepsilon_i},
\]
得到
\[
\log \kappa'(I)=\sum_{i=1}^{n-1}\varepsilon_i\log q_i.
\]
由于
\[
\EE[\varepsilon_i]=\frac12,
\qquad
\mathrm{Var}(\varepsilon_i)=\frac14,
\]
并且各 \(\varepsilon_i\) 相互独立，故
\[
\EE[\log \kappa'(I)]
=
\sum_{i=1}^{n-1}\EE[\varepsilon_i]\log q_i
=
\frac12\sum_{i=1}^{n-1}\log q_i,
\]
\[
\mathrm{Var}(\log \kappa'(I))
=
\sum_{i=1}^{n-1}\mathrm{Var}(\varepsilon_i)(\log q_i)^2
=
\frac14\sum_{i=1}^{n-1}(\log q_i)^2.
\]
第 \(2\) 条成立。

最后，由定理 \ref{thm:xi-chain-interior-godel-average-budget-optimality}，任意可寻址 Gödel 编码的平均对数复杂度都至少等于对应解码素数对数和的一半。将该定理应用于此处固定的解码素数族 \(q_1,\dots,q_{n-1}\)，便得
\[
\frac{1}{|\mathcal I_n|}
\sum_{I\in\mathcal I_n}\log E(I)
\ge
\frac12\sum_{i=1}^{n-1}\log q_i.
\]
而同一定理也已证明规范编码 \(\kappa'\) 在同一组解码素数下恰达到该下界。故第 \(3\) 条成立。证毕。
\end{proof}

\noindent
因此，链上规范平方自由 Gödel 编码的平均制度成本不仅在均匀源下具有完全显式的均值与方差，而且在固定解码素数族的全部可寻址编码中仍保持逐坐标意义下的精确最优性。

\endinput


在分解刚性与素数幂单链相图之外，有限分辨率可见算术的幂零层与约化核也可由同一 Fibonacci 模环结构完全恢复。

\begin{theorem}[有限分辨率可见算术的半单核与幂零厚度由 Fibonacci 模数的素因子结构完全决定]\label{thm:xi-visible-arithmetic-nilradical-semisimple-thickness}
\leanverified{paper\_xi\_visible\_arithmetic\_nilradical\_semisimple\_thickness}
沿用定理 \ref{thm:fold-congruence-mod-semirings} 与定理 \ref{thm:xi-visible-arithmetic-omega-controlled-splitting} 的记号，记
\[
R_m^{\mathrm{vis}}
:=
\Omega_m/\!\equiv_m
\cong
\ZZ/(N_m\ZZ),
\qquad
N_m:=F_{m+2},
\qquad
\operatorname{rad}(N_m):=\prod_{p\mid N_m}p.
\]
定义理想
\[
\mathfrak N_m:=\operatorname{rad}\bigl(R_m^{\mathrm{vis}}\bigr)
=
\bigl(\operatorname{rad}(N_m)\bigr)/(N_m)
\subset R_m^{\mathrm{vis}}.
\]
则：
\begin{enumerate}
  \item 幂零根基恰为 \(\mathfrak N_m\)，并且
  \[
  |\mathfrak N_m|
  =
  \frac{N_m}{\operatorname{rad}(N_m)}.
  \]
  \item 约化商满足
  \[
  R_m^{\mathrm{vis}}/\mathfrak N_m
  \cong
  \ZZ/\operatorname{rad}(N_m)\ZZ
  \cong
  \prod_{p\mid N_m}\FF_p.
  \]
  \item \(\mathfrak N_m\) 的最小幂零指数为
  \[
  \nu_m
  :=
  \min\{k\ge 1:\ \mathfrak N_m^k=0\}
  =
  \max_{p^\alpha\parallel N_m}\alpha.
  \]
  \item \(R_m^{\mathrm{vis}}\) 为局部半环，当且仅当 \(N_m\) 为素数幂；\(R_m^{\mathrm{vis}}\) 为约化半环，当且仅当 \(N_m\) 平方自由。
\end{enumerate}
\end{theorem}
\begin{proof}
将 \(N_m\) 分解为
\[
N_m=\prod_{j=1}^{r}p_j^{\alpha_j}.
\]
由中国剩余定理与定理 \ref{thm:fold-congruence-mod-semirings}，
\[
R_m^{\mathrm{vis}}
\cong
\ZZ/(N_m\ZZ)
\cong
\prod_{j=1}^{r}\ZZ/(p_j^{\alpha_j}\ZZ).
\]

对每个局部因子 \(\ZZ/(p_j^{\alpha_j}\ZZ)\)，其幂零元恰为理想 \((p_j)/(p_j^{\alpha_j})\)；故全局幂零根基正是这些局部幂零理想的直积，也即
\[
\mathfrak N_m
=
\bigl(\operatorname{rad}(N_m)\bigr)/(N_m).
\]
其基数因而为
\[
|\mathfrak N_m|
=
\frac{N_m}{\operatorname{rad}(N_m)}.
\]
这证明了第 \(1\) 条。

将每个局部因子模去其 maximal ideal \((p_j)\) 后，便得到剩余域 \(\FF_{p_j}\)。于是
\[
R_m^{\mathrm{vis}}/\mathfrak N_m
\cong
\prod_{j=1}^{r}\FF_{p_j}
\cong
\ZZ/\operatorname{rad}(N_m)\ZZ,
\]
从而第 \(2\) 条成立。

再看幂零指数。对任意 \(k\ge 1\)，条件 \(\mathfrak N_m^k=0\) 等价于
\[
\operatorname{rad}(N_m)^k\in (N_m),
\]
亦即对每个 \(j\) 都有 \(k\ge \alpha_j\)。因此最小此类 \(k\) 正是
\[
\max_j\alpha_j
=
\max_{p^\alpha\parallel N_m}\alpha,
\]
故第 \(3\) 条成立。

最后，\(R_m^{\mathrm{vis}}\) 为局部半环，当且仅当 CRT 分解中只有一个局部因子，亦即 \(r=1\)，等价于 \(N_m\) 为素数幂；而 \(R_m^{\mathrm{vis}}\) 为约化半环，当且仅当幂零根基为零，等价于所有 \(\alpha_j=1\)，亦即 \(N_m\) 平方自由。第 \(4\) 条成立。证毕。
\end{proof}

\endinput

\paragraph{分辨率算术几何：合同格、幂等元、幂零根与四相分类}
上一小节已经把固定分辨率 \(m\) 的折叠同余商
\[
R_m:=\Omega_m/\!\equiv_m
\]
识别为有限剩余类半环，并与 Fibonacci 模数层上的剩余类半环规范同构：
\[
R_m\cong \ZZ/(N_m\ZZ),
\qquad
N_m:=F_{m+2}.
\]
记
\[
\bar\pi_m:R_m\xrightarrow{\ \sim\ }\ZZ/(N_m\ZZ)
\]
为这一规范同构。于是，固定分辨率对象的代数性质不再依赖 \(\Omega_m\) 的具体组合呈现，而完全由 \(N_m\) 的素因子型控制。本段把这一事实组织为一层分辨率算术几何：全部粗粒化合同、投影型端同态、幂零方向、局部/分裂/约化相以及 window-\(6\) 的算术读出，均由单个 Fibonacci 模数的约数格与 Chinese remainder 结构统一决定。

\begin{theorem}[分辨率商半环的合同格与 Fibonacci 约数格]\label{thm:xi-fold-congruence-quotient-divisor-lattice}
\leanverified{paper\_xi\_fold\_congruence\_quotient\_divisor\_lattice}
对每个约数 \(d\mid N_m\)，令
\[
\rho_d:\ZZ/(N_m\ZZ)\longrightarrow \ZZ/(d\ZZ),
\qquad
\overline n\longmapsto \overline n \pmod d,
\]
并在 \(R_m\) 上定义关系
\[
x\,\Theta_d\,y
\quad\Longleftrightarrow\quad
\rho_d\!\bigl(\bar\pi_m(x)\bigr)=\rho_d\!\bigl(\bar\pi_m(y)\bigr).
\]
则 \(\Theta_d\) 是 \(R_m\) 的半环合同，且
\[
R_m/\Theta_d\cong \ZZ/(d\ZZ).
\]
此外，\(R_m\) 的任意半环合同都唯一地写成某个 \(\Theta_d\)；并且在偏序上有
\[
\Theta_{d_1}\le \Theta_{d_2}
\iff
d_2\mid d_1.
\]
换言之，
\[
\operatorname{Con}(R_m)
\cong
\operatorname{Idl}\bigl(\ZZ/(N_m\ZZ)\bigr)
\cong
\operatorname{Div}(N_m)^{\mathrm{op}}.
\]
\end{theorem}
\begin{proof}
由于 \(d\mid N_m\)，映射 \(\rho_d\) 良定义且保持加法与乘法；于是
\[
\rho_{m,d}:=\rho_d\circ \bar\pi_m:R_m\longrightarrow \ZZ/(d\ZZ)
\]
是满半环同态，而 \(\Theta_d\) 正是其核合同，从而 \(R_m/\Theta_d\cong \ZZ/(d\ZZ)\)。

现设 \(\Theta\) 为 \(R_m\) 上任意半环合同。经 \(\bar\pi_m\) 输运后，得到 \(\ZZ/(N_m\ZZ)\) 上的半环合同 \(\approx\)。由于 \(\ZZ/(N_m\ZZ)\) 是交换环，其任意半环商自动带有加法逆元，故 \(\approx\) 实际上就是环合同，并由某个理想唯一决定。另一方面，
\(\ZZ/(N_m\ZZ)\) 的全部理想恰为
\[
(d)/(N_m)
\qquad
(d\mid N_m),
\]
故 \(\Theta\) 必唯一对应某个 \(\Theta_d\)。

最后，\(\Theta_{d_1}\le \Theta_{d_2}\) 当且仅当对应理想满足
\[
(d_1)/(N_m)\subseteq (d_2)/(N_m),
\]
而这又等价于 \(d_2\mid d_1\)。于是合同格与约数格确为反同构。证毕。
\end{proof}

\begin{theorem}[分辨率商半环的非幺端同态与幺自同构刚性]\label{thm:xi-fold-congruence-unital-automorphism-rigidity}
\leanverified{paper\_xi\_fold\_congruence\_unital\_automorphism\_rigidity}
记
\[
\operatorname{End}_0(R_m)
:=
\{\phi:R_m\to R_m:\ \phi\ \text{保持 }0,+,\cdot\},
\]
\[
\operatorname{Idem}(R_m)
:=
\{e\in R_m:\ e^2=e\}.
\]
则映射
\[
\operatorname{End}_0(R_m)\longrightarrow \operatorname{Idem}(R_m),
\qquad
\phi\longmapsto \phi(1)
\]
为双射，其逆由
\[
e\longmapsto \phi_e,
\qquad
\phi_e(x):=ex
\]
给出。若
\[
N_m=\prod_{j=1}^{r}p_j^{a_j},
\qquad r=\omega(N_m),
\]
则
\[
|\operatorname{End}_0(R_m)|
=
|\operatorname{Idem}(R_m)|
=
2^r.
\]
特别地，
\[
\Aut_{\mathrm{unital\; semiring}}(R_m)=\{\mathrm{id}_{R_m}\}.
\]
\end{theorem}
\begin{proof}
经 \(\bar\pi_m\) 识别后，只需在 \(A:=\ZZ/(N_m\ZZ)\) 上证明。任取 \(\phi\in \operatorname{End}_0(A)\)，令
\[
e:=\phi(\overline 1).
\]
由乘法保持性，
\[
e^2=\phi(\overline 1)\phi(\overline 1)=\phi(\overline 1\cdot \overline 1)=\phi(\overline 1)=e,
\]
故 \(e\) 为幂等元。对任意 \(\overline k\in A\)，有
\[
\phi(\overline k)
=
\phi(k\cdot \overline 1)
=
k\,\phi(\overline 1)
=
\overline k\,e,
\]
因此 \(\phi=\phi_e\)，从而 \(\phi\mapsto \phi(\overline 1)\) 为单射。

反之，若 \(e^2=e\)，则
\[
\phi_e(\overline x+\overline y)=e(\overline x+\overline y)=e\overline x+e\overline y,
\]
\[
\phi_e(\overline x\,\overline y)=e\,\overline x\,\overline y
=
e^2\,\overline x\,\overline y
=
(e\overline x)(e\overline y)
=
\phi_e(\overline x)\phi_e(\overline y),
\]
故 \(\phi_e\) 为半环端同态。这就证明了所述双射。

再由 Chinese remainder 分解
\[
A\cong \prod_{j=1}^{r}\ZZ/(p_j^{a_j}\ZZ),
\]
每个局部因子中的幂等元只能是 \(0\) 或 \(1\)，故全体幂等元数恰为 \(2^r\)。

若 \(\phi\) 还是保持单位元的半环自同构，则 \(\phi(\overline 1)=\overline 1\)，于是对全部 \(\overline x\in A\) 都有
\[
\phi(\overline x)=\overline x\cdot \overline 1=\overline x.
\]
故 \(\phi=\mathrm{id}_A\)，传回 \(R_m\) 即得结论。证毕。
\end{proof}

\begin{theorem}[分辨率商半环的幂零根、约化核与平方自由判别]\label{thm:xi-fold-congruence-nilradical-reduced-quotient}
设
\[
\operatorname{rad}(N_m):=\prod_{p\mid N_m}p,
\qquad
\mathcal N_m:=\sqrt{(0)}\subset R_m.
\]
则
\[
\mathcal N_m
=
\bigl(\operatorname{rad}(N_m)\bigr)/(N_m),
\qquad
R_m/\mathcal N_m
\cong
\ZZ/\bigl(\operatorname{rad}(N_m)\ZZ\bigr)
\cong
\prod_{p\mid N_m}\FF_p.
\]
并且 \(R_m\) 约化当且仅当 \(N_m\) 平方自由；等价地，当且仅当 \(F_{m+2}\) 不含平方因子。
\end{theorem}
\begin{proof}
这正是定理 \ref{thm:xi-visible-arithmetic-nilradical-semisimple-thickness} 在记号 \(R_m=\Omega_m/\!\equiv_m\) 下的直接改写。该定理已给出幂零根基的显式理想描述、约化商的 CRT 闭式以及“约化当且仅当平方自由”的等价判别。证毕。
\end{proof}

\begin{theorem}[分辨率算术几何的四相分类]\label{thm:xi-fold-congruence-crt-primitive-idempotents}
设
\[
N_m=\prod_{j=1}^{r}p_j^{a_j}
\]
为 \(N_m\) 的素因子分解，其中 \(p_j\) 两两不同。则
\[
R_m\cong \prod_{j=1}^{r}\ZZ/(p_j^{a_j}\ZZ),
\]
并且恰落在下列四种互斥且穷尽的算术几何相之一：
\begin{enumerate}
  \item \textbf{域相}：\(N_m=p\) 为素数，此时
  \[
  R_m\cong \FF_p.
  \]
  \item \textbf{局部非约化相}：\(N_m=p^a\) 且 \(a\ge 2\)，此时 \(R_m\) 为局部对象，但含非零幂零元。
  \item \textbf{分裂约化相}：\(r\ge 2\) 且 \(a_1=\cdots=a_r=1\)，此时
  \[
  R_m\cong \prod_{j=1}^{r}\FF_{p_j}.
  \]
  \item \textbf{分裂非约化相}：\(r\ge 2\) 且至少有一个 \(a_j\ge 2\)，此时 \(R_m\) 既发生分裂，又在至少一个局部分量中携带幂零方向。
\end{enumerate}
此外，\(R_m\) 的原始中心幂等元恰有 \(r\) 个，而全部幂等元总数恰为 \(2^r\)。
\end{theorem}
\begin{proof}
Chinese remainder theorem 给出
\[
R_m
\cong
\ZZ/(N_m\ZZ)
\cong
\prod_{j=1}^{r}\ZZ/(p_j^{a_j}\ZZ).
\]
若 \(r=1\) 且 \(a_1=1\)，则 \(R_m\cong \FF_{p_1}\)，故为域；若 \(r=1\) 且 \(a_1\ge 2\)，则 \(R_m\cong \ZZ/(p_1^{a_1}\ZZ)\) 为局部环，并含非零幂零元 \(p_1\)。这给出前两相。

若 \(r\ge 2\) 且全部 \(a_j=1\)，则所有局部因子都约化，故
\[
R_m\cong \prod_{j=1}^{r}\FF_{p_j}
\]
约化且非局部；若 \(r\ge 2\) 且某个 \(a_j\ge 2\)，则整体仍非局部，但至少一个因子含非零幂零元，从而 \(R_m\) 非约化。这给出后三相，且四种情形显然互斥并且穷尽。

最后，由定理 \ref{thm:xi-fold-congruence-unital-automorphism-rigidity}，全部幂等元数等于 \(2^r\)。而原始中心幂等元正是各坐标投影
\[
e_j:=(0,\dots,0,1,0,\dots,0)
\qquad
(1\le j\le r),
\]
故恰有 \(r\) 个。证毕。
\end{proof}

\begin{proposition}[window-\texorpdfstring{$6$}{6} 的算术分裂与组合刚性应严格区分]\label{prop:xi-window6-resolution-arithmetic-split-separation}
当 \(m=6\) 时，
\[
N_6=F_8=21=3\cdot 7,
\qquad
R_6\cong \ZZ/(21\ZZ)\cong \FF_3\times \FF_7.
\]
因此 \(R_6\) 恰有 \(4\) 个幂等元，并且只对应 \(2\) 个非平凡互补中央幂等元，也即 \(2\) 个非平凡直因子。另一方面，window-\(6\) 的组合层刚性分裂为
\[
X_6=X_6^{\mathrm{cyc}}\sqcup X_6^{\mathrm{bdry}},
\qquad
|X_6|=21,\quad |X_6^{\mathrm{cyc}}|=18,\quad |X_6^{\mathrm{bdry}}|=3,
\]
而二进区间折叠的纤维大小分布则为
\[
d_6^{\mathrm{bin}}(w)\in\{2,3,4\}.
\]
因而，window-\(6\) 上的 \(4\) 个幂等元只参数化投影型非幺端同态的数目，并不对应于“四块 primitive 不变块”的组合分解。
\end{proposition}
\begin{proof}
由定理 \ref{thm:xi-fold-congruence-crt-primitive-idempotents}，\(R_6\) 处于分裂约化相，并且
\[
R_6\cong \FF_3\times \FF_7.
\]
于是由定理 \ref{thm:xi-fold-congruence-unital-automorphism-rigidity} 得
\[
|\operatorname{Idem}(R_6)|=2^{\omega(21)}=4.
\]
同时，由积环的坐标分解可知，其原始中心幂等元仅为 \((1,0)\) 与 \((0,1)\) 两个，故恰有两个非平凡互补直因子；更细的单位--零因子与幂等元显式代表则由定理 \ref{thm:xi-window6-residue-semiring-unit-zerodivisor-idempotent-trichotomy} 给出。

另一方面，window-\(6\) 的组合刚性分裂
\[
X_6=X_6^{\mathrm{cyc}}\sqcup X_6^{\mathrm{bdry}},
\qquad
|X_6^{\mathrm{cyc}}|=18,\quad |X_6^{\mathrm{bdry}}|=3,
\]
已由定理 \ref{thm:xi-window6-cyclic-kernel-bc3-root-datum} 记录，而推论 \ref{cor:fold-bin-m6-fiber-hist} 表明其纤维直方图只取三档 \(2,3,4\)。因此，算术分裂、纤维大小分布与组合刚性分裂分别编码不同层级的信息，不可彼此混同。证毕。
\end{proof}

\noindent 综上，固定分辨率层上的代数几何型完全由 Fibonacci 模数 \(N_m=F_{m+2}\) 的约数结构与指数向量决定；下一小节则把这种有限层刚性提升为跨分辨率兼容系统上的 Fibonacci-adic 逆极限环。

\endinput

\paragraph{局部化整数族的统一有限化障碍与连续时钟/素数账本裂解}
在定理 \ref{thm:xi-localized-integers-padic-kernel-rigidity}、
推论 \ref{cor:xi-localized-integers-prime-axis-increment-law} 与
定理 \ref{thm:xi-localized-integers-connected-rational-blindness} 已经把
\(\ZZ[S^{-1}]\) 的连通圆维部分与全不连通 \(p\)-进核部分严格分离之后，
还可进一步抽出两条结构性后果：其一是该对象族不存在统一的有限生成忠实素数账本，
其二是连续可见时钟冻结为一维，而离散审计通道数随素数支撑无界增长。

\begin{theorem}[素数账本的统一有限忠实化不可能性]\label{thm:xi-localized-integers-no-uniform-finitely-generated-ledger}
\leanverified{paper\_xi\_localized\_integers\_no\_uniform\_finitely\_generated\_ledger}
不存在固定有限生成交换群 \(L\) 与一族映射
\[
\Phi_S:G_S=\ZZ[S^{-1}]_{>0}\longrightarrow L
\qquad
(\,|S|<\infty\,)
\]
同时满足下列条件：
\begin{enumerate}
  \item \(\Phi_S(xy)=\Phi_S(x)+\Phi_S(y)\) 对任意 \(x,y\in G_S\) 成立；
  \item 记录在 prime-ledger 意义下忠实，即随着 \(S\) 增长所新增的素数支撑轴可由记录统一且无碰撞地恢复。
\end{enumerate}
换言之，有限素数局部化对象族不存在统一的有限生成可加账本，既能保持乘法对象的可组合性，又能对素数支撑给出忠实审计。
\end{theorem}
\begin{proof}
由推论 \ref{cor:xi-localized-integers-prime-axis-increment-law}，
每加入一个新素数 \(p\)，对偶端就恰增加一条新的 \(p\)-进轴
\[
\ZZ_p,
\]
而连通商 \(\TT\) 与有限秩圆维 \(\cdim_{\mathrm{fr}}=1\) 完全不变。
因此，若记录要忠实地恢复素数支撑，
它就必须能够区分任意大的有限素数集所对应的两两独立的 prime-axis 增量。

另一方面，固定有限生成交换群 \(L\) 只具有有限自由秩与有限 torsion 支撑，
不可能承载无界增长的彼此独立的素数账本方向。
故不存在满足上述两条要求的统一协议族 \((\Phi_S)_S\)。

从另一侧看，若把 \(\Phi_S\) 限制到正整数乘法幺半群，
则得到一类从乘法对象到有限生成可消去加法账本的忠实压缩；
这与定理 \ref{thm:killo-prime-freedom-non-finitizability} 所排除的有限生成 ledgerization 障碍一致。
故统一有限忠实化不可能。证毕。
\end{proof}

\begin{corollary}[一维连续时钟与无界离散账本道的裂解]\label{cor:xi-localized-integers-oneclock-unbounded-prime-ledgers}
对每个固定有限素数集 \(S\subset\mathbb P\)，都有
\[
\cdim_{\mathrm{fr}}(G_S)=1.
\]
但任何忠实的外置审计若要精确保留 \(S\)，至少需要 \(|S|\) 条两两独立的全不连通素数账本道；
等价地，在对象族
\[
\{\,\ZZ[S^{-1}]:\ |S|<\infty\,\}
\]
上，连续可见维数冻结为 \(1\)，而离散素数账本通道数无上界。
\end{corollary}
\begin{proof}
定理 \ref{thm:xi-localized-integers-padic-kernel-rigidity} 已给出
\[
(\widehat{G_S})^0\cong \TT,
\qquad
\cdim_{\mathrm{fr}}(G_S)=1,
\qquad
K_S=\prod_{p\in S}\ZZ_p.
\]
因此连续可见部分始终只有单一圆时钟。

另一方面，\(K_S\) 的每个因子 \(\ZZ_p\) 都对应一个不可与其它素数混淆的 \(p\)-主方向；
由定理 \ref{thm:xi-localized-integers-padic-kernel-rigidity}，
\[
p\in S
\iff
\text{\(K_S\) 的 \(p\)-主因子非平凡}.
\]
故若要忠实地保留 \(S\)，至少必须为每个 \(p\in S\) 保留一条独立的离散账本道。
这给出下界 \(|S|\)；当 \(S\) 在所有有限素数集上变化时，该下界无界。
因此“一维连续时钟/无界离散账本道”的裂解严格成立。证毕。
\end{proof}

\endinput

\paragraph{dyadic Fibonacci 完成层的局部分解、出生过滤与有限扇区障碍}
有限分辨率可见算术已经把每一层折叠同余商识别为单个 Fibonacci 模环，而定理 \ref{thm:xi-time-consistent-inverse-limit-readout} 与定理 \ref{thm:xi-rmin-is-zhat} 也已给出跨层兼容系统的逆极限接口。若把这两条链限制到 dyadic 子塔 \(m=2^k-2\)，则完成化对象表现出一种与通常 \(p\)-adic 增厚显著不同的局部结构：固定素数方向上的厚度在其首次出现后立即稳定，而全部无穷性仅由新素数方向的持续出生提供。

\begin{theorem}[dyadic Fibonacci 完成层的稳定局部分解]\label{thm:xi-dyadic-fibonacci-completion-stable-local-splitting}
沿用定理 \ref{thm:fold-congruence-mod-semirings} 的记号，并定义
\[
R_k:=\Omega_{2^k-2}/\!\equiv_{2^k-2}
\cong
\ZZ/(F_{2^k}\ZZ)
\qquad(k\ge 2),
\]
\[
\mathfrak R_{\mathrm{dy}}
:=
\varprojlim_{k\ge 2}R_k,
\qquad
\mathcal P_{\mathrm{dy}}
:=
\{p\in\mathbb P:\ \exists k\ge 2,\ p\mid F_{2^k}\}.
\]
对每个 \(p\in\mathcal P_{\mathrm{dy}}\)，记
\[
k_p:=\min\{k\ge 2:\ p\mid F_{2^k}\},
\qquad
e_p:=v_p(F_{2^{k_p}}).
\]
则存在与过渡映射相容的规范环同构
\[
R_k
\cong
\prod_{\substack{p\in\mathcal P_{\mathrm{dy}}\\ k_p\le k}}
\ZZ/(p^{e_p}\ZZ)
\qquad(k\ge 2),
\]
并且存在规范拓扑环同构
\[
\mathfrak R_{\mathrm{dy}}
\cong
\prod_{p\in\mathcal P_{\mathrm{dy}}}\ZZ/(p^{e_p}\ZZ).
\]
特别地，\(\mathfrak R_{\mathrm{dy}}\) 的无穷性并非来自任何固定素数方向上的无界 \(p\)-进增厚，而完全来自无穷多个新生素数方向的横向累积。
\end{theorem}
\begin{proof}
由定理 \ref{thm:fold-congruence-mod-semirings}，
\[
R_k\cong \ZZ/(F_{2^k}\ZZ).
\]
再由 Fibonacci 数的恒等式
\[
\gcd(F_a,F_b)=F_{\gcd(a,b)},
\]
可知 \(a\mid b\) 时有 \(F_a\mid F_b\)。因此 \((R_k)_{k\ge 2}\) 通过自然约化映射构成逆系统。

先注意 \(2,5\notin\mathcal P_{\mathrm{dy}}\)：\(F_n\) 为偶数当且仅当 \(3\mid n\)，而 \(5\mid F_n\) 当且仅当 \(5\mid n\)；但 \(2^k\) 既不被 \(3\) 整除，也不被 \(5\) 整除。故对任意 \(p\in\mathcal P_{\mathrm{dy}}\)，必有 \(p\neq 2,5\) 且 \(p\) 为奇素数。

记 \(\rho(p)\) 为 Fibonacci 数列模 \(p\) 的出现阶。由 Wall 的标准结果，
\[
p\mid F_n \iff \rho(p)\mid n,
\]
并且对 \(p\neq 2,5\) 有 \(p\)-进赋值公式\cite{Wall1960FibonacciModm,RowlandMedina2009PadicValuationFibonacci}
\[
v_p(F_n)=
\begin{cases}
v_p(n)+v_p(F_{\rho(p)}), & \rho(p)\mid n,\\[1mm]
0, & \rho(p)\nmid n.
\end{cases}
\]
由于 \(p\mid F_{2^{k_p}}\)，可知 \(\rho(p)\mid 2^{k_p}\)，故 \(\rho(p)=2^j\) 对某个 \(j\le k_p\)。若 \(j<k_p\)，则
\[
p\mid F_{\rho(p)}=F_{2^j},
\]
与 \(k_p\) 的极小性矛盾。因此
\[
\rho(p)=2^{k_p}.
\]
于是对任意 \(k\ge k_p\)，上述赋值公式给出
\[
v_p(F_{2^k})
=
v_p(2^{k-k_p})+v_p(F_{2^{k_p}})
=
e_p,
\]
因为 \(p\) 为奇素数，故 \(v_p(2^{k-k_p})=0\)。这说明每个素数方向在首次出现后其指数立即稳定。

对固定 \(k\)，中国剩余定理给出
\[
R_k
\cong
\prod_{p\mid F_{2^k}}\ZZ/(p^{v_p(F_{2^k})}\ZZ).
\]
而上一段已表明：若 \(p\mid F_{2^k}\)，则 \(k_p\le k\) 且 \(v_p(F_{2^k})=e_p\)；反之，若 \(k_p\le k\)，则 \(p\mid F_{2^k}\)。故
\[
R_k
\cong
\prod_{\substack{p\in\mathcal P_{\mathrm{dy}}\\ k_p\le k}}
\ZZ/(p^{e_p}\ZZ).
\]
在此识别下，过渡映射 \(R_{k+1}\to R_k\) 仅删除满足 \(k_p=k+1\) 的新生坐标，而对既有坐标保持恒等。于是映射
\[
\prod_{p\in\mathcal P_{\mathrm{dy}}}\ZZ/(p^{e_p}\ZZ)
\longrightarrow
\varprojlim_{k\ge 2}R_k,
\qquad
(a_p)_p\longmapsto
\bigl((a_p)_{k_p\le k}\bigr)_{k\ge 2}
\]
是一个拓扑环同构，证毕。
\end{proof}

\begin{corollary}[Cantor 型幂等骨架、非半局部性与非 Noether 性]\label{cor:xi-dyadic-fibonacci-completion-cantor-idempotent-skeleton}
在定理 \ref{thm:xi-dyadic-fibonacci-completion-stable-local-splitting} 的设定下，
\[
\operatorname{Idem}(\mathfrak R_{\mathrm{dy}})
\cong
\{0,1\}^{\mathcal P_{\mathrm{dy}}}.
\]
并且 \(\mathcal P_{\mathrm{dy}}\) 是可数无穷集，因此 \(\operatorname{Idem}(\mathfrak R_{\mathrm{dy}})\) 在乘积拓扑下同胚于经典 Cantor 集。特别地：
\begin{enumerate}
  \item \(\mathfrak R_{\mathrm{dy}}\) 具有 \(2^{\aleph_0}\) 个幂等元；
  \item \(\mathfrak R_{\mathrm{dy}}\) 不是半局部环；
  \item \(\mathfrak R_{\mathrm{dy}}\) 不是 Noether 环。
\end{enumerate}
\end{corollary}
\begin{proof}
记
\[
A_p:=\ZZ/(p^{e_p}\ZZ)
\qquad(p\in\mathcal P_{\mathrm{dy}}).
\]
由定理 \ref{thm:xi-dyadic-fibonacci-completion-stable-local-splitting}，
\[
\mathfrak R_{\mathrm{dy}}\cong \prod_{p\in\mathcal P_{\mathrm{dy}}}A_p.
\]
每个 \(A_p\) 都是有限局部环，故其幂等元只能是 \(0\) 或 \(1\)。因此
\[
\operatorname{Idem}(\mathfrak R_{\mathrm{dy}})
\cong
\prod_{p\in\mathcal P_{\mathrm{dy}}}\operatorname{Idem}(A_p)
\cong
\{0,1\}^{\mathcal P_{\mathrm{dy}}}.
\]

集合 \(\mathcal P_{\mathrm{dy}}\) 可数，因为它是素数集的子集。它又是无限集：当 \(k=3\) 时，\(7\mid F_8\) 且 \(7\nmid F_4\)，故已有一个新生奇素数；而对每个 \(k\ge 4\)，由推论 \ref{cor:xi-leyang-fold-infinite-mean-gamma-platform} 取 \(n=2^k\)，存在素数 \(p\) 满足 \(\rho(p)=2^k\)，于是 \(p\mid F_{2^k}\) 且 \(p\nmid F_{2^j}\)（\(j<k\)），从而 \(p\in\mathcal P_{\mathrm{dy}}\) 且 \(k_p=k\)。故 \(\mathcal P_{\mathrm{dy}}\) 为可数无穷集。任选枚举 \(\mathcal P_{\mathrm{dy}}=\{p_1,p_2,\dots\}\)，便得到
\[
\{0,1\}^{\mathcal P_{\mathrm{dy}}}\cong \{0,1\}^{\mathbb N},
\]
即 Cantor 集的标准模型。这给出枚举中的第一条。

对每个 \(p\in\mathcal P_{\mathrm{dy}}\)，理想
\[
\mathfrak m_p
:=
pA_p\times \prod_{q\in\mathcal P_{\mathrm{dy}}\setminus\{p\}}A_q
\]
的商同构于 \(\FF_p\)，故 \(\mathfrak m_p\) 是极大理想。不同 \(p\) 给出不同的极大理想，因此 \(\mathfrak R_{\mathrm{dy}}\) 至少有可数无穷多个极大理想，不能是半局部环。

最后取互异素数枚举 \(\mathcal P_{\mathrm{dy}}=\{p_1,p_2,\dots\}\)，并定义理想
\[
I_n
:=
\left(\prod_{j\le n}A_{p_j}\right)\times
\left(\prod_{j>n}0\right).
\]
则
\[
I_1\subsetneq I_2\subsetneq I_3\subsetneq \cdots,
\]
因为第 \(n+1\) 个原始坐标幂等元属于 \(I_{n+1}\setminus I_n\)。故 \(\mathfrak R_{\mathrm{dy}}\) 不是 Noether 环。证毕。
\end{proof}

\begin{theorem}[Jacobson 骨架与素数出生账本的半单恢复]\label{thm:xi-dyadic-fibonacci-completion-jacobson-semisimple-recovery}
在定理 \ref{thm:xi-dyadic-fibonacci-completion-stable-local-splitting} 的设定下，设
\[
A_p:=\ZZ/(p^{e_p}\ZZ)
\qquad(p\in\mathcal P_{\mathrm{dy}}).
\]
则
\[
J(\mathfrak R_{\mathrm{dy}})
=
\prod_{p\in\mathcal P_{\mathrm{dy}}}pA_p,
\qquad
\mathfrak R_{\mathrm{dy}}/J(\mathfrak R_{\mathrm{dy}})
\cong
\prod_{p\in\mathcal P_{\mathrm{dy}}}\FF_p.
\]
因而 \(\mathcal P_{\mathrm{dy}}\) 可由半单商本身完全恢复：
\[
p\in\mathcal P_{\mathrm{dy}}
\iff
\FF_p\ \text{作为直因子出现于}\ 
\mathfrak R_{\mathrm{dy}}/J(\mathfrak R_{\mathrm{dy}}).
\]
\end{theorem}
\begin{proof}
由定理 \ref{thm:xi-dyadic-fibonacci-completion-stable-local-splitting}，
\[
\mathfrak R_{\mathrm{dy}}\cong \prod_{p\in\mathcal P_{\mathrm{dy}}}A_p.
\]
每个 \(A_p\) 都是 Artin 局部环，其 Jacobson 根基为 \(pA_p\)，残余域为 \(\FF_p\)。

对任意环族 \((B_i)\)，有
\[
J\!\left(\prod_i B_i\right)=\prod_i J(B_i),
\]
因为 \((b_i)_i\) 在 \(\prod_i B_i\) 中为单位当且仅当每个 \(b_i\) 在 \(B_i\) 中都是单位。故
\[
J(\mathfrak R_{\mathrm{dy}})
=
\prod_{p\in\mathcal P_{\mathrm{dy}}}J(A_p)
=
\prod_{p\in\mathcal P_{\mathrm{dy}}}pA_p.
\]
再按坐标取商得到
\[
\mathfrak R_{\mathrm{dy}}/J(\mathfrak R_{\mathrm{dy}})
\cong
\prod_{p\in\mathcal P_{\mathrm{dy}}}A_p/J(A_p)
\cong
\prod_{p\in\mathcal P_{\mathrm{dy}}}\FF_p.
\]
由于不同素数对应的有限域 \(\FF_p\) 两两不同构，故 \(\FF_p\) 出现为直因子当且仅当 \(p\in\mathcal P_{\mathrm{dy}}\)。证毕。
\end{proof}

与定理 \ref{thm:xi-localized-phase-sampling-closed-form} 的局部化数系情形相比，这里的恢复机制更为刚性：对 \(\ZZ[S^{-1}]\) 而言，素数支撑 \(S\) 需经 Pontryagin 对偶短正合列中的 profinite 核 \(\prod_{p\in S}\ZZ_p\) 才能恢复；而在 dyadic Fibonacci 完成层中，半单商 \(\mathfrak R_{\mathrm{dy}}/J(\mathfrak R_{\mathrm{dy}})\) 已足以恢复全部素数出生支撑 \(\mathcal P_{\mathrm{dy}}\)。

\begin{theorem}[出生过滤与非 \texorpdfstring{$p$}{p}-adic 累厚]\label{thm:xi-dyadic-fibonacci-completion-birth-filtration}
\leanverified{paper\_xi\_dyadic\_fibonacci\_completion\_birth\_filtration}
在定理 \ref{thm:xi-dyadic-fibonacci-completion-stable-local-splitting} 的记号下，定义第 \(k\) 层的出生素数集
\[
\mathcal P_{\mathrm{dy}}^{\mathrm{birth}}(k)
:=
\{p\in\mathcal P_{\mathrm{dy}}:\ k_p=k\}.
\]
则对每个 \(n\ge 2\)，有规范环同构
\[
R_n
\cong
\prod_{2\le k\le n}\ \prod_{p\in\mathcal P_{\mathrm{dy}}^{\mathrm{birth}}(k)}
\ZZ/(p^{e_p}\ZZ),
\]
\[
\mathfrak R_{\mathrm{dy}}
\cong
\prod_{k\ge 2}\ \prod_{p\in\mathcal P_{\mathrm{dy}}^{\mathrm{birth}}(k)}
\ZZ/(p^{e_p}\ZZ).
\]
在此识别下，过渡映射 \(R_{n+1}\to R_n\) 精确等于删除全部满足 \(k_p=n+1\) 的新生坐标。并且对每个 \(k\ge 3\)，都有
\[
\mathcal P_{\mathrm{dy}}^{\mathrm{birth}}(k)\neq \varnothing.
\]
\end{theorem}
\begin{proof}
前两式只是把定理 \ref{thm:xi-dyadic-fibonacci-completion-stable-local-splitting} 中的局部分解按首次出现层级 \(k_p\) 重新分组。由于每个既有素数坐标的指数自出生之后即固定为 \(e_p\)，故过渡映射 \(R_{n+1}\to R_n\) 不会改变任何既有局部因子的指数，只会删除尚未在第 \(n\) 层出生的坐标；这正是所述的坐标删除。

对非空性，\(k=3\) 时有
\[
F_8=21=3\cdot 7,
\]
而 \(3\mid F_4\)、\(7\nmid F_4\)，故 \(7\in\mathcal P_{\mathrm{dy}}^{\mathrm{birth}}(3)\)。若 \(k\ge 4\)，则由推论 \ref{cor:xi-leyang-fold-infinite-mean-gamma-platform} 取 \(n=2^k\)，存在素数 \(p\) 使 \(\rho(p)=2^k\)。于是
\[
p\mid F_{2^k},
\qquad
p\nmid F_{2^j}\quad(j<k),
\]
故 \(p\in\mathcal P_{\mathrm{dy}}^{\mathrm{birth}}(k)\)。证毕。
\end{proof}

因此，dyadic Fibonacci 商塔的完成化并不是沿固定素数方向的 \(p\)-adic 累厚，而是一条按分辨率不断接入新生局部块的出生过滤。

\begin{theorem}[有限局部扇区上的忠实压缩不可能性]\label{thm:xi-dyadic-fibonacci-completion-no-semilocal-embedding}
不存在将 \(\mathfrak R_{\mathrm{dy}}\) 单射嵌入任何交换半局部环的环同态。特别地，不存在单射环同态
\[
\mathfrak R_{\mathrm{dy}}
\hookrightarrow
\prod_{i=1}^{r}B_i,
\]
其中每个 \(B_i\) 都是交换局部环。
\end{theorem}
\begin{proof}
由推论 \ref{cor:xi-dyadic-fibonacci-completion-cantor-idempotent-skeleton}，
\[
\left|\operatorname{Idem}(\mathfrak R_{\mathrm{dy}})\right|=2^{\aleph_0}.
\]
反设 \(f:\mathfrak R_{\mathrm{dy}}\hookrightarrow B\) 为某个交换半局部环 \(B\) 上的单射环同态。设 \(B\) 共有 \(r\) 个极大理想，则
\[
B/J(B)\cong \prod_{i=1}^{r}\kappa_i
\]
为有限个域的直积。交换环上幂等元模去 Jacobson 根基后可唯一提升，因此
\[
\operatorname{Idem}(B)\cong \operatorname{Idem}(B/J(B))
\cong
\{0,1\}^{r},
\]
从而 \(|\operatorname{Idem}(B)|=2^r<\infty\)。

另一方面，单射环同态保持不同元素不同，故不同幂等元在 \(f\) 下仍保持不同。这将迫使 \(B\) 至少含有 \(2^{\aleph_0}\) 个幂等元，与上式矛盾。故此类单射不存在。最后，有限个交换局部环的直积本身即为交换半局部环，故特别情形随即成立。证毕。
\end{proof}

\begin{conjecture}[dyadic Fibonacci 完成层的约化性]\label{conj:xi-dyadic-fibonacci-completion-reducedness}
在定理 \ref{thm:xi-dyadic-fibonacci-completion-stable-local-splitting} 的记号下，对每个 \(p\in\mathcal P_{\mathrm{dy}}\) 都有
\[
e_p=1.
\]
\end{conjecture}

\begin{proposition}[dyadic Fibonacci 完成层约化性的等价表述]\label{prop:xi-dyadic-fibonacci-completion-reduced-equivalences}
在定理 \ref{thm:xi-dyadic-fibonacci-completion-stable-local-splitting} 的记号下，下列四条等价：
\begin{enumerate}
  \item 对每个 \(p\in\mathcal P_{\mathrm{dy}}\)，都有 \(e_p=1\)；
  \item \(\mathfrak R_{\mathrm{dy}}\) 为约化环；
  \item \(J(\mathfrak R_{\mathrm{dy}})=0\)；
  \item
  \[
  \mathfrak R_{\mathrm{dy}}
  \cong
  \prod_{p\in\mathcal P_{\mathrm{dy}}}\FF_p.
  \]
\end{enumerate}
\end{proposition}
\begin{proof}
由定理 \ref{thm:xi-dyadic-fibonacci-completion-stable-local-splitting}，
\[
\mathfrak R_{\mathrm{dy}}
\cong
\prod_{p\in\mathcal P_{\mathrm{dy}}}\ZZ/(p^{e_p}\ZZ).
\]
对单个局部因子 \(A_p=\ZZ/(p^{e_p}\ZZ)\) 而言，以下条件等价：
\begin{enumerate}
  \item \(e_p=1\)；
  \item \(A_p\) 为域，亦即 \(A_p\cong \FF_p\)；
  \item \(A_p\) 为约化环；
  \item \(J(A_p)=pA_p=0\)。
\end{enumerate}
逐坐标应用这一判别便知：全部 \(e_p\) 同时等于 \(1\)，当且仅当整个直积环是约化的；也当且仅当其 Jacobson 根基
\[
J(\mathfrak R_{\mathrm{dy}})
=
\prod_{p\in\mathcal P_{\mathrm{dy}}}p\,\ZZ/(p^{e_p}\ZZ)
\]
为零；又这与直积退化为
\[
\prod_{p\in\mathcal P_{\mathrm{dy}}}\FF_p
\]
完全等价。证毕。
\end{proof}

该猜想把全部非约化障碍压缩到了 dyadic 出现阶支撑 \(\mathcal P_{\mathrm{dy}}\) 的局部指数上：若猜想成立，则 \(\mathfrak R_{\mathrm{dy}}\) 将退化为一个由可数无穷个残余域直积而成的纯残余对象，同时保留推论 \ref{cor:xi-dyadic-fibonacci-completion-cantor-idempotent-skeleton} 所揭示的 Cantor 型幂等骨架。

\endinput


\paragraph{Smith 谱分类、模核 \texorpdfstring{$\zeta$}{zeta} 主因子、整数椭圆原子分解与有限时域/单寄存器账本预算}
前述有限秩账本障碍之后，Smith 信息亏损谱的局部分类、模核计数的 Dirichlet 全局化、整数轴对齐椭圆的原子结构，以及有限时域与单寄存器口径下自然扩张的极小二进制外置预算还可继续压缩为下列结构性结论。

\begin{theorem}[\texorpdfstring{$p$}{p}-进信息亏损谱的完全分类]\label{thm:xi-smith-padic-loss-spectrum-classification}
设 \(f:\ZZ_{\ge 0}\to \ZZ_{\ge 0}\) 满足 \(f(0)=0\)，并对 \(k\ge 1\) 定义前向差分
\[
\Delta(k):=f(k)-f(k-1).
\]
则下列三条等价：
\begin{enumerate}
  \item 存在某个满行秩整数矩阵 \(A\in M_{m\times n}(\ZZ)\) 与素数 \(p\)，使得 \(f=f_p\)，其中 \(f_p\) 为定理 \ref{thm:killo-smith-loss-spectrum} 的 \(\,p\)-进信息亏损谱；
  \item 序列 \((\Delta(k))_{k\ge 1}\) 是非负、非增、有限支撑的整数列；
  \item \(f\) 是整数值、单调不减、离散凹并最终常数的函数。
\end{enumerate}
在这些等价条件成立时，约定对充分大的 \(k\) 有 \(\Delta(k)=0\)，并定义
\[
m_e:=\Delta(e)-\Delta(e+1)\qquad(e\ge 1).
\]
则 \(m_e\in\ZZ_{\ge 0}\)，且只有有限多个 \(m_e\) 非零，并且
\[
f(k)=\sum_{e\ge 1}m_e\min(k,e)\qquad(k\ge 0).
\]
该表示唯一。特别地，\(f\) 所对应的正 \(\,p\)-进 Smith 指数多重集由“指数 \(e\) 重复 \(m_e\) 次”唯一确定。
\leanverified{paper\_xi\_smith\_padic\_loss\_spectrum\_classification}
\end{theorem}
\begin{proof}
若 \(f=f_p\)，则定理 \ref{thm:killo-smith-loss-spectrum} 给出
\[
f(k)=\sum_{i=1}^{m}\min(k,e_i(p)),
\qquad
\Delta(k)=\#\{i:e_i(p)\ge k\}.
\]
对每个固定 \(e\)，函数 \(k\mapsto \min(k,e)\) 单调不减、离散凹且最终常数；其有限和 \(f\) 亦然。并且 \(\Delta(k)\) 显然是非负、非增且有限支撑的整数列，故第 \(1\) 条推出第 \(2\)、\(3\) 条。

第 \(2\) 条与第 \(3\) 条等价只需注意：整数值函数的离散凹性恰等价于前向差分单调不增；单调不减恰等价于前向差分非负；最终常数恰等价于前向差分最终为零。

现证第 \(2\) 条推出第 \(1\) 条。由 \(\Delta\) 非增且有限支撑，
\[
m_e=\Delta(e)-\Delta(e+1)\ge 0,
\]
并且只有有限多个 \(m_e\) 非零。对任意 \(k\ge 1\)，望远镜求和给出
\[
\sum_{e\ge k}m_e
=
\sum_{e\ge k}\bigl(\Delta(e)-\Delta(e+1)\bigr)
=
\Delta(k).
\]
于是
\[
f(k)
=
\sum_{j=1}^{k}\Delta(j)
=
\sum_{j=1}^{k}\sum_{e\ge j}m_e
=
\sum_{e\ge 1}m_e\sum_{j=1}^{k}\mathbf 1_{\{e\ge j\}}
=
\sum_{e\ge 1}m_e\min(k,e).
\]
若所有 \(m_e\) 都为零，则 \(f\equiv 0\)，取 \(A=(1)\) 即有 \(f=f_p\equiv 0\)。否则记
\[
r:=\sum_{e\ge 1}m_e\ge 1,
\]
并把每个整数 \(e\) 按重数 \(m_e\) 排成序列 \(e_1,\dots,e_r\)。固定任意素数 \(p\)，取对角矩阵
\[
A=\operatorname{diag}(p^{e_1},\dots,p^{e_r})\in M_r(\ZZ).
\]
由定理 \ref{thm:killo-smith-loss-spectrum}，
\[
f_p(k)=\sum_{\nu=1}^{r}\min(k,e_\nu)=\sum_{e\ge 1}m_e\min(k,e)=f(k),
\]
故第 \(1\) 条成立。

最后，唯一性由
\[
\Delta(k)=\sum_{e\ge k}m_e
\]
立即推出
\[
m_e=\Delta(e)-\Delta(e+1).
\]
因此正 \(\,p\)-进 Smith 指数多重集也随之唯一确定。证毕。
\end{proof}

\begin{corollary}[Smith 指数的二阶离散曲率恢复律]\label{cor:xi-smith-second-discrete-curvature-recovery}
\leanverified{paper\_xi\_smith\_second\_discrete\_curvature\_recovery}
沿用定理 \ref{thm:killo-smith-loss-spectrum} 的记号，并定义
\[
f_p(0):=0,
\qquad
\Delta_p(k):=f_p(k)-f_p(k-1)\qquad(k\ge 1).
\]
则对每个 \(k\ge 1\) 都有
\[
\#\{i:e_i(p)=k\}
=
\Delta_p(k)-\Delta_p(k+1)
=
2f_p(k)-f_p(k-1)-f_p(k+1).
\]
因此下列三个局部不变量均可由 \(f_p\) 无损恢复：
\[
r_p:=\#\{i:e_i(p)\ge 1\}=f_p(1),
\]
\[
\tau_p:=\max\{k:\Delta_p(k)>0\}=\max_i e_i(p),
\]
\[
\sigma_p:=\lim_{k\to\infty}f_p(k)=\sum_{i=1}^{m}e_i(p)=v_p\!\Bigl(\prod_{i=1}^{m}d_i\Bigr).
\]
若进一步 \(A\) 为方阵满秩，则对一切 \(k\ge \tau_p\) 都有
\[
\#\ker(A\bmod p^k)=p^{\sigma_p},
\]
即模 \(p^k\) 的核大小在阈值 \(\tau_p\) 之后精确稳定为 \(|\det A|\) 的 \(\,p\)-进部分。
\end{corollary}
\begin{proof}
定理 \ref{thm:killo-smith-loss-spectrum} 给出
\[
\Delta_p(k)=\#\{i:e_i(p)\ge k\}.
\]
于是
\[
\Delta_p(k)-\Delta_p(k+1)
=
\#\{i:e_i(p)\ge k\}-\#\{i:e_i(p)\ge k+1\}
=
\#\{i:e_i(p)=k\}.
\]
再由
\[
\Delta_p(k)-\Delta_p(k+1)
=
\bigl(f_p(k)-f_p(k-1)\bigr)-\bigl(f_p(k+1)-f_p(k)\bigr),
\]
即得二阶离散曲率公式。

三个局部不变量随即成立。第一式取 \(k=1\) 即得；第二式来自 \(\Delta_p(k)>0\iff \exists\,i,\ e_i(p)\ge k\)；第三式因为 \(f_p(k)=\sum_i\min(k,e_i(p))\) 最终稳定于 \(\sum_i e_i(p)\)。

若 \(A\) 为方阵满秩，则 \(m=n\)。对 \(k\ge \tau_p\)，一切 \(i\) 都满足 \(\min(k,e_i(p))=e_i(p)\)，故
\[
f_p(k)=\sum_{i=1}^{n}e_i(p)=\sigma_p.
\]
由定理 \ref{thm:killo-smith-loss-spectrum} 中
\[
K_p(k)=\#\ker(A\bmod p^k)=p^{\,k(n-m)+f_p(k)},
\]
便得
\[
\#\ker(A\bmod p^k)=p^{\sigma_p}.
\]
又因 \(\prod_{i=1}^{n}d_i=|\det A|\)，该稳定值正是 \(|\det A|\) 的 \(\,p\)-进部分。证毕。
\end{proof}

\input{subsubsec__xi-time-protocol-conclusions-part62ad-smith-pwindow-determinant-index}

\input{subsubsec__xi-time-protocol-conclusions-part62aa-smith-conjugate-volume-platform}

\begin{theorem}[模核 Dirichlet 生成函数的 \texorpdfstring{$\zeta$}{zeta} 主因子与有限 Euler 修正]\label{thm:xi-kernel-dirichlet-zeta-finite-euler-correction}
\leanverified{paper\_xi\_kernel\_dirichlet\_zeta\_finite\_euler\_correction}
设 \(A\in M_{m\times n}(\ZZ)\) 满足 \(\operatorname{rank}_{\QQ}(A)=m\)，其 Smith 不变量为
\[
d_1\mid d_2\mid\cdots\mid d_m.
\]
对每个 \(b\ge 1\)，记
\[
A_b:(\ZZ/b\ZZ)^n\longrightarrow(\ZZ/b\ZZ)^m
\]
为模 \(b\) 的诱导映射，并定义归一化模核计数
\[
\kappa_A(b):=\frac{\#\ker(A_b)}{b^{\,n-m}}
=
\prod_{i=1}^{m}\gcd(d_i,b).
\]
再定义 Dirichlet 生成函数
\[
Z_A(s):=\sum_{b\ge 1}\frac{\kappa_A(b)}{b^s},
\qquad
\Re s>1.
\]
令
\[
\Delta_A:=\prod_{i=1}^{m}d_i,
\]
并对每个素数 \(p\mid \Delta_A\) 记
\[
e_i(p):=v_p(d_i),\qquad
E_p:=\max_i e_i(p),\qquad
\sigma_p:=\sum_{i=1}^{m}e_i(p),
\]
\[
f_p(k):=\sum_{i=1}^{m}\min(k,e_i(p)).
\]
则成立精确 Euler 分解
\[
Z_A(s)=\zeta(s)\prod_{p\mid \Delta_A}P_{A,p}(p^{-s}),
\]
其中
\[
P_{A,p}(T)
:=
(1-T)\sum_{k=0}^{E_p-1}p^{f_p(k)}T^k
+p^{\sigma_p}T^{E_p}
\in\ZZ[T].
\]
特别地，\(Z_A(s)\) 在复平面上亚纯延拓，仅在 \(s=1\) 具有一个单极点，并且
\[
\operatorname*{Res}_{s=1}Z_A(s)
=
\prod_{p\mid \Delta_A}P_{A,p}(p^{-1}).
\]
\end{theorem}
\begin{proof}
推论 \ref{cor:killo-kernel-size-general-modulus} 已给出
\[
\kappa_A(b)=\prod_{i=1}^{m}\gcd(d_i,b).
\]
若 \((b_1,b_2)=1\)，则
\[
\gcd(d_i,b_1b_2)=\gcd(d_i,b_1)\gcd(d_i,b_2),
\]
从而 \(\kappa_A\) 为乘法函数。故在 \(\Re s>1\) 上，
\[
Z_A(s)=\prod_{p}Z_{A,p}(s),
\qquad
Z_{A,p}(s):=\sum_{k\ge 0}\frac{\kappa_A(p^k)}{p^{ks}}.
\]

再由同一推论与定理 \ref{thm:killo-smith-loss-spectrum}，
\[
\kappa_A(p^k)=\prod_{i=1}^{m}p^{\min(k,e_i(p))}=p^{f_p(k)}.
\]
若 \(k\ge E_p\)，则 \(\min(k,e_i(p))=e_i(p)\) 对一切 \(i\) 都成立，故 \(f_p(k)=\sigma_p\)。记 \(T:=p^{-s}\)，则
\[
Z_{A,p}(s)
=
\sum_{k=0}^{E_p-1}p^{f_p(k)}T^k+\sum_{k\ge E_p}p^{\sigma_p}T^k
=
\sum_{k=0}^{E_p-1}p^{f_p(k)}T^k+\frac{p^{\sigma_p}T^{E_p}}{1-T}.
\]
整理可得
\[
Z_{A,p}(s)=\frac{P_{A,p}(T)}{1-T}.
\]

当 \(p\nmid \Delta_A\) 时，一切 \(e_i(p)=0\)，于是 \(\kappa_A(p^k)=1\) 对所有 \(k\) 都成立，故
\[
Z_{A,p}(s)=\sum_{k\ge 0}p^{-ks}=\frac{1}{1-p^{-s}}.
\]
将这些素数上的因子合并，便得到
\[
Z_A(s)
=
\prod_{p\nmid \Delta_A}\frac{1}{1-p^{-s}}
\prod_{p\mid \Delta_A}\frac{P_{A,p}(p^{-s})}{1-p^{-s}}
=
\zeta(s)\prod_{p\mid \Delta_A}P_{A,p}(p^{-s}).
\]
右端是 \(\zeta(s)\) 乘以有限个整函数因子，因此仅在 \(s=1\) 保留 \(\zeta(s)\) 的单极点。又 \(\operatorname*{Res}_{s=1}\zeta(s)=1\)，故
\[
\operatorname*{Res}_{s=1}Z_A(s)
=
\prod_{p\mid \Delta_A}P_{A,p}(p^{-1}).
\]
证毕。
\end{proof}

\begin{theorem}[整数轴对齐椭圆半群的自由交换原子分解]\label{thm:xi-integer-ellipse-free-commutative-monoid}
\leanverified{paper\_xi\_integer\_ellipse\_free\_commutative\_monoid}
沿用定理 \ref{thm:killo-ellipse-diagonal-prime-register-equivalence} 的记号，记
\[
\mathcal A:=\{E_{p,1},E_{1,p}:p\in\mathcal P\}\subset \mathsf E_{\NN}.
\]
则：
\begin{enumerate}
  \item \((\mathsf E_{\NN},\odot)\) 是以 \(\mathcal A\) 为原子集的自由交换幺半群；
  \item 若约定 \(M^{\odot r}\) 表示 \(r\) 次 \(\odot\)-乘积，并取 \(M^{\odot 0}=E_{1,1}\)，则每个 \(E_{a,b}\in\mathsf E_{\NN}\) 都具有唯一原子分解
  \[
  E_{a,b}
  =
  \left(\bigodot_{p\in\mathcal P}E_{p,1}^{\odot v_p(a)}\right)
  \odot
  \left(\bigodot_{p\in\mathcal P}E_{1,p}^{\odot v_p(b)}\right),
  \]
  其中仅有限多个因子不同于单位元 \(E_{1,1}\)；
  \item 其幺半群自同构群满足
  \[
  \operatorname{Aut}_{\mathrm{Mon}}(\mathsf E_{\NN},\odot)
  \cong
  \mathrm{Sym}(\mathcal A)
  \cong
  \mathrm{Sym}(\mathcal P\times\{x,y\}).
  \]
\end{enumerate}
\end{theorem}
\begin{proof}
定理 \ref{thm:killo-ellipse-diagonal-prime-register-equivalence} 已给出幺半群同构
\[
(\NN^{(\mathcal P)}\times \NN^{(\mathcal P)},+)
\cong
(\mathsf E_{\NN},\odot),\qquad
(r_1,r_2)\longmapsto E_{N(r_1),N(r_2)}.
\]
设 \(\delta_p\in\NN^{(\mathcal P)}\) 为在素数 \(p\) 处取值 \(1\)、其余处取值 \(0\) 的基向量。则
\[
(\delta_p,0)\longmapsto E_{p,1},
\qquad
(0,\delta_p)\longmapsto E_{1,p}.
\]
而 \((\NN^{(\mathcal P)}\times \NN^{(\mathcal P)},+)\) 正是以这些基向量为自由基的交换幺半群，因此其原子（不可约元）恰为
\[
\{(\delta_p,0),(0,\delta_p):p\in\mathcal P\}.
\]
经同构传回，便得到第 \(1\) 条。

任取
\[
a=\prod_{p\in\mathcal P}p^{v_p(a)},
\qquad
b=\prod_{p\in\mathcal P}p^{v_p(b)},
\]
则在自由基坐标下，
\[
(r_1,r_2)=\sum_{p}v_p(a)(\delta_p,0)+\sum_{p}v_p(b)(0,\delta_p).
\]
把该式送回 \((\mathsf E_{\NN},\odot)\) 便得到所示原子分解。由于自由交换幺半群中的基坐标唯一，故该分解唯一，第 \(2\) 条成立。

最后，幺半群自同构保持不可约元，因此必将原子集 \(\mathcal A\) 置换到自身。反之，任一 \(\mathcal A\) 上的置换都可在自由交换幺半群上按乘法唯一延拓为自同构。于是
\[
\operatorname{Aut}_{\mathrm{Mon}}(\mathsf E_{\NN},\odot)
\cong
\mathrm{Sym}(\mathcal A).
\]
再由
\[
\mathcal A=\{E_{p,1},E_{1,p}:p\in\mathcal P\}
\cong
\mathcal P\times\{x,y\},
\]
即得最后一个同构。证毕。
\end{proof}

\begin{corollary}[保持面积的整数轴对齐椭圆幺半群自同构仅由各素数轴交换组成]\label{cor:xi-integer-ellipse-area-preserving-automorphism-product}
沿用定理 \ref{thm:xi-integer-ellipse-free-commutative-monoid} 的记号，并定义面积函数
\[
\operatorname{Area}(E_{a,b}):=\pi ab.
\]
记
\[
\operatorname{Aut}_{\mathrm{Mon},\operatorname{Area}}(\mathsf E_{\NN},\odot)
:=
\left\{
\Phi\in \operatorname{Aut}_{\mathrm{Mon}}(\mathsf E_{\NN},\odot):
\operatorname{Area}(\Phi(E))=\operatorname{Area}(E)\ \forall E\in\mathsf E_{\NN}
\right\}.
\]
则有规范同构
\[
\operatorname{Aut}_{\mathrm{Mon},\operatorname{Area}}(\mathsf E_{\NN},\odot)
\cong
\prod_{p\in\mathcal P}(\ZZ/2\ZZ),
\]
其中第 \(p\) 个因子由两原子
\[
E_{p,1}\longleftrightarrow E_{1,p}
\]
的交换生成。
\end{corollary}
\begin{proof}
由定理 \ref{thm:xi-integer-ellipse-free-commutative-monoid}，任意幺半群自同构都唯一对应于原子集
\[
\mathcal A=\{E_{p,1},E_{1,p}:p\in\mathcal P\}
\]
上的一个置换。若 \(\Phi\) 还保持面积逐点不变，则对每个素数 \(p\) 有
\[
\operatorname{Area}(E_{p,1})=\operatorname{Area}(E_{1,p})=\pi p.
\]
而当 \(q\neq p\) 时，
\[
\operatorname{Area}(E_{q,1})=\operatorname{Area}(E_{1,q})=\pi q\neq \pi p.
\]
故 \(\Phi(E_{p,1})\) 只能取 \(E_{p,1}\) 或 \(E_{1,p}\)，同理 \(\Phi(E_{1,p})\) 亦只能取该二元组中的另一项。于是 \(\Phi\) 必逐素数保持二点块
\[
\{E_{p,1},E_{1,p}\}
\]
并在每个块上仅有“恒等”与“交换”两种可能。

反之，对每个 \(p\) 独立选择是否交换 \(E_{p,1}\) 与 \(E_{1,p}\)，便得到原子集 \(\mathcal A\) 上的一个置换；由定理 \ref{thm:xi-integer-ellipse-free-commutative-monoid}，该置换唯一延拓为 \((\mathsf E_{\NN},\odot)\) 的幺半群自同构。由于这种延拓保持每个素数对 \(\{E_{p,1},E_{1,p}\}\) 的面积 \(\pi p\)，故它逐点保持全部椭圆的面积。
更具体地，若
\[
E_{a,b}
=
\left(\bigodot_{p}E_{p,1}^{\odot \alpha_p}\right)
\odot
\left(\bigodot_{p}E_{1,p}^{\odot \beta_p}\right),
\]
则任意这类自同构只会把每个固定 \(p\) 上的 \(\alpha_p,\beta_p\) 交换或保持不变，因此
\[
v_p(ab)=\alpha_p+\beta_p
\]
逐素数保持不变，从而
\[
\operatorname{Area}(E_{a,b})=\pi ab
\]
逐点保持不变。

因此，保持面积的自同构群恰为各素数上独立交换所成的直积
\[
\prod_{p\in\mathcal P}(\ZZ/2\ZZ).
\]
证毕。
\end{proof}

\begin{theorem}[有限时域自然扩张的最小二进制外置预算]\label{thm:xi-finite-time-natural-extension-active-branch-budget}
\leanverified{paper\_xi\_finite\_time\_natural\_extension\_active\_branch\_budget}
沿用定理 \ref{thm:xi-layered-prime-register-sharp-natural-extension} 的分层外置语义，但把第 \(j\) 层的 prime-slice 通道替换为有限字母表 \(A_j\)。设
\[
X_0\xrightarrow{\pi_0}X_1\xrightarrow{\pi_1}\cdots\xrightarrow{\pi_{N-1}}X_N
\]
为有限时域链，并假设每层纤维满足
\[
\#\pi_j^{-1}(y)\le 2
\qquad
(\forall\,y\in X_{j+1},\ 0\le j\le N-1).
\]
定义活跃分支层集合
\[
B_N:=\{\,j\in\{0,\dots,N-1\}:\exists\,y\in X_{j+1}\text{ 使 }\#\pi_j^{-1}(y)=2\,\}.
\]
考虑任意分层外置编码
\[
\Psi(x_0)=\bigl(x_N,h_0(x_0),\dots,h_{N-1}(x_0)\bigr)\in X_N\times\prod_{j=0}^{N-1}A_j
\]
并假设 \(\Psi\) 为单射。则对每个 \(j\in B_N\) 都必有
\[
|A_j|\ge 2.
\]
因此总二进制预算满足
\[
\sum_{j=0}^{N-1}\left\lceil \log_2|A_j|\right\rceil\ge |B_N|.
\]
反之，存在单射分层外置使得
\[
|A_j|=
\begin{cases}
2,& j\in B_N,\\
1,& j\notin B_N.
\end{cases}
\]
故最小总二进制外置预算恰等于 \(|B_N|\)。
\end{theorem}
\begin{proof}
对任意活跃层 \(j\in B_N\)，存在 \(y\in X_{j+1}\) 使 \(\#\pi_j^{-1}(y)=2\)。推论 \ref{cor:killo-layered-necessity-minimality} 的论证对该单层可逐字局部化：若第 \(j\) 层通道退化为单值，则可以构造两条具有相同终端宏态、仅在第 \(j\) 层分支选择不同的历史，使得分层外置无法保持单射。故每个活跃层都必须配备非平凡字母表，亦即 \(|A_j|\ge 2\)。对所有 \(j\in B_N\) 求和便得
\[
\sum_{j=0}^{N-1}\left\lceil \log_2|A_j|\right\rceil\ge |B_N|.
\]

反向上，对每个 \(j\in B_N\) 与每个 \(y\in X_{j+1}\)，由于 \(\#\pi_j^{-1}(y)\le 2\)，可取单射
\[
\iota_{j,y}:\pi_j^{-1}(y)\hookrightarrow \{0,1\}.
\]
若 \(j\notin B_N\)，则每条纤维都是单点集，可唯一地取
\[
\iota_{j,y}:\pi_j^{-1}(y)\to \{0\}.
\]
对任意 \(x_0\in X_0\)，递推定义
\[
x_{j+1}:=\pi_j(x_j)\qquad(0\le j\le N-1),
\]
并设
\[
h_j(x_0):=\iota_{j,x_{j+1}}(x_j)\in A_j,
\qquad
A_j:=
\begin{cases}
\{0,1\},& j\in B_N,\\
\{0\},& j\notin B_N.
\end{cases}
\]
于是
\[
\Psi(x_0)=\bigl(x_N,h_0(x_0),\dots,h_{N-1}(x_0)\bigr).
\]
给定 \(\Psi(x_0)\)，可由 \(x_N\) 与 \(h_{N-1}\) 先恢复唯一的 \(x_{N-1}\)，再由 \(x_{N-1}\) 与 \(h_{N-2}\) 恢复唯一的 \(x_{N-2}\)，如此逐层逆推，最终恢复唯一的 \(x_0\)。故 \(\Psi\) 单射。

该构造在每个活跃层恰使用一位二进制字母，在每个非活跃层使用平凡单字母，因此总预算正好等于 \(|B_N|\)。证毕。
\end{proof}

\begin{proposition}[自然扩张单寄存器实现的最小字母表与二进长度]\label{prop:xi-natural-extension-single-register-minimal-alphabet}
\leanverified{paper\_xi\_natural\_extension\_single\_register\_minimal\_alphabet}
设 \(T:X\twoheadrightarrow X\) 为有限纤维满射，并记
\[
M:=\sup_{y\in X}\#T^{-1}(y)<\infty.
\]
考虑定理 \ref{thm:killo-natural-extension-branch-register} 中把自然扩张编码为右移插入系统的单寄存器模型：取有限字母表 \(\mathcal A\)，并对每个 \(y\in X\) 选取单射
\[
\iota_y:T^{-1}(y)\hookrightarrow \mathcal A,
\qquad
\beta(x):=\iota_{T(x)}(x).
\]
则下列两条等价：
\begin{enumerate}
  \item 对全部 \(y\in X\)，上述纤维编码可同时取到；
  \item \(|\mathcal A|\ge M\)。
\end{enumerate}
因此，若记单个历史符号所需的最小二进制寄存器长度为
\[
L_{\min}(T),
\]
则有精确公式
\[
\boxed{\;
L_{\min}(T)=\left\lceil \log_2 M \right\rceil.\;}
\]
并且该下界可达：取
\[
\mathcal A=\{0,1,\dots,M-1\}
\]
并在每条纤维上任选单射 \(\iota_y\)，即可得到与定理 \ref{thm:killo-natural-extension-branch-register} 同型的右移插入双射模型而不丢失任何前史信息。
\end{proposition}
\begin{proof}
若存在某个 \(y\in X\) 使 \(\#T^{-1}(y)=M\)，而 \(|\mathcal A|<M\)，则不存在单射
\[
T^{-1}(y)\hookrightarrow \mathcal A.
\]
于是无法仅凭单一历史符号 \(\beta(x)\in\mathcal A\) 无歧义地区分该纤维内的全部前史分支。故必要条件为
\[
|\mathcal A|\ge M.
\]

反过来，若 \(|\mathcal A|\ge M\)，则对每个 \(y\in X\) 都可选取单射
\[
\iota_y:T^{-1}(y)\hookrightarrow \mathcal A.
\]
由此定义
\[
\beta(x):=\iota_{T(x)}(x)
\]
后，完全沿定理 \ref{thm:killo-natural-extension-branch-register} 的递推构造，即得编码
\[
C:X^{\leftarrow}\longrightarrow X\times \mathcal A^{\NN},
\qquad
C\bigl((x_0,x_{-1},x_{-2},\dots)\bigr)
=
\bigl(x_0,(\beta(x_{-1}),\beta(x_{-2}),\dots)\bigr),
\]
其为到像上的双射，并把自然扩张左移共轭为右移插入变换
\[
\widehat T(x,(a_1,a_2,\dots))=(T(x),(\beta(x),a_1,a_2,\dots)).
\]
故充分性也成立。

于是最小可行字母表大小恰为 \(M\)。把 \(M\) 个字母固定编码为二进制串所需的最小位长正是
\[
\left\lceil \log_2 M \right\rceil,
\]
故
\[
L_{\min}(T)=\left\lceil \log_2 M \right\rceil.
\]
证毕。
\end{proof}

\input{subsubsec__xi-time-protocol-conclusions-part62ac-natural-extension-minimal-d-alphabet}

\input{subsubsec__xi-time-protocol-conclusions-part62ab-smith-dirichlet-torsion-grothendieck}

\input{subsubsec__xi-time-protocol-conclusions-part62a-zg-branch-tower-field-separation}
\input{subsubsec__xi-time-protocol-conclusions-part62b-rankone-hologram-leyang-matrix-euler}
\input{subsubsec__xi-time-protocol-conclusions-part62b-dyadic-multiplicity-lp-obstruction}
\input{subsubsec__xi-time-protocol-conclusions-part62c-godel-tate-shift-ahlfors-oracle}
\input{subsubsec__xi-time-protocol-conclusions-part62ca-holographic-grothendieck-leyang-interface-rigidity}
\input{subsubsec__xi-time-protocol-conclusions-part62d-solenoid-rose-lissajous-prime-kernel}
\input{subsubsec__xi-time-protocol-conclusions-part62d-unitcircle-singular-ring-resultant-rh}
\input{subsubsec__xi-time-protocol-conclusions-part62d-grothendieck-ledger-hom-matrix}
\input{subsubsec__xi-time-protocol-conclusions-part62d-golden-shift-smith-periodic-finitization}
\input{subsubsec__xi-time-protocol-conclusions-part62d-canonical-monoid-tauberian-torsor}
\input{subsubsec__xi-time-protocol-conclusions-part62dc-hologram-bernoulli-exactdimension}
\input{subsubsec__xi-time-protocol-conclusions-part62dca-hologram-ifs-fixedpoint-exactdimension-ahlfors}
\input{subsubsec__xi-time-protocol-conclusions-part62dcb-hologram-ambient-riccati-piecewiseaffine}
\input{subsubsec__xi-time-protocol-conclusions-part62dd-addressable-godel-source-majorization}
\input{subsubsec__xi-time-protocol-conclusions-part62d-leyang-branch-oracle-wreath-gaussian}
\input{subsubsec__xi-time-protocol-conclusions-part62d-leyang-branch-graph-closures}
\input{subsubsec__xi-time-protocol-conclusions-part62dh-leyang-wreath-tate-density-closure}
\input{subsubsec__xi-time-protocol-conclusions-part62d-leyang-chebotarev-fiber-binary-rigidity}
\input{subsubsec__xi-time-protocol-conclusions-part62di-tripleclass-orthogonality-sixstep-renormalization}
\input{subsubsec__xi-time-protocol-conclusions-part62e-chebotarev-signcube-fullsplit-dilution}
\input{subsubsec__xi-time-protocol-conclusions-part62d-zg-rh-witt-normalform-rigidity}
\input{subsubsec__xi-time-protocol-conclusions-part62da-golden-cylinder-smith-localization-pontryagin}
\input{subsubsec__xi-time-protocol-conclusions-part62db-fiberphase-localizedcover-saturation}
\input{subsubsec__xi-time-protocol-conclusions-part62de-hologram-smith-coroot-closure}
\input{subsubsec__xi-time-protocol-conclusions-part62deb-hologram-fullshift-primitive-zeta-closure}
\input{subsubsec__xi-time-protocol-conclusions-part62debb-finite-localization-axis-obstruction}
\input{subsubsec__xi-time-protocol-conclusions-part62dea-localized-natural-extension-b3c3-rootpolytopes}
\input{subsubsec__xi-time-protocol-conclusions-part62df-fiber-capacity-primeshadow-cubic-synchrony}
\input{subsubsec__xi-time-protocol-conclusions-part62dg-leyang-cylinder-phase-activation-separation}
\input{subsubsec__xi-time-protocol-conclusions-part62dga-activated-escape-deflation-partition}
\input{subsubsec__xi-time-protocol-conclusions-part62dgb-hologram-cylinder-smith-singlepole}
\input{subsubsec__xi-time-protocol-conclusions-part62dgba-smith-kernel-finite-euler-leading-constant}
\input{subsubsec__xi-time-protocol-conclusions-part62dgc-zg-matrix-euler-fivecore-collapse}
\input{subsubsec__xi-time-protocol-conclusions-part62dgd-zg-jfraction-golden-cantor-separation}
\input{subsubsec__xi-time-protocol-conclusions-part63-schur-ghost-crt-hook-pressure}
\input{subsubsec__xi-time-protocol-conclusions-part63b-jet-gap-prime-smith-ellipse-rigidity}
\input{subsubsec__xi-time-protocol-conclusions-part64-peter-weyl-artin-cover-zeta-poles}

\endinput

\paragraph{Hankel 缺口商环、除子层级与高阶异常的交互阶刚性}
在过拟合递推的 Hankel--null 缺口已被识别为自由阿贝尔群，而常数层高阶交叉异常的 Möbius 分解与 moment 放大律也已闭合之后，还可把这两条链进一步压缩为一组更强的结构定理：前者给出原始缺口模的规范商环模型及其中间层级的除子格分类，后者给出多轴异常的最小交互阶刻画及其在 moment 编译下的严格保持。

\begin{theorem}[过拟合 Hankel 缺口模的规范商环模型]\label{thm:xi-time-part62af-hankel-gap-quotient-ring-model}
\leanverified{paper\_xi\_time\_part62af\_hankel\_gap\_quotient\_ring\_model}
沿用定理 \ref{thm:pom-hankel-syndrome-gap-rank-defect} 的记号。设
\[
P(\lambda)=P_{\min}(\lambda)\,R(\lambda),
\qquad
d_{\min}:=\deg P_{\min},
\qquad
\Delta:=\deg R,
\]
其中 \(P_{\min},R\in\ZZ[\lambda]\) 均首一，且 \(P_{\min}\) 为从偏移 \(m_0\) 起的首一最小湮灭多项式。
对任意 \(n\ge \deg P+1\)，记
\[
\mathsf Q_n(P;s)
:=
\ker_{\ZZ}\bigl((\mathsf H^{(m_0)}_n)^{\mathsf T}\bigr)\big/\mathcal M_n(P).
\]
则映射
\[
\Theta_R:\ZZ[\lambda]_{<\Delta}\longrightarrow \mathsf Q_n(P;s),\qquad
U(\lambda)\longmapsto
\Bigl[\operatorname{coeff}_{<n}\!\bigl(U(\lambda)P_{\min}(\lambda)\bigr)\Bigr]
\]
为同构。经由首一多项式 \(R\) 的唯一余式代表，仍记其诱导同构为
\[
\Theta_R:\ZZ[\lambda]/(R)\xrightarrow{\sim}\mathsf Q_n(P;s),
\]
并有规范识别
\[
\boxed{\;
\mathsf Q_n(P;s)\ \cong\ \ZZ[\lambda]/(R).\;}
\]
若记 \(\bar\sigma\) 为截断移位在 \(\mathsf Q_n(P;s)\) 上诱导的端同态，则在 \(\Theta_R\) 下，
\[
\bar\sigma
\quad\longleftrightarrow\quad
\mathsf T_R:\ZZ[\lambda]_{<\Delta}\to \ZZ[\lambda]_{<\Delta},
\qquad
\mathsf T_R(U):=\operatorname{rem}(\lambda U,R).
\]
因此 \(\bar\sigma\otimes_{\ZZ}\QQ\) 的特征多项式恰为 \(R(\lambda)\)。
\end{theorem}
\begin{proof}
由定理 \ref{thm:pom-hankel-syndrome-gap-rank-defect}，
\[
\mathsf Q_n(P;s)
\cong
\ZZ[\lambda]_{<n-d_{\min}}\Big/ R\cdot \ZZ[\lambda]_{<n-\deg P}.
\]
任取 \(A(\lambda)\in \ZZ[\lambda]_{<n-d_{\min}}\)。由于 \(R\) 首一，带余除法给出唯一分解
\[
A(\lambda)=R(\lambda)Q(\lambda)+U(\lambda),
\qquad
\deg U<\Delta.
\]
又因
\[
\deg A<n-d_{\min}=(n-\deg P)+\Delta,
\]
故必有 \(\deg Q<n-\deg P\)。因此上述商群的每个商类都由唯一余式 \(U\in \ZZ[\lambda]_{<\Delta}\) 代表，从而得到 \(\Theta_R\) 的双射性；这也正是
\(
\mathsf Q_n(P;s)\cong \ZZ[\lambda]/(R)
\)
的规范实现。

再看移位作用。把 \(\ZZ^n\) 与 \(\ZZ[\lambda]_{<n}\) 以系数向量同一，并令
\[
\sigma\bigl(\operatorname{coeff}_{<n}(A)\bigr):=
\operatorname{coeff}_{<n}(\lambda A).
\]
由于 \(\mathcal M_n(P_{\min})\) 与 \(\mathcal M_n(P)\) 都是截断倍乘子模，二者均对 \(\sigma\) 不变，故 \(\sigma\) 在商模 \(\mathsf Q_n(P;s)\) 上诱导出 \(\bar\sigma\)。对任意 \(U\in\ZZ[\lambda]_{<\Delta}\)，有
\[
\bar\sigma\bigl(\Theta_R(U)\bigr)
=
\Bigl[\operatorname{coeff}_{<n}\!\bigl(\lambda U P_{\min}\bigr)\Bigr].
\]
把 \(\lambda U\) 对 \(R\) 作带余除法，写成
\[
\lambda U=R Q+\operatorname{rem}(\lambda U,R).
\]
则
\[
\lambda U P_{\min}
=
Q P + \operatorname{rem}(\lambda U,R)\,P_{\min},
\]
其中 \(Q P\) 的截断系数向量落在 \(\mathcal M_n(P)\) 中，故在商模内消失，从而
\[
\bar\sigma\bigl(\Theta_R(U)\bigr)
=
\Theta_R\!\bigl(\operatorname{rem}(\lambda U,R)\bigr).
\]
这就证明了 \(\bar\sigma\) 与 \(\mathsf T_R\) 的共轭关系。最后，在基
\(
1,\lambda,\dots,\lambda^{\Delta-1}
\)
下，\(\mathsf T_R\otimes_{\ZZ}\QQ\) 的矩阵正是 \(R\) 的 Frobenius 伴随块，故其特征多项式恰为 \(R\)。证毕。
\end{proof}

\begin{theorem}[过拟合综合征层级的除子格分类]\label{thm:xi-time-part62af-hankel-gap-divisor-lattice}
\leanverified{paper\_xi\_time\_part62af\_hankel\_gap\_divisor\_lattice}
沿用定理 \ref{thm:xi-time-part62af-hankel-gap-quotient-ring-model} 的记号。对每个首一因子 \(D\mid R\)，定义中间综合征子模
\[
\mathcal N_n(D):=\mathcal M_n(P_{\min}D)
\subseteq
\ker_{\ZZ}\bigl((\mathsf H^{(m_0)}_n)^{\mathsf T}\bigr).
\]
则对任意首一因子 \(D_1,D_2\mid R\)，有
\[
\boxed{\;
D_1\mid D_2
\quad\Longleftrightarrow\quad
\mathcal N_n(D_2)\subseteq \mathcal N_n(D_1),\;}
\]
以及
\[
\boxed{\;
\mathcal N_n(D_1)\cap \mathcal N_n(D_2)
=
\mathcal N_n\bigl(\operatorname{lcm}(D_1,D_2)\bigr).\;}
\]
此外，在定理 \ref{thm:xi-time-part62af-hankel-gap-quotient-ring-model} 的同构
\[
\Theta_R:\ZZ[\lambda]/(R)\xrightarrow{\sim}\mathsf Q_n(P;s)
\]
下，有精确识别
\[
\boxed{\;
\Theta_R^{-1}\!\bigl(\mathcal N_n(D)/\mathcal M_n(P)\bigr)
=
(D)/(R)\subseteq \ZZ[\lambda]/(R).\;}
\]
因而
\[
\boxed{\;
\rank_{\ZZ}\bigl(\mathcal N_n(D)/\mathcal M_n(P)\bigr)
=
\Delta-\deg D,\;}
\qquad
\boxed{\;
\rank_{\ZZ}\Bigl(\mathsf Q_n(P;s)\big/\bigl(\mathcal N_n(D)/\mathcal M_n(P)\bigr)\Bigr)
=
\deg D.\;}
\]
\end{theorem}
\begin{proof}
由命题 \ref{prop:pom-hankel-syndrome-divisibility-reversal-lcm}，对
\[
P_1=P_{\min}D_1,
\qquad
P_2=P_{\min}D_2
\]
应用整除--包含反序对应，立得
\[
D_1\mid D_2
\iff
P_{\min}D_1\mid P_{\min}D_2
\iff
\mathcal N_n(D_2)\subseteq \mathcal N_n(D_1),
\]
且交公式同理由该命题给出。

再证商环识别。取 \(U\in \ZZ[\lambda]_{<\Delta}\)。由定理
\ref{thm:xi-time-part62af-hankel-gap-quotient-ring-model}，类
\[
\Theta_R(U)
=
\Bigl[\operatorname{coeff}_{<n}(U P_{\min})\Bigr]
\]
落入 \(\mathcal N_n(D)/\mathcal M_n(P)\) 当且仅当
\[
\operatorname{coeff}_{<n}(U P_{\min})\in \mathcal M_n(P_{\min}D).
\]
由于
\[
\deg(UP_{\min})<\Delta+d_{\min}=\deg P<n,
\]
此时截断无损；故上式等价于存在 \(Q\in\ZZ[\lambda]\) 使
\[
U P_{\min}=Q\,P_{\min}D.
\]
在整环 \(\ZZ[\lambda]\) 中约去首一多项式 \(P_{\min}\)，便得到
\[
U=Q D.
\]
也就是说，\(\Theta_R(U)\) 恰当且仅当 \(U\) 落在理想 \((D)\) 中。于是
\[
\Theta_R^{-1}\!\bigl(\mathcal N_n(D)/\mathcal M_n(P)\bigr)=(D)/(R).
\]

若写 \(R=D\widetilde R\)，则每个 \((D)/(R)\) 中的类都唯一写成
\[
D V\pmod R,
\qquad
V\in\ZZ[\lambda]_{<\deg\widetilde R},
\]
故其为秩 \(\deg\widetilde R=\Delta-\deg D\) 的自由阿贝尔群。后一个秩公式则由
\[
\rank_{\ZZ}\mathsf Q_n(P;s)=\Delta
\]
与前式相减即得。证毕。
\end{proof}

\begin{theorem}[高阶交叉异常的最小交互阶刻画]\label{thm:xi-time-part62af-cross-anom-interaction-degree}
\leanverified{paper\_xi\_time\_part62af\_cross\_anom\_interaction\_degree}
沿用定义 \ref{def:pom-higher-cross-anomaly} 与定理 \ref{thm:pom-mobius-complete-decomposition} 的记号，记
\[
f_\theta(\chi):=\log\mathfrak{M}_\chi(\theta),
\qquad
\chi\in\widehat G=\prod_{i=1}^{k}\widehat{G_i}.
\]
对 \(1\le r\le k\)，称 \(f_\theta\) 的交互阶不超过 \(r\)，若存在函数族
\[
L_T:\prod_{i\in T}\widehat{G_i}\to\RR
\qquad
\bigl(\varnothing\neq T\subseteq [k],\ |T|\le r\bigr)
\]
使得对一切 \(\chi\in\widehat G\) 都有
\[
f_\theta(\chi)
=
\sum_{\substack{\varnothing\neq T\subseteq [k]\\ |T|\le r}}
L_T(\chi_T).
\]
记满足该条件的最小 \(r\) 为 \(\deg_{\mathrm{int}}f_\theta\)，并约定当
\(
f_\theta\equiv 0
\)
时
\(
\deg_{\mathrm{int}}f_\theta:=0
\)
（等价地，\(\max\varnothing:=0\)）。
则对任意 \(1\le r\le k\)，成立严格等价
\[
\boxed{\;
\deg_{\mathrm{int}}f_\theta\le r
\quad\Longleftrightarrow\quad
\Delta_S(\cdot;\theta)\equiv 0
\ \text{对所有}\ |S|>r.\;}
\]
特别地，
\[
\boxed{\;
\deg_{\mathrm{int}}f_\theta
=
\max\bigl\{|S|:\ \Delta_S(\cdot;\theta)\not\equiv 0\bigr\}.\;}
\]
\end{theorem}
\begin{proof}
若 \(\deg_{\mathrm{int}}f_\theta\le r\)，则
\[
f_\theta=\sum_{|T|\le r}L_T(\chi_T).
\]
任取 \(|S|>r\)。对任一求和项 \(L_T(\chi_T)\)，总可取某个 \(i\in S\setminus T\)，于是定义
\ref{def:pom-higher-cross-anomaly} 证明中出现的差分算子 \(\delta_i\) 满足
\[
\delta_i L_T=0.
\]
因此
\[
\delta_S L_T=0
\qquad
(\forall\,|T|\le r),
\]
从而
\[
\Delta_S(\cdot;\theta)=\delta_S f_\theta\equiv 0.
\]

反之，若所有 \(|S|>r\) 的 \(\Delta_S(\cdot;\theta)\) 都恒为零，则由定理
\ref{thm:pom-mobius-complete-decomposition}，
\[
f_\theta(\chi)
=
\sum_{\varnothing\neq S\subseteq [k]}\Delta_S(\chi_S;\theta)
=
\sum_{\substack{\varnothing\neq S\subseteq [k]\\ |S|\le r}}
\Delta_S(\chi_S;\theta).
\]
取 \(L_S:=\Delta_S(\cdot;\theta)\) 即得到一个交互阶不超过 \(r\) 的表示。

最后一个公式只是上述等价的最小化表达：\(\deg_{\mathrm{int}}f_\theta\) 恰等于最高非零 Möbius 分量的阶数。证毕。
\end{proof}

\begin{theorem}[moment 编译下交互阶的刚性保持与纯三体失明]\label{thm:xi-time-part62af-moment-interaction-degree-rigidity}
设 \(G_1,\dots,G_k\) 为有限阿贝尔群，记
\[
\widehat G:=\prod_{i=1}^{k}\widehat{G_i},
\qquad
\chi^0:=(\chi_{1,0},\dots,\chi_{k,0}).
\]
取任意投影词 \(w\)，并以常数层角色坐标定义其异常函数
\[
a_w:\widehat G\to \RR,
\qquad
a_w(\chi):=[\mathrm{Anom}(w)]_\chi,
\qquad
a_w(\chi^0):=0.
\]
对每个 \(q\ge 2\)，记
\[
w^{\boxtimes q}:=\mathrm{MOM}_q\circ w.
\]
由命题 \ref{prop:pom-anom-gap-amplification} 与推论
\ref{cor:pom-anom-gap-amplification-completeness} 的坐标化读法，常数层异常满足
\[
a_{w^{\boxtimes q}}=q\,a_w.
\]
再按定义 \ref{def:pom-higher-cross-anomaly} 的同一 Möbius 差分记
\[
\Delta_S(w):=\Delta_S(a_w)
\qquad
(\varnothing\neq S\subseteq [k]),
\]
并定义
\[
\deg_{\mathrm{int}}(w):=\deg_{\mathrm{int}}(a_w).
\]
则对任意 \(q\ge 2\) 与任意非空 \(S\subseteq [k]\)，有
\[
\boxed{\;
\Delta_S(w^{\boxtimes q})=q\,\Delta_S(w).\;}
\]
特别地，
\[
\boxed{\;
\deg_{\mathrm{int}}(w^{\boxtimes q})
=
\deg_{\mathrm{int}}(w).\;}
\]
进一步，若 \(w\) 满足
\[
\Delta_{\{i,j\}}(w)\equiv 0
\qquad
(\forall\, i\neq j),
\]
且对某个三元组 \(i<j<k\) 有
\[
\Delta_{\{i,j,k\}}(w)\not\equiv 0,
\]
则对一切 \(q\ge 2\) 仍成立
\[
\Delta_{\{i,j\}}(w^{\boxtimes q})\equiv 0
\qquad
(\forall\, i\neq j),
\]
以及
\[
\Delta_{\{i,j,k\}}(w^{\boxtimes q})
=
q\,\Delta_{\{i,j,k\}}(w)
\not\equiv 0.
\]
因此，任何仅依赖两体交叉异常的全矩阶审计都无法恢复这类纯三体不可分性。
\end{theorem}
\begin{proof}
由
\(
a_{w^{\boxtimes q}}=q\,a_w
\)
及 Möbius 差分的线性性，立得
\[
\Delta_S(w^{\boxtimes q})
=
\Delta_S(q\,a_w)
=
q\,\Delta_S(w)
\qquad
(\varnothing\neq S\subseteq [k]).
\]
于是
\[
\Delta_S(w)\equiv 0
\iff
\Delta_S(w^{\boxtimes q})\equiv 0,
\]
故非零支撑集合在 moment 编译下保持不变。由定理
\ref{thm:xi-time-part62af-cross-anom-interaction-degree} 的最小交互阶公式，即得
\[
\deg_{\mathrm{int}}(w^{\boxtimes q})
=
\deg_{\mathrm{int}}(w).
\]

若 \(w\) 为纯三体型，则所有两体分量原本恒为零，而某个三体分量非零；把上式分别应用于 \(S=\{i,j\}\) 与 \(S=\{i,j,k\}\) 即得所示结论。命题
\ref{prop:pom-pairwise-cross-anom-insufficient} 进一步保证这种制度确实存在。证毕。
\end{proof}

\noindent
因此，过拟合递推所留下的全部原始 Hankel 缺口并非任意自由方向，而是由单一首一因子 \(R\) 的商环与除子格完全控制；与此同时，多轴异常的真实复杂度也并不由两体审计决定，而由最高非零 Möbius 分量的阶数决定，并在 moment 放大下保持严格不变。

\endinput

\paragraph{PRIMEGAME 的估值格点共轭、有限支撑分解与局部化 solenoid 宿主}
在自然扩张的外置预算、prime-register 账本语义与局部化 solenoid 终对象已经确立之后，固定 PRIMEGAME 微码的算术动力学还可进一步压缩为一条有限活跃支撑链：整数状态、first-fit 控制与局部化紧化在同一估值坐标中严格闭合。
沿用定义 \ref{def:pom-fractran}、命题 \ref{prop:pom-fractran-prime-translation} 与定理 \ref{thm:xi-time-part56-solenoid-terminal-prime-register-recovery} 的记号。固定一条十四指令 FRACTRAN 程序
\[
F=\left(\frac{a_1}{b_1},\dots,\frac{a_{14}}{b_{14}}\right),
\]
并称之为 PRIMEGAME。设其活跃素数集为
\[
S_*:=\{2,3,5,7,11,13,17,19,23,29\},
\]
且 PRIMEGAME 的全部分子与分母都只含 \(S_*\) 中的素因子。记
\[
\NN_{S_*}:=\left\{\prod_{p\in S_*}p^{x_p}:x_p\in \ZZ_{\ge 0}\right\},
\qquad
P_*:=\prod_{p\in S_*}p,
\qquad
G_{S_*}:=\ZZ[S_*^{-1}],
\]
并约定 \((\ZZ_{\ge 0})^{S_*}\) 与 \(\ZZ^{S_*}\) 上的偏序 \(x\ge y\) 总按逐坐标解释。

\begin{theorem}[PRIMEGAME 的精确 prime-register 共轭]\label{thm:xi-time-part62ae-primegame-exact-prime-register-conjugacy}
定义估值坐标
\[
\Psi_{S_*}:\NN_{S_*}\longrightarrow (\ZZ_{\ge 0})^{S_*},
\qquad
\Psi_{S_*}(n):=(v_p(n))_{p\in S_*}.
\]
对每个 \(1\le i\le 14\)，记
\[
\beta_i^{(*)}:=(v_p(b_i))_{p\in S_*}\in (\ZZ_{\ge 0})^{S_*},
\qquad
\delta_i^{(*)}:=(v_p(a_i)-v_p(b_i))_{p\in S_*}\in \ZZ^{S_*}.
\]
则 PRIMEGAME 的一步演化 \(T_F\) 在 \(\NN_{S_*}\) 上与估值格点上的分段平移
\[
\widetilde T_F(x):=x+\delta_{i(x)}^{(*)},
\qquad
i(x):=\min\{\,i:\ x\ge \beta_i^{(*)}\,\}
\]
严格共轭，即
\[
\Psi_{S_*}\circ T_F=\widetilde T_F\circ \Psi_{S_*}
\]
在 \(T_F\) 的定义域上恒成立。
\end{theorem}
\begin{proof}
命题 \ref{prop:pom-fractran-prime-translation} 已在全素数寄存器纤维上给出同一结论；此处只需把坐标限制到活跃支撑 \(S_*\)。对任意 \(n\in \NN_{S_*}\) 与任意指令 \(a_i/b_i\)，有
\[
\frac{a_i}{b_i}\ \text{可用于}\ n
\iff
b_i\mid n
\iff
v_p(n)\ge v_p(b_i)\quad(\forall\, p\in S_*),
\]
因为 PRIMEGAME 的全部分母都只含 \(S_*\)-素因子。故在 \(\Psi_{S_*}\)-坐标中，可行性恰由
\[
\Psi_{S_*}(n)\ge \beta_i^{(*)}
\]
刻画。若第 \(i\) 条指令被 first-fit 规则选中，则
\[
T_F(n)=n\frac{a_i}{b_i},
\]
从而对每个 \(p\in S_*\) 都有
\[
v_p(T_F(n))=v_p(n)+v_p(a_i)-v_p(b_i).
\]
于是
\[
\Psi_{S_*}(T_F(n))
=
\Psi_{S_*}(n)+\delta_i^{(*)}.
\]
再由 first-fit 选择定义，所取索引正是
\[
i(\Psi_{S_*}(n))
=
\min\{\,i:\ \Psi_{S_*}(n)\ge \beta_i^{(*)}\,\}.
\]
故共轭关系成立。证毕。
\end{proof}

\begin{theorem}[PRIMEGAME 的活跃支撑分解]\label{thm:xi-time-part62ae-primegame-active-support-splitting}
任取 \(n\in \ZZ_{>0}\)，则存在唯一分解
\[
n=u\,m,
\qquad
u\in \NN_{S_*},
\qquad
\gcd(m,P_*)=1.
\]
并且 PRIMEGAME 的一步演化满足：
\begin{enumerate}
  \item first-fit 选中的指令索引只由 \(u\) 决定，与 \(m\) 无关；
  \item 若 \(T_F(u)\) 有定义，则
  \[
  T_F(n)=m\,T_F(u);
  \]
  若 \(T_F(u)\) 无定义，则 \(T_F(n)\) 亦无定义。
\end{enumerate}
特别地，从初值 \(n_0=2\) 出发的 PRIMEGAME 轨道完全落在 \(\NN_{S_*}\) 内。
\end{theorem}
\begin{proof}
唯一分解来自整数的素因子分离：把 \(n\) 的 \(S_*\)-主部分并入 \(u\)，其余素因子并入 \(m\) 即得，而且此分解显然唯一。对任意指令 \(a_i/b_i\)，由于 \(\gcd(m,b_i)=1\) 且 \(b_i\) 的全部素因子都落在 \(S_*\)，有
\[
b_i\mid n
\iff
b_i\mid u.
\]
因此可行指令集合以及由此决定的 first-fit 索引只依赖于 \(u\)。若第 \(i\) 条指令适用于 \(u\)，则
\[
T_F(n)
=
n\frac{a_i}{b_i}
=
u\,m\frac{a_i}{b_i}
=
m\left(u\frac{a_i}{b_i}\right)
=
m\,T_F(u),
\]
其中最后一步利用了 PRIMEGAME 的全部分子也只含 \(S_*\)-素因子，故外部因子 \(m\) 完全惰性。若 \(u\) 上无可用指令，则 \(n\) 上也无可用指令。最后，由 \(2\in S_*\) 可知初值 \(n_0=2\) 对应 \(m=1\)；再由上式归纳即得整条轨道始终满足 \(m=1\)，从而完全落在 \(\NN_{S_*}\) 内。证毕。
\end{proof}

\begin{corollary}[PRIMEGAME 的有限阶段 solenoidal 宿主]\label{cor:xi-time-part62ae-primegame-finite-solenoidal-host}
PRIMEGAME 的一步控制律以及一切仅依赖 \(S_*\)-估值坐标的轨道观测，都先天地因子化到 \(G_{S_*}=\ZZ[S_*^{-1}]\) 的局部化阶段。其自然紧化宿主为有限阶段 arithmetic solenoid
\[
\Sigma_{S_*}:=\widehat{G_{S_*}},
\]
并满足 Pontryagin 对偶短正合列
\[
0\longrightarrow \prod_{p\in S_*}\ZZ_p
\longrightarrow \Sigma_{S_*}
\longrightarrow \TT
\longrightarrow 0.
\]
因此，对固定的 PRIMEGAME 微码而言，universal solenoid \(\widehat{\QQ}\) 并非操作层所必需；只有在把 PRIMEGAME 置入程序族、Mellin 读出或 \(\zeta\)-型完成时，才需要引入通用阶段。
\end{corollary}
\begin{proof}
定理 \ref{thm:xi-time-part62ae-primegame-active-support-splitting} 表明，PRIMEGAME 的一步选择与更新完全由 \(S_*\)-支撑部分 \(u\in \NN_{S_*}\) 决定，而与外部因子 \(m\) 解耦。因此任何只读取 \(S_*\)-估值坐标的轨道观测，都可经由 \(\NN_{S_*}\subset G_{S_*}\) 因子化。
另一方面，定理 \ref{thm:xi-time-part56-solenoid-terminal-prime-register-recovery} 与定理 \ref{thm:xi-localized-phase-sampling-closed-form} 已给出 \(G_{S_*}\) 的 Pontryagin 对偶紧化正是
\[
\Sigma_{S_*}=\widehat{G_{S_*}},
\]
并带有所示短正合列。故 PRIMEGAME 的自然紧化宿主已在有限局部化阶段闭合。证毕。
\end{proof}

\begin{corollary}[固定 FRACTRAN 微码的有限素数支撑刚性]\label{cor:xi-time-part62ae-fixed-fractran-finite-prime-support-rigidity}
任意固定长度的 FRACTRAN 程序
\[
F=\left(\frac{a_1}{b_1},\dots,\frac{a_k}{b_k}\right)
\]
都只激活有限个 prime-register 轴，即
\[
S(F):=\bigcup_{i=1}^k\{\,p\in \mathbb P:\ v_p(a_i b_i)>0\,\}
\]
为有限集，并且其一步 first-fit 控制与更新完全由该有限活跃支撑决定。因而单个固定程序表不能仅凭微码自身实现真正的无限素数控制；若要出现 universal-solenoidal 现象，必须把无穷性置于程序族、观测族或 Mellin/\(\zeta\) 完成过程中。
\end{corollary}
\begin{proof}
定理 \ref{thm:xi-time-part62ae-primegame-active-support-splitting} 的分解论证只用到了“程序表有限”与“全部分子分母的素因子支撑有限”这两点，因此对任意固定长度的 FRACTRAN 程序可逐字适用。把其中的 \(S_*\) 换成 \(S(F)\) 即得结论。证毕。
\end{proof}

因此，PRIMEGAME 在本文口径下可被精确识别为：哥德尔素数估值坐标给出状态空间，FRACTRAN 的 first-fit 规则给出微码执行，而有限局部化 solenoid 给出其紧化宿主。若继续推进谱侧分析，则自然对象将是 \(\Sigma_{S_*}\) 上的 transfer kernel、轨道 \(\zeta\) 与 Mellin 型读出。

\paragraph{单一 first-fit 微码的 classical 阻断、共尾素数耗尽与惰性谱扇区}
在 PRIMEGAME 的有限活跃支撑分解、局部化 solenoidal 宿主与有限素数账本刚性已经确定之后，若进一步要求单一 first-fit 微码直接抵达 classical \(\zeta\)/\(\Xi\) 的 Euler/Fredholm 层，则 prime-support、账本秩与有限阶段谱读出之间还会出现如下结构障碍。

\begin{theorem}[单一 first-fit 微码的经典 Euler 乘积阻断]\label{thm:xi-time-part62ae-single-firstfit-classical-euler-obstruction}
设
\[
F=\left(\frac{a_1}{b_1},\dots,\frac{a_k}{b_k}\right)
\]
为任意有限 FRACTRAN-compatible first-fit 微码，并记其活跃素数集为
\[
S(F):=\{\,p\in\mathbb P:\ v_p(a_i b_i)>0\ \text{for some }1\le i\le k\,\}.
\]
若某条证明路线坚持：其全部算术语义都由 \(F\) 的 prime-register 动力学，经结构保持的局部化 solenoidal 完成读出；并且该读出所诱导的 Euler/Fredholm 行列式在文稿既定的通用层语义下应与 classical \(\zeta\) 或完成函数 \(\Xi\) 一致，则该路线不可能由单一 \(F\) 实现。
\end{theorem}
\begin{proof}
记
\[
\NN_{S(F)}
:=
\left\{
\prod_{p\in S(F)}p^{x_p}:x_p\in\ZZ_{\ge 0}
\right\}.
\]
由推论 \ref{cor:xi-time-part62ae-fixed-fractran-finite-prime-support-rigidity}，\(F\) 的一步控制与更新完全由有限支撑 \(S(F)\) 决定。更具体地，若
\[
n=u\,m,
\qquad
u\in \NN_{S(F)},
\qquad
\gcd\!\Bigl(m,\prod_{p\in S(F)}p\Bigr)=1,
\]
则 first-fit 选中的指令索引只由 \(u\) 决定，且一旦 \(T_F(u)\) 有定义便有
\[
T_F(n)=m\,T_F(u).
\]
因此 \(F\) 的 prime-register 语义严格因子化到局部化阶段
\[
G_{S(F)}=\ZZ[S(F)^{-1}]
\]
之上。

再由定理 \ref{thm:killo-localization-circle-dimension-prime-ledger-splitting}，
\[
0\longrightarrow \prod_{p\in S(F)}\ZZ_p
\longrightarrow \widehat{G_{S(F)}}
\longrightarrow \TT
\longrightarrow 0,
\]
且 \(S(F)\) 可由核 \(\prod_{p\in S(F)}\ZZ_p\) 完全恢复。故任何仅从 \(\widehat{G_{S(F)}}\) 及其 Pontryagin 对偶账本 functorially 构造的对象，在结构层都不可能生成 \(S(F)\) 之外的新素数方向。

然而，文稿在 \S\ref{subsec:profinite-operator-l2} 中把“统一全部模数”的寄存器提升算子安置于 \(L^2(\widehat{\ZZ})\) 上，并在定理 \ref{thm:mellin-unitary-slice-unique}、定理 \ref{thm:app-horizon-toeplitz-lmi} 与定理 \ref{thm:xi-slice-horizon-carath-toeplitz} 中把 classical \(\Xi\)-切片的 Mellin--Plancherel 幺正读出及其 Toeplitz--Carath\'eodory 正性放在这一通用层上刻画。若单一有限微码已足以恢复 classical \(\zeta\) 或 \(\Xi\)，则该通用层中的全部素数尺度必须可由某个有限 \(S(F)\)-支撑对象结构性再现，矛盾。故结论成立。
\end{proof}

\begin{theorem}[共尾素数耗尽]\label{thm:xi-time-part62ae-cofinal-prime-exhaustion}
设 \(\mathscr S\) 为某一 FRACTRAN-compatible 路线中实际出现的全部有限 prime-support 之族。若该路线的目标是 classical \(\zeta\)/\(\Xi\) 而非固定有限核上的 RH 证书，则必存在一列有限素数集
\[
S_1\subset S_2\subset \cdots \subset \mathbb P,
\qquad
\bigcup_{n\ge 1}S_n=\mathbb P,
\]
并且每个 \(S_n\) 都可写成 \(\mathscr S\) 中有限多个支撑的并。特别地，通用 solenoidal 宿主并非附加修饰，而是由 prime-support 逻辑强制出的共尾极限对象。
\leanverified{paper\_xi\_time\_part62ae\_cofinal\_prime\_exhaustion}
\end{theorem}
\begin{proof}
若
\[
\bigcup_{S\in\mathscr S}S\neq \mathbb P,
\]
则可取某个素数 \(p_0\notin \bigcup_{S\in\mathscr S}S\)。于是对每个 \(S\in\mathscr S\)，定理 \ref{thm:killo-localization-circle-dimension-prime-ledger-splitting} 都给出
\[
\ker\!\bigl(\widehat{G_S}\to\TT\bigr)\cong \prod_{p\in S}\ZZ_p,
\]
其 \(p_0\)-主因子必为平凡。另一方面，定理 \ref{thm:xi-time-part62ac-localized-prime-adjoining-natural-extension} 表明：只有显式邻接素数 \(p_0\) 才会在自然扩张中新增一条 \(\ZZ_{p_0}\)-方向；推论 \ref{cor:xi-time-part62ac-prime-adjoining-order-independent-kernel-product} 又说明反复邻接后的最终对象只取决于实际并入的素数并。既然 \(p_0\) 从未出现，则任何由这些有限阶段及其共尾组合得到的对象都不可能在结构层看见 \(p_0\)。这与 classical Euler 乘积必须包含 \(p_0\)-尺度矛盾，故并集只能是 \(\mathbb P\)。

现取素数的固定枚举
\[
\mathbb P=\{p_1,p_2,\dots\}.
\]
对每个 \(n\ge 1\)，由上式并集等于 \(\mathbb P\)，可选 \(S^{(n)}\in\mathscr S\) 使 \(p_n\in S^{(n)}\)。定义
\[
S_n:=\bigcup_{j=1}^{n}S^{(j)}.
\]
则每个 \(S_n\) 皆有限，且
\[
S_1\subset S_2\subset\cdots,
\qquad
\bigcup_{n\ge 1}S_n=\mathbb P.
\]
再由推论 \ref{cor:xi-time-part62ac-prime-adjoining-order-independent-kernel-product}，每个有限阶段都对应规范短正合列
\[
0\longrightarrow \prod_{p\in S_n}\ZZ_p
\longrightarrow \Sigma_{S_n}
\longrightarrow \TT
\longrightarrow 0,
\qquad
\Sigma_{S_n}:=\widehat{G_{S_n}},
\]
并且各阶段之间的过渡正是逐素数邻接所诱导的共尾系统。故通用 solenoidal 宿主只能作为这些有限阶段的共尾极限出现，而不可能绕过 prime-support 账本被直接获得。证毕。
\end{proof}

\begin{theorem}[异常项与无穷秩账本的二择一]\label{thm:xi-time-part62ae-anomaly-infinite-rank-dichotomy}
设某条 Gödel--FRACTRAN 路线试图从 prime-register 微码抵达 classical \(\zeta\)/\(\Xi\)。若其 prime-support 最终覆盖全部素数，则至少发生下列两类机制之一：
\begin{enumerate}
  \item 承载乘法语义的有限阶段账本秩无界增长，并仅在共尾极限后进入通用层；
  \item 乘法到账本的传递在某一层失去严格同态性，而必须由异常签名、counterterm strictification 或普适可交换化群商承担补偿。
\end{enumerate}
特别地，不存在同时满足“prime-support 共尾覆盖全部素数、账本秩一致有界、乘法语义在各层皆严格同态”的证明链。
\end{theorem}
\begin{proof}
由定理 \ref{thm:xi-time-part62ae-cofinal-prime-exhaustion}，可取一列有限支撑
\[
S_1\subset S_2\subset\cdots\subset\mathbb P,
\qquad
\bigcup_{n\ge 1}S_n=\mathbb P.
\]
若在某一阶段中，\(S_n\) 内已有素数方向在账本中发生不可逆合并，从而不再可区分，则该路线上“乘法到账本”的传递已不再保持严格忠实语义，因而已经落入第二分支。故以下只需排除“各阶段都严格忠实同态化且账本秩一致有界”的可能性。

反设存在一致常数 \(R\)，使每一有限阶段的账本都可在秩至多 \(R\) 的加法/局部化对象上严格同态实现，并且在结构层忠实保留 \(S_n\) 中全部素数方向。由于 \(|S_n|\to\infty\)，可取 \(n\) 使
\[
|S_n|>R.
\]
记
\[
M_{S_n}:=\langle p:\ p\in S_n\rangle_{\times}\subset \NN^\times.
\]
则 \(M_{S_n}\cong (\NN^{|S_n|},+)\) 为由 \(|S_n|\) 个独立素数方向生成的自由交换幺半群。若其可在秩至多 \(R\) 的局部化账本中严格忠实实现，便与定理 \ref{thm:xi-time-part62debb-finite-prime-support-localized-axis-threshold} 及定理 \ref{thm:xi-localized-directsum-prime-truncation-minimal-channels} 矛盾。故只要严格同态性在各层保持，账本秩就不可能一致有界；这给出第一分支。

若反而坚持账本秩一致有界，则上面的矛盾表明：在某一有限阶段，乘法到账本的传递必不能继续保持严格同态。文稿内部对此类失配的可见补偿机制正是异常签名与 strictification 链：定义 \ref{def:pom-anom-signature} 给出异常对象，命题 \ref{prop:pom-counterterm-anom-cancel} 给出常数层对消律，命题 \ref{prop:pom-bc-quotient-strictify} 与命题 \ref{prop:pom-bc-quotient-universal} 给出普适可交换化群商，而定理 \ref{thm:dephys-offline-hair-curvature-reduction} 把这一失配压缩为“正规形—异常类—strictification”的统一归约。故任何绕开有限秩矛盾的路径都必须落入第二分支。特别地，三项条件不可能同时成立。证毕。
\end{proof}

\begin{corollary}[Ramanujan 窗口不可单独升级为 classical RH]\label{cor:xi-time-part62ae-ramanujan-window-no-classical-upgrade}
任意固定有限核家族上的 RH-sharp/Ramanujan 窗口结论，即便其边界由唯一临界参数与显式代数方程精确给出，也不能单独升级为 classical RH。特别地，命题 \ref{prop:real-input-40-collision-rhk-window} 给出的 \(4\times 4\) 碰撞位势窗口
\[
\Lambda(u)\le \lambda_1(u)^{1/2}
\iff
0<u\le u_R
\]
只构成一个有限核 RH 证书，而不构成 classical \(\Xi\) 的全局闭合。
\end{corollary}
\begin{proof}
命题 \ref{prop:real-input-40-collision-rhk-window} 已给出一个边界完全显式的有限维窗口，其参数 \(u_R\) 由唯一正实五次根确定。然 \S\ref{subsubsec:xi-horizon-ramanujan-feasible-reduction} 已明确指出：从 Ramanujan 机制出发逼近真实 RH，必须超出“固定有限维 \(+\) 单尺度 Dirichlet 化”的范式，进入增长维数的逼近序列或无穷维 Fredholm 行列式口径。

另一方面，定理 \ref{thm:xi-time-part62ae-single-firstfit-classical-euler-obstruction}、定理 \ref{thm:xi-time-part62ae-cofinal-prime-exhaustion} 与定理 \ref{thm:xi-time-part62ae-anomaly-infinite-rank-dichotomy} 表明：若目标是 classical \(\zeta\)/\(\Xi\)，则必须同时面对全素数共尾支撑、无界账本秩或异常补偿分支。单个固定有限核家族并不携带这些结构条件，故其 Ramanujan 窗口不能单独升级为 classical RH。证毕。
\end{proof}

\begin{proposition}[真实输入 \texorpdfstring{$40$}{40} 状态核的刚性惰性谱对]\label{prop:xi-time-part62ae-real-input40-rigid-inert-spectral-pair}
沿用定理 \ref{thm:real-input-40-det-uv-2plus8} 与推论 \ref{cor:real-input-40-zeta-uv-sqrtv-eigs} 的记号。则二维加权行列式
\[
\Delta(z;u,v)=\det(I-zM_{u,v})=(vz^2-1)P_8(z;u,v)
\]
对任意 \(u\) 恒带有一对参数刚性谱支
\[
\lambda_\pm=\pm\sqrt v.
\]
因此：
\begin{enumerate}
  \item 对每个整数 \(r\ge 1\)，都有
  \[
  \partial_u^{\,r}\log(vz^2-1)=0;
  \]
  \item 任何仅由 \(\partial_u^{\,r}\log\Delta(z;u,v)\) 所提取的谱统计，都只能读取 core 因子 \(P_8(z;u,v)\) 的变化，而对刚性谱对 \(\{\pm\sqrt v\}\) 完全失明；
  \item 因而任何仅以碰撞参数 \(u\) 的变形来读取“素数影子”的方案，都不可能在全谱层面保持忠实；若要与通用素数生成元层或 Mellin 幺正切片比对，必须先对该刚性扇区作商化或重整化。
\end{enumerate}
\end{proposition}
\begin{proof}
定理 \ref{thm:real-input-40-det-uv-2plus8} 与推论 \ref{cor:real-input-40-zeta-uv-sqrtv-eigs} 直接给出
\[
\Delta(z;u,v)=(vz^2-1)P_8(z;u,v),
\qquad
\lambda_\pm=\pm\sqrt v,
\]
其中 \(\lambda_\pm\) 与 \(u\) 无关。故对任意 \(r\ge 1\)，
\[
\partial_u^{\,r}\log(vz^2-1)=0,
\]
第一条成立。

进一步，
\[
\log\Delta(z;u,v)=\log(vz^2-1)+\log P_8(z;u,v),
\]
故一切 \(u\)-微分都自动杀掉刚性因子，只留下 \(P_8\) 的贡献。等价地，由定理 \ref{thm:xi-time-part73-twoparam-atomic-ghost-u-blindness} 与定理 \ref{thm:xi-time-part73-nonrigid-trace-newton-cayley-closure}，任何仅依赖 \(u\)-导数的 ghost/trace 审计都只能看到去除 \(\pm\sqrt v\) 后的非刚性迹谱 \(R_n(u,v)\)。

于是，若某个谱读出函子只使用 \(u\)-变形，则它会把“整条 determinant family”与“剔除刚性因子后的 core family”识别为同一对象；这在全谱层面不是忠实读出。故若要把该二维核接到通用素数生成元层或 Mellin 幺正切片的比较问题上，必须先将刚性惰性扇区 \((vz^2-1)\) 商去或重整化。证毕。
\end{proof}

\begin{conjecture}[Cofinal PRIMEGAME--\texorpdfstring{$\Xi$}{Xi} 极限定理]\label{conj:xi-time-part62ae-cofinal-primegame-xi-limit}
存在一族按有限素数集 \(S\subset\mathbb P\) 过滤的 FRACTRAN-compatible first-fit 微码 \(F_S\)，以及一族作用于 \(L^2(\widehat{\ZZ})\) 的迹类算子 \(\mathcal L_S\)，使得若记
\[
\widetilde\Delta_S(z):=\det(I-z\mathcal L_S),
\]
则满足：
\begin{enumerate}
  \item \(F_S\) 的 prime-register 语义因子化到
  \[
  G_S=\ZZ[S^{-1}];
  \]
  \item 若 \(S'\subset S\)，则 \(\widetilde\Delta_{S'}\) 与 \(\widetilde\Delta_S\) 在限制映射下相容；
  \item 当 \(S\nearrow \mathbb P\) 时，\(\mathcal L_S\) 在迹范数中收敛到某个迹类极限 \(\mathcal L_\infty\)，从而
  \[
  \widetilde\Delta_S(z)\longrightarrow
  \widetilde\Delta_\infty(z):=\det(I-z\mathcal L_\infty)
  \]
  于紧子集上一致成立；
  \item \(\widetilde\Delta_\infty\) 的 Mellin--Plancherel 幺正切片与圆盘完成化的 \(\Xi\)-截面一致。
\end{enumerate}
在此设定下，classical RH 等价于极限视界 Carath\'eodory 函数的 Toeplitz--PSD 层级成立。
\end{conjecture}

\paragraph{自然扩张层深、黄金共振 Schatten 排斥与单参数热力学塔的 RH 不可识别性}
在共尾 prime-support、有限阶段 solenoidal 宿主与 \(40\) 状态核的刚性谱扇区已经确定之后，还可把这三条链进一步压缩为一组彼此咬合的后果：自然扩张层深在局部化 prime-ledger 上具有严格线性下界；黄金共振振幅在算子侧排斥一切有限反项的 Schatten 正则化；而单参数热力学导数塔则不足以恢复圆盘完成化所需的 Toeplitz--Carath\'eodory 证书。

\begin{theorem}[自然扩张层深的 prime-slice 线性下界]\label{thm:xi-time-part62ae-natural-extension-prime-slice-linear-lower-bound}
沿用定理 \ref{thm:xi-layered-prime-register-sharp-natural-extension} 的设定。设
\[
\Psi_k(x_0):=(x_k,H_k(x_0)),
\qquad
H_k(x_0):=\prod_{j=0}^{k-1}h_j(x_0),
\]
且对每个 \(j<k\) 都存在两两不交的素数集 \(P_j\subset\mathbb P\)，使 \(h_j(x_0)\) 的全部素因子都落在 \(P_j\) 中。定义
\[
\operatorname{supp}_{\mathbb P}(h_j)
:=
\left\{
p\in\mathbb P:\ \exists\,x_0\in X_0,\ v_p\bigl(h_j(x_0)\bigr)>0
\right\}.
\]
若对每个 \(j<k\) 都存在 \(y_j\in X_{j+1}\) 使
\[
\#\pi_j^{-1}(y_j)\ge 2,
\]
且 \(\Psi_k\) 为单射，则有
\[
\#\Bigl(\bigcup_{j=0}^{k-1}\operatorname{supp}_{\mathbb P}(h_j)\Bigr)\ge k.
\]
更强地，对每个 \(j<k\) 都存在某个素数 \(p_j\in P_j\) 与某个微态 \(x_0^{(j)}\in X_0\) 使
\[
v_{p_j}\bigl(h_j(x_0^{(j)})\bigr)>0.
\]
\end{theorem}
\begin{proof}
由定理 \ref{thm:xi-layered-prime-register-sharp-natural-extension}，对每个 \(j<k\) 都存在某个微态 \(x_0^{(j)}\in X_0\) 使
\[
h_j\bigl(x_0^{(j)}\bigr)\neq 1.
\]
因此 \(h_j(x_0^{(j)})\) 至少含有一个素因子 \(p_j\in P_j\)，并满足
\[
v_{p_j}\bigl(h_j(x_0^{(j)})\bigr)>0.
\]
由于各 \(P_j\) 两两不交，故这些 \(p_j\) 两两不同，从而
\[
\#\Bigl(\bigcup_{j=0}^{k-1}\operatorname{supp}_{\mathbb P}(h_j)\Bigr)\ge k.
\]
证毕。
\end{proof}

\begin{corollary}[有限阶段 arithmetic solenoid 的分支历史容量]\label{cor:xi-time-part62ae-finite-solenoid-branch-history-capacity}
设 \(S\subset\mathbb P\) 为有限集，并记
\[
G_S=\ZZ[S^{-1}],
\qquad
0\longrightarrow \prod_{p\in S}\ZZ_p
\longrightarrow \widehat{G_S}
\longrightarrow \TT
\longrightarrow 0
\]
为定理 \ref{thm:xi-time-part56-solenoid-terminal-prime-register-recovery} 所给出的 arithmetic solenoid 短正合列。则任何 prime-register 方向受限于 \(S\) 的自然扩张忠实编码，至多只能承载 \(|S|\) 层非平凡分支历史。特别地，若某个有限纤维自映射 \(T:X\twoheadrightarrow X\) 的自然扩张在无穷多层上都存在真实分支，则不存在任何固定有限 \(S\) 使 \(\widehat{G_S}\) 忠实承载其全部前史；只能取一列
\[
S_1\subset S_2\subset\cdots,
\qquad
\bigcup_{n\ge 1}S_n=\mathbb P,
\]
并进入相应的共尾逆极限。
\end{corollary}
\begin{proof}
定理 \ref{thm:xi-time-part56-solenoid-terminal-prime-register-recovery} 与定理 \ref{thm:xi-localized-phase-sampling-closed-form} 表明：\(\widehat{G_S}\) 的全不连通 prime-ledger 核恰为
\[
\prod_{p\in S}\ZZ_p,
\]
因而其可用素数方向完全由 \(S\) 决定。若某个深度为 \(k\) 的自然扩张前史能在该有限阶段上被忠实编码，则它必然给出一个所有素因子都落在 \(S\) 中的单射分层 prime-register 外置。由定理 \ref{thm:xi-time-part62ae-natural-extension-prime-slice-linear-lower-bound}，此时至少需要 \(k\) 个互异素数方向，故 \(k\le |S|\)。

若 \(T\) 的自然扩张在无穷多层上均存在真实分支，则对任意 \(k\) 都可取长度 \(k\) 的前史截断，并由上段推出任何忠实编码都需要至少 \(k\) 个互异素数方向。于是任意固定有限 \(S\) 最终都会失效，只能取 \(S_n\nearrow\mathbb P\) 的共尾系统。证毕。
\end{proof}

\begin{theorem}[黄金共振振幅的 Hilbert--Schmidt 排斥]\label{thm:xi-time-part62ae-golden-resonance-hilbert-schmidt-obstruction}
设 \((e_u)_{u\ge 1}\) 为 Hilbert 空间 \(\mathcal H\) 中一组标准正交系，并令
\[
R
:=
\sum_{u\ge 1}C_\varphi(u)\,\langle\cdot,e_u\rangle e_u
\]
为其在自然极大定义域上的对角算子。则
\[
R\notin \mathcal S_2(\mathcal H).
\]
更强地，对任意 \(N\ge 0\)，尾算子
\[
R_{>N}
:=
\sum_{u>N}C_\varphi(u)\,\langle\cdot,e_u\rangle e_u
\]
亦不属于 \(\mathcal S_2(\mathcal H)\)。特别地，
\[
R\notin \mathcal S_1(\mathcal H).
\]
\end{theorem}
\begin{proof}
由定理 \ref{thm:fold-sigma-phi-diverges}，
\[
\Sigma_\varphi
:=
1+2\sum_{u=1}^{\infty}C_\varphi(u)^2
=+\infty,
\]
从而
\[
\sum_{u\ge 1}|C_\varphi(u)|^2=+\infty.
\]
这也正与推论 \ref{cor:fold-bernoulli-fourier-sampling-not-l2} 中 Fourier 采样序列的非平方可和性一致。对角算子的 Hilbert--Schmidt 范数满足
\[
\|R\|_{\mathcal S_2}^2
=
\sum_{u\ge 1}|C_\varphi(u)|^2,
\]
故 \(\|R\|_{\mathcal S_2}=+\infty\)，于是 \(R\notin\mathcal S_2(\mathcal H)\)。

同理，对任意固定 \(N\ge 0\)，有
\[
\|R_{>N}\|_{\mathcal S_2}^2
=
\sum_{u>N}|C_\varphi(u)|^2
=+\infty,
\]
故 \(R_{>N}\notin \mathcal S_2(\mathcal H)\)。最后，由 \(\mathcal S_1(\mathcal H)\subset \mathcal S_2(\mathcal H)\) 即得 \(R\notin \mathcal S_1(\mathcal H)\)。证毕。
\end{proof}

\begin{corollary}[有限反项 Fredholm 封口的不可能性]\label{cor:xi-time-part62ae-no-finite-counterterm-fredholm-closure}
在定理 \ref{thm:xi-time-part62ae-golden-resonance-hilbert-schmidt-obstruction} 的设定下，不存在任何有限秩算子 \(F\) 使
\[
R-F\in \mathcal S_1(\mathcal H).
\]
等价地，任何试图把原始黄金共振振幅直接作为 Mellin/Fourier 正交通道特征值的 Fredholm--Hilbert--P\'olya 封口，都不可能通过有限个 counterterm 获得迹类正则化。
\end{corollary}
\begin{proof}
若存在有限秩算子 \(F\) 使 \(R-F\in\mathcal S_1(\mathcal H)\)，则 \(R-F\in\mathcal S_2(\mathcal H)\)。另一方面，有限秩算子必属于 \(\mathcal S_2(\mathcal H)\)，故
\[
R=(R-F)+F\in\mathcal S_2(\mathcal H),
\]
与定理 \ref{thm:xi-time-part62ae-golden-resonance-hilbert-schmidt-obstruction} 矛盾。证毕。
\end{proof}

\begin{theorem}[单参数热力学塔对 RH 证书的不可识别性]\label{thm:xi-time-part62ae-oneparam-thermodynamic-tower-rh-indistinguishability}
固定 \(v>0\)。沿用定理 \ref{thm:real-input-40-det-uv-2plus8} 与推论 \ref{cor:real-input-40-zeta-uv-sqrtv-eigs} 的记号，定义
\[
\Delta_{\mathrm{full}}(z;u)
:=
\Delta(z;u,v)
=(vz^2-1)P_8(z;u,v),
\qquad
\Delta_{\mathrm{red}}(z;u):=P_8(z;u,v).
\]
对每个 \(r\ge 1\)，定义
\[
\mathcal C_r(z;u):=\partial_u^r\log\Delta(z;u,v).
\]
则对一切 \(r\ge 1\) 都有
\[
\partial_u^r\log \Delta_{\mathrm{full}}(z;u)
=
\partial_u^r\log \Delta_{\mathrm{red}}(z;u).
\]
因此，任何仅由 \(u\)-压力、\(u\)-矩、\(u\)-大偏差或其有限代数组合所生成、并最终只依赖这条 \(u\)-导数塔的热力学观察子，都无法区分 \(\Delta_{\mathrm{full}}\) 与 \(\Delta_{\mathrm{red}}\)。

然而，一旦把这两族对象送入文稿既定的圆盘完成化--Toeplitz--PSD 接口，则相应的 RH 正性证书并不相同。更一般地，若 \(B(z)\) 为任一与 \(u\) 无关的非常数完成因子，则 \(B(z)\Delta_{\mathrm{red}}(z;u)\) 与 \(\Delta_{\mathrm{red}}(z;u)\) 具有完全相同的全部 \(u\)-导数塔，但其 Toeplitz--Carath\'eodory 边界数据因
\[
\partial_z\log\!\bigl(B(z)\Delta_{\mathrm{red}}(z;u)\bigr)
=
\partial_z\log \Delta_{\mathrm{red}}(z;u)+\partial_z\log B(z)
\]
而被改写。特别地，取 \(B(z)=1-vz^2\) 时，\(\Delta_{\mathrm{full}}\) 与 \(\Delta_{\mathrm{red}}\) 在热力学塔上不可区分，却在完成化证书端保持可区分。
\end{theorem}
\begin{proof}
由
\[
\log\Delta_{\mathrm{full}}(z;u)
=
\log(vz^2-1)+\log P_8(z;u,v),
\]
并注意到 \(\log(vz^2-1)\) 与 \(u\) 无关，便对每个 \(r\ge 1\) 得到
\[
\partial_u^r\log\Delta_{\mathrm{full}}(z;u)
=
\partial_u^r\log P_8(z;u,v)
=
\partial_u^r\log\Delta_{\mathrm{red}}(z;u).
\]
因此任何仅由这条 \(u\)-导数塔构成的热力学观察子都无法区分二者。

另一方面，定理 \ref{thm:mellin-unitary-slice-unique} 与定理 \ref{thm:xi-slice-horizon-carath-toeplitz} 表明：classical RH 端所使用的是完成对象的 Mellin--Plancherel 幺正切片及其诱导的 Toeplitz--Carath\'eodory 边界数据，而非单一参数的热力学导数塔。若将完成对象乘上一个与 \(u\) 无关的非常数因子 \(B(z)\)，则全部 \(u\)-导数塔保持不变，但 \(\partial_z\log B(z)\) 会向边界系数额外注入一个非零序列，从而改变 Toeplitz 数据。取 \(B(z)=1-vz^2\) 即得所述具体例子。故从单参数热力学塔到 RH 证书的映射不是单射。证毕。
\end{proof}

\begin{corollary}[\texorpdfstring{$40$}{40} 状态核通向 classical RH 的最小额外输入]\label{cor:xi-time-part62ae-realinput40-minimal-extra-rh-input}
任何仅依赖 \(u\)-单参数热力学数据的 \(40\) 状态核--RH 方案，在逻辑上都不完备。若要抵达文稿中的 classical RH 终端证书，则至少必须补入下列三类信息之一：
\begin{enumerate}
  \item 一个与 \(u\) 独立的第二形变参数；
  \item 由分支历史或惯性谱扇区所决定的规范反因子；
  \item 直接给定的 Toeplitz--Carath\'eodory 边界数据。
\end{enumerate}
\end{corollary}
\begin{proof}
若单参数 \(u\)-热力学塔已足以恢复完成化 RH 证书，则定理 \ref{thm:xi-time-part62ae-oneparam-thermodynamic-tower-rh-indistinguishability} 中的 \(\Delta_{\mathrm{full}}\) 与 \(\Delta_{\mathrm{red}}\) 应在 RH 端不可区分；但该定理恰表明它们在全部 \(u\)-导数塔上不可区分，而在 Toeplitz--Carath\'eodory 证书端可区分，矛盾。故必须补入至少一类与单参数热力学塔独立的额外信息。证毕。
\end{proof}

\begin{conjecture}[分支惯性反因子与 reduced determinant 的共尾 \texorpdfstring{$\Xi$}{Xi} 极限]\label{conj:xi-time-part62ae-branch-inertia-counterfactor-cofinal-xi-limit}
猜想 \ref{conj:xi-time-part62ae-cofinal-primegame-xi-limit} 还应具有如下规范分解：存在一族按有限素数集 \(S\subset\mathbb P\) 过滤的 FRACTRAN-compatible / Fibonacci-compatible 内核 \(\Delta_S\)，以及与之伴随的一族规范反因子 \(B_S(z)\)，使得若记
\[
\Delta_S^{\mathrm{red}}(z):=\frac{\Delta_S(z)}{B_S(z)},
\]
则满足：
\begin{enumerate}
  \item \(B_S\) 恰吸收一切对所选热力学形变不可见、但对完成化证书可见的惯性谱支；
  \item \(\Delta_S^{\mathrm{red}}\) 对自然扩张分支历史保持忠实；
  \item 若 \(S'\subset S\)，则 \(B_{S'}\) 与 \(B_S\) 在限制映射下相容；
  \item 当 \(S\nearrow\mathbb P\) 时，\(\Delta_S^{\mathrm{red}}\) 在适当的 Fredholm--Mellin 口径下收敛到某个极限 \(\Delta_\infty\)，且 \(\Delta_\infty\) 的 Mellin--Plancherel 幺正切片与文稿圆盘完成化中的 \(\Xi\)-截面一致。
\end{enumerate}
在此设定下，classical RH 等价于 \(\Delta_\infty\) 所诱导的 Toeplitz--Carath\'eodory 数据全阶正性。
\end{conjecture}

综上，PRIMEGAME/FRACTRAN--solenoid--Mellin 路线在本文口径下呈现出一条明确的三重结构障碍：
\[
\boxed{\;
\text{分支历史要求共尾 prime-ledger，}\quad
\text{黄金共振排斥有限反项 Fredholm 封口，}\quad
\text{单参数热力学塔不能恢复 RH 的 Toeplitz 证书。}\;}
\]
因此，真正尚待建立的并非更多有限核实验，而是一套定义在共尾 prime-register 家族上的规范反因子理论，并要求其同时与自然扩张分支历史、有限阶段惯性谱扇区的剥离，以及 Mellin--Plancherel / Toeplitz--Carath\'eodory 终端接口严格相容。

\paragraph{有限分支前缀、Smith 曲率谱与几何对象的素数寄存器化}
在自然扩张的分支指标寄存器、有限深度 prime-slice 外置、Smith 信息亏损谱与整数轴对齐椭圆半群已经分别建立之后，还可把这四条线索压缩成一条更高阶的结构链：有限分支前缀给出 prime-slice 外置的规范前缀模型，Smith 曲率谱把全部模核损失压缩为离散边界函数，而几何对象的素数寄存器化则把这种刚性直接推送到结构层圆维。

\begin{theorem}[有限深度素数片外置是分支寄存器前缀的乘法封装]\label{thm:xi-finite-depth-primeslice-branch-prefix-compression}
设
\[
X_0\xrightarrow{\pi_0}X_1\xrightarrow{\pi_1}\cdots\xrightarrow{\pi_{k-1}}X_k
\]
为可数集合之间的有限纤维满射链。对每个 \(j=0,\dots,k-1\) 与每个 \(y\in X_{j+1}\)，固定单射
\[
\iota_{j,y}:\pi_j^{-1}(y)\hookrightarrow \NN.
\]
对 \(x_0\in X_0\) 递推定义
\[
x_{j+1}:=\pi_j(x_j),
\qquad
\beta_j(x_0):=\iota_{j,x_{j+1}}(x_j)\in\NN,
\]
并记有限深度分支寄存器
\[
\Psi_k(x_0):=(x_k,\beta_0(x_0),\dots,\beta_{k-1}(x_0))
\in X_k\times \NN^k.
\]
则 \(\Psi_k\) 为单射。

进一步，沿用定理 \ref{thm:killo-layered-prime-slices-finite-depth} 的两两不交 prime slices \((\mathcal P_j)_{0\le j\le k-1}\)，并为每层固定单射
\[
\alpha_j:X_{j+1}\times\NN\hookrightarrow \mathcal P_j.
\]
则存在一个定义在 \(\Psi_k(X_0)\) 上的规范映射
\[
\mathfrak M_k:\Psi_k(X_0)\longrightarrow X_k\times \ZZ_{>0}
\]
使得有限深度 prime-slice 外置
\[
\Phi_k(x_0):=
\left(
x_k,\ \prod_{j=0}^{k-1}\alpha_j(x_{j+1},\beta_j(x_0))
\right)
\]
满足
\[
\boxed{\;
\Phi_k=\mathfrak M_k\circ \Psi_k.\;}
\]
换言之，有限深度的 prime-slice 外置正是分支寄存器前缀的乘法封装；当该链来自某个有限纤维自映射的 \(k\)-步前史截断时，它恰是定理 \ref{thm:killo-natural-extension-branch-register} 与定理 \ref{thm:xi-prime-register-history-inverse-limit} 中自然扩张分支寄存器的长度 \(k\) 前缀压缩。
\end{theorem}
\begin{proof}
先证单射。若给定
\[
(x_k,a_0,\dots,a_{k-1})\in \Psi_k(X_0),
\]
则可自后向前唯一恢复：
\[
x_{k-1}=\iota_{k-1,x_k}^{-1}(a_{k-1}),
\qquad
x_{k-2}=\iota_{k-2,x_{k-1}}^{-1}(a_{k-2}),
\quad \dots,\quad
x_0=\iota_{0,x_1}^{-1}(a_0).
\]
因此 \(\Psi_k\) 在其像上可逆，从而为单射。

现定义 \(\mathfrak M_k\)。对任意 \((x_k,a_0,\dots,a_{k-1})\in \Psi_k(X_0)\)，先按上一步递推恢复唯一的中间态 \(x_{k-1},\dots,x_0\)，再置
\[
\mathfrak M_k(x_k,a_0,\dots,a_{k-1})
:=
\left(
x_k,\ \prod_{j=0}^{k-1}\alpha_j(x_{j+1},a_j)
\right).
\]
该定义良好，因为所有 \(x_j\) 已由 \((x_k,a_0,\dots,a_{k-1})\) 唯一确定。对真实输入 \(x_0\in X_0\)，代入 \(a_j=\beta_j(x_0)\) 即得
\[
(\mathfrak M_k\circ \Psi_k)(x_0)
=
\left(
x_k,\ \prod_{j=0}^{k-1}\alpha_j(x_{j+1},\beta_j(x_0))
\right)
=
\Phi_k(x_0).
\]
故 \(\Phi_k=\mathfrak M_k\circ \Psi_k\)。最后一条只是在 \(T\)-前史链情形下把 \(\Psi_k\) 识别为定理 \ref{thm:xi-prime-register-history-inverse-limit} 中截断分支码 \(\mathcal C_k\) 的一般链版本。证毕。
\end{proof}
\leanverified{paper\_xi\_finite\_depth\_primeslice\_branch\_prefix\_compression}

\begin{theorem}[分层素数片复杂度的精确最优值等于活跃分支深度]\label{thm:xi-layered-primeslice-complexity-active-depth}
在上述有限纤维链设定下，定义活跃分支层集合
\[
J:=
\left\{
\,0\le j\le k-1:\ \exists\,y\in X_{j+1}\ \text{使}\ \#\pi_j^{-1}(y)\ge 2
\right\}.
\]
在分层 prime-slice 外置语义下，定义其复杂度
\[
\mathfrak p(\pi_\bullet)
\]
为所有单射编码
\[
\Xi(x_0)=\bigl(x_k,H(x_0)\bigr),
\qquad
H(x_0)=\prod_{j=0}^{k-1}h_j(x_0),
\]
中非平凡 prime slices 的最小可能数目，其中每个 \(h_j(x_0)\) 的全部素因子均限制在某个预先指定的层片 \(\mathcal P_j\)。则有精确公式
\[
\boxed{\;
\mathfrak p(\pi_\bullet)=|J|.\;}
\]
特别地，若每一层都存在真正分支，即 \(J=\{0,\dots,k-1\}\)，则
\[
\boxed{\;
\mathfrak p(\pi_\bullet)=k.\;}
\]
\end{theorem}
\begin{proof}
这正是定理 \ref{thm:xi-prime-slice-nontrivial-layer-exact-minimality} 的内容。该定理证明：任一单射分层 prime-slice 编码的非平凡片数都不小于 \(|J|\)，而且存在构造恰好在每个活跃层使用一枚有效 prime slice、在非活跃层取平凡因子 \(1\) 的单射实现。故最小值正是 \(|J|\)。证毕。
\end{proof}

\begin{theorem}[Smith 信息亏损谱的离散曲率完全刻画]\label{thm:xi-smith-loss-spectrum-discrete-curvature-complete-characterization}
设 \(f:\ZZ_{\ge 0}\to \ZZ_{\ge 0}\) 满足 \(f(0)=0\)，并定义一阶差分
\[
\Delta(k):=f(k)-f(k-1)\qquad(k\ge 1).
\]
则下列两条等价：
\begin{enumerate}
\item 存在某个有限多重集 \(E\subset \ZZ_{\ge 1}\) 使
\[
f(k)=\sum_{e\in E}\min(k,e)
\qquad(k\ge 0);
\]
\item 序列 \((\Delta(k))_{k\ge 1}\) 非负、单调不增且最终为零。
\end{enumerate}
在这些条件成立时，\(E\) 由离散曲率唯一恢复：对每个 \(t\ge 1\)，其重数满足
\[
\boxed{\;
\#\{e\in E:\ e=t\}
=
\Delta(t)-\Delta(t+1)
=
2f(t)-f(t-1)-f(t+1).\;}
\]
因此，任何 \(\,p\)-进 Smith 信息亏损谱都可被视为一个有限 Young 边界函数，而确切的初等因子多重集正是该边界的离散曲率。
\leanverified{paper\_xi\_smith\_loss\_spectrum\_discrete\_curvature\_complete\_characterization}
\end{theorem}
\begin{proof}
定理 \ref{thm:xi-smith-padic-loss-spectrum-classification} 已证明：函数 \(f\) 可以写成
\[
f(k)=\sum_{e\ge 1}m_e\min(k,e)
\]
当且仅当其前向差分 \(\Delta(k)\) 非负、非增且有限支撑；并且此时重数 \(m_e\) 唯一确定。推论 \ref{cor:xi-smith-second-discrete-curvature-recovery} 进一步给出
\[
m_t=\Delta(t)-\Delta(t+1)=2f(t)-f(t-1)-f(t+1).
\]
将 \(m_t\) 解释为 \(E\) 中取值 \(t\) 的重数，即得全部结论。证毕。
\end{proof}

\begin{corollary}[Smith 核谱函数的离散凹性、平台截断与二阶差分反演]\label{cor:xi-smith-kernel-spectrum-discrete-concavity-plateau-inversion}
设 \(A\in M_{m\times n}(\ZZ)\) 满足 \(\operatorname{rank}_{\QQ}(A)=m\)，其 Smith 不变量为
\[
d_1\mid d_2\mid\cdots\mid d_m.
\]
固定素数 \(p\)，并记
\[
K_p(k):=\#\ker(A_{p^k})
\qquad(k\ge 1),
\]
\[
f_p(k):=\log_p K_p(k)-k(n-m)
\qquad(k\ge 1),
\qquad
f_p(0):=0.
\]
再记
\[
e_i(p):=v_p(d_i),
\qquad
E_p:=\max_{1\le i\le m}e_i(p),
\qquad
m_{p,\ell}:=\#\{i:\ e_i(p)=\ell\}
\quad(\ell\ge 1).
\]
则成立：
\begin{enumerate}
  \item \(f_p\) 是单调不减的离散凹函数，即
  \[
  f_p(k+1)-f_p(k)\ge f_p(k+2)-f_p(k+1)
  \qquad(k\ge 0);
  \]
  \item 当 \(k\ge E_p\) 时，\(f_p\) 进入平台：
  \[
  f_p(k)=\sum_{i=1}^{m}e_i(p);
  \]
  \item 对每个 \(\ell\ge 1\)，精确重数由二阶差分恢复：
  \[
  m_{p,\ell}
  =
  2f_p(\ell)-f_p(\ell-1)-f_p(\ell+1).
  \]
\end{enumerate}
\end{corollary}
\begin{proof}
由定理 \ref{thm:killo-smith-loss-spectrum}，
\[
f_p(k)=\sum_{i=1}^{m}\min\bigl(k,e_i(p)\bigr)
\qquad(k\ge 1),
\]
并且
\[
\Delta_p(k):=f_p(k)-f_p(k-1)=\#\{i:\ e_i(p)\ge k\}
\qquad(k\ge 1).
\]
因此 \(\Delta_p(k)\) 关于 \(k\) 非负且单调不增，故 \(f_p\) 单调不减且离散凹，这证明了第 \(1\) 条。

若 \(k\ge E_p\)，则 \(\min(k,e_i(p))=e_i(p)\) 对所有 \(i\) 都成立，故
\[
f_p(k)=\sum_{i=1}^{m}e_i(p),
\]
得到第 \(2\) 条。

最后，由上式直接有
\[
m_{p,\ell}
=
\#\{i:\ e_i(p)\ge \ell\}-\#\{i:\ e_i(p)\ge \ell+1\}
=
\Delta_p(\ell)-\Delta_p(\ell+1),
\]
从而
\[
m_{p,\ell}
=
\bigl(f_p(\ell)-f_p(\ell-1)\bigr)-\bigl(f_p(\ell+1)-f_p(\ell)\bigr)
=
2f_p(\ell)-f_p(\ell-1)-f_p(\ell+1).
\]
证毕。
\end{proof}

\begin{theorem}[模数信息亏损是 divisibility 格上的精确 valuation，且具有有理 Euler 因子]\label{thm:xi-kernel-loss-divisibility-valuation-rational-euler}
\leanverified{paper\_xi\_kernel\_loss\_divisibility\_valuation\_rational\_euler}
设 \(A\in M_{m\times n}(\ZZ)\) 满足 \(\operatorname{rank}_{\QQ}(A)=m\)，其 Smith 不变量为
\[
d_1\mid d_2\mid\cdots\mid d_m.
\]
对每个 \(b\ge 1\)，记
\[
A_b:(\ZZ/b\ZZ)^n\to(\ZZ/b\ZZ)^m
\]
为 \(A\) 的模 \(b\) 诱导映射。对任意正整数 \(b\)，定义归一化信息亏损
\[
L_A(b):=
\log \#\ker(A_b)-(n-m)\log b
=
\sum_{i=1}^{m}\log \gcd(d_i,b).
\]
则成立：
\begin{enumerate}
\item 对任意 \(b,c\ge 1\)，有精确 valuation 恒等式
\[
\boxed{\;
L_A(\operatorname{lcm}(b,c))+L_A(\gcd(b,c))
=
L_A(b)+L_A(c).\;}
\]
\item 定义 Dirichlet 生成函数
\[
\Xi_A(s):=\sum_{b\ge 1}\frac{\#\ker(A_b)}{b^{\,n-m+s}}
\qquad(\Re s>1).
\]
若记 \(e_i(p):=v_p(d_i)\)、\(E_p:=\max_i e_i(p)\)、\(\sigma_p:=\sum_i e_i(p)\)，则
\[
\boxed{\;
\Xi_A(s)=\prod_{p}\Xi_{A,p}(s),\qquad
\Xi_{A,p}(s)=\sum_{k=0}^{\infty}p^{-ks+\sum_{i=1}^{m}\min(k,e_i(p))}.\;}
\]
并且每个局部因子都是 \(p^{-s}\) 的有理函数，更精确地，
\[
\boxed{\;
\Xi_{A,p}(s)
=
\sum_{k=0}^{E_p-1}p^{-ks+\sum_i\min(k,e_i(p))}
\;+\;
\frac{p^{-E_p s+\sigma_p}}{1-p^{-s}}.\;}
\]
\end{enumerate}
\end{theorem}
\begin{proof}
由推论 \ref{cor:killo-kernel-size-general-modulus}，
\[
L_A(b)=\sum_{p}\left(\sum_{i=1}^{m}\min(v_p(b),e_i(p))\right)\log p.
\]
故要证第一式，只需逐素数验证
\[
\sum_i \min(\max\{u,v\},e_i)+\sum_i \min(\min\{u,v\},e_i)
=
\sum_i \min(u,e_i)+\sum_i \min(v,e_i),
\]
其中 \(u=v_p(b)\)、\(v=v_p(c)\)。而对每个固定 \(e\) 都有
\[
\min(\max\{u,v\},e)+\min(\min\{u,v\},e)=\min(u,e)+\min(v,e),
\]
逐项求和即得 valuation 恒等式。

第二式同样由推论 \ref{cor:killo-kernel-size-general-modulus} 与定理 \ref{thm:xi-kernel-dirichlet-zeta-finite-euler-correction} 直接推出。对 \(b=p^k\)，有
\[
\frac{\#\ker(A_{p^k})}{p^{k(n-m)}}
=
p^{\sum_i\min(k,e_i(p))},
\]
故局部因子为
\[
\Xi_{A,p}(s)=\sum_{k\ge 0}p^{-ks+\sum_i\min(k,e_i(p))}.
\]
当 \(k\ge E_p\) 时，\(\min(k,e_i(p))=e_i(p)\) 对所有 \(i\) 恒定，尾部即成为几何级数
\[
\sum_{k\ge E_p}p^{-ks+\sigma_p}
=
\frac{p^{-E_p s+\sigma_p}}{1-p^{-s}},
\]
于是得到所示有理表达。证毕。
\end{proof}

\begin{corollary}[模核大小在整除格上满足精确乘法恒等式]\label{cor:xi-kernel-size-gcd-lcm-exact-multiplicativity}
\leanverified{paper\_xi\_kernel\_size\_gcd\_lcm\_exact\_multiplicativity}
沿用定理 \ref{thm:xi-kernel-loss-divisibility-valuation-rational-euler} 的记号。对任意 \(b,c\ge 1\)，都有
\[
\#\ker(A_{\gcd(b,c)})\cdot \#\ker(A_{\operatorname{lcm}(b,c)})
=
\#\ker(A_b)\cdot \#\ker(A_c).
\]
等价地，模核大小函数在除去自由因子 \(b^{\,n-m}\) 之后，正是正整数整除格上的一个严格 valuation。
\end{corollary}
\begin{proof}
由定理 \ref{thm:xi-kernel-loss-divisibility-valuation-rational-euler}，
\[
L_A(\operatorname{lcm}(b,c))+L_A(\gcd(b,c))=L_A(b)+L_A(c).
\]
再由定义
\[
L_A(t)=\log\#\ker(A_t)-(n-m)\log t
\qquad(t\ge 1),
\]
以及恒等式
\[
\log\gcd(b,c)+\log\operatorname{lcm}(b,c)=\log b+\log c,
\]
可得
\[
\log\#\ker(A_{\gcd(b,c)})+\log\#\ker(A_{\operatorname{lcm}(b,c)})
=
\log\#\ker(A_b)+\log\#\ker(A_c).
\]
两边取指数即得所示精确乘法关系。证毕。
\end{proof}

\begin{theorem}[模核 Dirichlet 级数的 Euler--rational 形态与 Smith 完备性]\label{thm:xi-kernel-dirichlet-euler-rational-smith-completeness}
设 \(A\in M_{m\times n}(\ZZ)\) 满足 \(\operatorname{rank}_{\QQ}(A)=m\)，其 Smith 不变量为
\[
d_1\mid d_2\mid\cdots\mid d_m,\qquad d_i>0.
\]
对每个正整数 \(b\)，记
\[
A_b:(\ZZ/b\ZZ)^n\to(\ZZ/b\ZZ)^m
\]
为模 \(b\) 的诱导映射，并定义 Dirichlet 级数
\[
\mathcal K_A(s):=\sum_{b\ge 1}\frac{\#\ker(A_b)}{b^s}.
\]
则成立：
\[
\#\ker(A_b)=b^{\,n-m}\prod_{i=1}^{m}\gcd(d_i,b),
\]
从而 \(b\mapsto \#\ker(A_b)\) 为积性函数，故
\[
\mathcal K_A(s)=\prod_p \mathcal K_{A,p}(s),
\]
其中局部因子为
\[
\mathcal K_{A,p}(s)
=
\sum_{k\ge 0}p^{-k(s-(n-m))+f_p(k)},
\qquad
f_p(k):=\sum_{i=1}^{m}\min\bigl(e_i(p),k\bigr),
\]
并记
\[
e_i(p):=v_p(d_i),\qquad
E_p:=\max_i e_i(p),\qquad
\sigma_p:=\sum_{i=1}^{m}e_i(p).
\]
则更精确地，
\[
\mathcal K_{A,p}(s)
=
\sum_{k=0}^{E_p-1}p^{-k(s-(n-m))+f_p(k)}
\;+\;
\frac{p^{-E_p(s-(n-m))+\sigma_p}}{1-p^{-(s-(n-m))}},
\]
因此每个局部因子都是 \(p^{-s}\) 的有理函数。

若再记
\[
\Delta_A:=\prod_{i=1}^{m}d_i,
\]
则
\[
\mathcal K_A(s)
=
\zeta\!\bigl(s-(n-m)\bigr)
\prod_{p\mid \Delta_A}
\Bigl(1-p^{-(s-(n-m))}\Bigr)\mathcal K_{A,p}(s).
\]
因而 \(\mathcal K_A(s)\) 的收敛横坐标恰为
\[
\sigma_c=n-m+1,
\]
并在该点具有唯一简单极点。更进一步，在固定行秩 \(m\) 的口径下，\(\mathcal K_A(s)\) 唯一决定全部 Smith 不变量 \((d_1,\dots,d_m)\)。
\end{theorem}
\begin{proof}
由推论 \ref{cor:killo-kernel-size-general-modulus}，
\[
\#\ker(A_b)=b^{\,n-m}\prod_{i=1}^{m}\gcd(d_i,b).
\]
右端关于 \(b\) 为积性函数，故 \(\mathcal K_A\) 具有 Euler 乘积分解。对 \(b=p^k\)，再由同一推论与定理 \ref{thm:killo-smith-loss-spectrum}，
\[
\#\ker(A_{p^k})
=
p^{k(n-m)}\prod_{i=1}^{m}p^{\min(e_i(p),k)}
=
p^{k(n-m)+f_p(k)}.
\]
于是
\[
\mathcal K_{A,p}(s)
=
\sum_{k\ge 0}\frac{\#\ker(A_{p^k})}{p^{ks}}
=
\sum_{k\ge 0}p^{-k(s-(n-m))+f_p(k)}.
\]
当 \(k\ge E_p\) 时，\(\min(e_i(p),k)=e_i(p)\) 对所有 \(i\) 都成立，故 \(f_p(k)=\sigma_p\)，尾部因而化为几何级数
\[
\sum_{k\ge E_p}p^{-k(s-(n-m))+\sigma_p}
=
\frac{p^{-E_p(s-(n-m))+\sigma_p}}{1-p^{-(s-(n-m))}}.
\]
这就得到局部因子的有理表达式。

若 \(p\nmid \Delta_A\)，则全部 \(e_i(p)=0\)，从而
\[
\mathcal K_{A,p}(s)=\sum_{k\ge 0}p^{-k(s-(n-m))}
=
\frac{1}{1-p^{-(s-(n-m))}}.
\]
把这些好素数因子提出，即得
\[
\mathcal K_A(s)
=
\zeta\!\bigl(s-(n-m)\bigr)
\prod_{p\mid \Delta_A}
\Bigl(1-p^{-(s-(n-m))}\Bigr)\mathcal K_{A,p}(s).
\]
右侧修正项是有限 Euler 乘积，因此在 \(s=n-m+1\) 邻域内解析且非零；于是 \(\mathcal K_A\) 的收敛横坐标与 \(\zeta(s-(n-m))\) 相同，恰为 \(n-m+1\)，并在该点具有唯一简单极点。

最后证明 Smith 完备性。Dirichlet 级数 \(\mathcal K_A(s)\) 唯一决定全部系数
\[
a_b:=\#\ker(A_b).
\]
极点位置给出 \(n-m=\sigma_c-1\)。因此对每个素数 \(p\) 与每个 \(k\ge 1\)，可由
\[
a_{p^k}=p^{k(n-m)+f_p(k)}
\]
恢复
\[
f_p(k)=v_p(a_{p^k})-k(n-m).
\]
再由推论 \ref{cor:xi-smith-second-discrete-curvature-recovery}，
\[
\#\{i:\ e_i(p)=k\}
=
2f_p(k)-f_p(k-1)-f_p(k+1)
\qquad(k\ge 1),
\]
故每个素数 \(p\) 上的赋值多重集 \(\{e_i(p)\}_{i=1}^{m}\) 都被唯一确定。由于行秩 \(m\) 已固定，零赋值的个数亦随之确定。记
\[
0\le e_1^{\uparrow}(p)\le \cdots \le e_m^{\uparrow}(p)
\]
为该多重集的非降序排列，则逐项合并便唯一恢复整个 Smith 链
\[
d_i=\prod_p p^{e_i^{\uparrow}(p)},
\qquad(1\le i\le m).
\]
证毕。
\end{proof}

\begin{theorem}[整数轴对齐椭圆半群的结构层同态半圆维发散，而实现层仅为 \texorpdfstring{$1/2$}{1/2}]\label{thm:xi-integer-ellipse-half-circle-dimension-gap}
沿用定理 \ref{thm:killo-ellipse-diagonal-prime-register-equivalence} 的记号，令
\[
\mathsf E_{\NN}
:=
\left\{
E_{a,b}:\ \frac{x^2}{a^2}+\frac{y^2}{b^2}=1,\ a,b\in\NN_{>0}
\right\},
\qquad
E_{a,b}\odot E_{c,d}:=E_{ac,bd}.
\]
则
\[
\boxed{\;
\cdim^{\mathrm{hom}}_{+}(\mathsf E_{\NN},\odot)=+\infty,
\qquad
\cdim^{\mathrm{impl}}(\mathsf E_{\NN})=\frac12.\;}
\]
等价地，不存在任何有限加法寄存器系统 \((\NN^k,+)\) 能对整个整数轴对齐椭圆半群给出严格同态的完整编码；但作为纯集合，\(\mathsf E_{\NN}\) 仍只具有可数无限集的实现层复杂度。
\end{theorem}
\begin{proof}
对任意有限素数集 \(S\subset\mathbb P\)，定义子幺半群
\[
\mathsf E_S:=\{E_{a,1}:a\in \NN_S^\times\}\subset \mathsf E_{\NN}.
\]
映射
\[
\NN_S^\times\longrightarrow \mathsf E_S,\qquad a\longmapsto E_{a,1}
\]
是幺半群同构，故 \(\mathsf E_{\NN}\) 对任意 \(|S|=N\) 都包含一份 \(\NN_S^\times\) 的同态拷贝。于是由推论 \ref{cor:xi-unbounded-prime-directions-force-infinite-half-circle-dimension}，
\[
\cdim^{\mathrm{hom}}_{+}(\mathsf E_{\NN},\odot)=+\infty.
\]

另一方面，\(\mathsf E_{\NN}\) 与 \(\NN_{>0}\times \NN_{>0}\) 等势，因而是可数无限集。故存在集合单射
\[
\mathsf E_{\NN}\hookrightarrow \NN,
\]
从而
\[
\cdim^{\mathrm{impl}}(\mathsf E_{\NN})\le \frac12.
\]
又实现层半圆维为 \(0\) 当且仅当集合有限（见命题 \ref{prop:killo-implementation-vs-structural-half-circle-dimension} 的定义口径），故
\[
\cdim^{\mathrm{impl}}(\mathsf E_{\NN})=\frac12.
\]
证毕。
\end{proof}

\begin{corollary}[整数轴对齐椭圆半群不存在任何有限加法账本的忠实同态实现]\label{cor:xi-integer-ellipse-no-faithful-finite-additive-ledger}
沿用定理 \ref{thm:xi-integer-ellipse-half-circle-dimension-gap} 的记号，并记
\[
\mathsf E_x:=\{E_{a,1}:a\in\NN_{>0}\}\subset \mathsf E_{\NN}.
\]
则不存在任何有限整数 \(k\) 与单射幺半群同态
\[
\Phi:(\mathsf E_{\NN},\odot)\hookrightarrow (\NN^k,+).
\]
另一方面，\(\mathsf E_{\NN}\) 作为可数集合仍可单射进入 \(\NN\)。因此，对整数轴对齐椭圆半群而言，结构层的忠实同态实现与实现层的可寻址集合压缩发生严格分离。
\end{corollary}
\begin{proof}
映射
\[
\NN_{>0}\longrightarrow \mathsf E_x,\qquad
a\longmapsto E_{a,1}
\]
满足
\[
E_{a,1}\odot E_{c,1}=E_{ac,1},
\]
故给出幺半群同构
\[
(\NN_{>0},\times)\cong (\mathsf E_x,\odot).
\]
若存在某个有限 \(k\) 与单射幺半群同态
\[
\Phi:(\mathsf E_{\NN},\odot)\hookrightarrow (\NN^k,+),
\]
则其限制
\[
\Phi|_{\mathsf E_x}:(\NN_{>0},\times)\hookrightarrow (\NN^k,+)
\]
仍为单射幺半群同态。这与定理 \ref{thm:killo-prime-freedom-non-finitizability} 矛盾，因为 \((\NN^k,+)\) 是有限生成、交换且可消去的幺半群。故这样的 \(\Phi\) 不存在。

另一方面，\(\mathsf E_{\NN}\) 与 \(\NN_{>0}\times \NN_{>0}\) 等势，故是可数无限集，从而存在集合单射
\[
\mathsf E_{\NN}\hookrightarrow \NN.
\]
这就给出了所述的严格分离。证毕。
\end{proof}

\noindent
上述这些结果把附录 \(H\) 的若干局部构造压缩为一条统一刚性链：有限深度 branch register 是 prime-slice 外置的前缀模型，prime-slice 的最小复杂度由活跃分支层精确锁定；进入模数投影后，全部信息亏损不仅被 Smith 曲率谱完全分类，而且在整除格上形成精确 valuation 与原始模核乘法恒等式；最终，这条素数寄存器刚性可以直接投射到整数椭圆半群这样的显式几何对象，迫使其在结构层上表现出无限的加法寄存器秩，而实现层仍停留在可数可寻址的最低非零复杂度。

\endinput

\paragraph{自然扩张寄存器的逆极限规范性、Smith 局部亏损谱的自由化与 Euler 完备不变量}
在有限深度 prime-slice 前缀压缩、Smith 离散曲率恢复与整数椭圆原子分解之外，附录 \(H\) 中的 branch-register 语义还可继续压缩为逆极限、有限步过去的极小完备前缀、局部亏损谱的自由原子结构以及全模数核计数的 Smith 饱和阈值层的若干闭式结论。

\begin{theorem}[自然扩张分支寄存器系统是有限深度截断的严格逆极限]\label{thm:xi-branch-register-strict-inverse-limit}
\leanverified{paper\_xi\_branch\_register\_strict\_inverse\_limit}
沿用定理 \ref{thm:killo-natural-extension-branch-register} 的记号。对每个 \(k\ge 1\)，记
\[
X^{\leftarrow,k}
:=
\{(x_0,x_{-1},\dots,x_{-k})\in X^{k+1}:T(x_{-j})=x_{-j+1}\ \ (1\le j\le k)\},
\]
\[
\mathcal C_k(x_0,x_{-1},\dots,x_{-k})
:=
\bigl(x_0,(\beta(x_{-1}),\dots,\beta(x_{-k}))\bigr),
\]
\[
\widehat X_k^{\mathrm{br}}
:=
\mathcal C_k(X^{\leftarrow,k})
\subset X\times \NN^k.
\]
定义截断投影
\[
\pi_{k+1,k}^{\mathrm{br}}:\widehat X_{k+1}^{\mathrm{br}}\to \widehat X_k^{\mathrm{br}},
\qquad
\pi_{k+1,k}^{\mathrm{br}}(x,a_1,\dots,a_{k+1})=(x,a_1,\dots,a_k),
\]
以及有限深度更新
\[
\widehat T_k^{\mathrm{br}}(x,a_1,\dots,a_k)
:=
\begin{cases}
(T(x),\beta(x)),& k=1,\\
(T(x),\beta(x),a_1,\dots,a_{k-1}),& k\ge 2.
\end{cases}
\]
则 \((\widehat X_k^{\mathrm{br}},\pi_{k+1,k}^{\mathrm{br}})_{k\ge 1}\) 构成投影系统，并且存在规范同构
\[
\widehat X
\cong
\varprojlim_{k}\widehat X_k^{\mathrm{br}}
\]
把定理 \ref{thm:killo-natural-extension-branch-register} 中的 \(\widehat T\) 识别为
\[
\widehat T
=
\varprojlim_{k}\widehat T_k^{\mathrm{br}}.
\]
\end{theorem}
\begin{proof}
对每个 \(k\ge 1\)，定义截断映射
\[
\rho_k:\widehat X\to \widehat X_k^{\mathrm{br}},
\qquad
\rho_k(x,(a_1,a_2,\dots))=(x,a_1,\dots,a_k).
\]
由定理 \ref{thm:killo-natural-extension-branch-register}，\(\widehat X=\mathcal C(X^{\leftarrow})\)。若
\[
\mathcal C(x_0,x_{-1},x_{-2},\dots)=(x_0,(a_1,a_2,\dots)),
\]
则
\[
\rho_k(x_0,(a_1,a_2,\dots))
=
\mathcal C_k(x_0,x_{-1},\dots,x_{-k})\in \widehat X_k^{\mathrm{br}},
\]
故 \(\rho_k\) 良定义，且显然满足
\[
\pi_{k+1,k}^{\mathrm{br}}\circ \rho_{k+1}=\rho_k.
\]
于是得到自然映射
\[
\Theta:\widehat X\to \varprojlim_k \widehat X_k^{\mathrm{br}},
\qquad
\Theta(z):=(\rho_k(z))_{k\ge 1}.
\]

\(\Theta\) 的单射性来自前缀决定性：若 \(\Theta(z)=\Theta(z')\)，则二者全部有限前缀一致，从而其 \(\NN^\NN\) 坐标完全一致，故 \(z=z'\)。

再证满射。任取相容族
\[
(\xi_k)_{k\ge 1}\in \varprojlim_k \widehat X_k^{\mathrm{br}},
\qquad
\xi_k=(x^{(k)},a_1^{(k)},\dots,a_k^{(k)}).
\]
由相容性
\[
\pi_{k+1,k}^{\mathrm{br}}(\xi_{k+1})=\xi_k
\]
立得 \(x^{(k)}\) 对 \(k\) 稳定，且对每个固定 \(j\ge 1\)，当 \(k\ge j\) 时 \(a_j^{(k)}\) 亦稳定。记其稳定值为
\[
x_0,\qquad a_j\in \NN\ \ (j\ge 1).
\]
于是得到一个候选点
\[
z:=(x_0,(a_1,a_2,\dots))\in X\times \NN^\NN.
\]
另一方面，对每个 \(k\)，由 \(\xi_k\in \widehat X_k^{\mathrm{br}}\) 可知存在唯一有限历史
\[
(x_0,x_{-1}^{(k)},\dots,x_{-k}^{(k)})\in X^{\leftarrow,k}
\]
满足
\[
\mathcal C_k(x_0,x_{-1}^{(k)},\dots,x_{-k}^{(k)})
=
(x_0,a_1,\dots,a_k).
\]
由 \(\mathcal C_k\) 的递推反演，长度 \(k+1\) 与长度 \(k\) 的历史在前 \(k\) 层必然一致；因此这些有限历史彼此兼容，定义出唯一无限前史
\[
(x_0,x_{-1},x_{-2},\dots)\in X^{\leftarrow}.
\]
按构造有
\[
\mathcal C(x_0,x_{-1},x_{-2},\dots)=z,
\]
故 \(z\in \widehat X\) 且 \(\Theta(z)=(\xi_k)_k\)。因此 \(\Theta\) 为双射。

最后，对任意 \((x,(a_1,a_2,\dots))\in \widehat X\)，由定理 \ref{thm:killo-natural-extension-branch-register} 的右移插入公式，
\[
\widehat T(x,(a_1,a_2,\dots))=(T(x),(\beta(x),a_1,a_2,\dots)).
\]
截断到前 \(k\) 项便得到
\[
\rho_k\bigl(\widehat T(x,(a_1,a_2,\dots))\bigr)
=
\widehat T_k^{\mathrm{br}}(x,a_1,\dots,a_k).
\]
故 \(\Theta\) 将 \(\widehat T\) 识别为逆极限动力学 \(\varprojlim_k \widehat T_k^{\mathrm{br}}\)。证毕。
\end{proof}
\leanpartial{paper\_xi\_branch\_register\_strict\_inverse\_limit}{Skipped because the requested Lean signature packages only abstract propositions, which this round forbids for new paper theorems.}

\begin{theorem}[分支指标寄存器编码的规范共轭不变性]\label{thm:xi-branch-register-coding-conjugacy-invariance}\leanverified{paper\_xi\_branch\_register\_coding\_conjugacy\_invariance}
仍沿用定理 \ref{thm:killo-natural-extension-branch-register} 的设定。若另取一组纤维单射
\[
\iota_y':T^{-1}(y)\hookrightarrow \NN,
\qquad
\beta'(x):=\iota'_{T(x)}(x),
\]
并由此得到编码
\[
\mathcal C':X^{\leftarrow}\to X\times\NN^\NN,
\qquad
\widehat X':=\mathcal C'(X^{\leftarrow}),
\qquad
\widehat T':=\mathcal C'\circ \sigma\circ (\mathcal C')^{-1},
\]
则存在规范双射
\[
\Gamma:\widehat X\to \widehat X'
\]
满足
\[
\Gamma\circ \widehat T=\widehat T'\circ \Gamma.
\]
更具体地，若
\[
\mathcal C^{-1}(x_0,(a_1,a_2,\dots))=(x_0,x_{-1},x_{-2},\dots),
\]
则
\[
\Gamma(x_0,(a_1,a_2,\dots))
=
\bigl(x_0,(\beta'(x_{-1}),\beta'(x_{-2}),\dots)\bigr).
\]
等价地，改变纤维枚举 \((\iota_y)_y\) 只在各条纤维内部实施局部字母重命名，而不改变自然扩张的动力学共轭类。
\end{theorem}
\begin{proof}
定义
\[
\Gamma:=\mathcal C'\circ \mathcal C^{-1}:\widehat X\to \widehat X'.
\]
由于 \(\mathcal C,\mathcal C'\) 都是各自像空间上的双射，\(\Gamma\) 亦为双射。并且
\[
\Gamma\circ \widehat T
=
\mathcal C'\circ \mathcal C^{-1}\circ \mathcal C\circ \sigma\circ \mathcal C^{-1}
=
\mathcal C'\circ \sigma\circ \mathcal C^{-1}
=
\widehat T'\circ \mathcal C'\circ \mathcal C^{-1}
=
\widehat T'\circ \Gamma.
\]
这就得到所需共轭关系。

显式公式只需把 \((x_0,(a_1,a_2,\dots))\) 先经 \(\mathcal C^{-1}\) 恢复为唯一前史
\[
(x_0,x_{-1},x_{-2},\dots)\in X^{\leftarrow},
\]
再送入 \(\mathcal C'\)。因此
\[
\Gamma(x_0,(a_1,a_2,\dots))
=
\mathcal C'(x_0,x_{-1},x_{-2},\dots)
=
\bigl(x_0,(\beta'(x_{-1}),\beta'(x_{-2}),\dots)\bigr).
\]
上述映射完全由两套编码 \(\mathcal C,\mathcal C'\) 决定，因而给出了所需的规范共轭。证毕。
\end{proof}

\leanverified{paper\_xi\_branch\_register\_finite\_past\_minimal\_complete\_prefix}
\begin{theorem}[有限纤维满射的 \texorpdfstring{$n$}{n}-步过去由首端点与前 \texorpdfstring{$n$}{n} 个分支指标完全且唯一决定]\label{thm:xi-branch-register-finite-past-minimal-complete-prefix}
仍沿用定理 \ref{thm:killo-natural-extension-branch-register} 的设定。对每个 \(n\ge 1\)，记
\[
X^{\leftarrow,n}
:=
\{(x_0,x_{-1},\dots,x_{-n})\in X^{n+1}:T(x_{-j})=x_{-j+1}\ \ (1\le j\le n)\},
\]
\[
\mathcal C_n(x_0,x_{-1},\dots,x_{-n})
:=
\bigl(x_0,\beta(x_{-1}),\dots,\beta(x_{-n})\bigr),
\]
\[
\widehat X_n^{\mathrm{br}}
:=
\mathcal C_n(X^{\leftarrow,n})\subset X\times \NN^n.
\]
则 \(\mathcal C_n:X^{\leftarrow,n}\to \widehat X_n^{\mathrm{br}}\) 为双射。等价地，对任意 \(x\in X\)，映射
\[
T^{-n}(x)\longrightarrow \NN^n,
\qquad
y\longmapsto \bigl(\beta(T^{n-1}y),\beta(T^{n-2}y),\dots,\beta(y)\bigr)
\]
在其像上为双射。

特别地，\((x_0,\beta(x_{-1}),\dots,\beta(x_{-n}))\) 构成 \(n\)-步过去的极小完备前缀统计量：若
\[
F:X^{\leftarrow,n}\to Y
\]
满足存在映射 \(G:\widehat X_n^{\mathrm{br}}\to Y\) 使
\[
F=G\circ \mathcal C_n,
\]
并且 \(F\) 在 \(X^{\leftarrow,n}\) 上单射，则 \(G\) 必在 \(\widehat X_n^{\mathrm{br}}\) 上单射。因而不存在任何严格弱于 \(\mathcal C_n\) 而仍能在全部 \(n\)-步前史上保持单射的信息压缩。
\end{theorem}
\begin{proof}
由定理 \ref{thm:killo-natural-extension-branch-register}，无限分支寄存器编码
\[
\mathcal C(x_0,x_{-1},x_{-2},\dots)=\bigl(x_0,(\beta(x_{-1}),\beta(x_{-2}),\dots)\bigr)
\]
可按递推公式唯一反演：
\[
x_{-1}=\iota_{x_0}^{-1}\!\bigl(\beta(x_{-1})\bigr),\qquad
x_{-2}=\iota_{x_{-1}}^{-1}\!\bigl(\beta(x_{-2})\bigr),\qquad \dots
\]
故一旦给定
\[
\bigl(x_0,a_1,\dots,a_n\bigr)\in \widehat X_n^{\mathrm{br}},
\]
就可逐步唯一恢复
\[
x_{-1}=\iota_{x_0}^{-1}(a_1),\qquad
x_{-2}=\iota_{x_{-1}}^{-1}(a_2),\qquad \dots,\qquad
x_{-n}=\iota_{x_{-(n-1)}}^{-1}(a_n).
\]
因此 \(\mathcal C_n\) 在其像上可逆，从而为双射。

对固定 \(x\in X\)，映射
\[
T^{-n}(x)\longrightarrow X^{\leftarrow,n},
\qquad
y\longmapsto \bigl(x,T^{n-1}y,T^{n-2}y,\dots,y\bigr)
\]
显然为双射；再与 \(\mathcal C_n\) 复合，便得到所示从 \(T^{-n}(x)\) 到 \(\NN^n\) 的像上双射。

最后证明极小性。若 \(F=G\circ \mathcal C_n\) 且 \(F\) 单射，取任意
\[
u,v\in \widehat X_n^{\mathrm{br}}
\]
满足 \(G(u)=G(v)\)。由于 \(\mathcal C_n\) 对其像为双射，可取唯一
\[
z,z'\in X^{\leftarrow,n}
\]
使 \(\mathcal C_n(z)=u\)、\(\mathcal C_n(z')=v\)。于是
\[
F(z)=G(\mathcal C_n(z))=G(u)=G(v)=G(\mathcal C_n(z'))=F(z').
\]
由 \(F\) 的单射性得 \(z=z'\)，再由 \(\mathcal C_n\) 的单射性得 \(u=v\)。故 \(G\) 在 \(\widehat X_n^{\mathrm{br}}\) 上亦为单射，这正是所述极小完备性。证毕。
\end{proof}

\begin{corollary}[最坏情形下 \texorpdfstring{$n$}{n}-步过去的 prime-slice 可寻址维数恰等于深度]\label{cor:xi-branch-register-depth-primeslice-addressability}\leanverified{paper\_xi\_branch\_register\_depth\_primeslice\_addressability}
设 \(T:X\twoheadrightarrow X\) 为非单射有限纤维满射，并固定 \(n\ge 1\)。考虑所有单射分层乘法账本
\[
\Psi_n(y)=\bigl(T^n(y),H_n(y)\bigr),
\qquad
H_n(y)=\prod_{j=0}^{n-1}h_j(y),
\]
其中每个 \(h_j(y)\) 的全部素因子均限制在预先指定且两两不交的 prime slice \(\mathcal P_j\)。定义
\[
\operatorname{psupp}(\Psi_n):=\#\{\,0\le j\le n-1:\ h_j\not\equiv 1\,\}.
\]
则
\[
\inf_{\Psi_n\ \mathrm{injective}}\operatorname{psupp}(\Psi_n)=n.
\]
因此，自然扩张的深度不仅是分支寄存器前缀的最小完备长度，也正是最坏情形下不可压缩的 prime-slice 账本维数。
\end{corollary}
\begin{proof}
对长度 \(n\) 的前史塔取
\[
X_0=X_1=\cdots=X_n=X,
\qquad
\pi_0=\pi_1=\cdots=\pi_{n-1}=T.
\]
由于 \(T\) 非单射，存在某个 \(y\in X\) 使 \(\#T^{-1}(y)\ge 2\)，故该塔的每一层都满足“存在真正分支点”的最坏情形假设。于是定理 \ref{thm:xi-layered-prime-register-sharp-natural-extension} 立即给出：任意单射分层 prime-register 外置都必须在全部 \(n\) 层上保留非平凡的 prime slice，而且存在恰好每层写入一枚有效素数的单射实现。故上式成立。

再由定理 \ref{thm:xi-branch-register-finite-past-minimal-complete-prefix}，\(\mathcal C_n\) 已经给出 \(n\)-步过去的极小完备前缀统计量。因此同一深度 \(n\) 同时刻画了加性分支寄存器的最小完备长度与乘法 prime-slice 账本的最小可寻址维数。证毕。
\end{proof}

\leanverified{paper\_xi\_smith\_local\_loss\_profiles\_free\_commutative\_monoid}
\begin{theorem}[\texorpdfstring{$p$}{p}-局部信息亏损剖面的自由交换幺半群结构]\label{thm:xi-smith-local-loss-profiles-free-commutative-monoid}
固定素数 \(p\)。令 \(\mathfrak F_p\) 为所有由有限多重集 \(E\subset \ZZ_{>0}\) 诱导的函数
\[
f_E:\ZZ_{\ge 0}\to \ZZ_{\ge 0},
\qquad
f_E(k):=\sum_{e\in E}\min(k,e),
\]
组成的集合，其中 \(f_E(0)=0\)。赋予 \(\mathfrak F_p\) 逐点加法。对每个 \(t\ge 1\)，定义
\[
u_t(k):=\min(k,t)\qquad(k\ge 0).
\]
则 \((\mathfrak F_p,+)\) 是由 \(\{u_t\}_{t\ge 1}\) 自由生成的交换幺半群。更具体地，每个 \(f\in\mathfrak F_p\) 都唯一表示为
\[
f=\sum_{t\ge 1}m_tu_t,
\qquad
m_t\in \ZZ_{\ge 0},
\]
且仅有限多个 \(m_t\) 非零；其系数可由离散曲率唯一恢复：
\[
m_t=\Delta_f(t)-\Delta_f(t+1),
\qquad
\Delta_f(k):=f(k)-f(k-1).
\]
特别地，任意满行秩整数矩阵的 \(\,p\)-局部 Smith 信息亏损谱都属于 \(\mathfrak F_p\)。
\end{theorem}
\begin{proof}
若 \(E,F\) 为有限多重集，则
\[
f_E+f_F=f_{E\sqcup F},
\]
其中 \(E\sqcup F\) 表示按重数计的并。因此 \(\mathfrak F_p\) 在逐点加法下确为交换幺半群。又由定义
\[
f_E=\sum_{t\ge 1}m_tu_t,
\qquad
m_t:=\#\{e\in E:\ e=t\},
\]
可知 \(\{u_t\}_{t\ge 1}\) 生成 \(\mathfrak F_p\)。

唯一性可由前述 Smith 谱分类直接读出。事实上，定理 \ref{thm:xi-smith-padic-loss-spectrum-classification} 已说明，这类函数恰由非负、非增且最终为零的一阶差分所刻画；而推论 \ref{cor:xi-smith-second-discrete-curvature-recovery} 与定理 \ref{thm:xi-smith-loss-spectrum-discrete-curvature-complete-characterization} 给出
\[
\Delta_f(t)-\Delta_f(t+1)
=
2f(t)-f(t-1)-f(t+1),
\]
其值正是指数 \(t\) 的重数。故一旦
\[
f=\sum_{t\ge 1}m_tu_t=\sum_{t\ge 1}m_t'u_t,
\]
就必有
\[
m_t=\Delta_f(t)-\Delta_f(t+1)=m_t'
\qquad(\forall\,t\ge 1).
\]
这就证明了自由性。最后，若 \(f=f_p\) 为某个满行秩整数矩阵的 \(\,p\)-局部信息亏损谱，则由定理 \ref{thm:killo-smith-loss-spectrum}
\[
f_p(k)=\sum_{i=1}^{m}\min(k,e_i(p)),
\]
故 \(f_p\in \mathfrak F_p\)。证毕。
\end{proof}

\begin{theorem}[局部 Euler 因子对 Smith 正规形的完备性]\label{thm:xi-smith-local-euler-complete-invariant}
\leanverified{paper\_xi\_smith\_local\_euler\_complete\_invariant}
沿用定理 \ref{thm:xi-kernel-loss-divisibility-valuation-rational-euler} 的记号。对每个素数 \(p\)，把局部因子视为形式幂级数
\[
\Xi_{A,p}(T)
:=
\sum_{k\ge 0}p^{f_p(k)}T^k.
\]
则 \(\Xi_{A,p}(T)\) 唯一决定整个序列 \((f_p(k))_{k\ge 1}\)，从而唯一决定多重集 \(\{e_i(p)\}_{i=1}^{m}\)。因而全局 Dirichlet 级数
\[
\Xi_A(s)
:=
\sum_{b\ge 1}\frac{\#\ker(A_b)}{b^{\,n-m+s}}
\]
唯一决定 \(A\) 的全部 Smith 不变量。更一般地，若另一个满行秩整数矩阵
\[
B\in M_{m_B\times n_B}(\ZZ)
\]
满足
\[
\Xi_A=\Xi_B,
\]
则 \(A\) 与 \(B\) 具有相同的 Smith 正规形。
\end{theorem}
\begin{proof}
由定理 \ref{thm:xi-kernel-loss-divisibility-valuation-rational-euler}，
\[
\Xi_{A,p}(T)=\sum_{k\ge 0}p^{f_p(k)}T^k.
\]
因此对每个 \(k\ge 1\)，其 \(T^k\) 项系数恰为 \(p^{f_p(k)}\)，从而
\[
f_p(k)=\log_p\!\bigl([T^k]\Xi_{A,p}(T)\bigr).
\]
局部亏损谱一旦给定，定理 \ref{thm:killo-smith-loss-spectrum} 与定理 \ref{thm:xi-smith-loss-spectrum-discrete-curvature-complete-characterization} 即可无损恢复多重集 \(\{e_i(p)\}_{i=1}^{m}\)。

再看全局唯一性。若 \(\Xi_A=\Xi_B\) 作为 Dirichlet 级数相等，则对每个正整数 \(b\) 都有
\[
\frac{\#\ker(A_b)}{b^{\,n-m}}
=
\frac{\#\ker(B_b)}{b^{\,n_B-m_B}},
\]
亦即两者的归一化系数逐项相同。特别地，对任意素数幂 \(b=p^k\)，得到
\[
p^{f_{A,p}(k)}
=
\frac{\#\ker(A\bmod p^k)}{p^{k(n-m)}}
=
\frac{\#\ker(B\bmod p^k)}{p^{k(n_B-m_B)}}
=
p^{f_{B,p}(k)}.
\]
故
\[
f_{A,p}(k)=f_{B,p}(k)\qquad(\forall\,p,\ \forall\,k\ge 1).
\]
于是对每个素数 \(p\)，两矩阵的 \(\{e_i(p)\}\) 完全一致；逐素数合并便得全部 Smith 不变量一致，亦即 Smith 正规形相同。证毕。
\end{proof}

\begin{corollary}[自由原子结构在整数轴对齐椭圆半群中的几何实现]\label{cor:xi-free-atomic-structure-ellipse-realization}
定理 \ref{thm:xi-smith-local-loss-profiles-free-commutative-monoid} 与定理 \ref{thm:xi-integer-ellipse-free-commutative-monoid} 表明，\(\,p\)-局部信息亏损谱幺半群与整数轴对齐椭圆半群都属于以可数原子集为基的自由交换幺半群：前者的原子为 \(\{u_t\}_{t\ge 1}\)，后者的原子为
\[
\mathcal A=\{E_{p,1},E_{1,p}:p\in\mathcal P\}.
\]
因此
\[
\operatorname{Aut}_{\mathrm{Mon}}(\mathsf E_{\NN},\odot)\cong \mathrm{Sym}(\mathcal A)
\]
可视为附录 \(H\) 中自由原子机制在显式几何对象类上的完全对称实现。
\end{corollary}
\begin{proof}
定理 \ref{thm:xi-smith-local-loss-profiles-free-commutative-monoid} 给出 \(\mathfrak F_p\) 的自由原子分解，定理 \ref{thm:xi-integer-ellipse-free-commutative-monoid} 的第一条给出 \((\mathsf E_{\NN},\odot)\) 的自由原子分解，而其第三条正给出所述自同构群公式。证毕。
\end{proof}

\begin{corollary}[Smith 亏损谱的二阶差分、平台起点与总亏损统一恢复]\label{cor:xi-smith-loss-spectrum-plateau-start-total-loss}
\leanverified{paper\_xi\_smith\_loss\_spectrum\_plateau\_start\_total\_loss}
设 \(A\in M_{m\times n}(\ZZ)\) 满足 \(\operatorname{rank}_{\QQ}(A)=m\)。固定素数 \(p\)，并沿用定理 \ref{thm:killo-smith-loss-spectrum} 的记号：
\[
f_p(0):=0,\qquad
\Delta_p(k):=f_p(k)-f_p(k-1)\quad(k\ge 1),
\]
\[
e_i(p):=v_p(d_i),\qquad
m_{p,j}:=\#\{i:\ e_i(p)=j\}\quad(j\ge 1).
\]
则对每个 \(j\ge 1\) 都有
\[
m_{p,j}
=
\Delta_p(j)-\Delta_p(j+1)
=
2f_p(j)-f_p(j-1)-f_p(j+1).
\]
再记
\[
E_p:=\max_i e_i(p).
\]
则
\[
E_p
=
\max\bigl(\{k\ge 1:\Delta_p(k)>0\}\cup\{0\}\bigr)
=
\min\{k\ge 0:\ f_p(k)=f_p(k+1)\},
\]
并且
\[
\sum_{i=1}^{m}e_i(p)=\lim_{k\to\infty}f_p(k).
\]
换言之，\(f_p\) 是一条最终常值的非减离散凹函数，而其二阶差分、平台起点与平台值已经分别精确编码了各个 \(\,p\)-进高度的重数、最大高度与总亏损。
\end{corollary}
\begin{proof}
二阶差分公式与极限公式分别正是推论 \ref{cor:xi-smith-second-discrete-curvature-recovery} 与推论 \ref{cor:xi-smith-volume-recovered-from-padic-platform} 的内容。

再由推论 \ref{cor:xi-smith-kernel-spectrum-discrete-concavity-plateau-inversion}，\(f_p\) 为非减离散凹函数，并在 \(k\ge E_p\) 后进入平台。另一方面，
\[
\Delta_p(k)=\#\{i:\ e_i(p)\ge k\}
\]
表明 \(\Delta_p(k)>0\) 当且仅当存在某个 \(i\) 满足 \(e_i(p)\ge k\)，亦即 \(k\le E_p\)。故
\[
E_p=\max\bigl(\{k\ge 1:\Delta_p(k)>0\}\cup\{0\}\bigr).
\]
又因为
\[
f_p(k+1)-f_p(k)=\Delta_p(k+1),
\]
故 \(f_p(k)=f_p(k+1)\) 当且仅当 \(\Delta_p(k+1)=0\)，其最早发生恰在 \(k=E_p\)。这就得到第二个恢复公式。证毕。
\end{proof}

\begin{theorem}[全模数核计数的严格 Smith 饱和平台]\label{thm:xi-smith-global-kernel-dmax-saturation-platform}\leanverified{paper\_xi\_smith\_global\_kernel\_dmax\_saturation\_platform}
设 \(A\in M_{m\times n}(\ZZ)\) 满足 \(\operatorname{rank}_{\QQ}(A)=m\)，其 Smith 不变量为
\[
d_1\mid d_2\mid\cdots\mid d_m.
\]
对每个 \(N\ge 1\)，记
\[
A_N:(\ZZ/N\ZZ)^n\to(\ZZ/N\ZZ)^m,
\qquad
K_A(N):=\#\ker(A_N).
\]
若记 \(e_i(p):=v_p(d_i)\)、\(f_p(k):=\sum_{i=1}^{m}\min(k,e_i(p))\)，并设
\[
N=\prod_p p^{\nu_p(N)},
\qquad
\Delta_A:=\prod_{i=1}^{m}d_i,
\qquad
d_{\max}:=d_m,
\]
则
\[
K_A(N)
=
N^{\,n-m}\prod_{i=1}^{m}\gcd(d_i,N)
=
N^{\,n-m}\prod_{p}p^{\,f_p(\nu_p(N))}.
\]
特别地，
\[
K_A(N)\le N^{\,n-m}\Delta_A,
\]
并且等号成立当且仅当
\[
d_{\max}\mid N.
\]
因而全模数核计数在整除格上存在一条严格的 Smith 饱和平台：
\[
K_A(N)=N^{\,n-m}\Delta_A
\qquad(d_{\max}\mid N),
\]
并且 \(d_{\max}\) 正是使全部局部亏损同时饱和的最小阈值。
\end{theorem}
\begin{proof}
由推论 \ref{cor:killo-kernel-size-general-modulus}，
\[
K_A(N)=N^{\,n-m}\prod_{i=1}^{m}\gcd(d_i,N).
\]
再把每个 \(\gcd(d_i,N)\) 按素数分解，便得
\[
\prod_{i=1}^{m}\gcd(d_i,N)
=
\prod_{p}p^{\sum_{i=1}^{m}\min(\nu_p(N),e_i(p))}
=
\prod_{p}p^{f_p(\nu_p(N))}.
\]
这证明了第一式。

对每个 \(i\) 都有 \(\gcd(d_i,N)\le d_i\)，故
\[
K_A(N)\le N^{\,n-m}\prod_{i=1}^{m}d_i
=
N^{\,n-m}\Delta_A.
\]
并且取等当且仅当
\[
\gcd(d_i,N)=d_i
\qquad(1\le i\le m),
\]
亦即当且仅当每个 \(d_i\mid N\)。由于
\[
d_1\mid d_2\mid\cdots\mid d_m,
\]
上述条件又等价于 \(d_m\mid N\)，即 \(d_{\max}\mid N\)。

因此达到上界的模数恰为 \(d_{\max}\) 的倍数，故一旦模数吸收最大 Smith 不变量，全部局部亏损便同时饱和；反之若 \(d_{\max}\nmid N\)，则至少有一个因子 \(\gcd(d_i,N)\) 严格小于 \(d_i\)，从而上界不能取到。故 \(d_{\max}\) 正是所述最小饱和阈值。证毕。
\end{proof}

\noindent
由此，附录 \(H\) 中的结构链可进一步收束为
\[
\text{有限深度分支寄存器}
\Longrightarrow
\text{\(n\)-步过去的极小完备前缀}
\Longrightarrow
\text{自然扩张的严格逆极限},
\]
\[
\text{自然扩张深度}
\Longleftrightarrow
\text{最坏 prime-slice 账本维数},
\]
\[
\text{Smith 局部亏损谱}
\Longrightarrow
\text{二阶曲率与平台起点恢复}
\Longrightarrow
\text{自由原子分解}
\Longrightarrow
\text{Euler 完备不变量}
\Longrightarrow
\text{全模数饱和阈值},
\]
而推论 \ref{cor:xi-free-atomic-structure-ellipse-realization} 则说明同一自由原子机制在整数轴对齐椭圆半群上获得了显式几何实现与全对称自同构群；定理 \ref{thm:xi-smith-global-kernel-dmax-saturation-platform} 则进一步表明，这条局部到全局的链最终由单一最大 Smith 不变量锁定其整除阈值。

\endinput

\paragraph{分支指标像空间的坐标逆极限与全深度 prime-register 的无限秩障碍}
在定理 \ref{thm:xi-branch-register-strict-inverse-limit} 已把自然扩张的分支寄存器模型识别为有限深度截断的严格逆极限之后，还可把这一事实重写为像空间坐标本身的逆极限完备性；进一步，当每一层都存在真实分支时，忠实的全深度 prime-register 账本便不再是有限层 prime-slice 最优化的简单延长，而会在结构层上强制出现可数无穷个彼此独立的素数方向。

\begin{theorem}[分支指标像空间本身就是有限截断坐标系的规范逆极限]\label{thm:xi-branch-register-image-coordinate-inverse-limit}
\leanverified{paper\_xi\_branch\_register\_image\_coordinate\_inverse\_limit}
沿用定理 \ref{thm:killo-natural-extension-branch-register} 的记号。对每个 \(n\ge 1\)，定义前缀投影
\[
\mathrm{pr}_{\le n}:\widehat X\to X\times \NN^n,
\qquad
\mathrm{pr}_{\le n}\bigl(x,(a_1,a_2,\dots)\bigr):=(x,a_1,\dots,a_n),
\]
并记
\[
Y_n:=\mathrm{pr}_{\le n}(\widehat X)
=
\Bigl\{(x_0,a_1,\dots,a_n)\in X\times\NN^n:\ \exists\,\widehat x\in\widehat X,
\mathrm{pr}_{\le n}(\widehat x)=(x_0,a_1,\dots,a_n)\Bigr\}.
\]
定义截断映射
\[
\rho_{n+1,n}:Y_{n+1}\to Y_n,
\qquad
\rho_{n+1,n}(x_0,a_1,\dots,a_{n+1})=(x_0,a_1,\dots,a_n).
\]
则规范映射
\[
\Theta:\widehat X\to \varprojlim(Y_n,\rho_{n+1,n}),
\qquad
\Theta(\widehat x):=\bigl(\mathrm{pr}_{\le n}(\widehat x)\bigr)_{n\ge 1},
\]
为双射。

进一步，对每个 \(n\ge 1\) 定义有限深度右移插入系统
\[
\widehat T_n:Y_n\to Y_n,
\qquad
\widehat T_n(x_0,a_1,\dots,a_n):=
\begin{cases}
(T(x_0),\beta(x_0)),& n=1,\\
(T(x_0),\beta(x_0),a_1,\dots,a_{n-1}),& n\ge 2.
\end{cases}
\]
则在上述识别下，
\[
\widehat T=\varprojlim_n \widehat T_n.
\]
换言之，分支指标模型 \(\widehat X\) 本身就是有限深度右移插入系统的逆极限对象。
\end{theorem}
\begin{proof}
由定理 \ref{thm:xi-branch-register-finite-past-minimal-complete-prefix}，\(n\)-步过去的分支寄存器前缀
\[
(x_0,a_1,\dots,a_n)
\]
与唯一有限前史
\[
(x_0,x_{-1},\dots,x_{-n})\in X^{\leftarrow,n}
\]
一一对应。反过来，由 \(T\) 为满射，任意有限前史都可继续向后逐层选取原像，因而总能延拓为某个无限前史。故 \(Y_n\) 与定理 \ref{thm:xi-branch-register-strict-inverse-limit} 中的
\[
\widehat X_n^{\mathrm{br}}=\mathcal C_n(X^{\leftarrow,n})
\]
是同一对象，只是记号改写而已。

\(\Theta\) 的单射性是直接的：若两点在全部有限前缀上相同，则其 \(\widehat X\subset X\times\NN^\NN\) 中的全部坐标都相同，因而两点一致。

再证满射。任取一致族
\[
y^{(n)}=(x_0^{(n)},a_1^{(n)},\dots,a_n^{(n)})\in Y_n,
\qquad
\rho_{n+1,n}(y^{(n+1)})=y^{(n)}.
\]
由一致性可知 \(x_0^{(n)}\) 稳定，且对每个固定 \(j\ge 1\)，当 \(n\ge j\) 时 \(a_j^{(n)}\) 亦稳定。记稳定值为
\[
x_0,\qquad a_j\in\NN\quad (j\ge 1).
\]
由于每个 \(y^{(n)}\in Y_n\)，由定理 \ref{thm:xi-branch-register-finite-past-minimal-complete-prefix} 可恢复唯一有限前史
\[
(x_0,x_{-1}^{(n)},\dots,x_{-n}^{(n)})\in X^{\leftarrow,n}.
\]
而 \(\rho_{n+1,n}(y^{(n+1)})=y^{(n)}\) 又迫使这些有限前史在重叠部分完全一致，故它们定义出唯一无限前史
\[
(x_0,x_{-1},x_{-2},\dots)\in X^{\leftarrow}.
\]
将其送入定理 \ref{thm:killo-natural-extension-branch-register} 的编码映射，得到
\[
\widehat x:=\mathcal C(x_0,x_{-1},x_{-2},\dots)\in \widehat X,
\]
并满足
\[
\mathrm{pr}_{\le n}(\widehat x)=y^{(n)}\qquad(\forall n\ge 1).
\]
故 \(\Theta\) 为满射。

最后，由定理 \ref{thm:killo-natural-extension-branch-register} 的右移插入公式
\[
\widehat T(x,(a_1,a_2,\dots))=(T(x),(\beta(x),a_1,a_2,\dots)),
\]
对每个 \(n\) 取前缀便得到
\[
\mathrm{pr}_{\le n}\!\bigl(\widehat T(\widehat x)\bigr)
=
\widehat T_n\!\bigl(\mathrm{pr}_{\le n}(\widehat x)\bigr).
\]
因此 \(\Theta\) 将 \(\widehat T\) 精确识别为有限深度右移插入系统的逆极限动力学。证毕。
\end{proof}

\begin{theorem}[忠实的全深度 prime-register 实现必然具有无限素数秩，因而排斥任何有限生成加法账本]\label{thm:xi-infinite-depth-prime-ledger-infinite-rank-obstruction}
\leanverified{paper\_xi\_infinite\_depth\_prime\_ledger\_infinite\_rank\_obstruction}
设 \(T:X\twoheadrightarrow X\) 为有限纤维满射，并假设其逆向分支深度无界：对每个 \(N\ge 1\)，都存在长度 \(N\) 的前史片段
\[
x_{-N}\xrightarrow{\,T\,}x_{-N+1}\xrightarrow{\,T\,}\cdots\xrightarrow{\,T\,}x_{-1}\xrightarrow{\,T\,}x_0
\]
满足
\[
\#T^{-1}(x_{-j+1})\ge 2
\qquad(1\le j\le N).
\]
称一个全深度 prime-register 账本实现为忠实，若它对每个有限深度 \(N\) 的前缀截断都与定理 \ref{thm:xi-layered-prime-register-sharp-natural-extension} 的精确分层外置语义相容。

则任意忠实的全深度实现都必须使用无穷多个彼此独立的 prime slices。更具体地，若从每一层的有效 prime slice 中选取一枚代表素数
\[
q_j\in \mathcal P_j\qquad (j\ge 0),
\]
并定义自由交换幺半群
\[
M_\infty
:=
\left\{
\prod_{j\ge 0}q_j^{a_j}:\ a_j\in\ZZ_{\ge 0},\ a_j=0\ \text{仅有限多个例外}
\right\},
\]
则 \(M_\infty\) 在唯一分解意义下同构于 \(\NN_{\times}\)。

特别地，不存在任何有限生成交换群 \(L\) 使忠实的全深度 prime-register 账本能够严格同态化到
\[
(\ZZ\oplus L,+)
\]
这样的有限秩加法账本上。
\end{theorem}
\begin{proof}
第一部分正是定理 \ref{thm:xi-natural-extension-infinite-prime-register-necessity} 的内容：在逆向分支深度无界时，任何精确实现整个自然扩张的分层 prime-register 外置都不可能只使用有限多个 prime slices。结合上一定理，\(\widehat X\) 已被识别为全部有限深度前缀系统的逆极限，因此忠实的全深度账本若与每个有限前缀的精确恢复相容，就必须在每一层保留一个独立的有效 prime direction。

由 \((q_j)_{j\ge 0}\) 两两不同且都是素数，唯一分解给出幺半群同构
\[
\vartheta:\NN_{\times}\xrightarrow{\ \sim\ } M_\infty,
\]
其在素数生成元上由
\[
\vartheta(p_{j+1})=q_j
\qquad(j\ge 0)
\]
确定。

反设存在有限生成交换群 \(L\)，使忠实的全深度账本可严格同态化到 \((\ZZ\oplus L,+)\) 上。则其 restriction 在上述独立素数方向上给出一个幺半群单射
\[
M_\infty\hookrightarrow (\ZZ\oplus L,+).
\]
与 \(\vartheta\) 复合便得到
\[
\NN_{\times}\hookrightarrow (\ZZ\oplus L,+)
\]
的单射幺半群同态。
但 \((\ZZ\oplus L,+)\) 是有限生成交换可消去幺半群，这与定理 \ref{thm:killo-prime-freedom-non-finitizability} 矛盾；等价地，这也与定理 \ref{thm:killo-godel-compression-not-finite-rank-homologizable} 所排除的有限秩同态化账本机制相冲突。故此类有限生成加法账本不存在。证毕。
\end{proof}

\noindent
配合定理 \ref{thm:xi-terminal-fixed-freezing-tv-collapse}、定理 \ref{thm:xi-fixed-freezing-escort-bounded-observable-collapse} 以及推论 \ref{cor:xi-time-part57bb-freezing-groundstate-subset-equidistribution} 可见：自然扩张首先在分支指标像空间层面构成严格逆极限对象；冻结 escort 态随后在统计层面指数塌缩到最大纤维上的微正则分布；而一旦要求把这种全深度可逆账本忠实地外置为 prime-register，其结构层便不可避免地转入可数无穷素数秩，从而排斥一切有限秩加法线性化。

\endinput

\paragraph{自然扩张前缀的 Gödel 算术化与 \texorpdfstring{$p$}{p}-局部核增长的单尾极点}
在自然扩张分支寄存器的有限前缀、prime-slice 历史压缩与 Smith 局部亏损谱的 Euler 因子化都已建立之后，还可把有限前史地址与 \(p\)-局部核增长的解析尾部进一步压缩为下列两条接口。

\begin{corollary}[分层素数外置恰是自然扩张前缀地址的 Gödel 化]\label{cor:xi-time-part62ba-primeslice-natural-extension-prefix-godelization}
\leanverified{paper\_xi\_time\_part62ba\_primeslice\_natural\_extension\_prefix\_godelization}
设 \(X\) 为可数集合，\(T:X\twoheadrightarrow X\) 为有限纤维满射，并沿用定理 \ref{thm:killo-natural-extension-branch-register}、定理 \ref{thm:xi-prime-register-history-inverse-limit} 与定理 \ref{thm:xi-branch-register-finite-past-minimal-complete-prefix} 的记号。对任意 \(k\ge 1\)，在自映射前史塔
\[
X_0=X_1=\cdots=X_k=X,
\qquad
\pi_0=\pi_1=\cdots=\pi_{k-1}=T
\]
上，定义长度 \(k\) 的截断地址
\[
\mathcal C_k(x_0,x_{-1},\dots,x_{-k})
:=
\bigl(x_0,\beta(x_{-1}),\dots,\beta(x_{-k})\bigr).
\]
再取定定理 \ref{thm:xi-prime-register-history-inverse-limit} 中的分层单射
\[
\alpha_\ell:X^{\leftarrow,\ell-1}\times\NN\hookrightarrow \mathcal P_{\ell-1}
\qquad (1\le \ell\le k),
\]
并令
\[
G_k(x_0,\dots,x_{-k})
:=
\prod_{\ell=1}^{k}
\alpha_\ell\!\bigl((x_0,\dots,x_{-\ell+1}),\beta(x_{-\ell})\bigr).
\]
则对任意固定首端点 \(x_0\in X\)，长度 \(k\) 的逆轨道前缀
\[
(x_{-1},\dots,x_{-k})
\]
与整数 \(G_k(x_0,\dots,x_{-k})\) 之间构成双射。等价地，定理 \ref{thm:killo-layered-prime-slices-finite-depth} 的分层素数外置，在自映射前史塔上恰是长度 \(k\) 的自然扩张地址 \(\mathcal C_k\) 的 Gödel 编码。
\end{corollary}
\begin{proof}
定理 \ref{thm:xi-branch-register-finite-past-minimal-complete-prefix} 说明：对固定 \(x_0\) 与 \(k\)，截断地址 \(\mathcal C_k\) 与长度 \(k\) 的前史前缀彼此唯一决定。另一方面，定理 \ref{thm:xi-prime-register-history-inverse-limit} 已证明：由 \((x_0,G_k)\) 可确定性恢复唯一历史 \((x_0,\dots,x_{-k})\)，等价地可恢复唯一的 \(\mathcal C_k(x_0,\dots,x_{-k})\)。因此在固定首端点 \(x_0\) 后，整数 \(G_k\) 与长度 \(k\) 的逆轨道前缀双向确定，结论成立。证毕。
\end{proof}

\begin{theorem}[\texorpdfstring{$p$}{p}-局部核增长序列的局部 \texorpdfstring{$\zeta$}{zeta} 函数只有唯一尾极点]\label{thm:xi-time-part62ba-padic-kernel-growth-local-zeta-single-tail-pole}
\leanverified{paper\_xi\_time\_part62ba\_padic\_kernel\_growth\_local\_zeta\_single\_tail\_pole}
设 \(A\in M_{m\times n}(\ZZ)\) 满足 \(\operatorname{rank}_{\QQ}(A)=m\)，其 Smith 不变量为
\[
d_1\mid d_2\mid\cdots\mid d_m.
\]
固定素数 \(p\)，并定义
\[
K_p(k):=\#\ker(A\bmod p^k)
\qquad (k\ge 1),
\]
\[
K_p(0):=1,
\]
\[
e_i(p):=v_p(d_i)
\qquad (1\le i\le m),
\qquad
E_p:=\max_{1\le i\le m}e_i(p).
\]
再定义 \(p\)-局部核增长的 ordinary generating function
\[
\mathcal Z_{A,p}^{\mathrm{loc}}(z)
:=
\sum_{k\ge 0}K_p(k)z^k.
\]
则 \(\mathcal Z_{A,p}^{\mathrm{loc}}\) 为有理函数，并且具有显式分解
\[
\boxed{\;
\mathcal Z_{A,p}^{\mathrm{loc}}(z)
=
\sum_{k=0}^{E_p-1}K_p(k)z^k
+
\frac{K_p(E_p)\,z^{E_p}}{1-p^{\,n-m}z}.\;}
\]
换言之，\(p\)-局部核增长在有限窗口 \(\{0,1,\dots,E_p\}\) 之后进入严格几何尾部；其尾部唯一极点位于
\[
\boxed{\;
z=p^{-(n-m)}.\;}
\]
\end{theorem}
\begin{proof}
由定理 \ref{thm:killo-smith-loss-spectrum}，对每个 \(k\ge 1\) 都有
\[
K_p(k)
=
p^{\,k(n-m)+\sum_{i=1}^{m}\min(k,e_i(p))}.
\]
若 \(E_p=0\)，则一切 \(e_i(p)=0\)，从而
\[
K_p(k)=p^{k(n-m)}
\qquad (k\ge 0),
\]
于是
\[
\mathcal Z_{A,p}^{\mathrm{loc}}(z)=\frac{1}{1-p^{\,n-m}z},
\]
结论显然成立。以下设 \(E_p\ge 1\)。

当 \(k\ge E_p\) 时，对一切 \(i\) 都有 \(\min(k,e_i(p))=e_i(p)\)，故
\[
K_p(k)
=
p^{\,k(n-m)+\sum_{i=1}^{m}e_i(p)}
=
K_p(E_p)\,p^{(k-E_p)(n-m)}.
\]
于是尾部级数成为公比 \(p^{n-m}z\) 的几何级数：
\[
\sum_{k\ge E_p}K_p(k)z^k
=
K_p(E_p)z^{E_p}\sum_{r\ge 0}\bigl(p^{n-m}z\bigr)^r
=
\frac{K_p(E_p)\,z^{E_p}}{1-p^{\,n-m}z}.
\]
再加上前 \(E_p\) 项，即得所示“有限前缀 + 几何尾部”分解。由于分母只有 \(1-p^{\,n-m}z\) 这一项，故尾部唯一极点恰位于 \(z=p^{-(n-m)}\)。证毕。
\end{proof}

\noindent
因此，prime-slice 外置与自然扩张地址之间的关系并不是两套并列装置，而是同一有限前史证书的加性与乘法两种坐标；与之平行，Smith 的 \(p\)-局部核增长也把全部非平凡复杂性压缩到有限前缀，而把尾部解析结构严格收束为唯一一个几何极点。

\endinput

\paragraph{Fibonacci 递推 Fourier 采样、零点缩放均匀律与局部化 solenoid 的自覆盖分类}
Pisot--Bernoulli Fourier 采样、Fibonacci 共振整函数的零点缩放统计，以及有限素数局部化 solenoid 的内部自同态结构，还可进一步压缩为下列四条可直接调用的结论。

\begin{theorem}[Pisot--Bernoulli Fourier 采样在任意非零 Fibonacci 递推方向上不发生 Rajchman 衰减]\label{thm:xi-pisot-bernoulli-fibonacci-direction-nonrajchman}
\leanverified{paper\_xi\_pisot\_bernoulli\_fibonacci\_direction\_nonrajchman}
令 \(q:=\varphi^{-1}\in(0,1)\)，并令 \(\mu_q^{(2)}\) 为定理 \ref{thm:fold-pisot-bernoulli-convolution-representation} 中的 Pisot 参数 Bernoulli 卷积移位测度。取任意非零整数初值对
\[
(u_0,u_1)\in \ZZ^2\setminus\{(0,0)\},
\]
并令 \((u_n)_{n\ge 0}\) 满足 Fibonacci 递推
\[
u_{n+2}=u_{n+1}+u_n.
\]
再记
\[
\psi:=-\varphi^{-1},\qquad
\alpha:=\frac{u_1-u_0\psi}{\sqrt5},\qquad
\beta:=\frac{u_0\varphi-u_1}{\sqrt5}.
\]
则极限
\[
\lim_{n\to\infty}\Bigl|\widehat{\mu_q^{(2)}}(\pi u_n)\Bigr|
\]
存在且严格为正，并满足显式公式
\[
\boxed{\;
\lim_{n\to\infty}\Bigl|\widehat{\mu_q^{(2)}}(\pi u_n)\Bigr|
=
c(u_0,u_1)
:=
\Bigl|\cos(\pi\alpha)\Bigr|
\prod_{k=1}^{\infty}\Bigl|\cos(\pi\alpha\varphi^{-k})\Bigr|
\prod_{k=1}^{\infty}
\Bigl|\cos(\pi\beta\varphi^{-k})\Bigr|
>0.
\;}
\]
此外 \(|u_n|\to\infty\)。若 \(\nu\) 记 \(\mu_q^{(2)}\) 在 \(\mathbb T_2:=\RR/(2\ZZ)\) 上的模 \(2\) 周期化测度，则
\[
\Bigl|\widehat{\nu}(u_n)\Bigr|
=
\Bigl|\widehat{\mu_q^{(2)}}(\pi u_n)\Bigr|
\longrightarrow c(u_0,u_1)>0,
\]
从而 \(\nu\) 沿任意非平凡 Fibonacci 递推逃逸序列都不发生 Rajchman 衰减。
\end{theorem}
\begin{proof}
由定理 \ref{thm:fold-pisot-bernoulli-convolution-representation}，
\[
\Bigl|\widehat{\mu_q^{(2)}}(\pi u)\Bigr|=C_\varphi(u)
\qquad(u\in\ZZ).
\]
另一方面，定理 \ref{thm:fold-resonance-ladder-fibonacci-directional-limit} 对任意非零初值对给出
\[
\lim_{n\to\infty}C_\varphi(u_n)
=
\Bigl|\cos(\pi\alpha)\Bigr|
\prod_{k=1}^{\infty}\Bigl|\cos(\pi\alpha\varphi^{-k})\Bigr|
\prod_{k=1}^{\infty}
\Bigl|\cos(\pi\beta\varphi^{-k})\Bigr|
>0.
\]
两式合并即得所示 Fourier 采样极限公式。

再证 \(|u_n|\to\infty\)。由 Binet 公式，
\[
u_n=\alpha\varphi^n+\beta\psi^n.
\]
若 \(\alpha=0\)，则
\[
u_1=u_0\psi.
\]
由于 \(\psi=-\varphi^{-1}\notin \QQ\)，这与 \(u_0,u_1\in\ZZ\) 同时成立只能推出 \(u_0=u_1=0\)，与假设矛盾。故 \(\alpha\neq 0\)。于是
\[
u_n=\alpha\varphi^n\Bigl(1+\frac{\beta}{\alpha}(\psi/\varphi)^n\Bigr)
=
\alpha\varphi^n(1+o(1)),
\]
从而
\[
|u_n|=|\alpha|\varphi^n(1+o(1))\to\infty.
\]

最后，由推论 \ref{cor:fold-bernoulli-fourier-sampling-not-l2} 的证明，\(\nu\) 的 Fourier 系数满足
\[
\widehat{\nu}(u)=\widehat{\mu_q^{(2)}}(\pi u)
\qquad(u\in\ZZ).
\]
既然 \(|u_n|\to\infty\) 且 \(|\widehat{\mu_q^{(2)}}(\pi u_n)|\to c(u_0,u_1)>0\)，便得到
\[
\Bigl|\widehat{\nu}(u_n)\Bigr|\to c(u_0,u_1)>0,
\]
从而 \(\nu\) 沿该逃逸序列不发生 Rajchman 衰减。证毕。
\end{proof}

\begin{corollary}[任意非零 Fibonacci 递推方向上的对数级平方能量律]\label{cor:xi-pisot-bernoulli-fibonacci-direction-log-energy}
沿用定理 \ref{thm:xi-pisot-bernoulli-fibonacci-direction-nonrajchman} 的记号，并记
\[
c(u_0,u_1)
:=
\lim_{n\to\infty}\Bigl|\widehat{\mu_q^{(2)}}(\pi u_n)\Bigr|
>0.
\]
则当 \(N\to\infty\) 时，
\[
\boxed{\;
\sum_{\substack{n\ge 0\\ |u_n|\le N}}
\Bigl|\widehat{\mu_q^{(2)}}(\pi u_n)\Bigr|^2
=
\bigl(c(u_0,u_1)^2+o(1)\bigr)\log_\varphi N.
\;}
\]
特别地，
\[
\sum_{n\ge 0}\Bigl|\widehat{\mu_q^{(2)}}(\pi u_n)\Bigr|^2=+\infty.
\]
\end{corollary}
\begin{proof}
由定理 \ref{thm:xi-pisot-bernoulli-fibonacci-direction-nonrajchman}，
\[
u_n=\alpha\varphi^n+\beta\psi^n,
\qquad
\alpha\neq 0,
\]
因而
\[
|u_n|=|\alpha|\varphi^n(1+o(1)).
\]
设
\[
M(N):=\#\{n\ge 0:\ |u_n|\le N\}.
\]
则上式等价于
\[
M(N)=\log_\varphi N-\log_\varphi|\alpha|+O(1)
=
\log_\varphi N+O(1).
\]

另一方面，由
\[
\Bigl|\widehat{\mu_q^{(2)}}(\pi u_n)\Bigr|^2\to c(u_0,u_1)^2
\]
可知其部分和满足 Ces\`aro 渐近
\[
\sum_{n=0}^{M}\Bigl|\widehat{\mu_q^{(2)}}(\pi u_n)\Bigr|^2
=
c(u_0,u_1)^2\,M+o(M)
\qquad(M\to\infty).
\]
取 \(M=M(N)\) 即得
\[
\sum_{\substack{n\ge 0\\ |u_n|\le N}}
\Bigl|\widehat{\mu_q^{(2)}}(\pi u_n)\Bigr|^2
=
c(u_0,u_1)^2\,M(N)+o(M(N))
=
\bigl(c(u_0,u_1)^2+o(1)\bigr)\log_\varphi N.
\]
由于右端随 \(N\to\infty\) 对数发散，第二条随即成立。证毕。
\end{proof}

\begin{theorem}[Fibonacci 共振整函数的缩放零点计数测度趋于区间均匀分布]\label{thm:xi-fold-resonance-scaled-zero-measure-uniform}
沿用定理 \ref{thm:fold-resonance-entire-lp} 与定理 \ref{thm:xi-fold-resonance-zero-count-selfsimilar-linear} 的记号，记
\[
\mathcal F_\varphi(z):=\prod_{n=2}^{\infty}\cos(\pi z\varphi^{-n}),
\]
\[
N_\varphi(R):=\#\bigl(\mathcal Z(\mathcal F_\varphi)\cap[-R,R]\bigr),
\qquad
\nu_R:=\frac{1}{N_\varphi(R)}\sum_{\substack{\zeta\in\mathcal Z(\mathcal F_\varphi)\\ |\zeta|\le R}}\delta_{\zeta/R}.
\]
则当 \(R\to\infty\) 时，
\[
\boxed{\;
\nu_R \Longrightarrow \frac12\,\mathbf 1_{[-1,1]}(x)\,\dd x.
\;}
\]
特别地，
\[
\boxed{\;
N_\varphi(R)=2R+O(\log R).\;}
\]
\end{theorem}
\begin{proof}
由定理 \ref{thm:fold-resonance-entire-lp}，
\[
\mathcal Z(\mathcal F_\varphi)
=
\Bigl\{(k+\tfrac12)\varphi^n:\ k\in\ZZ,\ n\ge 2\Bigr\}.
\]
固定闭区间 \([a,b]\subset(-1,1)\)。对每个固定层 \(n\ge 2\)，落在区间 \([aR,bR]\) 中的零点个数等于满足
\[
aR\le (k+\tfrac12)\varphi^n\le bR
\]
的整数 \(k\) 的个数，因此为
\[
(b-a)R\varphi^{-n}+O(1),
\]
其中误差仅来自端点取整。只有当 \(\varphi^n\lesssim R\) 时这一层才可能贡献非零点，故有效层数为 \(O(\log R)\)。于是求和得到
\[
\#\bigl(\mathcal Z(\mathcal F_\varphi)\cap[aR,bR]\bigr)
=
(b-a)R\sum_{n=2}^{\infty}\varphi^{-n}+O(\log R).
\]
由于
\[
\sum_{n=2}^{\infty}\varphi^{-n}=1,
\]
遂得
\[
\#\bigl(\mathcal Z(\mathcal F_\varphi)\cap[aR,bR]\bigr)
=
(b-a)R+O(\log R).
\]
取 \([a,b]=[-1,1]\) 即得
\[
N_\varphi(R)=2R+O(\log R),
\]
这与定理 \ref{thm:xi-fold-resonance-zero-count-selfsimilar-linear} 的线性主项一致。

于是对任意 \([a,b]\subset(-1,1)\)，有
\[
\nu_R([a,b])
=
\frac{(b-a)R+O(\log R)}{2R+O(\log R)}
\longrightarrow
\frac{b-a}{2}.
\]
而 \(\frac12\mathbf 1_{[-1,1]}(x)\,\dd x\) 对区间端点不赋质量，故由 Portmanteau 定理可知
\[
\nu_R \Longrightarrow \frac12\,\mathbf 1_{[-1,1]}(x)\,\dd x.
\]
证毕。
\end{proof}

\begin{theorem}[局部化 solenoid 的连续自覆盖分类与精确度数公式]\label{thm:xi-localized-solenoid-self-cover-degree-classification}
设 \(S\subset\mathbb P\) 为有限素数集，并记
\[
G_S:=\ZZ[S^{-1}]
\]
及其 Pontryagin 对偶 \(\widehat{G_S}\)。则每个连续群自同态
\[
f:\widehat{G_S}\to \widehat{G_S}
\]
都唯一对应于某个 \(r\in G_S\)，其对偶离散同态为
\[
f^\vee(x)=rx
\qquad(x\in G_S).
\]
其中有限度连续自覆盖恰对应于 \(r\neq 0\) 的情形。对任意非零 \(r\in G_S\)，可唯一写成
\[
r=u\,n,
\qquad
u\in G_S^\times=\ZZ[S^{-1}]^\times,
\qquad
n\in\NN,
\qquad
\gcd\!\Bigl(n,\prod_{p\in S}p\Bigr)=1.
\]
若 \(f=f_r\) 为对应的连续自覆盖，则
\[
\boxed{\;
\ker(f)\cong \ZZ/n\ZZ,
\qquad
\deg(f)=n.
\;}
\]
特别地，
\[
\boxed{\;
\{\deg(f): f\ \text{为}\ \widehat{G_S}\ \text{的有限度连续自覆盖}\}
=
\{n\in\NN:\ p\mid n \Rightarrow p\notin S\},
\;}
\]
并且
\[
\boxed{\;
\operatorname{Aut}_{\mathrm{cts}}(\widehat{G_S})
\cong
\ZZ[S^{-1}]^\times
=
\left\{\pm\prod_{p\in S}p^{m_p}:m_p\in\ZZ\right\}.
\;}
\]
等价地，\(f\) 为同构当且仅当 \(n=1\)。
\end{theorem}
\begin{proof}
由定理 \ref{thm:xi-localized-endomorphism-ring-solenoid-indecomposable} 与推论 \ref{cor:xi-localized-integers-dual-automorphism-units}，
\[
\operatorname{End}_{\mathrm{cts}}(\widehat{G_S})
\cong
\operatorname{End}(G_S)^{\mathrm{op}}
\cong
G_S.
\]
因此每个连续自同态都唯一对应于某个 \(r\in G_S\)，并在对偶离散群上表现为乘法
\[
x\mapsto rx.
\]
若 \(r=0\)，则 \(f\) 为常值映射，显然不是自覆盖。以下设 \(r\neq 0\)。

由于 \(G_S=\ZZ[S^{-1}]\subset \QQ\)，\(r\) 的每个素数赋值都良定义；把属于 \(S\) 的全部素数赋值并入单位项，便得到唯一分解
\[
r=u\,n,
\]
其中 \(u\in G_S^\times\)，而 \(n\) 为正的 \(S\)-自由整数。由于乘以非零 \(r\) 在无挠群 \(G_S\subset\QQ\) 上单射，故 \(f^\vee\) 为单射。Pontryagin 对偶把单射送为满射，于是 \(f\) 为满同态。

又因 \(u\) 为单位，
\[
rG_S=u(nG_S)=nG_S.
\]
因此由定理 \ref{thm:xi-localized-finite-quotient-phase-layer-classification}，
\[
G_S/rG_S
\cong
G_S/nG_S
\cong
\ZZ/n\ZZ.
\]
对离散短正合列
\[
0\longrightarrow G_S\xrightarrow{\ \times r\ }G_S\longrightarrow G_S/rG_S\longrightarrow 0
\]
作 Pontryagin 对偶，得到
\[
0\longrightarrow (G_S/rG_S)^\wedge\longrightarrow \widehat{G_S}\xrightarrow{\ f\ }\widehat{G_S}\longrightarrow 0,
\]
从而
\[
\ker(f)\cong (G_S/rG_S)^\wedge\cong \ZZ/n\ZZ.
\]
特别地，\(\ker(f)\) 有限且
\[
\#\ker(f)=n.
\]
而紧交换群之间的连续满同态若核有限，则为有限度覆盖，其度数等于核的基数，故
\[
\deg(f)=\#\ker(f)=n.
\]

反之，对任意与 \(S\) 互素的正整数 \(n\)，乘法映射
\[
\times n:G_S\to G_S
\]
是单射，其对偶给出一个核同构于 \(\ZZ/n\ZZ\) 的连续自覆盖，因此上述度数集合恰好被全部 \(S\)-自由正整数实现。

最后，\(f\) 为同构当且仅当 \(\ker(f)\) 平凡，当且仅当 \(n=1\)。这等价于 \(r\in G_S^\times\)，故
\[
\operatorname{Aut}_{\mathrm{cts}}(\widehat{G_S})\cong G_S^\times
=
\left\{\pm\prod_{p\in S}p^{m_p}:m_p\in\ZZ\right\},
\]
与推论 \ref{cor:xi-localized-integers-dual-automorphism-units} 一致。证毕。
\end{proof}

\noindent
这些结果说明：Pisot--Bernoulli 卷积在任意非零 Fibonacci 递推方向上的 Fourier 采样都保留严格正的共振极限，因此非 Rajchman 障碍并非少数特殊子列现象；与之并行，Fibonacci 共振整函数虽由几何放大的半整数格层叠构成，其缩放零点计数测度却在宏观上趋于 \([-1,1]\) 上的均匀律；而有限素数局部化 solenoid 的全部有限度自覆盖，则完全由 \(S\)-自由整数度数控制，\(S\)-内素数只保留为零代价的单位自同构。

\endinput


\endinput


\endinput


\endinput
